\documentclass[11pt,a4paper]{report}
\begin{document}
\begin{titlepage}
\centering
{\Huge Dominium Conversation Corpus Book v0\par}
\vspace{1cm}
{\Large A Derived Reader, Synthesis, Decision, Promotion, and Reconciliation Guide\par}
\vspace{1cm}
{\large Status: DERIVED \\ Authority Class: advisory\_synthesis \\ Promotion Status: not\_promoted\par}
\vfill
{\large 2026-05-28\par}
\end{titlepage}
\tableofcontents
\clearpage
\chapter{Front Matter}

\chapter{Dominium Conversation Corpus Book v0}

\section{A Derived Reader, Synthesis, Decision, Promotion, and Reconciliation Guide}

Status: DERIVED

Last Reviewed: 2026-05-28

Supersedes: none

Superseded By: none

Stability: provisional

Authority Class: advisory\_synthesis

Source Root: docs/archive/conversations/

Promotion Status: not\_promoted

Canon impact: unchanged

Contract/schema impact: unchanged

Implementation impact: unchanged

Release impact: unchanged

Queue impact: unchanged



\section{Authority and Non-Promotion Notice}

This book is a publication artifact derived from docs/archive/conversations/.

It is not canon, not architecture doctrine, not a contract, not a schema, and

not an implementation plan.


The book does not promote archive claims into repo truth. Current repo authority

remains docs/canon/constitution\_v1.md, docs/canon/glossary\_v1.md,

AGENTS.md, current contracts/schema law, current queue state, and validated

repo artifacts.


\section{How To Read This Book}

\begin{itemize}
\item Read Part I for orientation and current guardrails.
\item Read Part II for the narrative synthesis.
\item Read Part III for authority and reconciliation.
\item Read Part IV for decisions that need human or future-queue handling.
\item Read Part V for promotion candidates that may become later microtasks.
\item Read Part VI for contradiction, staleness, and verification risk.
\item Use the appendices for dense source and register material.
\end{itemize}

\section{What Not To Conclude}

\begin{itemize}
\item This book does not open blocked queue areas.
\item This book does not resolve contradictions by itself.
\item This book does not replace current repo docs.
\item This book does not claim release readiness.
\item This book does not authorize canon, contract, schema, implementation, release,
\end{itemize}
  or queue changes.


\section{Source Hierarchy}

1. Canon and glossary.

2. AGENTS.md.

3. Scope-specific contracts, schema law, current queue, and validated repo

   artifacts.

4. Generated archive outputs with provenance.

5. Historical conversation material as advisory evidence.


\section{ID Conventions}

\begin{itemize}
\item DECISION-*: decision docket items.
\item PROMOTE-*: raw promotion candidates.
\item CONTRA-*: contradiction, staleness, or drift findings.
\item Source paths are repo-relative unless explicitly marked otherwise.
\end{itemize}


\chapter{Orientation}


{\small\raggedright SOURCE PATH: docs/\allowbreak{}archive/\allowbreak{}conversations/\allowbreak{}\_\allowbreak{}synthesis/\allowbreak{}READ\_\allowbreak{}THIS\_\allowbreak{}FIRST\_\allowbreak{}v0.md\par}


\section{Read This First v0}

Use this page as the shortest path through the conversation-corpus adjudication surfaces.


\subsection{If You Want Current Truth}

\begin{itemize}
\item Start with [README.md](../../../../README.md).
\item Then read [constitution\_v1.md](../../../canon/constitution\_v1.md), [glossary\_v1.md](../../../canon/glossary\_v1.md), [AGENTS.md](../../../../AGENTS.md), and [.aide/queue/current.toml](../../../../.aide/queue/current.toml).
\end{itemize}

\subsection{If You Want Archive Context}

\begin{itemize}
\item Read [PROJECT\_SYNTHESIS\_BOOK\_v0.md](PROJECT\_SYNTHESIS\_BOOK\_v0.md).
\item Then use [conversation\_reader\_index.md](../\_reader/conversation\_reader\_index.md) for individual conversations.
\end{itemize}

\subsection{If You Want Decisions}

\begin{itemize}
\item Read [DECISION\_DOCKET\_v0.md](../\_decision/DECISION\_DOCKET\_v0.md): 30 decision items.
\item Use [DEFERRED\_DECISIONS\_v0.md](../\_decision/DEFERRED\_DECISIONS\_v0.md) to avoid accidentally opening blocked scope.
\end{itemize}

\subsection{If You Want Promotion Candidates}

\begin{itemize}
\item Read [PROMOTION\_REVIEW\_BOARD\_v0.md](../\_promotion/PROMOTION\_REVIEW\_BOARD\_v0.md): 135 candidates.
\item Start with [PROMOTION\_WAVE\_1\_CANDIDATES\_v0.md](../\_promotion/PROMOTION\_WAVE\_1\_CANDIDATES\_v0.md): 18 docs-only review candidates.
\end{itemize}

\subsection{Hard Guardrails}

\begin{itemize}
\item Do not promote archive claims without a later reviewed task.
\item Do not patch canon, contracts, schema, implementation, release, or queue files from this atlas.
\item Do not open blocked queue scopes from old conversations.
\end{itemize}



{\small\raggedright SOURCE PATH: docs/\allowbreak{}archive/\allowbreak{}conversations/\allowbreak{}\_\allowbreak{}synthesis/\allowbreak{}CURRENT\_\allowbreak{}PROJECT\_\allowbreak{}ATLAS\_\allowbreak{}v0.md\par}


\section{Current Project Atlas v0}

This atlas is a reader map. It is not canon and does not promote archived conversation claims.


\subsection{Read Order}

1. README.md (../../../../README.md) for current project orientation.

2. constitution\_v1.md (../../../canon/constitution\_v1.md) and glossary\_v1.md (../../../canon/glossary\_v1.md) for binding meaning.

3. AGENTS.md (../../../../AGENTS.md) and AUTHORITY\_ORDER.md (../../../planning/AUTHORITY\_ORDER.md) for repo/agent authority rules.

4. .aide/queue/current.toml (../../../../.aide/queue/current.toml) for current scope gates.

5. DECISION\_DOCKET\_v0.md (../\_decision/DECISION\_DOCKET\_v0.md) and PROMOTION\_REVIEW\_BOARD\_v0.md (../\_promotion/PROMOTION\_REVIEW\_BOARD\_v0.md) for archive-derived adjudication.


\subsection{What Dominium Is}

Current repo truth: Dominium is a deterministic, contract-governed simulation game and operating environment built on the Domino deterministic substrate. It is about lawful simulation of invention, production, logistics, economics, settlement, trust, communication, and institutional power.


Conversation-derived context: the archive expands that picture into a long-horizon product ecosystem: engine substrate, official game/domain layer, Workbench, setup/launcher, content packs, portability, release identity, and repo-governed assistant workflows. That context is historical and advisory.


\subsection{What Domino Is}

Domino is the reusable deterministic substrate: execution ordering, storage, replay, ABI-facing boundaries, and engine mechanisms. Archived conversations consistently treat Domino as reusable infrastructure rather than a one-off game binary.


\subsection{Active Surfaces}

\begin{quote}\small Dense table content is summarized in this PDF rendering. See the Markdown source and HTML book for the full table.\end{quote}

\subsection{Current Work Gates}

\begin{itemize}
\item Foundation lock status signal: PASS\_WITH\_WARNINGS.
\item Current blocked queue scopes: broad\_workbench\_ui, gameplay, native\_gui, package\_runtime, provider\_runtime, release\_publication, renderer\_implementation, runtime\_module\_loader.
\item Safe now: archive reading, decision review, promotion review, reconciliation crosswalks, docs-only microtask preparation, and validators.
\item Not safe now: implementation work in blocked scopes, canon/schema/contract rewrites from archive claims, release publication, renderer implementation, gameplay, provider/package runtime, broad Workbench UI, native GUI.
\end{itemize}

\subsection{Unresolved Decisions}

\begin{itemize}
\item Decision docket items: 30.
\item Workbench/AIDE/Codex/tooling: 6.
\item architecture/boundaries: 12.
\item provider/content/packs: 5.
\item release/setup/launcher: 2.
\item renderer/UI/platform: 3.
\item world/time/civilization simulation: 2.
\end{itemize}

\subsection{Promotion State}

\begin{itemize}
\item Raw promotion candidates: 135.
\item Wave 1 docs-only candidates: 18.
\item architecture\_doc\_candidate: 16.
\item archive\_only: 9.
\item blocked\_by\_queue: 10.
\item contract\_candidate: 11.
\item docs\_clarification\_candidate: 1.
\item insufficient\_evidence: 53.
\item needs\_user\_decision: 12.
\item planning\_doc\_candidate: 1.
\item reject\_as\_noise: 19.
\item schema\_candidate: 3.
\end{itemize}

\subsection{What Requires Review Before Promotion}

\begin{itemize}
\item Any claim touching canon, glossary, authority order, contracts, schema law, current queue, release, implementation, or blocked scope.
\item Any old language/platform/baseline claim.
\item Any claim whose target path is only a wildcard such as docs/architecture/** or contracts/**.
\item Any Wave 1 candidate before source text and target docs are inspected together.
\end{itemize}

\subsection{Recommended Next Steps}

1. Manually review the decision docket and choose defer/preserve/promote-later dispositions.

2. Select a small subset of Wave 1 candidates for docs-only microtasks.

3. For each microtask, name the source claim ID, exact target doc, authority support, non-goals, validation, and promotion status update.

4. Keep implementation blocked unless the current queue opens it.



\chapter{Project Synthesis}


{\small\raggedright SOURCE PATH: docs/\allowbreak{}archive/\allowbreak{}conversations/\allowbreak{}\_\allowbreak{}synthesis/\allowbreak{}EXECUTIVE\_\allowbreak{}BRIEF\_\allowbreak{}v0.md\par}


\section{Executive Brief v0}

The conversation archive now supports a derived project picture: Dominium is a deterministic, contract-governed simulation game and operating environment on top of Domino. The old conversations consistently point toward a broad product: simulation engine, official game/domain layer, Workbench, setup/launcher, content packs, release identity, portability, and repo governance.


This brief is not authority. Current truth remains in README.md (../../../../README.md), constitution\_v1.md (../../../canon/constitution\_v1.md), glossary\_v1.md (../../../canon/glossary\_v1.md), AGENTS.md (../../../../AGENTS.md), contracts, schema law, and current queue state.


\subsection{Current Signal}

\begin{itemize}
\item Source conversations: 45
\item Source files: 604
\item Acceptance result: PASS\_WITH\_WARNINGS
\item Promotion candidates: 135
\item Audit findings: 227
\item Blocked current queue scopes: broad\_workbench\_ui, gameplay, native\_gui, package\_runtime, provider\_runtime, release\_publication, renderer\_implementation, runtime\_module\_loader
\end{itemize}

\subsection{Main Takeaway}

The archive is ready for understanding and synthesis. It is not ready for direct promotion. The next review surface should compare the synthesis against current repo authority and reduce the raw promotion queue into a smaller decision backlog.




{\small\raggedright SOURCE PATH: docs/\allowbreak{}archive/\allowbreak{}conversations/\allowbreak{}\_\allowbreak{}synthesis/\allowbreak{}PROJECT\_\allowbreak{}SYNTHESIS\_\allowbreak{}BOOK\_\allowbreak{}v0.md\par}


\section{Dominium Conversation Corpus Synthesis Book v0}

This book is a derived reading guide over the archived conversation corpus. It is not canon, not architecture doctrine, not a contract, and not an implementation plan.


\subsection{1. Status And Authority}

\begin{itemize}
\item Acceptance review result: PASS\_WITH\_WARNINGS.
\item Authority class: advisory\_synthesis.
\item Promotion status: not\_promoted.
\item Current repo artifacts outrank every conversation-derived statement in this book.
\end{itemize}

\subsection{2. Executive Overview}

The archive describes Dominium as a long-horizon deterministic simulation product: a game, engine-backed operating environment, validation workbench, content platform, and governance-heavy repository. Across the conversations, the recurring intent is not to build a renderer-owned game loop or a pile of UI screens. The recurring intent is to preserve lawful simulation truth, expose it through command/result/evidence surfaces, and let products project that truth without mutating it.


The current repo already establishes the strongest version of that principle through canon, glossary, README, and queue state. The conversation corpus adds historical design intent around world scale, Workbench, setup/launcher, release identity, provider/content boundaries, and future UI/editor experiences. Those claims remain advisory until reconciled and promoted.


\subsection{3. What Dominium Is Trying To Be}

Repo-established truth: Dominium is the official game/product/domain layer on top of Domino, the reusable deterministic substrate. It is concerned with invention, production, logistics, economics, settlement, trust, communication, and institutional power emerging from lawful simulation.


Conversation-derived intent: the archive repeatedly extends that picture toward a broad universe-scale simulation platform, with real-world defaults, arbitrary authored packs, robust tooling, deterministic evidence, long-term portability, and a Workbench that makes the project inspectable and operable.


\subsection{4. Core Mental Model}

The current README gives the clearest spine: intent -> command -> capability/refusal check -> service or deterministic process -> result/document/snapshot -> diagnostics/evidence -> view/action model -> projection -> shell. The conversations mostly orbit that same model, even when they discuss UI, renderer, platform, setup, or release packaging.


\subsection{5. Engine / Game / Runtime / Product Boundaries}

Repo truth keeps Engine, Game, Client, Server, Workbench, setup, launcher, tools, contracts, and content in separate roles. Conversation-derived material reinforces the boundary: renderer and UI project state, clients issue intents, server/game/engine law remains authoritative, and tools validate or inspect rather than silently mutate truth.


Representative sources: Dominium Advanced Simulation and Infrastructure Architecture (../\_reader/by\_chat/advanced\_simulation\_infrastructure.md), Dominium APP0 Runtime, Platform, and Renderer Architecture (../\_reader/by\_chat/app\_runtime\_platform\_renderers.md), Dominium Architecture, Application Layer, TestX, and CodeHygiene Planning (../\_reader/by\_chat/app\_testx\_codehygiene.md), Dominium/Domino Architecture and Codex Prompt Roadmap (../\_reader/by\_chat/architecture\_codex\_prompts.md), Dominium Architecture, UI, Providers, and Robot OS Strategy (../\_reader/by\_chat/architecture\_ui\_providers.md), Dominium Build and Future-Proofing Architecture (../\_reader/by\_chat/build\_and\_future\_proofing.md), Dominium Canonical Architecture, Repository Foundation, and Provider Model (../\_reader/by\_chat/canonical\_structure\_and\_framework.md), Dominium Chronology \& Celestial Systems (../\_reader/by\_chat/chronology\_celestial\_systems.md).


\subsection{6. Determinism, Replay, And Provenance}

The canon makes determinism primary: ordering, reductions, RNG streams, replay, hash partitions, and provenance are non-negotiable. The corpus adds historical emphasis on deterministic world scale, portable builds, replay proofs, verification artifacts, and avoiding hidden fallback behavior.


Representative sources: Dominium Advanced Simulation and Infrastructure Architecture (../\_reader/by\_chat/advanced\_simulation\_infrastructure.md), Dominium APP0 Runtime, Platform, and Renderer Architecture (../\_reader/by\_chat/app\_runtime\_platform\_renderers.md), Dominium Architecture, Application Layer, TestX, and CodeHygiene Planning (../\_reader/by\_chat/app\_testx\_codehygiene.md), Dominium/Domino Architecture and Codex Prompt Roadmap (../\_reader/by\_chat/architecture\_codex\_prompts.md), Dominium Architecture, UI, Providers, and Robot OS Strategy (../\_reader/by\_chat/architecture\_ui\_providers.md), Dominium Build and Future-Proofing Architecture (../\_reader/by\_chat/build\_and\_future\_proofing.md), Dominium Canonical Architecture, Repository Foundation, and Provider Model (../\_reader/by\_chat/canonical\_structure\_and\_framework.md), Dominium Development Routes and Continuity Preservation (../\_reader/by\_chat/development\_routes.md).


\subsection{7. Law, Authority, Refusal, And Capability Model}

Current repo truth says authority permits attempts, law decides accept/refuse/transform, and refusals must be deterministic and auditable. Conversations often use different local language, but the recurring direction is aligned: no convenience bypass, no generated output as truth, and no old prompt as authority.


Representative sources: Dominium Advanced Simulation and Infrastructure Architecture (../\_reader/by\_chat/advanced\_simulation\_infrastructure.md), Dominium APP0 Runtime, Platform, and Renderer Architecture (../\_reader/by\_chat/app\_runtime\_platform\_renderers.md), Dominium Architecture, Application Layer, TestX, and CodeHygiene Planning (../\_reader/by\_chat/app\_testx\_codehygiene.md), Dominium/Domino Architecture and Codex Prompt Roadmap (../\_reader/by\_chat/architecture\_codex\_prompts.md), Dominium Architecture, UI, Providers, and Robot OS Strategy (../\_reader/by\_chat/architecture\_ui\_providers.md), Dominium Build and Future-Proofing Architecture (../\_reader/by\_chat/build\_and\_future\_proofing.md), Dominium Canonical Architecture, Repository Foundation, and Provider Model (../\_reader/by\_chat/canonical\_structure\_and\_framework.md), Dominium Chronology \& Celestial Systems (../\_reader/by\_chat/chronology\_celestial\_systems.md).


\subsection{8. World, Reality, Scale, And Refinement Model}

Conversation-derived intent is expansive: worlds, planets, celestial systems, universe explorer concepts, chronology, spatial refinement, simulation domains, macro/micro transitions, and real-world defaults recur across the archive. This is a design-intent signal, not current implementation authority.


Representative sources: Dominium APP0 Runtime, Platform, and Renderer Architecture (../\_reader/by\_chat/app\_runtime\_platform\_renderers.md), Dominium Architecture, Application Layer, TestX, and CodeHygiene Planning (../\_reader/by\_chat/app\_testx\_codehygiene.md), Dominium/Domino Architecture and Codex Prompt Roadmap (../\_reader/by\_chat/architecture\_codex\_prompts.md), Dominium Architecture, UI, Providers, and Robot OS Strategy (../\_reader/by\_chat/architecture\_ui\_providers.md), Dominium Build and Future-Proofing Architecture (../\_reader/by\_chat/build\_and\_future\_proofing.md), Dominium Chronology \& Celestial Systems (../\_reader/by\_chat/chronology\_celestial\_systems.md), Dominium Development Routes and Continuity Preservation (../\_reader/by\_chat/development\_routes.md), Dominium Architecture I (../\_reader/by\_chat/dominium\_architecture\_i.md).


\subsection{9. Timekeeping, Chronology, And 2038 Resilience}

Time appears as both a simulation domain and a platform durability concern. The archive repeatedly points to chronology, calendars, celestial alignment, 2038 resilience, and timestamp policy. Current promotion must verify all such claims against current contracts and language/platform baselines.


Representative sources: Dominium Advanced Simulation and Infrastructure Architecture (../\_reader/by\_chat/advanced\_simulation\_infrastructure.md), Dominium APP0 Runtime, Platform, and Renderer Architecture (../\_reader/by\_chat/app\_runtime\_platform\_renderers.md), Dominium Architecture, Application Layer, TestX, and CodeHygiene Planning (../\_reader/by\_chat/app\_testx\_codehygiene.md), Dominium/Domino Architecture and Codex Prompt Roadmap (../\_reader/by\_chat/architecture\_codex\_prompts.md), Dominium Canonical Architecture, Repository Foundation, and Provider Model (../\_reader/by\_chat/canonical\_structure\_and\_framework.md), Dominium Chronology \& Celestial Systems (../\_reader/by\_chat/chronology\_celestial\_systems.md), Dominium Development Routes and Continuity Preservation (../\_reader/by\_chat/development\_routes.md), Documentation Standards, README Strategy, and Handoff Packaging (../\_reader/by\_chat/documentation\_standards\_readme.md).


\subsection{10. Civilization, Institutions, Economy, Logistics}

The corpus frames Dominium around emergent production, logistics, economics, settlement, communication, trust, institutions, and power. These are project identity themes; they are not authorization to open gameplay or runtime work while the current queue blocks it.


Representative sources: Dominium Advanced Simulation and Infrastructure Architecture (../\_reader/by\_chat/advanced\_simulation\_infrastructure.md), Dominium APP0 Runtime, Platform, and Renderer Architecture (../\_reader/by\_chat/app\_runtime\_platform\_renderers.md), Dominium Architecture, Application Layer, TestX, and CodeHygiene Planning (../\_reader/by\_chat/app\_testx\_codehygiene.md), Dominium/Domino Architecture and Codex Prompt Roadmap (../\_reader/by\_chat/architecture\_codex\_prompts.md), Dominium Architecture, UI, Providers, and Robot OS Strategy (../\_reader/by\_chat/architecture\_ui\_providers.md), Dominium Build and Future-Proofing Architecture (../\_reader/by\_chat/build\_and\_future\_proofing.md), Dominium Canonical Architecture, Repository Foundation, and Provider Model (../\_reader/by\_chat/canonical\_structure\_and\_framework.md), Dominium Chronology \& Celestial Systems (../\_reader/by\_chat/chronology\_celestial\_systems.md).


\subsection{11. UI, Renderer, Workbench, And Tools}

The conversations strongly prefer a rich Workbench and eventual editor/operator surfaces. Current authority is narrower: broad Workbench UI, renderer implementation, and native GUI remain blocked. Therefore synthesis may describe intent, but follow-up tasks must stay contract/planning scoped until the queue opens implementation scope.


Representative sources: Dominium Advanced Simulation and Infrastructure Architecture (../\_reader/by\_chat/advanced\_simulation\_infrastructure.md), Dominium Architecture, UI, Providers, and Robot OS Strategy (../\_reader/by\_chat/architecture\_ui\_providers.md), Dominium Canonical Architecture, Repository Foundation, and Provider Model (../\_reader/by\_chat/canonical\_structure\_and\_framework.md), Dominium Development Routes and Continuity Preservation (../\_reader/by\_chat/development\_routes.md), Dominium Architecture IV (../\_reader/by\_chat/dominium\_architecture\_iv.md), Dominium + Domino Codex Planning and Handoff (../\_reader/by\_chat/dominium\_domino\_codex\_planning.md), Dominium Workbench, AIDE, Presentation Architecture, Provider Strategy, Product-Spine Planning, and Preservation Companion (../\_reader/by\_chat/dominium\_full\_conversation.md), Domino Dominium Workbench (../\_reader/by\_chat/domino\_dominium\_workbench.md).


\subsection{12. AIDE, Codex, Governance, And Patch Workflow}

The archive treats AIDE/Codex/XStack as repo-control and patch-execution aids. Current repo truth agrees that they do not replace canon, contracts, review gates, or validation. Their correct role is bounded evidence-producing work, not direct semantic authority.


Representative sources: Dominium Architecture, Application Layer, TestX, and CodeHygiene Planning (../\_reader/by\_chat/app\_testx\_codehygiene.md), Dominium Canonical Architecture, Repository Foundation, and Provider Model (../\_reader/by\_chat/canonical\_structure\_and\_framework.md), Dominium Canon, Repository Alignment, and Portability Doctrine (../\_reader/by\_chat/dominium\_complete\_conversation.md), Dominium Workbench, AIDE, Presentation Architecture, Provider Strategy, Product-Spine Planning, and Preservation Companion (../\_reader/by\_chat/dominium\_full\_conversation.md), Dominium Foundation Lock, Workbench Spine, and Parallel Codex Handoff (../\_reader/by\_chat/foundation\_workbench\_codex.md), Dominium Language, Platform, and Architecture Baseline (../\_reader/by\_chat/language\_platform\_architecture.md), Dominium Modularity, AIDE Refactorability, and Future-Proof Architecture (../\_reader/by\_chat/modularity\_aide\_refactorability.md), Dominium Omega/Xi/Pi Architecture \& Future-Proofing Planning (../\_reader/by\_chat/omega\_xi\_pi\_architecture\_future.md).


\subsection{13. Release, Setup, Launcher, Platform, Portability}

Conversations cover setup, launcher, release identity, versioning, portability, platform support, and public distribution. Current queue still blocks release publication, so these remain planning and audit material unless a later reviewed task opens the scope.


Representative sources: Dominium Advanced Simulation and Infrastructure Architecture (../\_reader/by\_chat/advanced\_simulation\_infrastructure.md), Dominium APP0 Runtime, Platform, and Renderer Architecture (../\_reader/by\_chat/app\_runtime\_platform\_renderers.md), Dominium Architecture, Application Layer, TestX, and CodeHygiene Planning (../\_reader/by\_chat/app\_testx\_codehygiene.md), Dominium/Domino Architecture and Codex Prompt Roadmap (../\_reader/by\_chat/architecture\_codex\_prompts.md), Dominium Architecture, UI, Providers, and Robot OS Strategy (../\_reader/by\_chat/architecture\_ui\_providers.md), Dominium Build and Future-Proofing Architecture (../\_reader/by\_chat/build\_and\_future\_proofing.md), Dominium Canonical Architecture, Repository Foundation, and Provider Model (../\_reader/by\_chat/canonical\_structure\_and\_framework.md), Dominium Development Routes and Continuity Preservation (../\_reader/by\_chat/development\_routes.md).


\subsection{14. Content, Packs, Modding, Providers}

The corpus repeatedly returns to pack-driven composition, authored content, providers, binary/content boundaries, and modding. Canon requires pack-driven integration and explicit refusal for missing packs. Current queue blocks package runtime and provider runtime, so promotion must be careful.


Representative sources: Dominium Advanced Simulation and Infrastructure Architecture (../\_reader/by\_chat/advanced\_simulation\_infrastructure.md), Dominium APP0 Runtime, Platform, and Renderer Architecture (../\_reader/by\_chat/app\_runtime\_platform\_renderers.md), Dominium Architecture, Application Layer, TestX, and CodeHygiene Planning (../\_reader/by\_chat/app\_testx\_codehygiene.md), Dominium/Domino Architecture and Codex Prompt Roadmap (../\_reader/by\_chat/architecture\_codex\_prompts.md), Dominium Architecture, UI, Providers, and Robot OS Strategy (../\_reader/by\_chat/architecture\_ui\_providers.md), Dominium Build and Future-Proofing Architecture (../\_reader/by\_chat/build\_and\_future\_proofing.md), Dominium Canonical Architecture, Repository Foundation, and Provider Model (../\_reader/by\_chat/canonical\_structure\_and\_framework.md), Dominium Chronology \& Celestial Systems (../\_reader/by\_chat/chronology\_celestial\_systems.md).


\subsection{15. What Was Decided Across Conversations}

The strongest recurring conversation-level decisions are advisory: keep simulation truth deterministic, separate rendering from authority, keep Workbench a projection/inspection/control surface rather than truth owner, preserve source provenance, and avoid direct promotion from old chats. These are valuable because they align with current repo doctrine, but alignment is not the same as promotion.


\subsection{16. What Is Still Unresolved}

\begin{itemize}
\item Promotion candidate triage: The raw promotion queue contains useful leads but also preservation-process noise. Source: [../\_promotion/PROMOTION\_QUEUE.md](../\_promotion/PROMOTION\_QUEUE.md).
\item Workbench scope: Old conversations discuss Workbench as rich operator surface, but the current queue blocks broad Workbench UI. Source: [../\_wiki/topics/workbench.md](../\_wiki/topics/workbench.md).
\item Renderer and platform boundary: Conversations repeatedly discuss renderer/platform plans, while current authority keeps rendering presentational and implementation scope blocked. Source: [../\_wiki/topics/platform.md](../\_wiki/topics/platform.md).
\item Provider and content boundaries: Provider, pack, and content discussions need separation from runtime module loading and package runtime blocks. Source: [../\_wiki/topics/content.md](../\_wiki/topics/content.md).
\item World scale and fidelity roadmap: Universe, world, chronology, and simulation conversations need sequencing against current contracts and queue limits. Source: [../\_wiki/topics/worldgen.md](../\_wiki/topics/worldgen.md).
\item Release/publication meaning: Release/versioning conversations should remain planning until release publication authority opens. Source: [../\_wiki/topics/release.md](../\_wiki/topics/release.md).
\end{itemize}

\subsection{17. Contradictions And Drift}

The audit layer contains review triggers, not resolved contradictions. The main risk classes are conversation claims against current queue restrictions, stale external/platform claims, old docs/baseline drift, and conversation-vs-conversation disagreements.


\begin{itemize}
\item conversation\_vs\_conversation: 2 findings.
\item conversation\_vs\_current\_queue: 102 findings.
\item conversation\_vs\_docs: 2 findings.
\item stale\_external\_claim: 121 findings.
\end{itemize}

\subsection{18. Stale Claims And Verification Queue}

External platform, SDK, renderer, language baseline, release, provider, package runtime, and implementation claims must be rechecked. The synthesis should use those old claims to identify questions, not to assert current facts.


\subsection{19. Promotion Candidates}

The raw queue contains 135 generated candidates in PROMOTION\_QUEUE.md (../\_promotion/PROMOTION\_QUEUE.md). Acceptance review found the queue useful but noisy. It should be triaged before any candidate patch work.


\subsection{20. Recommended Reconciliation Roadmap}

1. Keep this synthesis advisory and archive-scoped.

2. Build a claim review matrix against canon, glossary, AGENTS, current queue, contracts, schema law, and targeted current docs.

3. Triage promotion candidates into serious, historical, stale, noisy, rejected, and needs-user-decision groups.

4. Patch live docs only through narrow promotion tasks with named source claims and validation.


\subsection{21. Appendices}

Current repo truth, as used by this synthesis:


\begin{itemize}
\item [README.md](../../../../README.md) describes Dominium as a deterministic, contract-governed simulation game and operating environment built on the Domino deterministic substrate.
\item [docs/canon/constitution\_v1.md](../../../canon/constitution\_v1.md) binds determinism, process-only mutation, law-gated authority, no runtime mode flags, pack-driven integration, explicit refusal, and truth/perceived/render separation.
\item [docs/canon/glossary\_v1.md](../../../canon/glossary\_v1.md) defines vocabulary such as Engine, Client, AuthorityContext, Domain, Derived artifact, and Contract.
\item [AGENTS.md](../../../../AGENTS.md) and [docs/planning/AUTHORITY\_ORDER.md](../../../planning/AUTHORITY\_ORDER.md) keep archived conversations below canon, glossary, governance, contracts, current queue, and validated repo artifacts.
\item [.aide/queue/current.toml](../../../../.aide/queue/current.toml) currently blocks broad feature work including: broad\_workbench\_ui, gameplay, native\_gui, package\_runtime, provider\_runtime, release\_publication, renderer\_implementation, runtime\_module\_loader.
\end{itemize}


\subsubsection{Topic Coverage}

\begin{itemize}
\item architecture (45 conversations): system boundaries, product shape, engine/game/client/server separation, and repository structure. Sources: [Dominium Advanced Simulation and Infrastructure Architecture](../\_reader/by\_chat/advanced\_simulation\_infrastructure.md), [Dominium APP0 Runtime, Platform, and Renderer Architecture](../\_reader/by\_chat/app\_runtime\_platform\_renderers.md), [Dominium Architecture, Application Layer, TestX, and CodeHygiene Planning](../\_reader/by\_chat/app\_testx\_codehygiene.md), [Dominium/Domino Architecture and Codex Prompt Roadmap](../\_reader/by\_chat/architecture\_codex\_prompts.md), [Dominium Architecture, UI, Providers, and Robot OS Strategy](../\_reader/by\_chat/architecture\_ui\_providers.md).
\item content (45 conversations): packs, mods, authored payload, GUI/content boundaries, and provider/content separation. Sources: [Dominium Advanced Simulation and Infrastructure Architecture](../\_reader/by\_chat/advanced\_simulation\_infrastructure.md), [Dominium APP0 Runtime, Platform, and Renderer Architecture](../\_reader/by\_chat/app\_runtime\_platform\_renderers.md), [Dominium Architecture, Application Layer, TestX, and CodeHygiene Planning](../\_reader/by\_chat/app\_testx\_codehygiene.md), [Dominium/Domino Architecture and Codex Prompt Roadmap](../\_reader/by\_chat/architecture\_codex\_prompts.md), [Dominium Architecture, UI, Providers, and Robot OS Strategy](../\_reader/by\_chat/architecture\_ui\_providers.md).
\item ui (45 conversations): presentation, tools, editor concepts, perceived model surfaces, and command/result display. Sources: [Dominium Advanced Simulation and Infrastructure Architecture](../\_reader/by\_chat/advanced\_simulation\_infrastructure.md), [Dominium APP0 Runtime, Platform, and Renderer Architecture](../\_reader/by\_chat/app\_runtime\_platform\_renderers.md), [Dominium Architecture, Application Layer, TestX, and CodeHygiene Planning](../\_reader/by\_chat/app\_testx\_codehygiene.md), [Dominium/Domino Architecture and Codex Prompt Roadmap](../\_reader/by\_chat/architecture\_codex\_prompts.md), [Dominium Architecture, UI, Providers, and Robot OS Strategy](../\_reader/by\_chat/architecture\_ui\_providers.md).
\item governance (44 conversations): authority order, canon, contracts, refusal, review gates, and agent operation. Sources: [Dominium Advanced Simulation and Infrastructure Architecture](../\_reader/by\_chat/advanced\_simulation\_infrastructure.md), [Dominium APP0 Runtime, Platform, and Renderer Architecture](../\_reader/by\_chat/app\_runtime\_platform\_renderers.md), [Dominium Architecture, Application Layer, TestX, and CodeHygiene Planning](../\_reader/by\_chat/app\_testx\_codehygiene.md), [Dominium/Domino Architecture and Codex Prompt Roadmap](../\_reader/by\_chat/architecture\_codex\_prompts.md), [Dominium Architecture, UI, Providers, and Robot OS Strategy](../\_reader/by\_chat/architecture\_ui\_providers.md).
\item simulation (44 conversations): deterministic domains, processes, machines, physical systems, economy, and civilization dynamics. Sources: [Dominium Advanced Simulation and Infrastructure Architecture](../\_reader/by\_chat/advanced\_simulation\_infrastructure.md), [Dominium APP0 Runtime, Platform, and Renderer Architecture](../\_reader/by\_chat/app\_runtime\_platform\_renderers.md), [Dominium Architecture, Application Layer, TestX, and CodeHygiene Planning](../\_reader/by\_chat/app\_testx\_codehygiene.md), [Dominium/Domino Architecture and Codex Prompt Roadmap](../\_reader/by\_chat/architecture\_codex\_prompts.md), [Dominium Architecture, UI, Providers, and Robot OS Strategy](../\_reader/by\_chat/architecture\_ui\_providers.md).
\item tooling (44 conversations): Codex/AIDE/XStack/TestX/RepoX/AuditX tooling and patch workflow. Sources: [Dominium Advanced Simulation and Infrastructure Architecture](../\_reader/by\_chat/advanced\_simulation\_infrastructure.md), [Dominium APP0 Runtime, Platform, and Renderer Architecture](../\_reader/by\_chat/app\_runtime\_platform\_renderers.md), [Dominium Architecture, Application Layer, TestX, and CodeHygiene Planning](../\_reader/by\_chat/app\_testx\_codehygiene.md), [Dominium/Domino Architecture and Codex Prompt Roadmap](../\_reader/by\_chat/architecture\_codex\_prompts.md), [Dominium Architecture, UI, Providers, and Robot OS Strategy](../\_reader/by\_chat/architecture\_ui\_providers.md).
\item release (43 conversations): version identity, packaging, setup, update, publication, and release-control concerns. Sources: [Dominium Advanced Simulation and Infrastructure Architecture](../\_reader/by\_chat/advanced\_simulation\_infrastructure.md), [Dominium APP0 Runtime, Platform, and Renderer Architecture](../\_reader/by\_chat/app\_runtime\_platform\_renderers.md), [Dominium Architecture, Application Layer, TestX, and CodeHygiene Planning](../\_reader/by\_chat/app\_testx\_codehygiene.md), [Dominium/Domino Architecture and Codex Prompt Roadmap](../\_reader/by\_chat/architecture\_codex\_prompts.md), [Dominium Architecture, UI, Providers, and Robot OS Strategy](../\_reader/by\_chat/architecture\_ui\_providers.md).
\item contracts\_schema (42 conversations): machine-readable contracts, schemas, compatibility, manifests, and registries. Sources: [Dominium Advanced Simulation and Infrastructure Architecture](../\_reader/by\_chat/advanced\_simulation\_infrastructure.md), [Dominium APP0 Runtime, Platform, and Renderer Architecture](../\_reader/by\_chat/app\_runtime\_platform\_renderers.md), [Dominium Architecture, Application Layer, TestX, and CodeHygiene Planning](../\_reader/by\_chat/app\_testx\_codehygiene.md), [Dominium/Domino Architecture and Codex Prompt Roadmap](../\_reader/by\_chat/architecture\_codex\_prompts.md), [Dominium Architecture, UI, Providers, and Robot OS Strategy](../\_reader/by\_chat/architecture\_ui\_providers.md).
\item determinism (39 conversations): replay, ordering, proof, fixed identity, RNG discipline, and thread-count invariance. Sources: [Dominium Advanced Simulation and Infrastructure Architecture](../\_reader/by\_chat/advanced\_simulation\_infrastructure.md), [Dominium APP0 Runtime, Platform, and Renderer Architecture](../\_reader/by\_chat/app\_runtime\_platform\_renderers.md), [Dominium Architecture, Application Layer, TestX, and CodeHygiene Planning](../\_reader/by\_chat/app\_testx\_codehygiene.md), [Dominium/Domino Architecture and Codex Prompt Roadmap](../\_reader/by\_chat/architecture\_codex\_prompts.md), [Dominium Architecture, UI, Providers, and Robot OS Strategy](../\_reader/by\_chat/architecture\_ui\_providers.md).
\item platform (39 conversations): runtime shell, platform adapter, renderer boundary, operating-system support, and portability. Sources: [Dominium Advanced Simulation and Infrastructure Architecture](../\_reader/by\_chat/advanced\_simulation\_infrastructure.md), [Dominium APP0 Runtime, Platform, and Renderer Architecture](../\_reader/by\_chat/app\_runtime\_platform\_renderers.md), [Dominium Architecture, Application Layer, TestX, and CodeHygiene Planning](../\_reader/by\_chat/app\_testx\_codehygiene.md), [Dominium Architecture, UI, Providers, and Robot OS Strategy](../\_reader/by\_chat/architecture\_ui\_providers.md), [Dominium Build and Future-Proofing Architecture](../\_reader/by\_chat/build\_and\_future\_proofing.md).
\item worldgen (39 conversations): world, universe, celestial, terrain, chronology, and large-scale reality modeling. Sources: [Dominium APP0 Runtime, Platform, and Renderer Architecture](../\_reader/by\_chat/app\_runtime\_platform\_renderers.md), [Dominium Architecture, Application Layer, TestX, and CodeHygiene Planning](../\_reader/by\_chat/app\_testx\_codehygiene.md), [Dominium/Domino Architecture and Codex Prompt Roadmap](../\_reader/by\_chat/architecture\_codex\_prompts.md), [Dominium Architecture, UI, Providers, and Robot OS Strategy](../\_reader/by\_chat/architecture\_ui\_providers.md), [Dominium Build and Future-Proofing Architecture](../\_reader/by\_chat/build\_and\_future\_proofing.md).
\item timekeeping (35 conversations): calendar, chronology, timestamp durability, 2038 resilience, and time-domain law. Sources: [Dominium Advanced Simulation and Infrastructure Architecture](../\_reader/by\_chat/advanced\_simulation\_infrastructure.md), [Dominium APP0 Runtime, Platform, and Renderer Architecture](../\_reader/by\_chat/app\_runtime\_platform\_renderers.md), [Dominium Architecture, Application Layer, TestX, and CodeHygiene Planning](../\_reader/by\_chat/app\_testx\_codehygiene.md), [Dominium/Domino Architecture and Codex Prompt Roadmap](../\_reader/by\_chat/architecture\_codex\_prompts.md), [Dominium Canonical Architecture, Repository Foundation, and Provider Model](../\_reader/by\_chat/canonical\_structure\_and\_framework.md).
\item setup\_launcher (25 conversations): setup, repair, rollback, launcher profiles, instances, and product orchestration. Sources: [Dominium APP0 Runtime, Platform, and Renderer Architecture](../\_reader/by\_chat/app\_runtime\_platform\_renderers.md), [Dominium Architecture, Application Layer, TestX, and CodeHygiene Planning](../\_reader/by\_chat/app\_testx\_codehygiene.md), [Dominium/Domino Architecture and Codex Prompt Roadmap](../\_reader/by\_chat/architecture\_codex\_prompts.md), [Dominium Canonical Architecture, Repository Foundation, and Provider Model](../\_reader/by\_chat/canonical\_structure\_and\_framework.md), [Documentation Standards, README Strategy, and Handoff Packaging](../\_reader/by\_chat/documentation\_standards\_readme.md).
\item workbench (20 conversations): operator-facing validation, evidence, inspection, and later authoring surfaces. Sources: [Dominium Advanced Simulation and Infrastructure Architecture](../\_reader/by\_chat/advanced\_simulation\_infrastructure.md), [Dominium Architecture, UI, Providers, and Robot OS Strategy](../\_reader/by\_chat/architecture\_ui\_providers.md), [Dominium Canonical Architecture, Repository Foundation, and Provider Model](../\_reader/by\_chat/canonical\_structure\_and\_framework.md), [Dominium Development Routes and Continuity Preservation](../\_reader/by\_chat/development\_routes.md), [Dominium Architecture IV](../\_reader/by\_chat/dominium\_architecture\_iv.md).
\item xstack\_aide (14 conversations): repo-control and assistant-operation systems around AIDE, XStack, AuditX, RepoX, and TestX. Sources: [Dominium Architecture, Application Layer, TestX, and CodeHygiene Planning](../\_reader/by\_chat/app\_testx\_codehygiene.md), [Dominium Canonical Architecture, Repository Foundation, and Provider Model](../\_reader/by\_chat/canonical\_structure\_and\_framework.md), [Dominium Canon, Repository Alignment, and Portability Doctrine](../\_reader/by\_chat/dominium\_complete\_conversation.md), [Dominium Workbench, AIDE, Presentation Architecture, Provider Strategy, Product-Spine Planning, and Preservation Companion](../\_reader/by\_chat/dominium\_full\_conversation.md), [Dominium Foundation Lock, Workbench Spine, and Parallel Codex Handoff](../\_reader/by\_chat/foundation\_workbench\_codex.md).
\end{itemize}



{\small\raggedright SOURCE PATH: docs/\allowbreak{}archive/\allowbreak{}conversations/\allowbreak{}\_\allowbreak{}synthesis/\allowbreak{}FULL\_\allowbreak{}PROJECT\_\allowbreak{}PICTURE\_\allowbreak{}v0.md\par}


\section{Full Project Picture v0}

This map combines current repo orientation with conversation-derived design intent. When they differ, current repo authority wins.


\subsection{Product Identity}

Domino is the deterministic substrate. Dominium is the official game, product, and domain layer. Workbench is a governed validation and inspection environment. Setup and launcher compose product instances. Tools validate, generate, audit, and migrate. Contracts and schema law define public identity and compatibility meaning.


\subsection{Boundary Map}

\begin{quote}\small Dense table content is summarized in this PDF rendering. See the Markdown source and HTML book for the full table.\end{quote}

\subsection{Topic Map}

\begin{itemize}
\item architecture (45 conversations): system boundaries, product shape, engine/game/client/server separation, and repository structure. Sources: [Dominium Advanced Simulation and Infrastructure Architecture](../\_reader/by\_chat/advanced\_simulation\_infrastructure.md), [Dominium APP0 Runtime, Platform, and Renderer Architecture](../\_reader/by\_chat/app\_runtime\_platform\_renderers.md), [Dominium Architecture, Application Layer, TestX, and CodeHygiene Planning](../\_reader/by\_chat/app\_testx\_codehygiene.md), [Dominium/Domino Architecture and Codex Prompt Roadmap](../\_reader/by\_chat/architecture\_codex\_prompts.md), [Dominium Architecture, UI, Providers, and Robot OS Strategy](../\_reader/by\_chat/architecture\_ui\_providers.md).
\item content (45 conversations): packs, mods, authored payload, GUI/content boundaries, and provider/content separation. Sources: [Dominium Advanced Simulation and Infrastructure Architecture](../\_reader/by\_chat/advanced\_simulation\_infrastructure.md), [Dominium APP0 Runtime, Platform, and Renderer Architecture](../\_reader/by\_chat/app\_runtime\_platform\_renderers.md), [Dominium Architecture, Application Layer, TestX, and CodeHygiene Planning](../\_reader/by\_chat/app\_testx\_codehygiene.md), [Dominium/Domino Architecture and Codex Prompt Roadmap](../\_reader/by\_chat/architecture\_codex\_prompts.md), [Dominium Architecture, UI, Providers, and Robot OS Strategy](../\_reader/by\_chat/architecture\_ui\_providers.md).
\item ui (45 conversations): presentation, tools, editor concepts, perceived model surfaces, and command/result display. Sources: [Dominium Advanced Simulation and Infrastructure Architecture](../\_reader/by\_chat/advanced\_simulation\_infrastructure.md), [Dominium APP0 Runtime, Platform, and Renderer Architecture](../\_reader/by\_chat/app\_runtime\_platform\_renderers.md), [Dominium Architecture, Application Layer, TestX, and CodeHygiene Planning](../\_reader/by\_chat/app\_testx\_codehygiene.md), [Dominium/Domino Architecture and Codex Prompt Roadmap](../\_reader/by\_chat/architecture\_codex\_prompts.md), [Dominium Architecture, UI, Providers, and Robot OS Strategy](../\_reader/by\_chat/architecture\_ui\_providers.md).
\item governance (44 conversations): authority order, canon, contracts, refusal, review gates, and agent operation. Sources: [Dominium Advanced Simulation and Infrastructure Architecture](../\_reader/by\_chat/advanced\_simulation\_infrastructure.md), [Dominium APP0 Runtime, Platform, and Renderer Architecture](../\_reader/by\_chat/app\_runtime\_platform\_renderers.md), [Dominium Architecture, Application Layer, TestX, and CodeHygiene Planning](../\_reader/by\_chat/app\_testx\_codehygiene.md), [Dominium/Domino Architecture and Codex Prompt Roadmap](../\_reader/by\_chat/architecture\_codex\_prompts.md), [Dominium Architecture, UI, Providers, and Robot OS Strategy](../\_reader/by\_chat/architecture\_ui\_providers.md).
\item simulation (44 conversations): deterministic domains, processes, machines, physical systems, economy, and civilization dynamics. Sources: [Dominium Advanced Simulation and Infrastructure Architecture](../\_reader/by\_chat/advanced\_simulation\_infrastructure.md), [Dominium APP0 Runtime, Platform, and Renderer Architecture](../\_reader/by\_chat/app\_runtime\_platform\_renderers.md), [Dominium Architecture, Application Layer, TestX, and CodeHygiene Planning](../\_reader/by\_chat/app\_testx\_codehygiene.md), [Dominium/Domino Architecture and Codex Prompt Roadmap](../\_reader/by\_chat/architecture\_codex\_prompts.md), [Dominium Architecture, UI, Providers, and Robot OS Strategy](../\_reader/by\_chat/architecture\_ui\_providers.md).
\item tooling (44 conversations): Codex/AIDE/XStack/TestX/RepoX/AuditX tooling and patch workflow. Sources: [Dominium Advanced Simulation and Infrastructure Architecture](../\_reader/by\_chat/advanced\_simulation\_infrastructure.md), [Dominium APP0 Runtime, Platform, and Renderer Architecture](../\_reader/by\_chat/app\_runtime\_platform\_renderers.md), [Dominium Architecture, Application Layer, TestX, and CodeHygiene Planning](../\_reader/by\_chat/app\_testx\_codehygiene.md), [Dominium/Domino Architecture and Codex Prompt Roadmap](../\_reader/by\_chat/architecture\_codex\_prompts.md), [Dominium Architecture, UI, Providers, and Robot OS Strategy](../\_reader/by\_chat/architecture\_ui\_providers.md).
\item release (43 conversations): version identity, packaging, setup, update, publication, and release-control concerns. Sources: [Dominium Advanced Simulation and Infrastructure Architecture](../\_reader/by\_chat/advanced\_simulation\_infrastructure.md), [Dominium APP0 Runtime, Platform, and Renderer Architecture](../\_reader/by\_chat/app\_runtime\_platform\_renderers.md), [Dominium Architecture, Application Layer, TestX, and CodeHygiene Planning](../\_reader/by\_chat/app\_testx\_codehygiene.md), [Dominium/Domino Architecture and Codex Prompt Roadmap](../\_reader/by\_chat/architecture\_codex\_prompts.md), [Dominium Architecture, UI, Providers, and Robot OS Strategy](../\_reader/by\_chat/architecture\_ui\_providers.md).
\item contracts\_schema (42 conversations): machine-readable contracts, schemas, compatibility, manifests, and registries. Sources: [Dominium Advanced Simulation and Infrastructure Architecture](../\_reader/by\_chat/advanced\_simulation\_infrastructure.md), [Dominium APP0 Runtime, Platform, and Renderer Architecture](../\_reader/by\_chat/app\_runtime\_platform\_renderers.md), [Dominium Architecture, Application Layer, TestX, and CodeHygiene Planning](../\_reader/by\_chat/app\_testx\_codehygiene.md), [Dominium/Domino Architecture and Codex Prompt Roadmap](../\_reader/by\_chat/architecture\_codex\_prompts.md), [Dominium Architecture, UI, Providers, and Robot OS Strategy](../\_reader/by\_chat/architecture\_ui\_providers.md).
\item determinism (39 conversations): replay, ordering, proof, fixed identity, RNG discipline, and thread-count invariance. Sources: [Dominium Advanced Simulation and Infrastructure Architecture](../\_reader/by\_chat/advanced\_simulation\_infrastructure.md), [Dominium APP0 Runtime, Platform, and Renderer Architecture](../\_reader/by\_chat/app\_runtime\_platform\_renderers.md), [Dominium Architecture, Application Layer, TestX, and CodeHygiene Planning](../\_reader/by\_chat/app\_testx\_codehygiene.md), [Dominium/Domino Architecture and Codex Prompt Roadmap](../\_reader/by\_chat/architecture\_codex\_prompts.md), [Dominium Architecture, UI, Providers, and Robot OS Strategy](../\_reader/by\_chat/architecture\_ui\_providers.md).
\item platform (39 conversations): runtime shell, platform adapter, renderer boundary, operating-system support, and portability. Sources: [Dominium Advanced Simulation and Infrastructure Architecture](../\_reader/by\_chat/advanced\_simulation\_infrastructure.md), [Dominium APP0 Runtime, Platform, and Renderer Architecture](../\_reader/by\_chat/app\_runtime\_platform\_renderers.md), [Dominium Architecture, Application Layer, TestX, and CodeHygiene Planning](../\_reader/by\_chat/app\_testx\_codehygiene.md), [Dominium Architecture, UI, Providers, and Robot OS Strategy](../\_reader/by\_chat/architecture\_ui\_providers.md), [Dominium Build and Future-Proofing Architecture](../\_reader/by\_chat/build\_and\_future\_proofing.md).
\item worldgen (39 conversations): world, universe, celestial, terrain, chronology, and large-scale reality modeling. Sources: [Dominium APP0 Runtime, Platform, and Renderer Architecture](../\_reader/by\_chat/app\_runtime\_platform\_renderers.md), [Dominium Architecture, Application Layer, TestX, and CodeHygiene Planning](../\_reader/by\_chat/app\_testx\_codehygiene.md), [Dominium/Domino Architecture and Codex Prompt Roadmap](../\_reader/by\_chat/architecture\_codex\_prompts.md), [Dominium Architecture, UI, Providers, and Robot OS Strategy](../\_reader/by\_chat/architecture\_ui\_providers.md), [Dominium Build and Future-Proofing Architecture](../\_reader/by\_chat/build\_and\_future\_proofing.md).
\item timekeeping (35 conversations): calendar, chronology, timestamp durability, 2038 resilience, and time-domain law. Sources: [Dominium Advanced Simulation and Infrastructure Architecture](../\_reader/by\_chat/advanced\_simulation\_infrastructure.md), [Dominium APP0 Runtime, Platform, and Renderer Architecture](../\_reader/by\_chat/app\_runtime\_platform\_renderers.md), [Dominium Architecture, Application Layer, TestX, and CodeHygiene Planning](../\_reader/by\_chat/app\_testx\_codehygiene.md), [Dominium/Domino Architecture and Codex Prompt Roadmap](../\_reader/by\_chat/architecture\_codex\_prompts.md), [Dominium Canonical Architecture, Repository Foundation, and Provider Model](../\_reader/by\_chat/canonical\_structure\_and\_framework.md).
\item setup\_launcher (25 conversations): setup, repair, rollback, launcher profiles, instances, and product orchestration. Sources: [Dominium APP0 Runtime, Platform, and Renderer Architecture](../\_reader/by\_chat/app\_runtime\_platform\_renderers.md), [Dominium Architecture, Application Layer, TestX, and CodeHygiene Planning](../\_reader/by\_chat/app\_testx\_codehygiene.md), [Dominium/Domino Architecture and Codex Prompt Roadmap](../\_reader/by\_chat/architecture\_codex\_prompts.md), [Dominium Canonical Architecture, Repository Foundation, and Provider Model](../\_reader/by\_chat/canonical\_structure\_and\_framework.md), [Documentation Standards, README Strategy, and Handoff Packaging](../\_reader/by\_chat/documentation\_standards\_readme.md).
\item workbench (20 conversations): operator-facing validation, evidence, inspection, and later authoring surfaces. Sources: [Dominium Advanced Simulation and Infrastructure Architecture](../\_reader/by\_chat/advanced\_simulation\_infrastructure.md), [Dominium Architecture, UI, Providers, and Robot OS Strategy](../\_reader/by\_chat/architecture\_ui\_providers.md), [Dominium Canonical Architecture, Repository Foundation, and Provider Model](../\_reader/by\_chat/canonical\_structure\_and\_framework.md), [Dominium Development Routes and Continuity Preservation](../\_reader/by\_chat/development\_routes.md), [Dominium Architecture IV](../\_reader/by\_chat/dominium\_architecture\_iv.md).
\item xstack\_aide (14 conversations): repo-control and assistant-operation systems around AIDE, XStack, AuditX, RepoX, and TestX. Sources: [Dominium Architecture, Application Layer, TestX, and CodeHygiene Planning](../\_reader/by\_chat/app\_testx\_codehygiene.md), [Dominium Canonical Architecture, Repository Foundation, and Provider Model](../\_reader/by\_chat/canonical\_structure\_and\_framework.md), [Dominium Canon, Repository Alignment, and Portability Doctrine](../\_reader/by\_chat/dominium\_complete\_conversation.md), [Dominium Workbench, AIDE, Presentation Architecture, Provider Strategy, Product-Spine Planning, and Preservation Companion](../\_reader/by\_chat/dominium\_full\_conversation.md), [Dominium Foundation Lock, Workbench Spine, and Parallel Codex Handoff](../\_reader/by\_chat/foundation\_workbench\_codex.md).
\end{itemize}

\subsection{Blocked Versus Open}

Open for derived archive work: reading, synthesis, crosswalks, contradiction review, promotion triage, and narrow docs-only planning. Blocked by current queue: broad Workbench UI, runtime module loader, provider runtime, package runtime, gameplay, renderer implementation, native GUI, and release publication.




{\small\raggedright SOURCE PATH: docs/\allowbreak{}archive/\allowbreak{}conversations/\allowbreak{}\_\allowbreak{}synthesis/\allowbreak{}CONTRADICTIONS\_\allowbreak{}TO\_\allowbreak{}RECONCILE\_\allowbreak{}v0.md\par}


\section{Contradictions To Reconcile v0}

This document summarizes contradiction classes that should feed the reconciliation crosswalk. It does not resolve them.


\subsection{Counts By Class}

\begin{itemize}
\item conversation\_vs\_conversation: 2
\item conversation\_vs\_current\_queue: 102
\item conversation\_vs\_docs: 2
\item stale\_external\_claim: 121
\end{itemize}

\subsection{Highest-Priority Classes}

\begin{itemize}
\item conversation\_vs\_current\_queue: old work suggestions that touch currently blocked scope.
\item conversation\_vs\_docs: old baseline or architecture statements that may drift from current README/canon.
\item stale\_external\_claim: external platform, SDK, language, renderer, or runtime claims needing current verification.
\item conversation\_vs\_conversation: disagreements among old conversations that need review before promotion.
\end{itemize}

\subsection{Sample Findings}

\begin{itemize}
\item CONTRA-0001 conversation\_vs\_current\_queue advanced\_simulation\_infrastructure: Archived conversation discusses work related to blocked scope gameplay.
\item CONTRA-0002 stale\_external\_claim advanced\_simulation\_infrastructure: Archived conversation contains potentially stale external or baseline claim: C89.
\item CONTRA-0003 stale\_external\_claim advanced\_simulation\_infrastructure: Archived conversation contains potentially stale external or baseline claim: C++98.
\item CONTRA-0004 stale\_external\_claim advanced\_simulation\_infrastructure: Archived conversation contains potentially stale external or baseline claim: ISO C89.
\item CONTRA-0005 conversation\_vs\_current\_queue app\_runtime\_platform\_renderers: Archived conversation discusses work related to blocked scope runtime\_module\_loader.
\item CONTRA-0006 conversation\_vs\_current\_queue app\_runtime\_platform\_renderers: Archived conversation discusses work related to blocked scope provider\_runtime.
\item CONTRA-0007 conversation\_vs\_current\_queue app\_runtime\_platform\_renderers: Archived conversation discusses work related to blocked scope gameplay.
\item CONTRA-0008 conversation\_vs\_current\_queue app\_runtime\_platform\_renderers: Archived conversation discusses work related to blocked scope renderer\_implementation.
\item CONTRA-0009 stale\_external\_claim app\_runtime\_platform\_renderers: Archived conversation contains potentially stale external or baseline claim: SDK.
\item CONTRA-0010 conversation\_vs\_current\_queue app\_testx\_codehygiene: Archived conversation discusses work related to blocked scope gameplay.
\item CONTRA-0011 conversation\_vs\_current\_queue app\_testx\_codehygiene: Archived conversation discusses work related to blocked scope renderer\_implementation.
\item CONTRA-0012 stale\_external\_claim app\_testx\_codehygiene: Archived conversation contains potentially stale external or baseline claim: C89.
\item CONTRA-0013 stale\_external\_claim app\_testx\_codehygiene: Archived conversation contains potentially stale external or baseline claim: C++98.
\item CONTRA-0014 stale\_external\_claim app\_testx\_codehygiene: Archived conversation contains potentially stale external or baseline claim: DX12.
\item CONTRA-0015 stale\_external\_claim app\_testx\_codehygiene: Archived conversation contains potentially stale external or baseline claim: iOS.
\item CONTRA-0016 conversation\_vs\_current\_queue architecture\_codex\_prompts: Archived conversation discusses work related to blocked scope provider\_runtime.
\item CONTRA-0017 conversation\_vs\_current\_queue architecture\_codex\_prompts: Archived conversation discusses work related to blocked scope gameplay.
\item CONTRA-0018 stale\_external\_claim architecture\_codex\_prompts: Archived conversation contains potentially stale external or baseline claim: C89.
\item CONTRA-0019 stale\_external\_claim architecture\_codex\_prompts: Archived conversation contains potentially stale external or baseline claim: C++98.
\item CONTRA-0020 stale\_external\_claim architecture\_codex\_prompts: Archived conversation contains potentially stale external or baseline claim: ISO C89.
\item CONTRA-0021 conversation\_vs\_current\_queue architecture\_ui\_providers: Archived conversation discusses work related to blocked scope broad\_workbench\_ui.
\item CONTRA-0022 conversation\_vs\_current\_queue architecture\_ui\_providers: Archived conversation discusses work related to blocked scope provider\_runtime.
\item CONTRA-0023 conversation\_vs\_current\_queue architecture\_ui\_providers: Archived conversation discusses work related to blocked scope gameplay.
\item CONTRA-0024 conversation\_vs\_current\_queue architecture\_ui\_providers: Archived conversation discusses work related to blocked scope renderer\_implementation.
\item CONTRA-0025 conversation\_vs\_current\_queue architecture\_ui\_providers: Archived conversation discusses work related to blocked scope native\_gui.
\item CONTRA-0026 stale\_external\_claim architecture\_ui\_providers: Archived conversation contains potentially stale external or baseline claim: C89.
\item CONTRA-0027 stale\_external\_claim architecture\_ui\_providers: Archived conversation contains potentially stale external or baseline claim: C++98.
\item CONTRA-0028 stale\_external\_claim architecture\_ui\_providers: Archived conversation contains potentially stale external or baseline claim: DX12.
\item CONTRA-0029 stale\_external\_claim architecture\_ui\_providers: Archived conversation contains potentially stale external or baseline claim: SDK.
\item CONTRA-0030 stale\_external\_claim Build\_and\_Future\_Proofing: Archived conversation contains potentially stale external or baseline claim: C89.
\end{itemize}

\subsection{Reconciliation Rule}

Resolve nothing by convenience. Check canon, glossary, AGENTS, authority order, current queue, contracts/schema law, and targeted current docs before promoting any claim.



\chapter{Authority and Reconciliation}


{\small\raggedright SOURCE PATH: docs/\allowbreak{}archive/\allowbreak{}conversations/\allowbreak{}\_\allowbreak{}reconciliation/\allowbreak{}REPO\_\allowbreak{}AUTHORITY\_\allowbreak{}CROSSWALK\_\allowbreak{}v0.md\par}


\section{Repo Authority Crosswalk v0}

This crosswalk compares the derived synthesis to current authority surfaces. It does not promote any conversation claim.


\begin{quote}\small Dense table content is summarized in this PDF rendering. See the Markdown source and HTML book for the full table.\end{quote}

\subsection{Classification Counts}

\begin{itemize}
\item blocked\_by\_current\_queue: 10
\item consistent\_with\_repo\_but\_not\_formalized: 41
\item insufficient\_evidence: 46
\item needs\_user\_decision: 12
\item rejected\_noise: 19
\item stale\_or\_superseded: 7
\end{itemize}



{\small\raggedright SOURCE PATH: docs/\allowbreak{}archive/\allowbreak{}conversations/\allowbreak{}\_\allowbreak{}reconciliation/\allowbreak{}CURRENT\_\allowbreak{}CANON\_\allowbreak{}ALIGNMENT\_\allowbreak{}v0.md\par}


\section{Current Canon Alignment v0}

This review identifies broad areas where the conversation-derived synthesis appears aligned with current canon, and where caution remains.


\subsection{Aligned Signals}

\begin{itemize}
\item Determinism, replay, and provenance recur in the corpus and are directly supported by constitution axioms.
\item Truth/perceived/render separation recurs in UI, renderer, and Workbench conversations and is explicitly canonical.
\item Process-only mutation and law-gated authority are compatible with the archive's repeated refusal of UI-owned truth.
\item Pack-driven integration recurs in content/provider/modding conversations and is canonical in high-level doctrine.
\item AIDE/Codex/XStack as bounded repo-control surfaces aligns with AGENTS.md when they remain non-authoritative.
\end{itemize}

\subsection{Caution Areas}

\begin{itemize}
\item Conversation claims about implementation sequencing must obey current queue constraints.
\item Claims about renderer, native GUI, provider runtime, package runtime, gameplay, and release publication are blocked until stronger current authority opens them.
\item Old language/platform/baseline claims require current verification before reuse.
\item Synthesis terms not present in the glossary cannot become canonical without controlled glossary work.
\end{itemize}



{\small\raggedright SOURCE PATH: docs/\allowbreak{}archive/\allowbreak{}conversations/\allowbreak{}\_\allowbreak{}reconciliation/\allowbreak{}BLOCKED\_\allowbreak{}SCOPE\_\allowbreak{}ALIGNMENT\_\allowbreak{}v0.md\par}


\section{Blocked Scope Alignment v0}

Current queue state blocks broad feature work. Conversation-derived claims that touch these areas remain historical unless a later reviewed task opens scope.


\begin{quote}\small Dense table content is summarized in this PDF rendering. See the Markdown source and HTML book for the full table.\end{quote}

No row authorizes implementation. This is a guardrail for later claim review.




{\small\raggedright SOURCE PATH: docs/\allowbreak{}archive/\allowbreak{}conversations/\allowbreak{}\_\allowbreak{}reconciliation/\allowbreak{}CLAIM\_\allowbreak{}REVIEW\_\allowbreak{}MATRIX\_\allowbreak{}v0.md\par}


\section{Claim Review Matrix v0}

Each row is a raw promotion candidate classified for review. None are promoted.



\begin{quote}\small Dense table summarized for the reader edition. See the source file, HTML output, or reference appendix source for full detail.\end{quote}

\begin{quote}\small Dense table content is summarized in this PDF rendering. See the Markdown source and HTML book for the full table.\end{quote}


\chapter{Decision Docket}


{\small\raggedright SOURCE PATH: docs/\allowbreak{}archive/\allowbreak{}conversations/\allowbreak{}\_\allowbreak{}decision/\allowbreak{}DECISION\_\allowbreak{}SUMMARY\_\allowbreak{}v0.md\par}


\section{Decision Summary v0}

Summary of the generated decision docket. No decision is accepted here.


\subsection{Counts}

\begin{itemize}
\item Total decisions: 30
\item Workbench/AIDE/Codex/tooling: 6
\item architecture/boundaries: 12
\item provider/content/packs: 5
\item release/setup/launcher: 2
\item renderer/UI/platform: 3
\item world/time/civilization simulation: 2
\end{itemize}

\subsection{Owner Counts}

\begin{itemize}
\item future\_queue\_decision: 18
\item user\_decision: 12
\end{itemize}

\subsection{Recommended Defaults}

\begin{itemize}
\item defer\_pending\_user\_decision: 12
\item defer\_until\_queue\_opens: 10
\item keep\_blocked: 8
\end{itemize}

\subsection{Top 10 Decisions}

\begin{itemize}
\item DECISION-0001 architecture/boundaries: What disposition should be chosen for this unresolved claim: The user then asked whether arbitrary placement was possible or whether the system was stuck to grids. The answer established that the engine should be grid-agnostic. [FACT] The...?
\item DECISION-0002 architecture/boundaries: Should this conversation-derived claim remain archive-only until gameplay is explicitly opened?
\item DECISION-0003 architecture/boundaries: What disposition should be chosen for this unresolved claim: Future work must define the client's allowed dependency surface and whether it can access any simulation-adjacent prediction layer.?
\item DECISION-0004 architecture/boundaries: Should this conversation-derived claim remain archive-only until provider\_runtime is explicitly opened?
\item DECISION-0005 provider/content/packs: Should this conversation-derived claim remain archive-only until provider\_runtime is explicitly opened?
\item DECISION-0006 world/time/civilization simulation: Should this conversation-derived claim become a future review item, remain historical, or be deferred: The assistant formalized this as the Perfect Earth Calendar, later called HPC-E. The final structure is clear, though the exact leap rule remains unresolved.?
\item DECISION-0007 architecture/boundaries: What disposition should be chosen for this unresolved claim: Before relying on this chat as project authority, a future assistant or human reviewer must confirm whether Dominium is actually intended to be simulation-heavy, whether strict...?
\item DECISION-0008 renderer/UI/platform: Should this conversation-derived claim become a future review item, remain historical, or be deferred: The key thing to remember: this chat produced the standards and prompts for future work; it did not verify or change the project. The next assistant must inspect the actual repo...?
\item DECISION-0009 architecture/boundaries: What disposition should be chosen for this unresolved claim: This chat was about preserving and re-grounding Dominium's architecture around a very old constitutional contract, then checking whether the current GitHub repository still foll...?
\item DECISION-0010 architecture/boundaries: Should this conversation-derived claim remain archive-only until provider\_runtime is explicitly opened?
\end{itemize}

\subsection{Safe Near-Term Docs-Only Clarifications}

\begin{itemize}
\item None in this docket. Current generated decisions are user decisions, queue decisions, or deferrals.
\item Items touching blocked queue scope must remain explanatory or deferred; no implementation scope is opened.
\end{itemize}

\subsection{Recommended Deferrals}

\begin{itemize}
\item DECISION-0001: What disposition should be chosen for this unresolved claim: The user then asked whether arbitrary placement was possible or whether the system was stuck to grids. The answer established that the engine should be grid-agnostic. [FACT] The...?
\item DECISION-0002: Should this conversation-derived claim remain archive-only until gameplay is explicitly opened?
\item DECISION-0003: What disposition should be chosen for this unresolved claim: Future work must define the client's allowed dependency surface and whether it can access any simulation-adjacent prediction layer.?
\item DECISION-0004: Should this conversation-derived claim remain archive-only until provider\_runtime is explicitly opened?
\item DECISION-0005: Should this conversation-derived claim remain archive-only until provider\_runtime is explicitly opened?
\item DECISION-0006: Should this conversation-derived claim become a future review item, remain historical, or be deferred: The assistant formalized this as the Perfect Earth Calendar, later called HPC-E. The final structure is clear, though the exact leap rule remains unresolved.?
\item DECISION-0007: What disposition should be chosen for this unresolved claim: Before relying on this chat as project authority, a future assistant or human reviewer must confirm whether Dominium is actually intended to be simulation-heavy, whether strict...?
\item DECISION-0008: Should this conversation-derived claim become a future review item, remain historical, or be deferred: The key thing to remember: this chat produced the standards and prompts for future work; it did not verify or change the project. The next assistant must inspect the actual repo...?
\item DECISION-0009: What disposition should be chosen for this unresolved claim: This chat was about preserving and re-grounding Dominium's architecture around a very old constitutional contract, then checking whether the current GitHub repository still foll...?
\item DECISION-0010: Should this conversation-derived claim remain archive-only until provider\_runtime is explicitly opened?
\item DECISION-0011: Should this conversation-derived claim remain archive-only until provider\_runtime is explicitly opened?
\item DECISION-0012: Should this conversation-derived claim remain archive-only until native\_gui is explicitly opened?
\end{itemize}



{\small\raggedright SOURCE PATH: docs/\allowbreak{}archive/\allowbreak{}conversations/\allowbreak{}\_\allowbreak{}decision/\allowbreak{}DECISION\_\allowbreak{}DOCKET\_\allowbreak{}v0.md\par}


\section{Decision Docket v0}

This docket turns unresolved conversation-derived claims into review questions. It does not decide or promote them.


Decision items: 30


\subsection{architecture/boundaries}

\subsubsection{DECISION-0001}

\begin{itemize}
\item Source promotion/claim IDs: PROMOTE-0001
\item Source conversations: advanced\_simulation\_infrastructure
\item Question: What disposition should be chosen for this unresolved claim: The user then asked whether arbitrary placement was possible or whether the system was stuck to grids. The answer established that the engine should be grid-agnostic. [FACT] The...?
\item Why it matters: The claim frames unresolved future work or a decision boundary.
\item Current repo authority status: conversation-derived advisory claim; current repo authority not modified
\item Queue status: not directly blocked by queue keyword match
\item Recommended default: defer\_pending\_user\_decision
\item Decision owner: user\_decision
\item Risk if unresolved: Semantic drift or blocked-scope leakage if treated as current truth too early.
\item Evidence needed: Check canon, glossary, AGENTS.md, current queue, target current docs, and source conversation before any promotion.
\item Options:
\end{itemize}
  - promote\_later\_after\_authority\_review

  - preserve\_as\_history

  - defer

\begin{itemize}
\item Consequences:
\end{itemize}
  - promote\_later\_after\_authority\_review: creates a future docs-only or planning review task; still does not authorize implementation.

  - preserve\_as\_history: keeps the archive readable but removes pressure to patch current docs.

  - defer: leaves the claim unresolved until stronger authority or user decision exists.


\subsubsection{DECISION-0002}

\begin{itemize}
\item Source promotion/claim IDs: PROMOTE-0005
\item Source conversations: app\_runtime\_platform\_renderers
\item Question: Should this conversation-derived claim remain archive-only until gameplay is explicitly opened?
\item Why it matters: The claim touches currently blocked scope: gameplay.
\item Current repo authority status: conversation-derived advisory claim; current repo authority not modified
\item Queue status: gameplay
\item Recommended default: defer\_until\_queue\_opens
\item Decision owner: future\_queue\_decision
\item Risk if unresolved: Semantic drift or blocked-scope leakage if treated as current truth too early.
\item Evidence needed: Check canon, glossary, AGENTS.md, current queue, target current docs, and source conversation before any promotion.
\item Options:
\end{itemize}
  - promote\_later\_after\_authority\_review

  - preserve\_as\_history

  - defer

\begin{itemize}
\item Consequences:
\end{itemize}
  - promote\_later\_after\_authority\_review: creates a future docs-only or planning review task; still does not authorize implementation.

  - preserve\_as\_history: keeps the archive readable but removes pressure to patch current docs.

  - defer: leaves the claim unresolved until stronger authority or user decision exists.


\subsubsection{DECISION-0003}

\begin{itemize}
\item Source promotion/claim IDs: PROMOTE-0006
\item Source conversations: app\_runtime\_platform\_renderers
\item Question: What disposition should be chosen for this unresolved claim: Future work must define the client's allowed dependency surface and whether it can access any simulation-adjacent prediction layer.?
\item Why it matters: The claim frames unresolved future work or a decision boundary.
\item Current repo authority status: conversation-derived advisory claim; current repo authority not modified
\item Queue status: not directly blocked by queue keyword match
\item Recommended default: defer\_pending\_user\_decision
\item Decision owner: user\_decision
\item Risk if unresolved: Semantic drift or blocked-scope leakage if treated as current truth too early.
\item Evidence needed: Check canon, glossary, AGENTS.md, current queue, target current docs, and source conversation before any promotion.
\item Options:
\end{itemize}
  - promote\_later\_after\_authority\_review

  - preserve\_as\_history

  - defer

\begin{itemize}
\item Consequences:
\end{itemize}
  - promote\_later\_after\_authority\_review: creates a future docs-only or planning review task; still does not authorize implementation.

  - preserve\_as\_history: keeps the archive readable but removes pressure to patch current docs.

  - defer: leaves the claim unresolved until stronger authority or user decision exists.


\subsubsection{DECISION-0004}

\begin{itemize}
\item Source promotion/claim IDs: PROMOTE-0015
\item Source conversations: architecture\_ui\_providers
\item Question: Should this conversation-derived claim remain archive-only until provider\_runtime is explicitly opened?
\item Why it matters: The claim touches currently blocked scope: provider\_runtime.
\item Current repo authority status: conversation-derived advisory claim; current repo authority not modified
\item Queue status: provider\_runtime
\item Recommended default: defer\_until\_queue\_opens
\item Decision owner: future\_queue\_decision
\item Risk if unresolved: Semantic drift or blocked-scope leakage if treated as current truth too early.
\item Evidence needed: Check canon, glossary, AGENTS.md, current queue, target current docs, and source conversation before any promotion.
\item Options:
\end{itemize}
  - promote\_later\_after\_authority\_review

  - preserve\_as\_history

  - defer

\begin{itemize}
\item Consequences:
\end{itemize}
  - promote\_later\_after\_authority\_review: creates a future docs-only or planning review task; still does not authorize implementation.

  - preserve\_as\_history: keeps the archive readable but removes pressure to patch current docs.

  - defer: leaves the claim unresolved until stronger authority or user decision exists.


\subsubsection{DECISION-0007}

\begin{itemize}
\item Source promotion/claim IDs: PROMOTE-0025
\item Source conversations: development\_routes
\item Question: What disposition should be chosen for this unresolved claim: Before relying on this chat as project authority, a future assistant or human reviewer must confirm whether Dominium is actually intended to be simulation-heavy, whether strict...?
\item Why it matters: The claim frames unresolved future work or a decision boundary.
\item Current repo authority status: conversation-derived advisory claim; current repo authority not modified
\item Queue status: not directly blocked by queue keyword match
\item Recommended default: defer\_pending\_user\_decision
\item Decision owner: user\_decision
\item Risk if unresolved: Semantic drift or blocked-scope leakage if treated as current truth too early.
\item Evidence needed: Check canon, glossary, AGENTS.md, current queue, target current docs, and source conversation before any promotion.
\item Options:
\end{itemize}
  - promote\_later\_after\_authority\_review

  - preserve\_as\_history

  - defer

\begin{itemize}
\item Consequences:
\end{itemize}
  - promote\_later\_after\_authority\_review: creates a future docs-only or planning review task; still does not authorize implementation.

  - preserve\_as\_history: keeps the archive readable but removes pressure to patch current docs.

  - defer: leaves the claim unresolved until stronger authority or user decision exists.


\subsubsection{DECISION-0009}

\begin{itemize}
\item Source promotion/claim IDs: PROMOTE-0043
\item Source conversations: Dominium\_Complete\_Conversation
\item Question: What disposition should be chosen for this unresolved claim: This chat was about preserving and re-grounding Dominium's architecture around a very old constitutional contract, then checking whether the current GitHub repository still foll...?
\item Why it matters: The claim frames unresolved future work or a decision boundary.
\item Current repo authority status: conversation-derived advisory claim; current repo authority not modified
\item Queue status: not directly blocked by queue keyword match
\item Recommended default: defer\_pending\_user\_decision
\item Decision owner: user\_decision
\item Risk if unresolved: Semantic drift or blocked-scope leakage if treated as current truth too early.
\item Evidence needed: Check canon, glossary, AGENTS.md, current queue, target current docs, and source conversation before any promotion.
\item Options:
\end{itemize}
  - promote\_later\_after\_authority\_review

  - preserve\_as\_history

  - defer

\begin{itemize}
\item Consequences:
\end{itemize}
  - promote\_later\_after\_authority\_review: creates a future docs-only or planning review task; still does not authorize implementation.

  - preserve\_as\_history: keeps the archive readable but removes pressure to patch current docs.

  - defer: leaves the claim unresolved until stronger authority or user decision exists.


\subsubsection{DECISION-0010}

\begin{itemize}
\item Source promotion/claim IDs: PROMOTE-0049
\item Source conversations: dominium\_full\_conversation
\item Question: Should this conversation-derived claim remain archive-only until provider\_runtime is explicitly opened?
\item Why it matters: The claim touches currently blocked scope: provider\_runtime.
\item Current repo authority status: conversation-derived advisory claim; current repo authority not modified
\item Queue status: provider\_runtime
\item Recommended default: defer\_until\_queue\_opens
\item Decision owner: future\_queue\_decision
\item Risk if unresolved: Semantic drift or blocked-scope leakage if treated as current truth too early.
\item Evidence needed: Check canon, glossary, AGENTS.md, current queue, target current docs, and source conversation before any promotion.
\item Options:
\end{itemize}
  - promote\_later\_after\_authority\_review

  - preserve\_as\_history

  - defer

\begin{itemize}
\item Consequences:
\end{itemize}
  - promote\_later\_after\_authority\_review: creates a future docs-only or planning review task; still does not authorize implementation.

  - preserve\_as\_history: keeps the archive readable but removes pressure to patch current docs.

  - defer: leaves the claim unresolved until stronger authority or user decision exists.


\subsubsection{DECISION-0014}

\begin{itemize}
\item Source promotion/claim IDs: PROMOTE-0069
\item Source conversations: Framework\_Open\_Source\_Provider
\item Question: Should this conversation-derived claim remain archive-only until provider\_runtime is explicitly opened?
\item Why it matters: The claim touches currently blocked scope: provider\_runtime.
\item Current repo authority status: conversation-derived advisory claim; current repo authority not modified
\item Queue status: provider\_runtime
\item Recommended default: defer\_until\_queue\_opens
\item Decision owner: future\_queue\_decision
\item Risk if unresolved: Semantic drift or blocked-scope leakage if treated as current truth too early.
\item Evidence needed: Check canon, glossary, AGENTS.md, current queue, target current docs, and source conversation before any promotion.
\item Options:
\end{itemize}
  - promote\_later\_after\_authority\_review

  - preserve\_as\_history

  - defer

\begin{itemize}
\item Consequences:
\end{itemize}
  - promote\_later\_after\_authority\_review: creates a future docs-only or planning review task; still does not authorize implementation.

  - preserve\_as\_history: keeps the archive readable but removes pressure to patch current docs.

  - defer: leaves the claim unresolved until stronger authority or user decision exists.


\subsubsection{DECISION-0015}

\begin{itemize}
\item Source promotion/claim IDs: PROMOTE-0075
\item Source conversations: Language\_Platform\_Architecture
\item Question: Should this conversation-derived claim remain archive-only until provider\_runtime is explicitly opened?
\item Why it matters: The claim touches currently blocked scope: provider\_runtime.
\item Current repo authority status: conversation-derived advisory claim; current repo authority not modified
\item Queue status: provider\_runtime
\item Recommended default: defer\_until\_queue\_opens
\item Decision owner: future\_queue\_decision
\item Risk if unresolved: Semantic drift or blocked-scope leakage if treated as current truth too early.
\item Evidence needed: Check canon, glossary, AGENTS.md, current queue, target current docs, and source conversation before any promotion.
\item Options:
\end{itemize}
  - promote\_later\_after\_authority\_review

  - preserve\_as\_history

  - defer

\begin{itemize}
\item Consequences:
\end{itemize}
  - promote\_later\_after\_authority\_review: creates a future docs-only or planning review task; still does not authorize implementation.

  - preserve\_as\_history: keeps the archive readable but removes pressure to patch current docs.

  - defer: leaves the claim unresolved until stronger authority or user decision exists.


\subsubsection{DECISION-0018}

\begin{itemize}
\item Source promotion/claim IDs: PROMOTE-0113
\item Source conversations: testx\_repox\_governance
\item Question: What disposition should be chosen for this unresolved claim: The user then asked whether the implementation was industry-accepted and what could be improved. The response framed the approach as closer to game engines, operating systems, a...?
\item Why it matters: The claim frames unresolved future work or a decision boundary.
\item Current repo authority status: conversation-derived advisory claim; current repo authority not modified
\item Queue status: not directly blocked by queue keyword match
\item Recommended default: defer\_pending\_user\_decision
\item Decision owner: user\_decision
\item Risk if unresolved: Semantic drift or blocked-scope leakage if treated as current truth too early.
\item Evidence needed: Check canon, glossary, AGENTS.md, current queue, target current docs, and source conversation before any promotion.
\item Options:
\end{itemize}
  - promote\_later\_after\_authority\_review

  - preserve\_as\_history

  - defer

\begin{itemize}
\item Consequences:
\end{itemize}
  - promote\_later\_after\_authority\_review: creates a future docs-only or planning review task; still does not authorize implementation.

  - preserve\_as\_history: keeps the archive readable but removes pressure to patch current docs.

  - defer: leaves the claim unresolved until stronger authority or user decision exists.


\subsubsection{DECISION-0019}

\begin{itemize}
\item Source promotion/claim IDs: PROMOTE-0120
\item Source conversations: UE6\_Domino\_Deterministic\_Universe
\item Question: Should this conversation-derived claim remain archive-only until gameplay is explicitly opened?
\item Why it matters: The claim touches currently blocked scope: gameplay.
\item Current repo authority status: conversation-derived advisory claim; current repo authority not modified
\item Queue status: gameplay
\item Recommended default: defer\_until\_queue\_opens
\item Decision owner: future\_queue\_decision
\item Risk if unresolved: Semantic drift or blocked-scope leakage if treated as current truth too early.
\item Evidence needed: Check canon, glossary, AGENTS.md, current queue, target current docs, and source conversation before any promotion.
\item Options:
\end{itemize}
  - promote\_later\_after\_authority\_review

  - preserve\_as\_history

  - defer

\begin{itemize}
\item Consequences:
\end{itemize}
  - promote\_later\_after\_authority\_review: creates a future docs-only or planning review task; still does not authorize implementation.

  - preserve\_as\_history: keeps the archive readable but removes pressure to patch current docs.

  - defer: leaves the claim unresolved until stronger authority or user decision exists.


\subsubsection{DECISION-0024}

\begin{itemize}
\item Source promotion/claim IDs: blocked\_scope:gameplay
\item Source conversations: current queue plus conversation audit findings
\item Question: Should gameplay remain blocked for all conversation-derived work until a future reviewed queue phase opens it?
\item Why it matters: The current queue marks gameplay as BLOCKED and the corpus produced 35 related candidate/finding signals.
\item Current repo authority status: .aide/queue/current.toml is current repo authority for this scope gate.
\item Queue status: gameplay: BLOCKED
\item Recommended default: keep\_blocked
\item Decision owner: future\_queue\_decision
\item Risk if unresolved: Implementation or live-doc scope could be opened by historical momentum rather than current authority.
\item Evidence needed: Review current queue, blocked-scope alignment, and any target prompt before changing disposition.
\item Options:
\end{itemize}
  - keep\_blocked

  - docs\_only\_crosswalk

  - defer

\begin{itemize}
\item Consequences:
\end{itemize}
  - keep\_blocked: prevents old conversation claims from opening implementation scope.

  - docs\_only\_crosswalk: allows explanatory archive/current-doc mapping without implementation or promotion.

  - defer: postpones judgment until the queue is revised.


\subsection{Workbench/AIDE/Codex/tooling}

\subsubsection{DECISION-0011}

\begin{itemize}
\item Source promotion/claim IDs: PROMOTE-0056
\item Source conversations: Domino\_Dominium\_Workbench
\item Question: Should this conversation-derived claim remain archive-only until provider\_runtime is explicitly opened?
\item Why it matters: The claim touches currently blocked scope: provider\_runtime.
\item Current repo authority status: conversation-derived advisory claim; current repo authority not modified
\item Queue status: provider\_runtime
\item Recommended default: defer\_until\_queue\_opens
\item Decision owner: future\_queue\_decision
\item Risk if unresolved: Semantic drift or blocked-scope leakage if treated as current truth too early.
\item Evidence needed: Check canon, glossary, AGENTS.md, current queue, target current docs, and source conversation before any promotion.
\item Options:
\end{itemize}
  - promote\_later\_after\_authority\_review

  - preserve\_as\_history

  - defer

\begin{itemize}
\item Consequences:
\end{itemize}
  - promote\_later\_after\_authority\_review: creates a future docs-only or planning review task; still does not authorize implementation.

  - preserve\_as\_history: keeps the archive readable but removes pressure to patch current docs.

  - defer: leaves the claim unresolved until stronger authority or user decision exists.


\subsubsection{DECISION-0012}

\begin{itemize}
\item Source promotion/claim IDs: PROMOTE-0057
\item Source conversations: Domino\_Dominium\_Workbench
\item Question: Should this conversation-derived claim remain archive-only until native\_gui is explicitly opened?
\item Why it matters: The claim touches currently blocked scope: native\_gui.
\item Current repo authority status: conversation-derived advisory claim; current repo authority not modified
\item Queue status: native\_gui
\item Recommended default: defer\_until\_queue\_opens
\item Decision owner: future\_queue\_decision
\item Risk if unresolved: Semantic drift or blocked-scope leakage if treated as current truth too early.
\item Evidence needed: Check canon, glossary, AGENTS.md, current queue, target current docs, and source conversation before any promotion.
\item Options:
\end{itemize}
  - promote\_later\_after\_authority\_review

  - preserve\_as\_history

  - defer

\begin{itemize}
\item Consequences:
\end{itemize}
  - promote\_later\_after\_authority\_review: creates a future docs-only or planning review task; still does not authorize implementation.

  - preserve\_as\_history: keeps the archive readable but removes pressure to patch current docs.

  - defer: leaves the claim unresolved until stronger authority or user decision exists.


\subsubsection{DECISION-0013}

\begin{itemize}
\item Source promotion/claim IDs: PROMOTE-0065
\item Source conversations: Foundation\_Workbench\_Codex
\item Question: Should this conversation-derived claim remain archive-only until provider\_runtime is explicitly opened?
\item Why it matters: The claim touches currently blocked scope: provider\_runtime.
\item Current repo authority status: conversation-derived advisory claim; current repo authority not modified
\item Queue status: provider\_runtime
\item Recommended default: defer\_until\_queue\_opens
\item Decision owner: future\_queue\_decision
\item Risk if unresolved: Semantic drift or blocked-scope leakage if treated as current truth too early.
\item Evidence needed: Check canon, glossary, AGENTS.md, current queue, target current docs, and source conversation before any promotion.
\item Options:
\end{itemize}
  - promote\_later\_after\_authority\_review

  - preserve\_as\_history

  - defer

\begin{itemize}
\item Consequences:
\end{itemize}
  - promote\_later\_after\_authority\_review: creates a future docs-only or planning review task; still does not authorize implementation.

  - preserve\_as\_history: keeps the archive readable but removes pressure to patch current docs.

  - defer: leaves the claim unresolved until stronger authority or user decision exists.


\subsubsection{DECISION-0020}

\begin{itemize}
\item Source promotion/claim IDs: PROMOTE-0121
\item Source conversations: ui\_editor\_tool\_editor\_planning
\item Question: What disposition should be chosen for this unresolved claim: The most important thing to remember is that this chat was a planning and prompt-generation chat, not proof of implementation. **UNCERTAIN / UNVERIFIED:** No generated prompt is...?
\item Why it matters: The claim frames unresolved future work or a decision boundary.
\item Current repo authority status: conversation-derived advisory claim; current repo authority not modified
\item Queue status: not directly blocked by queue keyword match
\item Recommended default: defer\_pending\_user\_decision
\item Decision owner: user\_decision
\item Risk if unresolved: Semantic drift or blocked-scope leakage if treated as current truth too early.
\item Evidence needed: Check canon, glossary, AGENTS.md, current queue, target current docs, and source conversation before any promotion.
\item Options:
\end{itemize}
  - promote\_later\_after\_authority\_review

  - preserve\_as\_history

  - defer

\begin{itemize}
\item Consequences:
\end{itemize}
  - promote\_later\_after\_authority\_review: creates a future docs-only or planning review task; still does not authorize implementation.

  - preserve\_as\_history: keeps the archive readable but removes pressure to patch current docs.

  - defer: leaves the claim unresolved until stronger authority or user decision exists.


\subsubsection{DECISION-0021}

\begin{itemize}
\item Source promotion/claim IDs: PROMOTE-0123
\item Source conversations: ui\_editor\_tool\_editor\_planning
\item Question: Should this conversation-derived claim become a future review item, remain historical, or be deferred: The user uploaded screenshot bundles for setup and launcher. **UNCERTAIN / UNVERIFIED:** The screenshots were not inspected in this chat. The assistant still produced a logical...?
\item Why it matters: The claim frames unresolved future work or a decision boundary.
\item Current repo authority status: conversation-derived advisory claim; current repo authority not modified
\item Queue status: not directly blocked by queue keyword match
\item Recommended default: defer\_pending\_user\_decision
\item Decision owner: user\_decision
\item Risk if unresolved: Semantic drift or blocked-scope leakage if treated as current truth too early.
\item Evidence needed: Check canon, glossary, AGENTS.md, current queue, target current docs, and source conversation before any promotion.
\item Options:
\end{itemize}
  - promote\_later\_after\_authority\_review

  - preserve\_as\_history

  - defer

\begin{itemize}
\item Consequences:
\end{itemize}
  - promote\_later\_after\_authority\_review: creates a future docs-only or planning review task; still does not authorize implementation.

  - preserve\_as\_history: keeps the archive readable but removes pressure to patch current docs.

  - defer: leaves the claim unresolved until stronger authority or user decision exists.


\subsubsection{DECISION-0023}

\begin{itemize}
\item Source promotion/claim IDs: blocked\_scope:broad\_workbench\_ui
\item Source conversations: current queue plus conversation audit findings
\item Question: Should broad\_workbench\_ui remain blocked for all conversation-derived work until a future reviewed queue phase opens it?
\item Why it matters: The current queue marks broad\_workbench\_ui as BLOCKED and the corpus produced 7 related candidate/finding signals.
\item Current repo authority status: .aide/queue/current.toml is current repo authority for this scope gate.
\item Queue status: broad\_workbench\_ui: BLOCKED
\item Recommended default: keep\_blocked
\item Decision owner: future\_queue\_decision
\item Risk if unresolved: Implementation or live-doc scope could be opened by historical momentum rather than current authority.
\item Evidence needed: Review current queue, blocked-scope alignment, and any target prompt before changing disposition.
\item Options:
\end{itemize}
  - keep\_blocked

  - docs\_only\_crosswalk

  - defer

\begin{itemize}
\item Consequences:
\end{itemize}
  - keep\_blocked: prevents old conversation claims from opening implementation scope.

  - docs\_only\_crosswalk: allows explanatory archive/current-doc mapping without implementation or promotion.

  - defer: postpones judgment until the queue is revised.


\subsection{renderer/UI/platform}

\subsubsection{DECISION-0008}

\begin{itemize}
\item Source promotion/claim IDs: PROMOTE-0028
\item Source conversations: documentation\_standards\_readme
\item Question: Should this conversation-derived claim become a future review item, remain historical, or be deferred: The key thing to remember: this chat produced the standards and prompts for future work; it did not verify or change the project. The next assistant must inspect the actual repo...?
\item Why it matters: The claim frames unresolved future work or a decision boundary.
\item Current repo authority status: conversation-derived advisory claim; current repo authority not modified
\item Queue status: not directly blocked by queue keyword match
\item Recommended default: defer\_pending\_user\_decision
\item Decision owner: user\_decision
\item Risk if unresolved: Semantic drift or blocked-scope leakage if treated as current truth too early.
\item Evidence needed: Check canon, glossary, AGENTS.md, current queue, target current docs, and source conversation before any promotion.
\item Options:
\end{itemize}
  - promote\_later\_after\_authority\_review

  - preserve\_as\_history

  - defer

\begin{itemize}
\item Consequences:
\end{itemize}
  - promote\_later\_after\_authority\_review: creates a future docs-only or planning review task; still does not authorize implementation.

  - preserve\_as\_history: keeps the archive readable but removes pressure to patch current docs.

  - defer: leaves the claim unresolved until stronger authority or user decision exists.


\subsubsection{DECISION-0025}

\begin{itemize}
\item Source promotion/claim IDs: blocked\_scope:native\_gui
\item Source conversations: current queue plus conversation audit findings
\item Question: Should native\_gui remain blocked for all conversation-derived work until a future reviewed queue phase opens it?
\item Why it matters: The current queue marks native\_gui as BLOCKED and the corpus produced 14 related candidate/finding signals.
\item Current repo authority status: .aide/queue/current.toml is current repo authority for this scope gate.
\item Queue status: native\_gui: BLOCKED
\item Recommended default: keep\_blocked
\item Decision owner: future\_queue\_decision
\item Risk if unresolved: Implementation or live-doc scope could be opened by historical momentum rather than current authority.
\item Evidence needed: Review current queue, blocked-scope alignment, and any target prompt before changing disposition.
\item Options:
\end{itemize}
  - keep\_blocked

  - docs\_only\_crosswalk

  - defer

\begin{itemize}
\item Consequences:
\end{itemize}
  - keep\_blocked: prevents old conversation claims from opening implementation scope.

  - docs\_only\_crosswalk: allows explanatory archive/current-doc mapping without implementation or promotion.

  - defer: postpones judgment until the queue is revised.


\subsubsection{DECISION-0029}

\begin{itemize}
\item Source promotion/claim IDs: blocked\_scope:renderer\_implementation
\item Source conversations: current queue plus conversation audit findings
\item Question: Should renderer\_implementation remain blocked for all conversation-derived work until a future reviewed queue phase opens it?
\item Why it matters: The current queue marks renderer\_implementation as BLOCKED and the corpus produced 18 related candidate/finding signals.
\item Current repo authority status: .aide/queue/current.toml is current repo authority for this scope gate.
\item Queue status: renderer\_implementation: BLOCKED
\item Recommended default: keep\_blocked
\item Decision owner: future\_queue\_decision
\item Risk if unresolved: Implementation or live-doc scope could be opened by historical momentum rather than current authority.
\item Evidence needed: Review current queue, blocked-scope alignment, and any target prompt before changing disposition.
\item Options:
\end{itemize}
  - keep\_blocked

  - docs\_only\_crosswalk

  - defer

\begin{itemize}
\item Consequences:
\end{itemize}
  - keep\_blocked: prevents old conversation claims from opening implementation scope.

  - docs\_only\_crosswalk: allows explanatory archive/current-doc mapping without implementation or promotion.

  - defer: postpones judgment until the queue is revised.


\subsection{provider/content/packs}

\subsubsection{DECISION-0005}

\begin{itemize}
\item Source promotion/claim IDs: PROMOTE-0020
\item Source conversations: canonical\_structure\_and\_framework
\item Question: Should this conversation-derived claim remain archive-only until provider\_runtime is explicitly opened?
\item Why it matters: The claim touches currently blocked scope: provider\_runtime.
\item Current repo authority status: conversation-derived advisory claim; current repo authority not modified
\item Queue status: provider\_runtime
\item Recommended default: defer\_until\_queue\_opens
\item Decision owner: future\_queue\_decision
\item Risk if unresolved: Semantic drift or blocked-scope leakage if treated as current truth too early.
\item Evidence needed: Check canon, glossary, AGENTS.md, current queue, target current docs, and source conversation before any promotion.
\item Options:
\end{itemize}
  - promote\_later\_after\_authority\_review

  - preserve\_as\_history

  - defer

\begin{itemize}
\item Consequences:
\end{itemize}
  - promote\_later\_after\_authority\_review: creates a future docs-only or planning review task; still does not authorize implementation.

  - preserve\_as\_history: keeps the archive readable but removes pressure to patch current docs.

  - defer: leaves the claim unresolved until stronger authority or user decision exists.


\subsubsection{DECISION-0017}

\begin{itemize}
\item Source promotion/claim IDs: PROMOTE-0097
\item Source conversations: Portability\_Assurance\_Future\_Proof
\item Question: Should this conversation-derived claim become a future review item, remain historical, or be deferred: The actual repo was not inspected. The exact accepted directory structure, API policy, DDAP profile, compatibility promises, and first pilot module remain unresolved.?
\item Why it matters: The claim frames unresolved future work or a decision boundary.
\item Current repo authority status: conversation-derived advisory claim; current repo authority not modified
\item Queue status: not directly blocked by queue keyword match
\item Recommended default: defer\_pending\_user\_decision
\item Decision owner: user\_decision
\item Risk if unresolved: Semantic drift or blocked-scope leakage if treated as current truth too early.
\item Evidence needed: Check canon, glossary, AGENTS.md, current queue, target current docs, and source conversation before any promotion.
\item Options:
\end{itemize}
  - promote\_later\_after\_authority\_review

  - preserve\_as\_history

  - defer

\begin{itemize}
\item Consequences:
\end{itemize}
  - promote\_later\_after\_authority\_review: creates a future docs-only or planning review task; still does not authorize implementation.

  - preserve\_as\_history: keeps the archive readable but removes pressure to patch current docs.

  - defer: leaves the claim unresolved until stronger authority or user decision exists.


\subsubsection{DECISION-0026}

\begin{itemize}
\item Source promotion/claim IDs: blocked\_scope:package\_runtime
\item Source conversations: current queue plus conversation audit findings
\item Question: Should package\_runtime remain blocked for all conversation-derived work until a future reviewed queue phase opens it?
\item Why it matters: The current queue marks package\_runtime as BLOCKED and the corpus produced 4 related candidate/finding signals.
\item Current repo authority status: .aide/queue/current.toml is current repo authority for this scope gate.
\item Queue status: package\_runtime: BLOCKED
\item Recommended default: keep\_blocked
\item Decision owner: future\_queue\_decision
\item Risk if unresolved: Implementation or live-doc scope could be opened by historical momentum rather than current authority.
\item Evidence needed: Review current queue, blocked-scope alignment, and any target prompt before changing disposition.
\item Options:
\end{itemize}
  - keep\_blocked

  - docs\_only\_crosswalk

  - defer

\begin{itemize}
\item Consequences:
\end{itemize}
  - keep\_blocked: prevents old conversation claims from opening implementation scope.

  - docs\_only\_crosswalk: allows explanatory archive/current-doc mapping without implementation or promotion.

  - defer: postpones judgment until the queue is revised.


\subsubsection{DECISION-0027}

\begin{itemize}
\item Source promotion/claim IDs: blocked\_scope:provider\_runtime
\item Source conversations: current queue plus conversation audit findings
\item Question: Should provider\_runtime remain blocked for all conversation-derived work until a future reviewed queue phase opens it?
\item Why it matters: The current queue marks provider\_runtime as BLOCKED and the corpus produced 21 related candidate/finding signals.
\item Current repo authority status: .aide/queue/current.toml is current repo authority for this scope gate.
\item Queue status: provider\_runtime: BLOCKED
\item Recommended default: keep\_blocked
\item Decision owner: future\_queue\_decision
\item Risk if unresolved: Implementation or live-doc scope could be opened by historical momentum rather than current authority.
\item Evidence needed: Review current queue, blocked-scope alignment, and any target prompt before changing disposition.
\item Options:
\end{itemize}
  - keep\_blocked

  - docs\_only\_crosswalk

  - defer

\begin{itemize}
\item Consequences:
\end{itemize}
  - keep\_blocked: prevents old conversation claims from opening implementation scope.

  - docs\_only\_crosswalk: allows explanatory archive/current-doc mapping without implementation or promotion.

  - defer: postpones judgment until the queue is revised.


\subsubsection{DECISION-0030}

\begin{itemize}
\item Source promotion/claim IDs: blocked\_scope:runtime\_module\_loader
\item Source conversations: current queue plus conversation audit findings
\item Question: Should runtime\_module\_loader remain blocked for all conversation-derived work until a future reviewed queue phase opens it?
\item Why it matters: The current queue marks runtime\_module\_loader as BLOCKED and the corpus produced 5 related candidate/finding signals.
\item Current repo authority status: .aide/queue/current.toml is current repo authority for this scope gate.
\item Queue status: runtime\_module\_loader: BLOCKED
\item Recommended default: keep\_blocked
\item Decision owner: future\_queue\_decision
\item Risk if unresolved: Implementation or live-doc scope could be opened by historical momentum rather than current authority.
\item Evidence needed: Review current queue, blocked-scope alignment, and any target prompt before changing disposition.
\item Options:
\end{itemize}
  - keep\_blocked

  - docs\_only\_crosswalk

  - defer

\begin{itemize}
\item Consequences:
\end{itemize}
  - keep\_blocked: prevents old conversation claims from opening implementation scope.

  - docs\_only\_crosswalk: allows explanatory archive/current-doc mapping without implementation or promotion.

  - defer: postpones judgment until the queue is revised.


\subsection{world/time/civilization simulation}

\subsubsection{DECISION-0006}

\begin{itemize}
\item Source promotion/claim IDs: PROMOTE-0024
\item Source conversations: Chronology\_Celestial\_Systems
\item Question: Should this conversation-derived claim become a future review item, remain historical, or be deferred: The assistant formalized this as the Perfect Earth Calendar, later called HPC-E. The final structure is clear, though the exact leap rule remains unresolved.?
\item Why it matters: The claim frames unresolved future work or a decision boundary.
\item Current repo authority status: conversation-derived advisory claim; current repo authority not modified
\item Queue status: not directly blocked by queue keyword match
\item Recommended default: defer\_pending\_user\_decision
\item Decision owner: user\_decision
\item Risk if unresolved: Semantic drift or blocked-scope leakage if treated as current truth too early.
\item Evidence needed: Check canon, glossary, AGENTS.md, current queue, target current docs, and source conversation before any promotion.
\item Options:
\end{itemize}
  - promote\_later\_after\_authority\_review

  - preserve\_as\_history

  - defer

\begin{itemize}
\item Consequences:
\end{itemize}
  - promote\_later\_after\_authority\_review: creates a future docs-only or planning review task; still does not authorize implementation.

  - preserve\_as\_history: keeps the archive readable but removes pressure to patch current docs.

  - defer: leaves the claim unresolved until stronger authority or user decision exists.


\subsubsection{DECISION-0022}

\begin{itemize}
\item Source promotion/claim IDs: PROMOTE-0125
\item Source conversations: Universe\_Explorer\_Planning
\item Question: Should this conversation-derived claim become a future review item, remain historical, or be deferred: A major theme was that useful local inventions must become portable, standardized, industrialized, and institutionally adopted. The repo now has a Formalization Chain spec and P...?
\item Why it matters: The claim frames unresolved future work or a decision boundary.
\item Current repo authority status: conversation-derived advisory claim; current repo authority not modified
\item Queue status: not directly blocked by queue keyword match
\item Recommended default: defer\_pending\_user\_decision
\item Decision owner: user\_decision
\item Risk if unresolved: Semantic drift or blocked-scope leakage if treated as current truth too early.
\item Evidence needed: Check canon, glossary, AGENTS.md, current queue, target current docs, and source conversation before any promotion.
\item Options:
\end{itemize}
  - promote\_later\_after\_authority\_review

  - preserve\_as\_history

  - defer

\begin{itemize}
\item Consequences:
\end{itemize}
  - promote\_later\_after\_authority\_review: creates a future docs-only or planning review task; still does not authorize implementation.

  - preserve\_as\_history: keeps the archive readable but removes pressure to patch current docs.

  - defer: leaves the claim unresolved until stronger authority or user decision exists.


\subsection{release/setup/launcher}

\subsubsection{DECISION-0016}

\begin{itemize}
\item Source promotion/claim IDs: PROMOTE-0076
\item Source conversations: launcher\_app\_layer
\item Question: Should this conversation-derived claim become a future review item, remain historical, or be deferred: Someone reading this report should understand one central thing: this chat is not about inventing new launcher features anymore. It is about preserving boundaries, making the la...?
\item Why it matters: The claim frames unresolved future work or a decision boundary.
\item Current repo authority status: conversation-derived advisory claim; current repo authority not modified
\item Queue status: not directly blocked by queue keyword match
\item Recommended default: defer\_pending\_user\_decision
\item Decision owner: user\_decision
\item Risk if unresolved: Semantic drift or blocked-scope leakage if treated as current truth too early.
\item Evidence needed: Check canon, glossary, AGENTS.md, current queue, target current docs, and source conversation before any promotion.
\item Options:
\end{itemize}
  - promote\_later\_after\_authority\_review

  - preserve\_as\_history

  - defer

\begin{itemize}
\item Consequences:
\end{itemize}
  - promote\_later\_after\_authority\_review: creates a future docs-only or planning review task; still does not authorize implementation.

  - preserve\_as\_history: keeps the archive readable but removes pressure to patch current docs.

  - defer: leaves the claim unresolved until stronger authority or user decision exists.


\subsubsection{DECISION-0028}

\begin{itemize}
\item Source promotion/claim IDs: blocked\_scope:release\_publication
\item Source conversations: current queue plus conversation audit findings
\item Question: Should release\_publication remain blocked for all conversation-derived work until a future reviewed queue phase opens it?
\item Why it matters: The current queue marks release\_publication as BLOCKED and the corpus produced 8 related candidate/finding signals.
\item Current repo authority status: .aide/queue/current.toml is current repo authority for this scope gate.
\item Queue status: release\_publication: BLOCKED
\item Recommended default: keep\_blocked
\item Decision owner: future\_queue\_decision
\item Risk if unresolved: Implementation or live-doc scope could be opened by historical momentum rather than current authority.
\item Evidence needed: Review current queue, blocked-scope alignment, and any target prompt before changing disposition.
\item Options:
\end{itemize}
  - keep\_blocked

  - docs\_only\_crosswalk

  - defer

\begin{itemize}
\item Consequences:
\end{itemize}
  - keep\_blocked: prevents old conversation claims from opening implementation scope.

  - docs\_only\_crosswalk: allows explanatory archive/current-doc mapping without implementation or promotion.

  - defer: postpones judgment until the queue is revised.




{\small\raggedright SOURCE PATH: docs/\allowbreak{}archive/\allowbreak{}conversations/\allowbreak{}\_\allowbreak{}decision/\allowbreak{}DECISION\_\allowbreak{}OPTIONS\_\allowbreak{}MATRIX\_\allowbreak{}v0.md\par}


\section{Decision Options Matrix v0}

Matrix view of available decision options. Tables support review but do not replace the docket explanations.



\begin{quote}\small Dense table summarized for the reader edition. See the source file, HTML output, or reference appendix source for full detail.\end{quote}

\begin{quote}\small Dense table content is summarized in this PDF rendering. See the Markdown source and HTML book for the full table.\end{quote}



{\small\raggedright SOURCE PATH: docs/\allowbreak{}archive/\allowbreak{}conversations/\allowbreak{}\_\allowbreak{}decision/\allowbreak{}DEFERRED\_\allowbreak{}DECISIONS\_\allowbreak{}v0.md\par}


\section{Deferred Decisions v0}

These decisions should stay deferred unless a human or stronger repo authority opens them.


\subsection{DECISION-0001}

\begin{itemize}
\item Group: architecture/boundaries
\item Question: What disposition should be chosen for this unresolved claim: The user then asked whether arbitrary placement was possible or whether the system was stuck to grids. The answer established that the engine should be grid-agnostic. [FACT] The...?
\item Recommended default: defer\_pending\_user\_decision
\item Queue status: not directly blocked by queue keyword match
\item Evidence that would resolve it: Check canon, glossary, AGENTS.md, current queue, target current docs, and source conversation before any promotion.
\end{itemize}

\subsection{DECISION-0002}

\begin{itemize}
\item Group: architecture/boundaries
\item Question: Should this conversation-derived claim remain archive-only until gameplay is explicitly opened?
\item Recommended default: defer\_until\_queue\_opens
\item Queue status: gameplay
\item Evidence that would resolve it: Check canon, glossary, AGENTS.md, current queue, target current docs, and source conversation before any promotion.
\end{itemize}

\subsection{DECISION-0003}

\begin{itemize}
\item Group: architecture/boundaries
\item Question: What disposition should be chosen for this unresolved claim: Future work must define the client's allowed dependency surface and whether it can access any simulation-adjacent prediction layer.?
\item Recommended default: defer\_pending\_user\_decision
\item Queue status: not directly blocked by queue keyword match
\item Evidence that would resolve it: Check canon, glossary, AGENTS.md, current queue, target current docs, and source conversation before any promotion.
\end{itemize}

\subsection{DECISION-0004}

\begin{itemize}
\item Group: architecture/boundaries
\item Question: Should this conversation-derived claim remain archive-only until provider\_runtime is explicitly opened?
\item Recommended default: defer\_until\_queue\_opens
\item Queue status: provider\_runtime
\item Evidence that would resolve it: Check canon, glossary, AGENTS.md, current queue, target current docs, and source conversation before any promotion.
\end{itemize}

\subsection{DECISION-0005}

\begin{itemize}
\item Group: provider/content/packs
\item Question: Should this conversation-derived claim remain archive-only until provider\_runtime is explicitly opened?
\item Recommended default: defer\_until\_queue\_opens
\item Queue status: provider\_runtime
\item Evidence that would resolve it: Check canon, glossary, AGENTS.md, current queue, target current docs, and source conversation before any promotion.
\end{itemize}

\subsection{DECISION-0006}

\begin{itemize}
\item Group: world/time/civilization simulation
\item Question: Should this conversation-derived claim become a future review item, remain historical, or be deferred: The assistant formalized this as the Perfect Earth Calendar, later called HPC-E. The final structure is clear, though the exact leap rule remains unresolved.?
\item Recommended default: defer\_pending\_user\_decision
\item Queue status: not directly blocked by queue keyword match
\item Evidence that would resolve it: Check canon, glossary, AGENTS.md, current queue, target current docs, and source conversation before any promotion.
\end{itemize}

\subsection{DECISION-0007}

\begin{itemize}
\item Group: architecture/boundaries
\item Question: What disposition should be chosen for this unresolved claim: Before relying on this chat as project authority, a future assistant or human reviewer must confirm whether Dominium is actually intended to be simulation-heavy, whether strict...?
\item Recommended default: defer\_pending\_user\_decision
\item Queue status: not directly blocked by queue keyword match
\item Evidence that would resolve it: Check canon, glossary, AGENTS.md, current queue, target current docs, and source conversation before any promotion.
\end{itemize}

\subsection{DECISION-0008}

\begin{itemize}
\item Group: renderer/UI/platform
\item Question: Should this conversation-derived claim become a future review item, remain historical, or be deferred: The key thing to remember: this chat produced the standards and prompts for future work; it did not verify or change the project. The next assistant must inspect the actual repo...?
\item Recommended default: defer\_pending\_user\_decision
\item Queue status: not directly blocked by queue keyword match
\item Evidence that would resolve it: Check canon, glossary, AGENTS.md, current queue, target current docs, and source conversation before any promotion.
\end{itemize}

\subsection{DECISION-0009}

\begin{itemize}
\item Group: architecture/boundaries
\item Question: What disposition should be chosen for this unresolved claim: This chat was about preserving and re-grounding Dominium's architecture around a very old constitutional contract, then checking whether the current GitHub repository still foll...?
\item Recommended default: defer\_pending\_user\_decision
\item Queue status: not directly blocked by queue keyword match
\item Evidence that would resolve it: Check canon, glossary, AGENTS.md, current queue, target current docs, and source conversation before any promotion.
\end{itemize}

\subsection{DECISION-0010}

\begin{itemize}
\item Group: architecture/boundaries
\item Question: Should this conversation-derived claim remain archive-only until provider\_runtime is explicitly opened?
\item Recommended default: defer\_until\_queue\_opens
\item Queue status: provider\_runtime
\item Evidence that would resolve it: Check canon, glossary, AGENTS.md, current queue, target current docs, and source conversation before any promotion.
\end{itemize}

\subsection{DECISION-0011}

\begin{itemize}
\item Group: Workbench/AIDE/Codex/tooling
\item Question: Should this conversation-derived claim remain archive-only until provider\_runtime is explicitly opened?
\item Recommended default: defer\_until\_queue\_opens
\item Queue status: provider\_runtime
\item Evidence that would resolve it: Check canon, glossary, AGENTS.md, current queue, target current docs, and source conversation before any promotion.
\end{itemize}

\subsection{DECISION-0012}

\begin{itemize}
\item Group: Workbench/AIDE/Codex/tooling
\item Question: Should this conversation-derived claim remain archive-only until native\_gui is explicitly opened?
\item Recommended default: defer\_until\_queue\_opens
\item Queue status: native\_gui
\item Evidence that would resolve it: Check canon, glossary, AGENTS.md, current queue, target current docs, and source conversation before any promotion.
\end{itemize}

\subsection{DECISION-0013}

\begin{itemize}
\item Group: Workbench/AIDE/Codex/tooling
\item Question: Should this conversation-derived claim remain archive-only until provider\_runtime is explicitly opened?
\item Recommended default: defer\_until\_queue\_opens
\item Queue status: provider\_runtime
\item Evidence that would resolve it: Check canon, glossary, AGENTS.md, current queue, target current docs, and source conversation before any promotion.
\end{itemize}

\subsection{DECISION-0014}

\begin{itemize}
\item Group: architecture/boundaries
\item Question: Should this conversation-derived claim remain archive-only until provider\_runtime is explicitly opened?
\item Recommended default: defer\_until\_queue\_opens
\item Queue status: provider\_runtime
\item Evidence that would resolve it: Check canon, glossary, AGENTS.md, current queue, target current docs, and source conversation before any promotion.
\end{itemize}

\subsection{DECISION-0015}

\begin{itemize}
\item Group: architecture/boundaries
\item Question: Should this conversation-derived claim remain archive-only until provider\_runtime is explicitly opened?
\item Recommended default: defer\_until\_queue\_opens
\item Queue status: provider\_runtime
\item Evidence that would resolve it: Check canon, glossary, AGENTS.md, current queue, target current docs, and source conversation before any promotion.
\end{itemize}

\subsection{DECISION-0016}

\begin{itemize}
\item Group: release/setup/launcher
\item Question: Should this conversation-derived claim become a future review item, remain historical, or be deferred: Someone reading this report should understand one central thing: this chat is not about inventing new launcher features anymore. It is about preserving boundaries, making the la...?
\item Recommended default: defer\_pending\_user\_decision
\item Queue status: not directly blocked by queue keyword match
\item Evidence that would resolve it: Check canon, glossary, AGENTS.md, current queue, target current docs, and source conversation before any promotion.
\end{itemize}

\subsection{DECISION-0017}

\begin{itemize}
\item Group: provider/content/packs
\item Question: Should this conversation-derived claim become a future review item, remain historical, or be deferred: The actual repo was not inspected. The exact accepted directory structure, API policy, DDAP profile, compatibility promises, and first pilot module remain unresolved.?
\item Recommended default: defer\_pending\_user\_decision
\item Queue status: not directly blocked by queue keyword match
\item Evidence that would resolve it: Check canon, glossary, AGENTS.md, current queue, target current docs, and source conversation before any promotion.
\end{itemize}

\subsection{DECISION-0018}

\begin{itemize}
\item Group: architecture/boundaries
\item Question: What disposition should be chosen for this unresolved claim: The user then asked whether the implementation was industry-accepted and what could be improved. The response framed the approach as closer to game engines, operating systems, a...?
\item Recommended default: defer\_pending\_user\_decision
\item Queue status: not directly blocked by queue keyword match
\item Evidence that would resolve it: Check canon, glossary, AGENTS.md, current queue, target current docs, and source conversation before any promotion.
\end{itemize}

\subsection{DECISION-0019}

\begin{itemize}
\item Group: architecture/boundaries
\item Question: Should this conversation-derived claim remain archive-only until gameplay is explicitly opened?
\item Recommended default: defer\_until\_queue\_opens
\item Queue status: gameplay
\item Evidence that would resolve it: Check canon, glossary, AGENTS.md, current queue, target current docs, and source conversation before any promotion.
\end{itemize}

\subsection{DECISION-0020}

\begin{itemize}
\item Group: Workbench/AIDE/Codex/tooling
\item Question: What disposition should be chosen for this unresolved claim: The most important thing to remember is that this chat was a planning and prompt-generation chat, not proof of implementation. **UNCERTAIN / UNVERIFIED:** No generated prompt is...?
\item Recommended default: defer\_pending\_user\_decision
\item Queue status: not directly blocked by queue keyword match
\item Evidence that would resolve it: Check canon, glossary, AGENTS.md, current queue, target current docs, and source conversation before any promotion.
\end{itemize}

\subsection{DECISION-0021}

\begin{itemize}
\item Group: Workbench/AIDE/Codex/tooling
\item Question: Should this conversation-derived claim become a future review item, remain historical, or be deferred: The user uploaded screenshot bundles for setup and launcher. **UNCERTAIN / UNVERIFIED:** The screenshots were not inspected in this chat. The assistant still produced a logical...?
\item Recommended default: defer\_pending\_user\_decision
\item Queue status: not directly blocked by queue keyword match
\item Evidence that would resolve it: Check canon, glossary, AGENTS.md, current queue, target current docs, and source conversation before any promotion.
\end{itemize}

\subsection{DECISION-0022}

\begin{itemize}
\item Group: world/time/civilization simulation
\item Question: Should this conversation-derived claim become a future review item, remain historical, or be deferred: A major theme was that useful local inventions must become portable, standardized, industrialized, and institutionally adopted. The repo now has a Formalization Chain spec and P...?
\item Recommended default: defer\_pending\_user\_decision
\item Queue status: not directly blocked by queue keyword match
\item Evidence that would resolve it: Check canon, glossary, AGENTS.md, current queue, target current docs, and source conversation before any promotion.
\end{itemize}

\subsection{DECISION-0023}

\begin{itemize}
\item Group: Workbench/AIDE/Codex/tooling
\item Question: Should broad\_workbench\_ui remain blocked for all conversation-derived work until a future reviewed queue phase opens it?
\item Recommended default: keep\_blocked
\item Queue status: broad\_workbench\_ui: BLOCKED
\item Evidence that would resolve it: Review current queue, blocked-scope alignment, and any target prompt before changing disposition.
\end{itemize}

\subsection{DECISION-0024}

\begin{itemize}
\item Group: architecture/boundaries
\item Question: Should gameplay remain blocked for all conversation-derived work until a future reviewed queue phase opens it?
\item Recommended default: keep\_blocked
\item Queue status: gameplay: BLOCKED
\item Evidence that would resolve it: Review current queue, blocked-scope alignment, and any target prompt before changing disposition.
\end{itemize}

\subsection{DECISION-0025}

\begin{itemize}
\item Group: renderer/UI/platform
\item Question: Should native\_gui remain blocked for all conversation-derived work until a future reviewed queue phase opens it?
\item Recommended default: keep\_blocked
\item Queue status: native\_gui: BLOCKED
\item Evidence that would resolve it: Review current queue, blocked-scope alignment, and any target prompt before changing disposition.
\end{itemize}

\subsection{DECISION-0026}

\begin{itemize}
\item Group: provider/content/packs
\item Question: Should package\_runtime remain blocked for all conversation-derived work until a future reviewed queue phase opens it?
\item Recommended default: keep\_blocked
\item Queue status: package\_runtime: BLOCKED
\item Evidence that would resolve it: Review current queue, blocked-scope alignment, and any target prompt before changing disposition.
\end{itemize}

\subsection{DECISION-0027}

\begin{itemize}
\item Group: provider/content/packs
\item Question: Should provider\_runtime remain blocked for all conversation-derived work until a future reviewed queue phase opens it?
\item Recommended default: keep\_blocked
\item Queue status: provider\_runtime: BLOCKED
\item Evidence that would resolve it: Review current queue, blocked-scope alignment, and any target prompt before changing disposition.
\end{itemize}

\subsection{DECISION-0028}

\begin{itemize}
\item Group: release/setup/launcher
\item Question: Should release\_publication remain blocked for all conversation-derived work until a future reviewed queue phase opens it?
\item Recommended default: keep\_blocked
\item Queue status: release\_publication: BLOCKED
\item Evidence that would resolve it: Review current queue, blocked-scope alignment, and any target prompt before changing disposition.
\end{itemize}

\subsection{DECISION-0029}

\begin{itemize}
\item Group: renderer/UI/platform
\item Question: Should renderer\_implementation remain blocked for all conversation-derived work until a future reviewed queue phase opens it?
\item Recommended default: keep\_blocked
\item Queue status: renderer\_implementation: BLOCKED
\item Evidence that would resolve it: Review current queue, blocked-scope alignment, and any target prompt before changing disposition.
\end{itemize}

\subsection{DECISION-0030}

\begin{itemize}
\item Group: provider/content/packs
\item Question: Should runtime\_module\_loader remain blocked for all conversation-derived work until a future reviewed queue phase opens it?
\item Recommended default: keep\_blocked
\item Queue status: runtime\_module\_loader: BLOCKED
\item Evidence that would resolve it: Review current queue, blocked-scope alignment, and any target prompt before changing disposition.
\end{itemize}


\chapter{Promotion Review}


{\small\raggedright SOURCE PATH: docs/\allowbreak{}archive/\allowbreak{}conversations/\allowbreak{}\_\allowbreak{}promotion/\allowbreak{}PROMOTION\_\allowbreak{}WAVE\_\allowbreak{}1\_\allowbreak{}CANDIDATES\_\allowbreak{}v0.md\par}


\section{Promotion Wave 1 Candidates v0}

Wave 1 candidates are docs-only review candidates that appear least blocked by current queue. They still require separate microtasks.


\subsection{PROMOTE-0003 - architecture\_doc\_candidate}

\begin{itemize}
\item Source conversation: advanced\_simulation\_infrastructure
\item Proposed target: docs/architecture/** after target selection
\item Validation required: source check, authority-order check, link check, generated output validator, FAST
\item Recommended microtask: DOC-PROMOTION-PROMOTE-0003
\item Claim: [INFERENCE] The conclusion was not a finalized implementation but a working architecture: each simulation domain should be deterministic, fixed-point, content-driven, versioned, replayable, and integrated through clear read/write boundaries.
\end{itemize}

\subsection{PROMOTE-0004 - architecture\_doc\_candidate}

\begin{itemize}
\item Source conversation: app\_runtime\_platform\_renderers
\item Proposed target: docs/architecture/** after target selection
\item Validation required: source check, authority-order check, link check, generated output validator, FAST
\item Recommended microtask: DOC-PROMOTION-PROMOTE-0004
\item Claim: The central topic was APP0: the application/runtime/platform/renderer layer of Dominium. This topic came up because the user provided a formal prompt defining what APP0 should do.
\end{itemize}

\subsection{PROMOTE-0032 - architecture\_doc\_candidate}

\begin{itemize}
\item Source conversation: Dominium\_Architecture\_I
\item Proposed target: docs/architecture/** after target selection
\item Validation required: source check, authority-order check, link check, generated output validator, FAST
\item Recommended microtask: DOC-PROMOTION-PROMOTE-0032
\item Claim: The user also made a clear versioning decision. **FACT:** The user wanted **major.minor.patch semantic versioning** for all components/packages and for the complete game package. The user explicitly did **not** want build numbers or build dates in filenames. Instead, those details should live in metadata. This became part of the future build/package design.
\end{itemize}

\subsection{PROMOTE-0035 - architecture\_doc\_candidate}

\begin{itemize}
\item Source conversation: Dominium\_Architecture\_II
\item Proposed target: docs/architecture/** after target selection
\item Validation required: source check, authority-order check, link check, generated output validator, FAST
\item Recommended microtask: DOC-PROMOTION-PROMOTE-0035
\item Claim: The entire conversation is anchored by the requirement that Dominium must run deterministically. This means the same initial state and same inputs must produce the same state across machines, compilers, operating systems, and eventually retro and modern platforms.
\end{itemize}

\subsection{PROMOTE-0038 - architecture\_doc\_candidate}

\begin{itemize}
\item Source conversation: Dominium\_Architecture\_III
\item Proposed target: docs/architecture/** after target selection
\item Validation required: source check, authority-order check, link check, generated output validator, FAST
\item Recommended microtask: DOC-PROMOTION-PROMOTE-0038
\item Claim: This context matters because the rest of the conversation happened inside those constraints. The launcher could not become a dumping ground for OS-specific behavior, sim mutation, renderer decisions, or nondeterministic game-state changes. It had to fit the project's deterministic layering.
\end{itemize}

\subsection{PROMOTE-0050 - architecture\_doc\_candidate}

\begin{itemize}
\item Source conversation: dominium\_full\_conversation
\item Proposed target: docs/architecture/** after target selection
\item Validation required: source check, authority-order check, link check, generated output validator, FAST
\item Recommended microtask: DOC-PROMOTION-PROMOTE-0050
\item Claim: The user then redirected the entire plan. They argued that native OS widget GUI tools already exist through Visual Studio, Xcode, etc. The project needed a cross-platform rendered tool environment that uses the same CLI, TUI, and rendered GUI systems as the client. The old UI Editor and Tool Editor were good exploratory ideas but bad final products. This produced the Dominium Workbench concept.
\end{itemize}

\subsection{PROMOTE-0054 - planning\_doc\_candidate}

\begin{itemize}
\item Source conversation: dominium\_setup
\item Proposed target: docs/planning/** after target selection
\item Validation required: source check, authority-order check, link check, generated output validator, FAST
\item Recommended microtask: DOC-PROMOTION-PROMOTE-0054
\item Claim: The user later asked where to create Visual Studio and Xcode apps after opening the repo as a folder. The assistant advised that Visual Studio and Xcode projects should be generated through CMake, not hand-authored or treated as authoritative source files. This decision became consistent with later Codex prompts.
\end{itemize}

\subsection{PROMOTE-0063 - architecture\_doc\_candidate}

\begin{itemize}
\item Source conversation: engine\_baseline\_architecture
\item Proposed target: docs/architecture/** after target selection
\item Validation required: source check, authority-order check, link check, generated output validator, FAST
\item Recommended microtask: DOC-PROMOTION-PROMOTE-0063
\item Claim: The central architecture distinction was that Domino should be a reusable deterministic simulation substrate, while Dominium should be one game/product layer built on it. This came up immediately when the user asked for a total description of the project. It mattered because the user wants the code to be reusable for future games and possibly other engine projects.
\end{itemize}

\subsection{PROMOTE-0073 - architecture\_doc\_candidate}

\begin{itemize}
\item Source conversation: Language\_Platform\_Architecture
\item Proposed target: docs/architecture/** after target selection
\item Validation required: source check, authority-order check, link check, generated output validator, FAST
\item Recommended microtask: DOC-PROMOTION-PROMOTE-0073
\item Claim: The 32-bit question followed. The answer was to make full native products 64-bit and keep 32-bit as constrained or projection support only. The key caveat was not to make native pointer size part of game law. Save, replay, network, IDs, and renderer IR must use fixed-width types and never serialize size\_t, long, uintptr\_t, or raw pointers.
\end{itemize}

\subsection{PROMOTE-0084 - architecture\_doc\_candidate}

\begin{itemize}
\item Source conversation: Modularity\_AIDE\_Refactorability
\item Proposed target: docs/architecture/** after target selection
\item Validation required: source check, authority-order check, link check, generated output validator, FAST
\item Recommended microtask: DOC-PROMOTION-PROMOTE-0084
\item Claim: The discussion proposed target-based CMake and module boundaries rather than path mythology. Public headers, private headers, allowed dependencies, and forbidden dependencies should be explicit. Apps must remain thin. Engine must not depend on game/apps/runtime UI. Runtime adapts host/platform/rendering without owning simulation truth.
\end{itemize}

\subsection{PROMOTE-0088 - architecture\_doc\_candidate}

\begin{itemize}
\item Source conversation: OS\_Interface\_Repo\_Architecture
\item Proposed target: docs/architecture/** after target selection
\item Validation required: source check, authority-order check, link check, generated output validator, FAST
\item Recommended microtask: DOC-PROMOTION-PROMOTE-0088
\item Claim: DECISION-01 - CLI mandatory, TUI expected, GUI modular.** This was part of the initial architecture baseline and was repeatedly reaffirmed. It matters because every product must remain operable in recovery/headless/automation contexts. GUI is allowed but not authoritative.
\end{itemize}

\subsection{PROMOTE-0100 - docs\_clarification\_candidate}

\begin{itemize}
\item Source conversation: readme\_ports\_determinism
\item Proposed target: README.md or docs/** after target selection
\item Validation required: source check, authority-order check, link check, generated output validator, FAST
\item Recommended microtask: DOC-PROMOTION-PROMOTE-0100
\item Claim: The user pasted the final README after those cleanup changes. At that point, the README had the current active form: deterministic constraints clarified, ports unified under one source hierarchy, /ports metadata-only if retained, Section 9 normative, lockstep canonical, content-lock mismatch fatal, and disk format versions immutable.
\end{itemize}

\subsection{PROMOTE-0102 - architecture\_doc\_candidate}

\begin{itemize}
\item Source conversation: readme\_ports\_determinism
\item Proposed target: docs/architecture/** after target selection
\item Validation required: source check, authority-order check, link check, generated output validator, FAST
\item Recommended microtask: DOC-PROMOTION-PROMOTE-0102
\item Claim: The central artifact was the root README.md. The README describes **Domino** as the deterministic engine core and **Dominium** as the official game/tooling/runtime layer. It is written as a high-level architecture document, not a low-level implementation spec.
\end{itemize}

\subsection{PROMOTE-0105 - architecture\_doc\_candidate}

\begin{itemize}
\item Source conversation: Refactor\_Architecture
\item Proposed target: docs/architecture/** after target selection
\item Validation required: source check, authority-order check, link check, generated output validator, FAST
\item Recommended microtask: DOC-PROMOTION-PROMOTE-0105
\item Claim: FACT:** All products should use shared dsys and dgfx. No product-specific platform or renderer code path.
\end{itemize}

\subsection{PROMOTE-0112 - architecture\_doc\_candidate}

\begin{itemize}
\item Source conversation: testx\_repox\_governance
\item Proposed target: docs/architecture/** after target selection
\item Validation required: source check, authority-order check, link check, generated output validator, FAST
\item Recommended microtask: DOC-PROMOTION-PROMOTE-0112
\item Claim: FACT:** The user is building Dominium / Domino as a deterministic universe engine + game, not a conventional game project. The visible chat repeatedly framed the project as long-lived, simulation-first, deterministic, modular, and designed to survive across many operating systems, toolchains, renderers, products, and distribution models.
\end{itemize}

\subsection{PROMOTE-0114 - architecture\_doc\_candidate}

\begin{itemize}
\item Source conversation: testx\_repox\_governance
\item Proposed target: docs/architecture/** after target selection
\item Validation required: source check, authority-order check, link check, generated output validator, FAST
\item Recommended microtask: DOC-PROMOTION-PROMOTE-0114
\item Claim: The user had another chat working on content and systems, and asked for a prompt to inform that chat of everything decided so far. A similar prompt was generated for the applications/platforms/renderers chat. These prompts established authoritative boundaries:
\end{itemize}

\subsection{PROMOTE-0115 - architecture\_doc\_candidate}

\begin{itemize}
\item Source conversation: Timekeeping\_and\_2038\_Resilience
\item Proposed target: docs/architecture/** after target selection
\item Validation required: source check, authority-order check, link check, generated output validator, FAST
\item Recommended microtask: DOC-PROMOTION-PROMOTE-0115
\item Claim: A future assistant should understand that this chat contributes a time architecture doctrine for Dominium: **ACT is authority; DSYS time is runtime-only; observer clocks are derived; civil/astronomical time is projection-only; wall-clock time must never drive authoritative ordering.
\end{itemize}

\subsection{PROMOTE-0131 - architecture\_doc\_candidate}

\begin{itemize}
\item Source conversation: World\_Architecture
\item Proposed target: docs/architecture/** after target selection
\item Validation required: source check, authority-order check, link check, generated output validator, FAST
\item Recommended microtask: DOC-PROMOTION-PROMOTE-0131
\item Claim: Before reading the rest of this report, remember this: this chat's central contribution is not one isolated feature. It is a whole design philosophy for Dominium's world engine. The game should be deterministic, fixed-point, sparse, modular, content-driven, and field-based. Chunks are not the world. Chunks are just how the world is cached, meshed, streamed, and saved.
\end{itemize}



{\small\raggedright SOURCE PATH: docs/\allowbreak{}archive/\allowbreak{}conversations/\allowbreak{}\_\allowbreak{}promotion/\allowbreak{}PROMOTION\_\allowbreak{}REVIEW\_\allowbreak{}BOARD\_\allowbreak{}v0.md\par}


\section{Promotion Review Board v0}

This board converts raw promotion candidates into bounded review items. It does not patch or promote live docs.


\subsection{Counts}

\begin{itemize}
\item Total candidates: 135
\item architecture\_doc\_candidate: 16
\item archive\_only: 9
\item blocked\_by\_queue: 10
\item contract\_candidate: 11
\item docs\_clarification\_candidate: 1
\item insufficient\_evidence: 53
\item needs\_user\_decision: 12
\item planning\_doc\_candidate: 1
\item reject\_as\_noise: 19
\item schema\_candidate: 3
\end{itemize}

\subsection{Review Items}

\subsubsection{PROMOTE-0001 - needs\_user\_decision}

\begin{itemize}
\item Source conversation: advanced\_simulation\_infrastructure
\item Proposed target path: docs/archive/conversations/\_decision/DECISION\_DOCKET\_v0.md
\item Current authority support: not yet verified against target docs
\item Current authority conflict: none detected by triage
\item Queue risk: none\_detected
\item Validation required: preserve as archive evidence; no live target validation
\item Recommended next action: hold\_or\_review
\item Claim: The user then asked whether arbitrary placement was possible or whether the system was stuck to grids. The answer established that the engine should be grid-agnostic. [FACT] The user's question directly made arbitrary placement a major topic. [INFERENCE] The answer was accepted as the current design direction because the user then asked for a Codex prompt plan to implement changes necessary to achieve it.
\end{itemize}

\subsubsection{PROMOTE-0002 - reject\_as\_noise}

\begin{itemize}
\item Source conversation: advanced\_simulation\_infrastructure
\item Proposed target path: none
\item Current authority support: not yet verified against target docs
\item Current authority conflict: none detected by triage
\item Queue risk: none\_detected
\item Validation required: preserve as archive evidence; no live target validation
\item Recommended next action: hold\_or\_review
\item Claim: Finally, the chat shifted into archival mode: it produced a maximum-fidelity context transfer packet, then a downloadable report package, then an in-chat reader version. Those later outputs preserved the discussion for future aggregation.
\end{itemize}

\subsubsection{PROMOTE-0003 - architecture\_doc\_candidate}

\begin{itemize}
\item Source conversation: advanced\_simulation\_infrastructure
\item Proposed target path: docs/architecture/** after target selection
\item Current authority support: not yet verified against target docs
\item Current authority conflict: none detected by triage
\item Queue risk: none\_detected
\item Validation required: source check, authority-order check, link check, generated output validator, FAST
\item Recommended next action: create\_microtask
\item Claim: [INFERENCE] The conclusion was not a finalized implementation but a working architecture: each simulation domain should be deterministic, fixed-point, content-driven, versioned, replayable, and integrated through clear read/write boundaries.
\end{itemize}

\subsubsection{PROMOTE-0004 - architecture\_doc\_candidate}

\begin{itemize}
\item Source conversation: app\_runtime\_platform\_renderers
\item Proposed target path: docs/architecture/** after target selection
\item Current authority support: not yet verified against target docs
\item Current authority conflict: none detected by triage
\item Queue risk: none\_detected
\item Validation required: source check, authority-order check, link check, generated output validator, FAST
\item Recommended next action: create\_microtask
\item Claim: The central topic was APP0: the application/runtime/platform/renderer layer of Dominium. This topic came up because the user provided a formal prompt defining what APP0 should do.
\end{itemize}

\subsubsection{PROMOTE-0005 - blocked\_by\_queue}

\begin{itemize}
\item Source conversation: app\_runtime\_platform\_renderers
\item Proposed target path: none
\item Current authority support: not yet verified against target docs
\item Current authority conflict: blocked queue scope
\item Queue risk: blocked
\item Validation required: preserve as archive evidence; no live target validation
\item Recommended next action: hold\_or\_review
\item Claim: The most important conclusion was that APP0 must remain outside the core simulation rules. **FACT:** the user explicitly forbade redesigning simulation rules, altering content definitions, changing life/civilization/economy logic, and introducing gameplay shortcuts. **FACT:** authoritative logic remains in engine + game.
\end{itemize}

\subsubsection{PROMOTE-0006 - needs\_user\_decision}

\begin{itemize}
\item Source conversation: app\_runtime\_platform\_renderers
\item Proposed target path: docs/archive/conversations/\_decision/DECISION\_DOCKET\_v0.md
\item Current authority support: not yet verified against target docs
\item Current authority conflict: none detected by triage
\item Queue risk: none\_detected
\item Validation required: preserve as archive evidence; no live target validation
\item Recommended next action: hold\_or\_review
\item Claim: Future work must define the client's allowed dependency surface and whether it can access any simulation-adjacent prediction layer.
\end{itemize}

\subsubsection{PROMOTE-0007 - schema\_candidate}

\begin{itemize}
\item Source conversation: app\_testx\_codehygiene
\item Proposed target path: schema/** or schemas/** after authority review
\item Current authority support: not yet verified against target docs
\item Current authority conflict: none detected by triage
\item Queue risk: none\_detected
\item Validation required: authority review plus targeted schema/contract parser and FAST
\item Recommended next action: hold\_or\_review
\item Claim: A large set of prompts was then generated to lock down architecture, determinism, performance, schema governance, rendering, epistemic UI, sharding, interest sets, and fidelity projection. This became the Phase 1 hardening layer. Additional audit prompts were introduced to ensure consistency before proceeding into life, civilization, war, content, agents, tools, mods, and final long-term policy.
\end{itemize}

\subsubsection{PROMOTE-0008 - reject\_as\_noise}

\begin{itemize}
\item Source conversation: app\_testx\_codehygiene
\item Proposed target path: none
\item Current authority support: not yet verified against target docs
\item Current authority conflict: none detected by triage
\item Queue risk: none\_detected
\item Validation required: preserve as archive evidence; no live target validation
\item Recommended next action: hold\_or\_review
\item Claim: Because the chat was huge, the user asked for a maximum-fidelity context transfer packet. Then they asked for downloadable report files. Then they asked for an in-chat reader. Finally, they asked for this human-readable explanatory report. This final report is meant to let a future assistant or human understand the substance without re-reading the whole conversation.
\end{itemize}

\subsubsection{PROMOTE-0009 - contract\_candidate}

\begin{itemize}
\item Source conversation: app\_testx\_codehygiene
\item Proposed target path: contracts/** or docs/reference/contracts/** after authority review
\item Current authority support: not yet verified against target docs
\item Current authority conflict: none detected by triage
\item Queue risk: none\_detected
\item Validation required: authority review plus targeted schema/contract parser and FAST
\item Recommended next action: hold\_or\_review
\item Claim: The central project is a deterministic universe engine/game. The user's ambition is not just a normal game but a world runtime that can represent everything from a single-room scenario to an AI-only universe to an MMO. The project is meant to support real-world defaults, arbitrary modding, procedural and defined content, strong epistemics, and long-term replayability.
\end{itemize}

\subsubsection{PROMOTE-0010 - insufficient\_evidence}

\begin{itemize}
\item Source conversation: architecture\_codex\_prompts
\item Proposed target path: none
\item Current authority support: not yet verified against target docs
\item Current authority conflict: none detected by triage
\item Queue risk: none\_detected
\item Validation required: preserve as archive evidence; no live target validation
\item Recommended next action: hold\_or\_review
\item Claim: The most important thing to remember is this: **this chat is the architectural and implementation-roadmap backbone for Dominium/Domino, but it is not an implementation log**. It should be preserved as a design/specification source and as a prompt library, while actual code and current facts must be verified separately.
\end{itemize}

\subsubsection{PROMOTE-0011 - insufficient\_evidence}

\begin{itemize}
\item Source conversation: architecture\_codex\_prompts
\item Proposed target path: none
\item Current authority support: not yet verified against target docs
\item Current authority conflict: none detected by triage
\item Queue risk: none\_detected
\item Validation required: preserve as archive evidence; no live target validation
\item Recommended next action: hold\_or\_review
\item Claim: FACT: The user opened by pasting an "Extended Master Starter Prompt - Dominium + Domino." That prompt established the role of the assistant as a senior engine architect and defined the main project philosophy. Domino was described as a deterministic, fixed-point, C89 engine portable across old and modern OS/architecture combinations. Dominium was described as the product suite built on top of Domino.
\end{itemize}

\subsubsection{PROMOTE-0012 - archive\_only}

\begin{itemize}
\item Source conversation: architecture\_codex\_prompts
\item Proposed target path: docs/archive/conversations/**
\item Current authority support: not yet verified against target docs
\item Current authority conflict: none detected by triage
\item Queue risk: none\_detected
\item Validation required: preserve as archive evidence; no live target validation
\item Recommended next action: hold\_or\_review
\item Claim: FACT: The user asked, succinctly, how a player would build a complete machine workshop. The assistant described a progression: extract raw materials, process them into industrial materials, produce mechanical components, produce electrical components, establish power distribution, establish logistics, assemble machine frames and machines, add measurement/control tooling, then scale horizontally and vertically.
\end{itemize}

\subsubsection{PROMOTE-0013 - insufficient\_evidence}

\begin{itemize}
\item Source conversation: architecture\_ui\_providers
\item Proposed target path: none
\item Current authority support: not yet verified against target docs
\item Current authority conflict: none detected by triage
\item Queue risk: none\_detected
\item Validation required: preserve as archive evidence; no live target validation
\item Recommended next action: hold\_or\_review
\item Claim: The language baseline decision matters because it changes what old renderers and toolchains matter. C17/C++17 reduces the need to design around C89/C++98 constraints but does not remove the need for C-compatible ABI boundaries. The system floor decision similarly moves many older renderers and OS targets into research/back-port categories.
\end{itemize}

\subsubsection{PROMOTE-0014 - insufficient\_evidence}

\begin{itemize}
\item Source conversation: architecture\_ui\_providers
\item Proposed target path: none
\item Current authority support: not yet verified against target docs
\item Current authority conflict: none detected by triage
\item Queue risk: none\_detected
\item Validation required: preserve as archive evidence; no live target validation
\item Recommended next action: hold\_or\_review
\item Claim: The Robot OS and robotic seed-civilisation decisions are the most important design decisions. They are explicit user statements, not merely assistant proposals. They should be formalized as product/game requirements. They affect UI, fog-of-war, spawning, respawn, automation, industry, tutorial/planner design, and MMO anti-cheat.
\end{itemize}

\subsubsection{PROMOTE-0015 - blocked\_by\_queue}

\begin{itemize}
\item Source conversation: architecture\_ui\_providers
\item Proposed target path: none
\item Current authority support: not yet verified against target docs
\item Current authority conflict: blocked queue scope
\item Queue risk: blocked
\item Validation required: preserve as archive evidence; no live target validation
\item Recommended next action: hold\_or\_review
\item Claim: The central tradeoff was speed versus purity. Raylib, SDL2, Lua, raygui, rlgl, rlsw, and related libraries can accelerate visible progress, but if they leak into game law, UI documents, saves, replays, or public contracts, they become architectural capture. The chosen compromise is to use them as providers behind Dominium-owned services and conformance tests.
\end{itemize}

\subsubsection{PROMOTE-0016 - archive\_only}

\begin{itemize}
\item Source conversation: Build\_and\_Future\_Proofing
\item Proposed target path: docs/archive/conversations/**
\item Current authority support: not yet verified against target docs
\item Current authority conflict: none detected by triage
\item Queue risk: none\_detected
\item Validation required: preserve as archive evidence; no live target validation
\item Recommended next action: hold\_or\_review
\item Claim: Plain-language limitation: this package is reliable as a preservation of the current visible chat. It should not be treated as a complete archive of every Dominium discussion in other chats. Where the report uses project or repository context, it is labelled as such. The most important caveat is that many items are recommendations, not final user decisions.
\end{itemize}

\subsubsection{PROMOTE-0017 - reject\_as\_noise}

\begin{itemize}
\item Source conversation: Build\_and\_Future\_Proofing
\item Proposed target path: none
\item Current authority support: not yet verified against target docs
\item Current authority conflict: none detected by triage
\item Queue risk: none\_detected
\item Validation required: preserve as archive evidence; no live target validation
\item Recommended next action: hold\_or\_review
\item Claim: The final uploaded prompt requested a complete preservation package for this chat: a human-readable explanation, structured registers, context transfer packet, spec sheet, aggregator packet, self-audit, and files/ZIP package. This response completes that export task and creates downloadable files for later reading and aggregation.
\end{itemize}

\subsubsection{PROMOTE-0018 - archive\_only}

\begin{itemize}
\item Source conversation: Build\_and\_Future\_Proofing
\item Proposed target path: docs/archive/conversations/**
\item Current authority support: not yet verified against target docs
\item Current authority conflict: none detected by triage
\item Queue risk: none\_detected
\item Validation required: preserve as archive evidence; no live target validation
\item Recommended next action: hold\_or\_review
\item Claim: The final explicit goal was preservation: create a maximum-fidelity package for this chat so the user can understand it later, ask questions in this chat, merge it with other old-chat reports, and eventually build a master project spec book.
\end{itemize}

\subsubsection{PROMOTE-0019 - insufficient\_evidence}

\begin{itemize}
\item Source conversation: canonical\_structure\_and\_framework
\item Proposed target path: none
\item Current authority support: not yet verified against target docs
\item Current authority conflict: none detected by triage
\item Queue risk: none\_detected
\item Validation required: preserve as archive evidence; no live target validation
\item Recommended next action: hold\_or\_review
\item Claim: This produced a pattern: generate a large Codex/AIDE task, the user ran it or reported a result, then the assistant evaluated what was truly fixed versus what was only validator/document churn. The user eventually demanded a one-shot "actual final cleanup" prompt because previous passes had sometimes added validators without moving directories. That prompt explicitly required real git mv routing, not just reports.
\end{itemize}

\subsubsection{PROMOTE-0020 - blocked\_by\_queue}

\begin{itemize}
\item Source conversation: canonical\_structure\_and\_framework
\item Proposed target path: none
\item Current authority support: not yet verified against target docs
\item Current authority conflict: blocked queue scope
\item Queue risk: blocked
\item Validation required: preserve as archive evidence; no live target validation
\item Recommended next action: hold\_or\_review
\item Claim: The final provider directory doctrine became service-first: runtime/<service>/providers/<provider>/, not runtime/<vendor>/<service>/. Provider choices belong in release/profiles/ or content/profiles/, not app path names. External code belongs under external/upstream/ or external/vendor/, but the repo should choose one convention.
\end{itemize}

\subsubsection{PROMOTE-0021 - schema\_candidate}

\begin{itemize}
\item Source conversation: canonical\_structure\_and\_framework
\item Proposed target path: schema/** or schemas/** after authority review
\item Current authority support: not yet verified against target docs
\item Current authority conflict: none detected by triage
\item Queue risk: none\_detected
\item Validation required: authority review plus targeted schema/contract parser and FAST
\item Recommended next action: hold\_or\_review
\item Claim: Packs are distributable authored payloads. The preferred pack layout became category-based, such as content/packs/<category>/<pack\_id>/. Pack IDs must not depend on paths. Saves and replays must use stable IDs, schema versions, canonical serialization, deterministic hashes, and explicit migration/refusal policies. Backward compatibility must be a compatibility corpus plus tests, not intention.
\end{itemize}

\subsubsection{PROMOTE-0022 - insufficient\_evidence}

\begin{itemize}
\item Source conversation: Chronology\_Celestial\_Systems
\item Proposed target path: none
\item Current authority support: not yet verified against target docs
\item Current authority conflict: none detected by triage
\item Queue risk: none\_detected
\item Validation required: preserve as archive evidence; no live target validation
\item Recommended next action: hold\_or\_review
\item Claim: This reinforced the core celestial-content rule: locations should exist explicitly when they change player decisions.
\end{itemize}

\subsubsection{PROMOTE-0023 - insufficient\_evidence}

\begin{itemize}
\item Source conversation: Chronology\_Celestial\_Systems
\item Proposed target path: none
\item Current authority support: not yet verified against target docs
\item Current authority conflict: none detected by triage
\item Queue risk: none\_detected
\item Validation required: preserve as archive evidence; no live target validation
\item Recommended next action: hold\_or\_review
\item Claim: This became the spatial foundation for later time/calendar decisions because every major body or region could potentially have its own local time, calendar, or operational standard.
\end{itemize}

\subsubsection{PROMOTE-0024 - needs\_user\_decision}

\begin{itemize}
\item Source conversation: Chronology\_Celestial\_Systems
\item Proposed target path: docs/archive/conversations/\_decision/DECISION\_DOCKET\_v0.md
\item Current authority support: not yet verified against target docs
\item Current authority conflict: none detected by triage
\item Queue risk: none\_detected
\item Validation required: preserve as archive evidence; no live target validation
\item Recommended next action: hold\_or\_review
\item Claim: The assistant formalized this as the Perfect Earth Calendar, later called HPC-E. The final structure is clear, though the exact leap rule remains unresolved.
\end{itemize}

\subsubsection{PROMOTE-0025 - needs\_user\_decision}

\begin{itemize}
\item Source conversation: development\_routes
\item Proposed target path: docs/archive/conversations/\_decision/DECISION\_DOCKET\_v0.md
\item Current authority support: not yet verified against target docs
\item Current authority conflict: none detected by triage
\item Queue risk: none\_detected
\item Validation required: preserve as archive evidence; no live target validation
\item Recommended next action: hold\_or\_review
\item Claim: Before relying on this chat as project authority, a future assistant or human reviewer must confirm whether Dominium is actually intended to be simulation-heavy, whether strict determinism matters, whether replay or multiplayer are goals, whether modding matters, and what the actual game loop is. Until then, the kernel-first plan should be treated as a strong provisional proposal, not a final specification.
\end{itemize}

\subsubsection{PROMOTE-0026 - insufficient\_evidence}

\begin{itemize}
\item Source conversation: development\_routes
\item Proposed target path: none
\item Current authority support: not yet verified against target docs
\item Current authority conflict: none detected by triage
\item Queue risk: none\_detected
\item Validation required: preserve as archive evidence; no live target validation
\item Recommended next action: hold\_or\_review
\item Claim: The first answer was assertive. It stated that Dominium must follow Route C and described the route as the only viable one for Dominium. Later parts of the chat corrected the status of that claim. FACT: the assistant made that recommendation. UNCERTAIN / UNVERIFIED: the recommendation was not validated with external sources in the chat and was not explicitly accepted by the user.
\end{itemize}

\subsubsection{PROMOTE-0027 - archive\_only}

\begin{itemize}
\item Source conversation: development\_routes
\item Proposed target path: docs/archive/conversations/**
\item Current authority support: not yet verified against target docs
\item Current authority conflict: none detected by triage
\item Queue risk: none\_detected
\item Validation required: preserve as archive evidence; no live target validation
\item Recommended next action: hold\_or\_review
\item Claim: The initial technical answer proposed a layered stack. It started with mathematical and temporal primitives, then moved into the deterministic simulation kernel, then world state, systems, persistence and replay, tooling, presentation, content, and finally distribution, modding, and multiplayer.
\end{itemize}

\subsubsection{PROMOTE-0028 - needs\_user\_decision}

\begin{itemize}
\item Source conversation: documentation\_standards\_readme
\item Proposed target path: docs/archive/conversations/\_decision/DECISION\_DOCKET\_v0.md
\item Current authority support: not yet verified against target docs
\item Current authority conflict: none detected by triage
\item Queue risk: none\_detected
\item Validation required: preserve as archive evidence; no live target validation
\item Recommended next action: hold\_or\_review
\item Claim: The key thing to remember: this chat produced the standards and prompts for future work; it did not verify or change the project. The next assistant must inspect the actual repository before writing repo-specific docs, README content, platform claims, license claims, or subsystem descriptions.
\end{itemize}

\subsubsection{PROMOTE-0029 - contract\_candidate}

\begin{itemize}
\item Source conversation: documentation\_standards\_readme
\item Proposed target path: contracts/** or docs/reference/contracts/** after authority review
\item Current authority support: not yet verified against target docs
\item Current authority conflict: none detected by triage
\item Queue risk: none\_detected
\item Validation required: authority review plus targeted schema/contract parser and FAST
\item Recommended next action: hold\_or\_review
\item Claim: FACT: The chat established that public API contracts should live in headers.
\end{itemize}

\subsubsection{PROMOTE-0030 - insufficient\_evidence}

\begin{itemize}
\item Source conversation: documentation\_standards\_readme
\item Proposed target path: none
\item Current authority support: not yet verified against target docs
\item Current authority conflict: none detected by triage
\item Queue risk: none\_detected
\item Validation required: preserve as archive evidence; no live target validation
\item Recommended next action: hold\_or\_review
\item Claim: FACT: Source files should document implementation rationale, non-obvious invariants, and hazards.
\end{itemize}

\subsubsection{PROMOTE-0031 - reject\_as\_noise}

\begin{itemize}
\item Source conversation: Dominium\_Architecture\_I
\item Proposed target path: none
\item Current authority support: not yet verified against target docs
\item Current authority conflict: none detected by triage
\item Queue risk: none\_detected
\item Validation required: preserve as archive evidence; no live target validation
\item Recommended next action: hold\_or\_review
\item Claim: During this stage, the work was still broadly about "the project spec." The user wanted a comprehensive description of the game and its technical systems. The assistant produced large design documents, but many of those earlier contents are not visible in the retained transcript. This is important because the final report package marks those earlier outputs as referenced but not fully recovered.
\end{itemize}

\subsubsection{PROMOTE-0032 - architecture\_doc\_candidate}

\begin{itemize}
\item Source conversation: Dominium\_Architecture\_I
\item Proposed target path: docs/architecture/** after target selection
\item Current authority support: not yet verified against target docs
\item Current authority conflict: none detected by triage
\item Queue risk: none\_detected
\item Validation required: source check, authority-order check, link check, generated output validator, FAST
\item Recommended next action: create\_microtask
\item Claim: The user also made a clear versioning decision. **FACT:** The user wanted **major.minor.patch semantic versioning** for all components/packages and for the complete game package. The user explicitly did **not** want build numbers or build dates in filenames. Instead, those details should live in metadata. This became part of the future build/package design.
\end{itemize}

\subsubsection{PROMOTE-0033 - insufficient\_evidence}

\begin{itemize}
\item Source conversation: Dominium\_Architecture\_I
\item Proposed target path: none
\item Current authority support: not yet verified against target docs
\item Current authority conflict: none detected by triage
\item Queue risk: none\_detected
\item Validation required: preserve as archive evidence; no live target validation
\item Recommended next action: hold\_or\_review
\item Claim: The key change came when the user decided that Codex could not read the entire conversation at once. **FACT:** The user said they would copy and paste the transcript with timestamps and would not edit it manually, but because Codex could not read the whole context at once, they would not bother asking Codex to compile a full book directly from the whole chat.
\end{itemize}

\subsubsection{PROMOTE-0034 - reject\_as\_noise}

\begin{itemize}
\item Source conversation: Dominium\_Architecture\_II
\item Proposed target path: none
\item Current authority support: not yet verified against target docs
\item Current authority conflict: none detected by triage
\item Queue risk: none\_detected
\item Validation required: preserve as archive evidence; no live target validation
\item Recommended next action: hold\_or\_review
\item Claim: The last part of the chat was handoff extraction. The user asked for discovery inventory, structured registers, a context transfer packet, and then a downloadable package. The generated package captured the chat's workstreams, decisions, tasks, constraints, open questions, artifacts, risks, and verification items. This current report is a plain-language explanation of that substance.
\end{itemize}

\subsubsection{PROMOTE-0035 - architecture\_doc\_candidate}

\begin{itemize}
\item Source conversation: Dominium\_Architecture\_II
\item Proposed target path: docs/architecture/** after target selection
\item Current authority support: not yet verified against target docs
\item Current authority conflict: none detected by triage
\item Queue risk: none\_detected
\item Validation required: source check, authority-order check, link check, generated output validator, FAST
\item Recommended next action: create\_microtask
\item Claim: The entire conversation is anchored by the requirement that Dominium must run deterministically. This means the same initial state and same inputs must produce the same state across machines, compilers, operating systems, and eventually retro and modern platforms.
\end{itemize}

\subsubsection{PROMOTE-0036 - insufficient\_evidence}

\begin{itemize}
\item Source conversation: Dominium\_Architecture\_II
\item Proposed target path: none
\item Current authority support: not yet verified against target docs
\item Current authority conflict: none detected by triage
\item Queue risk: none\_detected
\item Validation required: preserve as archive evidence; no live target validation
\item Recommended next action: hold\_or\_review
\item Claim: The chat discussed fixed tick phases, virtual lanes, merge order, command buffers, event queues, and deterministic serialization. The simulation must not depend on render FPS, wall-clock time, OS scheduler behaviour, GPU state, or floating-point quirks. This explains many later decisions: C89 core, integer/fixed-point math, no sim floats, no platform headers in simulation, and renderer as pure observer.
\end{itemize}

\subsubsection{PROMOTE-0037 - insufficient\_evidence}

\begin{itemize}
\item Source conversation: Dominium\_Architecture\_III
\item Proposed target path: none
\item Current authority support: not yet verified against target docs
\item Current authority conflict: none detected by triage
\item Queue risk: none\_detected
\item Validation required: preserve as archive evidence; no live target validation
\item Recommended next action: hold\_or\_review
\item Claim: The view system then expanded beyond the launcher. **FACT:** The user wants instant zoom to any scale, instant switching between top-down 2D and first-person 3D, and arbitrary cameras for free cam, map viewing, content creation, HUD cameras, CCTV, overlays, windows, and offscreen targets. These are client/render-side concerns, not simulation concerns.
\end{itemize}

\subsubsection{PROMOTE-0038 - architecture\_doc\_candidate}

\begin{itemize}
\item Source conversation: Dominium\_Architecture\_III
\item Proposed target path: docs/architecture/** after target selection
\item Current authority support: not yet verified against target docs
\item Current authority conflict: none detected by triage
\item Queue risk: none\_detected
\item Validation required: source check, authority-order check, link check, generated output validator, FAST
\item Recommended next action: create\_microtask
\item Claim: This context matters because the rest of the conversation happened inside those constraints. The launcher could not become a dumping ground for OS-specific behavior, sim mutation, renderer decisions, or nondeterministic game-state changes. It had to fit the project's deterministic layering.
\end{itemize}

\subsubsection{PROMOTE-0039 - insufficient\_evidence}

\begin{itemize}
\item Source conversation: Dominium\_Architecture\_III
\item Proposed target path: none
\item Current authority support: not yet verified against target docs
\item Current authority conflict: none detected by triage
\item Queue risk: none\_detected
\item Validation required: preserve as archive evidence; no live target validation
\item Recommended next action: hold\_or\_review
\item Claim: The user then uploaded dominium.7z, saying it was the git repository as of that moment. **FACT:** The prior assistant said it could not open .7z archives in that environment. That means the actual repository state remained unverified at that stage. This became important later, because many implementation ideas were discussed without confirmed access to the repo contents.
\end{itemize}

\subsubsection{PROMOTE-0040 - insufficient\_evidence}

\begin{itemize}
\item Source conversation: Dominium\_Architecture\_IV
\item Proposed target path: none
\item Current authority support: not yet verified against target docs
\item Current authority conflict: none detected by triage
\item Queue risk: none\_detected
\item Validation required: preserve as archive evidence; no live target validation
\item Recommended next action: hold\_or\_review
\item Claim: The reusable engine layer was named **Domino**. This became one of the most important naming and architectural decisions in the chat. Domino is the engine/core/platform/sim/mod layer, reusable for other games and projects. Dominium is the game/product layer built on it.
\end{itemize}

\subsubsection{PROMOTE-0041 - insufficient\_evidence}

\begin{itemize}
\item Source conversation: Dominium\_Architecture\_IV
\item Proposed target path: none
\item Current authority support: not yet verified against target docs
\item Current authority conflict: none detected by triage
\item Queue risk: none\_detected
\item Validation required: preserve as archive evidence; no live target validation
\item Recommended next action: hold\_or\_review
\item Claim: Blueprints are plans or diffs: place element, remove element, modify terrain, place machine, connect network, etc. When accepted, they generate jobs. Manual player actions and queued work for humans, robots, or other agents should use the same job pipeline. This matters because it keeps replay, AI, and automation consistent.
\end{itemize}

\subsubsection{PROMOTE-0042 - insufficient\_evidence}

\begin{itemize}
\item Source conversation: Dominium\_Architecture\_IV
\item Proposed target path: none
\item Current authority support: not yet verified against target docs
\item Current authority conflict: none detected by triage
\item Queue risk: none\_detected
\item Validation required: preserve as archive evidence; no live target validation
\item Recommended next action: hold\_or\_review
\item Claim: The user then asked how to structure Codex prompts. The final high-level roadmap became:
\end{itemize}

\subsubsection{PROMOTE-0043 - needs\_user\_decision}

\begin{itemize}
\item Source conversation: Dominium\_Complete\_Conversation
\item Proposed target path: docs/archive/conversations/\_decision/DECISION\_DOCKET\_v0.md
\item Current authority support: not yet verified against target docs
\item Current authority conflict: none detected by triage
\item Queue risk: none\_detected
\item Validation required: preserve as archive evidence; no live target validation
\item Recommended next action: hold\_or\_review
\item Claim: This chat was about preserving and re-grounding Dominium's architecture around a very old constitutional contract, then checking whether the current GitHub repository still follows that contract, and finally broadening the question into a platform-engineering doctrine for portability, modularity, extensibility, reuse, future-proofing, and long-term compatibility.
\end{itemize}

\subsubsection{PROMOTE-0044 - contract\_candidate}

\begin{itemize}
\item Source conversation: Dominium\_Complete\_Conversation
\item Proposed target path: contracts/** or docs/reference/contracts/** after authority review
\item Current authority support: not yet verified against target docs
\item Current authority conflict: none detected by triage
\item Queue risk: none\_detected
\item Validation required: authority review plus targeted schema/contract parser and FAST
\item Recommended next action: hold\_or\_review
\item Claim: The glossary bound terms such as Authority, Law, Process, Lens, SessionSpec, UniverseIdentity, Macro Capsule, SRZ, RepoX, TestX, AuditX, CompatX, SecureX, and related terms. It matters because future assistants must not use sloppy synonyms like "mode" where the canon requires ExperienceProfile or LawProfile. This topic directly supports modularity because stable vocabulary is part of stable architecture.
\end{itemize}

\subsubsection{PROMOTE-0045 - insufficient\_evidence}

\begin{itemize}
\item Source conversation: Dominium\_Complete\_Conversation
\item Proposed target path: none
\item Current authority support: not yet verified against target docs
\item Current authority conflict: none detected by triage
\item Queue risk: none\_detected
\item Validation required: preserve as archive evidence; no live target validation
\item Recommended next action: hold\_or\_review
\item Claim: Uncertainty: during this preservation pass, a potential inconsistency appeared: some docs claim product roots moved under apps/, while an earlier inspected code path under root client/ was available and apps/client was not fetched successfully. This must be verified against the current physical tree.
\end{itemize}

\subsubsection{PROMOTE-0046 - insufficient\_evidence}

\begin{itemize}
\item Source conversation: dominium\_domino\_codex\_planning
\item Proposed target path: none
\item Current authority support: not yet verified against target docs
\item Current authority conflict: none detected by triage
\item Queue risk: none\_detected
\item Validation required: preserve as archive evidence; no live target validation
\item Recommended next action: hold\_or\_review
\item Claim: INFERENCE:** The user accepted this direction in practice, because the next request was to generate Codex prompts to implement each step fully. The user did not explicitly say, "I formally approve the Milestone-0 plan," but they proceeded to ask for prompts based on it.
\end{itemize}

\subsubsection{PROMOTE-0047 - insufficient\_evidence}

\begin{itemize}
\item Source conversation: dominium\_domino\_codex\_planning
\item Proposed target path: none
\item Current authority support: not yet verified against target docs
\item Current authority conflict: none detected by triage
\item Queue risk: none\_detected
\item Validation required: preserve as archive evidence; no live target validation
\item Recommended next action: hold\_or\_review
\item Claim: The user then asked for recommendations. The assistant recommended a clean startup policy: explicit flags first, no environment/config override in v1, terminal detection through dsys, generic AUTO behavior, product-specific headless handling for the game, build-time GUI/TUI capabilities, and structural error codes for fallback. The user accepted this enough to request a "one big Codex prompt" to implement it.
\end{itemize}

\subsubsection{PROMOTE-0048 - reject\_as\_noise}

\begin{itemize}
\item Source conversation: dominium\_domino\_codex\_planning
\item Proposed target path: none
\item Current authority support: not yet verified against target docs
\item Current authority conflict: none detected by triage
\item Queue risk: none\_detected
\item Validation required: preserve as archive evidence; no live target validation
\item Recommended next action: hold\_or\_review
\item Claim: However, the later context transfer packet treated that inspection claim as **UNCERTAIN / UNVERIFIED**. This was careful and correct: within the visible chat, there was no reliable tool trace proving the repo had actually been inspected. Therefore, future assistants must verify dominium.zip or the active checkout before relying on the platform/render/API answer.
\end{itemize}

\subsubsection{PROMOTE-0049 - blocked\_by\_queue}

\begin{itemize}
\item Source conversation: dominium\_full\_conversation
\item Proposed target path: none
\item Current authority support: not yet verified against target docs
\item Current authority conflict: blocked queue scope
\item Queue risk: blocked
\item Validation required: preserve as archive evidence; no live target validation
\item Recommended next action: hold\_or\_review
\item Claim: A future assistant must understand that this chat is not simply about UI. It records the architecture for how Dominium should build itself: through contracts, commands, services, providers, documents, patches, modules, packs, apps, workspaces, diagnostics, evidence, AIDE, Codex, and eventually Workbench.
\end{itemize}

\subsubsection{PROMOTE-0050 - architecture\_doc\_candidate}

\begin{itemize}
\item Source conversation: dominium\_full\_conversation
\item Proposed target path: docs/architecture/** after target selection
\item Current authority support: not yet verified against target docs
\item Current authority conflict: none detected by triage
\item Queue risk: none\_detected
\item Validation required: source check, authority-order check, link check, generated output validator, FAST
\item Recommended next action: create\_microtask
\item Claim: The user then redirected the entire plan. They argued that native OS widget GUI tools already exist through Visual Studio, Xcode, etc. The project needed a cross-platform rendered tool environment that uses the same CLI, TUI, and rendered GUI systems as the client. The old UI Editor and Tool Editor were good exploratory ideas but bad final products. This produced the Dominium Workbench concept.
\end{itemize}

\subsubsection{PROMOTE-0051 - insufficient\_evidence}

\begin{itemize}
\item Source conversation: dominium\_full\_conversation
\item Proposed target path: none
\item Current authority support: not yet verified against target docs
\item Current authority conflict: none detected by triage
\item Queue risk: none\_detected
\item Validation required: preserve as archive evidence; no live target validation
\item Recommended next action: hold\_or\_review
\item Claim: The discussion then moved into AIDE. The user wanted to automate as much as possible through Codex and AIDE. We developed a task operating model: AIDE creates WorkUnits, tracks attempts, blockers, evidence, repairs, resumes, checkpoints, and promotion decisions. Codex executes bounded tasks. Development can continue with classified partials and repairable blockers; promotion to main requires evidence.
\end{itemize}

\subsubsection{PROMOTE-0052 - reject\_as\_noise}

\begin{itemize}
\item Source conversation: dominium\_setup
\item Proposed target path: none
\item Current authority support: not yet verified against target docs
\item Current authority conflict: none detected by triage
\item Queue risk: none\_detected
\item Validation required: preserve as archive evidence; no live target validation
\item Recommended next action: hold\_or\_review
\item Claim: Near the end, the chat shifted into preservation mode. The user asked for a maximum-fidelity context transfer packet, then for a downloadable report package, then for an in-chat reader version of that package. Those outputs preserved the decisions, tasks, constraints, artifacts, rejected options, and verification queue for future assistants or aggregation into a larger Project Spec Book.
\end{itemize}

\subsubsection{PROMOTE-0053 - archive\_only}

\begin{itemize}
\item Source conversation: dominium\_setup
\item Proposed target path: docs/archive/conversations/**
\item Current authority support: not yet verified against target docs
\item Current authority conflict: none detected by triage
\item Queue risk: none\_detected
\item Validation required: preserve as archive evidence; no live target validation
\item Recommended next action: hold\_or\_review
\item Claim: These directory trees were not final, but they mattered because they exposed the user's preference for simple, shallow, logical structures and native IDE workflows that do not become the source of truth.
\end{itemize}

\subsubsection{PROMOTE-0054 - planning\_doc\_candidate}

\begin{itemize}
\item Source conversation: dominium\_setup
\item Proposed target path: docs/planning/** after target selection
\item Current authority support: not yet verified against target docs
\item Current authority conflict: none detected by triage
\item Queue risk: none\_detected
\item Validation required: source check, authority-order check, link check, generated output validator, FAST
\item Recommended next action: create\_microtask
\item Claim: The user later asked where to create Visual Studio and Xcode apps after opening the repo as a folder. The assistant advised that Visual Studio and Xcode projects should be generated through CMake, not hand-authored or treated as authoritative source files. This decision became consistent with later Codex prompts.
\end{itemize}

\subsubsection{PROMOTE-0055 - contract\_candidate}

\begin{itemize}
\item Source conversation: Domino\_Dominium\_Workbench
\item Proposed target path: contracts/** or docs/reference/contracts/** after authority review
\item Current authority support: not yet verified against target docs
\item Current authority conflict: none detected by triage
\item Queue risk: none\_detected
\item Validation required: authority review plus targeted schema/contract parser and FAST
\item Recommended next action: hold\_or\_review
\item Claim: The original UI Editor / Tool Editor plan is now **superseded as a final product**. It is not lost: its useful parts become Workbench capabilities, especially Interface Studio, UI/HUD Sandbox, Theme Laboratory, TUI Studio, Rendered GUI Studio, Preview Matrix, validation panels, import/export tools, and document-patch workflows.
\end{itemize}

\subsubsection{PROMOTE-0056 - blocked\_by\_queue}

\begin{itemize}
\item Source conversation: Domino\_Dominium\_Workbench
\item Proposed target path: none
\item Current authority support: not yet verified against target docs
\item Current authority conflict: blocked queue scope
\item Queue risk: blocked
\item Validation required: preserve as archive evidence; no live target validation
\item Recommended next action: hold\_or\_review
\item Claim: The central decision is that Workbench must not become semantic authority. Workbench is a surface over contracts, commands, services, documents, patches, providers, packs, modules, apps, artifacts, diagnostics, evidence, tests, and replay/proof. It is the richest human/agent interface over the system, but it must not bypass command, capability, refusal, validation, document, patch, or proof law.
\end{itemize}

\subsubsection{PROMOTE-0057 - blocked\_by\_queue}

\begin{itemize}
\item Source conversation: Domino\_Dominium\_Workbench
\item Proposed target path: none
\item Current authority support: not yet verified against target docs
\item Current authority conflict: blocked queue scope
\item Queue risk: blocked
\item Validation required: preserve as archive evidence; no live target validation
\item Recommended next action: hold\_or\_review
\item Claim: The second central decision is that CLI, TUI, rendered GUI, OS-native GUI, and headless automation must be **projections of the same command/result/refusal/document/view spine**. A GUI button, a TUI hotkey, a CLI command, a headless JSON job, and a future OS-native widget action should route through the same command/service path and produce the same typed result, diagnostics, refusals, logs, and evidence.
\end{itemize}

\subsubsection{PROMOTE-0058 - insufficient\_evidence}

\begin{itemize}
\item Source conversation: domino\_engine\_refactor\_prompts
\item Proposed target path: none
\item Current authority support: not yet verified against target docs
\item Current authority conflict: none detected by triage
\item Queue risk: none\_detected
\item Validation required: preserve as archive evidence; no live target validation
\item Recommended next action: hold\_or\_review
\item Claim: FACT:** The Domino engine must be ISO C89. The Dominium UI/tools layer must be portable C++98. Engine code must not use OS APIs. Determinism matters, so floats are disallowed in deterministic paths. The engine must avoid platform-dependent behavior and hardcoded game semantics. Everything should be content-driven.
\end{itemize}

\subsubsection{PROMOTE-0059 - contract\_candidate}

\begin{itemize}
\item Source conversation: domino\_engine\_refactor\_prompts
\item Proposed target path: contracts/** or docs/reference/contracts/** after authority review
\item Current authority support: not yet verified against target docs
\item Current authority conflict: none detected by triage
\item Queue risk: none\_detected
\item Validation required: authority review plus targeted schema/contract parser and FAST
\item Recommended next action: hold\_or\_review
\item Claim: This became one of the central architectural ideas of the chat. It makes behavior extensible without giving minds direct write access to the world. A mod can add new sensors, observations, intents, capabilities, or actions, but state mutation still passes through a deterministic pipeline.
\end{itemize}

\subsubsection{PROMOTE-0060 - reject\_as\_noise}

\begin{itemize}
\item Source conversation: domino\_engine\_refactor\_prompts
\item Proposed target path: none
\item Current authority support: not yet verified against target docs
\item Current authority conflict: none detected by triage
\item Queue risk: none\_detected
\item Validation required: preserve as archive evidence; no live target validation
\item Recommended next action: hold\_or\_review
\item Claim: After the architecture and prompts were created, the user asked for a maximum-fidelity context transfer packet. The assistant produced one. The user then asked for a downloadable report package; the assistant generated package files and a ZIP. The user later asked for an in-chat reader version of that package, and finally requested this current human-readable explanatory report.
\end{itemize}

\subsubsection{PROMOTE-0061 - insufficient\_evidence}

\begin{itemize}
\item Source conversation: engine\_baseline\_architecture
\item Proposed target path: none
\item Current authority support: not yet verified against target docs
\item Current authority conflict: none detected by triage
\item Queue risk: none\_detected
\item Validation required: preserve as archive evidence; no live target validation
\item Recommended next action: hold\_or\_review
\item Claim: This corrected the plan. The final sequencing became: first **Milestone 0: Make the baseline path honest**. That means fixing server/runtime circular import, CLI forwarding, session\_create -> session\_boot, missing time-anchor policy registry, and making the strict local playtest validator pass. Only after that should the builder/destruction lab begin.
\end{itemize}

\subsubsection{PROMOTE-0062 - reject\_as\_noise}

\begin{itemize}
\item Source conversation: engine\_baseline\_architecture
\item Proposed target path: none
\item Current authority support: not yet verified against target docs
\item Current authority conflict: none detected by triage
\item Queue risk: none\_detected
\item Validation required: preserve as archive evidence; no live target validation
\item Recommended next action: hold\_or\_review
\item Claim: The final user action uploaded Pasted text.txt, a detailed instruction prompt requiring a full preservation report, structured registers, spec sheet, aggregator packet, self-audit, and downloadable file package for this chat. That request is the current task.
\end{itemize}

\subsubsection{PROMOTE-0063 - architecture\_doc\_candidate}

\begin{itemize}
\item Source conversation: engine\_baseline\_architecture
\item Proposed target path: docs/architecture/** after target selection
\item Current authority support: not yet verified against target docs
\item Current authority conflict: none detected by triage
\item Queue risk: none\_detected
\item Validation required: source check, authority-order check, link check, generated output validator, FAST
\item Recommended next action: create\_microtask
\item Claim: The central architecture distinction was that Domino should be a reusable deterministic simulation substrate, while Dominium should be one game/product layer built on it. This came up immediately when the user asked for a total description of the project. It mattered because the user wants the code to be reusable for future games and possibly other engine projects.
\end{itemize}

\subsubsection{PROMOTE-0064 - insufficient\_evidence}

\begin{itemize}
\item Source conversation: Foundation\_Workbench\_Codex
\item Proposed target path: none
\item Current authority support: not yet verified against target docs
\item Current authority conflict: none detected by triage
\item Queue risk: none\_detected
\item Validation required: preserve as archive evidence; no live target validation
\item Recommended next action: hold\_or\_review
\item Claim: After the root skeleton improved, the assistant and user recognized that the deeper problem was no longer top-level directories. The problem became semantic duplication and governance: what is public, what is private, what is stable, what is provisional, what is generated, what is a fixture, what must stay compatible, what can be replaced, and what proof is required.
\end{itemize}

\subsubsection{PROMOTE-0065 - blocked\_by\_queue}

\begin{itemize}
\item Source conversation: Foundation\_Workbench\_Codex
\item Proposed target path: none
\item Current authority support: not yet verified against target docs
\item Current authority conflict: blocked queue scope
\item Queue risk: blocked
\item Validation required: preserve as archive evidence; no live target validation
\item Recommended next action: hold\_or\_review
\item Claim: The key conclusion was that Workbench is not the general module system; it is one consumer of the module/command/service/provider/pack/artifact system. Workbench must not call private tools directly. It must route through registered commands and typed results, diagnostics, refusals, views, and evidence.
\end{itemize}

\subsubsection{PROMOTE-0066 - contract\_candidate}

\begin{itemize}
\item Source conversation: Foundation\_Workbench\_Codex
\item Proposed target path: contracts/** or docs/reference/contracts/** after authority review
\item Current authority support: not yet verified against target docs
\item Current authority conflict: none detected by triage
\item Queue risk: none\_detected
\item Validation required: authority review plus targeted schema/contract parser and FAST
\item Recommended next action: hold\_or\_review
\item Claim: The topic came up because the project needed a model that could support world simulation, modding, tooling, Workbench, release, portability, and future games without repeating endless refactors. The conclusion was that stable contracts and semantic IDs must be separated from replaceable private implementation.
\end{itemize}

\subsubsection{PROMOTE-0067 - reject\_as\_noise}

\begin{itemize}
\item Source conversation: Framework\_Open\_Source\_Provider
\item Proposed target path: none
\item Current authority support: not yet verified against target docs
\item Current authority conflict: none detected by triage
\item Queue risk: none\_detected
\item Validation required: preserve as archive evidence; no live target validation
\item Recommended next action: hold\_or\_review
\item Claim: Finally, the user uploaded a detailed preservation prompt and requested a maximum-fidelity report, registers, spec sheet, aggregator packet, audit, and downloadable files. This package is the response to that request.
\end{itemize}

\subsubsection{PROMOTE-0068 - insufficient\_evidence}

\begin{itemize}
\item Source conversation: Framework\_Open\_Source\_Provider
\item Proposed target path: none
\item Current authority support: not yet verified against target docs
\item Current authority conflict: none detected by triage
\item Queue risk: none\_detected
\item Validation required: preserve as archive evidence; no live target validation
\item Recommended next action: hold\_or\_review
\item Claim: This became the central architectural theme. The conversation refined earlier vendor-shaped paths into service-first paths. The stable pattern is:
\end{itemize}

\subsubsection{PROMOTE-0069 - blocked\_by\_queue}

\begin{itemize}
\item Source conversation: Framework\_Open\_Source\_Provider
\item Proposed target path: none
\item Current authority support: not yet verified against target docs
\item Current authority conflict: blocked queue scope
\item Queue risk: blocked
\item Validation required: preserve as archive evidence; no live target validation
\item Recommended next action: hold\_or\_review
\item Claim: The chat inspected julesc013/dominium and found evidence of existing abstractions such as system and graphics layers, stubs, and soft-backed backends. The finding was that raylib could fill concrete provider gaps without replacing the architecture. However, current repo facts are stale/uncertain and must be verified, especially the C17/C++17 vs C90/C++98 baseline contradiction.
\end{itemize}

\subsubsection{PROMOTE-0070 - insufficient\_evidence}

\begin{itemize}
\item Source conversation: gui\_binary\_content
\item Proposed target path: none
\item Current authority support: not yet verified against target docs
\item Current authority conflict: none detected by triage
\item Queue risk: none\_detected
\item Validation required: preserve as archive evidence; no live target validation
\item Recommended next action: hold\_or\_review
\item Claim: That was an important change of direction. It meant the assistant's polished prompt should not be treated as final. The work needed conceptual discussion first. The assistant then analyzed CONTENT0 and identified several issues that could cause bad assumptions to become embedded if Codex acted too soon.
\end{itemize}

\subsubsection{PROMOTE-0071 - insufficient\_evidence}

\begin{itemize}
\item Source conversation: gui\_binary\_content
\item Proposed target path: none
\item Current authority support: not yet verified against target docs
\item Current authority conflict: none detected by triage
\item Queue risk: none\_detected
\item Validation required: preserve as archive evidence; no live target validation
\item Recommended next action: hold\_or\_review
\item Claim: what start scenarios must explicitly exclude,
\end{itemize}

\subsubsection{PROMOTE-0072 - insufficient\_evidence}

\begin{itemize}
\item Source conversation: gui\_binary\_content
\item Proposed target path: none
\item Current authority support: not yet verified against target docs
\item Current authority conflict: none detected by triage
\item Queue risk: none\_detected
\item Validation required: preserve as archive evidence; no live target validation
\item Recommended next action: hold\_or\_review
\item Claim: This shifted the content work from "populate required files" to "design a scalable content representation system." It also made clear that procedural generation and defined data are not enemies in the user's vision. The desired architecture must allow them to coexist.
\end{itemize}

\subsubsection{PROMOTE-0073 - architecture\_doc\_candidate}

\begin{itemize}
\item Source conversation: Language\_Platform\_Architecture
\item Proposed target path: docs/architecture/** after target selection
\item Current authority support: not yet verified against target docs
\item Current authority conflict: none detected by triage
\item Queue risk: none\_detected
\item Validation required: source check, authority-order check, link check, generated output validator, FAST
\item Recommended next action: create\_microtask
\item Claim: The 32-bit question followed. The answer was to make full native products 64-bit and keep 32-bit as constrained or projection support only. The key caveat was not to make native pointer size part of game law. Save, replay, network, IDs, and renderer IR must use fixed-width types and never serialize size\_t, long, uintptr\_t, or raw pointers.
\end{itemize}

\subsubsection{PROMOTE-0074 - contract\_candidate}

\begin{itemize}
\item Source conversation: Language\_Platform\_Architecture
\item Proposed target path: contracts/** or docs/reference/contracts/** after authority review
\item Current authority support: not yet verified against target docs
\item Current authority conflict: none detected by triage
\item Queue risk: none\_detected
\item Validation required: authority review plus targeted schema/contract parser and FAST
\item Recommended next action: hold\_or\_review
\item Claim: The user then consolidated a broad future plan around Domino, Dominium, Workbench, AIDE, contracts, tests, replay, and evidence. The answer agreed that the plan was strong, but identified missing central pieces: composition resolver, lockfiles, compatibility corpus, trust/permissions, virtual filesystem roots, performance budgets, and stable public-surface promotion rules.
\end{itemize}

\subsubsection{PROMOTE-0075 - blocked\_by\_queue}

\begin{itemize}
\item Source conversation: Language\_Platform\_Architecture
\item Proposed target path: none
\item Current authority support: not yet verified against target docs
\item Current authority conflict: blocked queue scope
\item Queue risk: blocked
\item Validation required: preserve as archive evidence; no live target validation
\item Recommended next action: hold\_or\_review
\item Claim: Finally, the user pasted advice favoring C99 or C++11 due to raylib/SDL and legacy targets. The answer rejected a pivot. Raylib and SDL2 being C APIs only means provider boundaries should be C-callable; it does not force the whole engine or game to C99. The final recommendation remained C17 + C++17, with raylib/SDL2 treated as providers and with external deployment claims placed into the verification queue.
\end{itemize}

\subsubsection{PROMOTE-0076 - needs\_user\_decision}

\begin{itemize}
\item Source conversation: launcher\_app\_layer
\item Proposed target path: docs/archive/conversations/\_decision/DECISION\_DOCKET\_v0.md
\item Current authority support: not yet verified against target docs
\item Current authority conflict: none detected by triage
\item Queue risk: none\_detected
\item Validation required: preserve as archive evidence; no live target validation
\item Recommended next action: hold\_or\_review
\item Claim: Someone reading this report should understand one central thing: this chat is not about inventing new launcher features anymore. It is about preserving boundaries, making the launcher implementation explicit, and ensuring future work happens on top of verified code rather than vague architectural memory.
\end{itemize}

\subsubsection{PROMOTE-0077 - archive\_only}

\begin{itemize}
\item Source conversation: launcher\_app\_layer
\item Proposed target path: docs/archive/conversations/**
\item Current authority support: not yet verified against target docs
\item Current authority conflict: none detected by triage
\item Queue risk: none\_detected
\item Validation required: preserve as archive evidence; no live target validation
\item Recommended next action: hold\_or\_review
\item Claim: There was also an early suggestion to create a DUI facade, a Dominium UI abstraction, to support native widgets and fallback rendering. This idea was useful as a conceptual stepping stone but was not ultimately locked as a final requirement in the form originally suggested. Later application-layer canon emphasized **UI IR**, **command graphs**, and **binding validation** rather than a specific DUI facade design.
\end{itemize}

\subsubsection{PROMOTE-0078 - schema\_candidate}

\begin{itemize}
\item Source conversation: launcher\_app\_layer
\item Proposed target path: schema/** or schemas/** after authority review
\item Current authority support: not yet verified against target docs
\item Current authority conflict: none detected by triage
\item Queue risk: none\_detected
\item Validation required: authority review plus targeted schema/contract parser and FAST
\item Recommended next action: hold\_or\_review
\item Claim: This phase mattered because it framed the launcher not as a menu program but as a **control plane** for installed products, packs, profiles, compatibility, and launch contracts. However, later canon tightened the permitted communication routes: cross-product communication must go through schema/ and libs/contracts, not arbitrary plugin conventions.
\end{itemize}

\subsubsection{PROMOTE-0079 - reject\_as\_noise}

\begin{itemize}
\item Source conversation: Launcher\_Setup\_Architecture
\item Proposed target path: none
\item Current authority support: not yet verified against target docs
\item Current authority conflict: none detected by triage
\item Queue risk: none\_detected
\item Validation required: preserve as archive evidence; no live target validation
\item Recommended next action: hold\_or\_review
\item Claim: The chat also produced multiple Codex work-order prompts and finally a downloadable report package. Those artifacts are useful, but the substance is the architecture: keep engine deterministic, keep setup/launcher/runtime boundaries explicit, make the launcher optional and modular, use a strong process/instance model, and now implement the launcher over dsys/dgfx rather than ad hoc platform UI stacks.
\end{itemize}

\subsubsection{PROMOTE-0080 - insufficient\_evidence}

\begin{itemize}
\item Source conversation: Launcher\_Setup\_Architecture
\item Proposed target path: none
\item Current authority support: not yet verified against target docs
\item Current authority conflict: none detected by triage
\item Queue risk: none\_detected
\item Validation required: preserve as archive evidence; no live target validation
\item Recommended next action: hold\_or\_review
\item Claim: FACT:** The user's initial architecture text emphasized:
\end{itemize}

\subsubsection{PROMOTE-0081 - contract\_candidate}

\begin{itemize}
\item Source conversation: Launcher\_Setup\_Architecture
\item Proposed target path: contracts/** or docs/reference/contracts/** after authority review
\item Current authority support: not yet verified against target docs
\item Current authority conflict: none detected by triage
\item Queue risk: none\_detected
\item Validation required: authority review plus targeted schema/contract parser and FAST
\item Recommended next action: hold\_or\_review
\item Claim: The assistant responded by stress-testing this philosophy. It pointed out determinism risks around Lua, plugins, time sources, and ABI details, and suggested more explicit contracts for file headers, TLV sections, runtime CLI capabilities, plugin exported symbols, and setup/launcher boundaries. These were assistant proposals, not all user-stated decisions, but the user continued building on that direction.
\end{itemize}

\subsubsection{PROMOTE-0082 - reject\_as\_noise}

\begin{itemize}
\item Source conversation: Modularity\_AIDE\_Refactorability
\item Proposed target path: none
\item Current authority support: not yet verified against target docs
\item Current authority conflict: none detected by triage
\item Queue risk: none\_detected
\item Validation required: preserve as archive evidence; no live target validation
\item Recommended next action: hold\_or\_review
\item Claim: The final user action was uploading a preservation-package prompt. That prompt requested a maximum-fidelity preservation report, structured registers, context transfer packet, spec sheet, aggregator packet, self-audit, downloadable files, and a final in-chat reader. This package is the result.
\end{itemize}

\subsubsection{PROMOTE-0083 - insufficient\_evidence}

\begin{itemize}
\item Source conversation: Modularity\_AIDE\_Refactorability
\item Proposed target path: none
\item Current authority support: not yet verified against target docs
\item Current authority conflict: none detected by triage
\item Queue risk: none\_detected
\item Validation required: preserve as archive evidence; no live target validation
\item Recommended next action: hold\_or\_review
\item Claim: The user objected to the XStack/AuditX/RepoX/TestX framing and wanted AIDE adopted quickly. The resulting plan made AIDE the repo-native control plane for inventory, task queues, policies, move maps, salvage maps, evidence ledgers, validation, and refactor history. Existing tools should be recycled, not discarded. This is central to future refactorability.
\end{itemize}

\subsubsection{PROMOTE-0084 - architecture\_doc\_candidate}

\begin{itemize}
\item Source conversation: Modularity\_AIDE\_Refactorability
\item Proposed target path: docs/architecture/** after target selection
\item Current authority support: not yet verified against target docs
\item Current authority conflict: none detected by triage
\item Queue risk: none\_detected
\item Validation required: source check, authority-order check, link check, generated output validator, FAST
\item Recommended next action: create\_microtask
\item Claim: The discussion proposed target-based CMake and module boundaries rather than path mythology. Public headers, private headers, allowed dependencies, and forbidden dependencies should be explicit. Apps must remain thin. Engine must not depend on game/apps/runtime UI. Runtime adapts host/platform/rendering without owning simulation truth.
\end{itemize}

\subsubsection{PROMOTE-0085 - insufficient\_evidence}

\begin{itemize}
\item Source conversation: omega\_xi\_pi\_architecture\_future
\item Proposed target path: none
\item Current authority support: not yet verified against target docs
\item Current authority conflict: none detected by triage
\item Queue risk: none\_detected
\item Validation required: preserve as archive evidence; no live target validation
\item Recommended next action: hold\_or\_review
\item Claim: The final strategic direction before this preservation request was to run a ?-series: snapshot intake, reality extraction, blueprint reconciliation, foundation readiness, and final prompt synthesis. This is needed because plans must now be mapped onto current repo reality rather than executed abstractly. The next chat should pick up there.
\end{itemize}

\subsubsection{PROMOTE-0086 - insufficient\_evidence}

\begin{itemize}
\item Source conversation: omega\_xi\_pi\_architecture\_future
\item Proposed target path: none
\item Current authority support: not yet verified against target docs
\item Current authority conflict: none detected by triage
\item Queue risk: none\_detected
\item Validation required: preserve as archive evidence; no live target validation
\item Recommended next action: hold\_or\_review
\item Claim: The user later reported Xi completion. The enduring lesson is that prompt instructions are not enough; architecture must be machine-readable and enforced.
\end{itemize}

\subsubsection{PROMOTE-0087 - insufficient\_evidence}

\begin{itemize}
\item Source conversation: omega\_xi\_pi\_architecture\_future
\item Proposed target path: none
\item Current authority support: not yet verified against target docs
\item Current authority conflict: none detected by triage
\item Queue risk: none\_detected
\item Validation required: preserve as archive evidence; no live target validation
\item Recommended next action: hold\_or\_review
\item Claim: 4. Preserve all plans, tasks, constraints, risks, decisions, artifacts, and future directions across chats.
\end{itemize}

\subsubsection{PROMOTE-0088 - architecture\_doc\_candidate}

\begin{itemize}
\item Source conversation: OS\_Interface\_Repo\_Architecture
\item Proposed target path: docs/architecture/** after target selection
\item Current authority support: not yet verified against target docs
\item Current authority conflict: none detected by triage
\item Queue risk: none\_detected
\item Validation required: source check, authority-order check, link check, generated output validator, FAST
\item Recommended next action: create\_microtask
\item Claim: DECISION-01 - CLI mandatory, TUI expected, GUI modular.** This was part of the initial architecture baseline and was repeatedly reaffirmed. It matters because every product must remain operable in recovery/headless/automation contexts. GUI is allowed but not authoritative.
\end{itemize}

\subsubsection{PROMOTE-0089 - contract\_candidate}

\begin{itemize}
\item Source conversation: OS\_Interface\_Repo\_Architecture
\item Proposed target path: contracts/** or docs/reference/contracts/** after authority review
\item Current authority support: not yet verified against target docs
\item Current authority conflict: none detected by triage
\item Queue risk: none\_detected
\item Validation required: authority review plus targeted schema/contract parser and FAST
\item Recommended next action: hold\_or\_review
\item Claim: DECISION-02 - Thin GUI shells over shared contracts.** The chat consistently rejected GUI families becoming separate product architectures. This affects client, launcher, setup, server admin, tools, and Workbench modules.
\end{itemize}

\subsubsection{PROMOTE-0090 - contract\_candidate}

\begin{itemize}
\item Source conversation: OS\_Interface\_Repo\_Architecture
\item Proposed target path: contracts/** or docs/reference/contracts/** after authority review
\item Current authority support: not yet verified against target docs
\item Current authority conflict: none detected by triage
\item Queue risk: none\_detected
\item Validation required: authority review plus targeted schema/contract parser and FAST
\item Recommended next action: hold\_or\_review
\item Claim: DECISION-03 - Repo ownership layout.** The user pushed for repository convergence; the discussion established that folders should map to ownership and contract boundaries. This decision is final enough to guide work, but details remain subject to machine-readable layout contracts.
\end{itemize}

\subsubsection{PROMOTE-0091 - reject\_as\_noise}

\begin{itemize}
\item Source conversation: platform\_renderer\_api\_plan
\item Proposed target path: none
\item Current authority support: not yet verified against target docs
\item Current authority conflict: none detected by triage
\item Queue risk: none\_detected
\item Validation required: preserve as archive evidence; no live target validation
\item Recommended next action: hold\_or\_review
\item Claim: The chat ended by generating a maximum-fidelity transfer packet, then a downloadable report package, then an in-chat reader version. The main thing to remember is this: **the active artifact is the final 9-prompt post-master Codex plan, but it is a plan, not evidence that the code exists. The repo must be inspected before execution.
\end{itemize}

\subsubsection{PROMOTE-0092 - insufficient\_evidence}

\begin{itemize}
\item Source conversation: platform\_renderer\_api\_plan
\item Proposed target path: none
\item Current authority support: not yet verified against target docs
\item Current authority conflict: none detected by triage
\item Queue risk: none\_detected
\item Validation required: preserve as archive evidence; no live target validation
\item Recommended next action: hold\_or\_review
\item Claim: These ideas became the architectural heart of the final plan.
\end{itemize}

\subsubsection{PROMOTE-0093 - insufficient\_evidence}

\begin{itemize}
\item Source conversation: platform\_renderer\_api\_plan
\item Proposed target path: none
\item Current authority support: not yet verified against target docs
\item Current authority conflict: none detected by triage
\item Queue risk: none\_detected
\item Validation required: preserve as archive evidence; no live target validation
\item Recommended next action: hold\_or\_review
\item Claim: Stable ABI boundaries must use C ABI, POD structs, explicit function tables, and u32 abi\_version plus u32 struct\_size first.
\end{itemize}

\subsubsection{PROMOTE-0094 - insufficient\_evidence}

\begin{itemize}
\item Source conversation: platform\_support
\item Proposed target path: none
\item Current authority support: not yet verified against target docs
\item Current authority conflict: none detected by triage
\item Queue risk: none\_detected
\item Validation required: preserve as archive evidence; no live target validation
\item Recommended next action: hold\_or\_review
\item Claim: FACT: That early answer was generated by the assistant, not stated by the user as a final decision. It matters because it introduced a very broad portability ambition, but it was later superseded by more practical platform planning. It should now be treated as historical context or possible research scope, not as a controlling product commitment.
\end{itemize}

\subsubsection{PROMOTE-0095 - insufficient\_evidence}

\begin{itemize}
\item Source conversation: platform\_support
\item Proposed target path: none
\item Current authority support: not yet verified against target docs
\item Current authority conflict: none detected by triage
\item Queue risk: none\_detected
\item Validation required: preserve as archive evidence; no live target validation
\item Recommended next action: hold\_or\_review
\item Claim: The user then supplied a detailed inventory covering PlayStation, Xbox, Nintendo, PC handhelds, Android devices, Apple devices, Web platforms, AR, VR, and cross-cutting software targets. This inventory became the central artifact of the chat.
\end{itemize}

\subsubsection{PROMOTE-0096 - insufficient\_evidence}

\begin{itemize}
\item Source conversation: platform\_support
\item Proposed target path: none
\item Current authority support: not yet verified against target docs
\item Current authority conflict: none detected by triage
\item Queue risk: none\_detected
\item Validation required: preserve as archive evidence; no live target validation
\item Recommended next action: hold\_or\_review
\item Claim: FACT: The user's list included legacy and current consoles, handhelds, firmware/OS names, Android software variants, Apple OSes, Web runtimes, AR/VR hardware, XR platforms, and graphics/runtime APIs. FACT: It was a list of desired coverage or consideration, not a final statement that every item must receive full parity support.
\end{itemize}

\subsubsection{PROMOTE-0097 - needs\_user\_decision}

\begin{itemize}
\item Source conversation: Portability\_Assurance\_Future\_Proof
\item Proposed target path: docs/archive/conversations/\_decision/DECISION\_DOCKET\_v0.md
\item Current authority support: not yet verified against target docs
\item Current authority conflict: none detected by triage
\item Queue risk: none\_detected
\item Validation required: preserve as archive evidence; no live target validation
\item Recommended next action: hold\_or\_review
\item Claim: The actual repo was not inspected. The exact accepted directory structure, API policy, DDAP profile, compatibility promises, and first pilot module remain unresolved.
\end{itemize}

\subsubsection{PROMOTE-0098 - insufficient\_evidence}

\begin{itemize}
\item Source conversation: Portability\_Assurance\_Future\_Proof
\item Proposed target path: none
\item Current authority support: not yet verified against target docs
\item Current authority conflict: none detected by triage
\item Queue risk: none\_detected
\item Validation required: preserve as archive evidence; no live target validation
\item Recommended next action: hold\_or\_review
\item Claim: Status: Assistant recommendation; not separately accepted after answer.
\end{itemize}

\subsubsection{PROMOTE-0099 - insufficient\_evidence}

\begin{itemize}
\item Source conversation: Portability\_Assurance\_Future\_Proof
\item Proposed target path: none
\item Current authority support: not yet verified against target docs
\item Current authority conflict: none detected by triage
\item Queue risk: none\_detected
\item Validation required: preserve as archive evidence; no live target validation
\item Recommended next action: hold\_or\_review
\item Claim: Rationale: Replacement must be objectively testable.
\end{itemize}

\subsubsection{PROMOTE-0100 - docs\_clarification\_candidate}

\begin{itemize}
\item Source conversation: readme\_ports\_determinism
\item Proposed target path: README.md or docs/** after target selection
\item Current authority support: not yet verified against target docs
\item Current authority conflict: none detected by triage
\item Queue risk: none\_detected
\item Validation required: source check, authority-order check, link check, generated output validator, FAST
\item Recommended next action: create\_microtask
\item Claim: The user pasted the final README after those cleanup changes. At that point, the README had the current active form: deterministic constraints clarified, ports unified under one source hierarchy, /ports metadata-only if retained, Section 9 normative, lockstep canonical, content-lock mismatch fatal, and disk format versions immutable.
\end{itemize}

\subsubsection{PROMOTE-0101 - reject\_as\_noise}

\begin{itemize}
\item Source conversation: readme\_ports\_determinism
\item Proposed target path: none
\item Current authority support: not yet verified against target docs
\item Current authority conflict: none detected by triage
\item Queue risk: none\_detected
\item Validation required: preserve as archive evidence; no live target validation
\item Recommended next action: hold\_or\_review
\item Claim: After the README work, the user asked for a maximum-fidelity context transfer packet. The assistant produced a long state-transfer document with sections for project inventory, decisions, tasks, constraints, open questions, artifacts, risks, and next actions.
\end{itemize}

\subsubsection{PROMOTE-0102 - architecture\_doc\_candidate}

\begin{itemize}
\item Source conversation: readme\_ports\_determinism
\item Proposed target path: docs/architecture/** after target selection
\item Current authority support: not yet verified against target docs
\item Current authority conflict: none detected by triage
\item Queue risk: none\_detected
\item Validation required: source check, authority-order check, link check, generated output validator, FAST
\item Recommended next action: create\_microtask
\item Claim: The central artifact was the root README.md. The README describes **Domino** as the deterministic engine core and **Dominium** as the official game/tooling/runtime layer. It is written as a high-level architecture document, not a low-level implementation spec.
\end{itemize}

\subsubsection{PROMOTE-0103 - insufficient\_evidence}

\begin{itemize}
\item Source conversation: Refactor\_Architecture
\item Proposed target path: none
\item Current authority support: not yet verified against target docs
\item Current authority conflict: none detected by triage
\item Queue risk: none\_detected
\item Validation required: preserve as archive evidence; no live target validation
\item Recommended next action: hold\_or\_review
\item Claim: The most important thing to remember is that this chat did not implement the refactor. It designed the architecture and generated prompts and handoff material. Any future assistant must verify actual repository state before assuming files were moved, prompts were applied, or code was updated.
\end{itemize}

\subsubsection{PROMOTE-0104 - archive\_only}

\begin{itemize}
\item Source conversation: Refactor\_Architecture
\item Proposed target path: docs/archive/conversations/**
\item Current authority support: not yet verified against target docs
\item Current authority conflict: none detected by triage
\item Queue risk: none\_detected
\item Validation required: preserve as archive evidence; no live target validation
\item Recommended next action: hold\_or\_review
\item Claim: FACT:** This early discussion created the future UI/packs design context.
\end{itemize}

\subsubsection{PROMOTE-0105 - architecture\_doc\_candidate}

\begin{itemize}
\item Source conversation: Refactor\_Architecture
\item Proposed target path: docs/architecture/** after target selection
\item Current authority support: not yet verified against target docs
\item Current authority conflict: none detected by triage
\item Queue risk: none\_detected
\item Validation required: source check, authority-order check, link check, generated output validator, FAST
\item Recommended next action: create\_microtask
\item Claim: FACT:** All products should use shared dsys and dgfx. No product-specific platform or renderer code path.
\end{itemize}

\subsubsection{PROMOTE-0106 - insufficient\_evidence}

\begin{itemize}
\item Source conversation: Refactor\_Control\_Plane
\item Proposed target path: none
\item Current authority support: not yet verified against target docs
\item Current authority conflict: none detected by triage
\item Queue risk: none\_detected
\item Validation required: preserve as archive evidence; no live target validation
\item Recommended next action: hold\_or\_review
\item Claim: The user then asked about optimizing prompts for GPT-5.3/GPT-5.4 and later supplied advice about harness engineering. That advice argued that teams succeed by engineering the environment around agents: tools, docs, linters, feedback loops, memory, and observability. The conversation accepted the core point: the model is not the whole system; the harness is the multiplier.
\end{itemize}

\subsubsection{PROMOTE-0107 - insufficient\_evidence}

\begin{itemize}
\item Source conversation: Refactor\_Control\_Plane
\item Proposed target path: none
\item Current authority support: not yet verified against target docs
\item Current authority conflict: none detected by triage
\item Queue risk: none\_detected
\item Validation required: preserve as archive evidence; no live target validation
\item Recommended next action: hold\_or\_review
\item Claim: The Grug Brained Developer became a design filter: avoid complexity, avoid premature abstraction, do small safe refactors, prefer integration tests, invest in logging, and respect existing structures before tearing them down. This influenced the decision to narrow AIDE and not overbuild runtime/hosts too early.
\end{itemize}

\subsubsection{PROMOTE-0108 - archive\_only}

\begin{itemize}
\item Source conversation: Refactor\_Control\_Plane
\item Proposed target path: docs/archive/conversations/**
\item Current authority support: not yet verified against target docs
\item Current authority conflict: none detected by triage
\item Queue risk: none\_detected
\item Validation required: preserve as archive evidence; no live target validation
\item Recommended next action: hold\_or\_review
\item Claim: The user eventually named AIDE: Automated Integrated Development Environment. The accepted split became: XStack remains Dominium's strict local governance/proof profile; AIDE becomes the portable public layer. AIDE would live in its own repo, be usable as a standalone repo pack, later as CLI/app/extensions, and eventually perhaps a runtime/service/host system.
\end{itemize}

\subsubsection{PROMOTE-0109 - insufficient\_evidence}

\begin{itemize}
\item Source conversation: Release\_Identity\_and\_Versioning
\item Proposed target path: none
\item Current authority support: not yet verified against target docs
\item Current authority conflict: none detected by triage
\item Queue risk: none\_detected
\item Validation required: preserve as archive evidence; no live target validation
\item Recommended next action: hold\_or\_review
\item Claim: The user then objected to versioning policy drift in real products and expressed concern about SemVer's "1.x forever" failure mode. The first proposed answer was a four-part public version scheme, GEN.EPOCH.FEATURE.PATCH, meant to avoid fake major bumps. That idea was useful as an intermediate step but later became less central because XStack already has GBN for dense build history and separate build identity.
\end{itemize}

\subsubsection{PROMOTE-0110 - insufficient\_evidence}

\begin{itemize}
\item Source conversation: Release\_Identity\_and\_Versioning
\item Proposed target path: none
\item Current authority support: not yet verified against target docs
\item Current authority conflict: none detected by triage
\item Queue risk: none\_detected
\item Validation required: preserve as archive evidence; no live target validation
\item Recommended next action: hold\_or\_review
\item Claim: The conversation then moved into how this might fit XStack. The model shifted from "better public version number" to "layered identity model." The assistant recommended preserving per-product versions, a global build number, build IDs, compatibility versions, and suite versions as separate layers. This became the foundation for later decisions.
\end{itemize}

\subsubsection{PROMOTE-0111 - insufficient\_evidence}

\begin{itemize}
\item Source conversation: Release\_Identity\_and\_Versioning
\item Proposed target path: none
\item Current authority support: not yet verified against target docs
\item Current authority conflict: none detected by triage
\item Queue risk: none\_detected
\item Validation required: preserve as archive evidence; no live target validation
\item Recommended next action: hold\_or\_review
\item Claim: The user suggested that suites might use a separate consumer-facing or marketing version, while each component could use stricter SemVer. The chat accepted this as a mature pattern: suites are curated bundles, while components may have technical versions. The model was refined to avoid synchronized fake versioning and to avoid treating a suite major version as universal breakage.
\end{itemize}

\subsubsection{PROMOTE-0112 - architecture\_doc\_candidate}

\begin{itemize}
\item Source conversation: testx\_repox\_governance
\item Proposed target path: docs/architecture/** after target selection
\item Current authority support: not yet verified against target docs
\item Current authority conflict: none detected by triage
\item Queue risk: none\_detected
\item Validation required: source check, authority-order check, link check, generated output validator, FAST
\item Recommended next action: create\_microtask
\item Claim: FACT:** The user is building Dominium / Domino as a deterministic universe engine + game, not a conventional game project. The visible chat repeatedly framed the project as long-lived, simulation-first, deterministic, modular, and designed to survive across many operating systems, toolchains, renderers, products, and distribution models.
\end{itemize}

\subsubsection{PROMOTE-0113 - needs\_user\_decision}

\begin{itemize}
\item Source conversation: testx\_repox\_governance
\item Proposed target path: docs/archive/conversations/\_decision/DECISION\_DOCKET\_v0.md
\item Current authority support: not yet verified against target docs
\item Current authority conflict: none detected by triage
\item Queue risk: none\_detected
\item Validation required: preserve as archive evidence; no live target validation
\item Recommended next action: hold\_or\_review
\item Claim: The user then asked whether the implementation was industry-accepted and what could be improved. The response framed the approach as closer to game engines, operating systems, and long-lived infrastructure than typical game development. The important conclusion was that the design was not exotic, but it was unusually rigorous for games.
\end{itemize}

\subsubsection{PROMOTE-0114 - architecture\_doc\_candidate}

\begin{itemize}
\item Source conversation: testx\_repox\_governance
\item Proposed target path: docs/architecture/** after target selection
\item Current authority support: not yet verified against target docs
\item Current authority conflict: none detected by triage
\item Queue risk: none\_detected
\item Validation required: source check, authority-order check, link check, generated output validator, FAST
\item Recommended next action: create\_microtask
\item Claim: The user had another chat working on content and systems, and asked for a prompt to inform that chat of everything decided so far. A similar prompt was generated for the applications/platforms/renderers chat. These prompts established authoritative boundaries:
\end{itemize}

\subsubsection{PROMOTE-0115 - architecture\_doc\_candidate}

\begin{itemize}
\item Source conversation: Timekeeping\_and\_2038\_Resilience
\item Proposed target path: docs/architecture/** after target selection
\item Current authority support: not yet verified against target docs
\item Current authority conflict: none detected by triage
\item Queue risk: none\_detected
\item Validation required: source check, authority-order check, link check, generated output validator, FAST
\item Recommended next action: create\_microtask
\item Claim: A future assistant should understand that this chat contributes a time architecture doctrine for Dominium: **ACT is authority; DSYS time is runtime-only; observer clocks are derived; civil/astronomical time is projection-only; wall-clock time must never drive authoritative ordering.
\end{itemize}

\subsubsection{PROMOTE-0116 - insufficient\_evidence}

\begin{itemize}
\item Source conversation: Timekeeping\_and\_2038\_Resilience
\item Proposed target path: none
\item Current authority support: not yet verified against target docs
\item Current authority conflict: none detected by triage
\item Queue risk: none\_detected
\item Validation required: preserve as archive evidence; no live target validation
\item Recommended next action: hold\_or\_review
\item Claim: What remains uncertain is the precise target list of legacy machines/toolchains and how far compatibility must go, especially for 16-bit compilers with weak or absent 64-bit arithmetic.
\end{itemize}

\subsubsection{PROMOTE-0117 - reject\_as\_noise}

\begin{itemize}
\item Source conversation: Timekeeping\_and\_2038\_Resilience
\item Proposed target path: none
\item Current authority support: not yet verified against target docs
\item Current authority conflict: none detected by triage
\item Queue risk: none\_detected
\item Validation required: preserve as archive evidence; no live target validation
\item Recommended next action: hold\_or\_review
\item Claim: The final topic is this export task. The user requested a human-readable report first, followed by registers, spec prep, context-transfer packet, aggregator packet, self-audit, and files. This package is designed to let the chat be read later without reopening the full transcript and to support merger into a larger Dominium spec book.
\end{itemize}

\subsubsection{PROMOTE-0118 - reject\_as\_noise}

\begin{itemize}
\item Source conversation: UE6\_Domino\_Deterministic\_Universe
\item Proposed target path: none
\item Current authority support: not yet verified against target docs
\item Current authority conflict: none detected by triage
\item Queue risk: none\_detected
\item Validation required: preserve as archive evidence; no live target validation
\item Recommended next action: hold\_or\_review
\item Claim: FACT: This package preserves the currently visible chat about whether Dominium should use UE6, UE5, Domino, Unreal, or a custom engine layer for a highly ambitious deterministic solar-system-scale game. It also preserves the current uploaded prompt as an artifact because it defines this preservation task.
\end{itemize}

\subsubsection{PROMOTE-0119 - insufficient\_evidence}

\begin{itemize}
\item Source conversation: UE6\_Domino\_Deterministic\_Universe
\item Proposed target path: none
\item Current authority support: not yet verified against target docs
\item Current authority conflict: none detected by triage
\item Queue risk: none\_detected
\item Validation required: preserve as archive evidence; no live target validation
\item Recommended next action: hold\_or\_review
\item Claim: That response also introduced a rule that shaped the rest of the chat: Unreal should render and host Dominium, but Unreal should not define Dominium. That was the first major conceptual decision point. It reframed Dominium not as an engine project in the narrow sense, but as a portable game/simulation system with replaceable frontends.
\end{itemize}

\subsubsection{PROMOTE-0120 - blocked\_by\_queue}

\begin{itemize}
\item Source conversation: UE6\_Domino\_Deterministic\_Universe
\item Proposed target path: none
\item Current authority support: not yet verified against target docs
\item Current authority conflict: blocked queue scope
\item Queue risk: blocked
\item Validation required: preserve as archive evidence; no live target validation
\item Recommended next action: hold\_or\_review
\item Claim: Conclusion: machine gameplay must be designed as a scalable graph simulation, not as unrestricted physics.
\end{itemize}

\subsubsection{PROMOTE-0121 - needs\_user\_decision}

\begin{itemize}
\item Source conversation: ui\_editor\_tool\_editor\_planning
\item Proposed target path: docs/archive/conversations/\_decision/DECISION\_DOCKET\_v0.md
\item Current authority support: not yet verified against target docs
\item Current authority conflict: none detected by triage
\item Queue risk: none\_detected
\item Validation required: preserve as archive evidence; no live target validation
\item Recommended next action: hold\_or\_review
\item Claim: The most important thing to remember is that this chat was a planning and prompt-generation chat, not proof of implementation. **UNCERTAIN / UNVERIFIED:** No generated prompt is evidence that repository code was actually changed. The next assistant must first verify whether the prompts were executed, inspect the repo/source bundle if available, and avoid treating planned files as existing files.
\end{itemize}

\subsubsection{PROMOTE-0122 - insufficient\_evidence}

\begin{itemize}
\item Source conversation: ui\_editor\_tool\_editor\_planning
\item Proposed target path: none
\item Current authority support: not yet verified against target docs
\item Current authority conflict: none detected by triage
\item Queue risk: none\_detected
\item Validation required: preserve as archive evidence; no live target validation
\item Recommended next action: hold\_or\_review
\item Claim: This was one of the most important structure decisions in the chat. It prevented the initial implementation from becoming too broad while preserving the long-term goal. The UI Editor would be a bootstrap tool, not the final architecture.
\end{itemize}

\subsubsection{PROMOTE-0123 - needs\_user\_decision}

\begin{itemize}
\item Source conversation: ui\_editor\_tool\_editor\_planning
\item Proposed target path: docs/archive/conversations/\_decision/DECISION\_DOCKET\_v0.md
\item Current authority support: not yet verified against target docs
\item Current authority conflict: none detected by triage
\item Queue risk: none\_detected
\item Validation required: preserve as archive evidence; no live target validation
\item Recommended next action: hold\_or\_review
\item Claim: The user uploaded screenshot bundles for setup and launcher. **UNCERTAIN / UNVERIFIED:** The screenshots were not inspected in this chat. The assistant still produced a logical layout description, and the user accepted it as useful. Future work should inspect the bundles before claiming exact screenshot fidelity.
\end{itemize}

\subsubsection{PROMOTE-0124 - archive\_only}

\begin{itemize}
\item Source conversation: Universe\_Explorer\_Planning
\item Proposed target path: docs/archive/conversations/**
\item Current authority support: not yet verified against target docs
\item Current authority conflict: none detected by triage
\item Queue risk: none\_detected
\item Validation required: preserve as archive evidence; no live target validation
\item Recommended next action: hold\_or\_review
\item Claim: The pasted discussion on trees, screws, pottery, axes, chairs, tables, and machines established that the player should manipulate material, geometry, constraints, processes, tools, stations, and affordances. The conclusion is that "item classes" must not be the truth substrate. Recipes and blueprints can exist, but as higher-order formalizations.
\end{itemize}

\subsubsection{PROMOTE-0125 - needs\_user\_decision}

\begin{itemize}
\item Source conversation: Universe\_Explorer\_Planning
\item Proposed target path: docs/archive/conversations/\_decision/DECISION\_DOCKET\_v0.md
\item Current authority support: not yet verified against target docs
\item Current authority conflict: none detected by triage
\item Queue risk: none\_detected
\item Validation required: preserve as archive evidence; no live target validation
\item Recommended next action: hold\_or\_review
\item Claim: A major theme was that useful local inventions must become portable, standardized, industrialized, and institutionally adopted. The repo now has a Formalization Chain spec and Player Desire Acceptance Map that strongly preserve this. The future work is making this playable: drafting, measuring, testing, blueprinting, certifying, teaching, manufacturing, maintaining, and revising designs.
\end{itemize}

\subsubsection{PROMOTE-0126 - contract\_candidate}

\begin{itemize}
\item Source conversation: Universe\_Explorer\_Planning
\item Proposed target path: contracts/** or docs/reference/contracts/** after authority review
\item Current authority support: not yet verified against target docs
\item Current authority conflict: none detected by triage
\item Queue risk: none\_detected
\item Validation required: authority review plus targeted schema/contract parser and FAST
\item Recommended next action: hold\_or\_review
\item Claim: Workbench was treated as the likely first surface for the Universe Explorer, but only if it remains projection, not authority. The repo queue says PRESENTATION-CONTRACT-01 is next, with PROJECTION-CONFORMANCE-01 as alternate. The conclusion was that presentation/projection contracts must come before building UI or renderer features.
\end{itemize}

\subsubsection{PROMOTE-0127 - insufficient\_evidence}

\begin{itemize}
\item Source conversation: Workbench\_AIDE\_Product\_Spine
\item Proposed target path: none
\item Current authority support: not yet verified against target docs
\item Current authority conflict: none detected by triage
\item Queue risk: none\_detected
\item Validation required: preserve as archive evidence; no live target validation
\item Recommended next action: hold\_or\_review
\item Claim: The user noted that earlier prompts were hard to copy because they were not always contained in one code block, and that prompts should better handle dirty worktrees and concurrent tasks. From then on, generated prompts included detailed dirty worktree handling, allowed/forbidden paths, non-goals, validation, blocker classification, and commit/final-response formats.
\end{itemize}

\subsubsection{PROMOTE-0128 - insufficient\_evidence}

\begin{itemize}
\item Source conversation: Workbench\_AIDE\_Product\_Spine
\item Proposed target path: none
\item Current authority support: not yet verified against target docs
\item Current authority conflict: none detected by triage
\item Queue risk: none\_detected
\item Validation required: preserve as archive evidence; no live target validation
\item Recommended next action: hold\_or\_review
\item Claim: The final recommended sequence was: finish replay proof and barebones client, run product spine review, then begin limited parallel dev. Larger parallel waves should wait until minimum AIDE workflow law exists.
\end{itemize}

\subsubsection{PROMOTE-0129 - reject\_as\_noise}

\begin{itemize}
\item Source conversation: Workbench\_AIDE\_Product\_Spine
\item Proposed target path: none
\item Current authority support: not yet verified against target docs
\item Current authority conflict: none detected by triage
\item Queue risk: none\_detected
\item Validation required: preserve as archive evidence; no live target validation
\item Recommended next action: hold\_or\_review
\item Claim: The user then requested this full preservation package so a new chat could continue without re-explaining everything. This final handoff is itself part of the preservation output.
\end{itemize}

\subsubsection{PROMOTE-0130 - insufficient\_evidence}

\begin{itemize}
\item Source conversation: World\_Architecture
\item Proposed target path: none
\item Current authority support: not yet verified against target docs
\item Current authority conflict: none detected by triage
\item Queue risk: none\_detected
\item Validation required: preserve as archive evidence; no live target validation
\item Recommended next action: hold\_or\_review
\item Claim: The future relevance of this chat is high. It should feed directly into the future project spec book and into a corrected implementation prompt. But it should not be treated as final implementation detail everywhere. Many high-level decisions are settled, but solver formulas, exact file encoding details, build system integration, actual repository layout, and some numerical representations still require verification.
\end{itemize}

\subsubsection{PROMOTE-0131 - architecture\_doc\_candidate}

\begin{itemize}
\item Source conversation: World\_Architecture
\item Proposed target path: docs/architecture/** after target selection
\item Current authority support: not yet verified against target docs
\item Current authority conflict: none detected by triage
\item Queue risk: none\_detected
\item Validation required: source check, authority-order check, link check, generated output validator, FAST
\item Recommended next action: create\_microtask
\item Claim: Before reading the rest of this report, remember this: this chat's central contribution is not one isolated feature. It is a whole design philosophy for Dominium's world engine. The game should be deterministic, fixed-point, sparse, modular, content-driven, and field-based. Chunks are not the world. Chunks are just how the world is cached, meshed, streamed, and saved.
\end{itemize}

\subsubsection{PROMOTE-0132 - insufficient\_evidence}

\begin{itemize}
\item Source conversation: World\_Architecture
\item Proposed target path: none
\item Current authority support: not yet verified against target docs
\item Current authority conflict: none detected by triage
\item Queue risk: none\_detected
\item Validation required: preserve as archive evidence; no live target validation
\item Recommended next action: hold\_or\_review
\item Claim: This comparison was not a final technical dependency, but it helped clarify the design goal: Dominium should not simply copy any one of these systems. It should combine their useful elements while avoiding their limits.
\end{itemize}

\subsubsection{PROMOTE-0133 - insufficient\_evidence}

\begin{itemize}
\item Source conversation: xstack\_lab\_galaxy
\item Proposed target path: none
\item Current authority support: not yet verified against target docs
\item Current authority conflict: none detected by triage
\item Queue risk: none\_detected
\item Validation required: preserve as archive evidence; no live target validation
\item Recommended next action: hold\_or\_review
\item Claim: and XStack must serve development rather than development serving XStack.
\end{itemize}

\subsubsection{PROMOTE-0134 - reject\_as\_noise}

\begin{itemize}
\item Source conversation: xstack\_lab\_galaxy
\item Proposed target path: none
\item Current authority support: not yet verified against target docs
\item Current authority conflict: none detected by triage
\item Queue risk: none\_detected
\item Validation required: preserve as archive evidence; no live target validation
\item Recommended next action: hold\_or\_review
\item Claim: After the transcript appeared, the user asked for a maximum-fidelity Context Transfer Packet. The assistant produced one, with sections such as workstreams, decisions, tasks, constraints, open questions, risks, artifacts, and verification queue. Then the user asked to turn that into a downloadable report package. The assistant created Markdown/YAML files and a ZIP archive.
\end{itemize}

\subsubsection{PROMOTE-0135 - reject\_as\_noise}

\begin{itemize}
\item Source conversation: xstack\_lab\_galaxy
\item Proposed target path: none
\item Current authority support: not yet verified against target docs
\item Current authority conflict: none detected by triage
\item Queue risk: none\_detected
\item Validation required: preserve as archive evidence; no live target validation
\item Recommended next action: hold\_or\_review
\item Claim: Finally, the user asked not for another package, but for an in-chat reader version of the package. The assistant rendered a structured in-chat reader. The user then asked for this current response: a human-readable, intuitive narrative report.
\end{itemize}



{\small\raggedright SOURCE PATH: docs/\allowbreak{}archive/\allowbreak{}conversations/\allowbreak{}\_\allowbreak{}promotion/\allowbreak{}PROMOTION\_\allowbreak{}TARGET\_\allowbreak{}MAP\_\allowbreak{}v0.md\par}


\section{Promotion Target Map v0}

Target paths are candidate destinations only. They do not authorize edits.



\begin{quote}\small Dense table summarized for the reader edition. See the source file, HTML output, or reference appendix source for full detail.\end{quote}

\begin{quote}\small Dense table content is summarized in this PDF rendering. See the Markdown source and HTML book for the full table.\end{quote}



{\small\raggedright SOURCE PATH: docs/\allowbreak{}archive/\allowbreak{}conversations/\allowbreak{}\_\allowbreak{}promotion/\allowbreak{}PROMOTION\_\allowbreak{}NOT\_\allowbreak{}READY\_\allowbreak{}v0.md\par}


\section{Promotion Not Ready v0}

These candidates are not ready for Wave 1 because they are blocked, noisy, stale, under-evidenced, or require user decisions.


\begin{itemize}
\item PROMOTE-0001 needs\_user\_decision advanced\_simulation\_infrastructure: The claim frames unresolved future work or a decision boundary.
\item PROMOTE-0002 reject\_as\_noise advanced\_simulation\_infrastructure: The candidate describes preservation or archival-process mechanics rather than a clean project claim.
\item PROMOTE-0005 blocked\_by\_queue app\_runtime\_platform\_renderers: The claim touches currently blocked scope: gameplay.
\item PROMOTE-0006 needs\_user\_decision app\_runtime\_platform\_renderers: The claim frames unresolved future work or a decision boundary.
\item PROMOTE-0007 schema\_candidate app\_testx\_codehygiene: The claim appears compatible with current repo doctrine but has not been checked against target docs.
\item PROMOTE-0008 reject\_as\_noise app\_testx\_codehygiene: The candidate describes preservation or archival-process mechanics rather than a clean project claim.
\item PROMOTE-0009 contract\_candidate app\_testx\_codehygiene: The claim appears compatible with current repo doctrine but has not been checked against target docs.
\item PROMOTE-0010 insufficient\_evidence architecture\_codex\_prompts: The candidate is too broad or underspecified for a patch proposal.
\item PROMOTE-0011 insufficient\_evidence architecture\_codex\_prompts: The claim contains an external, language, platform, or baseline term that may be stale.
\item PROMOTE-0012 archive\_only architecture\_codex\_prompts: The claim appears compatible with current repo doctrine but has not been checked against target docs.
\item PROMOTE-0013 insufficient\_evidence architecture\_ui\_providers: The claim contains an external, language, platform, or baseline term that may be stale.
\item PROMOTE-0014 insufficient\_evidence architecture\_ui\_providers: The candidate is too broad or underspecified for a patch proposal.
\item PROMOTE-0015 blocked\_by\_queue architecture\_ui\_providers: The claim touches currently blocked scope: provider\_runtime.
\item PROMOTE-0016 archive\_only Build\_and\_Future\_Proofing: The claim appears compatible with current repo doctrine but has not been checked against target docs.
\item PROMOTE-0017 reject\_as\_noise Build\_and\_Future\_Proofing: The candidate describes preservation or archival-process mechanics rather than a clean project claim.
\item PROMOTE-0018 archive\_only Build\_and\_Future\_Proofing: The claim appears compatible with current repo doctrine but has not been checked against target docs.
\item PROMOTE-0019 insufficient\_evidence canonical\_structure\_and\_framework: The candidate is too broad or underspecified for a patch proposal.
\item PROMOTE-0020 blocked\_by\_queue canonical\_structure\_and\_framework: The claim touches currently blocked scope: provider\_runtime.
\item PROMOTE-0021 schema\_candidate canonical\_structure\_and\_framework: The claim appears compatible with current repo doctrine but has not been checked against target docs.
\item PROMOTE-0022 insufficient\_evidence Chronology\_Celestial\_Systems: The candidate is too broad or underspecified for a patch proposal.
\item PROMOTE-0023 insufficient\_evidence Chronology\_Celestial\_Systems: The candidate is too broad or underspecified for a patch proposal.
\item PROMOTE-0024 needs\_user\_decision Chronology\_Celestial\_Systems: The claim frames unresolved future work or a decision boundary.
\item PROMOTE-0025 needs\_user\_decision development\_routes: The claim frames unresolved future work or a decision boundary.
\item PROMOTE-0026 insufficient\_evidence development\_routes: The candidate is too broad or underspecified for a patch proposal.
\item PROMOTE-0027 archive\_only development\_routes: The claim appears compatible with current repo doctrine but has not been checked against target docs.
\item PROMOTE-0028 needs\_user\_decision documentation\_standards\_readme: The claim frames unresolved future work or a decision boundary.
\item PROMOTE-0029 contract\_candidate documentation\_standards\_readme: The claim appears compatible with current repo doctrine but has not been checked against target docs.
\item PROMOTE-0030 insufficient\_evidence documentation\_standards\_readme: The candidate is too broad or underspecified for a patch proposal.
\item PROMOTE-0031 reject\_as\_noise Dominium\_Architecture\_I: The candidate describes preservation or archival-process mechanics rather than a clean project claim.
\item PROMOTE-0033 insufficient\_evidence Dominium\_Architecture\_I: The candidate is too broad or underspecified for a patch proposal.
\item PROMOTE-0034 reject\_as\_noise Dominium\_Architecture\_II: The candidate describes preservation or archival-process mechanics rather than a clean project claim.
\item PROMOTE-0036 insufficient\_evidence Dominium\_Architecture\_II: The claim contains an external, language, platform, or baseline term that may be stale.
\item PROMOTE-0037 insufficient\_evidence Dominium\_Architecture\_III: The candidate is too broad or underspecified for a patch proposal.
\item PROMOTE-0039 insufficient\_evidence Dominium\_Architecture\_III: The candidate is too broad or underspecified for a patch proposal.
\item PROMOTE-0040 insufficient\_evidence Dominium\_Architecture\_IV: The candidate is too broad or underspecified for a patch proposal.
\item PROMOTE-0041 insufficient\_evidence Dominium\_Architecture\_IV: The candidate is too broad or underspecified for a patch proposal.
\item PROMOTE-0042 insufficient\_evidence Dominium\_Architecture\_IV: The candidate is too broad or underspecified for a patch proposal.
\item PROMOTE-0043 needs\_user\_decision Dominium\_Complete\_Conversation: The claim frames unresolved future work or a decision boundary.
\item PROMOTE-0044 contract\_candidate Dominium\_Complete\_Conversation: The claim appears compatible with current repo doctrine but has not been checked against target docs.
\item PROMOTE-0045 insufficient\_evidence Dominium\_Complete\_Conversation: The candidate is too broad or underspecified for a patch proposal.
\item PROMOTE-0046 insufficient\_evidence dominium\_domino\_codex\_planning: The candidate is too broad or underspecified for a patch proposal.
\item PROMOTE-0047 insufficient\_evidence dominium\_domino\_codex\_planning: The candidate is too broad or underspecified for a patch proposal.
\item PROMOTE-0048 reject\_as\_noise dominium\_domino\_codex\_planning: The candidate describes preservation or archival-process mechanics rather than a clean project claim.
\item PROMOTE-0049 blocked\_by\_queue dominium\_full\_conversation: The claim touches currently blocked scope: provider\_runtime.
\item PROMOTE-0051 insufficient\_evidence dominium\_full\_conversation: The candidate is too broad or underspecified for a patch proposal.
\item PROMOTE-0052 reject\_as\_noise dominium\_setup: The candidate describes preservation or archival-process mechanics rather than a clean project claim.
\item PROMOTE-0053 archive\_only dominium\_setup: The claim appears compatible with current repo doctrine but has not been checked against target docs.
\item PROMOTE-0055 contract\_candidate Domino\_Dominium\_Workbench: The claim appears compatible with current repo doctrine but has not been checked against target docs.
\item PROMOTE-0056 blocked\_by\_queue Domino\_Dominium\_Workbench: The claim touches currently blocked scope: provider\_runtime.
\item PROMOTE-0057 blocked\_by\_queue Domino\_Dominium\_Workbench: The claim touches currently blocked scope: native\_gui.
\item PROMOTE-0058 insufficient\_evidence domino\_engine\_refactor\_prompts: The claim contains an external, language, platform, or baseline term that may be stale.
\item PROMOTE-0059 contract\_candidate domino\_engine\_refactor\_prompts: The claim appears compatible with current repo doctrine but has not been checked against target docs.
\item PROMOTE-0060 reject\_as\_noise domino\_engine\_refactor\_prompts: The candidate describes preservation or archival-process mechanics rather than a clean project claim.
\item PROMOTE-0061 insufficient\_evidence engine\_baseline\_architecture: The candidate is too broad or underspecified for a patch proposal.
\item PROMOTE-0062 reject\_as\_noise engine\_baseline\_architecture: The candidate describes preservation or archival-process mechanics rather than a clean project claim.
\item PROMOTE-0064 insufficient\_evidence Foundation\_Workbench\_Codex: The candidate is too broad or underspecified for a patch proposal.
\item PROMOTE-0065 blocked\_by\_queue Foundation\_Workbench\_Codex: The claim touches currently blocked scope: provider\_runtime.
\item PROMOTE-0066 contract\_candidate Foundation\_Workbench\_Codex: The claim appears compatible with current repo doctrine but has not been checked against target docs.
\item PROMOTE-0067 reject\_as\_noise Framework\_Open\_Source\_Provider: The candidate describes preservation or archival-process mechanics rather than a clean project claim.
\item PROMOTE-0068 insufficient\_evidence Framework\_Open\_Source\_Provider: The candidate is too broad or underspecified for a patch proposal.
\item PROMOTE-0069 blocked\_by\_queue Framework\_Open\_Source\_Provider: The claim touches currently blocked scope: provider\_runtime.
\item PROMOTE-0070 insufficient\_evidence gui\_binary\_content: The candidate is too broad or underspecified for a patch proposal.
\item PROMOTE-0071 insufficient\_evidence gui\_binary\_content: The claim contains an external, language, platform, or baseline term that may be stale.
\item PROMOTE-0072 insufficient\_evidence gui\_binary\_content: The candidate is too broad or underspecified for a patch proposal.
\item PROMOTE-0074 contract\_candidate Language\_Platform\_Architecture: The claim appears compatible with current repo doctrine but has not been checked against target docs.
\item PROMOTE-0075 blocked\_by\_queue Language\_Platform\_Architecture: The claim touches currently blocked scope: provider\_runtime.
\item PROMOTE-0076 needs\_user\_decision launcher\_app\_layer: The claim frames unresolved future work or a decision boundary.
\item PROMOTE-0077 archive\_only launcher\_app\_layer: The claim appears compatible with current repo doctrine but has not been checked against target docs.
\item PROMOTE-0078 schema\_candidate launcher\_app\_layer: The claim appears compatible with current repo doctrine but has not been checked against target docs.
\item PROMOTE-0079 reject\_as\_noise Launcher\_Setup\_Architecture: The candidate describes preservation or archival-process mechanics rather than a clean project claim.
\item PROMOTE-0080 insufficient\_evidence Launcher\_Setup\_Architecture: The candidate is too broad or underspecified for a patch proposal.
\item PROMOTE-0081 contract\_candidate Launcher\_Setup\_Architecture: The claim appears compatible with current repo doctrine but has not been checked against target docs.
\item PROMOTE-0082 reject\_as\_noise Modularity\_AIDE\_Refactorability: The candidate describes preservation or archival-process mechanics rather than a clean project claim.
\item PROMOTE-0083 insufficient\_evidence Modularity\_AIDE\_Refactorability: The candidate is too broad or underspecified for a patch proposal.
\item PROMOTE-0085 insufficient\_evidence omega\_xi\_pi\_architecture\_future: The candidate is too broad or underspecified for a patch proposal.
\item PROMOTE-0086 insufficient\_evidence omega\_xi\_pi\_architecture\_future: The candidate is too broad or underspecified for a patch proposal.
\item PROMOTE-0087 insufficient\_evidence omega\_xi\_pi\_architecture\_future: The candidate is too broad or underspecified for a patch proposal.
\item PROMOTE-0089 contract\_candidate OS\_Interface\_Repo\_Architecture: The claim appears compatible with current repo doctrine but has not been checked against target docs.
\item PROMOTE-0090 contract\_candidate OS\_Interface\_Repo\_Architecture: The claim appears compatible with current repo doctrine but has not been checked against target docs.
\item PROMOTE-0091 reject\_as\_noise platform\_renderer\_api\_plan: The candidate describes preservation or archival-process mechanics rather than a clean project claim.
\item PROMOTE-0092 insufficient\_evidence platform\_renderer\_api\_plan: The candidate is too broad or underspecified for a patch proposal.
\item PROMOTE-0093 insufficient\_evidence platform\_renderer\_api\_plan: The candidate is too broad or underspecified for a patch proposal.
\item PROMOTE-0094 insufficient\_evidence platform\_support: The candidate is too broad or underspecified for a patch proposal.
\item PROMOTE-0095 insufficient\_evidence platform\_support: The claim contains an external, language, platform, or baseline term that may be stale.
\item PROMOTE-0096 insufficient\_evidence platform\_support: The claim contains an external, language, platform, or baseline term that may be stale.
\item PROMOTE-0097 needs\_user\_decision Portability\_Assurance\_Future\_Proof: The claim frames unresolved future work or a decision boundary.
\item PROMOTE-0098 insufficient\_evidence Portability\_Assurance\_Future\_Proof: The candidate is too broad or underspecified for a patch proposal.
\item PROMOTE-0099 insufficient\_evidence Portability\_Assurance\_Future\_Proof: The candidate is too broad or underspecified for a patch proposal.
\item PROMOTE-0101 reject\_as\_noise readme\_ports\_determinism: The candidate describes preservation or archival-process mechanics rather than a clean project claim.
\item PROMOTE-0103 insufficient\_evidence Refactor\_Architecture: The candidate is too broad or underspecified for a patch proposal.
\item PROMOTE-0104 archive\_only Refactor\_Architecture: The claim appears compatible with current repo doctrine but has not been checked against target docs.
\item PROMOTE-0106 insufficient\_evidence Refactor\_Control\_Plane: The candidate is too broad or underspecified for a patch proposal.
\item PROMOTE-0107 insufficient\_evidence Refactor\_Control\_Plane: The candidate is too broad or underspecified for a patch proposal.
\item PROMOTE-0108 archive\_only Refactor\_Control\_Plane: The claim appears compatible with current repo doctrine but has not been checked against target docs.
\item PROMOTE-0109 insufficient\_evidence Release\_Identity\_and\_Versioning: The candidate is too broad or underspecified for a patch proposal.
\item PROMOTE-0110 insufficient\_evidence Release\_Identity\_and\_Versioning: The candidate is too broad or underspecified for a patch proposal.
\item PROMOTE-0111 insufficient\_evidence Release\_Identity\_and\_Versioning: The candidate is too broad or underspecified for a patch proposal.
\item PROMOTE-0113 needs\_user\_decision testx\_repox\_governance: The claim frames unresolved future work or a decision boundary.
\item PROMOTE-0116 insufficient\_evidence Timekeeping\_and\_2038\_Resilience: The candidate is too broad or underspecified for a patch proposal.
\item PROMOTE-0117 reject\_as\_noise Timekeeping\_and\_2038\_Resilience: The candidate describes preservation or archival-process mechanics rather than a clean project claim.
\item PROMOTE-0118 reject\_as\_noise UE6\_Domino\_Deterministic\_Universe: The candidate describes preservation or archival-process mechanics rather than a clean project claim.
\item PROMOTE-0119 insufficient\_evidence UE6\_Domino\_Deterministic\_Universe: The candidate is too broad or underspecified for a patch proposal.
\item PROMOTE-0120 blocked\_by\_queue UE6\_Domino\_Deterministic\_Universe: The claim touches currently blocked scope: gameplay.
\item PROMOTE-0121 needs\_user\_decision ui\_editor\_tool\_editor\_planning: The claim frames unresolved future work or a decision boundary.
\item PROMOTE-0122 insufficient\_evidence ui\_editor\_tool\_editor\_planning: The candidate is too broad or underspecified for a patch proposal.
\item PROMOTE-0123 needs\_user\_decision ui\_editor\_tool\_editor\_planning: The claim frames unresolved future work or a decision boundary.
\item PROMOTE-0124 archive\_only Universe\_Explorer\_Planning: The claim appears compatible with current repo doctrine but has not been checked against target docs.
\item PROMOTE-0125 needs\_user\_decision Universe\_Explorer\_Planning: The claim frames unresolved future work or a decision boundary.
\item PROMOTE-0126 contract\_candidate Universe\_Explorer\_Planning: The claim appears compatible with current repo doctrine but has not been checked against target docs.
\item PROMOTE-0127 insufficient\_evidence Workbench\_AIDE\_Product\_Spine: The candidate is too broad or underspecified for a patch proposal.
\item PROMOTE-0128 insufficient\_evidence Workbench\_AIDE\_Product\_Spine: The candidate is too broad or underspecified for a patch proposal.
\item PROMOTE-0129 reject\_as\_noise Workbench\_AIDE\_Product\_Spine: The candidate describes preservation or archival-process mechanics rather than a clean project claim.
\item PROMOTE-0130 insufficient\_evidence World\_Architecture: The candidate is too broad or underspecified for a patch proposal.
\item PROMOTE-0132 insufficient\_evidence World\_Architecture: The candidate is too broad or underspecified for a patch proposal.
\item PROMOTE-0133 insufficient\_evidence xstack\_lab\_galaxy: The candidate is too broad or underspecified for a patch proposal.
\item PROMOTE-0134 reject\_as\_noise xstack\_lab\_galaxy: The candidate describes preservation or archival-process mechanics rather than a clean project claim.
\item PROMOTE-0135 reject\_as\_noise xstack\_lab\_galaxy: The candidate describes preservation or archival-process mechanics rather than a clean project claim.
\end{itemize}



{\small\raggedright SOURCE PATH: docs/\allowbreak{}archive/\allowbreak{}conversations/\allowbreak{}\_\allowbreak{}promotion/\allowbreak{}PROMOTION\_\allowbreak{}TRIAGE\_\allowbreak{}v0.md\par}


\section{Promotion Triage v0}

This triage reduces the raw promotion queue into review classes. It does not promote claims.


\subsection{Counts}

\begin{itemize}
\item blocked\_by\_current\_queue: 10
\item consistent\_with\_repo\_but\_not\_formalized: 41
\item insufficient\_evidence: 46
\item needs\_user\_decision: 12
\item rejected\_noise: 19
\item stale\_or\_superseded: 7
\end{itemize}

\subsection{Recommended Review Order}

1. Review TOP\_PROMOTION\_CANDIDATES\_v0.md (TOP\_PROMOTION\_CANDIDATES\_v0.md).

2. Review NEEDS\_USER\_DECISION\_v0.md (NEEDS\_USER\_DECISION\_v0.md).

3. Preserve REJECTED\_NOISE\_v0.md (REJECTED\_NOISE\_v0.md) as historical/non-promoted extraction noise.

4. Use CLAIM\_REVIEW\_MATRIX\_v0.md (../\_reconciliation/CLAIM\_REVIEW\_MATRIX\_v0.md) before any live-doc patch task.



\chapter{Contradictions, Staleness, and Verification}


{\small\raggedright SOURCE PATH: docs/\allowbreak{}archive/\allowbreak{}conversations/\allowbreak{}\_\allowbreak{}audit/\allowbreak{}CONTRADICTION\_\allowbreak{}MATRIX.md\par}


\section{Contradiction Matrix}

Findings are review queues, not resolved doctrine.



\begin{quote}\small Dense table summarized for the reader edition. See the source file, HTML output, or reference appendix source for full detail.\end{quote}

\begin{quote}\small Dense table content is summarized in this PDF rendering. See the Markdown source and HTML book for the full table.\end{quote}



{\small\raggedright SOURCE PATH: docs/\allowbreak{}archive/\allowbreak{}conversations/\allowbreak{}\_\allowbreak{}audit/\allowbreak{}STALENESS\_\allowbreak{}AND\_\allowbreak{}VERIFICATION.md\par}


\section{Staleness And Verification}

External platform, SDK, release, language-baseline, and implementation claims are stale until verified.


\subsection{Packages With Explicit Uncertainty}

\subsubsection{advanced\_simulation\_infrastructure}

\begin{itemize}
\item [UNCERTAIN / UNVERIFIED] No repository files were inspected in this chat. No code was implemented. Proposed paths, module names, command names, file names, VM concepts, and data schemas are design proposals until checked against the actual codebase.
\item [UNCERTAIN / UNVERIFIED] Exact solvers, fixed-point formats, module paths, data schemas, and tick ordering remain to be verified and specified.
\item [UNCERTAIN / UNVERIFIED] The exact gameplay rules runtime was not specified. The command examples were conceptual.
\item [UNCERTAIN / UNVERIFIED] The exact cross-section schema, attachment schema, and packing rules remain unresolved.
\item [UNCERTAIN / UNVERIFIED] Microsegment length, VM design, storage layout, budget values, and actual code integration remain open.
\item [UNCERTAIN / UNVERIFIED] Actual standards packs, rule schema, semantic key format, localization, and AI interaction remain unresolved.
\item Then it produced a prompt for another GPT-5.2 chat that already had a refactor/optimization plan. That prompt told the other chat to amend its existing plan instead of restarting, and to verify coverage of arbitrary placement, unified spatial primitives, co-location, signage, buildings, and determinism/performance.
\item [UNCERTAIN / UNVERIFIED] The architecture has not been verified against the actual repository. The exact implementation details remain unresolved: Q formats, orientation math, command schemas, VM instruction set, microsegment length, carrier ownership, junction archetypes, facility modules, and DECOR/device promotion policy.
\end{itemize}

\subsubsection{app\_runtime\_platform\_renderers}

\begin{itemize}
\item What remains uncertain is whether the current repository already enforces these boundaries mechanically. The actual code needs inspection.
\item The unresolved issue is sharding semantics. APP0 requires sharding hooks, but the chat did not decide whether shards are spatial, logical, temporal, jurisdictional, or something else. The report preserved the idea that APP0 can provide sharding plumbing without inventing simulation semantics.
\item The unresolved issue is the trust and integrity model for setup. The chat did not decide whether content verification should use hashes, signatures, local manifests, network verification, or another model.
\item The unresolved issue is the exact elevation/audit model. The chat did not decide how tools authenticate, request write permissions, log actions, or interact with law enforcement modules.
\item The unresolved issue is the exact renderer capability schema. The chat proposed ideas but did not finalize fields, priorities, or target matrices.
\item The user asked to look at the old code snapshot in project attachments. A prior assistant claimed to inspect it, but later preservation documents marked that claim as **UNCERTAIN / UNVERIFIED**.
\item This became one of the most important unresolved issues. Almost every implementation question depends on the real code: directory layout, build system, languages, existing platform/render layers, client/server dependencies, and engine/game public APIs. The next real step is repository verification.
\item INFERENCE:** The user wants future assistants to be epistemically careful: not inventing repository facts, not treating proposals as decisions, and not losing uncertainty.
\end{itemize}

\subsubsection{app\_testx\_codehygiene}

\begin{itemize}
\item The main unresolved goal is practical execution: which prompts have actually been run, which files exist, and what implementation should happen next.
\item Setup handles install, verify, repair, rollback, upgrade, downgrade, uninstall, and status. Launcher can invoke Setup, but must not replicate its mutation logic. This prevents two products from disagreeing about installation state.
\item First, verify whether libs/contracts and libs/appcore already exist. Then implement the pure contract headers and appcore skeleton. This should include:
\item Before implementing app RepoX integration, verify actual RepoX metadata paths, TestX invocation, VALIDATE-0 commands, and BUILD-ID-0 files. Without that, an implementation prompt would guess.
\item Use labels like FACT, INFERENCE, UNCERTAIN, PROJECT-CONTEXT in reports.
\item From the user profile and instructions, future assistants should verify time-sensitive/current facts with web where relevant. This chat itself is mostly internal project planning, so external citations are not central.
\item EXEC, ECSX, KERN, ADOPT, DIST, HWCAPS, EXIST, DOMAIN, TRAVEL, TIME, OMNI, LIFE, CIV, WAR, AGENT, TOOL, MOD, FINAL prompt families were generated earlier. They are important as historical design artifacts, but actual repo execution is unverified.
\item The chat produced downloadable handoff files and later an in-chat reader. Those files are useful for aggregation but are not themselves project source code. The package's main caveat is that it does not verify actual repository state.
\end{itemize}

\subsubsection{architecture\_codex\_prompts}

\begin{itemize}
\item UNCERTAIN / UNVERIFIED: No repository code was inspected or modified in this chat. The prompts and architecture are plans and artifacts, not proof that implementation exists. The actual repository state, build system, existing code quality, GUI backend availability, TLV schemas, and whether any prompts have been applied remain unverified.
\item Uncertainty remains about actual repository code and build system, but the layering itself is a strong established principle.
\item Unresolved issues include exact TLV schema field tags, exact registry implementation, base pack structure, and whether engine comments/docs may use examples with domain words.
\item The main uncertainty is implementation: the prompts define it, but no code was verified.
\item Unresolved details include exact process IO schemas, machine runtime state, and the relationship between autonomous machines and agent-operated jobs.
\item Unresolved: exact nesting policy, allowed packing modes, and how far to model real packing constraints.
\item Uncertainty remains about freeform geometry versus grid shell as the long-term building representation.
\item The exact destruction and structural-load models remain unresolved and part of Path D.
\end{itemize}

\subsubsection{architecture\_ui\_providers}

\begin{itemize}
\item The exact first playable slice remains unresolved. The first provider wedge is clear, but exact sequence between client/workbench gameplay, Robot OS, and gameplay mechanics still needs a formal milestone. The exact Lua version, Linux baseline, current repo status, and external library/version support need verification. The exact degree of Unreal involvement remains unclear.
\item 1. Verify and pin external library/toolchain facts.
\item It remains uncertain how much the user wants Unreal in the near-term after raylib-first discussion. It also remains uncertain whether Lua 5.4 or 5.5 is preferred; the user appears to value pinning and stability more than "latest" for script ABI.
\end{itemize}

\subsubsection{Build\_and\_Future\_Proofing}

\begin{itemize}
\item Plain-language limitation: this package is reliable as a preservation of the current visible chat. It should not be treated as a complete archive of every Dominium discussion in other chats. Where the report uses project or repository context, it is labelled as such. The most important caveat is that many items are recommendations, not final user decisions.
\item The unresolved goals are implementation goals: deciding which recommendations become canon, adding public-surface/dependency/build contracts, validating the current live repo state, and applying the proposed cleanup tasks. The preservation package is complete, but the engineering work it describes remains pending.
\item 2. **Verify current repo state.** Confirm the actual current HEAD, build proof, smoke tests, and outstanding warnings before acting.
\item 10. **Feed this report into the master spec book aggregator.** Preserve uncertainty labels and avoid merging suggestions as decisions.
\item The preservation output must distinguish FACT, INFERENCE, UNCERTAIN/UNVERIFIED, and PROJECT-CONTEXT.
\item A future assistant may rely on stale repo state. Mitigation: verify live HEAD, docs, and CI before implementation.
\item What live repo facts should I verify before implementing these tasks?
\item The final uploaded prompt asked for this preservation package. The key thing to preserve is that most architecture proposals are recommendations, not accepted user decisions yet. The next best action is to decide which recommendations become canon, then implement either STRUCTURE-01: Public Surface Registry or the build tuple contract work, after verifying the current live repo state.
\end{itemize}

\subsubsection{canonical\_structure\_and\_framework}

\begin{itemize}
\item INFERENCE: The user wanted a durable method for deciding future structure changes, not just a one-time layout. They also wanted the assistant to become more rigorous: distinguish decisions from brainstorms, stop over-planning without execution, and preserve uncertainty. The user also wanted a structure that supports other Domino-based games and possibly different engine implementations.
\item A persistent uncertainty is the exact current repo state. The user supplied many status reports and directory exports, but several exports were internally inconsistent or stale at various points. The chat therefore repeatedly emphasized structure report integrity and verification before trusting tree analyses.
\item 1. Verify fast-strict/RepoX/structure report status from a clean current export.
\item The user dislikes vague reassurance and wants uncertainty labelled.
\item UNCERTAIN: The final exact convention for external/upstream versus external/vendor should be checked against current repo policy.
\item UNCERTAIN: Whether Lua should be pinned to 5.4 or 5.5 is not settled in this chat.
\item UNCERTAIN: Whether pack-internal content/ directories are legitimate pack law or legacy remains to be verified.
\item User-reported commits and status summaries: these include commits like 6e0dd93, 1406490, 3243fab, ce9ca, and others. Treat them as FACTs reported in the chat, but verify live repo state before acting.
\end{itemize}

\subsubsection{Chronology\_Celestial\_Systems}

\begin{itemize}
\item The assistant formalized this as the Perfect Earth Calendar, later called HPC-E. The final structure is clear, though the exact leap rule remains unresolved.
\item What remains uncertain is the exact verified object list. The assistant generated a real Milky Way system list, but that list must be checked before becoming data.
\item The unresolved issue is exact fidelity. The chat specified what should exist conceptually, but not the final data schema, map resolution, object attributes, or implementation priority.
\item What remains unresolved is the exact mathematical implementation of these frames and how much relativistic fidelity should exist in phase one.
\item The uncertainty is high relative to Earth. Many names and structures were assistant-generated and should be reviewed before implementation.
\item The unresolved issue is implementation reality. The actual Plan G/The Game codebase was not visible here, so any implementation plan must be reconciled with real architecture.
\item The only major unresolved part is the leap rule.
\item The user wants rigorous, structured, fact-aware assistance. The user specifically requested preservation of uncertainty labels and does not want assistant suggestions treated as user decisions.
\end{itemize}

\subsubsection{development\_routes}

\begin{itemize}
\item The first answer was assertive. It stated that Dominium must follow Route C and described the route as the only viable one for Dominium. Later parts of the chat corrected the status of that claim. FACT: the assistant made that recommendation. UNCERTAIN / UNVERIFIED: the recommendation was not validated with external sources in the chat and was not explicitly accepted by the user.
\item The assistant produced a large Context Transfer Packet. It preserved the technical route proposal, the caveats around its uncertainty, the proposed phase order, the open questions, the risks, the artifacts, and the next actions.
\item The assistant created those files and reported that the chat label was inferred as **Dominium Development Routes**. The package was said to be safe for aggregation **with caveats**. The main caveats were that the substantive project transcript was short, Route C was not user-confirmed, and Dominium's genre, core loop, stack, implementation state, files, platforms, team, and timeline were not established.
\item UNCERTAIN / UNVERIFIED: Route C was not accepted by the user as final.
\item UNCERTAIN / UNVERIFIED: the required level of determinism for Dominium was not established. It might need full cross-platform deterministic replay, or it might only need ordinary save/load correctness, or something between those extremes.
\item The outputs included a Context Transfer Packet and a downloadable package with multiple files. The files are less important than the purpose: preserve the exact status of the discussion, including uncertainty and assistant-proposal status, so future assistants do not accidentally distort the project.
\item The second explicit goal was to preserve the chat before retirement. The user did not want a normal summary. They wanted a maximum-fidelity transfer packet that would prevent loss of decisions, constraints, unresolved issues, rationale, and next actions. This was addressed by the Context Transfer Packet.
\item Some goals changed over time. The chat began as architecture guidance, became a state-transfer task, became a packaging task, became an in-chat inspection task, and now becomes a plain-English briefing task. The underlying continuity goal stayed consistent: preserve the useful substance without losing uncertainty.
\end{itemize}

\subsubsection{documentation\_standards\_readme}

\begin{itemize}
\item The key thing to remember: this chat produced the standards and prompts for future work; it did not verify or change the project. The next assistant must inspect the actual repository before writing repo-specific docs, README content, platform claims, license claims, or subsystem descriptions.
\item The chat began with the user asking how to have Codex generate documentation for every file in the project. The user wanted metadata, functions, definitions, dependencies, and other maintainability information documented. The core uncertainty was placement: should documentation live in separate sister files, or should it live in file headers and function headers?
\item UNCERTAIN / UNVERIFIED: Whether the project has similarly named files or modules is unknown.
\item What remains uncertain is how the actual project currently organizes public headers, internal headers, and docs. That must be verified from the repository.
\item What remains unverified is the actual project name, supported platforms, license, maturity, and current repository layout. Those cannot be responsibly written without inspecting the repo.
\item The decision depends on an assumption: that the project's public APIs are declared in headers. That is likely for a C/C++ codebase, but the actual public/private boundary is still unverified.
\item What must happen first: verify source roots, exclusions, CMake structure, CI environment variables, and Python availability.
\item What could go wrong: treating assistant suggestions as final requirements or merging contradictory chat claims without preserving uncertainty.
\end{itemize}

\subsubsection{Dominium\_Architecture\_I}

\begin{itemize}
\item What remains uncertain is the exact final scope of all simulation systems. Many systems were described in ambitious terms, but not all details are final or internally consistent. The systems layer especially needs canonicalisation before coding.
\item What remains uncertain is the current real-world status and capabilities of "Codex 5.1 Max." That is an external tool fact and requires verification before future use.
\item The key later decision was to skip top-level .txt devspec files because original Markdown files already existed in /docs/.... The visible chat does not show the full content of those docs. They should be preserved as referenced artifacts, but their exact content is **UNCERTAIN / UNVERIFIED**.
\item This affects dlocale, font rendering, platform path handling, data locale JSON files, mod APIs, UI, and documentation. The unresolved part is exactly how much retro support is actually feasible. Claims about old OSes, toolchains, SDL2, OpenGL, Direct3D, and macOS versions require external verification.
\item A large portion of the chat consisted of producing specs for those files. Each spec was meant to tell Codex how to implement that file. The conclusion is that the repository architecture is broad but structured. The unresolved issue is that many later files remain pending.
\item The unresolved issue is API consistency. Generated specs for memory and serialization are inconsistent across files, and some platform/render specs contain outdated or unverified external assumptions.
\item The most important unresolved issue is duplication and contradiction. dweather, dhydro, and dai\_core were generated in more than one form. The package does not choose a final version. A future assistant should not implement any of those blindly.
\item The user decided to skip the top-level .txt files because original Markdown docs already exist in /docs/.... This decision is final for the current workflow, but it depends on those docs actually existing and being current. That is **UNCERTAIN / UNVERIFIED** because the docs were not inspected in this chat.
\end{itemize}

\subsubsection{Dominium\_Architecture\_II}

\begin{itemize}
\item What remains uncertain is the exact code-level API and component layouts. The chat produced specs and prompts, not final source code.
\item This is an important implementation detail, but some of it remains unverified. Lua 5.4.4 portability to old compilers needs testing. The exact Lua data schemas still need definition. The docs also still have JSON-oriented format sections, so the Lua-vs-JSON policy needs reconciliation.
\item The biggest unresolved goals are practical:
\item Verify DX9/Windows 2000, SDL2, Lua/toolchain compatibility.
\item FACT: This was answered in response to Codex's clarifying questions. It is a concrete implementation direction, though Lua 5.4.4 portability remains unverified.
\item 1. **Verify and apply V4 docs.
\item The user repeatedly requested maximum-fidelity reports, registers, handoff packets, and human-readable summaries. The user wants labelled uncertainty and does not want future assistants to invent missing context.
\item or forget that Lua-vs-JSON remains unresolved.
\end{itemize}

\subsubsection{Dominium\_Architecture\_III}

\begin{itemize}
\item The user then uploaded dominium.7z, saying it was the git repository as of that moment. **FACT:** The prior assistant said it could not open .7z archives in that environment. That means the actual repository state remained unverified at that stage. This became important later, because many implementation ideas were discussed without confirmed access to the repo contents.
\item This created one of the main unresolved issues. DX12 was discussed earlier, but omitted from the latest renderer list. X11 and Wayland were discussed earlier, but the latest platform categories are POSIX, SDL1, SDL2, and Native. Therefore, the final status is:
\item DX12 is **UNCERTAIN / UNVERIFIED** and should not be assumed.
\item X11/Wayland placement is unresolved.
\item The prior assistant summarized these files as containing two different launcher concepts: an engine-rendered launcher view and a native Win32 launcher. It also claimed the CMake file built the Win32 source for non-Windows targets. However, this report must preserve that as **UNCERTAIN / UNVERIFIED**, because the files were not re-inspected during final packaging.
\item This is final as a user-stated product goal. Feasibility still requires verification.
\item Those outputs are important artifacts, but they are not the substance of the design. Their purpose is to preserve this chat's decisions, uncertainty, and future work for later aggregation into a master Project Spec Book.
\item What remains uncertain is not the constraint itself, but how existing repository code currently enforces it. That requires code inspection.
\end{itemize}

\subsubsection{Dominium\_Architecture\_IV}

\begin{itemize}
\item The conclusion was clear and accepted: **Domino is the reusable engine; Dominium is the game/product layer**. What remains uncertain is the actual repository state. The chat planned the split, but did not verify whether the repo already matches it.
\item The unresolved issue is that the actual repo was not inspected in this chat. The planned structure may not match the current files.
\item The unresolved issue is practical feasibility, especially for retro targets and Carbon OS. Their actual toolchains and system calls remain unverified.
\item The unresolved issue is the exact IR payload layout and how far limited renderers should go. A CGA or MDA backend cannot realistically support modern 3D in the same way as Vulkan. The agreed approach is explicit capability flags and graceful fallback or no-op behaviour.
\item The unresolved issue is the exact manifest/schema format. JSON, TOML, and INI-like examples appeared in prompts, but no final schema choice was made.
\item The unresolved issue is that external packaging details require verification. WiX, macOS notarization, Linux packaging, and AppImage tooling are all external-world facts and may be stale.
\item The package and this report should help future assistants continue without re-reading the whole chat. However, they must preserve uncertainty: actual repo state is unverified, some transcript sections were skipped, external tooling facts may be stale, and assistant suggestions are not user decisions unless accepted or built upon.
\item It is also reasonable to infer that the user wants the future project spec book to preserve rationale, not just lists of decisions. The later handoff requests emphasised visible rationale, rejected options, uncertainty, and changes of direction.
\end{itemize}

\subsubsection{Dominium\_Complete\_Conversation}

\begin{itemize}
\item Uncertainty: the raw pasted prompt is not the current repo authority unless materialized in repo canon. The later audit found that it had in fact been materialized under docs/canon/constitution\_v1.md.
\item Uncertainty: during this preservation pass, a potential inconsistency appeared: some docs claim product roots moved under apps/, while an earlier inspected code path under root client/ was available and apps/client was not fetched successfully. This must be verified against the current physical tree.
\item Unresolved goals include verifying the latest physical repo tree, proving runtime implementation maturity, formalizing public ABI/API boundaries, completing layout/naming cleanup, and building a second Domino-based product to prove reuse.
\item 1. Verify the current physical repo tree against layout docs and contracts/repo/layout.contract.toml.
\item The user likely wants future assistants to challenge weak framing, preserve uncertainty, avoid shallow "best practice" lists, and focus on decisions that survive future rewrites.
\item "Create a gated plan for verifying current repo layout and exceptions."
\end{itemize}

\subsubsection{dominium\_domino\_codex\_planning}

\begin{itemize}
\item The most important things to remember are: the project's strict architecture constraints, the Milestone-0 prompt sequence, the later shared CLI/TUI/GUI/input/render prompt sequences, the missing pack-system prompt, the unified startup policy, the unresolved repo-verification issue, and the need to treat generated prompts as plans rather than proof of implementation.
\item The assistant generated the first two prompts. The third was never produced because the conversation changed direction. That missing third prompt remains an important unresolved item.
\item However, the later context transfer packet treated that inspection claim as **UNCERTAIN / UNVERIFIED**. This was careful and correct: within the visible chat, there was no reliable tool trace proving the repo had actually been inspected. Therefore, future assistants must verify dominium.zip or the active checkout before relying on the platform/render/API answer.
\item What remains uncertain is how much of this architecture already exists in the current repo and how much is still planned. The chat produced prompts and plans, but it did not prove execution.
\item The unresolved issue is whether those prompts were ever run and whether the generated code, if any, conforms to the project's strict C89/C++98 and determinism requirements.
\item The unresolved complication is C89 portability. Standard C89 does not provide long long, but u64, i64, and q48\_16 require 64-bit representation. Future work must decide whether to allow compiler extensions, emulate 64-bit values for strict retro targets, or define platform tiers with conditional support.
\item The unresolved issue is that the prompt's minimal manifest parser and hardcoded platform path are only a starting point. A real implementation must generalize product discovery, avoid duplicated path logic, and preserve compatibility/versioning rules.
\item The unresolved issue is whether the prompts are too broad and whether the existing repo already has overlapping modules. Future work should inspect the repo before applying these prompts.
\end{itemize}

\subsubsection{dominium\_full\_conversation}

\begin{itemize}
\item Verify live repo state, then review/run FULL-GATE-LEGACY-TEST-ROUTE-01, then generate PACK-INTERNAL-LAYOUT-CANON-01 if the queue/status supports the maintenance lane.
\item Preserve uncertainty labels.
\item Future chats could easily misunderstand this conversation by treating Workbench as a GUI framework, treating raylib as the engine, resuming broad structure cleanup, skipping projection conformance, or assuming live repo status from stale pasted reports. The mitigation is to verify repo state first, preserve the contract/service/projection/provider distinctions, and follow the task queue.
\item The best next action is to verify live repo state, then review or run FULL-GATE-LEGACY-TEST-ROUTE-01, then generate PACK-INTERNAL-LAYOUT-CANON-01 if status supports it.
\end{itemize}

\subsubsection{dominium\_setup}

\begin{itemize}
\item FACT: The setup system was repeatedly defined as the only component allowed to install files, modify system/installation state, own installed-file metadata, repair installs, verify artifacts, uninstall, downgrade, upgrade, and roll back. This is the most important architectural rule from the chat.
\item What remains uncertain is the actual current implementation state in the repository after the applied Codex prompts. The design is clear, but the repo must still be inspected.
\item FACT: The final canonical setup layout uses setup/core/\{fetch,verify,install,rollback\}, setup/include/\{dsk,dsu\}, setup/packages, setup/platform, setup/tests, and setup/ui.
\item The unresolved part is whether the actual repository now matches this canonical layout.
\item UNCERTAIN / UNVERIFIED: Exact Visual Studio 2026 behavior and toolchain details should be verified in the current environment before relying on them.
\item UNCERTAIN / UNVERIFIED: The user stated the prompts were applied, but this chat did not show the final repository tree or build results.
\item UNCERTAIN / UNVERIFIED: The actual schema files under schema/setup/ need to be checked.
\item Several goals remain unresolved because the actual repo was not inspected here. We do not know whether the applied Codex prompts fully succeeded, whether setup builds, whether schemas exist, whether libs/contracts is complete, or whether setup/launcher CLI smoke tests pass.
\end{itemize}

\subsubsection{Domino\_Dominium\_Workbench}

\begin{itemize}
\item Reliability note:** This is a human-readable consolidation, not a live repository audit. Repo-current statements that came from pasted material or prior-chat excerpts remain **UNVERIFIED** until checked against the live julesc013/dominium repo and current AIDE/validator outputs.
\item Verify current live repo status and task gate state.
\item > Verify repo state, then generate COMMAND-RESULT-VIEW-SLICE-01.
\end{itemize}

\subsubsection{domino\_engine\_refactor\_prompts}

\begin{itemize}
\item The answer also introduced a deterministic bytecode VM as the preferred way to support modded behavior. A question was raised about whether native-code plugins should also be allowed, but the user did not answer. That remains unresolved.
\item What remains uncertain is the actual repository state. The chat designed architecture and prompts, but did not inspect source code. Therefore the plan needs repo verification before implementation.
\item The unresolved work is to generate detailed subsystem implementation prompts for heat, power, fluids, atmospheres, vehicles, and structural loads after the core engine spine is implemented.
\item The remaining uncertainty is plugin policy. A deterministic VM was proposed, but whether native-code plugins are allowed remains unresolved.
\item What remains uncertain is the exact thresholds and content policies for when things promote or demote.
\item The unresolved issue is how much of STRUCT compilation should be implemented first. The prompt plan is broad; a real implementation may need a smaller first milestone.
\item ignoring FACT/INFERENCE/UNCERTAIN labels;
\item UNCERTAIN / UNVERIFIED:** The package files were generated in the sandbox during the prior step, but repository files were not created or inspected.
\end{itemize}

\subsubsection{engine\_baseline\_architecture}

\begin{itemize}
\item Uncertainty: the repo has many modules and some historical/legacy surfaces, so exact implementation maturity varies by subsystem. The distinction should remain a formal requirement in a master spec.
\item Uncertainty: the full determinism envelope across all current modules was not independently re-run in this chat. User-supplied local validation results should be preserved but verified before being used as audit proof.
\item Uncertainty: this remains strategic architecture, not an implemented decision.
\item Unresolved goals include:
\item User wants uncertainty labelled and framing corrected when evidence disagrees.
\item UNCERTAIN: whether the user wants Unreal integration at all as a client/editor.
\item UNCERTAIN: exact tolerance for using external geometry libraries.
\item UNCERTAIN: whether broad repo restructuring will be desired after Milestone 0.
\end{itemize}

\subsubsection{Foundation\_Workbench\_Codex}

\begin{itemize}
\item What remains uncertain is how quickly this architecture will move from contracts and fixtures into runtime implementations and product slices. The next chat should not assume runtime provider resolvers, package runtime, or Workbench shell exist.
\item verify latest live repo state after user-pasted Wave 1 completion;
\item 1. Verify live repo state.
\item The user wants direct, evidence-grounded, audit-ready answers; timestamps/model labels; explicit uncertainty; no repeated clarification unless necessary; large continuous Codex prompt blocks when generating prompts; targeted tests instead of full suites; and rapid progress toward Workbench/code.
\item The user is likely to reject slow, report-only, overcautious process. Future assistants should produce executable next prompts quickly after verifying state.
\item It is uncertain how many concurrent Codex workers the user will actually run at once and whether the user wants coordinator prompts or worker prompts next.
\item The most important continuation is: verify current origin/main, confirm Wave 1 and Workbench validation state, then generate or run COMMAND-RESULT-VIEW-SLICE-01.
\end{itemize}

\subsubsection{Framework\_Open\_Source\_Provider}

\begin{itemize}
\item The chat inspected julesc013/dominium and found evidence of existing abstractions such as system and graphics layers, stubs, and soft-backed backends. The finding was that raylib could fill concrete provider gaps without replacing the architecture. However, current repo facts are stale/uncertain and must be verified, especially the C17/C++17 vs C90/C++98 baseline contradiction.
\item The exact first implementation branch, versions, and dependency pins are unresolved. The exact Domino Framework ABI is unresolved. The exact Lua version is unresolved. The current repository baseline must be verified. The formal policy for GPL/LGPL/unclear license material is unresolved. The sparse simulation and CAD systems need formal specs and prototypes.
\item The raylib decision was also directionally accepted. The user repeatedly expressed enthusiasm for raylib and asked about using as much of raylib infrastructure as possible. The caveat is that raylib is a provider suite, not architecture. This affects rendering, audio, input, asset preview, and Workbench bootstrap.
\item 1. Verify repo baseline and dependency versions.
\item The user requires source-grounded, audit-ready, uncertainty-labelled responses. The preservation prompt explicitly requires FACT/INFERENCE/UNCERTAIN/PROJECT-CONTEXT labels, no invented facts, no treating brainstorms as decisions, and no over-compression. The prompt also requires a human-readable report first and downloadable files if possible ?filecite?turn29file0?.
\item The most blocking unresolved issues are repo baseline verification, exact dependency pins, minimal framework ABI, provider profile order, deterministic simulation spec, and license policy.
\item Verify the current julesc013/dominium CMake language baseline.
\end{itemize}

\subsubsection{gui\_binary\_content}

\begin{itemize}
\item GPT-5.5 Pro - 2026-05-27 00:00:00 Australia/Melbourne; time UNVERIFIED
\item whether knowledge systems should include uncertainty, false beliefs, or lost knowledge,
\item This stage established a pattern that continues throughout the chat: the user does not want impressive-looking prompts that hide unresolved design choices. The user wants the assumptions exposed before anything is formalized.
\item The assistant agreed with the general direction, but added caveats. The most important were:
\item Those caveats were not final user decisions, but they are important warnings.
\item The conclusion was not "run the prompt now." The conclusion was that CONTENT0 is a strong foundation but needs discussion before finalization. The assistant-generated rewrite is useful but not final. The unresolved issues need to be resolved before Codex or any implementation work touches the repository.
\item What remains uncertain is how exactly to encode these ideas in schemas and data. For example, it is not yet decided how deterministic seed hierarchies work, how macro content refines into micro content, or how timeline causality should be represented.
\item The unresolved issue is the shared backend/UI contract. Without that, the GUI plan is only a list of technologies.
\end{itemize}

\subsubsection{Language\_Platform\_Architecture}

\begin{itemize}
\item The 32-bit question followed. The answer was to make full native products 64-bit and keep 32-bit as constrained or projection support only. The key caveat was not to make native pointer size part of game law. Save, replay, network, IDs, and renderer IR must use fixed-width types and never serialize size\_t, long, uintptr\_t, or raw pointers.
\item Exact tolerance for exceptions/RTTI inside private C++17 code remains unsettled. Exact platform floors and raylib/SDL2 priority remain unresolved.
\item Verify the Windows 7 and macOS 10.9 toolchain/library assumptions using official sources.
\end{itemize}

\subsubsection{launcher\_app\_layer}

\begin{itemize}
\item These prompts matter because they superseded earlier advice. They also created a new practical priority: verify that the applied prompts actually succeeded.
\item This changed the correct future behavior. The next assistant must not redesign. It must implement, audit, maintain, document, verify, and harden.
\item This matters because it prevents platform-specific drift. If CLI has verify, GUI must not implement a different hidden version of verify. If GUI has a pack enable flow, it should map to the same command/intent as CLI and TUI.
\item The final purity prompt required those to be moved or quarantined. However, actual repo success remains unverified in this chat.
\item The important change is that the conversation moved from "What architecture should we use?" to "Architecture is locked; how do we verify and harden the implementation?"
\item The largest unresolved goal is verifying the actual repo. The user stated that two Codex prompts were applied, but this chat did not independently verify the filesystem, build, tests, scripts, or launcher docs.
\item Another unresolved goal is confirming whether the launcher hardening prompt has been applied.
\item A third unresolved goal is confirming whether command graph, UI IR, binding validation, zero-pack tests, RepoX changelog display, and BUILD-ID mismatch refusal are implemented.
\end{itemize}

\subsubsection{Launcher\_Setup\_Architecture}

\begin{itemize}
\item Uncertainty:** The exact repository layout and existing implementation were not inspected in this chat.
\item Uncertainty:** The final language boundary for setup remains unresolved after the later C89/dsys/dgfx launcher direction.
\item Uncertainty:** The exact mapping into runtime binaries remains unverified.
\item Uncertainty:** Accounts/social/wiki/forum/direct messaging are future-facing and were not converted into the final five-tab implementation plan.
\item Uncertainty:** Earlier setup designs used JSON manifests. Later dsys/dgfx architecture leaned toward TLV .dmeta files.
\item Uncertainty:** Actual plugin ABI details and dynamic loading support depend on dsys/repo APIs.
\item Uncertainty:** Actual dsys/dgfx APIs were not inspected.
\item Conclusion reached:** This chat is now documented and can be merged later, but with caveats.
\end{itemize}

\subsubsection{Modularity\_AIDE\_Refactorability}

\begin{itemize}
\item The live repo still needs verification. The exact AIDE files, contracts, validators, and migration tasks must be implemented. It remains unresolved which existing roots and tools map to which final homes, which APIs become stable, which code is reusable across products/games/projects, and which prior assistant recommendations should become formal requirements after user review.
\item 1. **TASK-09** - Verify live repo baseline: current branch, root inventory, validators, CTest status, generated artifacts.
\item The user prefers source-grounded, audit-ready decisions; bounded uncertainty; explicit task prompts; and reusable engineering doctrine. Future assistants should avoid shallow affirmation and should distinguish accepted user decisions from assistant recommendations.
\item Important caveat: the only current uploaded file available for this task is Pasted text.txt, which contains the preservation-package prompt. Earlier generated handoff files, ZIP packages, or uploaded docs referenced in prior conversation are not available in this chat unless the user re-uploads them.
\item The most serious risk is that a future assistant treats this chat as if it verified the live repo. It did not. The second serious risk is that it hardens every assistant suggestion into a requirement. Several items are strong recommendations aligned with user goals, but still require implementation review. Future chats must preserve labels and verify stale facts.
\item Future assistants must verify live repo state before implementing cleanup.
\end{itemize}

\subsubsection{omega\_xi\_pi\_architecture\_future}

\begin{itemize}
\item Preserve FACT / INFERENCE / UNCERTAIN labels.
\item Do not invent or flatten uncertainty.
\item 7. "List all artifacts I should verify in GitHub main."
\item 9. "What should I verify before starting Sigma?"
\end{itemize}

\subsubsection{OS\_Interface\_Repo\_Architecture}

\begin{itemize}
\item The unresolved goals are implementation-level. The repo still needs command dispatch unification, product boot proof, portable projection proof, AppShell rendered-mode law update, Workbench module contracts, document/patch/result/refusal schemas, and a boot-to-replay MVP. It also needs continued exception retirement, CTest/RepoX remediation, and verification of current build status.
\item 2. Repair canonical verify test discovery.
\item The user wants uncertainty labels and correction of incorrect framing.
\item This preservation task specifically required FACT / INFERENCE / UNCERTAIN / PROJECT-CONTEXT labels, stable IDs, registers, spec sheets, an aggregator packet, self-audit, and downloadable files if tools are available.
\item UNCERTAIN: Whether the user wants to rename engine conceptually to kernel in code, or only use kernel as a conceptual layer.
\item UNCERTAIN: Whether the final public product name should be Dominium Workbench, Dominium Studio, Domino Workbench, or something else.
\item UNCERTAIN: How aggressive the next actual code refactor should be versus staged adapters and codegen.
\end{itemize}

\subsubsection{platform\_renderer\_api\_plan}

\begin{itemize}
\item This is not a fatal problem for the prompt pack, because the prompts themselves repeatedly tell Codex to discover files with rg, tree, type, CMake commands, and target listing. But it is a major caveat for any future assistant: do not assume the actual repo matches every file path or API name unless checked.
\item UNCERTAIN / UNVERIFIED:** The conversation did not prove that these prompts were actually applied to the repository. They are an assumed prerequisite, not verified code state.
\item Uncertainty remains about actual repo state. The prompts are designed to discover the repo state, but the chat itself did not inspect the archive.
\item The final active prompt pack narrowed runtime completion to Tier-1 on Windows and compile-gated correctness for Tier-2. This is a practical compromise because Codex runs on Windows. But it also creates a caveat: if "fully complete" later means fully running X11/Wayland/Cocoa/Metal/Vulkan, more prompts or native CI validation will be required.
\item This became one of the final "done" gates, along with scripts\textbackslash{}build\_codex\_verify.bat.
\item This matters because the conversation was long and contained multiple plans, superseded options, caveats, and active artifacts. The report package and this plain-language report exist to prevent future assistants from restarting, re-asking, or losing distinctions between decisions, brainstorms, and uncertainties.
\item The user also seems to want future assistants to operate with high continuity and not lose subtle distinctions. This is inferred from the repeated requests for maximum-fidelity transfer, stable IDs, labels, uncertainty tracking, and report packaging.
\item Another unresolved goal is whether the master engine prompts 1-14 are truly implemented. The later plan depends on them, but this chat did not verify the repo.
\end{itemize}

\subsubsection{platform\_support}

\begin{itemize}
\item UNCERTAIN / UNVERIFIED: The name "Domino" was never confirmed by the user. Future work should treat "Domino" as an assistant-created placeholder unless the user explicitly accepts it. The useful idea is the distinction between full product support and reduced engine/core support, not necessarily the name.
\item The main unresolved issue is the exact definition of support. Does Tier-0 mean identical feature parity, simultaneous release, identical save compatibility, identical UI, equal QA coverage, or some platform-specific differences? The assistant inferred that Tier-0 should mean feature parity, full QA, long-term support, and high release priority, but the exact contractual meaning still needs formal definition.
\item What remains unresolved is the exact PC baseline. The chat did not decide whether Windows 10, Windows 11, or both are required. It did not decide whether Windows ARM64 matters. It did not decide whether macOS support includes Intel Macs, Apple Silicon Macs, or both. It did not define Linux distributions, glibc requirements, Flatpak/AppImage/deb/rpm packaging, X11/Wayland policy, or Steam Deck verification.
\item UNCERTAIN / UNVERIFIED: The chat did not establish minimum Android API level, ABI list, target graphics APIs, Play Store rules, Android TV status, Android Go status, Automotive status, Chromebook status, or vendor ROM support. These need future decisions and current-source verification.
\item The user also listed iPod Touch, Apple TV, Apple Watch, Mac, macOS, tvOS, watchOS, and Apple Vision Pro indirectly through AR. These were not classified in detail. INFERENCE: macOS belongs under PC. UNCERTAIN / UNVERIFIED: tvOS, watchOS, and visionOS remain undecided. Full Dominium on Apple Watch is not established and should not be assumed.
\item The unresolved issues are substantial. The chat did not decide whether WebGPU is required, whether WebGL fallback is mandatory, whether the game must work offline as a PWA, whether WASM threading/SIMD is required, how browser storage will work, how asset loading will be handled, or what browser versions will be supported.
\item UNCERTAIN / UNVERIFIED: Current console platform facts were not verified in this chat. Future work must verify official PlayStation, Xbox, and Nintendo developer requirements from current official sources before implementation planning.
\item INFERENCE: PC handhelds should be subtargets under PC, with their own QA and UI profiles. UNCERTAIN / UNVERIFIED: Steam Deck's exact tier was not finally decided.
\end{itemize}

\subsubsection{Portability\_Assurance\_Future\_Proof}

\begin{itemize}
\item The chat proposed using standards such as DO-178C, NIST SSDF, OWASP ASVS, SLSA, and SPDX as design inputs only. The proposed internal version is DDAP v0 with DIL levels. Uncertainty: user has not explicitly ratified DDAP.
\item The actual repo was not inspected. The exact accepted directory structure, API policy, DDAP profile, compatibility promises, and first pilot module remain unresolved.
\item Caveat: carry forward as INFERENCE / assistant recommendation.
\item Caveat: carry forward as FACT goal + INFERENCE architecture.
\item Caveat: carry forward as assistant recommendation.
\item Implications: Safe for aggregation only with caveats.
\item Caveat: carry forward as FACT / UNCERTAIN scope.
\item 10. Feed this package into the master Project Spec Book with caveats.
\end{itemize}

\subsubsection{readme\_ports\_determinism}

\begin{itemize}
\item What remains uncertain is whether the actual repository README.md matches the final pasted version. The chat only saw pasted text and generated prompts; it did not inspect a Git repository.
\item The unresolved issue is whether a metadata-only /ports directory still violates the user's preference. The final README keeps /ports in a limited form, but the user's wording could mean no /ports directory at all. That should be clarified before writing the directory contract.
\item The unresolved work is to define the actual network protocol, prediction model, rollback window, reconciliation logic, and state-hash validation.
\item The major unresolved goal is deciding exactly what "no separate ports directory/system" means physically in the repository. The final README allows /ports as metadata-only. That may satisfy the user, but it is not confirmed. If the user meant no /ports directory at all, the README still needs another revision.
\item Another unresolved goal is making the README and normative specs align. The README says specs are authoritative, but those specs were not inspected in this chat.
\item UNCERTAIN / UNVERIFIED: /ports as metadata-only is the current README state, but may not be fully accepted.
\item A related rejected idea was **retro build flows under /ports/<target> containing source or alternative implementations**. The final README superseded this with metadata-only /ports, assuming /ports is retained at all. This rejection is final for source code and behavior, but the existence of /ports itself remains uncertain.
\item The first next step is to verify the actual repository README.md against the final pasted README. The chat only used pasted text. If the repository differs, future work should start from the actual file, not from memory.
\end{itemize}

\subsubsection{Refactor\_Architecture}

\begin{itemize}
\item Reliability note: Repository state, Codex execution, exact version numbers, and actual file modifications are **UNCERTAIN / UNVERIFIED** unless explicitly stated otherwise.
\item The most important thing to remember is that this chat did not implement the refactor. It designed the architecture and generated prompts and handoff material. Any future assistant must verify actual repository state before assuming files were moved, prompts were applied, or code was updated.
\item UNCERTAIN / UNVERIFIED:** It did not establish that this system already exists in code.
\item What remains uncertain is the exact degree to which existing code currently respects that boundary. The user provided a repo tree, but this chat did not inspect or modify code directly.
\item Unresolved details include exact default DOMINIUM\_HOME paths per OS and exact implementation of import/GC.
\item Unresolved details include exact schemas, parser choice, tests, and actual code implementation.
\item The key unresolved point is implementation timing. The report preserves this architecture, but the initial Codex refactor was not supposed to rewrite the entire UI system immediately.
\item Actual implementation is unresolved. This chat generated architecture and prompts, not a verified changed codebase.
\end{itemize}

\subsubsection{Refactor\_Control\_Plane}

\begin{itemize}
\item Unresolved goals include finishing AIDE Q35-Q57, stabilizing Dominium CTest/RepoX blockers, defining final AIDE install/upgrade bundles, implementing root recycling machinery, deciding when product boot proof can proceed, and later deciding whether AIDE Runtime/Gateway/Hosts are truly needed.
\item 1. Verify current Dominium HEAD and POST-CONVERGE status.
\item The user values direct, source-grounded, audit-ready, detailed answers. The user asked for preservation packages that are human-readable and machine-aggregatable. The user repeatedly prefers not to lose nuance, rejected options, uncertainties, or artifacts. The user wants current facts verified when possible and wants future assistants to avoid re-asking answered questions.
\item UNCERTAIN: Exact public-facing naming for AIDE Core/Kernel remains partly open. Exact threshold for accepting CTest-warning status remains a user/repo governance decision. Exact final AIDE installation packaging details remain pending.
\item Future assistants must verify live repo state before implementing.
\item "Merge this with another chat's preservation package without losing uncertainty labels."
\end{itemize}

\subsubsection{Release\_Identity\_and\_Versioning}

\begin{itemize}
\item The chat rebuilt SemVer from first principles and decided that strict SemVer should apply only to components with declared public contracts. Candidate true-SemVer entities include SDKs, engine libraries, tool APIs/CLIs, protocol libraries, plugin hosts, schema libraries, and reusable runtime libraries. It remains unresolved which exact repo modules qualify.
\item GBN should identify CI/public/internal builds and appear as provenance, e.g. +gbn.7137, but it must not be used as SemVer precedence. BII should primarily be structured manifest metadata, because it may encode lane, target, kind, configuration, and other build identity fields. The exact BII schema remains unresolved.
\item The main unresolved goals are to produce formal specs: Release Constitution, SemVer Component Inventory, Suite Release Policy, Build Identity Spec, Channel/Lifecycle Spec, Artifact Naming Spec, Manifest Schema, Target Taxonomy, and Capability Registry.
\item It is still uncertain how strict the user wants suite X.Y.Z field meanings to be, whether capabilities should fully replace compatibility profiles or coexist with them, and how much old-platform support should be real versus aspirational.
\item The most important unresolved issue is formal classification. Without knowing which entities are strict SemVer components, suite versions, product release IDs, build fields, or capabilities, the rest of the policy cannot be enforced safely.
\item Platform targeting remains unresolved. The chat concluded that coarse families like WinNT5, WinNT10, MacOSX4, Linux5, or DOS5 are useful support labels, but exact binary compatibility needs target baselines and runtime/toolchain profiles. Linux kernel major alone is not enough. Windows 9x/NT/NT5/NT6/NT10 likely need separate build lanes or at least explicit tested baselines.
\end{itemize}

\subsubsection{testx\_repox\_governance}

\begin{itemize}
\item This topic remains unresolved in implementation because the repo snapshot showed .dominium\_build\_number and update\_build\_number.cmake, but the chat did not inspect their contents. Those files may still implement older build-number behavior. That must be verified.
\item The unresolved part is implementation verification. The repo has Sol/Earth content in data and game/content; the engine must not assume those exist.
\item The unresolved part is how much of TESTX is actually implemented in the repo. The repo contains many tests and scripts, but this chat did not inspect them. A rules-to-checks audit is needed.
\item This remains conceptually accepted in the chat, but implementation is unverified.
\item The unresolved goals are mostly implementation and verification:
\item The repo already has UI IR/codegen infrastructure. The next step is to verify whether GUI/TUI truly bind to the CLI command graph and whether UI is data rather than logic.
\item The user has Windows 10 + VS2022. The next step is to verify v141/v141\_xp/SDK components and create isolated CMake presets for XP/Win7/Win8/Win10/Win11 builds.
\item Acting before verifying repo state.
\end{itemize}

\subsubsection{Timekeeping\_and\_2038\_Resilience}

\begin{itemize}
\item What remains uncertain is the current exact status of specific libc/OS/toolchain combinations, because those facts can change and were not reverified during this preservation turn.
\item What remains uncertain is the precise target list of legacy machines/toolchains and how far compatibility must go, especially for 16-bit compilers with weak or absent 64-bit arithmetic.
\item The main unresolved goals are backend audit, ACT serialization audit, civil/astronomical projection design, and verification of external platform facts.
\item The main uncertainty is completeness of repo inspection. The assistant read key files, but not all serialization paths and not all platform backends.
\item 6. Verify current platform/library facts before using them in formal public documentation.
\item The preservation prompt explicitly requires a human-readable report first, not only a machine-readable handoff. It requires uncertainty labels, stable registers, preservation of artifacts, no invented facts, and downloadable files if tools are available. ?filecite?turn30file0?
\item The user prefers source-grounded, audit-ready technical reasoning, direct structure, and careful distinction between facts, recommendations, and unresolved issues. PROJECT-CONTEXT indicates Dominium values C89/C++98 portability, deterministic architecture, CLI/TUI support, and clean platform/runtime separation.
\item The strongest unresolved technical task is auditing actual backend timers and persistence formats.
\end{itemize}

\subsubsection{UE6\_Domino\_Deterministic\_Universe}

\begin{itemize}
\item UNCERTAIN / UNVERIFIED: This is not a complete reconstruction of every Dominium conversation ever held. The project context shows that many other chats exist about platforms, renderers, setup, game architecture, ecology, ECS, launcher design, and development planning, but their full transcripts are not visible here. Those items are labelled PROJECT-CONTEXT when referenced.
\item Uncertainty: public UE6 technical capabilities and requirements remain time-sensitive and need verification before any future hard commitment.
\item The chat rejected the idea that clients can authoritatively simulate resource production, hidden state, enemy bases, combat, or persistent terrain changes in an MMO. Clients can assist, but the server must verify and commit.
\item REJECTED-02: Waiting for UE6 before beginning serious work. Deprioritised because UE6 availability and capabilities remain uncertain.
\item The user wants direct, source-grounded, audit-ready answers; explicit uncertainty; bounded estimates; and correction when framing is wrong. The preservation prompt explicitly requires human-readable explanation first, structured registers second, uncertainty labels, no invented facts, no hidden chain-of-thought, and downloadable files if tools are available.
\item UNCERTAIN: The final intended role of Domino remains unclear in the visible transcript. It may be a custom engine, target runtime, renderer, or broader platform plan, but the current chat does not fully define it.
\end{itemize}

\subsubsection{ui\_editor\_tool\_editor\_planning}

\begin{itemize}
\item The most important thing to remember is that this chat was a planning and prompt-generation chat, not proof of implementation. **UNCERTAIN / UNVERIFIED:** No generated prompt is evidence that repository code was actually changed. The next assistant must first verify whether the prompts were executed, inspect the repo/source bundle if available, and avoid treating planned files as existing files.
\item The user uploaded screenshot bundles for setup and launcher. **UNCERTAIN / UNVERIFIED:** The screenshots were not inspected in this chat. The assistant still produced a logical layout description, and the user accepted it as useful. Future work should inspect the bundles before claiming exact screenshot fidelity.
\item What remains uncertain is the actual state of the code. The user named files, but the chat did not inspect the repository. A future assistant must verify exact paths, build targets, TLV parser implementation, and current launcher behavior.
\item Unresolved details include the actual TLV wire format, migration strategy from launcher\_ui\_v1.tlv, and how much legacy data can be preserved during import.
\item The unresolved risk is implementation complexity. Real native tab controls, splitter behavior, and scroll panels can be nontrivial in a child-HWND Win32 UI.
\item This is a strong planned design, but some details remain assistant recommendations rather than explicit user-confirmed decisions. Future implementation should preserve that uncertainty if writing a formal spec.
\item The unresolved issue is that the uploaded screenshot bundles were not inspected. The logical specs are accepted planning artifacts, not verified screenshot extractions.
\item The exact scope of IDE-native live editing remains unresolved. The assistant proposed preview hosts, but it is not confirmed whether that fully satisfies the user's desire to use Visual Studio, Xcode, and Linux GUI tools.
\end{itemize}

\subsubsection{Universe\_Explorer\_Planning}

\begin{itemize}
\item Uncertain: whether all seven universal layers are now formally captured under exactly that naming. The repo has equivalent or related doctrine, but not every early phrase appears as a single artifact.
\item Uncertain: the repo contains pieces of this through affordance/capability/domain docs, but the assistant did not find one single master Deep Primitives spec in the docs it read.
\item The most important unresolved issue is that the Universe Explorer plan is still a proposed/recommended next strategy, not yet a completed repo task. The next best action is either to draft PRESENTATION-CONTRACT-01 / UNIVERSE-EXPLORER-CONTRACT-01, or to preserve this chat and merge it with other old-chat reports into a master Project Spec Book.
\end{itemize}

\subsubsection{Workbench\_AIDE\_Product\_Spine}

\begin{itemize}
\item The largest uncertainty is not the conceptual architecture; that has converged. The largest uncertainty is live execution state: whether prompts generated near the end of the chat have already been run locally, committed, or pushed. The next chat should first inspect .aide/queue/current.toml, recent commits, and relevant audit files.
\item Remaining uncertainty: exact implementation of erosion/ecology/evolution proxies remains future domain work.
\item 1. Verify current repo state.
\item Explicitly label uncertainty.
\item The next chat should verify repo state first.
\item Preserve FACT / INFERENCE / UNCERTAIN / PROJECT-CONTEXT labels.
\item Verify live repo state before acting.
\item safe\_for\_aggregation: "Yes, with caveats"
\end{itemize}

\subsubsection{World\_Architecture}

\begin{itemize}
\item The main unresolved issue is implementation detail. The design says Q16.16 and Q4.12, but the earlier prompt used signed types incorrectly. The future implementation must use unsigned local coordinates or a centered signed design.
\item The unresolved parts are the exact meshing algorithm, the fixed-point scaling of ?, and the sample resolution. A 32? sample grid per 16m chunk was recommended but not confirmed as final.
\item The exact solvers remain unresolved. The architecture is clear, but formulas and resolutions are still future work.
\item The save format remains conceptual in places. Exact binary headers, endian policy, compression, and migration strategy are unresolved.
\item The crucial caveat is implementation: unsigned or centered representations are needed. The design intent is final, but the exact typedefs must be fixed. If the implementation uses signed types incorrectly, the whole coordinate system breaks.
\item This is not yet an implementation-level final spec because solver formulas are unresolved.
\item FACT:** Core target is C89. This is non-negotiable in principle, but exact portability mechanics are unresolved because strict C89 lacks standard fixed-width integer types.
\item The user wants direct, detailed, rigorous, critical responses. The user prefers source labels, uncertainty preservation, and avoiding invented facts. The user does not want assistant suggestions silently upgraded to decisions.
\end{itemize}

\subsubsection{xstack\_lab\_galaxy}

\begin{itemize}
\item But this chat did not verify the repository. That became one of the most important unresolved issues.
\item This is one of the most important outputs associated with this chat, but it is also one of the most uncertain because it is user-reported from another chat and not verified directly here. The user explicitly wanted it audited for completion and consistency.
\item INFERENCE: The user also wanted to prevent future assistants from flattening nuance. They repeatedly asked for maximum fidelity, uncertainty labels, rejected options, contradictions, and rationale. This implies a strong preference for preserving complexity rather than oversimplifying.
\item INFERENCE: The user likely wants to resume productive development after verifying the substrate. The next likely feature area is survival or Lab Galaxy refinement, but the user has not made a final decision in this visible chat.
\item 4. From documentation to verifying a new 13-prompt substrate.
\item The biggest unresolved goal is verifying the actual repository. The package says what was reported, but not what is proven. Another unresolved goal is deciding what comes next: survival, Lab Galaxy UX, distributed SRZ, or documentation polish.
\item The main assumption behind all of this is that the repo can actually enforce these contracts. That remains unverified in this chat.
\item The best next step is to verify the actual repository state. This matters because the long 13-prompt transcript is not proof. The user explicitly does not know whether everything was truly implemented.
\end{itemize}

\subsection{Stale Claim Findings}

\begin{itemize}
\item CONTRA-0002 advanced\_simulation\_infrastructure: Archived conversation contains potentially stale external or baseline claim: C89.
\item CONTRA-0003 advanced\_simulation\_infrastructure: Archived conversation contains potentially stale external or baseline claim: C++98.
\item CONTRA-0004 advanced\_simulation\_infrastructure: Archived conversation contains potentially stale external or baseline claim: ISO C89.
\item CONTRA-0009 app\_runtime\_platform\_renderers: Archived conversation contains potentially stale external or baseline claim: SDK.
\item CONTRA-0012 app\_testx\_codehygiene: Archived conversation contains potentially stale external or baseline claim: C89.
\item CONTRA-0013 app\_testx\_codehygiene: Archived conversation contains potentially stale external or baseline claim: C++98.
\item CONTRA-0014 app\_testx\_codehygiene: Archived conversation contains potentially stale external or baseline claim: DX12.
\item CONTRA-0015 app\_testx\_codehygiene: Archived conversation contains potentially stale external or baseline claim: iOS.
\item CONTRA-0018 architecture\_codex\_prompts: Archived conversation contains potentially stale external or baseline claim: C89.
\item CONTRA-0019 architecture\_codex\_prompts: Archived conversation contains potentially stale external or baseline claim: C++98.
\item CONTRA-0020 architecture\_codex\_prompts: Archived conversation contains potentially stale external or baseline claim: ISO C89.
\item CONTRA-0026 architecture\_ui\_providers: Archived conversation contains potentially stale external or baseline claim: C89.
\item CONTRA-0027 architecture\_ui\_providers: Archived conversation contains potentially stale external or baseline claim: C++98.
\item CONTRA-0028 architecture\_ui\_providers: Archived conversation contains potentially stale external or baseline claim: DX12.
\item CONTRA-0029 architecture\_ui\_providers: Archived conversation contains potentially stale external or baseline claim: SDK.
\item CONTRA-0030 Build\_and\_Future\_Proofing: Archived conversation contains potentially stale external or baseline claim: C89.
\item CONTRA-0031 Build\_and\_Future\_Proofing: Archived conversation contains potentially stale external or baseline claim: C++98.
\item CONTRA-0032 Build\_and\_Future\_Proofing: Archived conversation contains potentially stale external or baseline claim: SDK.
\item CONTRA-0038 canonical\_structure\_and\_framework: Archived conversation contains potentially stale external or baseline claim: C89.
\item CONTRA-0039 canonical\_structure\_and\_framework: Archived conversation contains potentially stale external or baseline claim: C++98.
\item CONTRA-0040 canonical\_structure\_and\_framework: Archived conversation contains potentially stale external or baseline claim: SDK.
\item CONTRA-0043 development\_routes: Archived conversation contains potentially stale external or baseline claim: iOS.
\item CONTRA-0044 documentation\_standards\_readme: Archived conversation contains potentially stale external or baseline claim: C89.
\item CONTRA-0045 documentation\_standards\_readme: Archived conversation contains potentially stale external or baseline claim: C++98.
\item CONTRA-0046 documentation\_standards\_readme: Archived conversation contains potentially stale external or baseline claim: ISO C89.
\item CONTRA-0047 documentation\_standards\_readme: Archived conversation contains potentially stale external or baseline claim: iOS.
\item CONTRA-0050 Dominium\_Architecture\_I: Archived conversation contains potentially stale external or baseline claim: C89.
\item CONTRA-0051 Dominium\_Architecture\_I: Archived conversation contains potentially stale external or baseline claim: C++98.
\item CONTRA-0054 Dominium\_Architecture\_II: Archived conversation contains potentially stale external or baseline claim: C89.
\item CONTRA-0055 Dominium\_Architecture\_II: Archived conversation contains potentially stale external or baseline claim: Windows 98.
\item CONTRA-0056 Dominium\_Architecture\_II: Archived conversation contains potentially stale external or baseline claim: Windows NT 2000.
\item CONTRA-0057 Dominium\_Architecture\_II: Archived conversation contains potentially stale external or baseline claim: DX12.
\item CONTRA-0058 Dominium\_Architecture\_II: Archived conversation contains potentially stale external or baseline claim: Android.
\item CONTRA-0061 Dominium\_Architecture\_III: Archived conversation contains potentially stale external or baseline claim: C89.
\item CONTRA-0062 Dominium\_Architecture\_III: Archived conversation contains potentially stale external or baseline claim: C++98.
\item CONTRA-0063 Dominium\_Architecture\_III: Archived conversation contains potentially stale external or baseline claim: Win98.
\item CONTRA-0064 Dominium\_Architecture\_III: Archived conversation contains potentially stale external or baseline claim: Windows 98.
\item CONTRA-0065 Dominium\_Architecture\_III: Archived conversation contains potentially stale external or baseline claim: Windows NT 2000.
\item CONTRA-0066 Dominium\_Architecture\_III: Archived conversation contains potentially stale external or baseline claim: Mac OS 9.
\item CONTRA-0067 Dominium\_Architecture\_III: Archived conversation contains potentially stale external or baseline claim: Mac OS 8.
\item CONTRA-0068 Dominium\_Architecture\_III: Archived conversation contains potentially stale external or baseline claim: DX12.
\item CONTRA-0072 Dominium\_Architecture\_IV: Archived conversation contains potentially stale external or baseline claim: CP/M.
\item CONTRA-0073 Dominium\_Architecture\_IV: Archived conversation contains potentially stale external or baseline claim: Android.
\item CONTRA-0074 Dominium\_Architecture\_IV: Archived conversation contains potentially stale external or baseline claim: SDK.
\item CONTRA-0076 Dominium\_Complete\_Conversation: Archived conversation contains potentially stale external or baseline claim: C89.
\item CONTRA-0077 Dominium\_Complete\_Conversation: Archived conversation contains potentially stale external or baseline claim: C++98.
\item CONTRA-0078 dominium\_domino\_codex\_planning: Archived conversation contains potentially stale external or baseline claim: C89.
\item CONTRA-0079 dominium\_domino\_codex\_planning: Archived conversation contains potentially stale external or baseline claim: C++98.
\item CONTRA-0080 dominium\_domino\_codex\_planning: Archived conversation contains potentially stale external or baseline claim: ISO C89.
\item CONTRA-0081 dominium\_domino\_codex\_planning: Archived conversation contains potentially stale external or baseline claim: CP/M.
\item CONTRA-0082 dominium\_domino\_codex\_planning: Archived conversation contains potentially stale external or baseline claim: Android.
\item CONTRA-0088 dominium\_full\_conversation: Archived conversation contains potentially stale external or baseline claim: SDK.
\item CONTRA-0091 dominium\_setup: Archived conversation contains potentially stale external or baseline claim: SDK.
\item CONTRA-0096 Domino\_Dominium\_Workbench: Archived conversation contains potentially stale external or baseline claim: SDK.
\item CONTRA-0098 domino\_engine\_refactor\_prompts: Archived conversation contains potentially stale external or baseline claim: C89.
\item CONTRA-0099 domino\_engine\_refactor\_prompts: Archived conversation contains potentially stale external or baseline claim: C++98.
\item CONTRA-0100 domino\_engine\_refactor\_prompts: Archived conversation contains potentially stale external or baseline claim: ISO C89.
\item CONTRA-0111 Foundation\_Workbench\_Codex: Archived conversation contains potentially stale external or baseline claim: C89.
\item CONTRA-0112 Foundation\_Workbench\_Codex: Archived conversation contains potentially stale external or baseline claim: C++98.
\item CONTRA-0113 Foundation\_Workbench\_Codex: Archived conversation contains potentially stale external or baseline claim: SDK.
\item CONTRA-0118 Framework\_Open\_Source\_Provider: Archived conversation contains potentially stale external or baseline claim: C++98.
\item CONTRA-0119 Framework\_Open\_Source\_Provider: Archived conversation contains potentially stale external or baseline claim: SDK.
\item CONTRA-0122 gui\_binary\_content: Archived conversation contains potentially stale external or baseline claim: iOS.
\item CONTRA-0123 gui\_binary\_content: Archived conversation contains potentially stale external or baseline claim: Android.
\item CONTRA-0124 gui\_binary\_content: Archived conversation contains potentially stale external or baseline claim: SDK.
\item CONTRA-0129 Language\_Platform\_Architecture: Archived conversation contains potentially stale external or baseline claim: C89.
\item CONTRA-0130 Language\_Platform\_Architecture: Archived conversation contains potentially stale external or baseline claim: C++98.
\item CONTRA-0131 Language\_Platform\_Architecture: Archived conversation contains potentially stale external or baseline claim: Windows 98.
\item CONTRA-0132 Language\_Platform\_Architecture: Archived conversation contains potentially stale external or baseline claim: Mac OS 9.
\item CONTRA-0133 Language\_Platform\_Architecture: Archived conversation contains potentially stale external or baseline claim: iOS.
\item CONTRA-0134 Language\_Platform\_Architecture: Archived conversation contains potentially stale external or baseline claim: Android.
\item CONTRA-0135 Language\_Platform\_Architecture: Archived conversation contains potentially stale external or baseline claim: SDK.
\item CONTRA-0137 launcher\_app\_layer: Archived conversation contains potentially stale external or baseline claim: SDK.
\item CONTRA-0140 Launcher\_Setup\_Architecture: Archived conversation contains potentially stale external or baseline claim: C89.
\item CONTRA-0141 Launcher\_Setup\_Architecture: Archived conversation contains potentially stale external or baseline claim: C++98.
\item CONTRA-0142 Launcher\_Setup\_Architecture: Archived conversation contains potentially stale external or baseline claim: CP/M.
\item CONTRA-0143 Launcher\_Setup\_Architecture: Archived conversation contains potentially stale external or baseline claim: SDK.
\item CONTRA-0146 Modularity\_AIDE\_Refactorability: Archived conversation contains potentially stale external or baseline claim: C89.
\item CONTRA-0152 OS\_Interface\_Repo\_Architecture: Archived conversation contains potentially stale external or baseline claim: Mac OS 8.
\item CONTRA-0153 OS\_Interface\_Repo\_Architecture: Archived conversation contains potentially stale external or baseline claim: Android.
\item CONTRA-0154 OS\_Interface\_Repo\_Architecture: Archived conversation contains potentially stale external or baseline claim: SDK.
\item CONTRA-0158 platform\_renderer\_api\_plan: Archived conversation contains potentially stale external or baseline claim: C89.
\item CONTRA-0159 platform\_renderer\_api\_plan: Archived conversation contains potentially stale external or baseline claim: C++98.
\item CONTRA-0160 platform\_renderer\_api\_plan: Archived conversation contains potentially stale external or baseline claim: ISO C89.
\item CONTRA-0161 platform\_renderer\_api\_plan: Archived conversation contains potentially stale external or baseline claim: CP/M.
\item CONTRA-0162 platform\_renderer\_api\_plan: Archived conversation contains potentially stale external or baseline claim: Android.
\item CONTRA-0163 platform\_renderer\_api\_plan: Archived conversation contains potentially stale external or baseline claim: SDK.
\item CONTRA-0167 platform\_support: Archived conversation contains potentially stale external or baseline claim: CP/M.
\item CONTRA-0168 platform\_support: Archived conversation contains potentially stale external or baseline claim: iOS.
\item CONTRA-0169 platform\_support: Archived conversation contains potentially stale external or baseline claim: Android.
\item CONTRA-0170 platform\_support: Archived conversation contains potentially stale external or baseline claim: WebGPU.
\item CONTRA-0171 platform\_support: Archived conversation contains potentially stale external or baseline claim: console program.
\item CONTRA-0172 platform\_support: Archived conversation contains potentially stale external or baseline claim: SDK.
\item CONTRA-0174 Portability\_Assurance\_Future\_Proof: Archived conversation contains potentially stale external or baseline claim: C89.
\item CONTRA-0175 Portability\_Assurance\_Future\_Proof: Archived conversation contains potentially stale external or baseline claim: C++98.
\item CONTRA-0176 readme\_ports\_determinism: Archived conversation contains potentially stale external or baseline claim: C89.
\item CONTRA-0177 readme\_ports\_determinism: Archived conversation contains potentially stale external or baseline claim: C++98.
\item CONTRA-0178 readme\_ports\_determinism: Archived conversation contains potentially stale external or baseline claim: CP/M.
\item CONTRA-0179 readme\_ports\_determinism: Archived conversation contains potentially stale external or baseline claim: 286.
\item CONTRA-0180 Refactor\_Architecture: Archived conversation contains potentially stale external or baseline claim: C89.
\item CONTRA-0181 Refactor\_Architecture: Archived conversation contains potentially stale external or baseline claim: CP/M.
\item CONTRA-0182 Refactor\_Architecture: Archived conversation contains potentially stale external or baseline claim: Android.
\item CONTRA-0183 Refactor\_Architecture: Archived conversation contains potentially stale external or baseline claim: SDK.
\item CONTRA-0185 Refactor\_Control\_Plane: Archived conversation contains potentially stale external or baseline claim: SDK.
\item CONTRA-0187 Release\_Identity\_and\_Versioning: Archived conversation contains potentially stale external or baseline claim: SDK.
\item CONTRA-0192 testx\_repox\_governance: Archived conversation contains potentially stale external or baseline claim: C89.
\item CONTRA-0193 testx\_repox\_governance: Archived conversation contains potentially stale external or baseline claim: C++98.
\item CONTRA-0194 testx\_repox\_governance: Archived conversation contains potentially stale external or baseline claim: iOS.
\item CONTRA-0195 testx\_repox\_governance: Archived conversation contains potentially stale external or baseline claim: SDK.
\item CONTRA-0196 Timekeeping\_and\_2038\_Resilience: Archived conversation contains potentially stale external or baseline claim: C89.
\item CONTRA-0197 Timekeeping\_and\_2038\_Resilience: Archived conversation contains potentially stale external or baseline claim: C++98.
\item CONTRA-0200 UE6\_Domino\_Deterministic\_Universe: Archived conversation contains potentially stale external or baseline claim: C89.
\item CONTRA-0201 UE6\_Domino\_Deterministic\_Universe: Archived conversation contains potentially stale external or baseline claim: C++98.
\item CONTRA-0204 ui\_editor\_tool\_editor\_planning: Archived conversation contains potentially stale external or baseline claim: C++98.
\item CONTRA-0205 ui\_editor\_tool\_editor\_planning: Archived conversation contains potentially stale external or baseline claim: Android.
\item CONTRA-0216 Workbench\_AIDE\_Product\_Spine: Archived conversation contains potentially stale external or baseline claim: C89.
\item CONTRA-0217 Workbench\_AIDE\_Product\_Spine: Archived conversation contains potentially stale external or baseline claim: C++98.
\item CONTRA-0219 World\_Architecture: Archived conversation contains potentially stale external or baseline claim: C89.
\item CONTRA-0220 World\_Architecture: Archived conversation contains potentially stale external or baseline claim: C++98.
\item CONTRA-0221 World\_Architecture: Archived conversation contains potentially stale external or baseline claim: ISO C89.
\item CONTRA-0222 World\_Architecture: Archived conversation contains potentially stale external or baseline claim: SDK.
\end{itemize}



{\small\raggedright SOURCE PATH: docs/\allowbreak{}archive/\allowbreak{}conversations/\allowbreak{}\_\allowbreak{}audit/\allowbreak{}UNCERTAINTY\_\allowbreak{}REGISTER.md\par}


\section{Uncertainty Register}

Uncertainty entries are extracted from archived reports and require current repo verification.


\subsection{advanced\_simulation\_infrastructure}

\begin{itemize}
\item Source file: docs/archive/conversations/advanced\_simulation\_infrastructure/dominium\_advanced\_simulation\_infrastructure\_architecture\_\_human\_readable\_chat\_report.txt
\item [UNCERTAIN / UNVERIFIED] No repository files were inspected in this chat. No code was implemented. Proposed paths, module names, command names, file names, VM concepts, and data schemas are design proposals until checked against the actual codebase.
\item [UNCERTAIN / UNVERIFIED] Exact solvers, fixed-point formats, module paths, data schemas, and tick ordering remain to be verified and specified.
\item [UNCERTAIN / UNVERIFIED] The exact gameplay rules runtime was not specified. The command examples were conceptual.
\item [UNCERTAIN / UNVERIFIED] The exact cross-section schema, attachment schema, and packing rules remain unresolved.
\item [UNCERTAIN / UNVERIFIED] Microsegment length, VM design, storage layout, budget values, and actual code integration remain open.
\item [UNCERTAIN / UNVERIFIED] Actual standards packs, rule schema, semantic key format, localization, and AI interaction remain unresolved.
\item Then it produced a prompt for another GPT-5.2 chat that already had a refactor/optimization plan. That prompt told the other chat to amend its existing plan instead of restarting, and to verify coverage of arbitrary placement, unified spatial primitives, co-location, signage, buildings, and determinism/performance.
\item [UNCERTAIN / UNVERIFIED] The architecture has not been verified against the actual repository. The exact implementation details remain unresolved: Q formats, orientation math, command schemas, VM instruction set, microsegment length, carrier ownership, junction archetypes, facility modules, and DECOR/device promotion policy.
\item [UNCERTAIN / UNVERIFIED] Exact Q formats were not finalized.
\item [UNCERTAIN / UNVERIFIED] The exact VM is unresolved.
\item The most important practical next step is to inspect the current repository. This chat proposed module names and file paths, but did not verify them.
\item [FACT] The user asked later reports to preserve uncertainty labels: FACT, INFERENCE, UNCERTAIN / UNVERIFIED, and PROJECT-CONTEXT.
\end{itemize}

\subsection{app\_runtime\_platform\_renderers}

\begin{itemize}
\item Source file: docs/archive/conversations/app\_runtime\_platform\_renderers/dominium\_app0\_runtime\_platform\_renderers\_\_human\_readable\_chat\_report.txt
\item What remains uncertain is whether the current repository already enforces these boundaries mechanically. The actual code needs inspection.
\item The unresolved issue is sharding semantics. APP0 requires sharding hooks, but the chat did not decide whether shards are spatial, logical, temporal, jurisdictional, or something else. The report preserved the idea that APP0 can provide sharding plumbing without inventing simulation semantics.
\item The unresolved issue is the trust and integrity model for setup. The chat did not decide whether content verification should use hashes, signatures, local manifests, network verification, or another model.
\item The unresolved issue is the exact elevation/audit model. The chat did not decide how tools authenticate, request write permissions, log actions, or interact with law enforcement modules.
\item The unresolved issue is the exact renderer capability schema. The chat proposed ideas but did not finalize fields, priorities, or target matrices.
\item The user asked to look at the old code snapshot in project attachments. A prior assistant claimed to inspect it, but later preservation documents marked that claim as **UNCERTAIN / UNVERIFIED**.
\item This became one of the most important unresolved issues. Almost every implementation question depends on the real code: directory layout, build system, languages, existing platform/render layers, client/server dependencies, and engine/game public APIs. The next real step is repository verification.
\item INFERENCE:** The user wants future assistants to be epistemically careful: not inventing repository facts, not treating proposals as decisions, and not losing uncertainty.
\item The exact module architecture remains unresolved. The exact platform/renderer support matrix remains unresolved. The exact relationship to the current repository remains unresolved because the old snapshot has not been verified. The exact setup trust model, tool audit model, sharding semantics, and client prediction policy remain unresolved.
\item This affects the client, server, launcher, setup, and tools. The client can display and interact; the server can host authority; the launcher can start sessions; setup can install and verify; tools can inspect and audit. But none of them should invent unauthorized simulation outcomes.
\item The assumption behind this decision is that engine/game code can be hosted through stable public APIs. Whether the current repo already supports that is **UNCERTAIN / UNVERIFIED**.
\item The unresolved caveat is prediction. If client-side prediction is later needed, it must be designed carefully as non-authoritative and probably separated from authority-committing simulation code.
\end{itemize}

\subsection{app\_testx\_codehygiene}

\begin{itemize}
\item Source file: docs/archive/conversations/app\_testx\_codehygiene/dominium\_app\_testx\_codehygiene\_handoff\_\_human\_readable\_chat\_report.txt
\item The main unresolved goal is practical execution: which prompts have actually been run, which files exist, and what implementation should happen next.
\item Setup handles install, verify, repair, rollback, upgrade, downgrade, uninstall, and status. Launcher can invoke Setup, but must not replicate its mutation logic. This prevents two products from disagreeing about installation state.
\item First, verify whether libs/contracts and libs/appcore already exist. Then implement the pure contract headers and appcore skeleton. This should include:
\item Before implementing app RepoX integration, verify actual RepoX metadata paths, TestX invocation, VALIDATE-0 commands, and BUILD-ID-0 files. Without that, an implementation prompt would guess.
\item Use labels like FACT, INFERENCE, UNCERTAIN, PROJECT-CONTEXT in reports.
\item From the user profile and instructions, future assistants should verify time-sensitive/current facts with web where relevant. This chat itself is mostly internal project planning, so external citations are not central.
\item EXEC, ECSX, KERN, ADOPT, DIST, HWCAPS, EXIST, DOMAIN, TRAVEL, TIME, OMNI, LIFE, CIV, WAR, AGENT, TOOL, MOD, FINAL prompt families were generated earlier. They are important as historical design artifacts, but actual repo execution is unverified.
\item The chat produced downloadable handoff files and later an in-chat reader. Those files are useful for aggregation but are not themselves project source code. The package's main caveat is that it does not verify actual repository state.
\item The most important unresolved issue is actual repository state. We do not know which prompts were run, which files exist, or which systems are implemented.
\item The latest project-context says CANON\_INDEX.md is the single source of truth. The package does not verify its existence. This should be checked before using docs as canon.
\item The versioning model is stated as canon, but actual implementation is unverified.
\item The project is too complex for vague summaries. Future assistants should preserve IDs, artifacts, and caveats.
\end{itemize}

\subsection{architecture\_codex\_prompts}

\begin{itemize}
\item Source file: docs/archive/conversations/architecture\_codex\_prompts/dominium\_domino\_architecture\_codex\_prompts\_\_human\_readable\_chat\_report.txt
\item UNCERTAIN / UNVERIFIED: No repository code was inspected or modified in this chat. The prompts and architecture are plans and artifacts, not proof that implementation exists. The actual repository state, build system, existing code quality, GUI backend availability, TLV schemas, and whether any prompts have been applied remain unverified.
\item Uncertainty remains about actual repository code and build system, but the layering itself is a strong established principle.
\item Unresolved issues include exact TLV schema field tags, exact registry implementation, base pack structure, and whether engine comments/docs may use examples with domain words.
\item The main uncertainty is implementation: the prompts define it, but no code was verified.
\item Unresolved details include exact process IO schemas, machine runtime state, and the relationship between autonomous machines and agent-operated jobs.
\item Unresolved: exact nesting policy, allowed packing modes, and how far to model real packing constraints.
\item Uncertainty remains about freeform geometry versus grid shell as the long-term building representation.
\item The exact destruction and structural-load models remain unresolved and part of Path D.
\item Unverified: actual save file format and current repo serialization code.
\item The most important caveat is that actual repo state is unknown.
\item Unverified: actual hash design and validators have not been implemented in this chat.
\item UNCERTAIN / UNVERIFIED: Actual implementation remains unresolved. The prompts are ready, but no code was verified.
\end{itemize}

\subsection{architecture\_ui\_providers}

\begin{itemize}
\item Source file: docs/archive/conversations/architecture\_ui\_providers/dominium\_architecture\_ui\_providers\_robot\_os\_strategy\_\_01\_human\_readable\_report.md
\item The exact first playable slice remains unresolved. The first provider wedge is clear, but exact sequence between client/workbench gameplay, Robot OS, and gameplay mechanics still needs a formal milestone. The exact Lua version, Linux baseline, current repo status, and external library/version support need verification. The exact degree of Unreal involvement remains unclear.
\item 1. Verify and pin external library/toolchain facts.
\item It remains uncertain how much the user wants Unreal in the near-term after raylib-first discussion. It also remains uncertain whether Lua 5.4 or 5.5 is preferred; the user appears to value pinning and stability more than "latest" for script ABI.
\end{itemize}

\subsection{Build\_and\_Future\_Proofing}

\begin{itemize}
\item Source file: docs/archive/conversations/Build\_and\_Future\_Proofing/Dominium\_Build\_and\_Future\_Proofing\_Architecture\_\_01\_human\_readable\_report.md
\item Plain-language limitation: this package is reliable as a preservation of the current visible chat. It should not be treated as a complete archive of every Dominium discussion in other chats. Where the report uses project or repository context, it is labelled as such. The most important caveat is that many items are recommendations, not final user decisions.
\item The unresolved goals are implementation goals: deciding which recommendations become canon, adding public-surface/dependency/build contracts, validating the current live repo state, and applying the proposed cleanup tasks. The preservation package is complete, but the engineering work it describes remains pending.
\item 2. **Verify current repo state.** Confirm the actual current HEAD, build proof, smoke tests, and outstanding warnings before acting.
\item 10. **Feed this report into the master spec book aggregator.** Preserve uncertainty labels and avoid merging suggestions as decisions.
\item The preservation output must distinguish FACT, INFERENCE, UNCERTAIN/UNVERIFIED, and PROJECT-CONTEXT.
\item A future assistant may rely on stale repo state. Mitigation: verify live HEAD, docs, and CI before implementation.
\item What live repo facts should I verify before implementing these tasks?
\item The final uploaded prompt asked for this preservation package. The key thing to preserve is that most architecture proposals are recommendations, not accepted user decisions yet. The next best action is to decide which recommendations become canon, then implement either STRUCTURE-01: Public Surface Registry or the build tuple contract work, after verifying the current live repo state.
\end{itemize}

\subsection{canonical\_structure\_and\_framework}

\begin{itemize}
\item Source file: docs/archive/conversations/canonical\_structure\_and\_framework/dominium\_canonical\_architecture\_repo\_foundation\_\_01\_human\_readable\_report.md
\item INFERENCE: The user wanted a durable method for deciding future structure changes, not just a one-time layout. They also wanted the assistant to become more rigorous: distinguish decisions from brainstorms, stop over-planning without execution, and preserve uncertainty. The user also wanted a structure that supports other Domino-based games and possibly different engine implementations.
\item A persistent uncertainty is the exact current repo state. The user supplied many status reports and directory exports, but several exports were internally inconsistent or stale at various points. The chat therefore repeatedly emphasized structure report integrity and verification before trusting tree analyses.
\item 1. Verify fast-strict/RepoX/structure report status from a clean current export.
\item The user dislikes vague reassurance and wants uncertainty labelled.
\item UNCERTAIN: The final exact convention for external/upstream versus external/vendor should be checked against current repo policy.
\item UNCERTAIN: Whether Lua should be pinned to 5.4 or 5.5 is not settled in this chat.
\item UNCERTAIN: Whether pack-internal content/ directories are legitimate pack law or legacy remains to be verified.
\item User-reported commits and status summaries: these include commits like 6e0dd93, 1406490, 3243fab, ce9ca, and others. Treat them as FACTs reported in the chat, but verify live repo state before acting.
\item Is full CTest currently green? The chat repeatedly says no or not run; verify before any release/trust claim.
\item Are stale AuditX/identity evidence and launcher marker debt still blocking fast strict? User reports changed over time; verify live.
\item Is the final live queue pointing to PROJECTION-CONFORMANCE-01, PRESENTATION-CONTRACT-01, or a maintenance task? Verify current .aide/queue/current.toml.
\item Should external/upstream or external/vendor be canonical? Verify current repo convention.
\end{itemize}

\subsection{Chronology\_Celestial\_Systems}

\begin{itemize}
\item Source file: docs/archive/conversations/Chronology\_Celestial\_Systems/Dominium\_Chronology\_Celestial\_Systems\_\_human\_readable\_chat\_report.txt
\item The assistant formalized this as the Perfect Earth Calendar, later called HPC-E. The final structure is clear, though the exact leap rule remains unresolved.
\item What remains uncertain is the exact verified object list. The assistant generated a real Milky Way system list, but that list must be checked before becoming data.
\item The unresolved issue is exact fidelity. The chat specified what should exist conceptually, but not the final data schema, map resolution, object attributes, or implementation priority.
\item What remains unresolved is the exact mathematical implementation of these frames and how much relativistic fidelity should exist in phase one.
\item The uncertainty is high relative to Earth. Many names and structures were assistant-generated and should be reviewed before implementation.
\item The unresolved issue is implementation reality. The actual Plan G/The Game codebase was not visible here, so any implementation plan must be reconciled with real architecture.
\item The only major unresolved part is the leap rule.
\item The user wants rigorous, structured, fact-aware assistance. The user specifically requested preservation of uncertainty labels and does not want assistant suggestions treated as user decisions.
\item Future assistants should not invent Plan G details, assume codebase language constraints, hardcode unverified astronomy, treat all assistant-generated calendars as final, flatten tentative items into decisions, or reveal time/date in UI without in-world justification.
\item The most important unresolved issue is the real Plan G/The Game architecture. This chat could not inspect it. Any implementation plan must first compare this design to the actual project.
\item The second major unresolved issue is the exact HPC-E leap rule. The calendar structure is set, but the leap cadence is not. This must be decided before implementation.
\item The third unresolved issue is the exact time-scale mapping for Gregorian January 1, 2000 AD. The user specified the civil date, but not the exact technical instant. This affects deterministic world start.
\end{itemize}

\subsection{development\_routes}

\begin{itemize}
\item Source file: docs/archive/conversations/development\_routes/dominium\_development\_routes\_\_human\_readable\_chat\_report.txt
\item The first answer was assertive. It stated that Dominium must follow Route C and described the route as the only viable one for Dominium. Later parts of the chat corrected the status of that claim. FACT: the assistant made that recommendation. UNCERTAIN / UNVERIFIED: the recommendation was not validated with external sources in the chat and was not explicitly accepted by the user.
\item The assistant produced a large Context Transfer Packet. It preserved the technical route proposal, the caveats around its uncertainty, the proposed phase order, the open questions, the risks, the artifacts, and the next actions.
\item The assistant created those files and reported that the chat label was inferred as **Dominium Development Routes**. The package was said to be safe for aggregation **with caveats**. The main caveats were that the substantive project transcript was short, Route C was not user-confirmed, and Dominium's genre, core loop, stack, implementation state, files, platforms, team, and timeline were not established.
\item UNCERTAIN / UNVERIFIED: Route C was not accepted by the user as final.
\item UNCERTAIN / UNVERIFIED: the required level of determinism for Dominium was not established. It might need full cross-platform deterministic replay, or it might only need ordinary save/load correctness, or something between those extremes.
\item The outputs included a Context Transfer Packet and a downloadable package with multiple files. The files are less important than the purpose: preserve the exact status of the discussion, including uncertainty and assistant-proposal status, so future assistants do not accidentally distort the project.
\item The second explicit goal was to preserve the chat before retirement. The user did not want a normal summary. They wanted a maximum-fidelity transfer packet that would prevent loss of decisions, constraints, unresolved issues, rationale, and next actions. This was addressed by the Context Transfer Packet.
\item Some goals changed over time. The chat began as architecture guidance, became a state-transfer task, became a packaging task, became an in-chat inspection task, and now becomes a plain-English briefing task. The underlying continuity goal stayed consistent: preserve the useful substance without losing uncertainty.
\item The chat-specific reporting decisions were more settled. The user explicitly required that the package be limited to this chat, use labels for fact/inference/uncertainty, preserve assistant suggestions as suggestions, and avoid inventing facts. These are firm process decisions because they came directly from the user.
\item The user's strongest visible concern is epistemic discipline: do not invent, do not silently infer, do not promote suggestions into decisions, and do not erase uncertainty.
\item Before building anything, the project needs a clearer definition of Dominium. The first unresolved technical step is to define the genre and core player loop. This matters because a grand strategy game, colony sim, city-builder, RTS, political simulator, and economic simulation all imply different state models and time scales.
\item The largest unresolved issue is whether Dominium is actually simulation-heavy enough to justify the proposed kernel-first deterministic route. The assistant assumed that kind of complexity, but the visible conversation does not prove it. The user needs to confirm or correct that assumption.
\end{itemize}

\subsection{documentation\_standards\_readme}

\begin{itemize}
\item Source file: docs/archive/conversations/documentation\_standards\_readme/documentation\_standards\_readme\_handoff\_\_human\_readable\_chat\_report.txt
\item The key thing to remember: this chat produced the standards and prompts for future work; it did not verify or change the project. The next assistant must inspect the actual repository before writing repo-specific docs, README content, platform claims, license claims, or subsystem descriptions.
\item The chat began with the user asking how to have Codex generate documentation for every file in the project. The user wanted metadata, functions, definitions, dependencies, and other maintainability information documented. The core uncertainty was placement: should documentation live in separate sister files, or should it live in file headers and function headers?
\item UNCERTAIN / UNVERIFIED: Whether the project has similarly named files or modules is unknown.
\item What remains uncertain is how the actual project currently organizes public headers, internal headers, and docs. That must be verified from the repository.
\item What remains unverified is the actual project name, supported platforms, license, maturity, and current repository layout. Those cannot be responsibly written without inspecting the repo.
\item The decision depends on an assumption: that the project's public APIs are declared in headers. That is likely for a C/C++ codebase, but the actual public/private boundary is still unverified.
\item What must happen first: verify source roots, exclusions, CMake structure, CI environment variables, and Python availability.
\item What could go wrong: treating assistant suggestions as final requirements or merging contradictory chat claims without preserving uncertainty.
\item FACT: The user wants facts checked and uncertainty preserved.
\item INFERENCE: The user values preserving rejected options and uncertainty because they expect to aggregate many old chats later.
\item UNCERTAIN / UNVERIFIED: Actual repository language standards, compiler settings, source roots, and assembly presence still need verification.
\item This is one of the most concrete artifacts. It describes how Codex should create a Python checker, integrate it into CMake, document it, and verify behavior.
\end{itemize}

\subsection{Dominium\_Architecture\_I}

\begin{itemize}
\item Source file: docs/archive/conversations/Dominium\_Architecture\_I/Dominium\_Architecture\_I\_\_human\_readable\_chat\_report.txt
\item What remains uncertain is the exact final scope of all simulation systems. Many systems were described in ambitious terms, but not all details are final or internally consistent. The systems layer especially needs canonicalisation before coding.
\item What remains uncertain is the current real-world status and capabilities of "Codex 5.1 Max." That is an external tool fact and requires verification before future use.
\item The key later decision was to skip top-level .txt devspec files because original Markdown files already existed in /docs/.... The visible chat does not show the full content of those docs. They should be preserved as referenced artifacts, but their exact content is **UNCERTAIN / UNVERIFIED**.
\item This affects dlocale, font rendering, platform path handling, data locale JSON files, mod APIs, UI, and documentation. The unresolved part is exactly how much retro support is actually feasible. Claims about old OSes, toolchains, SDL2, OpenGL, Direct3D, and macOS versions require external verification.
\item A large portion of the chat consisted of producing specs for those files. Each spec was meant to tell Codex how to implement that file. The conclusion is that the repository architecture is broad but structured. The unresolved issue is that many later files remain pending.
\item The unresolved issue is API consistency. Generated specs for memory and serialization are inconsistent across files, and some platform/render specs contain outdated or unverified external assumptions.
\item The most important unresolved issue is duplication and contradiction. dweather, dhydro, and dai\_core were generated in more than one form. The package does not choose a final version. A future assistant should not implement any of those blindly.
\item The user decided to skip the top-level .txt files because original Markdown docs already exist in /docs/.... This decision is final for the current workflow, but it depends on those docs actually existing and being current. That is **UNCERTAIN / UNVERIFIED** because the docs were not inspected in this chat.
\item The user explicitly required that the final handoff/report preserve facts, decisions, preferences, rejected options, unresolved issues, artifacts, and rationale. The user required labels for FACT, INFERENCE, UNCERTAIN / UNVERIFIED, and PROJECT-CONTEXT. The user required that assistant suggestions not be treated as user decisions unless accepted.
\item All artifacts should feed into the future Project Spec Book only after uncertainty and contradictions are handled. The file-spec prompts are especially important, but they require API cleanup before implementation.
\item The most important unresolved issue is canonicalisation. dweather, dhydro, and dai\_core each have duplicate conflicting specs. A future assistant needs to compare the versions and either select one or merge them carefully.
\item The next unresolved issue is core API consistency. dmem and dserialize are used inconsistently across generated specs. This must be fixed before coding.
\end{itemize}

\subsection{Dominium\_Architecture\_II}

\begin{itemize}
\item Source file: docs/archive/conversations/Dominium\_Architecture\_II/Dominium\_Architecture\_II\_\_human\_readable\_chat\_report.txt
\item What remains uncertain is the exact code-level API and component layouts. The chat produced specs and prompts, not final source code.
\item This is an important implementation detail, but some of it remains unverified. Lua 5.4.4 portability to old compilers needs testing. The exact Lua data schemas still need definition. The docs also still have JSON-oriented format sections, so the Lua-vs-JSON policy needs reconciliation.
\item The biggest unresolved goals are practical:
\item Verify DX9/Windows 2000, SDL2, Lua/toolchain compatibility.
\item FACT: This was answered in response to Codex's clarifying questions. It is a concrete implementation direction, though Lua 5.4.4 portability remains unverified.
\item 1. **Verify and apply V4 docs.
\item The user repeatedly requested maximum-fidelity reports, registers, handoff packets, and human-readable summaries. The user wants labelled uncertainty and does not want future assistants to invent missing context.
\item or forget that Lua-vs-JSON remains unresolved.
\item The most important unresolved issue is **actual repo state**. The chat generated many documents, but it is unknown whether they were applied to disk. Before Codex implements anything, the repo must be inspected.
\item Another major unresolved issue is **Lua data versus JSON formats**. The user explicitly wants Lua for all MVP data, but generated DATA\_FORMATS.md and data specs include JSON-style formats. This may be acceptable if JSON is long-term and Lua is MVP, but the docs must say that clearly.
\item The exact MVP component set is also unresolved. The chat implies Transform, occupancy, power node, fluid node, data node, worker/agent, building/device, prefab IDs, inventory/job state, and possibly simple vehicle/path components, but exact fields are not final.
\item Major unresolved issues are actual repo state, Lua-vs-JSON policy, render command/camera API, external compatibility, MVP component fields, and licence validity.
\end{itemize}

\subsection{Dominium\_Architecture\_III}

\begin{itemize}
\item Source file: docs/archive/conversations/Dominium\_Architecture\_III/Dominium\_Architecture\_III\_\_human\_readable\_chat\_report.txt
\item The user then uploaded dominium.7z, saying it was the git repository as of that moment. **FACT:** The prior assistant said it could not open .7z archives in that environment. That means the actual repository state remained unverified at that stage. This became important later, because many implementation ideas were discussed without confirmed access to the repo contents.
\item This created one of the main unresolved issues. DX12 was discussed earlier, but omitted from the latest renderer list. X11 and Wayland were discussed earlier, but the latest platform categories are POSIX, SDL1, SDL2, and Native. Therefore, the final status is:
\item DX12 is **UNCERTAIN / UNVERIFIED** and should not be assumed.
\item X11/Wayland placement is unresolved.
\item The prior assistant summarized these files as containing two different launcher concepts: an engine-rendered launcher view and a native Win32 launcher. It also claimed the CMake file built the Win32 source for non-Windows targets. However, this report must preserve that as **UNCERTAIN / UNVERIFIED**, because the files were not re-inspected during final packaging.
\item This is final as a user-stated product goal. Feasibility still requires verification.
\item Those outputs are important artifacts, but they are not the substance of the design. Their purpose is to preserve this chat's decisions, uncertainty, and future work for later aggregation into a master Project Spec Book.
\item What remains uncertain is not the constraint itself, but how existing repository code currently enforces it. That requires code inspection.
\item Uncertainty remains around how the uploaded launcher code currently maps to this desired UI. The prior assistant said the uploaded native launcher had different tabs, including Console and missing Accounts, but this must be verified.
\item Uncertainty remains around aliases. Earlier names such as dom-launcher and dom-setup might be kept as aliases, but that was not finalized.
\item The major unresolved issue is taxonomy: how exactly POSIX, SDL1, SDL2, and Native map to Win32, Cocoa, Carbon, X11, Wayland, headless, and retro platforms.
\item The unresolved part is feasibility. Supporting Linux 2.4, Mac OS 8.5, Windows 3, and MS-DOS requires separate toolchain research. The chat did not verify that every target can be built.
\end{itemize}

\subsection{Dominium\_Architecture\_IV}

\begin{itemize}
\item Source file: docs/archive/conversations/Dominium\_Architecture\_IV/Dominium\_Architecture\_IV\_\_human\_readable\_chat\_report.txt
\item The conclusion was clear and accepted: **Domino is the reusable engine; Dominium is the game/product layer**. What remains uncertain is the actual repository state. The chat planned the split, but did not verify whether the repo already matches it.
\item The unresolved issue is that the actual repo was not inspected in this chat. The planned structure may not match the current files.
\item The unresolved issue is practical feasibility, especially for retro targets and Carbon OS. Their actual toolchains and system calls remain unverified.
\item The unresolved issue is the exact IR payload layout and how far limited renderers should go. A CGA or MDA backend cannot realistically support modern 3D in the same way as Vulkan. The agreed approach is explicit capability flags and graceful fallback or no-op behaviour.
\item The unresolved issue is the exact manifest/schema format. JSON, TOML, and INI-like examples appeared in prompts, but no final schema choice was made.
\item The unresolved issue is that external packaging details require verification. WiX, macOS notarization, Linux packaging, and AppImage tooling are all external-world facts and may be stale.
\item The package and this report should help future assistants continue without re-reading the whole chat. However, they must preserve uncertainty: actual repo state is unverified, some transcript sections were skipped, external tooling facts may be stale, and assistant suggestions are not user decisions unless accepted or built upon.
\item It is also reasonable to infer that the user wants the future project spec book to preserve rationale, not just lists of decisions. The later handoff requests emphasised visible rationale, rejected options, uncertainty, and changes of direction.
\item Before doing it, verify the repo state. The generated prompts assume specific directories and may need adaptation.
\item Important items labelled as fact, inference, or uncertainty.
\item The **Phase 4 setup prompts** describe setup core and wrappers for Windows, macOS, Linux, and retro targets. They should be used only after verifying current external packaging tool details.
\item The chat also created downloadable report files and handoff packages. For understanding the substance of the chat, the important thing is not the files themselves but the information they preserved: roadmap, decisions, tasks, constraints, risks, and unresolved issues.
\end{itemize}

\subsection{Dominium\_Complete\_Conversation}

\begin{itemize}
\item Source file: docs/archive/conversations/Dominium\_Complete\_Conversation/Accompanying\_Report/Dominium\_Canon\_Portability\_Audit\_\_01\_human\_readable\_report.md
\item Uncertainty: the raw pasted prompt is not the current repo authority unless materialized in repo canon. The later audit found that it had in fact been materialized under docs/canon/constitution\_v1.md.
\item Uncertainty: during this preservation pass, a potential inconsistency appeared: some docs claim product roots moved under apps/, while an earlier inspected code path under root client/ was available and apps/client was not fetched successfully. This must be verified against the current physical tree.
\item Unresolved goals include verifying the latest physical repo tree, proving runtime implementation maturity, formalizing public ABI/API boundaries, completing layout/naming cleanup, and building a second Domino-based product to prove reuse.
\item 1. Verify the current physical repo tree against layout docs and contracts/repo/layout.contract.toml.
\item The user likely wants future assistants to challenge weak framing, preserve uncertainty, avoid shallow "best practice" lists, and focus on decisions that survive future rewrites.
\item "Create a gated plan for verifying current repo layout and exceptions."
\end{itemize}

\subsection{dominium\_domino\_codex\_planning}

\begin{itemize}
\item Source file: docs/archive/conversations/dominium\_domino\_codex\_planning/dominium\_domino\_codex\_planning\_\_human\_readable\_chat\_report.txt
\item The most important things to remember are: the project's strict architecture constraints, the Milestone-0 prompt sequence, the later shared CLI/TUI/GUI/input/render prompt sequences, the missing pack-system prompt, the unified startup policy, the unresolved repo-verification issue, and the need to treat generated prompts as plans rather than proof of implementation.
\item The assistant generated the first two prompts. The third was never produced because the conversation changed direction. That missing third prompt remains an important unresolved item.
\item However, the later context transfer packet treated that inspection claim as **UNCERTAIN / UNVERIFIED**. This was careful and correct: within the visible chat, there was no reliable tool trace proving the repo had actually been inspected. Therefore, future assistants must verify dominium.zip or the active checkout before relying on the platform/render/API answer.
\item What remains uncertain is how much of this architecture already exists in the current repo and how much is still planned. The chat produced prompts and plans, but it did not prove execution.
\item The unresolved issue is whether those prompts were ever run and whether the generated code, if any, conforms to the project's strict C89/C++98 and determinism requirements.
\item The unresolved complication is C89 portability. Standard C89 does not provide long long, but u64, i64, and q48\_16 require 64-bit representation. Future work must decide whether to allow compiler extensions, emulate 64-bit values for strict retro targets, or define platform tiers with conditional support.
\item The unresolved issue is that the prompt's minimal manifest parser and hardcoded platform path are only a starting point. A real implementation must generalize product discovery, avoid duplicated path logic, and preserve compatibility/versioning rules.
\item The unresolved issue is whether the prompts are too broad and whether the existing repo already has overlapping modules. Future work should inspect the repo before applying these prompts.
\item What remains uncertain is how much of this already exists in the repo and how realistic cross-platform IME is for early implementation. IME is often platform-specific and complex, so the generated prompt should probably be applied carefully or split.
\item The unresolved issue is API alignment. The prompt proposed new APIs, but the actual repo's dgfx API may already differ. The final report therefore recommends verifying current render headers and backend code before applying the prompt.
\item The unresolved work is to generate this prompt, ideally after inspecting actual content/pack docs and source files. The prompt should probably be staged, because pack integration touches content loading, dependency resolution, audio, rendering, shader compilation, GUI themes, vector icons, and streaming.
\item This topic remains important but unresolved in implementation. The policy is strong, but the implementation prompt needs repair to avoid linking problems. The safer approach is to keep the generic mode-selection algorithm independent of product-specific symbols, likely through callback tables or product-local dispatchers.
\end{itemize}

\subsection{dominium\_full\_conversation}

\begin{itemize}
\item Source file: docs/archive/conversations/dominium\_full\_conversation/companion\_report/dominium\_full\_conversation\_companion\_report\_\_01\_complete\_human\_report.md
\item Verify live repo state, then review/run FULL-GATE-LEGACY-TEST-ROUTE-01, then generate PACK-INTERNAL-LAYOUT-CANON-01 if the queue/status supports the maintenance lane.
\item Preserve uncertainty labels.
\item Future chats could easily misunderstand this conversation by treating Workbench as a GUI framework, treating raylib as the engine, resuming broad structure cleanup, skipping projection conformance, or assuming live repo status from stale pasted reports. The mitigation is to verify repo state first, preserve the contract/service/projection/provider distinctions, and follow the task queue.
\item The best next action is to verify live repo state, then review or run FULL-GATE-LEGACY-TEST-ROUTE-01, then generate PACK-INTERNAL-LAYOUT-CANON-01 if status supports it.
\end{itemize}

\subsection{dominium\_setup}

\begin{itemize}
\item Source file: docs/archive/conversations/dominium\_setup/dominium\_setup\_handoff\_\_human\_readable\_chat\_report.txt
\item FACT: The setup system was repeatedly defined as the only component allowed to install files, modify system/installation state, own installed-file metadata, repair installs, verify artifacts, uninstall, downgrade, upgrade, and roll back. This is the most important architectural rule from the chat.
\item What remains uncertain is the actual current implementation state in the repository after the applied Codex prompts. The design is clear, but the repo must still be inspected.
\item FACT: The final canonical setup layout uses setup/core/\{fetch,verify,install,rollback\}, setup/include/\{dsk,dsu\}, setup/packages, setup/platform, setup/tests, and setup/ui.
\item The unresolved part is whether the actual repository now matches this canonical layout.
\item UNCERTAIN / UNVERIFIED: Exact Visual Studio 2026 behavior and toolchain details should be verified in the current environment before relying on them.
\item UNCERTAIN / UNVERIFIED: The user stated the prompts were applied, but this chat did not show the final repository tree or build results.
\item UNCERTAIN / UNVERIFIED: The actual schema files under schema/setup/ need to be checked.
\item Several goals remain unresolved because the actual repo was not inspected here. We do not know whether the applied Codex prompts fully succeeded, whether setup builds, whether schemas exist, whether libs/contracts is complete, or whether setup/launcher CLI smoke tests pass.
\item FACT: The chat repeatedly established that setup is the only component allowed to install, upgrade, downgrade, repair, verify, uninstall, roll back, and own installed-file metadata.
\item FACT: The final setup layout is setup/core/\{fetch,verify,install,rollback\}, setup/include/\{dsk,dsu\}, setup/packages, setup/platform, setup/tests, and setup/ui.
\item This affects future application work, but some exact implementation paths remain unverified.
\item Another tradeoff was between broad future support and current implementation scope. The architecture allows legacy platforms, storefronts, and offline/online acquisition, but many implementation details remain unverified. The chat preserved those as goals without pretending they are already implemented.
\end{itemize}

\subsection{Domino\_Dominium\_Workbench}

\begin{itemize}
\item Source file: docs/archive/conversations/Domino\_Dominium\_Workbench/dominium\_workbench\_full\_conversation\_companion\_\_human\_readable\_detailed\_summary\_report.md
\item Reliability note:** This is a human-readable consolidation, not a live repository audit. Repo-current statements that came from pasted material or prior-chat excerpts remain **UNVERIFIED** until checked against the live julesc013/dominium repo and current AIDE/validator outputs.
\item Verify current live repo status and task gate state.
\item > Verify repo state, then generate COMMAND-RESULT-VIEW-SLICE-01.
\end{itemize}

\subsection{domino\_engine\_refactor\_prompts}

\begin{itemize}
\item Source file: docs/archive/conversations/domino\_engine\_refactor\_prompts/dominium\_domino\_engine\_refactor\_prompts\_\_human\_readable\_chat\_report.txt
\item The answer also introduced a deterministic bytecode VM as the preferred way to support modded behavior. A question was raised about whether native-code plugins should also be allowed, but the user did not answer. That remains unresolved.
\item What remains uncertain is the actual repository state. The chat designed architecture and prompts, but did not inspect source code. Therefore the plan needs repo verification before implementation.
\item The unresolved work is to generate detailed subsystem implementation prompts for heat, power, fluids, atmospheres, vehicles, and structural loads after the core engine spine is implemented.
\item The remaining uncertainty is plugin policy. A deterministic VM was proposed, but whether native-code plugins are allowed remains unresolved.
\item What remains uncertain is the exact thresholds and content policies for when things promote or demote.
\item The unresolved issue is how much of STRUCT compilation should be implemented first. The prompt plan is broad; a real implementation may need a smaller first milestone.
\item ignoring FACT/INFERENCE/UNCERTAIN labels;
\item UNCERTAIN / UNVERIFIED:** The package files were generated in the sandbox during the prior step, but repository files were not created or inspected.
\item The prompts are detailed, but repository state is unverified. Blindly applying them could create duplicate systems or wrong paths.
\item Before merging this into a master spec, preserve uncertainty labels and verify against repository state and later decisions.
\item The most important unresolved issues are VM/native plugin policy, exact Q formats, repository structure, existing TLV/scheduler state, and derived cache hash policy.
\item 9. "What should I verify in the repository before running Codex Prompt 1?"
\end{itemize}

\subsection{engine\_baseline\_architecture}

\begin{itemize}
\item Source file: docs/archive/conversations/engine\_baseline\_architecture/domino\_dominium\_engine\_baseline\_architecture\_\_01\_human\_readable\_report.md
\item Uncertainty: the repo has many modules and some historical/legacy surfaces, so exact implementation maturity varies by subsystem. The distinction should remain a formal requirement in a master spec.
\item Uncertainty: the full determinism envelope across all current modules was not independently re-run in this chat. User-supplied local validation results should be preserved but verified before being used as audit proof.
\item Uncertainty: this remains strategic architecture, not an implemented decision.
\item Unresolved goals include:
\item User wants uncertainty labelled and framing corrected when evidence disagrees.
\item UNCERTAIN: whether the user wants Unreal integration at all as a client/editor.
\item UNCERTAIN: exact tolerance for using external geometry libraries.
\item UNCERTAIN: whether broad repo restructuring will be desired after Milestone 0.
\item UNCERTAIN: exact visual style, UI stack, or programming language evolution.
\end{itemize}

\subsection{Foundation\_Workbench\_Codex}

\begin{itemize}
\item Source file: docs/archive/conversations/Foundation\_Workbench\_Codex/Dominium\_Foundation\_Workbench\_Parallel\_Codex\_Preservation\_\_01\_human\_readable\_report.md
\item What remains uncertain is how quickly this architecture will move from contracts and fixtures into runtime implementations and product slices. The next chat should not assume runtime provider resolvers, package runtime, or Workbench shell exist.
\item verify latest live repo state after user-pasted Wave 1 completion;
\item 1. Verify live repo state.
\item The user wants direct, evidence-grounded, audit-ready answers; timestamps/model labels; explicit uncertainty; no repeated clarification unless necessary; large continuous Codex prompt blocks when generating prompts; targeted tests instead of full suites; and rapid progress toward Workbench/code.
\item The user is likely to reject slow, report-only, overcautious process. Future assistants should produce executable next prompts quickly after verifying state.
\item It is uncertain how many concurrent Codex workers the user will actually run at once and whether the user wants coordinator prompts or worker prompts next.
\item The most important continuation is: verify current origin/main, confirm Wave 1 and Workbench validation state, then generate or run COMMAND-RESULT-VIEW-SLICE-01.
\end{itemize}

\subsection{Framework\_Open\_Source\_Provider}

\begin{itemize}
\item Source file: docs/archive/conversations/Framework\_Open\_Source\_Provider/Domino\_Framework\_Open\_Source\_Provider\_Architecture\_\_01\_human\_readable\_report.md
\item The chat inspected julesc013/dominium and found evidence of existing abstractions such as system and graphics layers, stubs, and soft-backed backends. The finding was that raylib could fill concrete provider gaps without replacing the architecture. However, current repo facts are stale/uncertain and must be verified, especially the C17/C++17 vs C90/C++98 baseline contradiction.
\item The exact first implementation branch, versions, and dependency pins are unresolved. The exact Domino Framework ABI is unresolved. The exact Lua version is unresolved. The current repository baseline must be verified. The formal policy for GPL/LGPL/unclear license material is unresolved. The sparse simulation and CAD systems need formal specs and prototypes.
\item The raylib decision was also directionally accepted. The user repeatedly expressed enthusiasm for raylib and asked about using as much of raylib infrastructure as possible. The caveat is that raylib is a provider suite, not architecture. This affects rendering, audio, input, asset preview, and Workbench bootstrap.
\item 1. Verify repo baseline and dependency versions.
\item The user requires source-grounded, audit-ready, uncertainty-labelled responses. The preservation prompt explicitly requires FACT/INFERENCE/UNCERTAIN/PROJECT-CONTEXT labels, no invented facts, no treating brainstorms as decisions, and no over-compression. The prompt also requires a human-readable report first and downloadable files if possible ?filecite?turn29file0?.
\item The most blocking unresolved issues are repo baseline verification, exact dependency pins, minimal framework ABI, provider profile order, deterministic simulation spec, and license policy.
\item Verify the current julesc013/dominium CMake language baseline.
\end{itemize}

\subsection{gui\_binary\_content}

\begin{itemize}
\item Source file: docs/archive/conversations/gui\_binary\_content/dominium\_gui\_binary\_content\_handoff\_\_human\_readable\_chat\_report.txt
\item GPT-5.5 Pro - 2026-05-27 00:00:00 Australia/Melbourne; time UNVERIFIED
\item whether knowledge systems should include uncertainty, false beliefs, or lost knowledge,
\item This stage established a pattern that continues throughout the chat: the user does not want impressive-looking prompts that hide unresolved design choices. The user wants the assumptions exposed before anything is formalized.
\item The assistant agreed with the general direction, but added caveats. The most important were:
\item Those caveats were not final user decisions, but they are important warnings.
\item The conclusion was not "run the prompt now." The conclusion was that CONTENT0 is a strong foundation but needs discussion before finalization. The assistant-generated rewrite is useful but not final. The unresolved issues need to be resolved before Codex or any implementation work touches the repository.
\item What remains uncertain is how exactly to encode these ideas in schemas and data. For example, it is not yet decided how deterministic seed hierarchies work, how macro content refines into micro content, or how timeline causality should be represented.
\item The unresolved issue is the shared backend/UI contract. Without that, the GUI plan is only a list of technologies.
\item The conversation did not yet define how those product backends are exposed. That remains the most important unresolved design issue.
\item The assistant caveat was that old .NET support must not be faked. Supporting .NET 4.0-era machines is not the same as targeting .NET 4.8. This remains a technical verification item, not a fully resolved decision.
\item The assistant caveat was that old macOS support requires serious toolchain planning. A current Xcode version may not target macOS 10.9. AppKit also does not necessarily need to be Intel-only. These are important but not fully verified within this chat.
\item The unresolved question is whether macOS 10.9 is a hard requirement or an aspirational compatibility goal.
\end{itemize}

\subsection{Language\_Platform\_Architecture}

\begin{itemize}
\item Source file: docs/archive/conversations/Language\_Platform\_Architecture/Dominium\_Language\_Platform\_Architecture\_Baseline\_\_01\_human\_readable\_report.md
\item The 32-bit question followed. The answer was to make full native products 64-bit and keep 32-bit as constrained or projection support only. The key caveat was not to make native pointer size part of game law. Save, replay, network, IDs, and renderer IR must use fixed-width types and never serialize size\_t, long, uintptr\_t, or raw pointers.
\item Exact tolerance for exceptions/RTTI inside private C++17 code remains unsettled. Exact platform floors and raylib/SDL2 priority remain unresolved.
\item Verify the Windows 7 and macOS 10.9 toolchain/library assumptions using official sources.
\end{itemize}

\subsection{launcher\_app\_layer}

\begin{itemize}
\item Source file: docs/archive/conversations/launcher\_app\_layer/dominium\_launcher\_app\_layer\_handoff\_\_human\_readable\_chat\_report.txt
\item These prompts matter because they superseded earlier advice. They also created a new practical priority: verify that the applied prompts actually succeeded.
\item This changed the correct future behavior. The next assistant must not redesign. It must implement, audit, maintain, document, verify, and harden.
\item This matters because it prevents platform-specific drift. If CLI has verify, GUI must not implement a different hidden version of verify. If GUI has a pack enable flow, it should map to the same command/intent as CLI and TUI.
\item The final purity prompt required those to be moved or quarantined. However, actual repo success remains unverified in this chat.
\item The important change is that the conversation moved from "What architecture should we use?" to "Architecture is locked; how do we verify and harden the implementation?"
\item The largest unresolved goal is verifying the actual repo. The user stated that two Codex prompts were applied, but this chat did not independently verify the filesystem, build, tests, scripts, or launcher docs.
\item Another unresolved goal is confirming whether the launcher hardening prompt has been applied.
\item A third unresolved goal is confirming whether command graph, UI IR, binding validation, zero-pack tests, RepoX changelog display, and BUILD-ID mismatch refusal are implemented.
\item This decision depends on CMake and script enforcement. It remains operationally unverified until the current repo is checked.
\item The launcher may expose repair or verify entry points in the sense that it can invoke setup through a contract, but it must not perform install mutation itself. This separation protects auditability and prevents hidden mutation.
\item This decision reflects the state of the project: after major refactors and applied prompts, the next risk is planning on assumptions. The correct move is to audit, document, verify, and harden the current launcher implementation.
\item stale or unverified repository assumptions,
\end{itemize}

\subsection{Launcher\_Setup\_Architecture}

\begin{itemize}
\item Source file: docs/archive/conversations/Launcher\_Setup\_Architecture/Dominium\_Launcher\_Setup\_Architecture\_\_human\_readable\_chat\_report.txt
\item Uncertainty:** The exact repository layout and existing implementation were not inspected in this chat.
\item Uncertainty:** The final language boundary for setup remains unresolved after the later C89/dsys/dgfx launcher direction.
\item Uncertainty:** The exact mapping into runtime binaries remains unverified.
\item Uncertainty:** Accounts/social/wiki/forum/direct messaging are future-facing and were not converted into the final five-tab implementation plan.
\item Uncertainty:** Earlier setup designs used JSON manifests. Later dsys/dgfx architecture leaned toward TLV .dmeta files.
\item Uncertainty:** Actual plugin ABI details and dynamic loading support depend on dsys/repo APIs.
\item Uncertainty:** Actual dsys/dgfx APIs were not inspected.
\item Conclusion reached:** This chat is now documented and can be merged later, but with caveats.
\item INFERENCE:** The user wanted to avoid later assistant confusion by preserving rejected options, uncertainties, and changes of direction.
\item The unresolved goals are mostly implementation-definition questions:
\item The main assumption throughout is that the underlying project either already has, or will have, the shared libraries and platform abstractions described. That assumption is unverified in this chat.
\item The user wants strict engineering tone, detailed reasoning, and preservation of uncertainty. In this specific request, the user asked for prose rather than a register dump.
\end{itemize}

\subsection{Modularity\_AIDE\_Refactorability}

\begin{itemize}
\item Source file: docs/archive/conversations/Modularity\_AIDE\_Refactorability/Dominium\_Modularity\_AIDE\_Refactorability\_Architecture\_\_01\_human\_readable\_report.md
\item The live repo still needs verification. The exact AIDE files, contracts, validators, and migration tasks must be implemented. It remains unresolved which existing roots and tools map to which final homes, which APIs become stable, which code is reusable across products/games/projects, and which prior assistant recommendations should become formal requirements after user review.
\item 1. **TASK-09** - Verify live repo baseline: current branch, root inventory, validators, CTest status, generated artifacts.
\item The user prefers source-grounded, audit-ready decisions; bounded uncertainty; explicit task prompts; and reusable engineering doctrine. Future assistants should avoid shallow affirmation and should distinguish accepted user decisions from assistant recommendations.
\item Important caveat: the only current uploaded file available for this task is Pasted text.txt, which contains the preservation-package prompt. Earlier generated handoff files, ZIP packages, or uploaded docs referenced in prior conversation are not available in this chat unless the user re-uploads them.
\item The most serious risk is that a future assistant treats this chat as if it verified the live repo. It did not. The second serious risk is that it hardens every assistant suggestion into a requirement. Several items are strong recommendations aligned with user goals, but still require implementation review. Future chats must preserve labels and verify stale facts.
\item Future assistants must verify live repo state before implementing cleanup.
\end{itemize}

\subsection{omega\_xi\_pi\_architecture\_future}

\begin{itemize}
\item Source file: docs/archive/conversations/omega\_xi\_pi\_architecture\_future/dominium\_omega\_xi\_pi\_architecture\_future\_proofing\_planning\_\_01\_human\_readable\_report.md
\item Preserve FACT / INFERENCE / UNCERTAIN labels.
\item Do not invent or flatten uncertainty.
\item 7. "List all artifacts I should verify in GitHub main."
\item 9. "What should I verify before starting Sigma?"
\end{itemize}

\subsection{OS\_Interface\_Repo\_Architecture}

\begin{itemize}
\item Source file: docs/archive/conversations/OS\_Interface\_Repo\_Architecture/Dominium\_OS\_Interface\_Repo\_Architecture\_\_01\_human\_readable\_report.md
\item The unresolved goals are implementation-level. The repo still needs command dispatch unification, product boot proof, portable projection proof, AppShell rendered-mode law update, Workbench module contracts, document/patch/result/refusal schemas, and a boot-to-replay MVP. It also needs continued exception retirement, CTest/RepoX remediation, and verification of current build status.
\item 2. Repair canonical verify test discovery.
\item The user wants uncertainty labels and correction of incorrect framing.
\item This preservation task specifically required FACT / INFERENCE / UNCERTAIN / PROJECT-CONTEXT labels, stable IDs, registers, spec sheets, an aggregator packet, self-audit, and downloadable files if tools are available.
\item UNCERTAIN: Whether the user wants to rename engine conceptually to kernel in code, or only use kernel as a conceptual layer.
\item UNCERTAIN: Whether the final public product name should be Dominium Workbench, Dominium Studio, Domino Workbench, or something else.
\item UNCERTAIN: How aggressive the next actual code refactor should be versus staged adapters and codegen.
\end{itemize}

\subsection{platform\_renderer\_api\_plan}

\begin{itemize}
\item Source file: docs/archive/conversations/platform\_renderer\_api\_plan/dominium\_codex\_platform\_renderer\_api\_plan\_\_human\_readable\_chat\_report.txt
\item This is not a fatal problem for the prompt pack, because the prompts themselves repeatedly tell Codex to discover files with rg, tree, type, CMake commands, and target listing. But it is a major caveat for any future assistant: do not assume the actual repo matches every file path or API name unless checked.
\item UNCERTAIN / UNVERIFIED:** The conversation did not prove that these prompts were actually applied to the repository. They are an assumed prerequisite, not verified code state.
\item Uncertainty remains about actual repo state. The prompts are designed to discover the repo state, but the chat itself did not inspect the archive.
\item The final active prompt pack narrowed runtime completion to Tier-1 on Windows and compile-gated correctness for Tier-2. This is a practical compromise because Codex runs on Windows. But it also creates a caveat: if "fully complete" later means fully running X11/Wayland/Cocoa/Metal/Vulkan, more prompts or native CI validation will be required.
\item This became one of the final "done" gates, along with scripts\textbackslash{}build\_codex\_verify.bat.
\item This matters because the conversation was long and contained multiple plans, superseded options, caveats, and active artifacts. The report package and this plain-language report exist to prevent future assistants from restarting, re-asking, or losing distinctions between decisions, brainstorms, and uncertainties.
\item The user also seems to want future assistants to operate with high continuity and not lose subtle distinctions. This is inferred from the repeated requests for maximum-fidelity transfer, stable IDs, labels, uncertainty tracking, and report packaging.
\item Another unresolved goal is whether the master engine prompts 1-14 are truly implemented. The later plan depends on them, but this chat did not verify the repo.
\item This is a decision with a caveat. Because Codex runs on Windows, the final prompt pack treats non-Windows Tier-2 platforms as host-gated and compile-correct when built on their native hosts, not fully runtime-validated on Windows. This makes practical sense, but it may not satisfy a stronger interpretation of "fully complete" if the user later demands runtime Tier-2 validation.
\item Fifth, the user wanted **future spec-book construction**. That means uncertainty and provenance matter. The chat repeatedly preserved labels such as FACT, INFERENCE, and UNCERTAIN / UNVERIFIED. This prevents future aggregation from treating assistant suggestions or brainstorms as binding decisions.
\item cmake -S . -B build\textbackslash{}verify\_initial -G Ninja -DCMAKE\_BUILD\_TYPE=Debug
\item cmake --build build\textbackslash{}verify\_initial --target help
\end{itemize}

\subsection{platform\_support}

\begin{itemize}
\item Source file: docs/archive/conversations/platform\_support/dominium\_platform\_support\_planning\_\_human\_readable\_chat\_report.txt
\item UNCERTAIN / UNVERIFIED: The name "Domino" was never confirmed by the user. Future work should treat "Domino" as an assistant-created placeholder unless the user explicitly accepts it. The useful idea is the distinction between full product support and reduced engine/core support, not necessarily the name.
\item The main unresolved issue is the exact definition of support. Does Tier-0 mean identical feature parity, simultaneous release, identical save compatibility, identical UI, equal QA coverage, or some platform-specific differences? The assistant inferred that Tier-0 should mean feature parity, full QA, long-term support, and high release priority, but the exact contractual meaning still needs formal definition.
\item What remains unresolved is the exact PC baseline. The chat did not decide whether Windows 10, Windows 11, or both are required. It did not decide whether Windows ARM64 matters. It did not decide whether macOS support includes Intel Macs, Apple Silicon Macs, or both. It did not define Linux distributions, glibc requirements, Flatpak/AppImage/deb/rpm packaging, X11/Wayland policy, or Steam Deck verification.
\item UNCERTAIN / UNVERIFIED: The chat did not establish minimum Android API level, ABI list, target graphics APIs, Play Store rules, Android TV status, Android Go status, Automotive status, Chromebook status, or vendor ROM support. These need future decisions and current-source verification.
\item The user also listed iPod Touch, Apple TV, Apple Watch, Mac, macOS, tvOS, watchOS, and Apple Vision Pro indirectly through AR. These were not classified in detail. INFERENCE: macOS belongs under PC. UNCERTAIN / UNVERIFIED: tvOS, watchOS, and visionOS remain undecided. Full Dominium on Apple Watch is not established and should not be assumed.
\item The unresolved issues are substantial. The chat did not decide whether WebGPU is required, whether WebGL fallback is mandatory, whether the game must work offline as a PWA, whether WASM threading/SIMD is required, how browser storage will work, how asset loading will be handled, or what browser versions will be supported.
\item UNCERTAIN / UNVERIFIED: Current console platform facts were not verified in this chat. Future work must verify official PlayStation, Xbox, and Nintendo developer requirements from current official sources before implementation planning.
\item INFERENCE: PC handhelds should be subtargets under PC, with their own QA and UI profiles. UNCERTAIN / UNVERIFIED: Steam Deck's exact tier was not finally decided.
\item UNCERTAIN / UNVERIFIED: XR was never promoted to top-tier. Its priority remains unresolved. Future work should not let XR quietly expand the scope of the main product unless the user explicitly elevates it.
\item INFERENCE: The user was also trying to prevent vague "support everything" language from hiding real engineering constraints. The repeated requests for decisions, constraints, unresolved issues, verification items, and future actions imply that the user wanted a plan that could later become part of a formal project specification.
\item The biggest unresolved goal is turning the platform direction into a real, implementable matrix. The chat did not define exact minimum versions, APIs, architectures, release timelines, engine/toolchain, build system, store strategy, or QA device list.
\item The decision depends on an unresolved assumption: that the engine/toolchain can support Android well enough for the intended Dominium experience. That assumption needs verification after engine/toolchain choice.
\end{itemize}

\subsection{Portability\_Assurance\_Future\_Proof}

\begin{itemize}
\item Source file: docs/archive/conversations/Portability\_Assurance\_Future\_Proof/Domino\_Dominium\_Portability\_Assurance\_Future\_Proof\_Architecture\_\_01\_human\_readable\_report.md
\item The chat proposed using standards such as DO-178C, NIST SSDF, OWASP ASVS, SLSA, and SPDX as design inputs only. The proposed internal version is DDAP v0 with DIL levels. Uncertainty: user has not explicitly ratified DDAP.
\item The actual repo was not inspected. The exact accepted directory structure, API policy, DDAP profile, compatibility promises, and first pilot module remain unresolved.
\item Caveat: carry forward as INFERENCE / assistant recommendation.
\item Caveat: carry forward as FACT goal + INFERENCE architecture.
\item Caveat: carry forward as assistant recommendation.
\item Implications: Safe for aggregation only with caveats.
\item Caveat: carry forward as FACT / UNCERTAIN scope.
\item 10. Feed this package into the master Project Spec Book with caveats.
\item The user explicitly required human-readable preservation, uncertainty labels, no invented facts, no silent inference, and no treating assistant recommendations as user decisions. The user explicitly values portability, modularity, extensibility, reuse, replaceability, future-proofing, and proper long-lived engineering.
\item Resolution path: Verify against official/current sources before normative text.
\item The uploaded prompt requested this preservation package. The main caveat is that the package preserves visible chat context and the uploaded prompt, not a guaranteed hidden full transcript or inspected repository. The best next action is to inspect the actual repo and turn the strongest recommendations into a small set of enforceable policies and one conformance-test pilot.
\end{itemize}

\subsection{readme\_ports\_determinism}

\begin{itemize}
\item Source file: docs/archive/conversations/readme\_ports\_determinism/dominium\_readme\_ports\_determinism\_\_human\_readable\_chat\_report.txt
\item What remains uncertain is whether the actual repository README.md matches the final pasted version. The chat only saw pasted text and generated prompts; it did not inspect a Git repository.
\item The unresolved issue is whether a metadata-only /ports directory still violates the user's preference. The final README keeps /ports in a limited form, but the user's wording could mean no /ports directory at all. That should be clarified before writing the directory contract.
\item The unresolved work is to define the actual network protocol, prediction model, rollback window, reconciliation logic, and state-hash validation.
\item The major unresolved goal is deciding exactly what "no separate ports directory/system" means physically in the repository. The final README allows /ports as metadata-only. That may satisfy the user, but it is not confirmed. If the user meant no /ports directory at all, the README still needs another revision.
\item Another unresolved goal is making the README and normative specs align. The README says specs are authoritative, but those specs were not inspected in this chat.
\item UNCERTAIN / UNVERIFIED: /ports as metadata-only is the current README state, but may not be fully accepted.
\item A related rejected idea was **retro build flows under /ports/<target> containing source or alternative implementations**. The final README superseded this with metadata-only /ports, assuming /ports is retained at all. This rejection is final for source code and behavior, but the existence of /ports itself remains uncertain.
\item The first next step is to verify the actual repository README.md against the final pasted README. The chat only used pasted text. If the repository differs, future work should start from the actual file, not from memory.
\item Possible blockers include the unresolved /ports question, missing access to actual repository files, and unspecified algorithms for RNG, state hashing, CRCs, and content hashing.
\item The reporting constraints later in the chat were also explicit: preserve facts, inferences, uncertainty, rejected options, contradictions, artifacts, prompts, and rationale. Do not invent. Do not silently infer. Do not treat assistant suggestions as user decisions unless accepted.
\item The user's profile and instructions emphasize bluntness, rigor, directness, and fact-checking. In this chat, they repeatedly asked for high-fidelity state preservation and explicitly rejected over-compressed summaries. Future assistants should avoid vague reassurance and instead give concrete analysis, caveats, and next actions.
\item The ninth artifact was the final README pasted by the user. It is the current textual baseline, but it is unverified against the actual repository.
\end{itemize}

\subsection{Refactor\_Architecture}

\begin{itemize}
\item Source file: docs/archive/conversations/Refactor\_Architecture/Dominium\_Domino\_Refactor\_Architecture\_\_human\_readable\_chat\_report.txt
\item Reliability note: Repository state, Codex execution, exact version numbers, and actual file modifications are **UNCERTAIN / UNVERIFIED** unless explicitly stated otherwise.
\item The most important thing to remember is that this chat did not implement the refactor. It designed the architecture and generated prompts and handoff material. Any future assistant must verify actual repository state before assuming files were moved, prompts were applied, or code was updated.
\item UNCERTAIN / UNVERIFIED:** It did not establish that this system already exists in code.
\item What remains uncertain is the exact degree to which existing code currently respects that boundary. The user provided a repo tree, but this chat did not inspect or modify code directly.
\item Unresolved details include exact default DOMINIUM\_HOME paths per OS and exact implementation of import/GC.
\item Unresolved details include exact schemas, parser choice, tests, and actual code implementation.
\item The key unresolved point is implementation timing. The report preserves this architecture, but the initial Codex refactor was not supposed to rewrite the entire UI system immediately.
\item Actual implementation is unresolved. This chat generated architecture and prompts, not a verified changed codebase.
\item Specific unresolved goals include:
\item The immediate plan is not to redesign everything again. The immediate plan is to verify current repo state and then implement the structural refactor in a controlled way.
\item This matters because the chat produced prompts, but did not verify that they were applied. If Codex has not run, the next step is to run or refine the master refactor prompt. If Codex has run, the next step is the consistency pass.
\item Preserve FACT / INFERENCE / UNCERTAIN labels in transfer/report contexts.
\end{itemize}

\subsection{Refactor\_Control\_Plane}

\begin{itemize}
\item Source file: docs/archive/conversations/Refactor\_Control\_Plane/AIDE\_XStack\_Dominium\_Refactor\_Control\_Plane\_\_01\_human\_readable\_report.md
\item Unresolved goals include finishing AIDE Q35-Q57, stabilizing Dominium CTest/RepoX blockers, defining final AIDE install/upgrade bundles, implementing root recycling machinery, deciding when product boot proof can proceed, and later deciding whether AIDE Runtime/Gateway/Hosts are truly needed.
\item 1. Verify current Dominium HEAD and POST-CONVERGE status.
\item The user values direct, source-grounded, audit-ready, detailed answers. The user asked for preservation packages that are human-readable and machine-aggregatable. The user repeatedly prefers not to lose nuance, rejected options, uncertainties, or artifacts. The user wants current facts verified when possible and wants future assistants to avoid re-asking answered questions.
\item UNCERTAIN: Exact public-facing naming for AIDE Core/Kernel remains partly open. Exact threshold for accepting CTest-warning status remains a user/repo governance decision. Exact final AIDE installation packaging details remain pending.
\item Future assistants must verify live repo state before implementing.
\item "Merge this with another chat's preservation package without losing uncertainty labels."
\end{itemize}

\subsection{Release\_Identity\_and\_Versioning}

\begin{itemize}
\item Source file: docs/archive/conversations/Release\_Identity\_and\_Versioning/Dominium\_XStack\_Release\_Identity\_and\_Versioning\_\_01\_human\_readable\_report.md
\item The chat rebuilt SemVer from first principles and decided that strict SemVer should apply only to components with declared public contracts. Candidate true-SemVer entities include SDKs, engine libraries, tool APIs/CLIs, protocol libraries, plugin hosts, schema libraries, and reusable runtime libraries. It remains unresolved which exact repo modules qualify.
\item GBN should identify CI/public/internal builds and appear as provenance, e.g. +gbn.7137, but it must not be used as SemVer precedence. BII should primarily be structured manifest metadata, because it may encode lane, target, kind, configuration, and other build identity fields. The exact BII schema remains unresolved.
\item The main unresolved goals are to produce formal specs: Release Constitution, SemVer Component Inventory, Suite Release Policy, Build Identity Spec, Channel/Lifecycle Spec, Artifact Naming Spec, Manifest Schema, Target Taxonomy, and Capability Registry.
\item It is still uncertain how strict the user wants suite X.Y.Z field meanings to be, whether capabilities should fully replace compatibility profiles or coexist with them, and how much old-platform support should be real versus aspirational.
\item The most important unresolved issue is formal classification. Without knowing which entities are strict SemVer components, suite versions, product release IDs, build fields, or capabilities, the rest of the policy cannot be enforced safely.
\item Platform targeting remains unresolved. The chat concluded that coarse families like WinNT5, WinNT10, MacOSX4, Linux5, or DOS5 are useful support labels, but exact binary compatibility needs target baselines and runtime/toolchain profiles. Linux kernel major alone is not enough. Windows 9x/NT/NT5/NT6/NT10 likely need separate build lanes or at least explicit tested baselines.
\end{itemize}

\subsection{testx\_repox\_governance}

\begin{itemize}
\item Source file: docs/archive/conversations/testx\_repox\_governance/dominium\_testx\_repox\_governance\_handoff\_\_human\_readable\_chat\_report.txt
\item This topic remains unresolved in implementation because the repo snapshot showed .dominium\_build\_number and update\_build\_number.cmake, but the chat did not inspect their contents. Those files may still implement older build-number behavior. That must be verified.
\item The unresolved part is implementation verification. The repo has Sol/Earth content in data and game/content; the engine must not assume those exist.
\item The unresolved part is how much of TESTX is actually implemented in the repo. The repo contains many tests and scripts, but this chat did not inspect them. A rules-to-checks audit is needed.
\item This remains conceptually accepted in the chat, but implementation is unverified.
\item The unresolved goals are mostly implementation and verification:
\item The repo already has UI IR/codegen infrastructure. The next step is to verify whether GUI/TUI truly bind to the CLI command graph and whether UI is data rather than logic.
\item The user has Windows 10 + VS2022. The next step is to verify v141/v141\_xp/SDK components and create isolated CMake presets for XP/Win7/Win8/Win10/Win11 builds.
\item Acting before verifying repo state.
\item The user also explicitly required the report style to preserve facts, inferences, uncertainties, rejected options, artifacts, timelines, risks, and visible rationale.
\item This recommendation is useful for immediate practical builds but requires verification of installed components and current toolchain facts.
\item Manual product SemVer is required, but the actual file layout is unverified.
\item The repo has UI tooling, but this chat did not verify that GUI/TUI bind to CLI command graph or that tools are read-only by default.
\end{itemize}

\subsection{Timekeeping\_and\_2038\_Resilience}

\begin{itemize}
\item Source file: docs/archive/conversations/Timekeeping\_and\_2038\_Resilience/Dominium\_Timekeeping\_and\_2038\_Resilience\_\_01\_human\_readable\_report.md
\item What remains uncertain is the current exact status of specific libc/OS/toolchain combinations, because those facts can change and were not reverified during this preservation turn.
\item What remains uncertain is the precise target list of legacy machines/toolchains and how far compatibility must go, especially for 16-bit compilers with weak or absent 64-bit arithmetic.
\item The main unresolved goals are backend audit, ACT serialization audit, civil/astronomical projection design, and verification of external platform facts.
\item The main uncertainty is completeness of repo inspection. The assistant read key files, but not all serialization paths and not all platform backends.
\item 6. Verify current platform/library facts before using them in formal public documentation.
\item The preservation prompt explicitly requires a human-readable report first, not only a machine-readable handoff. It requires uncertainty labels, stable registers, preservation of artifacts, no invented facts, and downloadable files if tools are available. ?filecite?turn30file0?
\item The user prefers source-grounded, audit-ready technical reasoning, direct structure, and careful distinction between facts, recommendations, and unresolved issues. PROJECT-CONTEXT indicates Dominium values C89/C++98 portability, deterministic architecture, CLI/TUI support, and clean platform/runtime separation.
\item The strongest unresolved technical task is auditing actual backend timers and persistence formats.
\item The main conclusion is that Dominium does not need a conceptual rewrite for time. It needs boundary hardening. Future work should audit all DSYS backends, freeze ACT units and serialization, ban wall-clock time from authority paths, make civil/astronomical time a derived projection layer, and verify any external platform facts before making them formal documentation.
\end{itemize}

\subsection{UE6\_Domino\_Deterministic\_Universe}

\begin{itemize}
\item Source file: docs/archive/conversations/UE6\_Domino\_Deterministic\_Universe/Dominium\_UE6\_Domino\_Deterministic\_Universe\_\_01\_human\_readable\_report.md
\item UNCERTAIN / UNVERIFIED: This is not a complete reconstruction of every Dominium conversation ever held. The project context shows that many other chats exist about platforms, renderers, setup, game architecture, ecology, ECS, launcher design, and development planning, but their full transcripts are not visible here. Those items are labelled PROJECT-CONTEXT when referenced.
\item Uncertainty: public UE6 technical capabilities and requirements remain time-sensitive and need verification before any future hard commitment.
\item The chat rejected the idea that clients can authoritatively simulate resource production, hidden state, enemy bases, combat, or persistent terrain changes in an MMO. Clients can assist, but the server must verify and commit.
\item REJECTED-02: Waiting for UE6 before beginning serious work. Deprioritised because UE6 availability and capabilities remain uncertain.
\item The user wants direct, source-grounded, audit-ready answers; explicit uncertainty; bounded estimates; and correction when framing is wrong. The preservation prompt explicitly requires human-readable explanation first, structured registers second, uncertainty labels, no invented facts, no hidden chain-of-thought, and downloadable files if tools are available.
\item UNCERTAIN: The final intended role of Domino remains unclear in the visible transcript. It may be a custom engine, target runtime, renderer, or broader platform plan, but the current chat does not fully define it.
\end{itemize}

\subsection{ui\_editor\_tool\_editor\_planning}

\begin{itemize}
\item Source file: docs/archive/conversations/ui\_editor\_tool\_editor\_planning/dominium\_ui\_editor\_tool\_editor\_planning\_\_human\_readable\_chat\_report.txt
\item The most important thing to remember is that this chat was a planning and prompt-generation chat, not proof of implementation. **UNCERTAIN / UNVERIFIED:** No generated prompt is evidence that repository code was actually changed. The next assistant must first verify whether the prompts were executed, inspect the repo/source bundle if available, and avoid treating planned files as existing files.
\item The user uploaded screenshot bundles for setup and launcher. **UNCERTAIN / UNVERIFIED:** The screenshots were not inspected in this chat. The assistant still produced a logical layout description, and the user accepted it as useful. Future work should inspect the bundles before claiming exact screenshot fidelity.
\item What remains uncertain is the actual state of the code. The user named files, but the chat did not inspect the repository. A future assistant must verify exact paths, build targets, TLV parser implementation, and current launcher behavior.
\item Unresolved details include the actual TLV wire format, migration strategy from launcher\_ui\_v1.tlv, and how much legacy data can be preserved during import.
\item The unresolved risk is implementation complexity. Real native tab controls, splitter behavior, and scroll panels can be nontrivial in a child-HWND Win32 UI.
\item This is a strong planned design, but some details remain assistant recommendations rather than explicit user-confirmed decisions. Future implementation should preserve that uncertainty if writing a formal spec.
\item The unresolved issue is that the uploaded screenshot bundles were not inspected. The logical specs are accepted planning artifacts, not verified screenshot extractions.
\item The exact scope of IDE-native live editing remains unresolved. The assistant proposed preview hosts, but it is not confirmed whether that fully satisfies the user's desire to use Visual Studio, Xcode, and Linux GUI tools.
\item The exact first implementation status is unresolved. The chat generated prompts but did not show Codex execution results.
\item UNCERTAIN / INFERENCE:** The assistant generated a preview-host strategy. It was a reasonable way to integrate Visual Studio, Xcode, and Linux tools while preserving DUI TLV as canonical. But the user did not confirm whether this satisfies "utilising all GUI tools available in each IDE."
\item The sixth tradeoff was IDE-native editing. Visual Studio and Xcode have their own native designer ecosystems, but those do not naturally edit DUI TLV. The preview-host proposal preserves DUI as canonical while using IDE build/run/watch loops. This is practical but may not fully satisfy a desire for direct IDE designer manipulation. That is why it remains unresolved.
\item What must happen before it: verify launcher/setup build targets and loader integration points.
\end{itemize}

\subsection{Universe\_Explorer\_Planning}

\begin{itemize}
\item Source file: docs/archive/conversations/Universe\_Explorer\_Planning/Dominium\_Universe\_Explorer\_Planning\_Handoff\_\_01\_human\_readable\_report.md
\item Uncertain: whether all seven universal layers are now formally captured under exactly that naming. The repo has equivalent or related doctrine, but not every early phrase appears as a single artifact.
\item Uncertain: the repo contains pieces of this through affordance/capability/domain docs, but the assistant did not find one single master Deep Primitives spec in the docs it read.
\item The most important unresolved issue is that the Universe Explorer plan is still a proposed/recommended next strategy, not yet a completed repo task. The next best action is either to draft PRESENTATION-CONTRACT-01 / UNIVERSE-EXPLORER-CONTRACT-01, or to preserve this chat and merge it with other old-chat reports into a master Project Spec Book.
\end{itemize}

\subsection{Workbench\_AIDE\_Product\_Spine}

\begin{itemize}
\item Source file: docs/archive/conversations/Workbench\_AIDE\_Product\_Spine/Dominium\_Architecture\_Workbench\_AIDE\_Product\_Spine\_Planning\_\_01\_human\_readable\_report.md
\item The largest uncertainty is not the conceptual architecture; that has converged. The largest uncertainty is live execution state: whether prompts generated near the end of the chat have already been run locally, committed, or pushed. The next chat should first inspect .aide/queue/current.toml, recent commits, and relevant audit files.
\item Remaining uncertainty: exact implementation of erosion/ecology/evolution proxies remains future domain work.
\item 1. Verify current repo state.
\item Explicitly label uncertainty.
\item The next chat should verify repo state first.
\item Preserve FACT / INFERENCE / UNCERTAIN / PROJECT-CONTEXT labels.
\item Verify live repo state before acting.
\item safe\_for\_aggregation: "Yes, with caveats"
\item "UNCERTAIN / UNVERIFIED"
\item "Verify current queue and commits before acting."
\item "Preserve uncertainty and rejected options."
\item uncertain\_or\_not\_established:
\end{itemize}

\subsection{World\_Architecture}

\begin{itemize}
\item Source file: docs/archive/conversations/World\_Architecture/Dominium\_World\_Architecture\_\_human\_readable\_chat\_report.txt
\item The main unresolved issue is implementation detail. The design says Q16.16 and Q4.12, but the earlier prompt used signed types incorrectly. The future implementation must use unsigned local coordinates or a centered signed design.
\item The unresolved parts are the exact meshing algorithm, the fixed-point scaling of ?, and the sample resolution. A 32? sample grid per 16m chunk was recommended but not confirmed as final.
\item The exact solvers remain unresolved. The architecture is clear, but formulas and resolutions are still future work.
\item The save format remains conceptual in places. Exact binary headers, endian policy, compression, and migration strategy are unresolved.
\item The crucial caveat is implementation: unsigned or centered representations are needed. The design intent is final, but the exact typedefs must be fixed. If the implementation uses signed types incorrectly, the whole coordinate system breaks.
\item This is not yet an implementation-level final spec because solver formulas are unresolved.
\item FACT:** Core target is C89. This is non-negotiable in principle, but exact portability mechanics are unresolved because strict C89 lacks standard fixed-width integer types.
\item The user wants direct, detailed, rigorous, critical responses. The user prefers source labels, uncertainty preservation, and avoiding invented facts. The user does not want assistant suggestions silently upgraded to decisions.
\item The most important unresolved issue is the fixed-point implementation detail. The user chose Q16.16 runtime and Q4.12 save-local, but the implementation must choose the exact signed/unsigned representation. The prior prompt's signed version is invalid. This must be resolved before code.
\item The second major unresolved issue is C89 portability. The user wants C89, but the architecture needs 64-bit IDs, RNG states, and accumulators. Strict ISO C89 does not standardize stdint.h or 64-bit integer types. A portability layer or compiler-extension policy is required.
\item The third unresolved issue is exact save encoding. TLV, Region files, content locks, and sparse saves were designed conceptually, but endian policy, padding, alignment, compression, and path grammar still need specification.
\item The fourth unresolved issue is solver math. FluidSpaces, weather, hydrology, oil pressure, ruptures, hydraulics, thermal systems, and electrical networks were designed architecturally, but formulas and fixed-point scales are not final. These do not block v0 core implementation but must be resolved before those systems are coded.
\end{itemize}

\subsection{xstack\_lab\_galaxy}

\begin{itemize}
\item Source file: docs/archive/conversations/xstack\_lab\_galaxy/dominium\_xstack\_lab\_galaxy\_handoff\_\_human\_readable\_chat\_report.txt
\item But this chat did not verify the repository. That became one of the most important unresolved issues.
\item This is one of the most important outputs associated with this chat, but it is also one of the most uncertain because it is user-reported from another chat and not verified directly here. The user explicitly wanted it audited for completion and consistency.
\item INFERENCE: The user also wanted to prevent future assistants from flattening nuance. They repeatedly asked for maximum fidelity, uncertainty labels, rejected options, contradictions, and rationale. This implies a strong preference for preserving complexity rather than oversimplifying.
\item INFERENCE: The user likely wants to resume productive development after verifying the substrate. The next likely feature area is survival or Lab Galaxy refinement, but the user has not made a final decision in this visible chat.
\item 4. From documentation to verifying a new 13-prompt substrate.
\item The biggest unresolved goal is verifying the actual repository. The package says what was reported, but not what is proven. Another unresolved goal is deciding what comes next: survival, Lab Galaxy UX, distributed SRZ, or documentation polish.
\item The main assumption behind all of this is that the repo can actually enforce these contracts. That remains unverified in this chat.
\item The best next step is to verify the actual repository state. This matters because the long 13-prompt transcript is not proof. The user explicitly does not know whether everything was truly implemented.
\item Because the user cares about XStack portability, the new assistant should verify that runtime can build/run without tools/xstack and that engine/game/client/server do not import or depend on it.
\item FACT: The user asked for FACT / INFERENCE / UNCERTAIN labels.
\item This matters because future work depends on whether the 13-prompt changes are committed, pushed, or still local. The transcript reports a clean branch ahead by 10 commits, but this chat did not verify it.
\item Before merging into a master spec, the aggregator should preserve uncertainty and verify the actual repository.
\end{itemize}



{\small\raggedright SOURCE PATH: docs/\allowbreak{}archive/\allowbreak{}conversations/\allowbreak{}\_\allowbreak{}audit/\allowbreak{}DOC\_\allowbreak{}DRIFT\_\allowbreak{}REPORT.md\par}


\section{Document Drift Report}

This report flags archived claims that may drift from current README, queue, language baseline, or authority order.


\subsection{CONTRA-0001 - conversation\_vs\_current\_queue}

\begin{itemize}
\item Source conversation: advanced\_simulation\_infrastructure
\item Source file: docs/archive/conversations/advanced\_simulation\_infrastructure/dominium\_advanced\_simulation\_infrastructure\_architecture\_\_human\_readable\_chat\_report.txt
\item Claim: Archived conversation discusses work related to blocked scope gameplay.
\item Conflict target: .aide/queue/current.toml
\item Recommended disposition: preserve\_as\_history; require current queue authority before action
\end{itemize}

\subsection{CONTRA-0002 - stale\_external\_claim}

\begin{itemize}
\item Source conversation: advanced\_simulation\_infrastructure
\item Source file: docs/archive/conversations/advanced\_simulation\_infrastructure/dominium\_advanced\_simulation\_infrastructure\_architecture\_\_human\_readable\_chat\_report.txt
\item Claim: Archived conversation contains potentially stale external or baseline claim: C89.
\item Conflict target: README.md / current platform and language baseline
\item Recommended disposition: quarantined until current repo and external facts are checked
\end{itemize}

\subsection{CONTRA-0003 - stale\_external\_claim}

\begin{itemize}
\item Source conversation: advanced\_simulation\_infrastructure
\item Source file: docs/archive/conversations/advanced\_simulation\_infrastructure/dominium\_advanced\_simulation\_infrastructure\_architecture\_\_human\_readable\_chat\_report.txt
\item Claim: Archived conversation contains potentially stale external or baseline claim: C++98.
\item Conflict target: README.md / current platform and language baseline
\item Recommended disposition: quarantined until current repo and external facts are checked
\end{itemize}

\subsection{CONTRA-0004 - stale\_external\_claim}

\begin{itemize}
\item Source conversation: advanced\_simulation\_infrastructure
\item Source file: docs/archive/conversations/advanced\_simulation\_infrastructure/dominium\_advanced\_simulation\_infrastructure\_architecture\_\_human\_readable\_chat\_report.txt
\item Claim: Archived conversation contains potentially stale external or baseline claim: ISO C89.
\item Conflict target: README.md / current platform and language baseline
\item Recommended disposition: quarantined until current repo and external facts are checked
\end{itemize}

\subsection{CONTRA-0005 - conversation\_vs\_current\_queue}

\begin{itemize}
\item Source conversation: app\_runtime\_platform\_renderers
\item Source file: docs/archive/conversations/app\_runtime\_platform\_renderers/dominium\_app0\_runtime\_platform\_renderers\_\_human\_readable\_chat\_report.txt
\item Claim: Archived conversation discusses work related to blocked scope runtime\_module\_loader.
\item Conflict target: .aide/queue/current.toml
\item Recommended disposition: preserve\_as\_history; require current queue authority before action
\end{itemize}

\subsection{CONTRA-0006 - conversation\_vs\_current\_queue}

\begin{itemize}
\item Source conversation: app\_runtime\_platform\_renderers
\item Source file: docs/archive/conversations/app\_runtime\_platform\_renderers/dominium\_app0\_runtime\_platform\_renderers\_\_human\_readable\_chat\_report.txt
\item Claim: Archived conversation discusses work related to blocked scope provider\_runtime.
\item Conflict target: .aide/queue/current.toml
\item Recommended disposition: preserve\_as\_history; require current queue authority before action
\end{itemize}

\subsection{CONTRA-0007 - conversation\_vs\_current\_queue}

\begin{itemize}
\item Source conversation: app\_runtime\_platform\_renderers
\item Source file: docs/archive/conversations/app\_runtime\_platform\_renderers/dominium\_app0\_runtime\_platform\_renderers\_\_human\_readable\_chat\_report.txt
\item Claim: Archived conversation discusses work related to blocked scope gameplay.
\item Conflict target: .aide/queue/current.toml
\item Recommended disposition: preserve\_as\_history; require current queue authority before action
\end{itemize}

\subsection{CONTRA-0008 - conversation\_vs\_current\_queue}

\begin{itemize}
\item Source conversation: app\_runtime\_platform\_renderers
\item Source file: docs/archive/conversations/app\_runtime\_platform\_renderers/dominium\_app0\_runtime\_platform\_renderers\_\_human\_readable\_chat\_report.txt
\item Claim: Archived conversation discusses work related to blocked scope renderer\_implementation.
\item Conflict target: .aide/queue/current.toml
\item Recommended disposition: preserve\_as\_history; require current queue authority before action
\end{itemize}

\subsection{CONTRA-0009 - stale\_external\_claim}

\begin{itemize}
\item Source conversation: app\_runtime\_platform\_renderers
\item Source file: docs/archive/conversations/app\_runtime\_platform\_renderers/dominium\_app0\_runtime\_platform\_renderers\_\_human\_readable\_chat\_report.txt
\item Claim: Archived conversation contains potentially stale external or baseline claim: SDK.
\item Conflict target: README.md / current platform and language baseline
\item Recommended disposition: quarantined until current repo and external facts are checked
\end{itemize}

\subsection{CONTRA-0010 - conversation\_vs\_current\_queue}

\begin{itemize}
\item Source conversation: app\_testx\_codehygiene
\item Source file: docs/archive/conversations/app\_testx\_codehygiene/dominium\_app\_testx\_codehygiene\_handoff\_\_human\_readable\_chat\_report.txt
\item Claim: Archived conversation discusses work related to blocked scope gameplay.
\item Conflict target: .aide/queue/current.toml
\item Recommended disposition: preserve\_as\_history; require current queue authority before action
\end{itemize}

\subsection{CONTRA-0011 - conversation\_vs\_current\_queue}

\begin{itemize}
\item Source conversation: app\_testx\_codehygiene
\item Source file: docs/archive/conversations/app\_testx\_codehygiene/dominium\_app\_testx\_codehygiene\_handoff\_\_human\_readable\_chat\_report.txt
\item Claim: Archived conversation discusses work related to blocked scope renderer\_implementation.
\item Conflict target: .aide/queue/current.toml
\item Recommended disposition: preserve\_as\_history; require current queue authority before action
\end{itemize}

\subsection{CONTRA-0012 - stale\_external\_claim}

\begin{itemize}
\item Source conversation: app\_testx\_codehygiene
\item Source file: docs/archive/conversations/app\_testx\_codehygiene/dominium\_app\_testx\_codehygiene\_handoff\_\_human\_readable\_chat\_report.txt
\item Claim: Archived conversation contains potentially stale external or baseline claim: C89.
\item Conflict target: README.md / current platform and language baseline
\item Recommended disposition: quarantined until current repo and external facts are checked
\end{itemize}

\subsection{CONTRA-0013 - stale\_external\_claim}

\begin{itemize}
\item Source conversation: app\_testx\_codehygiene
\item Source file: docs/archive/conversations/app\_testx\_codehygiene/dominium\_app\_testx\_codehygiene\_handoff\_\_human\_readable\_chat\_report.txt
\item Claim: Archived conversation contains potentially stale external or baseline claim: C++98.
\item Conflict target: README.md / current platform and language baseline
\item Recommended disposition: quarantined until current repo and external facts are checked
\end{itemize}

\subsection{CONTRA-0014 - stale\_external\_claim}

\begin{itemize}
\item Source conversation: app\_testx\_codehygiene
\item Source file: docs/archive/conversations/app\_testx\_codehygiene/dominium\_app\_testx\_codehygiene\_handoff\_\_human\_readable\_chat\_report.txt
\item Claim: Archived conversation contains potentially stale external or baseline claim: DX12.
\item Conflict target: README.md / current platform and language baseline
\item Recommended disposition: quarantined until current repo and external facts are checked
\end{itemize}

\subsection{CONTRA-0015 - stale\_external\_claim}

\begin{itemize}
\item Source conversation: app\_testx\_codehygiene
\item Source file: docs/archive/conversations/app\_testx\_codehygiene/dominium\_app\_testx\_codehygiene\_handoff\_\_human\_readable\_chat\_report.txt
\item Claim: Archived conversation contains potentially stale external or baseline claim: iOS.
\item Conflict target: README.md / current platform and language baseline
\item Recommended disposition: quarantined until current repo and external facts are checked
\end{itemize}

\subsection{CONTRA-0016 - conversation\_vs\_current\_queue}

\begin{itemize}
\item Source conversation: architecture\_codex\_prompts
\item Source file: docs/archive/conversations/architecture\_codex\_prompts/dominium\_domino\_architecture\_codex\_prompts\_\_human\_readable\_chat\_report.txt
\item Claim: Archived conversation discusses work related to blocked scope provider\_runtime.
\item Conflict target: .aide/queue/current.toml
\item Recommended disposition: preserve\_as\_history; require current queue authority before action
\end{itemize}

\subsection{CONTRA-0017 - conversation\_vs\_current\_queue}

\begin{itemize}
\item Source conversation: architecture\_codex\_prompts
\item Source file: docs/archive/conversations/architecture\_codex\_prompts/dominium\_domino\_architecture\_codex\_prompts\_\_human\_readable\_chat\_report.txt
\item Claim: Archived conversation discusses work related to blocked scope gameplay.
\item Conflict target: .aide/queue/current.toml
\item Recommended disposition: preserve\_as\_history; require current queue authority before action
\end{itemize}

\subsection{CONTRA-0018 - stale\_external\_claim}

\begin{itemize}
\item Source conversation: architecture\_codex\_prompts
\item Source file: docs/archive/conversations/architecture\_codex\_prompts/dominium\_domino\_architecture\_codex\_prompts\_\_human\_readable\_chat\_report.txt
\item Claim: Archived conversation contains potentially stale external or baseline claim: C89.
\item Conflict target: README.md / current platform and language baseline
\item Recommended disposition: quarantined until current repo and external facts are checked
\end{itemize}

\subsection{CONTRA-0019 - stale\_external\_claim}

\begin{itemize}
\item Source conversation: architecture\_codex\_prompts
\item Source file: docs/archive/conversations/architecture\_codex\_prompts/dominium\_domino\_architecture\_codex\_prompts\_\_human\_readable\_chat\_report.txt
\item Claim: Archived conversation contains potentially stale external or baseline claim: C++98.
\item Conflict target: README.md / current platform and language baseline
\item Recommended disposition: quarantined until current repo and external facts are checked
\end{itemize}

\subsection{CONTRA-0020 - stale\_external\_claim}

\begin{itemize}
\item Source conversation: architecture\_codex\_prompts
\item Source file: docs/archive/conversations/architecture\_codex\_prompts/dominium\_domino\_architecture\_codex\_prompts\_\_human\_readable\_chat\_report.txt
\item Claim: Archived conversation contains potentially stale external or baseline claim: ISO C89.
\item Conflict target: README.md / current platform and language baseline
\item Recommended disposition: quarantined until current repo and external facts are checked
\end{itemize}

\subsection{CONTRA-0021 - conversation\_vs\_current\_queue}

\begin{itemize}
\item Source conversation: architecture\_ui\_providers
\item Source file: docs/archive/conversations/architecture\_ui\_providers/dominium\_architecture\_ui\_providers\_robot\_os\_strategy\_\_01\_human\_readable\_report.md
\item Claim: Archived conversation discusses work related to blocked scope broad\_workbench\_ui.
\item Conflict target: .aide/queue/current.toml
\item Recommended disposition: preserve\_as\_history; require current queue authority before action
\end{itemize}

\subsection{CONTRA-0022 - conversation\_vs\_current\_queue}

\begin{itemize}
\item Source conversation: architecture\_ui\_providers
\item Source file: docs/archive/conversations/architecture\_ui\_providers/dominium\_architecture\_ui\_providers\_robot\_os\_strategy\_\_01\_human\_readable\_report.md
\item Claim: Archived conversation discusses work related to blocked scope provider\_runtime.
\item Conflict target: .aide/queue/current.toml
\item Recommended disposition: preserve\_as\_history; require current queue authority before action
\end{itemize}

\subsection{CONTRA-0023 - conversation\_vs\_current\_queue}

\begin{itemize}
\item Source conversation: architecture\_ui\_providers
\item Source file: docs/archive/conversations/architecture\_ui\_providers/dominium\_architecture\_ui\_providers\_robot\_os\_strategy\_\_01\_human\_readable\_report.md
\item Claim: Archived conversation discusses work related to blocked scope gameplay.
\item Conflict target: .aide/queue/current.toml
\item Recommended disposition: preserve\_as\_history; require current queue authority before action
\end{itemize}

\subsection{CONTRA-0024 - conversation\_vs\_current\_queue}

\begin{itemize}
\item Source conversation: architecture\_ui\_providers
\item Source file: docs/archive/conversations/architecture\_ui\_providers/dominium\_architecture\_ui\_providers\_robot\_os\_strategy\_\_01\_human\_readable\_report.md
\item Claim: Archived conversation discusses work related to blocked scope renderer\_implementation.
\item Conflict target: .aide/queue/current.toml
\item Recommended disposition: preserve\_as\_history; require current queue authority before action
\end{itemize}

\subsection{CONTRA-0025 - conversation\_vs\_current\_queue}

\begin{itemize}
\item Source conversation: architecture\_ui\_providers
\item Source file: docs/archive/conversations/architecture\_ui\_providers/dominium\_architecture\_ui\_providers\_robot\_os\_strategy\_\_01\_human\_readable\_report.md
\item Claim: Archived conversation discusses work related to blocked scope native\_gui.
\item Conflict target: .aide/queue/current.toml
\item Recommended disposition: preserve\_as\_history; require current queue authority before action
\end{itemize}

\subsection{CONTRA-0026 - stale\_external\_claim}

\begin{itemize}
\item Source conversation: architecture\_ui\_providers
\item Source file: docs/archive/conversations/architecture\_ui\_providers/dominium\_architecture\_ui\_providers\_robot\_os\_strategy\_\_01\_human\_readable\_report.md
\item Claim: Archived conversation contains potentially stale external or baseline claim: C89.
\item Conflict target: README.md / current platform and language baseline
\item Recommended disposition: quarantined until current repo and external facts are checked
\end{itemize}

\subsection{CONTRA-0027 - stale\_external\_claim}

\begin{itemize}
\item Source conversation: architecture\_ui\_providers
\item Source file: docs/archive/conversations/architecture\_ui\_providers/dominium\_architecture\_ui\_providers\_robot\_os\_strategy\_\_01\_human\_readable\_report.md
\item Claim: Archived conversation contains potentially stale external or baseline claim: C++98.
\item Conflict target: README.md / current platform and language baseline
\item Recommended disposition: quarantined until current repo and external facts are checked
\end{itemize}

\subsection{CONTRA-0028 - stale\_external\_claim}

\begin{itemize}
\item Source conversation: architecture\_ui\_providers
\item Source file: docs/archive/conversations/architecture\_ui\_providers/dominium\_architecture\_ui\_providers\_robot\_os\_strategy\_\_01\_human\_readable\_report.md
\item Claim: Archived conversation contains potentially stale external or baseline claim: DX12.
\item Conflict target: README.md / current platform and language baseline
\item Recommended disposition: quarantined until current repo and external facts are checked
\end{itemize}

\subsection{CONTRA-0029 - stale\_external\_claim}

\begin{itemize}
\item Source conversation: architecture\_ui\_providers
\item Source file: docs/archive/conversations/architecture\_ui\_providers/dominium\_architecture\_ui\_providers\_robot\_os\_strategy\_\_01\_human\_readable\_report.md
\item Claim: Archived conversation contains potentially stale external or baseline claim: SDK.
\item Conflict target: README.md / current platform and language baseline
\item Recommended disposition: quarantined until current repo and external facts are checked
\end{itemize}

\subsection{CONTRA-0030 - stale\_external\_claim}

\begin{itemize}
\item Source conversation: Build\_and\_Future\_Proofing
\item Source file: docs/archive/conversations/Build\_and\_Future\_Proofing/Dominium\_Build\_and\_Future\_Proofing\_Architecture\_\_01\_human\_readable\_report.md
\item Claim: Archived conversation contains potentially stale external or baseline claim: C89.
\item Conflict target: README.md / current platform and language baseline
\item Recommended disposition: quarantined until current repo and external facts are checked
\end{itemize}

\subsection{CONTRA-0031 - stale\_external\_claim}

\begin{itemize}
\item Source conversation: Build\_and\_Future\_Proofing
\item Source file: docs/archive/conversations/Build\_and\_Future\_Proofing/Dominium\_Build\_and\_Future\_Proofing\_Architecture\_\_01\_human\_readable\_report.md
\item Claim: Archived conversation contains potentially stale external or baseline claim: C++98.
\item Conflict target: README.md / current platform and language baseline
\item Recommended disposition: quarantined until current repo and external facts are checked
\end{itemize}

\subsection{CONTRA-0032 - stale\_external\_claim}

\begin{itemize}
\item Source conversation: Build\_and\_Future\_Proofing
\item Source file: docs/archive/conversations/Build\_and\_Future\_Proofing/Dominium\_Build\_and\_Future\_Proofing\_Architecture\_\_01\_human\_readable\_report.md
\item Claim: Archived conversation contains potentially stale external or baseline claim: SDK.
\item Conflict target: README.md / current platform and language baseline
\item Recommended disposition: quarantined until current repo and external facts are checked
\end{itemize}

\subsection{CONTRA-0033 - conversation\_vs\_current\_queue}

\begin{itemize}
\item Source conversation: canonical\_structure\_and\_framework
\item Source file: docs/archive/conversations/canonical\_structure\_and\_framework/dominium\_canonical\_architecture\_repo\_foundation\_\_01\_human\_readable\_report.md
\item Claim: Archived conversation discusses work related to blocked scope provider\_runtime.
\item Conflict target: .aide/queue/current.toml
\item Recommended disposition: preserve\_as\_history; require current queue authority before action
\end{itemize}

\subsection{CONTRA-0034 - conversation\_vs\_current\_queue}

\begin{itemize}
\item Source conversation: canonical\_structure\_and\_framework
\item Source file: docs/archive/conversations/canonical\_structure\_and\_framework/dominium\_canonical\_architecture\_repo\_foundation\_\_01\_human\_readable\_report.md
\item Claim: Archived conversation discusses work related to blocked scope package\_runtime.
\item Conflict target: .aide/queue/current.toml
\item Recommended disposition: preserve\_as\_history; require current queue authority before action
\end{itemize}

\subsection{CONTRA-0035 - conversation\_vs\_current\_queue}

\begin{itemize}
\item Source conversation: canonical\_structure\_and\_framework
\item Source file: docs/archive/conversations/canonical\_structure\_and\_framework/dominium\_canonical\_architecture\_repo\_foundation\_\_01\_human\_readable\_report.md
\item Claim: Archived conversation discusses work related to blocked scope gameplay.
\item Conflict target: .aide/queue/current.toml
\item Recommended disposition: preserve\_as\_history; require current queue authority before action
\end{itemize}

\subsection{CONTRA-0036 - conversation\_vs\_current\_queue}

\begin{itemize}
\item Source conversation: canonical\_structure\_and\_framework
\item Source file: docs/archive/conversations/canonical\_structure\_and\_framework/dominium\_canonical\_architecture\_repo\_foundation\_\_01\_human\_readable\_report.md
\item Claim: Archived conversation discusses work related to blocked scope renderer\_implementation.
\item Conflict target: .aide/queue/current.toml
\item Recommended disposition: preserve\_as\_history; require current queue authority before action
\end{itemize}

\subsection{CONTRA-0037 - conversation\_vs\_current\_queue}

\begin{itemize}
\item Source conversation: canonical\_structure\_and\_framework
\item Source file: docs/archive/conversations/canonical\_structure\_and\_framework/dominium\_canonical\_architecture\_repo\_foundation\_\_01\_human\_readable\_report.md
\item Claim: Archived conversation discusses work related to blocked scope native\_gui.
\item Conflict target: .aide/queue/current.toml
\item Recommended disposition: preserve\_as\_history; require current queue authority before action
\end{itemize}

\subsection{CONTRA-0038 - stale\_external\_claim}

\begin{itemize}
\item Source conversation: canonical\_structure\_and\_framework
\item Source file: docs/archive/conversations/canonical\_structure\_and\_framework/dominium\_canonical\_architecture\_repo\_foundation\_\_01\_human\_readable\_report.md
\item Claim: Archived conversation contains potentially stale external or baseline claim: C89.
\item Conflict target: README.md / current platform and language baseline
\item Recommended disposition: quarantined until current repo and external facts are checked
\end{itemize}

\subsection{CONTRA-0039 - stale\_external\_claim}

\begin{itemize}
\item Source conversation: canonical\_structure\_and\_framework
\item Source file: docs/archive/conversations/canonical\_structure\_and\_framework/dominium\_canonical\_architecture\_repo\_foundation\_\_01\_human\_readable\_report.md
\item Claim: Archived conversation contains potentially stale external or baseline claim: C++98.
\item Conflict target: README.md / current platform and language baseline
\item Recommended disposition: quarantined until current repo and external facts are checked
\end{itemize}

\subsection{CONTRA-0040 - stale\_external\_claim}

\begin{itemize}
\item Source conversation: canonical\_structure\_and\_framework
\item Source file: docs/archive/conversations/canonical\_structure\_and\_framework/dominium\_canonical\_architecture\_repo\_foundation\_\_01\_human\_readable\_report.md
\item Claim: Archived conversation contains potentially stale external or baseline claim: SDK.
\item Conflict target: README.md / current platform and language baseline
\item Recommended disposition: quarantined until current repo and external facts are checked
\end{itemize}

\subsection{CONTRA-0041 - conversation\_vs\_current\_queue}

\begin{itemize}
\item Source conversation: Chronology\_Celestial\_Systems
\item Source file: docs/archive/conversations/Chronology\_Celestial\_Systems/Dominium\_Chronology\_Celestial\_Systems\_\_human\_readable\_chat\_report.txt
\item Claim: Archived conversation discusses work related to blocked scope gameplay.
\item Conflict target: .aide/queue/current.toml
\item Recommended disposition: preserve\_as\_history; require current queue authority before action
\end{itemize}

\subsection{CONTRA-0042 - conversation\_vs\_current\_queue}

\begin{itemize}
\item Source conversation: development\_routes
\item Source file: docs/archive/conversations/development\_routes/dominium\_development\_routes\_\_human\_readable\_chat\_report.txt
\item Claim: Archived conversation discusses work related to blocked scope gameplay.
\item Conflict target: .aide/queue/current.toml
\item Recommended disposition: preserve\_as\_history; require current queue authority before action
\end{itemize}

\subsection{CONTRA-0043 - stale\_external\_claim}

\begin{itemize}
\item Source conversation: development\_routes
\item Source file: docs/archive/conversations/development\_routes/dominium\_development\_routes\_\_human\_readable\_chat\_report.txt
\item Claim: Archived conversation contains potentially stale external or baseline claim: iOS.
\item Conflict target: README.md / current platform and language baseline
\item Recommended disposition: quarantined until current repo and external facts are checked
\end{itemize}

\subsection{CONTRA-0044 - stale\_external\_claim}

\begin{itemize}
\item Source conversation: documentation\_standards\_readme
\item Source file: docs/archive/conversations/documentation\_standards\_readme/documentation\_standards\_readme\_handoff\_\_human\_readable\_chat\_report.txt
\item Claim: Archived conversation contains potentially stale external or baseline claim: C89.
\item Conflict target: README.md / current platform and language baseline
\item Recommended disposition: quarantined until current repo and external facts are checked
\end{itemize}

\subsection{CONTRA-0045 - stale\_external\_claim}

\begin{itemize}
\item Source conversation: documentation\_standards\_readme
\item Source file: docs/archive/conversations/documentation\_standards\_readme/documentation\_standards\_readme\_handoff\_\_human\_readable\_chat\_report.txt
\item Claim: Archived conversation contains potentially stale external or baseline claim: C++98.
\item Conflict target: README.md / current platform and language baseline
\item Recommended disposition: quarantined until current repo and external facts are checked
\end{itemize}

\subsection{CONTRA-0046 - stale\_external\_claim}

\begin{itemize}
\item Source conversation: documentation\_standards\_readme
\item Source file: docs/archive/conversations/documentation\_standards\_readme/documentation\_standards\_readme\_handoff\_\_human\_readable\_chat\_report.txt
\item Claim: Archived conversation contains potentially stale external or baseline claim: ISO C89.
\item Conflict target: README.md / current platform and language baseline
\item Recommended disposition: quarantined until current repo and external facts are checked
\end{itemize}

\subsection{CONTRA-0047 - stale\_external\_claim}

\begin{itemize}
\item Source conversation: documentation\_standards\_readme
\item Source file: docs/archive/conversations/documentation\_standards\_readme/documentation\_standards\_readme\_handoff\_\_human\_readable\_chat\_report.txt
\item Claim: Archived conversation contains potentially stale external or baseline claim: iOS.
\item Conflict target: README.md / current platform and language baseline
\item Recommended disposition: quarantined until current repo and external facts are checked
\end{itemize}

\subsection{CONTRA-0048 - conversation\_vs\_current\_queue}

\begin{itemize}
\item Source conversation: Dominium\_Architecture\_I
\item Source file: docs/archive/conversations/Dominium\_Architecture\_I/Dominium\_Architecture\_I\_\_human\_readable\_chat\_report.txt
\item Claim: Archived conversation discusses work related to blocked scope gameplay.
\item Conflict target: .aide/queue/current.toml
\item Recommended disposition: preserve\_as\_history; require current queue authority before action
\end{itemize}

\subsection{CONTRA-0049 - conversation\_vs\_current\_queue}

\begin{itemize}
\item Source conversation: Dominium\_Architecture\_I
\item Source file: docs/archive/conversations/Dominium\_Architecture\_I/Dominium\_Architecture\_I\_\_human\_readable\_chat\_report.txt
\item Claim: Archived conversation discusses work related to blocked scope renderer\_implementation.
\item Conflict target: .aide/queue/current.toml
\item Recommended disposition: preserve\_as\_history; require current queue authority before action
\end{itemize}

\subsection{CONTRA-0050 - stale\_external\_claim}

\begin{itemize}
\item Source conversation: Dominium\_Architecture\_I
\item Source file: docs/archive/conversations/Dominium\_Architecture\_I/Dominium\_Architecture\_I\_\_human\_readable\_chat\_report.txt
\item Claim: Archived conversation contains potentially stale external or baseline claim: C89.
\item Conflict target: README.md / current platform and language baseline
\item Recommended disposition: quarantined until current repo and external facts are checked
\end{itemize}

\subsection{CONTRA-0051 - stale\_external\_claim}

\begin{itemize}
\item Source conversation: Dominium\_Architecture\_I
\item Source file: docs/archive/conversations/Dominium\_Architecture\_I/Dominium\_Architecture\_I\_\_human\_readable\_chat\_report.txt
\item Claim: Archived conversation contains potentially stale external or baseline claim: C++98.
\item Conflict target: README.md / current platform and language baseline
\item Recommended disposition: quarantined until current repo and external facts are checked
\end{itemize}

\subsection{CONTRA-0052 - conversation\_vs\_current\_queue}

\begin{itemize}
\item Source conversation: Dominium\_Architecture\_II
\item Source file: docs/archive/conversations/Dominium\_Architecture\_II/Dominium\_Architecture\_II\_\_human\_readable\_chat\_report.txt
\item Claim: Archived conversation discusses work related to blocked scope gameplay.
\item Conflict target: .aide/queue/current.toml
\item Recommended disposition: preserve\_as\_history; require current queue authority before action
\end{itemize}

\subsection{CONTRA-0053 - conversation\_vs\_current\_queue}

\begin{itemize}
\item Source conversation: Dominium\_Architecture\_II
\item Source file: docs/archive/conversations/Dominium\_Architecture\_II/Dominium\_Architecture\_II\_\_human\_readable\_chat\_report.txt
\item Claim: Archived conversation discusses work related to blocked scope renderer\_implementation.
\item Conflict target: .aide/queue/current.toml
\item Recommended disposition: preserve\_as\_history; require current queue authority before action
\end{itemize}

\subsection{CONTRA-0054 - stale\_external\_claim}

\begin{itemize}
\item Source conversation: Dominium\_Architecture\_II
\item Source file: docs/archive/conversations/Dominium\_Architecture\_II/Dominium\_Architecture\_II\_\_human\_readable\_chat\_report.txt
\item Claim: Archived conversation contains potentially stale external or baseline claim: C89.
\item Conflict target: README.md / current platform and language baseline
\item Recommended disposition: quarantined until current repo and external facts are checked
\end{itemize}

\subsection{CONTRA-0055 - stale\_external\_claim}

\begin{itemize}
\item Source conversation: Dominium\_Architecture\_II
\item Source file: docs/archive/conversations/Dominium\_Architecture\_II/Dominium\_Architecture\_II\_\_human\_readable\_chat\_report.txt
\item Claim: Archived conversation contains potentially stale external or baseline claim: Windows 98.
\item Conflict target: README.md / current platform and language baseline
\item Recommended disposition: quarantined until current repo and external facts are checked
\end{itemize}

\subsection{CONTRA-0056 - stale\_external\_claim}

\begin{itemize}
\item Source conversation: Dominium\_Architecture\_II
\item Source file: docs/archive/conversations/Dominium\_Architecture\_II/Dominium\_Architecture\_II\_\_human\_readable\_chat\_report.txt
\item Claim: Archived conversation contains potentially stale external or baseline claim: Windows NT 2000.
\item Conflict target: README.md / current platform and language baseline
\item Recommended disposition: quarantined until current repo and external facts are checked
\end{itemize}

\subsection{CONTRA-0057 - stale\_external\_claim}

\begin{itemize}
\item Source conversation: Dominium\_Architecture\_II
\item Source file: docs/archive/conversations/Dominium\_Architecture\_II/Dominium\_Architecture\_II\_\_human\_readable\_chat\_report.txt
\item Claim: Archived conversation contains potentially stale external or baseline claim: DX12.
\item Conflict target: README.md / current platform and language baseline
\item Recommended disposition: quarantined until current repo and external facts are checked
\end{itemize}

\subsection{CONTRA-0058 - stale\_external\_claim}

\begin{itemize}
\item Source conversation: Dominium\_Architecture\_II
\item Source file: docs/archive/conversations/Dominium\_Architecture\_II/Dominium\_Architecture\_II\_\_human\_readable\_chat\_report.txt
\item Claim: Archived conversation contains potentially stale external or baseline claim: Android.
\item Conflict target: README.md / current platform and language baseline
\item Recommended disposition: quarantined until current repo and external facts are checked
\end{itemize}

\subsection{CONTRA-0059 - conversation\_vs\_current\_queue}

\begin{itemize}
\item Source conversation: Dominium\_Architecture\_III
\item Source file: docs/archive/conversations/Dominium\_Architecture\_III/Dominium\_Architecture\_III\_\_human\_readable\_chat\_report.txt
\item Claim: Archived conversation discusses work related to blocked scope renderer\_implementation.
\item Conflict target: .aide/queue/current.toml
\item Recommended disposition: preserve\_as\_history; require current queue authority before action
\end{itemize}

\subsection{CONTRA-0060 - conversation\_vs\_current\_queue}

\begin{itemize}
\item Source conversation: Dominium\_Architecture\_III
\item Source file: docs/archive/conversations/Dominium\_Architecture\_III/Dominium\_Architecture\_III\_\_human\_readable\_chat\_report.txt
\item Claim: Archived conversation discusses work related to blocked scope release\_publication.
\item Conflict target: .aide/queue/current.toml
\item Recommended disposition: preserve\_as\_history; require current queue authority before action
\end{itemize}

\subsection{CONTRA-0061 - stale\_external\_claim}

\begin{itemize}
\item Source conversation: Dominium\_Architecture\_III
\item Source file: docs/archive/conversations/Dominium\_Architecture\_III/Dominium\_Architecture\_III\_\_human\_readable\_chat\_report.txt
\item Claim: Archived conversation contains potentially stale external or baseline claim: C89.
\item Conflict target: README.md / current platform and language baseline
\item Recommended disposition: quarantined until current repo and external facts are checked
\end{itemize}

\subsection{CONTRA-0062 - stale\_external\_claim}

\begin{itemize}
\item Source conversation: Dominium\_Architecture\_III
\item Source file: docs/archive/conversations/Dominium\_Architecture\_III/Dominium\_Architecture\_III\_\_human\_readable\_chat\_report.txt
\item Claim: Archived conversation contains potentially stale external or baseline claim: C++98.
\item Conflict target: README.md / current platform and language baseline
\item Recommended disposition: quarantined until current repo and external facts are checked
\end{itemize}

\subsection{CONTRA-0063 - stale\_external\_claim}

\begin{itemize}
\item Source conversation: Dominium\_Architecture\_III
\item Source file: docs/archive/conversations/Dominium\_Architecture\_III/Dominium\_Architecture\_III\_\_human\_readable\_chat\_report.txt
\item Claim: Archived conversation contains potentially stale external or baseline claim: Win98.
\item Conflict target: README.md / current platform and language baseline
\item Recommended disposition: quarantined until current repo and external facts are checked
\end{itemize}

\subsection{CONTRA-0064 - stale\_external\_claim}

\begin{itemize}
\item Source conversation: Dominium\_Architecture\_III
\item Source file: docs/archive/conversations/Dominium\_Architecture\_III/Dominium\_Architecture\_III\_\_human\_readable\_chat\_report.txt
\item Claim: Archived conversation contains potentially stale external or baseline claim: Windows 98.
\item Conflict target: README.md / current platform and language baseline
\item Recommended disposition: quarantined until current repo and external facts are checked
\end{itemize}

\subsection{CONTRA-0065 - stale\_external\_claim}

\begin{itemize}
\item Source conversation: Dominium\_Architecture\_III
\item Source file: docs/archive/conversations/Dominium\_Architecture\_III/Dominium\_Architecture\_III\_\_human\_readable\_chat\_report.txt
\item Claim: Archived conversation contains potentially stale external or baseline claim: Windows NT 2000.
\item Conflict target: README.md / current platform and language baseline
\item Recommended disposition: quarantined until current repo and external facts are checked
\end{itemize}

\subsection{CONTRA-0066 - stale\_external\_claim}

\begin{itemize}
\item Source conversation: Dominium\_Architecture\_III
\item Source file: docs/archive/conversations/Dominium\_Architecture\_III/Dominium\_Architecture\_III\_\_human\_readable\_chat\_report.txt
\item Claim: Archived conversation contains potentially stale external or baseline claim: Mac OS 9.
\item Conflict target: README.md / current platform and language baseline
\item Recommended disposition: quarantined until current repo and external facts are checked
\end{itemize}

\subsection{CONTRA-0067 - stale\_external\_claim}

\begin{itemize}
\item Source conversation: Dominium\_Architecture\_III
\item Source file: docs/archive/conversations/Dominium\_Architecture\_III/Dominium\_Architecture\_III\_\_human\_readable\_chat\_report.txt
\item Claim: Archived conversation contains potentially stale external or baseline claim: Mac OS 8.
\item Conflict target: README.md / current platform and language baseline
\item Recommended disposition: quarantined until current repo and external facts are checked
\end{itemize}

\subsection{CONTRA-0068 - stale\_external\_claim}

\begin{itemize}
\item Source conversation: Dominium\_Architecture\_III
\item Source file: docs/archive/conversations/Dominium\_Architecture\_III/Dominium\_Architecture\_III\_\_human\_readable\_chat\_report.txt
\item Claim: Archived conversation contains potentially stale external or baseline claim: DX12.
\item Conflict target: README.md / current platform and language baseline
\item Recommended disposition: quarantined until current repo and external facts are checked
\end{itemize}

\subsection{CONTRA-0069 - conversation\_vs\_current\_queue}

\begin{itemize}
\item Source conversation: Dominium\_Architecture\_IV
\item Source file: docs/archive/conversations/Dominium\_Architecture\_IV/Dominium\_Architecture\_IV\_\_human\_readable\_chat\_report.txt
\item Claim: Archived conversation discusses work related to blocked scope gameplay.
\item Conflict target: .aide/queue/current.toml
\item Recommended disposition: preserve\_as\_history; require current queue authority before action
\end{itemize}

\subsection{CONTRA-0070 - conversation\_vs\_current\_queue}

\begin{itemize}
\item Source conversation: Dominium\_Architecture\_IV
\item Source file: docs/archive/conversations/Dominium\_Architecture\_IV/Dominium\_Architecture\_IV\_\_human\_readable\_chat\_report.txt
\item Claim: Archived conversation discusses work related to blocked scope renderer\_implementation.
\item Conflict target: .aide/queue/current.toml
\item Recommended disposition: preserve\_as\_history; require current queue authority before action
\end{itemize}

\subsection{CONTRA-0071 - conversation\_vs\_current\_queue}

\begin{itemize}
\item Source conversation: Dominium\_Architecture\_IV
\item Source file: docs/archive/conversations/Dominium\_Architecture\_IV/Dominium\_Architecture\_IV\_\_human\_readable\_chat\_report.txt
\item Claim: Archived conversation discusses work related to blocked scope native\_gui.
\item Conflict target: .aide/queue/current.toml
\item Recommended disposition: preserve\_as\_history; require current queue authority before action
\end{itemize}

\subsection{CONTRA-0072 - stale\_external\_claim}

\begin{itemize}
\item Source conversation: Dominium\_Architecture\_IV
\item Source file: docs/archive/conversations/Dominium\_Architecture\_IV/Dominium\_Architecture\_IV\_\_human\_readable\_chat\_report.txt
\item Claim: Archived conversation contains potentially stale external or baseline claim: CP/M.
\item Conflict target: README.md / current platform and language baseline
\item Recommended disposition: quarantined until current repo and external facts are checked
\end{itemize}

\subsection{CONTRA-0073 - stale\_external\_claim}

\begin{itemize}
\item Source conversation: Dominium\_Architecture\_IV
\item Source file: docs/archive/conversations/Dominium\_Architecture\_IV/Dominium\_Architecture\_IV\_\_human\_readable\_chat\_report.txt
\item Claim: Archived conversation contains potentially stale external or baseline claim: Android.
\item Conflict target: README.md / current platform and language baseline
\item Recommended disposition: quarantined until current repo and external facts are checked
\end{itemize}

\subsection{CONTRA-0074 - stale\_external\_claim}

\begin{itemize}
\item Source conversation: Dominium\_Architecture\_IV
\item Source file: docs/archive/conversations/Dominium\_Architecture\_IV/Dominium\_Architecture\_IV\_\_human\_readable\_chat\_report.txt
\item Claim: Archived conversation contains potentially stale external or baseline claim: SDK.
\item Conflict target: README.md / current platform and language baseline
\item Recommended disposition: quarantined until current repo and external facts are checked
\end{itemize}

\subsection{CONTRA-0075 - conversation\_vs\_current\_queue}

\begin{itemize}
\item Source conversation: Dominium\_Complete\_Conversation
\item Source file: docs/archive/conversations/Dominium\_Complete\_Conversation/Accompanying\_Report/Dominium\_Canon\_Portability\_Audit\_\_01\_human\_readable\_report.md
\item Claim: Archived conversation discusses work related to blocked scope gameplay.
\item Conflict target: .aide/queue/current.toml
\item Recommended disposition: preserve\_as\_history; require current queue authority before action
\end{itemize}

\subsection{CONTRA-0076 - stale\_external\_claim}

\begin{itemize}
\item Source conversation: Dominium\_Complete\_Conversation
\item Source file: docs/archive/conversations/Dominium\_Complete\_Conversation/Accompanying\_Report/Dominium\_Canon\_Portability\_Audit\_\_01\_human\_readable\_report.md
\item Claim: Archived conversation contains potentially stale external or baseline claim: C89.
\item Conflict target: README.md / current platform and language baseline
\item Recommended disposition: quarantined until current repo and external facts are checked
\end{itemize}

\subsection{CONTRA-0077 - stale\_external\_claim}

\begin{itemize}
\item Source conversation: Dominium\_Complete\_Conversation
\item Source file: docs/archive/conversations/Dominium\_Complete\_Conversation/Accompanying\_Report/Dominium\_Canon\_Portability\_Audit\_\_01\_human\_readable\_report.md
\item Claim: Archived conversation contains potentially stale external or baseline claim: C++98.
\item Conflict target: README.md / current platform and language baseline
\item Recommended disposition: quarantined until current repo and external facts are checked
\end{itemize}

\subsection{CONTRA-0078 - stale\_external\_claim}

\begin{itemize}
\item Source conversation: dominium\_domino\_codex\_planning
\item Source file: docs/archive/conversations/dominium\_domino\_codex\_planning/dominium\_domino\_codex\_planning\_\_human\_readable\_chat\_report.txt
\item Claim: Archived conversation contains potentially stale external or baseline claim: C89.
\item Conflict target: README.md / current platform and language baseline
\item Recommended disposition: quarantined until current repo and external facts are checked
\end{itemize}

\subsection{CONTRA-0079 - stale\_external\_claim}

\begin{itemize}
\item Source conversation: dominium\_domino\_codex\_planning
\item Source file: docs/archive/conversations/dominium\_domino\_codex\_planning/dominium\_domino\_codex\_planning\_\_human\_readable\_chat\_report.txt
\item Claim: Archived conversation contains potentially stale external or baseline claim: C++98.
\item Conflict target: README.md / current platform and language baseline
\item Recommended disposition: quarantined until current repo and external facts are checked
\end{itemize}

\subsection{CONTRA-0080 - stale\_external\_claim}

\begin{itemize}
\item Source conversation: dominium\_domino\_codex\_planning
\item Source file: docs/archive/conversations/dominium\_domino\_codex\_planning/dominium\_domino\_codex\_planning\_\_human\_readable\_chat\_report.txt
\item Claim: Archived conversation contains potentially stale external or baseline claim: ISO C89.
\item Conflict target: README.md / current platform and language baseline
\item Recommended disposition: quarantined until current repo and external facts are checked
\end{itemize}

\subsection{CONTRA-0081 - stale\_external\_claim}

\begin{itemize}
\item Source conversation: dominium\_domino\_codex\_planning
\item Source file: docs/archive/conversations/dominium\_domino\_codex\_planning/dominium\_domino\_codex\_planning\_\_human\_readable\_chat\_report.txt
\item Claim: Archived conversation contains potentially stale external or baseline claim: CP/M.
\item Conflict target: README.md / current platform and language baseline
\item Recommended disposition: quarantined until current repo and external facts are checked
\end{itemize}

\subsection{CONTRA-0082 - stale\_external\_claim}

\begin{itemize}
\item Source conversation: dominium\_domino\_codex\_planning
\item Source file: docs/archive/conversations/dominium\_domino\_codex\_planning/dominium\_domino\_codex\_planning\_\_human\_readable\_chat\_report.txt
\item Claim: Archived conversation contains potentially stale external or baseline claim: Android.
\item Conflict target: README.md / current platform and language baseline
\item Recommended disposition: quarantined until current repo and external facts are checked
\end{itemize}

\subsection{CONTRA-0083 - conversation\_vs\_current\_queue}

\begin{itemize}
\item Source conversation: dominium\_full\_conversation
\item Source file: docs/archive/conversations/dominium\_full\_conversation/companion\_report/dominium\_full\_conversation\_companion\_report\_\_01\_complete\_human\_report.md
\item Claim: Archived conversation discusses work related to blocked scope runtime\_module\_loader.
\item Conflict target: .aide/queue/current.toml
\item Recommended disposition: preserve\_as\_history; require current queue authority before action
\end{itemize}

\subsection{CONTRA-0084 - conversation\_vs\_current\_queue}

\begin{itemize}
\item Source conversation: dominium\_full\_conversation
\item Source file: docs/archive/conversations/dominium\_full\_conversation/companion\_report/dominium\_full\_conversation\_companion\_report\_\_01\_complete\_human\_report.md
\item Claim: Archived conversation discusses work related to blocked scope provider\_runtime.
\item Conflict target: .aide/queue/current.toml
\item Recommended disposition: preserve\_as\_history; require current queue authority before action
\end{itemize}

\subsection{CONTRA-0085 - conversation\_vs\_current\_queue}

\begin{itemize}
\item Source conversation: dominium\_full\_conversation
\item Source file: docs/archive/conversations/dominium\_full\_conversation/companion\_report/dominium\_full\_conversation\_companion\_report\_\_01\_complete\_human\_report.md
\item Claim: Archived conversation discusses work related to blocked scope package\_runtime.
\item Conflict target: .aide/queue/current.toml
\item Recommended disposition: preserve\_as\_history; require current queue authority before action
\end{itemize}

\subsection{CONTRA-0086 - conversation\_vs\_current\_queue}

\begin{itemize}
\item Source conversation: dominium\_full\_conversation
\item Source file: docs/archive/conversations/dominium\_full\_conversation/companion\_report/dominium\_full\_conversation\_companion\_report\_\_01\_complete\_human\_report.md
\item Claim: Archived conversation discusses work related to blocked scope gameplay.
\item Conflict target: .aide/queue/current.toml
\item Recommended disposition: preserve\_as\_history; require current queue authority before action
\end{itemize}

\subsection{CONTRA-0087 - conversation\_vs\_current\_queue}

\begin{itemize}
\item Source conversation: dominium\_full\_conversation
\item Source file: docs/archive/conversations/dominium\_full\_conversation/companion\_report/dominium\_full\_conversation\_companion\_report\_\_01\_complete\_human\_report.md
\item Claim: Archived conversation discusses work related to blocked scope native\_gui.
\item Conflict target: .aide/queue/current.toml
\item Recommended disposition: preserve\_as\_history; require current queue authority before action
\end{itemize}

\subsection{CONTRA-0088 - stale\_external\_claim}

\begin{itemize}
\item Source conversation: dominium\_full\_conversation
\item Source file: docs/archive/conversations/dominium\_full\_conversation/companion\_report/dominium\_full\_conversation\_companion\_report\_\_01\_complete\_human\_report.md
\item Claim: Archived conversation contains potentially stale external or baseline claim: SDK.
\item Conflict target: README.md / current platform and language baseline
\item Recommended disposition: quarantined until current repo and external facts are checked
\end{itemize}

\subsection{CONTRA-0089 - conversation\_vs\_current\_queue}

\begin{itemize}
\item Source conversation: dominium\_setup
\item Source file: docs/archive/conversations/dominium\_setup/dominium\_setup\_handoff\_\_human\_readable\_chat\_report.txt
\item Claim: Archived conversation discusses work related to blocked scope gameplay.
\item Conflict target: .aide/queue/current.toml
\item Recommended disposition: preserve\_as\_history; require current queue authority before action
\end{itemize}

\subsection{CONTRA-0090 - conversation\_vs\_current\_queue}

\begin{itemize}
\item Source conversation: dominium\_setup
\item Source file: docs/archive/conversations/dominium\_setup/dominium\_setup\_handoff\_\_human\_readable\_chat\_report.txt
\item Claim: Archived conversation discusses work related to blocked scope release\_publication.
\item Conflict target: .aide/queue/current.toml
\item Recommended disposition: preserve\_as\_history; require current queue authority before action
\end{itemize}

\subsection{CONTRA-0091 - stale\_external\_claim}

\begin{itemize}
\item Source conversation: dominium\_setup
\item Source file: docs/archive/conversations/dominium\_setup/dominium\_setup\_handoff\_\_human\_readable\_chat\_report.txt
\item Claim: Archived conversation contains potentially stale external or baseline claim: SDK.
\item Conflict target: README.md / current platform and language baseline
\item Recommended disposition: quarantined until current repo and external facts are checked
\end{itemize}

\subsection{CONTRA-0092 - conversation\_vs\_current\_queue}

\begin{itemize}
\item Source conversation: Domino\_Dominium\_Workbench
\item Source file: docs/archive/conversations/Domino\_Dominium\_Workbench/dominium\_workbench\_full\_conversation\_companion\_\_human\_readable\_detailed\_summary\_report.md
\item Claim: Archived conversation discusses work related to blocked scope broad\_workbench\_ui.
\item Conflict target: .aide/queue/current.toml
\item Recommended disposition: preserve\_as\_history; require current queue authority before action
\end{itemize}

\subsection{CONTRA-0093 - conversation\_vs\_current\_queue}

\begin{itemize}
\item Source conversation: Domino\_Dominium\_Workbench
\item Source file: docs/archive/conversations/Domino\_Dominium\_Workbench/dominium\_workbench\_full\_conversation\_companion\_\_human\_readable\_detailed\_summary\_report.md
\item Claim: Archived conversation discusses work related to blocked scope provider\_runtime.
\item Conflict target: .aide/queue/current.toml
\item Recommended disposition: preserve\_as\_history; require current queue authority before action
\end{itemize}

\subsection{CONTRA-0094 - conversation\_vs\_current\_queue}

\begin{itemize}
\item Source conversation: Domino\_Dominium\_Workbench
\item Source file: docs/archive/conversations/Domino\_Dominium\_Workbench/dominium\_workbench\_full\_conversation\_companion\_\_human\_readable\_detailed\_summary\_report.md
\item Claim: Archived conversation discusses work related to blocked scope gameplay.
\item Conflict target: .aide/queue/current.toml
\item Recommended disposition: preserve\_as\_history; require current queue authority before action
\end{itemize}

\subsection{CONTRA-0095 - conversation\_vs\_current\_queue}

\begin{itemize}
\item Source conversation: Domino\_Dominium\_Workbench
\item Source file: docs/archive/conversations/Domino\_Dominium\_Workbench/dominium\_workbench\_full\_conversation\_companion\_\_human\_readable\_detailed\_summary\_report.md
\item Claim: Archived conversation discusses work related to blocked scope native\_gui.
\item Conflict target: .aide/queue/current.toml
\item Recommended disposition: preserve\_as\_history; require current queue authority before action
\end{itemize}

\subsection{CONTRA-0096 - stale\_external\_claim}

\begin{itemize}
\item Source conversation: Domino\_Dominium\_Workbench
\item Source file: docs/archive/conversations/Domino\_Dominium\_Workbench/dominium\_workbench\_full\_conversation\_companion\_\_human\_readable\_detailed\_summary\_report.md
\item Claim: Archived conversation contains potentially stale external or baseline claim: SDK.
\item Conflict target: README.md / current platform and language baseline
\item Recommended disposition: quarantined until current repo and external facts are checked
\end{itemize}

\subsection{CONTRA-0097 - conversation\_vs\_current\_queue}

\begin{itemize}
\item Source conversation: domino\_engine\_refactor\_prompts
\item Source file: docs/archive/conversations/domino\_engine\_refactor\_prompts/dominium\_domino\_engine\_refactor\_prompts\_\_human\_readable\_chat\_report.txt
\item Claim: Archived conversation discusses work related to blocked scope gameplay.
\item Conflict target: .aide/queue/current.toml
\item Recommended disposition: preserve\_as\_history; require current queue authority before action
\end{itemize}

\subsection{CONTRA-0098 - stale\_external\_claim}

\begin{itemize}
\item Source conversation: domino\_engine\_refactor\_prompts
\item Source file: docs/archive/conversations/domino\_engine\_refactor\_prompts/dominium\_domino\_engine\_refactor\_prompts\_\_human\_readable\_chat\_report.txt
\item Claim: Archived conversation contains potentially stale external or baseline claim: C89.
\item Conflict target: README.md / current platform and language baseline
\item Recommended disposition: quarantined until current repo and external facts are checked
\end{itemize}

\subsection{CONTRA-0099 - stale\_external\_claim}

\begin{itemize}
\item Source conversation: domino\_engine\_refactor\_prompts
\item Source file: docs/archive/conversations/domino\_engine\_refactor\_prompts/dominium\_domino\_engine\_refactor\_prompts\_\_human\_readable\_chat\_report.txt
\item Claim: Archived conversation contains potentially stale external or baseline claim: C++98.
\item Conflict target: README.md / current platform and language baseline
\item Recommended disposition: quarantined until current repo and external facts are checked
\end{itemize}

\subsection{CONTRA-0100 - stale\_external\_claim}

\begin{itemize}
\item Source conversation: domino\_engine\_refactor\_prompts
\item Source file: docs/archive/conversations/domino\_engine\_refactor\_prompts/dominium\_domino\_engine\_refactor\_prompts\_\_human\_readable\_chat\_report.txt
\item Claim: Archived conversation contains potentially stale external or baseline claim: ISO C89.
\item Conflict target: README.md / current platform and language baseline
\item Recommended disposition: quarantined until current repo and external facts are checked
\end{itemize}

\subsection{CONTRA-0101 - conversation\_vs\_current\_queue}

\begin{itemize}
\item Source conversation: engine\_baseline\_architecture
\item Source file: docs/archive/conversations/engine\_baseline\_architecture/domino\_dominium\_engine\_baseline\_architecture\_\_01\_human\_readable\_report.md
\item Claim: Archived conversation discusses work related to blocked scope provider\_runtime.
\item Conflict target: .aide/queue/current.toml
\item Recommended disposition: preserve\_as\_history; require current queue authority before action
\end{itemize}

\subsection{CONTRA-0102 - conversation\_vs\_current\_queue}

\begin{itemize}
\item Source conversation: engine\_baseline\_architecture
\item Source file: docs/archive/conversations/engine\_baseline\_architecture/domino\_dominium\_engine\_baseline\_architecture\_\_01\_human\_readable\_report.md
\item Claim: Archived conversation discusses work related to blocked scope gameplay.
\item Conflict target: .aide/queue/current.toml
\item Recommended disposition: preserve\_as\_history; require current queue authority before action
\end{itemize}

\subsection{CONTRA-0103 - conversation\_vs\_current\_queue}

\begin{itemize}
\item Source conversation: Foundation\_Workbench\_Codex
\item Source file: docs/archive/conversations/Foundation\_Workbench\_Codex/Dominium\_Foundation\_Workbench\_Parallel\_Codex\_Preservation\_\_01\_human\_readable\_report.md
\item Claim: Archived conversation discusses work related to blocked scope broad\_workbench\_ui.
\item Conflict target: .aide/queue/current.toml
\item Recommended disposition: preserve\_as\_history; require current queue authority before action
\end{itemize}

\subsection{CONTRA-0104 - conversation\_vs\_current\_queue}

\begin{itemize}
\item Source conversation: Foundation\_Workbench\_Codex
\item Source file: docs/archive/conversations/Foundation\_Workbench\_Codex/Dominium\_Foundation\_Workbench\_Parallel\_Codex\_Preservation\_\_01\_human\_readable\_report.md
\item Claim: Archived conversation discusses work related to blocked scope runtime\_module\_loader.
\item Conflict target: .aide/queue/current.toml
\item Recommended disposition: preserve\_as\_history; require current queue authority before action
\end{itemize}

\subsection{CONTRA-0105 - conversation\_vs\_current\_queue}

\begin{itemize}
\item Source conversation: Foundation\_Workbench\_Codex
\item Source file: docs/archive/conversations/Foundation\_Workbench\_Codex/Dominium\_Foundation\_Workbench\_Parallel\_Codex\_Preservation\_\_01\_human\_readable\_report.md
\item Claim: Archived conversation discusses work related to blocked scope provider\_runtime.
\item Conflict target: .aide/queue/current.toml
\item Recommended disposition: preserve\_as\_history; require current queue authority before action
\end{itemize}

\subsection{CONTRA-0106 - conversation\_vs\_current\_queue}

\begin{itemize}
\item Source conversation: Foundation\_Workbench\_Codex
\item Source file: docs/archive/conversations/Foundation\_Workbench\_Codex/Dominium\_Foundation\_Workbench\_Parallel\_Codex\_Preservation\_\_01\_human\_readable\_report.md
\item Claim: Archived conversation discusses work related to blocked scope package\_runtime.
\item Conflict target: .aide/queue/current.toml
\item Recommended disposition: preserve\_as\_history; require current queue authority before action
\end{itemize}

\subsection{CONTRA-0107 - conversation\_vs\_current\_queue}

\begin{itemize}
\item Source conversation: Foundation\_Workbench\_Codex
\item Source file: docs/archive/conversations/Foundation\_Workbench\_Codex/Dominium\_Foundation\_Workbench\_Parallel\_Codex\_Preservation\_\_01\_human\_readable\_report.md
\item Claim: Archived conversation discusses work related to blocked scope gameplay.
\item Conflict target: .aide/queue/current.toml
\item Recommended disposition: preserve\_as\_history; require current queue authority before action
\end{itemize}

\subsection{CONTRA-0108 - conversation\_vs\_current\_queue}

\begin{itemize}
\item Source conversation: Foundation\_Workbench\_Codex
\item Source file: docs/archive/conversations/Foundation\_Workbench\_Codex/Dominium\_Foundation\_Workbench\_Parallel\_Codex\_Preservation\_\_01\_human\_readable\_report.md
\item Claim: Archived conversation discusses work related to blocked scope renderer\_implementation.
\item Conflict target: .aide/queue/current.toml
\item Recommended disposition: preserve\_as\_history; require current queue authority before action
\end{itemize}

\subsection{CONTRA-0109 - conversation\_vs\_current\_queue}

\begin{itemize}
\item Source conversation: Foundation\_Workbench\_Codex
\item Source file: docs/archive/conversations/Foundation\_Workbench\_Codex/Dominium\_Foundation\_Workbench\_Parallel\_Codex\_Preservation\_\_01\_human\_readable\_report.md
\item Claim: Archived conversation discusses work related to blocked scope native\_gui.
\item Conflict target: .aide/queue/current.toml
\item Recommended disposition: preserve\_as\_history; require current queue authority before action
\end{itemize}

\subsection{CONTRA-0110 - conversation\_vs\_current\_queue}

\begin{itemize}
\item Source conversation: Foundation\_Workbench\_Codex
\item Source file: docs/archive/conversations/Foundation\_Workbench\_Codex/Dominium\_Foundation\_Workbench\_Parallel\_Codex\_Preservation\_\_01\_human\_readable\_report.md
\item Claim: Archived conversation discusses work related to blocked scope release\_publication.
\item Conflict target: .aide/queue/current.toml
\item Recommended disposition: preserve\_as\_history; require current queue authority before action
\end{itemize}

\subsection{CONTRA-0111 - stale\_external\_claim}

\begin{itemize}
\item Source conversation: Foundation\_Workbench\_Codex
\item Source file: docs/archive/conversations/Foundation\_Workbench\_Codex/Dominium\_Foundation\_Workbench\_Parallel\_Codex\_Preservation\_\_01\_human\_readable\_report.md
\item Claim: Archived conversation contains potentially stale external or baseline claim: C89.
\item Conflict target: README.md / current platform and language baseline
\item Recommended disposition: quarantined until current repo and external facts are checked
\end{itemize}

\subsection{CONTRA-0112 - stale\_external\_claim}

\begin{itemize}
\item Source conversation: Foundation\_Workbench\_Codex
\item Source file: docs/archive/conversations/Foundation\_Workbench\_Codex/Dominium\_Foundation\_Workbench\_Parallel\_Codex\_Preservation\_\_01\_human\_readable\_report.md
\item Claim: Archived conversation contains potentially stale external or baseline claim: C++98.
\item Conflict target: README.md / current platform and language baseline
\item Recommended disposition: quarantined until current repo and external facts are checked
\end{itemize}

\subsection{CONTRA-0113 - stale\_external\_claim}

\begin{itemize}
\item Source conversation: Foundation\_Workbench\_Codex
\item Source file: docs/archive/conversations/Foundation\_Workbench\_Codex/Dominium\_Foundation\_Workbench\_Parallel\_Codex\_Preservation\_\_01\_human\_readable\_report.md
\item Claim: Archived conversation contains potentially stale external or baseline claim: SDK.
\item Conflict target: README.md / current platform and language baseline
\item Recommended disposition: quarantined until current repo and external facts are checked
\end{itemize}

\subsection{CONTRA-0114 - conversation\_vs\_current\_queue}

\begin{itemize}
\item Source conversation: Framework\_Open\_Source\_Provider
\item Source file: docs/archive/conversations/Framework\_Open\_Source\_Provider/Domino\_Framework\_Open\_Source\_Provider\_Architecture\_\_01\_human\_readable\_report.md
\item Claim: Archived conversation discusses work related to blocked scope broad\_workbench\_ui.
\item Conflict target: .aide/queue/current.toml
\item Recommended disposition: preserve\_as\_history; require current queue authority before action
\end{itemize}

\subsection{CONTRA-0115 - conversation\_vs\_current\_queue}

\begin{itemize}
\item Source conversation: Framework\_Open\_Source\_Provider
\item Source file: docs/archive/conversations/Framework\_Open\_Source\_Provider/Domino\_Framework\_Open\_Source\_Provider\_Architecture\_\_01\_human\_readable\_report.md
\item Claim: Archived conversation discusses work related to blocked scope provider\_runtime.
\item Conflict target: .aide/queue/current.toml
\item Recommended disposition: preserve\_as\_history; require current queue authority before action
\end{itemize}

\subsection{CONTRA-0116 - conversation\_vs\_current\_queue}

\begin{itemize}
\item Source conversation: Framework\_Open\_Source\_Provider
\item Source file: docs/archive/conversations/Framework\_Open\_Source\_Provider/Domino\_Framework\_Open\_Source\_Provider\_Architecture\_\_01\_human\_readable\_report.md
\item Claim: Archived conversation discusses work related to blocked scope gameplay.
\item Conflict target: .aide/queue/current.toml
\item Recommended disposition: preserve\_as\_history; require current queue authority before action
\end{itemize}

\subsection{CONTRA-0117 - conversation\_vs\_current\_queue}

\begin{itemize}
\item Source conversation: Framework\_Open\_Source\_Provider
\item Source file: docs/archive/conversations/Framework\_Open\_Source\_Provider/Domino\_Framework\_Open\_Source\_Provider\_Architecture\_\_01\_human\_readable\_report.md
\item Claim: Archived conversation discusses work related to blocked scope renderer\_implementation.
\item Conflict target: .aide/queue/current.toml
\item Recommended disposition: preserve\_as\_history; require current queue authority before action
\end{itemize}

\subsection{CONTRA-0118 - stale\_external\_claim}

\begin{itemize}
\item Source conversation: Framework\_Open\_Source\_Provider
\item Source file: docs/archive/conversations/Framework\_Open\_Source\_Provider/Domino\_Framework\_Open\_Source\_Provider\_Architecture\_\_01\_human\_readable\_report.md
\item Claim: Archived conversation contains potentially stale external or baseline claim: C++98.
\item Conflict target: README.md / current platform and language baseline
\item Recommended disposition: quarantined until current repo and external facts are checked
\end{itemize}

\subsection{CONTRA-0119 - stale\_external\_claim}

\begin{itemize}
\item Source conversation: Framework\_Open\_Source\_Provider
\item Source file: docs/archive/conversations/Framework\_Open\_Source\_Provider/Domino\_Framework\_Open\_Source\_Provider\_Architecture\_\_01\_human\_readable\_report.md
\item Claim: Archived conversation contains potentially stale external or baseline claim: SDK.
\item Conflict target: README.md / current platform and language baseline
\item Recommended disposition: quarantined until current repo and external facts are checked
\end{itemize}

\subsection{CONTRA-0120 - conversation\_vs\_current\_queue}

\begin{itemize}
\item Source conversation: gui\_binary\_content
\item Source file: docs/archive/conversations/gui\_binary\_content/dominium\_gui\_binary\_content\_handoff\_\_human\_readable\_chat\_report.txt
\item Claim: Archived conversation discusses work related to blocked scope gameplay.
\item Conflict target: .aide/queue/current.toml
\item Recommended disposition: preserve\_as\_history; require current queue authority before action
\end{itemize}

\subsection{CONTRA-0121 - conversation\_vs\_current\_queue}

\begin{itemize}
\item Source conversation: gui\_binary\_content
\item Source file: docs/archive/conversations/gui\_binary\_content/dominium\_gui\_binary\_content\_handoff\_\_human\_readable\_chat\_report.txt
\item Claim: Archived conversation discusses work related to blocked scope native\_gui.
\item Conflict target: .aide/queue/current.toml
\item Recommended disposition: preserve\_as\_history; require current queue authority before action
\end{itemize}

\subsection{CONTRA-0122 - stale\_external\_claim}

\begin{itemize}
\item Source conversation: gui\_binary\_content
\item Source file: docs/archive/conversations/gui\_binary\_content/dominium\_gui\_binary\_content\_handoff\_\_human\_readable\_chat\_report.txt
\item Claim: Archived conversation contains potentially stale external or baseline claim: iOS.
\item Conflict target: README.md / current platform and language baseline
\item Recommended disposition: quarantined until current repo and external facts are checked
\end{itemize}

\subsection{CONTRA-0123 - stale\_external\_claim}

\begin{itemize}
\item Source conversation: gui\_binary\_content
\item Source file: docs/archive/conversations/gui\_binary\_content/dominium\_gui\_binary\_content\_handoff\_\_human\_readable\_chat\_report.txt
\item Claim: Archived conversation contains potentially stale external or baseline claim: Android.
\item Conflict target: README.md / current platform and language baseline
\item Recommended disposition: quarantined until current repo and external facts are checked
\end{itemize}

\subsection{CONTRA-0124 - stale\_external\_claim}

\begin{itemize}
\item Source conversation: gui\_binary\_content
\item Source file: docs/archive/conversations/gui\_binary\_content/dominium\_gui\_binary\_content\_handoff\_\_human\_readable\_chat\_report.txt
\item Claim: Archived conversation contains potentially stale external or baseline claim: SDK.
\item Conflict target: README.md / current platform and language baseline
\item Recommended disposition: quarantined until current repo and external facts are checked
\end{itemize}

\subsection{CONTRA-0125 - conversation\_vs\_current\_queue}

\begin{itemize}
\item Source conversation: Language\_Platform\_Architecture
\item Source file: docs/archive/conversations/Language\_Platform\_Architecture/Dominium\_Language\_Platform\_Architecture\_Baseline\_\_01\_human\_readable\_report.md
\item Claim: Archived conversation discusses work related to blocked scope broad\_workbench\_ui.
\item Conflict target: .aide/queue/current.toml
\item Recommended disposition: preserve\_as\_history; require current queue authority before action
\end{itemize}

\subsection{CONTRA-0126 - conversation\_vs\_current\_queue}

\begin{itemize}
\item Source conversation: Language\_Platform\_Architecture
\item Source file: docs/archive/conversations/Language\_Platform\_Architecture/Dominium\_Language\_Platform\_Architecture\_Baseline\_\_01\_human\_readable\_report.md
\item Claim: Archived conversation discusses work related to blocked scope provider\_runtime.
\item Conflict target: .aide/queue/current.toml
\item Recommended disposition: preserve\_as\_history; require current queue authority before action
\end{itemize}

\subsection{CONTRA-0127 - conversation\_vs\_current\_queue}

\begin{itemize}
\item Source conversation: Language\_Platform\_Architecture
\item Source file: docs/archive/conversations/Language\_Platform\_Architecture/Dominium\_Language\_Platform\_Architecture\_Baseline\_\_01\_human\_readable\_report.md
\item Claim: Archived conversation discusses work related to blocked scope gameplay.
\item Conflict target: .aide/queue/current.toml
\item Recommended disposition: preserve\_as\_history; require current queue authority before action
\end{itemize}

\subsection{CONTRA-0128 - conversation\_vs\_current\_queue}

\begin{itemize}
\item Source conversation: Language\_Platform\_Architecture
\item Source file: docs/archive/conversations/Language\_Platform\_Architecture/Dominium\_Language\_Platform\_Architecture\_Baseline\_\_01\_human\_readable\_report.md
\item Claim: Archived conversation discusses work related to blocked scope renderer\_implementation.
\item Conflict target: .aide/queue/current.toml
\item Recommended disposition: preserve\_as\_history; require current queue authority before action
\end{itemize}

\subsection{CONTRA-0129 - stale\_external\_claim}

\begin{itemize}
\item Source conversation: Language\_Platform\_Architecture
\item Source file: docs/archive/conversations/Language\_Platform\_Architecture/Dominium\_Language\_Platform\_Architecture\_Baseline\_\_01\_human\_readable\_report.md
\item Claim: Archived conversation contains potentially stale external or baseline claim: C89.
\item Conflict target: README.md / current platform and language baseline
\item Recommended disposition: quarantined until current repo and external facts are checked
\end{itemize}

\subsection{CONTRA-0130 - stale\_external\_claim}

\begin{itemize}
\item Source conversation: Language\_Platform\_Architecture
\item Source file: docs/archive/conversations/Language\_Platform\_Architecture/Dominium\_Language\_Platform\_Architecture\_Baseline\_\_01\_human\_readable\_report.md
\item Claim: Archived conversation contains potentially stale external or baseline claim: C++98.
\item Conflict target: README.md / current platform and language baseline
\item Recommended disposition: quarantined until current repo and external facts are checked
\end{itemize}

\subsection{CONTRA-0131 - stale\_external\_claim}

\begin{itemize}
\item Source conversation: Language\_Platform\_Architecture
\item Source file: docs/archive/conversations/Language\_Platform\_Architecture/Dominium\_Language\_Platform\_Architecture\_Baseline\_\_01\_human\_readable\_report.md
\item Claim: Archived conversation contains potentially stale external or baseline claim: Windows 98.
\item Conflict target: README.md / current platform and language baseline
\item Recommended disposition: quarantined until current repo and external facts are checked
\end{itemize}

\subsection{CONTRA-0132 - stale\_external\_claim}

\begin{itemize}
\item Source conversation: Language\_Platform\_Architecture
\item Source file: docs/archive/conversations/Language\_Platform\_Architecture/Dominium\_Language\_Platform\_Architecture\_Baseline\_\_01\_human\_readable\_report.md
\item Claim: Archived conversation contains potentially stale external or baseline claim: Mac OS 9.
\item Conflict target: README.md / current platform and language baseline
\item Recommended disposition: quarantined until current repo and external facts are checked
\end{itemize}

\subsection{CONTRA-0133 - stale\_external\_claim}

\begin{itemize}
\item Source conversation: Language\_Platform\_Architecture
\item Source file: docs/archive/conversations/Language\_Platform\_Architecture/Dominium\_Language\_Platform\_Architecture\_Baseline\_\_01\_human\_readable\_report.md
\item Claim: Archived conversation contains potentially stale external or baseline claim: iOS.
\item Conflict target: README.md / current platform and language baseline
\item Recommended disposition: quarantined until current repo and external facts are checked
\end{itemize}

\subsection{CONTRA-0134 - stale\_external\_claim}

\begin{itemize}
\item Source conversation: Language\_Platform\_Architecture
\item Source file: docs/archive/conversations/Language\_Platform\_Architecture/Dominium\_Language\_Platform\_Architecture\_Baseline\_\_01\_human\_readable\_report.md
\item Claim: Archived conversation contains potentially stale external or baseline claim: Android.
\item Conflict target: README.md / current platform and language baseline
\item Recommended disposition: quarantined until current repo and external facts are checked
\end{itemize}

\subsection{CONTRA-0135 - stale\_external\_claim}

\begin{itemize}
\item Source conversation: Language\_Platform\_Architecture
\item Source file: docs/archive/conversations/Language\_Platform\_Architecture/Dominium\_Language\_Platform\_Architecture\_Baseline\_\_01\_human\_readable\_report.md
\item Claim: Archived conversation contains potentially stale external or baseline claim: SDK.
\item Conflict target: README.md / current platform and language baseline
\item Recommended disposition: quarantined until current repo and external facts are checked
\end{itemize}

\subsection{CONTRA-0136 - conversation\_vs\_current\_queue}

\begin{itemize}
\item Source conversation: launcher\_app\_layer
\item Source file: docs/archive/conversations/launcher\_app\_layer/dominium\_launcher\_app\_layer\_handoff\_\_human\_readable\_chat\_report.txt
\item Claim: Archived conversation discusses work related to blocked scope gameplay.
\item Conflict target: .aide/queue/current.toml
\item Recommended disposition: preserve\_as\_history; require current queue authority before action
\end{itemize}

\subsection{CONTRA-0137 - stale\_external\_claim}

\begin{itemize}
\item Source conversation: launcher\_app\_layer
\item Source file: docs/archive/conversations/launcher\_app\_layer/dominium\_launcher\_app\_layer\_handoff\_\_human\_readable\_chat\_report.txt
\item Claim: Archived conversation contains potentially stale external or baseline claim: SDK.
\item Conflict target: README.md / current platform and language baseline
\item Recommended disposition: quarantined until current repo and external facts are checked
\end{itemize}

\subsection{CONTRA-0138 - conversation\_vs\_current\_queue}

\begin{itemize}
\item Source conversation: Launcher\_Setup\_Architecture
\item Source file: docs/archive/conversations/Launcher\_Setup\_Architecture/Dominium\_Launcher\_Setup\_Architecture\_\_human\_readable\_chat\_report.txt
\item Claim: Archived conversation discusses work related to blocked scope gameplay.
\item Conflict target: .aide/queue/current.toml
\item Recommended disposition: preserve\_as\_history; require current queue authority before action
\end{itemize}

\subsection{CONTRA-0139 - conversation\_vs\_current\_queue}

\begin{itemize}
\item Source conversation: Launcher\_Setup\_Architecture
\item Source file: docs/archive/conversations/Launcher\_Setup\_Architecture/Dominium\_Launcher\_Setup\_Architecture\_\_human\_readable\_chat\_report.txt
\item Claim: Archived conversation discusses work related to blocked scope renderer\_implementation.
\item Conflict target: .aide/queue/current.toml
\item Recommended disposition: preserve\_as\_history; require current queue authority before action
\end{itemize}

\subsection{CONTRA-0140 - stale\_external\_claim}

\begin{itemize}
\item Source conversation: Launcher\_Setup\_Architecture
\item Source file: docs/archive/conversations/Launcher\_Setup\_Architecture/Dominium\_Launcher\_Setup\_Architecture\_\_human\_readable\_chat\_report.txt
\item Claim: Archived conversation contains potentially stale external or baseline claim: C89.
\item Conflict target: README.md / current platform and language baseline
\item Recommended disposition: quarantined until current repo and external facts are checked
\end{itemize}

\subsection{CONTRA-0141 - stale\_external\_claim}

\begin{itemize}
\item Source conversation: Launcher\_Setup\_Architecture
\item Source file: docs/archive/conversations/Launcher\_Setup\_Architecture/Dominium\_Launcher\_Setup\_Architecture\_\_human\_readable\_chat\_report.txt
\item Claim: Archived conversation contains potentially stale external or baseline claim: C++98.
\item Conflict target: README.md / current platform and language baseline
\item Recommended disposition: quarantined until current repo and external facts are checked
\end{itemize}

\subsection{CONTRA-0142 - stale\_external\_claim}

\begin{itemize}
\item Source conversation: Launcher\_Setup\_Architecture
\item Source file: docs/archive/conversations/Launcher\_Setup\_Architecture/Dominium\_Launcher\_Setup\_Architecture\_\_human\_readable\_chat\_report.txt
\item Claim: Archived conversation contains potentially stale external or baseline claim: CP/M.
\item Conflict target: README.md / current platform and language baseline
\item Recommended disposition: quarantined until current repo and external facts are checked
\end{itemize}

\subsection{CONTRA-0143 - stale\_external\_claim}

\begin{itemize}
\item Source conversation: Launcher\_Setup\_Architecture
\item Source file: docs/archive/conversations/Launcher\_Setup\_Architecture/Dominium\_Launcher\_Setup\_Architecture\_\_human\_readable\_chat\_report.txt
\item Claim: Archived conversation contains potentially stale external or baseline claim: SDK.
\item Conflict target: README.md / current platform and language baseline
\item Recommended disposition: quarantined until current repo and external facts are checked
\end{itemize}

\subsection{CONTRA-0144 - conversation\_vs\_current\_queue}

\begin{itemize}
\item Source conversation: Modularity\_AIDE\_Refactorability
\item Source file: docs/archive/conversations/Modularity\_AIDE\_Refactorability/Dominium\_Modularity\_AIDE\_Refactorability\_Architecture\_\_01\_human\_readable\_report.md
\item Claim: Archived conversation discusses work related to blocked scope renderer\_implementation.
\item Conflict target: .aide/queue/current.toml
\item Recommended disposition: preserve\_as\_history; require current queue authority before action
\end{itemize}

\subsection{CONTRA-0145 - conversation\_vs\_current\_queue}

\begin{itemize}
\item Source conversation: Modularity\_AIDE\_Refactorability
\item Source file: docs/archive/conversations/Modularity\_AIDE\_Refactorability/Dominium\_Modularity\_AIDE\_Refactorability\_Architecture\_\_01\_human\_readable\_report.md
\item Claim: Archived conversation discusses work related to blocked scope native\_gui.
\item Conflict target: .aide/queue/current.toml
\item Recommended disposition: preserve\_as\_history; require current queue authority before action
\end{itemize}

\subsection{CONTRA-0146 - stale\_external\_claim}

\begin{itemize}
\item Source conversation: Modularity\_AIDE\_Refactorability
\item Source file: docs/archive/conversations/Modularity\_AIDE\_Refactorability/Dominium\_Modularity\_AIDE\_Refactorability\_Architecture\_\_01\_human\_readable\_report.md
\item Claim: Archived conversation contains potentially stale external or baseline claim: C89.
\item Conflict target: README.md / current platform and language baseline
\item Recommended disposition: quarantined until current repo and external facts are checked
\end{itemize}

\subsection{CONTRA-0147 - conversation\_vs\_current\_queue}

\begin{itemize}
\item Source conversation: omega\_xi\_pi\_architecture\_future
\item Source file: docs/archive/conversations/omega\_xi\_pi\_architecture\_future/dominium\_omega\_xi\_pi\_architecture\_future\_proofing\_planning\_\_01\_human\_readable\_report.md
\item Claim: Archived conversation discusses work related to blocked scope runtime\_module\_loader.
\item Conflict target: .aide/queue/current.toml
\item Recommended disposition: preserve\_as\_history; require current queue authority before action
\end{itemize}

\subsection{CONTRA-0148 - conversation\_vs\_current\_queue}

\begin{itemize}
\item Source conversation: omega\_xi\_pi\_architecture\_future
\item Source file: docs/archive/conversations/omega\_xi\_pi\_architecture\_future/dominium\_omega\_xi\_pi\_architecture\_future\_proofing\_planning\_\_01\_human\_readable\_report.md
\item Claim: Archived conversation discusses work related to blocked scope native\_gui.
\item Conflict target: .aide/queue/current.toml
\item Recommended disposition: preserve\_as\_history; require current queue authority before action
\end{itemize}

\subsection{CONTRA-0149 - conversation\_vs\_current\_queue}

\begin{itemize}
\item Source conversation: OS\_Interface\_Repo\_Architecture
\item Source file: docs/archive/conversations/OS\_Interface\_Repo\_Architecture/Dominium\_OS\_Interface\_Repo\_Architecture\_\_01\_human\_readable\_report.md
\item Claim: Archived conversation discusses work related to blocked scope gameplay.
\item Conflict target: .aide/queue/current.toml
\item Recommended disposition: preserve\_as\_history; require current queue authority before action
\end{itemize}

\subsection{CONTRA-0150 - conversation\_vs\_current\_queue}

\begin{itemize}
\item Source conversation: OS\_Interface\_Repo\_Architecture
\item Source file: docs/archive/conversations/OS\_Interface\_Repo\_Architecture/Dominium\_OS\_Interface\_Repo\_Architecture\_\_01\_human\_readable\_report.md
\item Claim: Archived conversation discusses work related to blocked scope native\_gui.
\item Conflict target: .aide/queue/current.toml
\item Recommended disposition: preserve\_as\_history; require current queue authority before action
\end{itemize}

\subsection{CONTRA-0151 - conversation\_vs\_current\_queue}

\begin{itemize}
\item Source conversation: OS\_Interface\_Repo\_Architecture
\item Source file: docs/archive/conversations/OS\_Interface\_Repo\_Architecture/Dominium\_OS\_Interface\_Repo\_Architecture\_\_01\_human\_readable\_report.md
\item Claim: Archived conversation discusses work related to blocked scope release\_publication.
\item Conflict target: .aide/queue/current.toml
\item Recommended disposition: preserve\_as\_history; require current queue authority before action
\end{itemize}

\subsection{CONTRA-0152 - stale\_external\_claim}

\begin{itemize}
\item Source conversation: OS\_Interface\_Repo\_Architecture
\item Source file: docs/archive/conversations/OS\_Interface\_Repo\_Architecture/Dominium\_OS\_Interface\_Repo\_Architecture\_\_01\_human\_readable\_report.md
\item Claim: Archived conversation contains potentially stale external or baseline claim: Mac OS 8.
\item Conflict target: README.md / current platform and language baseline
\item Recommended disposition: quarantined until current repo and external facts are checked
\end{itemize}

\subsection{CONTRA-0153 - stale\_external\_claim}

\begin{itemize}
\item Source conversation: OS\_Interface\_Repo\_Architecture
\item Source file: docs/archive/conversations/OS\_Interface\_Repo\_Architecture/Dominium\_OS\_Interface\_Repo\_Architecture\_\_01\_human\_readable\_report.md
\item Claim: Archived conversation contains potentially stale external or baseline claim: Android.
\item Conflict target: README.md / current platform and language baseline
\item Recommended disposition: quarantined until current repo and external facts are checked
\end{itemize}

\subsection{CONTRA-0154 - stale\_external\_claim}

\begin{itemize}
\item Source conversation: OS\_Interface\_Repo\_Architecture
\item Source file: docs/archive/conversations/OS\_Interface\_Repo\_Architecture/Dominium\_OS\_Interface\_Repo\_Architecture\_\_01\_human\_readable\_report.md
\item Claim: Archived conversation contains potentially stale external or baseline claim: SDK.
\item Conflict target: README.md / current platform and language baseline
\item Recommended disposition: quarantined until current repo and external facts are checked
\end{itemize}

\subsection{CONTRA-0155 - conversation\_vs\_current\_queue}

\begin{itemize}
\item Source conversation: platform\_renderer\_api\_plan
\item Source file: docs/archive/conversations/platform\_renderer\_api\_plan/dominium\_codex\_platform\_renderer\_api\_plan\_\_human\_readable\_chat\_report.txt
\item Claim: Archived conversation discusses work related to blocked scope gameplay.
\item Conflict target: .aide/queue/current.toml
\item Recommended disposition: preserve\_as\_history; require current queue authority before action
\end{itemize}

\subsection{CONTRA-0156 - conversation\_vs\_current\_queue}

\begin{itemize}
\item Source conversation: platform\_renderer\_api\_plan
\item Source file: docs/archive/conversations/platform\_renderer\_api\_plan/dominium\_codex\_platform\_renderer\_api\_plan\_\_human\_readable\_chat\_report.txt
\item Claim: Archived conversation discusses work related to blocked scope renderer\_implementation.
\item Conflict target: .aide/queue/current.toml
\item Recommended disposition: preserve\_as\_history; require current queue authority before action
\end{itemize}

\subsection{CONTRA-0157 - conversation\_vs\_current\_queue}

\begin{itemize}
\item Source conversation: platform\_renderer\_api\_plan
\item Source file: docs/archive/conversations/platform\_renderer\_api\_plan/dominium\_codex\_platform\_renderer\_api\_plan\_\_human\_readable\_chat\_report.txt
\item Claim: Archived conversation discusses work related to blocked scope native\_gui.
\item Conflict target: .aide/queue/current.toml
\item Recommended disposition: preserve\_as\_history; require current queue authority before action
\end{itemize}

\subsection{CONTRA-0158 - stale\_external\_claim}

\begin{itemize}
\item Source conversation: platform\_renderer\_api\_plan
\item Source file: docs/archive/conversations/platform\_renderer\_api\_plan/dominium\_codex\_platform\_renderer\_api\_plan\_\_human\_readable\_chat\_report.txt
\item Claim: Archived conversation contains potentially stale external or baseline claim: C89.
\item Conflict target: README.md / current platform and language baseline
\item Recommended disposition: quarantined until current repo and external facts are checked
\end{itemize}

\subsection{CONTRA-0159 - stale\_external\_claim}

\begin{itemize}
\item Source conversation: platform\_renderer\_api\_plan
\item Source file: docs/archive/conversations/platform\_renderer\_api\_plan/dominium\_codex\_platform\_renderer\_api\_plan\_\_human\_readable\_chat\_report.txt
\item Claim: Archived conversation contains potentially stale external or baseline claim: C++98.
\item Conflict target: README.md / current platform and language baseline
\item Recommended disposition: quarantined until current repo and external facts are checked
\end{itemize}

\subsection{CONTRA-0160 - stale\_external\_claim}

\begin{itemize}
\item Source conversation: platform\_renderer\_api\_plan
\item Source file: docs/archive/conversations/platform\_renderer\_api\_plan/dominium\_codex\_platform\_renderer\_api\_plan\_\_human\_readable\_chat\_report.txt
\item Claim: Archived conversation contains potentially stale external or baseline claim: ISO C89.
\item Conflict target: README.md / current platform and language baseline
\item Recommended disposition: quarantined until current repo and external facts are checked
\end{itemize}

\subsection{CONTRA-0161 - stale\_external\_claim}

\begin{itemize}
\item Source conversation: platform\_renderer\_api\_plan
\item Source file: docs/archive/conversations/platform\_renderer\_api\_plan/dominium\_codex\_platform\_renderer\_api\_plan\_\_human\_readable\_chat\_report.txt
\item Claim: Archived conversation contains potentially stale external or baseline claim: CP/M.
\item Conflict target: README.md / current platform and language baseline
\item Recommended disposition: quarantined until current repo and external facts are checked
\end{itemize}

\subsection{CONTRA-0162 - stale\_external\_claim}

\begin{itemize}
\item Source conversation: platform\_renderer\_api\_plan
\item Source file: docs/archive/conversations/platform\_renderer\_api\_plan/dominium\_codex\_platform\_renderer\_api\_plan\_\_human\_readable\_chat\_report.txt
\item Claim: Archived conversation contains potentially stale external or baseline claim: Android.
\item Conflict target: README.md / current platform and language baseline
\item Recommended disposition: quarantined until current repo and external facts are checked
\end{itemize}

\subsection{CONTRA-0163 - stale\_external\_claim}

\begin{itemize}
\item Source conversation: platform\_renderer\_api\_plan
\item Source file: docs/archive/conversations/platform\_renderer\_api\_plan/dominium\_codex\_platform\_renderer\_api\_plan\_\_human\_readable\_chat\_report.txt
\item Claim: Archived conversation contains potentially stale external or baseline claim: SDK.
\item Conflict target: README.md / current platform and language baseline
\item Recommended disposition: quarantined until current repo and external facts are checked
\end{itemize}

\subsection{CONTRA-0164 - conversation\_vs\_current\_queue}

\begin{itemize}
\item Source conversation: platform\_support
\item Source file: docs/archive/conversations/platform\_support/dominium\_platform\_support\_planning\_\_human\_readable\_chat\_report.txt
\item Claim: Archived conversation discusses work related to blocked scope gameplay.
\item Conflict target: .aide/queue/current.toml
\item Recommended disposition: preserve\_as\_history; require current queue authority before action
\end{itemize}

\subsection{CONTRA-0165 - conversation\_vs\_current\_queue}

\begin{itemize}
\item Source conversation: platform\_support
\item Source file: docs/archive/conversations/platform\_support/dominium\_platform\_support\_planning\_\_human\_readable\_chat\_report.txt
\item Claim: Archived conversation discusses work related to blocked scope renderer\_implementation.
\item Conflict target: .aide/queue/current.toml
\item Recommended disposition: preserve\_as\_history; require current queue authority before action
\end{itemize}

\subsection{CONTRA-0166 - conversation\_vs\_current\_queue}

\begin{itemize}
\item Source conversation: platform\_support
\item Source file: docs/archive/conversations/platform\_support/dominium\_platform\_support\_planning\_\_human\_readable\_chat\_report.txt
\item Claim: Archived conversation discusses work related to blocked scope release\_publication.
\item Conflict target: .aide/queue/current.toml
\item Recommended disposition: preserve\_as\_history; require current queue authority before action
\end{itemize}

\subsection{CONTRA-0167 - stale\_external\_claim}

\begin{itemize}
\item Source conversation: platform\_support
\item Source file: docs/archive/conversations/platform\_support/dominium\_platform\_support\_planning\_\_human\_readable\_chat\_report.txt
\item Claim: Archived conversation contains potentially stale external or baseline claim: CP/M.
\item Conflict target: README.md / current platform and language baseline
\item Recommended disposition: quarantined until current repo and external facts are checked
\end{itemize}

\subsection{CONTRA-0168 - stale\_external\_claim}

\begin{itemize}
\item Source conversation: platform\_support
\item Source file: docs/archive/conversations/platform\_support/dominium\_platform\_support\_planning\_\_human\_readable\_chat\_report.txt
\item Claim: Archived conversation contains potentially stale external or baseline claim: iOS.
\item Conflict target: README.md / current platform and language baseline
\item Recommended disposition: quarantined until current repo and external facts are checked
\end{itemize}

\subsection{CONTRA-0169 - stale\_external\_claim}

\begin{itemize}
\item Source conversation: platform\_support
\item Source file: docs/archive/conversations/platform\_support/dominium\_platform\_support\_planning\_\_human\_readable\_chat\_report.txt
\item Claim: Archived conversation contains potentially stale external or baseline claim: Android.
\item Conflict target: README.md / current platform and language baseline
\item Recommended disposition: quarantined until current repo and external facts are checked
\end{itemize}

\subsection{CONTRA-0170 - stale\_external\_claim}

\begin{itemize}
\item Source conversation: platform\_support
\item Source file: docs/archive/conversations/platform\_support/dominium\_platform\_support\_planning\_\_human\_readable\_chat\_report.txt
\item Claim: Archived conversation contains potentially stale external or baseline claim: WebGPU.
\item Conflict target: README.md / current platform and language baseline
\item Recommended disposition: quarantined until current repo and external facts are checked
\end{itemize}

\subsection{CONTRA-0171 - stale\_external\_claim}

\begin{itemize}
\item Source conversation: platform\_support
\item Source file: docs/archive/conversations/platform\_support/dominium\_platform\_support\_planning\_\_human\_readable\_chat\_report.txt
\item Claim: Archived conversation contains potentially stale external or baseline claim: console program.
\item Conflict target: README.md / current platform and language baseline
\item Recommended disposition: quarantined until current repo and external facts are checked
\end{itemize}

\subsection{CONTRA-0172 - stale\_external\_claim}

\begin{itemize}
\item Source conversation: platform\_support
\item Source file: docs/archive/conversations/platform\_support/dominium\_platform\_support\_planning\_\_human\_readable\_chat\_report.txt
\item Claim: Archived conversation contains potentially stale external or baseline claim: SDK.
\item Conflict target: README.md / current platform and language baseline
\item Recommended disposition: quarantined until current repo and external facts are checked
\end{itemize}

\subsection{CONTRA-0173 - conversation\_vs\_current\_queue}

\begin{itemize}
\item Source conversation: Portability\_Assurance\_Future\_Proof
\item Source file: docs/archive/conversations/Portability\_Assurance\_Future\_Proof/Domino\_Dominium\_Portability\_Assurance\_Future\_Proof\_Architecture\_\_01\_human\_readable\_report.md
\item Claim: Archived conversation discusses work related to blocked scope gameplay.
\item Conflict target: .aide/queue/current.toml
\item Recommended disposition: preserve\_as\_history; require current queue authority before action
\end{itemize}

\subsection{CONTRA-0174 - stale\_external\_claim}

\begin{itemize}
\item Source conversation: Portability\_Assurance\_Future\_Proof
\item Source file: docs/archive/conversations/Portability\_Assurance\_Future\_Proof/Domino\_Dominium\_Portability\_Assurance\_Future\_Proof\_Architecture\_\_01\_human\_readable\_report.md
\item Claim: Archived conversation contains potentially stale external or baseline claim: C89.
\item Conflict target: README.md / current platform and language baseline
\item Recommended disposition: quarantined until current repo and external facts are checked
\end{itemize}

\subsection{CONTRA-0175 - stale\_external\_claim}

\begin{itemize}
\item Source conversation: Portability\_Assurance\_Future\_Proof
\item Source file: docs/archive/conversations/Portability\_Assurance\_Future\_Proof/Domino\_Dominium\_Portability\_Assurance\_Future\_Proof\_Architecture\_\_01\_human\_readable\_report.md
\item Claim: Archived conversation contains potentially stale external or baseline claim: C++98.
\item Conflict target: README.md / current platform and language baseline
\item Recommended disposition: quarantined until current repo and external facts are checked
\end{itemize}

\subsection{CONTRA-0176 - stale\_external\_claim}

\begin{itemize}
\item Source conversation: readme\_ports\_determinism
\item Source file: docs/archive/conversations/readme\_ports\_determinism/dominium\_readme\_ports\_determinism\_\_human\_readable\_chat\_report.txt
\item Claim: Archived conversation contains potentially stale external or baseline claim: C89.
\item Conflict target: README.md / current platform and language baseline
\item Recommended disposition: quarantined until current repo and external facts are checked
\end{itemize}

\subsection{CONTRA-0177 - stale\_external\_claim}

\begin{itemize}
\item Source conversation: readme\_ports\_determinism
\item Source file: docs/archive/conversations/readme\_ports\_determinism/dominium\_readme\_ports\_determinism\_\_human\_readable\_chat\_report.txt
\item Claim: Archived conversation contains potentially stale external or baseline claim: C++98.
\item Conflict target: README.md / current platform and language baseline
\item Recommended disposition: quarantined until current repo and external facts are checked
\end{itemize}

\subsection{CONTRA-0178 - stale\_external\_claim}

\begin{itemize}
\item Source conversation: readme\_ports\_determinism
\item Source file: docs/archive/conversations/readme\_ports\_determinism/dominium\_readme\_ports\_determinism\_\_human\_readable\_chat\_report.txt
\item Claim: Archived conversation contains potentially stale external or baseline claim: CP/M.
\item Conflict target: README.md / current platform and language baseline
\item Recommended disposition: quarantined until current repo and external facts are checked
\end{itemize}

\subsection{CONTRA-0179 - stale\_external\_claim}

\begin{itemize}
\item Source conversation: readme\_ports\_determinism
\item Source file: docs/archive/conversations/readme\_ports\_determinism/dominium\_readme\_ports\_determinism\_\_human\_readable\_chat\_report.txt
\item Claim: Archived conversation contains potentially stale external or baseline claim: 286.
\item Conflict target: README.md / current platform and language baseline
\item Recommended disposition: quarantined until current repo and external facts are checked
\end{itemize}

\subsection{CONTRA-0180 - stale\_external\_claim}

\begin{itemize}
\item Source conversation: Refactor\_Architecture
\item Source file: docs/archive/conversations/Refactor\_Architecture/Dominium\_Domino\_Refactor\_Architecture\_\_human\_readable\_chat\_report.txt
\item Claim: Archived conversation contains potentially stale external or baseline claim: C89.
\item Conflict target: README.md / current platform and language baseline
\item Recommended disposition: quarantined until current repo and external facts are checked
\end{itemize}

\subsection{CONTRA-0181 - stale\_external\_claim}

\begin{itemize}
\item Source conversation: Refactor\_Architecture
\item Source file: docs/archive/conversations/Refactor\_Architecture/Dominium\_Domino\_Refactor\_Architecture\_\_human\_readable\_chat\_report.txt
\item Claim: Archived conversation contains potentially stale external or baseline claim: CP/M.
\item Conflict target: README.md / current platform and language baseline
\item Recommended disposition: quarantined until current repo and external facts are checked
\end{itemize}

\subsection{CONTRA-0182 - stale\_external\_claim}

\begin{itemize}
\item Source conversation: Refactor\_Architecture
\item Source file: docs/archive/conversations/Refactor\_Architecture/Dominium\_Domino\_Refactor\_Architecture\_\_human\_readable\_chat\_report.txt
\item Claim: Archived conversation contains potentially stale external or baseline claim: Android.
\item Conflict target: README.md / current platform and language baseline
\item Recommended disposition: quarantined until current repo and external facts are checked
\end{itemize}

\subsection{CONTRA-0183 - stale\_external\_claim}

\begin{itemize}
\item Source conversation: Refactor\_Architecture
\item Source file: docs/archive/conversations/Refactor\_Architecture/Dominium\_Domino\_Refactor\_Architecture\_\_human\_readable\_chat\_report.txt
\item Claim: Archived conversation contains potentially stale external or baseline claim: SDK.
\item Conflict target: README.md / current platform and language baseline
\item Recommended disposition: quarantined until current repo and external facts are checked
\end{itemize}

\subsection{CONTRA-0184 - conversation\_vs\_current\_queue}

\begin{itemize}
\item Source conversation: Refactor\_Control\_Plane
\item Source file: docs/archive/conversations/Refactor\_Control\_Plane/AIDE\_XStack\_Dominium\_Refactor\_Control\_Plane\_\_01\_human\_readable\_report.md
\item Claim: Archived conversation discusses work related to blocked scope provider\_runtime.
\item Conflict target: .aide/queue/current.toml
\item Recommended disposition: preserve\_as\_history; require current queue authority before action
\end{itemize}

\subsection{CONTRA-0185 - stale\_external\_claim}

\begin{itemize}
\item Source conversation: Refactor\_Control\_Plane
\item Source file: docs/archive/conversations/Refactor\_Control\_Plane/AIDE\_XStack\_Dominium\_Refactor\_Control\_Plane\_\_01\_human\_readable\_report.md
\item Claim: Archived conversation contains potentially stale external or baseline claim: SDK.
\item Conflict target: README.md / current platform and language baseline
\item Recommended disposition: quarantined until current repo and external facts are checked
\end{itemize}

\subsection{CONTRA-0186 - conversation\_vs\_current\_queue}

\begin{itemize}
\item Source conversation: Release\_Identity\_and\_Versioning
\item Source file: docs/archive/conversations/Release\_Identity\_and\_Versioning/Dominium\_XStack\_Release\_Identity\_and\_Versioning\_\_01\_human\_readable\_report.md
\item Claim: Archived conversation discusses work related to blocked scope release\_publication.
\item Conflict target: .aide/queue/current.toml
\item Recommended disposition: preserve\_as\_history; require current queue authority before action
\end{itemize}

\subsection{CONTRA-0187 - stale\_external\_claim}

\begin{itemize}
\item Source conversation: Release\_Identity\_and\_Versioning
\item Source file: docs/archive/conversations/Release\_Identity\_and\_Versioning/Dominium\_XStack\_Release\_Identity\_and\_Versioning\_\_01\_human\_readable\_report.md
\item Claim: Archived conversation contains potentially stale external or baseline claim: SDK.
\item Conflict target: README.md / current platform and language baseline
\item Recommended disposition: quarantined until current repo and external facts are checked
\end{itemize}

\subsection{CONTRA-0188 - conversation\_vs\_current\_queue}

\begin{itemize}
\item Source conversation: testx\_repox\_governance
\item Source file: docs/archive/conversations/testx\_repox\_governance/dominium\_testx\_repox\_governance\_handoff\_\_human\_readable\_chat\_report.txt
\item Claim: Archived conversation discusses work related to blocked scope provider\_runtime.
\item Conflict target: .aide/queue/current.toml
\item Recommended disposition: preserve\_as\_history; require current queue authority before action
\end{itemize}

\subsection{CONTRA-0189 - conversation\_vs\_current\_queue}

\begin{itemize}
\item Source conversation: testx\_repox\_governance
\item Source file: docs/archive/conversations/testx\_repox\_governance/dominium\_testx\_repox\_governance\_handoff\_\_human\_readable\_chat\_report.txt
\item Claim: Archived conversation discusses work related to blocked scope gameplay.
\item Conflict target: .aide/queue/current.toml
\item Recommended disposition: preserve\_as\_history; require current queue authority before action
\end{itemize}

\subsection{CONTRA-0190 - conversation\_vs\_current\_queue}

\begin{itemize}
\item Source conversation: testx\_repox\_governance
\item Source file: docs/archive/conversations/testx\_repox\_governance/dominium\_testx\_repox\_governance\_handoff\_\_human\_readable\_chat\_report.txt
\item Claim: Archived conversation discusses work related to blocked scope renderer\_implementation.
\item Conflict target: .aide/queue/current.toml
\item Recommended disposition: preserve\_as\_history; require current queue authority before action
\end{itemize}

\subsection{CONTRA-0191 - conversation\_vs\_current\_queue}

\begin{itemize}
\item Source conversation: testx\_repox\_governance
\item Source file: docs/archive/conversations/testx\_repox\_governance/dominium\_testx\_repox\_governance\_handoff\_\_human\_readable\_chat\_report.txt
\item Claim: Archived conversation discusses work related to blocked scope release\_publication.
\item Conflict target: .aide/queue/current.toml
\item Recommended disposition: preserve\_as\_history; require current queue authority before action
\end{itemize}

\subsection{CONTRA-0192 - stale\_external\_claim}

\begin{itemize}
\item Source conversation: testx\_repox\_governance
\item Source file: docs/archive/conversations/testx\_repox\_governance/dominium\_testx\_repox\_governance\_handoff\_\_human\_readable\_chat\_report.txt
\item Claim: Archived conversation contains potentially stale external or baseline claim: C89.
\item Conflict target: README.md / current platform and language baseline
\item Recommended disposition: quarantined until current repo and external facts are checked
\end{itemize}

\subsection{CONTRA-0193 - stale\_external\_claim}

\begin{itemize}
\item Source conversation: testx\_repox\_governance
\item Source file: docs/archive/conversations/testx\_repox\_governance/dominium\_testx\_repox\_governance\_handoff\_\_human\_readable\_chat\_report.txt
\item Claim: Archived conversation contains potentially stale external or baseline claim: C++98.
\item Conflict target: README.md / current platform and language baseline
\item Recommended disposition: quarantined until current repo and external facts are checked
\end{itemize}

\subsection{CONTRA-0194 - stale\_external\_claim}

\begin{itemize}
\item Source conversation: testx\_repox\_governance
\item Source file: docs/archive/conversations/testx\_repox\_governance/dominium\_testx\_repox\_governance\_handoff\_\_human\_readable\_chat\_report.txt
\item Claim: Archived conversation contains potentially stale external or baseline claim: iOS.
\item Conflict target: README.md / current platform and language baseline
\item Recommended disposition: quarantined until current repo and external facts are checked
\end{itemize}

\subsection{CONTRA-0195 - stale\_external\_claim}

\begin{itemize}
\item Source conversation: testx\_repox\_governance
\item Source file: docs/archive/conversations/testx\_repox\_governance/dominium\_testx\_repox\_governance\_handoff\_\_human\_readable\_chat\_report.txt
\item Claim: Archived conversation contains potentially stale external or baseline claim: SDK.
\item Conflict target: README.md / current platform and language baseline
\item Recommended disposition: quarantined until current repo and external facts are checked
\end{itemize}

\subsection{CONTRA-0196 - stale\_external\_claim}

\begin{itemize}
\item Source conversation: Timekeeping\_and\_2038\_Resilience
\item Source file: docs/archive/conversations/Timekeeping\_and\_2038\_Resilience/Dominium\_Timekeeping\_and\_2038\_Resilience\_\_01\_human\_readable\_report.md
\item Claim: Archived conversation contains potentially stale external or baseline claim: C89.
\item Conflict target: README.md / current platform and language baseline
\item Recommended disposition: quarantined until current repo and external facts are checked
\end{itemize}

\subsection{CONTRA-0197 - stale\_external\_claim}

\begin{itemize}
\item Source conversation: Timekeeping\_and\_2038\_Resilience
\item Source file: docs/archive/conversations/Timekeeping\_and\_2038\_Resilience/Dominium\_Timekeeping\_and\_2038\_Resilience\_\_01\_human\_readable\_report.md
\item Claim: Archived conversation contains potentially stale external or baseline claim: C++98.
\item Conflict target: README.md / current platform and language baseline
\item Recommended disposition: quarantined until current repo and external facts are checked
\end{itemize}

\subsection{CONTRA-0198 - conversation\_vs\_current\_queue}

\begin{itemize}
\item Source conversation: UE6\_Domino\_Deterministic\_Universe
\item Source file: docs/archive/conversations/UE6\_Domino\_Deterministic\_Universe/Dominium\_UE6\_Domino\_Deterministic\_Universe\_\_01\_human\_readable\_report.md
\item Claim: Archived conversation discusses work related to blocked scope broad\_workbench\_ui.
\item Conflict target: .aide/queue/current.toml
\item Recommended disposition: preserve\_as\_history; require current queue authority before action
\end{itemize}

\subsection{CONTRA-0199 - conversation\_vs\_current\_queue}

\begin{itemize}
\item Source conversation: UE6\_Domino\_Deterministic\_Universe
\item Source file: docs/archive/conversations/UE6\_Domino\_Deterministic\_Universe/Dominium\_UE6\_Domino\_Deterministic\_Universe\_\_01\_human\_readable\_report.md
\item Claim: Archived conversation discusses work related to blocked scope gameplay.
\item Conflict target: .aide/queue/current.toml
\item Recommended disposition: preserve\_as\_history; require current queue authority before action
\end{itemize}

\subsection{CONTRA-0200 - stale\_external\_claim}

\begin{itemize}
\item Source conversation: UE6\_Domino\_Deterministic\_Universe
\item Source file: docs/archive/conversations/UE6\_Domino\_Deterministic\_Universe/Dominium\_UE6\_Domino\_Deterministic\_Universe\_\_01\_human\_readable\_report.md
\item Claim: Archived conversation contains potentially stale external or baseline claim: C89.
\item Conflict target: README.md / current platform and language baseline
\item Recommended disposition: quarantined until current repo and external facts are checked
\end{itemize}

\subsection{CONTRA-0201 - stale\_external\_claim}

\begin{itemize}
\item Source conversation: UE6\_Domino\_Deterministic\_Universe
\item Source file: docs/archive/conversations/UE6\_Domino\_Deterministic\_Universe/Dominium\_UE6\_Domino\_Deterministic\_Universe\_\_01\_human\_readable\_report.md
\item Claim: Archived conversation contains potentially stale external or baseline claim: C++98.
\item Conflict target: README.md / current platform and language baseline
\item Recommended disposition: quarantined until current repo and external facts are checked
\end{itemize}

\subsection{CONTRA-0202 - conversation\_vs\_current\_queue}

\begin{itemize}
\item Source conversation: ui\_editor\_tool\_editor\_planning
\item Source file: docs/archive/conversations/ui\_editor\_tool\_editor\_planning/dominium\_ui\_editor\_tool\_editor\_planning\_\_human\_readable\_chat\_report.txt
\item Claim: Archived conversation discusses work related to blocked scope renderer\_implementation.
\item Conflict target: .aide/queue/current.toml
\item Recommended disposition: preserve\_as\_history; require current queue authority before action
\end{itemize}

\subsection{CONTRA-0203 - conversation\_vs\_current\_queue}

\begin{itemize}
\item Source conversation: ui\_editor\_tool\_editor\_planning
\item Source file: docs/archive/conversations/ui\_editor\_tool\_editor\_planning/dominium\_ui\_editor\_tool\_editor\_planning\_\_human\_readable\_chat\_report.txt
\item Claim: Archived conversation discusses work related to blocked scope native\_gui.
\item Conflict target: .aide/queue/current.toml
\item Recommended disposition: preserve\_as\_history; require current queue authority before action
\end{itemize}

\subsection{CONTRA-0204 - stale\_external\_claim}

\begin{itemize}
\item Source conversation: ui\_editor\_tool\_editor\_planning
\item Source file: docs/archive/conversations/ui\_editor\_tool\_editor\_planning/dominium\_ui\_editor\_tool\_editor\_planning\_\_human\_readable\_chat\_report.txt
\item Claim: Archived conversation contains potentially stale external or baseline claim: C++98.
\item Conflict target: README.md / current platform and language baseline
\item Recommended disposition: quarantined until current repo and external facts are checked
\end{itemize}

\subsection{CONTRA-0205 - stale\_external\_claim}

\begin{itemize}
\item Source conversation: ui\_editor\_tool\_editor\_planning
\item Source file: docs/archive/conversations/ui\_editor\_tool\_editor\_planning/dominium\_ui\_editor\_tool\_editor\_planning\_\_human\_readable\_chat\_report.txt
\item Claim: Archived conversation contains potentially stale external or baseline claim: Android.
\item Conflict target: README.md / current platform and language baseline
\item Recommended disposition: quarantined until current repo and external facts are checked
\end{itemize}

\subsection{CONTRA-0206 - conversation\_vs\_current\_queue}

\begin{itemize}
\item Source conversation: Universe\_Explorer\_Planning
\item Source file: docs/archive/conversations/Universe\_Explorer\_Planning/Dominium\_Universe\_Explorer\_Planning\_Handoff\_\_01\_human\_readable\_report.md
\item Claim: Archived conversation discusses work related to blocked scope broad\_workbench\_ui.
\item Conflict target: .aide/queue/current.toml
\item Recommended disposition: preserve\_as\_history; require current queue authority before action
\end{itemize}

\subsection{CONTRA-0207 - conversation\_vs\_current\_queue}

\begin{itemize}
\item Source conversation: Universe\_Explorer\_Planning
\item Source file: docs/archive/conversations/Universe\_Explorer\_Planning/Dominium\_Universe\_Explorer\_Planning\_Handoff\_\_01\_human\_readable\_report.md
\item Claim: Archived conversation discusses work related to blocked scope provider\_runtime.
\item Conflict target: .aide/queue/current.toml
\item Recommended disposition: preserve\_as\_history; require current queue authority before action
\end{itemize}

\subsection{CONTRA-0208 - conversation\_vs\_current\_queue}

\begin{itemize}
\item Source conversation: Universe\_Explorer\_Planning
\item Source file: docs/archive/conversations/Universe\_Explorer\_Planning/Dominium\_Universe\_Explorer\_Planning\_Handoff\_\_01\_human\_readable\_report.md
\item Claim: Archived conversation discusses work related to blocked scope gameplay.
\item Conflict target: .aide/queue/current.toml
\item Recommended disposition: preserve\_as\_history; require current queue authority before action
\end{itemize}

\subsection{CONTRA-0209 - conversation\_vs\_current\_queue}

\begin{itemize}
\item Source conversation: Workbench\_AIDE\_Product\_Spine
\item Source file: docs/archive/conversations/Workbench\_AIDE\_Product\_Spine/Dominium\_Architecture\_Workbench\_AIDE\_Product\_Spine\_Planning\_\_01\_human\_readable\_report.md
\item Claim: Archived conversation discusses work related to blocked scope runtime\_module\_loader.
\item Conflict target: .aide/queue/current.toml
\item Recommended disposition: preserve\_as\_history; require current queue authority before action
\end{itemize}

\subsection{CONTRA-0210 - conversation\_vs\_current\_queue}

\begin{itemize}
\item Source conversation: Workbench\_AIDE\_Product\_Spine
\item Source file: docs/archive/conversations/Workbench\_AIDE\_Product\_Spine/Dominium\_Architecture\_Workbench\_AIDE\_Product\_Spine\_Planning\_\_01\_human\_readable\_report.md
\item Claim: Archived conversation discusses work related to blocked scope provider\_runtime.
\item Conflict target: .aide/queue/current.toml
\item Recommended disposition: preserve\_as\_history; require current queue authority before action
\end{itemize}

\subsection{CONTRA-0211 - conversation\_vs\_current\_queue}

\begin{itemize}
\item Source conversation: Workbench\_AIDE\_Product\_Spine
\item Source file: docs/archive/conversations/Workbench\_AIDE\_Product\_Spine/Dominium\_Architecture\_Workbench\_AIDE\_Product\_Spine\_Planning\_\_01\_human\_readable\_report.md
\item Claim: Archived conversation discusses work related to blocked scope package\_runtime.
\item Conflict target: .aide/queue/current.toml
\item Recommended disposition: preserve\_as\_history; require current queue authority before action
\end{itemize}

\subsection{CONTRA-0212 - conversation\_vs\_current\_queue}

\begin{itemize}
\item Source conversation: Workbench\_AIDE\_Product\_Spine
\item Source file: docs/archive/conversations/Workbench\_AIDE\_Product\_Spine/Dominium\_Architecture\_Workbench\_AIDE\_Product\_Spine\_Planning\_\_01\_human\_readable\_report.md
\item Claim: Archived conversation discusses work related to blocked scope gameplay.
\item Conflict target: .aide/queue/current.toml
\item Recommended disposition: preserve\_as\_history; require current queue authority before action
\end{itemize}

\subsection{CONTRA-0213 - conversation\_vs\_current\_queue}

\begin{itemize}
\item Source conversation: Workbench\_AIDE\_Product\_Spine
\item Source file: docs/archive/conversations/Workbench\_AIDE\_Product\_Spine/Dominium\_Architecture\_Workbench\_AIDE\_Product\_Spine\_Planning\_\_01\_human\_readable\_report.md
\item Claim: Archived conversation discusses work related to blocked scope renderer\_implementation.
\item Conflict target: .aide/queue/current.toml
\item Recommended disposition: preserve\_as\_history; require current queue authority before action
\end{itemize}

\subsection{CONTRA-0214 - conversation\_vs\_current\_queue}

\begin{itemize}
\item Source conversation: Workbench\_AIDE\_Product\_Spine
\item Source file: docs/archive/conversations/Workbench\_AIDE\_Product\_Spine/Dominium\_Architecture\_Workbench\_AIDE\_Product\_Spine\_Planning\_\_01\_human\_readable\_report.md
\item Claim: Archived conversation discusses work related to blocked scope native\_gui.
\item Conflict target: .aide/queue/current.toml
\item Recommended disposition: preserve\_as\_history; require current queue authority before action
\end{itemize}

\subsection{CONTRA-0215 - conversation\_vs\_current\_queue}

\begin{itemize}
\item Source conversation: Workbench\_AIDE\_Product\_Spine
\item Source file: docs/archive/conversations/Workbench\_AIDE\_Product\_Spine/Dominium\_Architecture\_Workbench\_AIDE\_Product\_Spine\_Planning\_\_01\_human\_readable\_report.md
\item Claim: Archived conversation discusses work related to blocked scope release\_publication.
\item Conflict target: .aide/queue/current.toml
\item Recommended disposition: preserve\_as\_history; require current queue authority before action
\end{itemize}

\subsection{CONTRA-0216 - stale\_external\_claim}

\begin{itemize}
\item Source conversation: Workbench\_AIDE\_Product\_Spine
\item Source file: docs/archive/conversations/Workbench\_AIDE\_Product\_Spine/Dominium\_Architecture\_Workbench\_AIDE\_Product\_Spine\_Planning\_\_01\_human\_readable\_report.md
\item Claim: Archived conversation contains potentially stale external or baseline claim: C89.
\item Conflict target: README.md / current platform and language baseline
\item Recommended disposition: quarantined until current repo and external facts are checked
\end{itemize}

\subsection{CONTRA-0217 - stale\_external\_claim}

\begin{itemize}
\item Source conversation: Workbench\_AIDE\_Product\_Spine
\item Source file: docs/archive/conversations/Workbench\_AIDE\_Product\_Spine/Dominium\_Architecture\_Workbench\_AIDE\_Product\_Spine\_Planning\_\_01\_human\_readable\_report.md
\item Claim: Archived conversation contains potentially stale external or baseline claim: C++98.
\item Conflict target: README.md / current platform and language baseline
\item Recommended disposition: quarantined until current repo and external facts are checked
\end{itemize}

\subsection{CONTRA-0218 - conversation\_vs\_current\_queue}

\begin{itemize}
\item Source conversation: World\_Architecture
\item Source file: docs/archive/conversations/World\_Architecture/Dominium\_World\_Architecture\_\_human\_readable\_chat\_report.txt
\item Claim: Archived conversation discusses work related to blocked scope gameplay.
\item Conflict target: .aide/queue/current.toml
\item Recommended disposition: preserve\_as\_history; require current queue authority before action
\end{itemize}

\subsection{CONTRA-0219 - stale\_external\_claim}

\begin{itemize}
\item Source conversation: World\_Architecture
\item Source file: docs/archive/conversations/World\_Architecture/Dominium\_World\_Architecture\_\_human\_readable\_chat\_report.txt
\item Claim: Archived conversation contains potentially stale external or baseline claim: C89.
\item Conflict target: README.md / current platform and language baseline
\item Recommended disposition: quarantined until current repo and external facts are checked
\end{itemize}

\subsection{CONTRA-0220 - stale\_external\_claim}

\begin{itemize}
\item Source conversation: World\_Architecture
\item Source file: docs/archive/conversations/World\_Architecture/Dominium\_World\_Architecture\_\_human\_readable\_chat\_report.txt
\item Claim: Archived conversation contains potentially stale external or baseline claim: C++98.
\item Conflict target: README.md / current platform and language baseline
\item Recommended disposition: quarantined until current repo and external facts are checked
\end{itemize}

\subsection{CONTRA-0221 - stale\_external\_claim}

\begin{itemize}
\item Source conversation: World\_Architecture
\item Source file: docs/archive/conversations/World\_Architecture/Dominium\_World\_Architecture\_\_human\_readable\_chat\_report.txt
\item Claim: Archived conversation contains potentially stale external or baseline claim: ISO C89.
\item Conflict target: README.md / current platform and language baseline
\item Recommended disposition: quarantined until current repo and external facts are checked
\end{itemize}

\subsection{CONTRA-0222 - stale\_external\_claim}

\begin{itemize}
\item Source conversation: World\_Architecture
\item Source file: docs/archive/conversations/World\_Architecture/Dominium\_World\_Architecture\_\_human\_readable\_chat\_report.txt
\item Claim: Archived conversation contains potentially stale external or baseline claim: SDK.
\item Conflict target: README.md / current platform and language baseline
\item Recommended disposition: quarantined until current repo and external facts are checked
\end{itemize}

\subsection{CONTRA-0223 - conversation\_vs\_current\_queue}

\begin{itemize}
\item Source conversation: xstack\_lab\_galaxy
\item Source file: docs/archive/conversations/xstack\_lab\_galaxy/dominium\_xstack\_lab\_galaxy\_handoff\_\_human\_readable\_chat\_report.txt
\item Claim: Archived conversation discusses work related to blocked scope gameplay.
\item Conflict target: .aide/queue/current.toml
\item Recommended disposition: preserve\_as\_history; require current queue authority before action
\end{itemize}

\subsection{CONTRA-0224 - conversation\_vs\_docs}

\begin{itemize}
\item Source conversation: advanced\_simulation\_infrastructure, app\_testx\_codehygiene, architecture\_codex\_prompts, architecture\_ui\_providers, Build\_and\_Future\_Proofing, canonical\_structure\_and\_framework, documentation\_standards\_readme, Dominium\_Architecture\_I
\item Source file: multiple primary sources
\item Claim: Corpus contains both C89 and C17 claims.
\item Conflict target: current repo authority and conversation cross-check
\item Recommended disposition: quarantined until reviewed
\end{itemize}

\subsection{CONTRA-0225 - conversation\_vs\_docs}

\begin{itemize}
\item Source conversation: advanced\_simulation\_infrastructure, app\_testx\_codehygiene, architecture\_codex\_prompts, architecture\_ui\_providers, Build\_and\_Future\_Proofing, canonical\_structure\_and\_framework, documentation\_standards\_readme, Dominium\_Architecture\_I
\item Source file: multiple primary sources
\item Claim: Corpus contains both C++98 and C++17 claims.
\item Conflict target: current repo authority and conversation cross-check
\item Recommended disposition: quarantined until reviewed
\end{itemize}



{\small\raggedright SOURCE PATH: docs/\allowbreak{}archive/\allowbreak{}conversations/\allowbreak{}\_\allowbreak{}audit/\allowbreak{}COVERAGE\_\allowbreak{}GAPS.md\par}


\section{Coverage Gaps}

Coverage gaps identify incomplete or structurally unclear archived packages.


\begin{quote}\small Dense table content is summarized in this PDF rendering. See the Markdown source and HTML book for the full table.\end{quote}



{\small\raggedright SOURCE PATH: docs/\allowbreak{}archive/\allowbreak{}conversations/\allowbreak{}\_\allowbreak{}audit/\allowbreak{}INTAKE\_\allowbreak{}ACCEPTANCE\_\allowbreak{}REVIEW.md\par}


\section{Intake Acceptance Review}

Result: PASS\_WITH\_WARNINGS


This review assesses whether the generated conversation-corpus intake is complete and useful enough to support derived synthesis. It does not promote conversation claims.


\subsection{Summary}

\begin{quote}\small Dense table content is summarized in this PDF rendering. See the Markdown source and HTML book for the full table.\end{quote}

\subsection{Acceptance Findings}

\begin{itemize}
\item All top-level source packages represented: True.
\item All source files hashed: True.
\item Manifest re-scan matches committed manifest: True.
\item Reader pages are present and advisory-scoped: True.
\item Wiki topic pages cover manifest topics: True.
\item Link integrity is clean for generated Markdown: True.
\item Promotion queue is useful as raw intake, but not clean enough for direct promotion: True.
\item Contradiction findings are actionable as a review backlog, but remain heuristic and need triage before use as doctrine evidence.
\end{itemize}

\subsection{Warnings Before Synthesis}

\begin{itemize}
\item Noisy or archival-process promotion candidates: 17.
\item Overlong promotion candidates: 44.
\item Promotion candidates with repo conflict still not\_checked: 135.
\item Reader placeholder counts: \{'no\_stable\_decision\_lines': 0, 'no\_explicit\_uncertainty\_lines': 0, 'no\_explicit\_future\_work\_lines': 0, 'no\_rejected\_or\_superseded\_lines': 1\}.
\item Contradiction findings are broad and keyword-assisted; use them as review triggers, not resolved conclusions.
\end{itemize}

\subsection{Decision}

PASS\_WITH\_WARNINGS


The corpus is ready for derived synthesis if synthesis cites reader/wiki/audit sources directly, keeps repo truth separate from conversation-derived intent, and treats the promotion queue as raw review material. It is not ready for direct live-doc promotion.


\subsection{Recommended Fixes Before Promotion}

\begin{itemize}
\item Triage promotion candidates into serious, historical, stale, noisy, and needs-user-decision buckets.
\item Reconcile high-value claims against canon, current queue, contracts, schema law, and implementation only in a separate task.
\item Keep the current acceptance result visible in the synthesis book front matter.
\end{itemize}

\subsection{Recommended Next Task}

CONVERSATION-SYNTHESIS-BOOK-01



\chapter{Wiki and Navigation Digest}


{\small\raggedright SOURCE PATH: docs/\allowbreak{}archive/\allowbreak{}conversations/\allowbreak{}\_\allowbreak{}wiki/\allowbreak{}index.md\par}


\section{Conversation Corpus Wiki}

This repo-local wiki is a navigation surface over the archived conversation corpus. It is not a GitHub Wiki and it is not authoritative doctrine.


\subsection{Browse}

\begin{itemize}
\item [Topics](topics/index.md)
\item [Workstreams](workstreams/index.md)
\item [Decisions](decisions/index.md)
\item [Tasks](tasks/index.md)
\item [Artifacts](artifacts/index.md)
\item [Open Questions](open\_questions/index.md)
\item [Risks](risks/index.md)
\end{itemize}

\subsection{Scope}

\begin{itemize}
\item Source conversations: 45
\item Authority class: advisory\_only
\item Promotion status: not\_reviewed
\end{itemize}

Use \_promotion/PROMOTION\_QUEUE.md before proposing any live-doc update.




{\small\raggedright SOURCE PATH: docs/\allowbreak{}archive/\allowbreak{}conversations/\allowbreak{}\_\allowbreak{}wiki/\allowbreak{}topics/\allowbreak{}index.md\par}


\section{Topic Index}

\begin{itemize}
\item [architecture](architecture.md): 45 conversations
\item [content](content.md): 45 conversations
\item [contracts\_schema](contracts\_schema.md): 42 conversations
\item [determinism](determinism.md): 39 conversations
\item [governance](governance.md): 44 conversations
\item [platform](platform.md): 39 conversations
\item [release](release.md): 43 conversations
\item [setup\_launcher](setup\_launcher.md): 25 conversations
\item [simulation](simulation.md): 44 conversations
\item [timekeeping](timekeeping.md): 35 conversations
\item [tooling](tooling.md): 44 conversations
\item [ui](ui.md): 45 conversations
\item [workbench](workbench.md): 20 conversations
\item [worldgen](worldgen.md): 39 conversations
\item [xstack\_aide](xstack\_aide.md): 14 conversations
\end{itemize}



{\small\raggedright SOURCE PATH: docs/\allowbreak{}archive/\allowbreak{}conversations/\allowbreak{}\_\allowbreak{}wiki/\allowbreak{}workstreams/\allowbreak{}index.md\par}


\section{Workstreams}

\subsection{architecture}

\begin{itemize}
\item [Dominium Advanced Simulation and Infrastructure Architecture](../../\_reader/by\_chat/advanced\_simulation\_infrastructure.md)
\item [Dominium APP0 Runtime, Platform, and Renderer Architecture](../../\_reader/by\_chat/app\_runtime\_platform\_renderers.md)
\item [Dominium Architecture, Application Layer, TestX, and CodeHygiene Planning](../../\_reader/by\_chat/app\_testx\_codehygiene.md)
\item [Dominium/Domino Architecture and Codex Prompt Roadmap](../../\_reader/by\_chat/architecture\_codex\_prompts.md)
\item [Dominium Architecture, UI, Providers, and Robot OS Strategy](../../\_reader/by\_chat/architecture\_ui\_providers.md)
\item [Dominium Build and Future-Proofing Architecture](../../\_reader/by\_chat/build\_and\_future\_proofing.md)
\item [Dominium Canonical Architecture, Repository Foundation, and Provider Model](../../\_reader/by\_chat/canonical\_structure\_and\_framework.md)
\item [Dominium Chronology \& Celestial Systems](../../\_reader/by\_chat/chronology\_celestial\_systems.md)
\item [Dominium Development Routes and Continuity Preservation](../../\_reader/by\_chat/development\_routes.md)
\item [Documentation Standards, README Strategy, and Handoff Packaging](../../\_reader/by\_chat/documentation\_standards\_readme.md)
\item [Dominium Architecture I](../../\_reader/by\_chat/dominium\_architecture\_i.md)
\item [Dominium Architecture II](../../\_reader/by\_chat/dominium\_architecture\_ii.md)
\item [Dominium Architecture III: Launcher, Platform, Renderer, and Handoff Architecture](../../\_reader/by\_chat/dominium\_architecture\_iii.md)
\item [Dominium Architecture IV](../../\_reader/by\_chat/dominium\_architecture\_iv.md)
\item [Dominium Canon, Repository Alignment, and Portability Doctrine](../../\_reader/by\_chat/dominium\_complete\_conversation.md)
\item [Dominium + Domino Codex Planning and Handoff](../../\_reader/by\_chat/dominium\_domino\_codex\_planning.md)
\item [Dominium Workbench, AIDE, Presentation Architecture, Provider Strategy, Product-Spine Planning, and Preservation Companion](../../\_reader/by\_chat/dominium\_full\_conversation.md)
\item [Dominium Setup Architecture and Handoff](../../\_reader/by\_chat/dominium\_setup.md)
\item [Domino Dominium Workbench](../../\_reader/by\_chat/domino\_dominium\_workbench.md)
\item [Dominium/Domino Engine Refactor Planning](../../\_reader/by\_chat/domino\_engine\_refactor\_prompts.md)
\item [Domino/Dominium Engine Baseline, Architecture, and Feasibility](../../\_reader/by\_chat/engine\_baseline\_architecture.md)
\item [Dominium Foundation Lock, Workbench Spine, and Parallel Codex Handoff](../../\_reader/by\_chat/foundation\_workbench\_codex.md)
\item [Domino Framework and Open-Source Provider Architecture](../../\_reader/by\_chat/framework\_open\_source\_provider.md)
\item [Dominium Content and GUI Rebuild Planning](../../\_reader/by\_chat/gui\_binary\_content.md)
\item [Dominium Language, Platform, and Architecture Baseline](../../\_reader/by\_chat/language\_platform\_architecture.md)
\item [Dominium Launcher Application-Layer Handoff](../../\_reader/by\_chat/launcher\_app\_layer.md)
\item [Dominium Launcher and Setup Architecture](../../\_reader/by\_chat/launcher\_setup\_architecture.md)
\item [Dominium Modularity, AIDE Refactorability, and Future-Proof Architecture](../../\_reader/by\_chat/modularity\_aide\_refactorability.md)
\item [Dominium Omega/Xi/Pi Architecture \& Future-Proofing Planning](../../\_reader/by\_chat/omega\_xi\_pi\_architecture\_future.md)
\item [Dominium OS-like Architecture, Repository Convergence, and Interface Operating Layer](../../\_reader/by\_chat/os\_interface\_repo\_architecture.md)
\item [Dominium Codex Platform Renderer API Plan](../../\_reader/by\_chat/platform\_renderer\_api\_plan.md)
\item [Dominium Platform Support Planning](../../\_reader/by\_chat/platform\_support.md)
\item [Domino/Dominium Portability, Assurance, and Future-Proof Architecture](../../\_reader/by\_chat/portability\_assurance\_future\_proof.md)
\item [Dominium README Architecture, Ports, and Determinism](../../\_reader/by\_chat/readme\_ports\_determinism.md)
\item [Dominium + Domino Refactor Architecture](../../\_reader/by\_chat/refactor\_architecture.md)
\item [AIDE, XStack, and Dominium Refactor Control Plane](../../\_reader/by\_chat/refactor\_control\_plane.md)
\item [Dominium XStack Release Identity and Versioning](../../\_reader/by\_chat/release\_identity\_and\_versioning.md)
\item [Dominium TestX/RepoX Governance and Handoff Chat](../../\_reader/by\_chat/testx\_repox\_governance.md)
\item [Dominium Timekeeping and 2038 Resilience](../../\_reader/by\_chat/timekeeping\_and\_2038\_resilience.md)
\item [Dominium UE6, Domino, and Deterministic Universe Feasibility](../../\_reader/by\_chat/ue6\_domino\_deterministic\_universe.md)
\item [Dominium UI Editor and Tool Editor Planning](../../\_reader/by\_chat/ui\_editor\_tool\_editor\_planning.md)
\item [Dominium Universe Explorer Planning and Repo Handoff](../../\_reader/by\_chat/universe\_explorer\_planning.md)
\item [Dominium Architecture, Workbench, AIDE, and Product-Spine Planning](../../\_reader/by\_chat/workbench\_aide\_product\_spine.md)
\item [Dominium World Architecture](../../\_reader/by\_chat/world\_architecture.md)
\item [Dominium XStack and Lab Galaxy Handoff](../../\_reader/by\_chat/xstack\_lab\_galaxy.md)
\end{itemize}

\subsection{content}

\begin{itemize}
\item [Dominium Advanced Simulation and Infrastructure Architecture](../../\_reader/by\_chat/advanced\_simulation\_infrastructure.md)
\item [Dominium APP0 Runtime, Platform, and Renderer Architecture](../../\_reader/by\_chat/app\_runtime\_platform\_renderers.md)
\item [Dominium Architecture, Application Layer, TestX, and CodeHygiene Planning](../../\_reader/by\_chat/app\_testx\_codehygiene.md)
\item [Dominium/Domino Architecture and Codex Prompt Roadmap](../../\_reader/by\_chat/architecture\_codex\_prompts.md)
\item [Dominium Architecture, UI, Providers, and Robot OS Strategy](../../\_reader/by\_chat/architecture\_ui\_providers.md)
\item [Dominium Build and Future-Proofing Architecture](../../\_reader/by\_chat/build\_and\_future\_proofing.md)
\item [Dominium Canonical Architecture, Repository Foundation, and Provider Model](../../\_reader/by\_chat/canonical\_structure\_and\_framework.md)
\item [Dominium Chronology \& Celestial Systems](../../\_reader/by\_chat/chronology\_celestial\_systems.md)
\item [Dominium Development Routes and Continuity Preservation](../../\_reader/by\_chat/development\_routes.md)
\item [Documentation Standards, README Strategy, and Handoff Packaging](../../\_reader/by\_chat/documentation\_standards\_readme.md)
\item [Dominium Architecture I](../../\_reader/by\_chat/dominium\_architecture\_i.md)
\item [Dominium Architecture II](../../\_reader/by\_chat/dominium\_architecture\_ii.md)
\item [Dominium Architecture III: Launcher, Platform, Renderer, and Handoff Architecture](../../\_reader/by\_chat/dominium\_architecture\_iii.md)
\item [Dominium Architecture IV](../../\_reader/by\_chat/dominium\_architecture\_iv.md)
\item [Dominium Canon, Repository Alignment, and Portability Doctrine](../../\_reader/by\_chat/dominium\_complete\_conversation.md)
\item [Dominium + Domino Codex Planning and Handoff](../../\_reader/by\_chat/dominium\_domino\_codex\_planning.md)
\item [Dominium Workbench, AIDE, Presentation Architecture, Provider Strategy, Product-Spine Planning, and Preservation Companion](../../\_reader/by\_chat/dominium\_full\_conversation.md)
\item [Dominium Setup Architecture and Handoff](../../\_reader/by\_chat/dominium\_setup.md)
\item [Domino Dominium Workbench](../../\_reader/by\_chat/domino\_dominium\_workbench.md)
\item [Dominium/Domino Engine Refactor Planning](../../\_reader/by\_chat/domino\_engine\_refactor\_prompts.md)
\item [Domino/Dominium Engine Baseline, Architecture, and Feasibility](../../\_reader/by\_chat/engine\_baseline\_architecture.md)
\item [Dominium Foundation Lock, Workbench Spine, and Parallel Codex Handoff](../../\_reader/by\_chat/foundation\_workbench\_codex.md)
\item [Domino Framework and Open-Source Provider Architecture](../../\_reader/by\_chat/framework\_open\_source\_provider.md)
\item [Dominium Content and GUI Rebuild Planning](../../\_reader/by\_chat/gui\_binary\_content.md)
\item [Dominium Language, Platform, and Architecture Baseline](../../\_reader/by\_chat/language\_platform\_architecture.md)
\item [Dominium Launcher Application-Layer Handoff](../../\_reader/by\_chat/launcher\_app\_layer.md)
\item [Dominium Launcher and Setup Architecture](../../\_reader/by\_chat/launcher\_setup\_architecture.md)
\item [Dominium Modularity, AIDE Refactorability, and Future-Proof Architecture](../../\_reader/by\_chat/modularity\_aide\_refactorability.md)
\item [Dominium Omega/Xi/Pi Architecture \& Future-Proofing Planning](../../\_reader/by\_chat/omega\_xi\_pi\_architecture\_future.md)
\item [Dominium OS-like Architecture, Repository Convergence, and Interface Operating Layer](../../\_reader/by\_chat/os\_interface\_repo\_architecture.md)
\item [Dominium Codex Platform Renderer API Plan](../../\_reader/by\_chat/platform\_renderer\_api\_plan.md)
\item [Dominium Platform Support Planning](../../\_reader/by\_chat/platform\_support.md)
\item [Domino/Dominium Portability, Assurance, and Future-Proof Architecture](../../\_reader/by\_chat/portability\_assurance\_future\_proof.md)
\item [Dominium README Architecture, Ports, and Determinism](../../\_reader/by\_chat/readme\_ports\_determinism.md)
\item [Dominium + Domino Refactor Architecture](../../\_reader/by\_chat/refactor\_architecture.md)
\item [AIDE, XStack, and Dominium Refactor Control Plane](../../\_reader/by\_chat/refactor\_control\_plane.md)
\item [Dominium XStack Release Identity and Versioning](../../\_reader/by\_chat/release\_identity\_and\_versioning.md)
\item [Dominium TestX/RepoX Governance and Handoff Chat](../../\_reader/by\_chat/testx\_repox\_governance.md)
\item [Dominium Timekeeping and 2038 Resilience](../../\_reader/by\_chat/timekeeping\_and\_2038\_resilience.md)
\item [Dominium UE6, Domino, and Deterministic Universe Feasibility](../../\_reader/by\_chat/ue6\_domino\_deterministic\_universe.md)
\item [Dominium UI Editor and Tool Editor Planning](../../\_reader/by\_chat/ui\_editor\_tool\_editor\_planning.md)
\item [Dominium Universe Explorer Planning and Repo Handoff](../../\_reader/by\_chat/universe\_explorer\_planning.md)
\item [Dominium Architecture, Workbench, AIDE, and Product-Spine Planning](../../\_reader/by\_chat/workbench\_aide\_product\_spine.md)
\item [Dominium World Architecture](../../\_reader/by\_chat/world\_architecture.md)
\item [Dominium XStack and Lab Galaxy Handoff](../../\_reader/by\_chat/xstack\_lab\_galaxy.md)
\end{itemize}

\subsection{contracts\_schema}

\begin{itemize}
\item [Dominium Advanced Simulation and Infrastructure Architecture](../../\_reader/by\_chat/advanced\_simulation\_infrastructure.md)
\item [Dominium APP0 Runtime, Platform, and Renderer Architecture](../../\_reader/by\_chat/app\_runtime\_platform\_renderers.md)
\item [Dominium Architecture, Application Layer, TestX, and CodeHygiene Planning](../../\_reader/by\_chat/app\_testx\_codehygiene.md)
\item [Dominium/Domino Architecture and Codex Prompt Roadmap](../../\_reader/by\_chat/architecture\_codex\_prompts.md)
\item [Dominium Architecture, UI, Providers, and Robot OS Strategy](../../\_reader/by\_chat/architecture\_ui\_providers.md)
\item [Dominium Build and Future-Proofing Architecture](../../\_reader/by\_chat/build\_and\_future\_proofing.md)
\item [Dominium Canonical Architecture, Repository Foundation, and Provider Model](../../\_reader/by\_chat/canonical\_structure\_and\_framework.md)
\item [Dominium Chronology \& Celestial Systems](../../\_reader/by\_chat/chronology\_celestial\_systems.md)
\item [Dominium Development Routes and Continuity Preservation](../../\_reader/by\_chat/development\_routes.md)
\item [Documentation Standards, README Strategy, and Handoff Packaging](../../\_reader/by\_chat/documentation\_standards\_readme.md)
\item [Dominium Architecture I](../../\_reader/by\_chat/dominium\_architecture\_i.md)
\item [Dominium Architecture II](../../\_reader/by\_chat/dominium\_architecture\_ii.md)
\item [Dominium Architecture IV](../../\_reader/by\_chat/dominium\_architecture\_iv.md)
\item [Dominium Canon, Repository Alignment, and Portability Doctrine](../../\_reader/by\_chat/dominium\_complete\_conversation.md)
\item [Dominium + Domino Codex Planning and Handoff](../../\_reader/by\_chat/dominium\_domino\_codex\_planning.md)
\item [Dominium Workbench, AIDE, Presentation Architecture, Provider Strategy, Product-Spine Planning, and Preservation Companion](../../\_reader/by\_chat/dominium\_full\_conversation.md)
\item [Dominium Setup Architecture and Handoff](../../\_reader/by\_chat/dominium\_setup.md)
\item [Domino Dominium Workbench](../../\_reader/by\_chat/domino\_dominium\_workbench.md)
\item [Dominium/Domino Engine Refactor Planning](../../\_reader/by\_chat/domino\_engine\_refactor\_prompts.md)
\item [Domino/Dominium Engine Baseline, Architecture, and Feasibility](../../\_reader/by\_chat/engine\_baseline\_architecture.md)
\item [Dominium Foundation Lock, Workbench Spine, and Parallel Codex Handoff](../../\_reader/by\_chat/foundation\_workbench\_codex.md)
\item [Domino Framework and Open-Source Provider Architecture](../../\_reader/by\_chat/framework\_open\_source\_provider.md)
\item [Dominium Content and GUI Rebuild Planning](../../\_reader/by\_chat/gui\_binary\_content.md)
\item [Dominium Language, Platform, and Architecture Baseline](../../\_reader/by\_chat/language\_platform\_architecture.md)
\item [Dominium Launcher Application-Layer Handoff](../../\_reader/by\_chat/launcher\_app\_layer.md)
\item [Dominium Launcher and Setup Architecture](../../\_reader/by\_chat/launcher\_setup\_architecture.md)
\item [Dominium Modularity, AIDE Refactorability, and Future-Proof Architecture](../../\_reader/by\_chat/modularity\_aide\_refactorability.md)
\item [Dominium Omega/Xi/Pi Architecture \& Future-Proofing Planning](../../\_reader/by\_chat/omega\_xi\_pi\_architecture\_future.md)
\item [Dominium OS-like Architecture, Repository Convergence, and Interface Operating Layer](../../\_reader/by\_chat/os\_interface\_repo\_architecture.md)
\item [Dominium Codex Platform Renderer API Plan](../../\_reader/by\_chat/platform\_renderer\_api\_plan.md)
\item [Dominium Platform Support Planning](../../\_reader/by\_chat/platform\_support.md)
\item [Domino/Dominium Portability, Assurance, and Future-Proof Architecture](../../\_reader/by\_chat/portability\_assurance\_future\_proof.md)
\item [Dominium README Architecture, Ports, and Determinism](../../\_reader/by\_chat/readme\_ports\_determinism.md)
\item [Dominium + Domino Refactor Architecture](../../\_reader/by\_chat/refactor\_architecture.md)
\item [AIDE, XStack, and Dominium Refactor Control Plane](../../\_reader/by\_chat/refactor\_control\_plane.md)
\item [Dominium XStack Release Identity and Versioning](../../\_reader/by\_chat/release\_identity\_and\_versioning.md)
\item [Dominium Timekeeping and 2038 Resilience](../../\_reader/by\_chat/timekeeping\_and\_2038\_resilience.md)
\item [Dominium UI Editor and Tool Editor Planning](../../\_reader/by\_chat/ui\_editor\_tool\_editor\_planning.md)
\item [Dominium Universe Explorer Planning and Repo Handoff](../../\_reader/by\_chat/universe\_explorer\_planning.md)
\item [Dominium Architecture, Workbench, AIDE, and Product-Spine Planning](../../\_reader/by\_chat/workbench\_aide\_product\_spine.md)
\item [Dominium World Architecture](../../\_reader/by\_chat/world\_architecture.md)
\item [Dominium XStack and Lab Galaxy Handoff](../../\_reader/by\_chat/xstack\_lab\_galaxy.md)
\end{itemize}

\subsection{determinism}

\begin{itemize}
\item [Dominium Advanced Simulation and Infrastructure Architecture](../../\_reader/by\_chat/advanced\_simulation\_infrastructure.md)
\item [Dominium APP0 Runtime, Platform, and Renderer Architecture](../../\_reader/by\_chat/app\_runtime\_platform\_renderers.md)
\item [Dominium Architecture, Application Layer, TestX, and CodeHygiene Planning](../../\_reader/by\_chat/app\_testx\_codehygiene.md)
\item [Dominium/Domino Architecture and Codex Prompt Roadmap](../../\_reader/by\_chat/architecture\_codex\_prompts.md)
\item [Dominium Architecture, UI, Providers, and Robot OS Strategy](../../\_reader/by\_chat/architecture\_ui\_providers.md)
\item [Dominium Build and Future-Proofing Architecture](../../\_reader/by\_chat/build\_and\_future\_proofing.md)
\item [Dominium Canonical Architecture, Repository Foundation, and Provider Model](../../\_reader/by\_chat/canonical\_structure\_and\_framework.md)
\item [Dominium Development Routes and Continuity Preservation](../../\_reader/by\_chat/development\_routes.md)
\item [Documentation Standards, README Strategy, and Handoff Packaging](../../\_reader/by\_chat/documentation\_standards\_readme.md)
\item [Dominium Architecture I](../../\_reader/by\_chat/dominium\_architecture\_i.md)
\item [Dominium Architecture II](../../\_reader/by\_chat/dominium\_architecture\_ii.md)
\item [Dominium Architecture III: Launcher, Platform, Renderer, and Handoff Architecture](../../\_reader/by\_chat/dominium\_architecture\_iii.md)
\item [Dominium Architecture IV](../../\_reader/by\_chat/dominium\_architecture\_iv.md)
\item [Dominium Canon, Repository Alignment, and Portability Doctrine](../../\_reader/by\_chat/dominium\_complete\_conversation.md)
\item [Dominium + Domino Codex Planning and Handoff](../../\_reader/by\_chat/dominium\_domino\_codex\_planning.md)
\item [Dominium Workbench, AIDE, Presentation Architecture, Provider Strategy, Product-Spine Planning, and Preservation Companion](../../\_reader/by\_chat/dominium\_full\_conversation.md)
\item [Dominium Setup Architecture and Handoff](../../\_reader/by\_chat/dominium\_setup.md)
\item [Domino Dominium Workbench](../../\_reader/by\_chat/domino\_dominium\_workbench.md)
\item [Dominium/Domino Engine Refactor Planning](../../\_reader/by\_chat/domino\_engine\_refactor\_prompts.md)
\item [Domino/Dominium Engine Baseline, Architecture, and Feasibility](../../\_reader/by\_chat/engine\_baseline\_architecture.md)
\item [Dominium Foundation Lock, Workbench Spine, and Parallel Codex Handoff](../../\_reader/by\_chat/foundation\_workbench\_codex.md)
\item [Domino Framework and Open-Source Provider Architecture](../../\_reader/by\_chat/framework\_open\_source\_provider.md)
\item [Dominium Content and GUI Rebuild Planning](../../\_reader/by\_chat/gui\_binary\_content.md)
\item [Dominium Language, Platform, and Architecture Baseline](../../\_reader/by\_chat/language\_platform\_architecture.md)
\item [Dominium Launcher Application-Layer Handoff](../../\_reader/by\_chat/launcher\_app\_layer.md)
\item [Dominium Launcher and Setup Architecture](../../\_reader/by\_chat/launcher\_setup\_architecture.md)
\item [Dominium Modularity, AIDE Refactorability, and Future-Proof Architecture](../../\_reader/by\_chat/modularity\_aide\_refactorability.md)
\item [Dominium OS-like Architecture, Repository Convergence, and Interface Operating Layer](../../\_reader/by\_chat/os\_interface\_repo\_architecture.md)
\item [Dominium Codex Platform Renderer API Plan](../../\_reader/by\_chat/platform\_renderer\_api\_plan.md)
\item [Domino/Dominium Portability, Assurance, and Future-Proof Architecture](../../\_reader/by\_chat/portability\_assurance\_future\_proof.md)
\item [Dominium README Architecture, Ports, and Determinism](../../\_reader/by\_chat/readme\_ports\_determinism.md)
\item [Dominium XStack Release Identity and Versioning](../../\_reader/by\_chat/release\_identity\_and\_versioning.md)
\item [Dominium TestX/RepoX Governance and Handoff Chat](../../\_reader/by\_chat/testx\_repox\_governance.md)
\item [Dominium Timekeeping and 2038 Resilience](../../\_reader/by\_chat/timekeeping\_and\_2038\_resilience.md)
\item [Dominium UE6, Domino, and Deterministic Universe Feasibility](../../\_reader/by\_chat/ue6\_domino\_deterministic\_universe.md)
\item [Dominium UI Editor and Tool Editor Planning](../../\_reader/by\_chat/ui\_editor\_tool\_editor\_planning.md)
\item [Dominium Architecture, Workbench, AIDE, and Product-Spine Planning](../../\_reader/by\_chat/workbench\_aide\_product\_spine.md)
\item [Dominium World Architecture](../../\_reader/by\_chat/world\_architecture.md)
\item [Dominium XStack and Lab Galaxy Handoff](../../\_reader/by\_chat/xstack\_lab\_galaxy.md)
\end{itemize}

\subsection{governance}

\begin{itemize}
\item [Dominium Advanced Simulation and Infrastructure Architecture](../../\_reader/by\_chat/advanced\_simulation\_infrastructure.md)
\item [Dominium APP0 Runtime, Platform, and Renderer Architecture](../../\_reader/by\_chat/app\_runtime\_platform\_renderers.md)
\item [Dominium Architecture, Application Layer, TestX, and CodeHygiene Planning](../../\_reader/by\_chat/app\_testx\_codehygiene.md)
\item [Dominium/Domino Architecture and Codex Prompt Roadmap](../../\_reader/by\_chat/architecture\_codex\_prompts.md)
\item [Dominium Architecture, UI, Providers, and Robot OS Strategy](../../\_reader/by\_chat/architecture\_ui\_providers.md)
\item [Dominium Build and Future-Proofing Architecture](../../\_reader/by\_chat/build\_and\_future\_proofing.md)
\item [Dominium Canonical Architecture, Repository Foundation, and Provider Model](../../\_reader/by\_chat/canonical\_structure\_and\_framework.md)
\item [Dominium Chronology \& Celestial Systems](../../\_reader/by\_chat/chronology\_celestial\_systems.md)
\item [Dominium Development Routes and Continuity Preservation](../../\_reader/by\_chat/development\_routes.md)
\item [Documentation Standards, README Strategy, and Handoff Packaging](../../\_reader/by\_chat/documentation\_standards\_readme.md)
\item [Dominium Architecture I](../../\_reader/by\_chat/dominium\_architecture\_i.md)
\item [Dominium Architecture II](../../\_reader/by\_chat/dominium\_architecture\_ii.md)
\item [Dominium Architecture III: Launcher, Platform, Renderer, and Handoff Architecture](../../\_reader/by\_chat/dominium\_architecture\_iii.md)
\item [Dominium Architecture IV](../../\_reader/by\_chat/dominium\_architecture\_iv.md)
\item [Dominium Canon, Repository Alignment, and Portability Doctrine](../../\_reader/by\_chat/dominium\_complete\_conversation.md)
\item [Dominium + Domino Codex Planning and Handoff](../../\_reader/by\_chat/dominium\_domino\_codex\_planning.md)
\item [Dominium Workbench, AIDE, Presentation Architecture, Provider Strategy, Product-Spine Planning, and Preservation Companion](../../\_reader/by\_chat/dominium\_full\_conversation.md)
\item [Dominium Setup Architecture and Handoff](../../\_reader/by\_chat/dominium\_setup.md)
\item [Domino Dominium Workbench](../../\_reader/by\_chat/domino\_dominium\_workbench.md)
\item [Dominium/Domino Engine Refactor Planning](../../\_reader/by\_chat/domino\_engine\_refactor\_prompts.md)
\item [Domino/Dominium Engine Baseline, Architecture, and Feasibility](../../\_reader/by\_chat/engine\_baseline\_architecture.md)
\item [Dominium Foundation Lock, Workbench Spine, and Parallel Codex Handoff](../../\_reader/by\_chat/foundation\_workbench\_codex.md)
\item [Domino Framework and Open-Source Provider Architecture](../../\_reader/by\_chat/framework\_open\_source\_provider.md)
\item [Dominium Content and GUI Rebuild Planning](../../\_reader/by\_chat/gui\_binary\_content.md)
\item [Dominium Language, Platform, and Architecture Baseline](../../\_reader/by\_chat/language\_platform\_architecture.md)
\item [Dominium Launcher Application-Layer Handoff](../../\_reader/by\_chat/launcher\_app\_layer.md)
\item [Dominium Launcher and Setup Architecture](../../\_reader/by\_chat/launcher\_setup\_architecture.md)
\item [Dominium Modularity, AIDE Refactorability, and Future-Proof Architecture](../../\_reader/by\_chat/modularity\_aide\_refactorability.md)
\item [Dominium Omega/Xi/Pi Architecture \& Future-Proofing Planning](../../\_reader/by\_chat/omega\_xi\_pi\_architecture\_future.md)
\item [Dominium OS-like Architecture, Repository Convergence, and Interface Operating Layer](../../\_reader/by\_chat/os\_interface\_repo\_architecture.md)
\item [Dominium Codex Platform Renderer API Plan](../../\_reader/by\_chat/platform\_renderer\_api\_plan.md)
\item [Dominium Platform Support Planning](../../\_reader/by\_chat/platform\_support.md)
\item [Domino/Dominium Portability, Assurance, and Future-Proof Architecture](../../\_reader/by\_chat/portability\_assurance\_future\_proof.md)
\item [Dominium README Architecture, Ports, and Determinism](../../\_reader/by\_chat/readme\_ports\_determinism.md)
\item [AIDE, XStack, and Dominium Refactor Control Plane](../../\_reader/by\_chat/refactor\_control\_plane.md)
\item [Dominium XStack Release Identity and Versioning](../../\_reader/by\_chat/release\_identity\_and\_versioning.md)
\item [Dominium TestX/RepoX Governance and Handoff Chat](../../\_reader/by\_chat/testx\_repox\_governance.md)
\item [Dominium Timekeeping and 2038 Resilience](../../\_reader/by\_chat/timekeeping\_and\_2038\_resilience.md)
\item [Dominium UE6, Domino, and Deterministic Universe Feasibility](../../\_reader/by\_chat/ue6\_domino\_deterministic\_universe.md)
\item [Dominium UI Editor and Tool Editor Planning](../../\_reader/by\_chat/ui\_editor\_tool\_editor\_planning.md)
\item [Dominium Universe Explorer Planning and Repo Handoff](../../\_reader/by\_chat/universe\_explorer\_planning.md)
\item [Dominium Architecture, Workbench, AIDE, and Product-Spine Planning](../../\_reader/by\_chat/workbench\_aide\_product\_spine.md)
\item [Dominium World Architecture](../../\_reader/by\_chat/world\_architecture.md)
\item [Dominium XStack and Lab Galaxy Handoff](../../\_reader/by\_chat/xstack\_lab\_galaxy.md)
\end{itemize}

\subsection{platform}

\begin{itemize}
\item [Dominium Advanced Simulation and Infrastructure Architecture](../../\_reader/by\_chat/advanced\_simulation\_infrastructure.md)
\item [Dominium APP0 Runtime, Platform, and Renderer Architecture](../../\_reader/by\_chat/app\_runtime\_platform\_renderers.md)
\item [Dominium Architecture, Application Layer, TestX, and CodeHygiene Planning](../../\_reader/by\_chat/app\_testx\_codehygiene.md)
\item [Dominium Architecture, UI, Providers, and Robot OS Strategy](../../\_reader/by\_chat/architecture\_ui\_providers.md)
\item [Dominium Build and Future-Proofing Architecture](../../\_reader/by\_chat/build\_and\_future\_proofing.md)
\item [Dominium Canonical Architecture, Repository Foundation, and Provider Model](../../\_reader/by\_chat/canonical\_structure\_and\_framework.md)
\item [Dominium Chronology \& Celestial Systems](../../\_reader/by\_chat/chronology\_celestial\_systems.md)
\item [Documentation Standards, README Strategy, and Handoff Packaging](../../\_reader/by\_chat/documentation\_standards\_readme.md)
\item [Dominium Architecture I](../../\_reader/by\_chat/dominium\_architecture\_i.md)
\item [Dominium Architecture II](../../\_reader/by\_chat/dominium\_architecture\_ii.md)
\item [Dominium Architecture III: Launcher, Platform, Renderer, and Handoff Architecture](../../\_reader/by\_chat/dominium\_architecture\_iii.md)
\item [Dominium Architecture IV](../../\_reader/by\_chat/dominium\_architecture\_iv.md)
\item [Dominium Canon, Repository Alignment, and Portability Doctrine](../../\_reader/by\_chat/dominium\_complete\_conversation.md)
\item [Dominium + Domino Codex Planning and Handoff](../../\_reader/by\_chat/dominium\_domino\_codex\_planning.md)
\item [Dominium Workbench, AIDE, Presentation Architecture, Provider Strategy, Product-Spine Planning, and Preservation Companion](../../\_reader/by\_chat/dominium\_full\_conversation.md)
\item [Dominium Setup Architecture and Handoff](../../\_reader/by\_chat/dominium\_setup.md)
\item [Domino Dominium Workbench](../../\_reader/by\_chat/domino\_dominium\_workbench.md)
\item [Dominium/Domino Engine Refactor Planning](../../\_reader/by\_chat/domino\_engine\_refactor\_prompts.md)
\item [Domino/Dominium Engine Baseline, Architecture, and Feasibility](../../\_reader/by\_chat/engine\_baseline\_architecture.md)
\item [Dominium Foundation Lock, Workbench Spine, and Parallel Codex Handoff](../../\_reader/by\_chat/foundation\_workbench\_codex.md)
\item [Domino Framework and Open-Source Provider Architecture](../../\_reader/by\_chat/framework\_open\_source\_provider.md)
\item [Dominium Content and GUI Rebuild Planning](../../\_reader/by\_chat/gui\_binary\_content.md)
\item [Dominium Language, Platform, and Architecture Baseline](../../\_reader/by\_chat/language\_platform\_architecture.md)
\item [Dominium Launcher Application-Layer Handoff](../../\_reader/by\_chat/launcher\_app\_layer.md)
\item [Dominium Launcher and Setup Architecture](../../\_reader/by\_chat/launcher\_setup\_architecture.md)
\item [Dominium Modularity, AIDE Refactorability, and Future-Proof Architecture](../../\_reader/by\_chat/modularity\_aide\_refactorability.md)
\item [Dominium Omega/Xi/Pi Architecture \& Future-Proofing Planning](../../\_reader/by\_chat/omega\_xi\_pi\_architecture\_future.md)
\item [Dominium OS-like Architecture, Repository Convergence, and Interface Operating Layer](../../\_reader/by\_chat/os\_interface\_repo\_architecture.md)
\item [Dominium Codex Platform Renderer API Plan](../../\_reader/by\_chat/platform\_renderer\_api\_plan.md)
\item [Dominium Platform Support Planning](../../\_reader/by\_chat/platform\_support.md)
\item [Dominium README Architecture, Ports, and Determinism](../../\_reader/by\_chat/readme\_ports\_determinism.md)
\item [Dominium + Domino Refactor Architecture](../../\_reader/by\_chat/refactor\_architecture.md)
\item [Dominium XStack Release Identity and Versioning](../../\_reader/by\_chat/release\_identity\_and\_versioning.md)
\item [Dominium TestX/RepoX Governance and Handoff Chat](../../\_reader/by\_chat/testx\_repox\_governance.md)
\item [Dominium Timekeeping and 2038 Resilience](../../\_reader/by\_chat/timekeeping\_and\_2038\_resilience.md)
\item [Dominium UE6, Domino, and Deterministic Universe Feasibility](../../\_reader/by\_chat/ue6\_domino\_deterministic\_universe.md)
\item [Dominium UI Editor and Tool Editor Planning](../../\_reader/by\_chat/ui\_editor\_tool\_editor\_planning.md)
\item [Dominium Universe Explorer Planning and Repo Handoff](../../\_reader/by\_chat/universe\_explorer\_planning.md)
\item [Dominium Architecture, Workbench, AIDE, and Product-Spine Planning](../../\_reader/by\_chat/workbench\_aide\_product\_spine.md)
\end{itemize}

\subsection{release}

\begin{itemize}
\item [Dominium Advanced Simulation and Infrastructure Architecture](../../\_reader/by\_chat/advanced\_simulation\_infrastructure.md)
\item [Dominium APP0 Runtime, Platform, and Renderer Architecture](../../\_reader/by\_chat/app\_runtime\_platform\_renderers.md)
\item [Dominium Architecture, Application Layer, TestX, and CodeHygiene Planning](../../\_reader/by\_chat/app\_testx\_codehygiene.md)
\item [Dominium/Domino Architecture and Codex Prompt Roadmap](../../\_reader/by\_chat/architecture\_codex\_prompts.md)
\item [Dominium Architecture, UI, Providers, and Robot OS Strategy](../../\_reader/by\_chat/architecture\_ui\_providers.md)
\item [Dominium Build and Future-Proofing Architecture](../../\_reader/by\_chat/build\_and\_future\_proofing.md)
\item [Dominium Canonical Architecture, Repository Foundation, and Provider Model](../../\_reader/by\_chat/canonical\_structure\_and\_framework.md)
\item [Dominium Development Routes and Continuity Preservation](../../\_reader/by\_chat/development\_routes.md)
\item [Documentation Standards, README Strategy, and Handoff Packaging](../../\_reader/by\_chat/documentation\_standards\_readme.md)
\item [Dominium Architecture I](../../\_reader/by\_chat/dominium\_architecture\_i.md)
\item [Dominium Architecture II](../../\_reader/by\_chat/dominium\_architecture\_ii.md)
\item [Dominium Architecture III: Launcher, Platform, Renderer, and Handoff Architecture](../../\_reader/by\_chat/dominium\_architecture\_iii.md)
\item [Dominium Architecture IV](../../\_reader/by\_chat/dominium\_architecture\_iv.md)
\item [Dominium Canon, Repository Alignment, and Portability Doctrine](../../\_reader/by\_chat/dominium\_complete\_conversation.md)
\item [Dominium + Domino Codex Planning and Handoff](../../\_reader/by\_chat/dominium\_domino\_codex\_planning.md)
\item [Dominium Workbench, AIDE, Presentation Architecture, Provider Strategy, Product-Spine Planning, and Preservation Companion](../../\_reader/by\_chat/dominium\_full\_conversation.md)
\item [Dominium Setup Architecture and Handoff](../../\_reader/by\_chat/dominium\_setup.md)
\item [Dominium/Domino Engine Refactor Planning](../../\_reader/by\_chat/domino\_engine\_refactor\_prompts.md)
\item [Domino/Dominium Engine Baseline, Architecture, and Feasibility](../../\_reader/by\_chat/engine\_baseline\_architecture.md)
\item [Dominium Foundation Lock, Workbench Spine, and Parallel Codex Handoff](../../\_reader/by\_chat/foundation\_workbench\_codex.md)
\item [Domino Framework and Open-Source Provider Architecture](../../\_reader/by\_chat/framework\_open\_source\_provider.md)
\item [Dominium Content and GUI Rebuild Planning](../../\_reader/by\_chat/gui\_binary\_content.md)
\item [Dominium Language, Platform, and Architecture Baseline](../../\_reader/by\_chat/language\_platform\_architecture.md)
\item [Dominium Launcher Application-Layer Handoff](../../\_reader/by\_chat/launcher\_app\_layer.md)
\item [Dominium Launcher and Setup Architecture](../../\_reader/by\_chat/launcher\_setup\_architecture.md)
\item [Dominium Modularity, AIDE Refactorability, and Future-Proof Architecture](../../\_reader/by\_chat/modularity\_aide\_refactorability.md)
\item [Dominium Omega/Xi/Pi Architecture \& Future-Proofing Planning](../../\_reader/by\_chat/omega\_xi\_pi\_architecture\_future.md)
\item [Dominium OS-like Architecture, Repository Convergence, and Interface Operating Layer](../../\_reader/by\_chat/os\_interface\_repo\_architecture.md)
\item [Dominium Codex Platform Renderer API Plan](../../\_reader/by\_chat/platform\_renderer\_api\_plan.md)
\item [Dominium Platform Support Planning](../../\_reader/by\_chat/platform\_support.md)
\item [Domino/Dominium Portability, Assurance, and Future-Proof Architecture](../../\_reader/by\_chat/portability\_assurance\_future\_proof.md)
\item [Dominium README Architecture, Ports, and Determinism](../../\_reader/by\_chat/readme\_ports\_determinism.md)
\item [Dominium + Domino Refactor Architecture](../../\_reader/by\_chat/refactor\_architecture.md)
\item [AIDE, XStack, and Dominium Refactor Control Plane](../../\_reader/by\_chat/refactor\_control\_plane.md)
\item [Dominium XStack Release Identity and Versioning](../../\_reader/by\_chat/release\_identity\_and\_versioning.md)
\item [Dominium TestX/RepoX Governance and Handoff Chat](../../\_reader/by\_chat/testx\_repox\_governance.md)
\item [Dominium Timekeeping and 2038 Resilience](../../\_reader/by\_chat/timekeeping\_and\_2038\_resilience.md)
\item [Dominium UE6, Domino, and Deterministic Universe Feasibility](../../\_reader/by\_chat/ue6\_domino\_deterministic\_universe.md)
\item [Dominium UI Editor and Tool Editor Planning](../../\_reader/by\_chat/ui\_editor\_tool\_editor\_planning.md)
\item [Dominium Universe Explorer Planning and Repo Handoff](../../\_reader/by\_chat/universe\_explorer\_planning.md)
\item [Dominium Architecture, Workbench, AIDE, and Product-Spine Planning](../../\_reader/by\_chat/workbench\_aide\_product\_spine.md)
\item [Dominium World Architecture](../../\_reader/by\_chat/world\_architecture.md)
\item [Dominium XStack and Lab Galaxy Handoff](../../\_reader/by\_chat/xstack\_lab\_galaxy.md)
\end{itemize}

\subsection{setup\_launcher}

\begin{itemize}
\item [Dominium APP0 Runtime, Platform, and Renderer Architecture](../../\_reader/by\_chat/app\_runtime\_platform\_renderers.md)
\item [Dominium Architecture, Application Layer, TestX, and CodeHygiene Planning](../../\_reader/by\_chat/app\_testx\_codehygiene.md)
\item [Dominium/Domino Architecture and Codex Prompt Roadmap](../../\_reader/by\_chat/architecture\_codex\_prompts.md)
\item [Dominium Canonical Architecture, Repository Foundation, and Provider Model](../../\_reader/by\_chat/canonical\_structure\_and\_framework.md)
\item [Documentation Standards, README Strategy, and Handoff Packaging](../../\_reader/by\_chat/documentation\_standards\_readme.md)
\item [Dominium Architecture I](../../\_reader/by\_chat/dominium\_architecture\_i.md)
\item [Dominium Architecture III: Launcher, Platform, Renderer, and Handoff Architecture](../../\_reader/by\_chat/dominium\_architecture\_iii.md)
\item [Dominium Architecture IV](../../\_reader/by\_chat/dominium\_architecture\_iv.md)
\item [Dominium + Domino Codex Planning and Handoff](../../\_reader/by\_chat/dominium\_domino\_codex\_planning.md)
\item [Dominium Workbench, AIDE, Presentation Architecture, Provider Strategy, Product-Spine Planning, and Preservation Companion](../../\_reader/by\_chat/dominium\_full\_conversation.md)
\item [Dominium Setup Architecture and Handoff](../../\_reader/by\_chat/dominium\_setup.md)
\item [Domino Dominium Workbench](../../\_reader/by\_chat/domino\_dominium\_workbench.md)
\item [Dominium Content and GUI Rebuild Planning](../../\_reader/by\_chat/gui\_binary\_content.md)
\item [Dominium Launcher Application-Layer Handoff](../../\_reader/by\_chat/launcher\_app\_layer.md)
\item [Dominium Launcher and Setup Architecture](../../\_reader/by\_chat/launcher\_setup\_architecture.md)
\item [Dominium Omega/Xi/Pi Architecture \& Future-Proofing Planning](../../\_reader/by\_chat/omega\_xi\_pi\_architecture\_future.md)
\item [Dominium OS-like Architecture, Repository Convergence, and Interface Operating Layer](../../\_reader/by\_chat/os\_interface\_repo\_architecture.md)
\item [Dominium Codex Platform Renderer API Plan](../../\_reader/by\_chat/platform\_renderer\_api\_plan.md)
\item [Domino/Dominium Portability, Assurance, and Future-Proof Architecture](../../\_reader/by\_chat/portability\_assurance\_future\_proof.md)
\item [Dominium + Domino Refactor Architecture](../../\_reader/by\_chat/refactor\_architecture.md)
\item [Dominium XStack Release Identity and Versioning](../../\_reader/by\_chat/release\_identity\_and\_versioning.md)
\item [Dominium TestX/RepoX Governance and Handoff Chat](../../\_reader/by\_chat/testx\_repox\_governance.md)
\item [Dominium UI Editor and Tool Editor Planning](../../\_reader/by\_chat/ui\_editor\_tool\_editor\_planning.md)
\item [Dominium Architecture, Workbench, AIDE, and Product-Spine Planning](../../\_reader/by\_chat/workbench\_aide\_product\_spine.md)
\item [Dominium World Architecture](../../\_reader/by\_chat/world\_architecture.md)
\end{itemize}

\subsection{simulation}

\begin{itemize}
\item [Dominium Advanced Simulation and Infrastructure Architecture](../../\_reader/by\_chat/advanced\_simulation\_infrastructure.md)
\item [Dominium APP0 Runtime, Platform, and Renderer Architecture](../../\_reader/by\_chat/app\_runtime\_platform\_renderers.md)
\item [Dominium Architecture, Application Layer, TestX, and CodeHygiene Planning](../../\_reader/by\_chat/app\_testx\_codehygiene.md)
\item [Dominium/Domino Architecture and Codex Prompt Roadmap](../../\_reader/by\_chat/architecture\_codex\_prompts.md)
\item [Dominium Architecture, UI, Providers, and Robot OS Strategy](../../\_reader/by\_chat/architecture\_ui\_providers.md)
\item [Dominium Build and Future-Proofing Architecture](../../\_reader/by\_chat/build\_and\_future\_proofing.md)
\item [Dominium Canonical Architecture, Repository Foundation, and Provider Model](../../\_reader/by\_chat/canonical\_structure\_and\_framework.md)
\item [Dominium Chronology \& Celestial Systems](../../\_reader/by\_chat/chronology\_celestial\_systems.md)
\item [Dominium Development Routes and Continuity Preservation](../../\_reader/by\_chat/development\_routes.md)
\item [Documentation Standards, README Strategy, and Handoff Packaging](../../\_reader/by\_chat/documentation\_standards\_readme.md)
\item [Dominium Architecture I](../../\_reader/by\_chat/dominium\_architecture\_i.md)
\item [Dominium Architecture II](../../\_reader/by\_chat/dominium\_architecture\_ii.md)
\item [Dominium Architecture III: Launcher, Platform, Renderer, and Handoff Architecture](../../\_reader/by\_chat/dominium\_architecture\_iii.md)
\item [Dominium Architecture IV](../../\_reader/by\_chat/dominium\_architecture\_iv.md)
\item [Dominium Canon, Repository Alignment, and Portability Doctrine](../../\_reader/by\_chat/dominium\_complete\_conversation.md)
\item [Dominium + Domino Codex Planning and Handoff](../../\_reader/by\_chat/dominium\_domino\_codex\_planning.md)
\item [Dominium Workbench, AIDE, Presentation Architecture, Provider Strategy, Product-Spine Planning, and Preservation Companion](../../\_reader/by\_chat/dominium\_full\_conversation.md)
\item [Dominium Setup Architecture and Handoff](../../\_reader/by\_chat/dominium\_setup.md)
\item [Domino Dominium Workbench](../../\_reader/by\_chat/domino\_dominium\_workbench.md)
\item [Dominium/Domino Engine Refactor Planning](../../\_reader/by\_chat/domino\_engine\_refactor\_prompts.md)
\item [Domino/Dominium Engine Baseline, Architecture, and Feasibility](../../\_reader/by\_chat/engine\_baseline\_architecture.md)
\item [Dominium Foundation Lock, Workbench Spine, and Parallel Codex Handoff](../../\_reader/by\_chat/foundation\_workbench\_codex.md)
\item [Domino Framework and Open-Source Provider Architecture](../../\_reader/by\_chat/framework\_open\_source\_provider.md)
\item [Dominium Content and GUI Rebuild Planning](../../\_reader/by\_chat/gui\_binary\_content.md)
\item [Dominium Language, Platform, and Architecture Baseline](../../\_reader/by\_chat/language\_platform\_architecture.md)
\item [Dominium Launcher Application-Layer Handoff](../../\_reader/by\_chat/launcher\_app\_layer.md)
\item [Dominium Launcher and Setup Architecture](../../\_reader/by\_chat/launcher\_setup\_architecture.md)
\item [Dominium Modularity, AIDE Refactorability, and Future-Proof Architecture](../../\_reader/by\_chat/modularity\_aide\_refactorability.md)
\item [Dominium Omega/Xi/Pi Architecture \& Future-Proofing Planning](../../\_reader/by\_chat/omega\_xi\_pi\_architecture\_future.md)
\item [Dominium OS-like Architecture, Repository Convergence, and Interface Operating Layer](../../\_reader/by\_chat/os\_interface\_repo\_architecture.md)
\item [Dominium Codex Platform Renderer API Plan](../../\_reader/by\_chat/platform\_renderer\_api\_plan.md)
\item [Dominium Platform Support Planning](../../\_reader/by\_chat/platform\_support.md)
\item [Domino/Dominium Portability, Assurance, and Future-Proof Architecture](../../\_reader/by\_chat/portability\_assurance\_future\_proof.md)
\item [Dominium README Architecture, Ports, and Determinism](../../\_reader/by\_chat/readme\_ports\_determinism.md)
\item [Dominium + Domino Refactor Architecture](../../\_reader/by\_chat/refactor\_architecture.md)
\item [AIDE, XStack, and Dominium Refactor Control Plane](../../\_reader/by\_chat/refactor\_control\_plane.md)
\item [Dominium XStack Release Identity and Versioning](../../\_reader/by\_chat/release\_identity\_and\_versioning.md)
\item [Dominium TestX/RepoX Governance and Handoff Chat](../../\_reader/by\_chat/testx\_repox\_governance.md)
\item [Dominium Timekeeping and 2038 Resilience](../../\_reader/by\_chat/timekeeping\_and\_2038\_resilience.md)
\item [Dominium UE6, Domino, and Deterministic Universe Feasibility](../../\_reader/by\_chat/ue6\_domino\_deterministic\_universe.md)
\item [Dominium UI Editor and Tool Editor Planning](../../\_reader/by\_chat/ui\_editor\_tool\_editor\_planning.md)
\item [Dominium Universe Explorer Planning and Repo Handoff](../../\_reader/by\_chat/universe\_explorer\_planning.md)
\item [Dominium Architecture, Workbench, AIDE, and Product-Spine Planning](../../\_reader/by\_chat/workbench\_aide\_product\_spine.md)
\item [Dominium World Architecture](../../\_reader/by\_chat/world\_architecture.md)
\end{itemize}

\subsection{timekeeping}

\begin{itemize}
\item [Dominium Advanced Simulation and Infrastructure Architecture](../../\_reader/by\_chat/advanced\_simulation\_infrastructure.md)
\item [Dominium APP0 Runtime, Platform, and Renderer Architecture](../../\_reader/by\_chat/app\_runtime\_platform\_renderers.md)
\item [Dominium Architecture, Application Layer, TestX, and CodeHygiene Planning](../../\_reader/by\_chat/app\_testx\_codehygiene.md)
\item [Dominium/Domino Architecture and Codex Prompt Roadmap](../../\_reader/by\_chat/architecture\_codex\_prompts.md)
\item [Dominium Canonical Architecture, Repository Foundation, and Provider Model](../../\_reader/by\_chat/canonical\_structure\_and\_framework.md)
\item [Dominium Chronology \& Celestial Systems](../../\_reader/by\_chat/chronology\_celestial\_systems.md)
\item [Dominium Development Routes and Continuity Preservation](../../\_reader/by\_chat/development\_routes.md)
\item [Documentation Standards, README Strategy, and Handoff Packaging](../../\_reader/by\_chat/documentation\_standards\_readme.md)
\item [Dominium Architecture I](../../\_reader/by\_chat/dominium\_architecture\_i.md)
\item [Dominium Architecture II](../../\_reader/by\_chat/dominium\_architecture\_ii.md)
\item [Dominium Architecture IV](../../\_reader/by\_chat/dominium\_architecture\_iv.md)
\item [Dominium Canon, Repository Alignment, and Portability Doctrine](../../\_reader/by\_chat/dominium\_complete\_conversation.md)
\item [Dominium + Domino Codex Planning and Handoff](../../\_reader/by\_chat/dominium\_domino\_codex\_planning.md)
\item [Dominium Workbench, AIDE, Presentation Architecture, Provider Strategy, Product-Spine Planning, and Preservation Companion](../../\_reader/by\_chat/dominium\_full\_conversation.md)
\item [Dominium Setup Architecture and Handoff](../../\_reader/by\_chat/dominium\_setup.md)
\item [Domino Dominium Workbench](../../\_reader/by\_chat/domino\_dominium\_workbench.md)
\item [Domino/Dominium Engine Baseline, Architecture, and Feasibility](../../\_reader/by\_chat/engine\_baseline\_architecture.md)
\item [Dominium Foundation Lock, Workbench Spine, and Parallel Codex Handoff](../../\_reader/by\_chat/foundation\_workbench\_codex.md)
\item [Dominium Content and GUI Rebuild Planning](../../\_reader/by\_chat/gui\_binary\_content.md)
\item [Dominium Language, Platform, and Architecture Baseline](../../\_reader/by\_chat/language\_platform\_architecture.md)
\item [Dominium Launcher and Setup Architecture](../../\_reader/by\_chat/launcher\_setup\_architecture.md)
\item [Dominium Modularity, AIDE Refactorability, and Future-Proof Architecture](../../\_reader/by\_chat/modularity\_aide\_refactorability.md)
\item [Dominium Omega/Xi/Pi Architecture \& Future-Proofing Planning](../../\_reader/by\_chat/omega\_xi\_pi\_architecture\_future.md)
\item [Dominium Codex Platform Renderer API Plan](../../\_reader/by\_chat/platform\_renderer\_api\_plan.md)
\item [Dominium Platform Support Planning](../../\_reader/by\_chat/platform\_support.md)
\item [Domino/Dominium Portability, Assurance, and Future-Proof Architecture](../../\_reader/by\_chat/portability\_assurance\_future\_proof.md)
\item [Dominium README Architecture, Ports, and Determinism](../../\_reader/by\_chat/readme\_ports\_determinism.md)
\item [Dominium + Domino Refactor Architecture](../../\_reader/by\_chat/refactor\_architecture.md)
\item [AIDE, XStack, and Dominium Refactor Control Plane](../../\_reader/by\_chat/refactor\_control\_plane.md)
\item [Dominium XStack Release Identity and Versioning](../../\_reader/by\_chat/release\_identity\_and\_versioning.md)
\item [Dominium Timekeeping and 2038 Resilience](../../\_reader/by\_chat/timekeeping\_and\_2038\_resilience.md)
\item [Dominium UE6, Domino, and Deterministic Universe Feasibility](../../\_reader/by\_chat/ue6\_domino\_deterministic\_universe.md)
\item [Dominium Architecture, Workbench, AIDE, and Product-Spine Planning](../../\_reader/by\_chat/workbench\_aide\_product\_spine.md)
\item [Dominium World Architecture](../../\_reader/by\_chat/world\_architecture.md)
\item [Dominium XStack and Lab Galaxy Handoff](../../\_reader/by\_chat/xstack\_lab\_galaxy.md)
\end{itemize}

\subsection{tooling}

\begin{itemize}
\item [Dominium Advanced Simulation and Infrastructure Architecture](../../\_reader/by\_chat/advanced\_simulation\_infrastructure.md)
\item [Dominium APP0 Runtime, Platform, and Renderer Architecture](../../\_reader/by\_chat/app\_runtime\_platform\_renderers.md)
\item [Dominium Architecture, Application Layer, TestX, and CodeHygiene Planning](../../\_reader/by\_chat/app\_testx\_codehygiene.md)
\item [Dominium/Domino Architecture and Codex Prompt Roadmap](../../\_reader/by\_chat/architecture\_codex\_prompts.md)
\item [Dominium Architecture, UI, Providers, and Robot OS Strategy](../../\_reader/by\_chat/architecture\_ui\_providers.md)
\item [Dominium Build and Future-Proofing Architecture](../../\_reader/by\_chat/build\_and\_future\_proofing.md)
\item [Dominium Canonical Architecture, Repository Foundation, and Provider Model](../../\_reader/by\_chat/canonical\_structure\_and\_framework.md)
\item [Dominium Chronology \& Celestial Systems](../../\_reader/by\_chat/chronology\_celestial\_systems.md)
\item [Dominium Development Routes and Continuity Preservation](../../\_reader/by\_chat/development\_routes.md)
\item [Documentation Standards, README Strategy, and Handoff Packaging](../../\_reader/by\_chat/documentation\_standards\_readme.md)
\item [Dominium Architecture I](../../\_reader/by\_chat/dominium\_architecture\_i.md)
\item [Dominium Architecture II](../../\_reader/by\_chat/dominium\_architecture\_ii.md)
\item [Dominium Architecture III: Launcher, Platform, Renderer, and Handoff Architecture](../../\_reader/by\_chat/dominium\_architecture\_iii.md)
\item [Dominium Architecture IV](../../\_reader/by\_chat/dominium\_architecture\_iv.md)
\item [Dominium Canon, Repository Alignment, and Portability Doctrine](../../\_reader/by\_chat/dominium\_complete\_conversation.md)
\item [Dominium + Domino Codex Planning and Handoff](../../\_reader/by\_chat/dominium\_domino\_codex\_planning.md)
\item [Dominium Workbench, AIDE, Presentation Architecture, Provider Strategy, Product-Spine Planning, and Preservation Companion](../../\_reader/by\_chat/dominium\_full\_conversation.md)
\item [Dominium Setup Architecture and Handoff](../../\_reader/by\_chat/dominium\_setup.md)
\item [Domino Dominium Workbench](../../\_reader/by\_chat/domino\_dominium\_workbench.md)
\item [Dominium/Domino Engine Refactor Planning](../../\_reader/by\_chat/domino\_engine\_refactor\_prompts.md)
\item [Domino/Dominium Engine Baseline, Architecture, and Feasibility](../../\_reader/by\_chat/engine\_baseline\_architecture.md)
\item [Dominium Foundation Lock, Workbench Spine, and Parallel Codex Handoff](../../\_reader/by\_chat/foundation\_workbench\_codex.md)
\item [Domino Framework and Open-Source Provider Architecture](../../\_reader/by\_chat/framework\_open\_source\_provider.md)
\item [Dominium Content and GUI Rebuild Planning](../../\_reader/by\_chat/gui\_binary\_content.md)
\item [Dominium Language, Platform, and Architecture Baseline](../../\_reader/by\_chat/language\_platform\_architecture.md)
\item [Dominium Launcher Application-Layer Handoff](../../\_reader/by\_chat/launcher\_app\_layer.md)
\item [Dominium Launcher and Setup Architecture](../../\_reader/by\_chat/launcher\_setup\_architecture.md)
\item [Dominium Modularity, AIDE Refactorability, and Future-Proof Architecture](../../\_reader/by\_chat/modularity\_aide\_refactorability.md)
\item [Dominium Omega/Xi/Pi Architecture \& Future-Proofing Planning](../../\_reader/by\_chat/omega\_xi\_pi\_architecture\_future.md)
\item [Dominium OS-like Architecture, Repository Convergence, and Interface Operating Layer](../../\_reader/by\_chat/os\_interface\_repo\_architecture.md)
\item [Dominium Codex Platform Renderer API Plan](../../\_reader/by\_chat/platform\_renderer\_api\_plan.md)
\item [Domino/Dominium Portability, Assurance, and Future-Proof Architecture](../../\_reader/by\_chat/portability\_assurance\_future\_proof.md)
\item [Dominium README Architecture, Ports, and Determinism](../../\_reader/by\_chat/readme\_ports\_determinism.md)
\item [Dominium + Domino Refactor Architecture](../../\_reader/by\_chat/refactor\_architecture.md)
\item [AIDE, XStack, and Dominium Refactor Control Plane](../../\_reader/by\_chat/refactor\_control\_plane.md)
\item [Dominium XStack Release Identity and Versioning](../../\_reader/by\_chat/release\_identity\_and\_versioning.md)
\item [Dominium TestX/RepoX Governance and Handoff Chat](../../\_reader/by\_chat/testx\_repox\_governance.md)
\item [Dominium Timekeeping and 2038 Resilience](../../\_reader/by\_chat/timekeeping\_and\_2038\_resilience.md)
\item [Dominium UE6, Domino, and Deterministic Universe Feasibility](../../\_reader/by\_chat/ue6\_domino\_deterministic\_universe.md)
\item [Dominium UI Editor and Tool Editor Planning](../../\_reader/by\_chat/ui\_editor\_tool\_editor\_planning.md)
\item [Dominium Universe Explorer Planning and Repo Handoff](../../\_reader/by\_chat/universe\_explorer\_planning.md)
\item [Dominium Architecture, Workbench, AIDE, and Product-Spine Planning](../../\_reader/by\_chat/workbench\_aide\_product\_spine.md)
\item [Dominium World Architecture](../../\_reader/by\_chat/world\_architecture.md)
\item [Dominium XStack and Lab Galaxy Handoff](../../\_reader/by\_chat/xstack\_lab\_galaxy.md)
\end{itemize}

\subsection{ui}

\begin{itemize}
\item [Dominium Advanced Simulation and Infrastructure Architecture](../../\_reader/by\_chat/advanced\_simulation\_infrastructure.md)
\item [Dominium APP0 Runtime, Platform, and Renderer Architecture](../../\_reader/by\_chat/app\_runtime\_platform\_renderers.md)
\item [Dominium Architecture, Application Layer, TestX, and CodeHygiene Planning](../../\_reader/by\_chat/app\_testx\_codehygiene.md)
\item [Dominium/Domino Architecture and Codex Prompt Roadmap](../../\_reader/by\_chat/architecture\_codex\_prompts.md)
\item [Dominium Architecture, UI, Providers, and Robot OS Strategy](../../\_reader/by\_chat/architecture\_ui\_providers.md)
\item [Dominium Build and Future-Proofing Architecture](../../\_reader/by\_chat/build\_and\_future\_proofing.md)
\item [Dominium Canonical Architecture, Repository Foundation, and Provider Model](../../\_reader/by\_chat/canonical\_structure\_and\_framework.md)
\item [Dominium Chronology \& Celestial Systems](../../\_reader/by\_chat/chronology\_celestial\_systems.md)
\item [Dominium Development Routes and Continuity Preservation](../../\_reader/by\_chat/development\_routes.md)
\item [Documentation Standards, README Strategy, and Handoff Packaging](../../\_reader/by\_chat/documentation\_standards\_readme.md)
\item [Dominium Architecture I](../../\_reader/by\_chat/dominium\_architecture\_i.md)
\item [Dominium Architecture II](../../\_reader/by\_chat/dominium\_architecture\_ii.md)
\item [Dominium Architecture III: Launcher, Platform, Renderer, and Handoff Architecture](../../\_reader/by\_chat/dominium\_architecture\_iii.md)
\item [Dominium Architecture IV](../../\_reader/by\_chat/dominium\_architecture\_iv.md)
\item [Dominium Canon, Repository Alignment, and Portability Doctrine](../../\_reader/by\_chat/dominium\_complete\_conversation.md)
\item [Dominium + Domino Codex Planning and Handoff](../../\_reader/by\_chat/dominium\_domino\_codex\_planning.md)
\item [Dominium Workbench, AIDE, Presentation Architecture, Provider Strategy, Product-Spine Planning, and Preservation Companion](../../\_reader/by\_chat/dominium\_full\_conversation.md)
\item [Dominium Setup Architecture and Handoff](../../\_reader/by\_chat/dominium\_setup.md)
\item [Domino Dominium Workbench](../../\_reader/by\_chat/domino\_dominium\_workbench.md)
\item [Dominium/Domino Engine Refactor Planning](../../\_reader/by\_chat/domino\_engine\_refactor\_prompts.md)
\item [Domino/Dominium Engine Baseline, Architecture, and Feasibility](../../\_reader/by\_chat/engine\_baseline\_architecture.md)
\item [Dominium Foundation Lock, Workbench Spine, and Parallel Codex Handoff](../../\_reader/by\_chat/foundation\_workbench\_codex.md)
\item [Domino Framework and Open-Source Provider Architecture](../../\_reader/by\_chat/framework\_open\_source\_provider.md)
\item [Dominium Content and GUI Rebuild Planning](../../\_reader/by\_chat/gui\_binary\_content.md)
\item [Dominium Language, Platform, and Architecture Baseline](../../\_reader/by\_chat/language\_platform\_architecture.md)
\item [Dominium Launcher Application-Layer Handoff](../../\_reader/by\_chat/launcher\_app\_layer.md)
\item [Dominium Launcher and Setup Architecture](../../\_reader/by\_chat/launcher\_setup\_architecture.md)
\item [Dominium Modularity, AIDE Refactorability, and Future-Proof Architecture](../../\_reader/by\_chat/modularity\_aide\_refactorability.md)
\item [Dominium Omega/Xi/Pi Architecture \& Future-Proofing Planning](../../\_reader/by\_chat/omega\_xi\_pi\_architecture\_future.md)
\item [Dominium OS-like Architecture, Repository Convergence, and Interface Operating Layer](../../\_reader/by\_chat/os\_interface\_repo\_architecture.md)
\item [Dominium Codex Platform Renderer API Plan](../../\_reader/by\_chat/platform\_renderer\_api\_plan.md)
\item [Dominium Platform Support Planning](../../\_reader/by\_chat/platform\_support.md)
\item [Domino/Dominium Portability, Assurance, and Future-Proof Architecture](../../\_reader/by\_chat/portability\_assurance\_future\_proof.md)
\item [Dominium README Architecture, Ports, and Determinism](../../\_reader/by\_chat/readme\_ports\_determinism.md)
\item [Dominium + Domino Refactor Architecture](../../\_reader/by\_chat/refactor\_architecture.md)
\item [AIDE, XStack, and Dominium Refactor Control Plane](../../\_reader/by\_chat/refactor\_control\_plane.md)
\item [Dominium XStack Release Identity and Versioning](../../\_reader/by\_chat/release\_identity\_and\_versioning.md)
\item [Dominium TestX/RepoX Governance and Handoff Chat](../../\_reader/by\_chat/testx\_repox\_governance.md)
\item [Dominium Timekeeping and 2038 Resilience](../../\_reader/by\_chat/timekeeping\_and\_2038\_resilience.md)
\item [Dominium UE6, Domino, and Deterministic Universe Feasibility](../../\_reader/by\_chat/ue6\_domino\_deterministic\_universe.md)
\item [Dominium UI Editor and Tool Editor Planning](../../\_reader/by\_chat/ui\_editor\_tool\_editor\_planning.md)
\item [Dominium Universe Explorer Planning and Repo Handoff](../../\_reader/by\_chat/universe\_explorer\_planning.md)
\item [Dominium Architecture, Workbench, AIDE, and Product-Spine Planning](../../\_reader/by\_chat/workbench\_aide\_product\_spine.md)
\item [Dominium World Architecture](../../\_reader/by\_chat/world\_architecture.md)
\item [Dominium XStack and Lab Galaxy Handoff](../../\_reader/by\_chat/xstack\_lab\_galaxy.md)
\end{itemize}

\subsection{workbench}

\begin{itemize}
\item [Dominium Advanced Simulation and Infrastructure Architecture](../../\_reader/by\_chat/advanced\_simulation\_infrastructure.md)
\item [Dominium Architecture, UI, Providers, and Robot OS Strategy](../../\_reader/by\_chat/architecture\_ui\_providers.md)
\item [Dominium Canonical Architecture, Repository Foundation, and Provider Model](../../\_reader/by\_chat/canonical\_structure\_and\_framework.md)
\item [Dominium Development Routes and Continuity Preservation](../../\_reader/by\_chat/development\_routes.md)
\item [Dominium Architecture IV](../../\_reader/by\_chat/dominium\_architecture\_iv.md)
\item [Dominium + Domino Codex Planning and Handoff](../../\_reader/by\_chat/dominium\_domino\_codex\_planning.md)
\item [Dominium Workbench, AIDE, Presentation Architecture, Provider Strategy, Product-Spine Planning, and Preservation Companion](../../\_reader/by\_chat/dominium\_full\_conversation.md)
\item [Domino Dominium Workbench](../../\_reader/by\_chat/domino\_dominium\_workbench.md)
\item [Domino/Dominium Engine Baseline, Architecture, and Feasibility](../../\_reader/by\_chat/engine\_baseline\_architecture.md)
\item [Dominium Foundation Lock, Workbench Spine, and Parallel Codex Handoff](../../\_reader/by\_chat/foundation\_workbench\_codex.md)
\item [Domino Framework and Open-Source Provider Architecture](../../\_reader/by\_chat/framework\_open\_source\_provider.md)
\item [Dominium Content and GUI Rebuild Planning](../../\_reader/by\_chat/gui\_binary\_content.md)
\item [Dominium Language, Platform, and Architecture Baseline](../../\_reader/by\_chat/language\_platform\_architecture.md)
\item [Dominium OS-like Architecture, Repository Convergence, and Interface Operating Layer](../../\_reader/by\_chat/os\_interface\_repo\_architecture.md)
\item [Dominium Codex Platform Renderer API Plan](../../\_reader/by\_chat/platform\_renderer\_api\_plan.md)
\item [Dominium + Domino Refactor Architecture](../../\_reader/by\_chat/refactor\_architecture.md)
\item [Dominium UI Editor and Tool Editor Planning](../../\_reader/by\_chat/ui\_editor\_tool\_editor\_planning.md)
\item [Dominium Universe Explorer Planning and Repo Handoff](../../\_reader/by\_chat/universe\_explorer\_planning.md)
\item [Dominium Architecture, Workbench, AIDE, and Product-Spine Planning](../../\_reader/by\_chat/workbench\_aide\_product\_spine.md)
\item [Dominium World Architecture](../../\_reader/by\_chat/world\_architecture.md)
\end{itemize}

\subsection{worldgen}

\begin{itemize}
\item [Dominium APP0 Runtime, Platform, and Renderer Architecture](../../\_reader/by\_chat/app\_runtime\_platform\_renderers.md)
\item [Dominium Architecture, Application Layer, TestX, and CodeHygiene Planning](../../\_reader/by\_chat/app\_testx\_codehygiene.md)
\item [Dominium/Domino Architecture and Codex Prompt Roadmap](../../\_reader/by\_chat/architecture\_codex\_prompts.md)
\item [Dominium Architecture, UI, Providers, and Robot OS Strategy](../../\_reader/by\_chat/architecture\_ui\_providers.md)
\item [Dominium Build and Future-Proofing Architecture](../../\_reader/by\_chat/build\_and\_future\_proofing.md)
\item [Dominium Chronology \& Celestial Systems](../../\_reader/by\_chat/chronology\_celestial\_systems.md)
\item [Dominium Development Routes and Continuity Preservation](../../\_reader/by\_chat/development\_routes.md)
\item [Dominium Architecture I](../../\_reader/by\_chat/dominium\_architecture\_i.md)
\item [Dominium Architecture II](../../\_reader/by\_chat/dominium\_architecture\_ii.md)
\item [Dominium Architecture III: Launcher, Platform, Renderer, and Handoff Architecture](../../\_reader/by\_chat/dominium\_architecture\_iii.md)
\item [Dominium Architecture IV](../../\_reader/by\_chat/dominium\_architecture\_iv.md)
\item [Dominium Canon, Repository Alignment, and Portability Doctrine](../../\_reader/by\_chat/dominium\_complete\_conversation.md)
\item [Dominium + Domino Codex Planning and Handoff](../../\_reader/by\_chat/dominium\_domino\_codex\_planning.md)
\item [Dominium Workbench, AIDE, Presentation Architecture, Provider Strategy, Product-Spine Planning, and Preservation Companion](../../\_reader/by\_chat/dominium\_full\_conversation.md)
\item [Dominium Setup Architecture and Handoff](../../\_reader/by\_chat/dominium\_setup.md)
\item [Domino Dominium Workbench](../../\_reader/by\_chat/domino\_dominium\_workbench.md)
\item [Dominium/Domino Engine Refactor Planning](../../\_reader/by\_chat/domino\_engine\_refactor\_prompts.md)
\item [Dominium Foundation Lock, Workbench Spine, and Parallel Codex Handoff](../../\_reader/by\_chat/foundation\_workbench\_codex.md)
\item [Domino Framework and Open-Source Provider Architecture](../../\_reader/by\_chat/framework\_open\_source\_provider.md)
\item [Dominium Content and GUI Rebuild Planning](../../\_reader/by\_chat/gui\_binary\_content.md)
\item [Dominium Language, Platform, and Architecture Baseline](../../\_reader/by\_chat/language\_platform\_architecture.md)
\item [Dominium Launcher Application-Layer Handoff](../../\_reader/by\_chat/launcher\_app\_layer.md)
\item [Dominium Launcher and Setup Architecture](../../\_reader/by\_chat/launcher\_setup\_architecture.md)
\item [Dominium Modularity, AIDE Refactorability, and Future-Proof Architecture](../../\_reader/by\_chat/modularity\_aide\_refactorability.md)
\item [Dominium Omega/Xi/Pi Architecture \& Future-Proofing Planning](../../\_reader/by\_chat/omega\_xi\_pi\_architecture\_future.md)
\item [Dominium Codex Platform Renderer API Plan](../../\_reader/by\_chat/platform\_renderer\_api\_plan.md)
\item [Dominium Platform Support Planning](../../\_reader/by\_chat/platform\_support.md)
\item [Domino/Dominium Portability, Assurance, and Future-Proof Architecture](../../\_reader/by\_chat/portability\_assurance\_future\_proof.md)
\item [Dominium README Architecture, Ports, and Determinism](../../\_reader/by\_chat/readme\_ports\_determinism.md)
\item [AIDE, XStack, and Dominium Refactor Control Plane](../../\_reader/by\_chat/refactor\_control\_plane.md)
\item [Dominium XStack Release Identity and Versioning](../../\_reader/by\_chat/release\_identity\_and\_versioning.md)
\item [Dominium TestX/RepoX Governance and Handoff Chat](../../\_reader/by\_chat/testx\_repox\_governance.md)
\item [Dominium Timekeeping and 2038 Resilience](../../\_reader/by\_chat/timekeeping\_and\_2038\_resilience.md)
\item [Dominium UE6, Domino, and Deterministic Universe Feasibility](../../\_reader/by\_chat/ue6\_domino\_deterministic\_universe.md)
\item [Dominium UI Editor and Tool Editor Planning](../../\_reader/by\_chat/ui\_editor\_tool\_editor\_planning.md)
\item [Dominium Universe Explorer Planning and Repo Handoff](../../\_reader/by\_chat/universe\_explorer\_planning.md)
\item [Dominium Architecture, Workbench, AIDE, and Product-Spine Planning](../../\_reader/by\_chat/workbench\_aide\_product\_spine.md)
\item [Dominium World Architecture](../../\_reader/by\_chat/world\_architecture.md)
\item [Dominium XStack and Lab Galaxy Handoff](../../\_reader/by\_chat/xstack\_lab\_galaxy.md)
\end{itemize}

\subsection{xstack\_aide}

\begin{itemize}
\item [Dominium Architecture, Application Layer, TestX, and CodeHygiene Planning](../../\_reader/by\_chat/app\_testx\_codehygiene.md)
\item [Dominium Canonical Architecture, Repository Foundation, and Provider Model](../../\_reader/by\_chat/canonical\_structure\_and\_framework.md)
\item [Dominium Canon, Repository Alignment, and Portability Doctrine](../../\_reader/by\_chat/dominium\_complete\_conversation.md)
\item [Dominium Workbench, AIDE, Presentation Architecture, Provider Strategy, Product-Spine Planning, and Preservation Companion](../../\_reader/by\_chat/dominium\_full\_conversation.md)
\item [Dominium Foundation Lock, Workbench Spine, and Parallel Codex Handoff](../../\_reader/by\_chat/foundation\_workbench\_codex.md)
\item [Dominium Language, Platform, and Architecture Baseline](../../\_reader/by\_chat/language\_platform\_architecture.md)
\item [Dominium Modularity, AIDE Refactorability, and Future-Proof Architecture](../../\_reader/by\_chat/modularity\_aide\_refactorability.md)
\item [Dominium Omega/Xi/Pi Architecture \& Future-Proofing Planning](../../\_reader/by\_chat/omega\_xi\_pi\_architecture\_future.md)
\item [AIDE, XStack, and Dominium Refactor Control Plane](../../\_reader/by\_chat/refactor\_control\_plane.md)
\item [Dominium XStack Release Identity and Versioning](../../\_reader/by\_chat/release\_identity\_and\_versioning.md)
\item [Dominium TestX/RepoX Governance and Handoff Chat](../../\_reader/by\_chat/testx\_repox\_governance.md)
\item [Dominium Universe Explorer Planning and Repo Handoff](../../\_reader/by\_chat/universe\_explorer\_planning.md)
\item [Dominium Architecture, Workbench, AIDE, and Product-Spine Planning](../../\_reader/by\_chat/workbench\_aide\_product\_spine.md)
\item [Dominium XStack and Lab Galaxy Handoff](../../\_reader/by\_chat/xstack\_lab\_galaxy.md)
\end{itemize}



{\small\raggedright SOURCE PATH: docs/\allowbreak{}archive/\allowbreak{}conversations/\allowbreak{}\_\allowbreak{}wiki/\allowbreak{}decisions/\allowbreak{}index.md\par}


\section{Decisions}

Decision-like entries are extracted from archived text and are not promoted.


\subsection{[Dominium Advanced Simulation and Infrastructure Architecture](../../\_reader/by\_chat/advanced\_simulation\_infrastructure.md)}

\begin{itemize}
\item The user then asked whether arbitrary placement was possible or whether the system was stuck to grids. The answer established that the engine should be grid-agnostic. [FACT] The user's question directly made arbitrary placement a major topic. [INFERENCE] The answer was accepted as the current design direction because the user then asked for a Codex prompt plan to implement changes necessary to achieve it.
\item Finally, the chat shifted into archival mode: it produced a maximum-fidelity context transfer packet, then a downloadable report package, then an in-chat reader version. Those later outputs preserved the discussion for future aggregation.
\item [INFERENCE] The conclusion was not a finalized implementation but a working architecture: each simulation domain should be deterministic, fixed-point, content-driven, versioned, replayable, and integrated through clear read/write boundaries.
\item [INFERENCE] This was one of the strongest architecture decisions because later topics repeatedly built on it.
\item This topic matters because it defines the game's construction ambition. The engine must be flexible enough to represent real infrastructure, but the UI must not become a full CAD system.
\item The key principle was that not every system should depend on decor directly. Decorative signs and decals are cheap visual objects. If they affect simulation or gameplay, they must be promoted to devices with proper STRUCT/TRANS state and ports.
\item The final phase was preservation. The user asked for a maximum-fidelity context transfer packet, then a downloadable report package, then an in-chat report reader. These artifacts were about preserving the chat for future use and aggregation, not about new architecture decisions.
\item The chat started as a simulation architecture task. It then became a gameplay layering task. It then became an infrastructure representation task. It then became a construction UX and arbitrary placement task. Finally, it became an implementation/handoff/reporting task.
\end{itemize}

\subsection{[Dominium APP0 Runtime, Platform, and Renderer Architecture](../../\_reader/by\_chat/app\_runtime\_platform\_renderers.md)}

\begin{itemize}
\item The central topic was APP0: the application/runtime/platform/renderer layer of Dominium. This topic came up because the user provided a formal prompt defining what APP0 should do.
\item The most important conclusion was that APP0 must remain outside the core simulation rules. **FACT:** the user explicitly forbade redesigning simulation rules, altering content definitions, changing life/civilization/economy logic, and introducing gameplay shortcuts. **FACT:** authoritative logic remains in engine + game.
\item Future work must define the client's allowed dependency surface and whether it can access any simulation-adjacent prediction layer.
\item The server was defined as the authoritative runtime host. It handles simulation, sharding, scheduling, law enforcement, persistence, integrity, and anti-cheat. It must be headless-capable and support AI-only autorun and MMO scale.
\item The key conclusion is that the server should not depend on graphics or windowing. This matters because dedicated servers, cloud servers, and AI-only autorun must run without a GPU or display.
\item The key conclusion is that the launcher must not install content or alter simulation state. That belongs to setup. This separation matters because launchers often become bloated, state-mutating, or inconsistent with installer/versioning logic.
\item Tools were discussed as world inspectors, economy inspectors, replay viewers, validators, optional editors, and profilers. The important point is that tools must use the same authority rules, never bypass law, and be auditable.
\item This is a distinctive part of Dominium's philosophy. Many projects treat tools as privileged by default. This chat preserved the opposite idea: tools may need elevated abilities, but those abilities must be explicit, authorized, and logged.
\end{itemize}

\subsection{[Dominium Architecture, Application Layer, TestX, and CodeHygiene Planning](../../\_reader/by\_chat/app\_testx\_codehygiene.md)}

\begin{itemize}
\item A large set of prompts was then generated to lock down architecture, determinism, performance, schema governance, rendering, epistemic UI, sharding, interest sets, and fidelity projection. This became the Phase 1 hardening layer. Additional audit prompts were introduced to ensure consistency before proceeding into life, civilization, war, content, agents, tools, mods, and final long-term policy.
\item Because the chat was huge, the user asked for a maximum-fidelity context transfer packet. Then they asked for downloadable report files. Then they asked for an in-chat reader. Finally, they asked for this human-readable explanatory report. This final report is meant to let a future assistant or human understand the substance without re-reading the whole conversation.
\item The central project is a deterministic universe engine/game. The user's ambition is not just a normal game but a world runtime that can represent everything from a single-room scenario to an AI-only universe to an MMO. The project is meant to support real-world defaults, arbitrary modding, procedural and defined content, strong epistemics, and long-term replayability.
\item The user provided a canonical repository structure and non-negotiable engine/game dependency direction. The engine is reusable and independent. The game depends on engine public APIs only. Engine internals must not leak into game/tools. Rendering backends belong to the engine, not the client. Launcher/setup logic must not enter engine/game. Client/server are executables that bind engine and game.
\item This topic is central to long-term scalability. It allows future CPUs, GPUs, NPUs, heterogeneous cores, cache structures, and distributed servers to be supported by adding backend policies rather than rewriting gameplay.
\item Rendering decisions were framed around explicit API first: Vulkan/DX12/Metal inform core contracts; legacy renderers are sinks or downgraded consumers. Rendering never affects authoritative simulation.
\item The user also wanted AI-only civilizations that could exist outside bounds of play but become visitable and indistinguishable from real civilizations. The answer was existence states and refinement contracts. A macro civilization can be alive without micro-simulating every citizen, but if it is reachable and visitable, it must refine deterministically into micro detail consistent with its macro history.
\item Large prompt families were generated for life/death/continuity, civilization/economy/governance, war/conflict, canonical content, agents, tools, mods, and final long-term maintenance. These are extensive background artifacts. Later, the chat shifted away from implementing those systems here because another bootstrap declared the core design complete and locked.
\end{itemize}

\subsection{[Dominium/Domino Architecture and Codex Prompt Roadmap](../../\_reader/by\_chat/architecture\_codex\_prompts.md)}

\begin{itemize}
\item The most important thing to remember is this: **this chat is the architectural and implementation-roadmap backbone for Dominium/Domino, but it is not an implementation log**. It should be preserved as a design/specification source and as a prompt library, while actual code and current facts must be verified separately.
\item FACT: The user opened by pasting an "Extended Master Starter Prompt - Dominium + Domino." That prompt established the role of the assistant as a senior engine architect and defined the main project philosophy. Domino was described as a deterministic, fixed-point, C89 engine portable across old and modern OS/architecture combinations. Dominium was described as the product suite built on top of Domino.
\item FACT: The user asked, succinctly, how a player would build a complete machine workshop. The assistant described a progression: extract raw materials, process them into industrial materials, produce mechanical components, produce electrical components, establish power distribution, establish logistics, assemble machine frames and machines, add measurement/control tooling, then scale horizontally and vertically.
\item FACT: The user then said they wanted inserters and belts, but "much more realistic" and with "WAY less processing power." They wanted conveyors/inserters to make sense in real-world use cases, and suggested inserters might be more generic robot arms. They also wanted belts constructed as splines like roads and rails.
\item FACT: The user explicitly wanted packaging into crates, bags, pallets, containers, jars, cans, boxes, etc. The stated purpose was to reduce the number of physical items simulated while preserving realistic packing/unpacking mechanics and obeying mass, volume, density, and other physical constraints.
\item FACT: The user then asked how hydrology, atmology, lithology, light, rain, wind, temperature, pressure, pollution, magnetic fields, radiation, gases, and similar phenomena should work with buildings and enclosed spaces. They also asked how to support blueprints/templates and how to build realistically on sloped terrain without ugly foundation blocks or forced leveling.
\item FACT: The user asked whether the whole system could be made more extensible and modular. The assistant proposed applying the same pattern everywhere: core + registry + models + data. This produced the repeated architecture model used later in prompts: each subsystem has a core API, model family/vtables, data prototypes, runtime instances, registries, TLV schemas, and save/load hooks.
\item FACT: The user asked to summarize the entire system and then asked for the hierarchy of data definitions/types and how to implement actual items/functions/features. The assistant produced major architectural summaries: IDs, protos, models, instances, systems; resource/env/build/trans/items/struct/vehicle/blueprints/jobs/worldgen/save.
\end{itemize}

\subsection{[Dominium Architecture, UI, Providers, and Robot OS Strategy](../../\_reader/by\_chat/architecture\_ui\_providers.md)}

\begin{itemize}
\item The language baseline decision matters because it changes what old renderers and toolchains matter. C17/C++17 reduces the need to design around C89/C++98 constraints but does not remove the need for C-compatible ABI boundaries. The system floor decision similarly moves many older renderers and OS targets into research/back-port categories.
\item The Robot OS and robotic seed-civilisation decisions are the most important design decisions. They are explicit user statements, not merely assistant proposals. They should be formalized as product/game requirements. They affect UI, fog-of-war, spawning, respawn, automation, industry, tutorial/planner design, and MMO anti-cheat.
\item The central tradeoff was speed versus purity. Raylib, SDL2, Lua, raygui, rlgl, rlsw, and related libraries can accelerate visible progress, but if they leak into game law, UI documents, saves, replays, or public contracts, they become architectural capture. The chosen compromise is to use them as providers behind Dominium-owned services and conformance tests.
\item The final game direction is robotic seed civilisation, not labour management.
\item 2. "Summarize the final renderer plan after the C17/C++17 transition."
\item 4. "Which decisions in this chat were explicit user decisions versus assistant synthesis?"
\item 5. "Which renderer/platform decisions are final, and which are still provisional?"
\end{itemize}

\subsection{[Dominium Build and Future-Proofing Architecture](../../\_reader/by\_chat/build\_and\_future\_proofing.md)}

\begin{itemize}
\item Plain-language limitation: this package is reliable as a preservation of the current visible chat. It should not be treated as a complete archive of every Dominium discussion in other chats. Where the report uses project or repository context, it is labelled as such. The most important caveat is that many items are recommendations, not final user decisions.
\item The final uploaded prompt requested a complete preservation package for this chat: a human-readable explanation, structured registers, context transfer packet, spec sheet, aggregator packet, self-audit, and files/ZIP package. This response completes that export task and creates downloadable files for later reading and aggregation.
\item The final explicit goal was preservation: create a maximum-fidelity package for this chat so the user can understand it later, ask questions in this chat, merge it with other old-chat reports, and eventually build a master project spec book.
\item The most important caution is that only DECISION-01 is clearly a user-stated constraint in this visible chat. The other items are strong recommendations produced in response to the user's requests for best practice; they should not be merged into canon unless the user later accepts them.
\item Manual drag-and-drop repo moves were rejected in earlier context and preserved here as a general risk; moves must go through migration maps, validators, shims, and proof.
\item 10. **Feed this report into the master spec book aggregator.** Preserve uncertainty labels and avoid merging suggestions as decisions.
\item Determinism, fixed-point behavior, and no hidden behavior are central constraints.
\item Multi-floor builds must be separate artifacts.
\end{itemize}

\subsection{[Dominium Canonical Architecture, Repository Foundation, and Provider Model](../../\_reader/by\_chat/canonical\_structure\_and\_framework.md)}

\begin{itemize}
\item This produced a pattern: generate a large Codex/AIDE task, the user ran it or reported a result, then the assistant evaluated what was truly fixed versus what was only validator/document churn. The user eventually demanded a one-shot "actual final cleanup" prompt because previous passes had sometimes added validators without moving directories. That prompt explicitly required real git mv routing, not just reports.
\item The final provider directory doctrine became service-first: runtime/<service>/providers/<provider>/, not runtime/<vendor>/<service>/. Provider choices belong in release/profiles/ or content/profiles/, not app path names. External code belongs under external/upstream/ or external/vendor/, but the repo should choose one convention.
\item Packs are distributable authored payloads. The preferred pack layout became category-based, such as content/packs/<category>/<pack\_id>/. Pack IDs must not depend on paths. Saves and replays must use stable IDs, schema versions, canonical serialization, deterministic hashes, and explicit migration/refusal policies. Backward compatibility must be a compatibility corpus plus tests, not intention.
\item INFERENCE: The user wanted a durable method for deciding future structure changes, not just a one-time layout. They also wanted the assistant to become more rigorous: distinguish decisions from brainstorms, stop over-planning without execution, and preserve uncertainty. The user also wanted a structure that supports other Domino-based games and possibly different engine implementations.
\item The closed root model was the most emotionally and technically important decision. The user repeatedly objected to root proliferation; the final root set became the anchor for all later advice. This decision affects every future proposal: if a concept can fit under an existing root, it should not become a new top-level root.
\item The no-src rule was also critical. The user considered repeated src/ and source/ folders as a central symptom of bad architecture. The final doctrine is that ownership directories are already source roots; implementation should live under meaningful module/subsystem folders, not generic src wrappers.
\item The framework-boundary decision rejected a top-level framework/ root. The accepted model is that Domino Framework exists as public surfaces, ABI headers, contracts, and provider/service law. This prevents root sprawl while still supporting SDK/export concepts later.
\item The provider model decision was a major response to raylib/SDL/Lua discussions. Third-party libraries should be providers behind first-party service contracts. They can accelerate visible progress but must not own saved data, game law, public ABI, or content schemas.
\end{itemize}

\subsection{[Dominium Chronology \& Celestial Systems](../../\_reader/by\_chat/chronology\_celestial\_systems.md)}

\begin{itemize}
\item This reinforced the core celestial-content rule: locations should exist explicitly when they change player decisions.
\item This became the spatial foundation for later time/calendar decisions because every major body or region could potentially have its own local time, calendar, or operational standard.
\item The assistant formalized this as the Perfect Earth Calendar, later called HPC-E. The final structure is clear, though the exact leap rule remains unresolved.
\item This material is useful but less final than the Earth calendar. The user asked to apply the same approach broadly, but did not individually confirm every name or number. The final package correctly marks these as needing review.
\item The final state is that this chat became both a design discussion and a packaged source document for future aggregation.
\item The conclusion was not "simulate everything." The conclusion was to classify content by relevance. Explicit objects are justified when they create gameplay, physics, navigation, history, strategy, or meaningful decisions. Otherwise, they should be procedural or statistical.
\item What remains uncertain is the exact verified object list. The assistant generated a real Milky Way system list, but that list must be checked before becoming data.
\item The conclusion was that every explicit celestial feature should change player decisions in at least one way. If it does not, it probably should not be explicit.
\end{itemize}

\subsection{[Dominium Development Routes and Continuity Preservation](../../\_reader/by\_chat/development\_routes.md)}

\begin{itemize}
\item Before relying on this chat as project authority, a future assistant or human reviewer must confirm whether Dominium is actually intended to be simulation-heavy, whether strict determinism matters, whether replay or multiplayer are goals, whether modding matters, and what the actual game loop is. Until then, the kernel-first plan should be treated as a strong provisional proposal, not a final specification.
\item The first answer was assertive. It stated that Dominium must follow Route C and described the route as the only viable one for Dominium. Later parts of the chat corrected the status of that claim. FACT: the assistant made that recommendation. UNCERTAIN / UNVERIFIED: the recommendation was not validated with external sources in the chat and was not explicitly accepted by the user.
\item The initial technical answer proposed a layered stack. It started with mathematical and temporal primitives, then moved into the deterministic simulation kernel, then world state, systems, persistence and replay, tooling, presentation, content, and finally distribution, modding, and multiplayer.
\item The user then asked to turn the Context Transfer Packet and visible chat context into a final downloadable, shareable, reusable report package. This was a packaging and normalization task. The user specified exact output files: a full chat report, YAML spec sheet, aggregator packet, registers, reader brief, verification/audit file, manifest, and ZIP package.
\item After the downloadable package was generated, the user asked not to create another handoff package but to render the package contents into a readable, navigable in-chat report. This led to an "IN-CHAT REPORT READER" response. That response explained the package contents, the workstreams, decisions, tasks, constraints, artifacts, risks, verification queue, and next questions.
\item The final state of the conversation is therefore: the chat has produced an architectural proposal for Dominium, converted that proposal into preservation artifacts, rendered those artifacts back into an in-chat reader format, and now requires a plain-English briefing that explains the substance and significance of the conversation.
\item FACT: these three routes were discussed.
\item FACT: Route C was recommended by the assistant.
\end{itemize}

\subsection{[Documentation Standards, README Strategy, and Handoff Packaging](../../\_reader/by\_chat/documentation\_standards\_readme.md)}

\begin{itemize}
\item The key thing to remember: this chat produced the standards and prompts for future work; it did not verify or change the project. The next assistant must inspect the actual repository before writing repo-specific docs, README content, platform claims, license claims, or subsystem descriptions.
\item FACT: The chat established that public API contracts should live in headers.
\item FACT: Source files should document implementation rationale, non-obvious invariants, and hazards.
\item FACT: External docs should cover architecture, dependency rules, extension recipes, and subsystem-level explanations.
\item FACT: The examples were generated in the chat.
\item FACT: They were not claimed to exist in the repository.
\item This was a major decision: README should not be the full spec. It should orient readers and point them to the deeper docs.
\item The final phase of the chat moved away from project documentation and into preserving the chat itself. The user requested a maximum-fidelity Context Transfer Packet, then requested that it be turned into a final downloadable/shareable report package. That package was created and included a full report, YAML spec sheet, aggregator packet, registers, reader brief, verification/audit file, manifest, and ZIP archive.
\end{itemize}

\subsection{[Dominium Architecture I](../../\_reader/by\_chat/dominium\_architecture\_i.md)}

\begin{itemize}
\item During this stage, the work was still broadly about "the project spec." The user wanted a comprehensive description of the game and its technical systems. The assistant produced large design documents, but many of those earlier contents are not visible in the retained transcript. This is important because the final report package marks those earlier outputs as referenced but not fully recovered.
\item The user also made a clear versioning decision. **FACT:** The user wanted **major.minor.patch semantic versioning** for all components/packages and for the complete game package. The user explicitly did **not** want build numbers or build dates in filenames. Instead, those details should live in metadata. This became part of the future build/package design.
\item The key change came when the user decided that Codex could not read the entire conversation at once. **FACT:** The user said they would copy and paste the transcript with timestamps and would not edit it manually, but because Codex could not read the whole context at once, they would not bother asking Codex to compile a full book directly from the whole chat.
\item The user provided a large devspec/ file tree. It included top-level docs, engine, platform, render, audio, systems, game, UI, launcher, mods, tools, data, tests, and scripts. Then the user changed one important thing: **FACT:** they decided to skip the top-level .txt files because the original Markdown docs already existed in /docs/....
\item The game layer began with g\_core.h, g\_core.cpp, and g\_world.h. The user then shifted to handoff/report work before g\_world.cpp was produced. **FACT:** The next strict-order file is therefore game/g\_world.cpp.txt.
\item The final generated package contains seven report files and a ZIP archive. The latest user request asks for a human-readable explanation rather than another machine-readable packet. This current report is therefore the narrative version: it explains what happened, what mattered, and what should be remembered.
\item FACT:** The whole chat revolves around the Dominium Game project. Dominium is described through its architecture: a deterministic, cross-platform simulation game with industrial, economic, logistics, climate, terrain, AI, research, and modding systems. The user was not merely asking for isolated code snippets. They were building a full software architecture.
\item The most important conclusion is that determinism is central. Systems are repeatedly specified to avoid hidden randomness, OS timers, platform dependence, uncontrolled floating point, and runtime allocation during simulation ticks. The visible rationale is that the project needs replay, save/load, and likely lockstep multiplayer compatibility. This requirement affects nearly every generated spec.
\end{itemize}

\subsection{[Dominium Architecture II](../../\_reader/by\_chat/dominium\_architecture\_ii.md)}

\begin{itemize}
\item The last part of the chat was handoff extraction. The user asked for discovery inventory, structured registers, a context transfer packet, and then a downloadable package. The generated package captured the chat's workstreams, decisions, tasks, constraints, open questions, artifacts, risks, and verification items. This current report is a plain-language explanation of that substance.
\item The entire conversation is anchored by the requirement that Dominium must run deterministically. This means the same initial state and same inputs must produce the same state across machines, compilers, operating systems, and eventually retro and modern platforms.
\item The chat discussed fixed tick phases, virtual lanes, merge order, command buffers, event queues, and deterministic serialization. The simulation must not depend on render FPS, wall-clock time, OS scheduler behaviour, GPU state, or floating-point quirks. This explains many later decisions: C89 core, integer/fixed-point math, no sim floats, no platform headers in simulation, and renderer as pure observer.
\item What remains uncertain is the exact code-level API and component layouts. The chat produced specs and prompts, not final source code.
\item The decision to use on-rails orbital mechanics mattered because it avoids the complexity of full n-body simulation. Instead of integrating every orbital body continuously, the game can use deterministic graph-like orbital states and transitions. This fits the rest of the engine: graphs, fields, frames, and state machines.
\item Railways began as a major topic but later became part of the larger transport framework. Roads, waterways, airways, and spaceways were all discussed. The final architecture treats transport as graph/corridor/path systems. Rail and road use alignments and nodes. Water can use depth fields plus navigation graphs. Air can use flight corridors, runways, and holding patterns. Space uses orbital graphs.
\item The final portion of the chat created a formal package of reports: full chat report, YAML spec sheet, aggregator packet, registers, reader brief, audit file, and manifest. These are not design decisions themselves, but they are important artifacts for later Project aggregation. The current response exists because the user wanted a human-readable explanation of the chat rather than more machine-readable outputs.
\item The building-system goal also evolved. The user originally asked vector versus voxel. The final model became hybrid: discrete blocks/cells/faces/edges as sim truth, vector as authoring/generation/visual layer.
\end{itemize}

\subsection{[Dominium Architecture III: Launcher, Platform, Renderer, and Handoff Architecture](../../\_reader/by\_chat/dominium\_architecture\_iii.md)}

\begin{itemize}
\item The view system then expanded beyond the launcher. **FACT:** The user wants instant zoom to any scale, instant switching between top-down 2D and first-person 3D, and arbitrary cameras for free cam, map viewing, content creation, HUD cameras, CCTV, overlays, windows, and offscreen targets. These are client/render-side concerns, not simulation concerns.
\item This context matters because the rest of the conversation happened inside those constraints. The launcher could not become a dumping ground for OS-specific behavior, sim mutation, renderer decisions, or nondeterministic game-state changes. It had to fit the project's deterministic layering.
\item The user then uploaded dominium.7z, saying it was the git repository as of that moment. **FACT:** The prior assistant said it could not open .7z archives in that environment. That means the actual repository state remained unverified at that stage. This became important later, because many implementation ideas were discussed without confirmed access to the repo contents.
\item The user clarified: **FACT:** "The settings selected in the launcher will be used to execute client or server binaries."
\item The user then specified another important constraint: **FACT:** the launcher backend should be C89, while frontend code for each binary can use C++98. The user also said client/server binaries can have single-system support.
\item A **C89 launcher backend/core** owns logical decisions.
\item The final direction became:
\item The conversation then refined the renderer design. The assistant had discussed vector\_soft, but the user rejected the idea of having a separate vector-only renderer. **FACT:** The user asked whether instead Dominium could have "just software renderers" capable of vector graphics and graphics using texture files on all systems as a fallback.
\end{itemize}

\subsection{[Dominium Architecture IV](../../\_reader/by\_chat/dominium\_architecture\_iv.md)}

\begin{itemize}
\item The reusable engine layer was named **Domino**. This became one of the most important naming and architectural decisions in the chat. Domino is the engine/core/platform/sim/mod layer, reusable for other games and projects. Dominium is the game/product layer built on it.
\item Blueprints are plans or diffs: place element, remove element, modify terrain, place machine, connect network, etc. When accepted, they generate jobs. Manual player actions and queued work for humans, robots, or other agents should use the same job pipeline. This matters because it keeps replay, AI, and automation consistent.
\item The user then asked how to structure Codex prompts. The final high-level roadmap became:
\item The assistant generated detailed Phase 1 prompts, platform backend templates, renderer backend templates, setup prompts, launcher prompts, tools prompts, and finally a Phase 7+ roadmap for the game. The game phase itself was not expanded into the same full prompt detail as Phases 1-6. That remains a future task.
\item At the end, the user shifted from architecture planning to chat retirement. The chat produced a discovery inventory, a context transfer packet, a final downloadable report package, and then an in-chat reader version. The current request asks for a human-readable narrative explanation rather than more registers or machine-oriented handoff files.
\item The conclusion was clear and accepted: **Domino is the reusable engine; Dominium is the game/product layer**. What remains uncertain is the actual repository state. The chat planned the split, but did not verify whether the repo already matches it.
\item The conclusion was to define UI modes separately from UI backends. This is one of the central product architecture decisions.
\item The package/instance system is central to the whole ecosystem. The user wants installed versions, official base/DLC/mod content, downloaded third-party mods and packs, shared data, instance-specific overrides, and arbitrary combinations without duplicating every possible mix.
\end{itemize}

\subsection{[Dominium Canon, Repository Alignment, and Portability Doctrine](../../\_reader/by\_chat/dominium\_complete\_conversation.md)}

\begin{itemize}
\item This chat was about preserving and re-grounding Dominium's architecture around a very old constitutional contract, then checking whether the current GitHub repository still follows that contract, and finally broadening the question into a platform-engineering doctrine for portability, modularity, extensibility, reuse, future-proofing, and long-term compatibility.
\item The glossary bound terms such as Authority, Law, Process, Lens, SessionSpec, UniverseIdentity, Macro Capsule, SRZ, RepoX, TestX, AuditX, CompatX, SecureX, and related terms. It matters because future assistants must not use sloppy synonyms like "mode" where the canon requires ExperienceProfile or LawProfile. This topic directly supports modularity because stable vocabulary is part of stable architecture.
\item Uncertainty: during this preservation pass, a potential inconsistency appeared: some docs claim product roots moved under apps/, while an earlier inspected code path under root client/ was available and apps/client was not fetched successfully. This must be verified against the current physical tree.
\item DECISION-01:** The old v1 contract/glossary were accepted as the baseline for this chat because the user supplied them as the old ground truth and asked later questions against them.
\item DECISION-02:** The live repo audit refined authority: the repo-materialized canon files and canon index are stronger for current repository work than raw prompts.
\item DECISION-03:** Ownership-root layout was preferred because it encodes who owns a file or subsystem, while generic src/source hides responsibility.
\item DECISION-04:** Stable contracts plus replaceable implementations best satisfy the user's desire for rewrites/refactors without losing compatibility.
\item DECISION-05:** The repo should not be described as fully implementing the old v1 vision yet; many systems are still declarative, scaffolded, or deferred.
\end{itemize}

\subsection{[Dominium + Domino Codex Planning and Handoff](../../\_reader/by\_chat/dominium\_domino\_codex\_planning.md)}

\begin{itemize}
\item INFERENCE:** The user accepted this direction in practice, because the next request was to generate Codex prompts to implement each step fully. The user did not explicitly say, "I formally approve the Milestone-0 plan," but they proceeded to ask for prompts based on it.
\item The user then asked for recommendations. The assistant recommended a clean startup policy: explicit flags first, no environment/config override in v1, terminal detection through dsys, generic AUTO behavior, product-specific headless handling for the game, build-time GUI/TUI capabilities, and structural error codes for fallback. The user accepted this enough to request a "one big Codex prompt" to implement it.
\item However, the later context transfer packet treated that inspection claim as **UNCERTAIN / UNVERIFIED**. This was careful and correct: within the visible chat, there was no reliable tool trace proving the repo had actually been inspected. Therefore, future assistants must verify dominium.zip or the active checkout before relying on the platform/render/API answer.
\item The core topic was the architecture of the Dominium + Domino project. **FACT:** The user defined Domino as a deterministic, fixed-point, C89 engine layer and Dominium as a C++98 product suite. This matters because every implementation decision had to respect those boundaries.
\item The conclusions were clear: future code must not bypass dsys or dgfx, deterministic sim must not use host floating-point, product code should not directly call OS APIs, and old content must not silently break. These are not just preferences; they are the governing rules of the project.
\item The numeric core was central because deterministic simulation depends on it. The generated prompt specified fixed-point types, deterministic RNG, saturating arithmetic, and a numeric test harness.
\item The unresolved complication is C89 portability. Standard C89 does not provide long long, but u64, i64, and q48\_16 require 64-bit representation. Future work must decide whether to allow compiler extensions, emulate 64-bit values for strict retro targets, or define platform tiers with conditional support.
\item Prompt 4 translated this into a minimal repo manifest system. It proposed product manifests, compatibility profile fields, and a launcher path that loads a primary game product. This was important because the launcher->game relationship is central to Dominium as a product suite.
\end{itemize}

\subsection{[Dominium Workbench, AIDE, Presentation Architecture, Provider Strategy, Product-Spine Planning, and Preservation Companion](../../\_reader/by\_chat/dominium\_full\_conversation.md)}

\begin{itemize}
\item A future assistant must understand that this chat is not simply about UI. It records the architecture for how Dominium should build itself: through contracts, commands, services, providers, documents, patches, modules, packs, apps, workspaces, diagnostics, evidence, AIDE, Codex, and eventually Workbench.
\item The user then redirected the entire plan. They argued that native OS widget GUI tools already exist through Visual Studio, Xcode, etc. The project needed a cross-platform rendered tool environment that uses the same CLI, TUI, and rendered GUI systems as the client. The old UI Editor and Tool Editor were good exploratory ideas but bad final products. This produced the Dominium Workbench concept.
\item The discussion then moved into AIDE. The user wanted to automate as much as possible through Codex and AIDE. We developed a task operating model: AIDE creates WorkUnits, tracks attempts, blockers, evidence, repairs, resumes, checkpoints, and promotion decisions. Codex executes bounded tasks. Development can continue with classified partials and repairable blockers; promotion to main requires evidence.
\item The final part of the conversation dealt with repo status and task queue. The user pasted status reports indicating that PRESENTATION-CONTRACT-01 completed with warnings, and then chose to generate six maintenance prompts before replanning. FULL-GATE-LEGACY-TEST-ROUTE-01 was generated. This preservation task followed.
\item The early UI Editor plan was a structured attempt to fix the launcher UI and build a tool for UI authoring. It covered DUI/TLV, native widgets, Win32 preview, deterministic TLV/JSON, layout, validation, codegen, ops.json, import/export, and launcher/setup generation. The conclusion was that these are useful components but not the final product architecture. They should become Workbench modules and services.
\item Workbench became the central product concept. It is a cross-platform rendered production environment, not a monolithic editor. It hosts workspaces and modules over contracts, commands, documents, patches, services, providers, packs, artifacts, evidence, and tests. It eventually self-hosts by editing its own safe artifacts first, then product UIs, packs, modules, apps, and AIDE work units.
\item The rendered GUI must be backend-agnostic. Layouts, controls, widgets, themes, styles, views, and workspaces should be modular and extensible. Themes include first-party primitive-only profiles and OEM+ mimic profiles. The GUI should be task-first, command-backed, evidence-visible, and projection-neutral.
\item Preserve all decisions, tasks, prompts, constraints, files, and artifacts in a downloadable package.
\end{itemize}

\subsection{[Dominium Setup Architecture and Handoff](../../\_reader/by\_chat/dominium\_setup.md)}

\begin{itemize}
\item Near the end, the chat shifted into preservation mode. The user asked for a maximum-fidelity context transfer packet, then for a downloadable report package, then for an in-chat reader version of that package. Those outputs preserved the decisions, tasks, constraints, artifacts, rejected options, and verification queue for future assistants or aggregation into a larger Project Spec Book.
\item These directory trees were not final, but they mattered because they exposed the user's preference for simple, shallow, logical structures and native IDE workflows that do not become the source of truth.
\item The user later asked where to create Visual Studio and Xcode apps after opening the repo as a folder. The assistant advised that Visual Studio and Xcode projects should be generated through CMake, not hand-authored or treated as authoritative source files. This decision became consistent with later Codex prompts.
\item Finally, the user retired the chat and asked for a maximum-fidelity context transfer packet. The assistant produced one. The user then asked for a downloadable shareable report package. The assistant created Markdown/YAML/report files and a ZIP. The user then asked for an in-chat reader version of the package, and the assistant rendered the package contents directly into the conversation.
\item The central topic was the setup system itself. It was discussed because the user needed a robust installer/updater/verifier capable of supporting many platforms, distribution channels, legacy environments, package managers, and future update flows.
\item FACT: The setup system was repeatedly defined as the only component allowed to install files, modify system/installation state, own installed-file metadata, repair installs, verify artifacts, uninstall, downgrade, upgrade, and roll back. This is the most important architectural rule from the chat.
\item The reason this mattered is that Dominium / Domino is a long-lived deterministic system. If launchers, engines, tools, package scripts, or platform installers each had their own install logic, the result would be untestable, non-deterministic, and difficult to audit. The setup system needed to be a central authority so that installed state, ownership, rollback, and verification all remain coherent.
\item What remains uncertain is the actual current implementation state in the repository after the applied Codex prompts. The design is clear, but the repo must still be inspected.
\end{itemize}

\subsection{[Domino Dominium Workbench](../../\_reader/by\_chat/domino\_dominium\_workbench.md)}

\begin{itemize}
\item The original UI Editor / Tool Editor plan is now **superseded as a final product**. It is not lost: its useful parts become Workbench capabilities, especially Interface Studio, UI/HUD Sandbox, Theme Laboratory, TUI Studio, Rendered GUI Studio, Preview Matrix, validation panels, import/export tools, and document-patch workflows.
\item The central decision is that Workbench must not become semantic authority. Workbench is a surface over contracts, commands, services, documents, patches, providers, packs, modules, apps, artifacts, diagnostics, evidence, tests, and replay/proof. It is the richest human/agent interface over the system, but it must not bypass command, capability, refusal, validation, document, patch, or proof law.
\item The second central decision is that CLI, TUI, rendered GUI, OS-native GUI, and headless automation must be **projections of the same command/result/refusal/document/view spine**. A GUI button, a TUI hotkey, a CLI command, a headless JSON job, and a future OS-native widget action should route through the same command/service path and produce the same typed result, diagnostics, refusals, logs, and evidence.
\item The user then rejected the old UI Editor / Tool Editor as final products. The user wanted cross-platform tools that use the same CLI, TUI, and rendered GUI systems as the client, not tools dependent on OS-native widgets. OS-native tools like Visual Studio and Xcode remain useful, but they are not the core presentation platform. This changed the plan from a UI editor to a **Workbench Platform**.
\item The final major expansion connected Workbench to the actual game direction. Universe Explorer was framed as a lawful inspection/materialization/reference-frame proof surface, not a renderer/free-camera demo. It should prove scale continuity, no modal loading, reference frames, materialization, representation ladders, interest sets, fidelity, provenance, renderer purity, and evidence.
\item Decision:** Old UI Editor and Tool Editor are superseded as final products.
\item Decision:** Workbench is not authority.
\item Reason:** If Workbench mutates files, truth, or product behavior directly, it will drift from CLI/headless/TUI/runtime semantics. Workbench must issue commands, produce patches, preview, validate, commit transactions, and record evidence.
\end{itemize}

\subsection{[Dominium/Domino Engine Refactor Planning](../../\_reader/by\_chat/domino\_engine\_refactor\_prompts.md)}

\begin{itemize}
\item FACT:** The Domino engine must be ISO C89. The Dominium UI/tools layer must be portable C++98. Engine code must not use OS APIs. Determinism matters, so floats are disallowed in deterministic paths. The engine must avoid platform-dependent behavior and hardcoded game semantics. Everything should be content-driven.
\item This became one of the central architectural ideas of the chat. It makes behavior extensible without giving minds direct write access to the world. A mod can add new sensors, observations, intents, capabilities, or actions, but state mutation still passes through a deterministic pipeline.
\item After the architecture and prompts were created, the user asked for a maximum-fidelity context transfer packet. The assistant produced one. The user then asked for a downloadable report package; the assistant generated package files and a ZIP. The user later asked for an in-chat reader version of that package, and finally requested this current human-readable explanatory report.
\item The final state of the conversation is therefore not just an architecture. It is an architecture plus an implementation roadmap plus preservation artifacts. The important substance remains the deterministic engine design and the prompt plan, not the packaging files.
\item The core idea is that deterministic simulation cannot tolerate platform-dependent behavior, nondeterministic iteration, floating-point inconsistencies, OS dependencies, or UI-driven state mutation. The engine must be portable, reproducible, and replayable. Therefore the architecture repeatedly favored fixed-size integer IDs, fixed-point math, TLV schemas, sorted iteration, deterministic queues, and replay hashes.
\item The conclusion was not a final physics implementation. It was a framework direction: each physics family should be a deterministic family plugged into the scheduler, using content-defined TLV prototypes, fixed iteration solvers, stable vtables, and replay hooks.
\item The user asked whether the same methods could apply across the whole engine. The answer was yes. This became a central theme.
\item Instead of storing final geometry as truth, buildings store footprints, volumes, enclosure definitions, surfaces, sockets, carrier intents, and edit records. Derived compilers produce occupancy, voids, enclosure graphs, surface graphs, support graphs, and carrier artifacts. These are caches, not authority.
\end{itemize}

\subsection{[Domino/Dominium Engine Baseline, Architecture, and Feasibility](../../\_reader/by\_chat/engine\_baseline\_architecture.md)}

\begin{itemize}
\item This corrected the plan. The final sequencing became: first **Milestone 0: Make the baseline path honest**. That means fixing server/runtime circular import, CLI forwarding, session\_create -> session\_boot, missing time-anchor policy registry, and making the strict local playtest validator pass. Only after that should the builder/destruction lab begin.
\item The final user action uploaded Pasted text.txt, a detailed instruction prompt requiring a full preservation report, structured registers, spec sheet, aggregator packet, self-audit, and downloadable file package for this chat. That request is the current task.
\item The central architecture distinction was that Domino should be a reusable deterministic simulation substrate, while Dominium should be one game/product layer built on it. This came up immediately when the user asked for a total description of the project. It mattered because the user wants the code to be reusable for future games and possibly other engine projects.
\item The visible conclusion was that the engine has real deterministic substrate pieces, but these must be hardened through small playable slices and proof/replay tests. Determinism is not only a mathematical preference; it enables multiplayer authority, replay, audit, mod compatibility, and long-term artifact preservation.
\item This became the practical pivot. The user pasted evidence that the repo has compiled targets and some passing tests but lacks a finished playable path. The final decision was that before game features, the team must create one canonical repo-local playable baseline command that passes strict validation.
\item Uncertainty: this remains strategic architecture, not an implemented decision.
\item The user explicitly accepted the correction that before a builder/destruction lab, the first deliverable should be one canonical repo-local playable baseline command passing strict validation. This is the most important sequencing decision in the chat.
\item This was a recommended architectural conclusion after reading live repo docs. It is not a single user-typed acceptance sentence, but the user continued building on it. It should be treated as a strong strategic direction, not a final immutable spec.
\end{itemize}

\subsection{[Dominium Foundation Lock, Workbench Spine, and Parallel Codex Handoff](../../\_reader/by\_chat/foundation\_workbench\_codex.md)}

\begin{itemize}
\item After the root skeleton improved, the assistant and user recognized that the deeper problem was no longer top-level directories. The problem became semantic duplication and governance: what is public, what is private, what is stable, what is provisional, what is generated, what is a fixture, what must stay compatible, what can be replaced, and what proof is required.
\item The key conclusion was that Workbench is not the general module system; it is one consumer of the module/command/service/provider/pack/artifact system. Workbench must not call private tools directly. It must route through registered commands and typed results, diagnostics, refusals, views, and evidence.
\item The topic came up because the project needed a model that could support world simulation, modding, tooling, Workbench, release, portability, and future games without repeating endless refactors. The conclusion was that stable contracts and semantic IDs must be separated from replaceable private implementation.
\item The repo root cleanup dominated much of the conversation. The old root structure was visibly unacceptable and triggered intense user frustration. The conversation explored cautious moves, gates, root inventories, salvage maps, bulk routing, and deterministic quarantine. The final stance was that the root skeleton is now mostly settled and should not be re-litigated unless validators show a live violation.
\item The language policy shifted from historical C89/C++98 to C17/C++17. The mainline native architecture policy shifted toward 64-bit source-native: x86\_64 and arm64. Public ABI remains C-compatible. Persisted formats must be fixed-width, explicit little-endian, and pointer-width independent.
\item Workbench is a future product surface, not the center of authority. The user wants to get to Workbench and code soon, but Workbench must route through commands/services/views/evidence rather than private direct mutation. The first Workbench slice was narrow validation. The next needed bridge is command-result-view projection.
\item Workbench's final shape is a shell hosting modules and workspaces. It eventually includes Project Graph Explorer, Interface Studio, Module/Pack Foundry, App Composer, Release Forge, Agent Work Board, Pack Browser, Renderer/Theme Laboratory, Replay/Trace Viewer, and related tools.
\item This became the immediate next frontier. The system must avoid separate CLI, TUI, Workbench, rendered, native, and headless implementations of the same behavior. The next slice should prove a command result can become a semantic view projected across multiple surfaces, preserving the same result schema, refusal codes, diagnostics, and evidence.
\end{itemize}

\subsection{[Domino Framework and Open-Source Provider Architecture](../../\_reader/by\_chat/framework\_open\_source\_provider.md)}

\begin{itemize}
\item Finally, the user uploaded a detailed preservation prompt and requested a maximum-fidelity report, registers, spec sheet, aggregator packet, audit, and downloadable files. This package is the response to that request.
\item This became the central architectural theme. The conversation refined earlier vendor-shaped paths into service-first paths. The stable pattern is:
\item The chat inspected julesc013/dominium and found evidence of existing abstractions such as system and graphics layers, stubs, and soft-backed backends. The finding was that raylib could fill concrete provider gaps without replacing the architecture. However, current repo facts are stale/uncertain and must be verified, especially the C17/C++17 vs C90/C++98 baseline contradiction.
\item The user's large game vision requires an architecture where the world is infinite by addressability but finite by active simulation. The chat proposed activity states such as cold, warm, scheduled, active, and hot. It also proposed cells/islands/constructs as units of scheduling and authority. Clients can contribute compute, but state changes must be host-verified. This topic should become a formal spec chapter.
\item The final topic is this package. The user requested a preservation report that is human-readable first, with registers, spec sheet, aggregator packet, self-audit, and downloadable files. This report preserves the visible chat, not any inaccessible past-chat transcripts.
\item The exact first implementation branch, versions, and dependency pins are unresolved. The exact Domino Framework ABI is unresolved. The exact Lua version is unresolved. The current repository baseline must be verified. The formal policy for GPL/LGPL/unclear license material is unresolved. The sparse simulation and CAD systems need formal specs and prototypes.
\item The main accepted direction was the framework/provider approach. The user explicitly said they liked the framework approach and asked whether it could be used to make a Domino framework provided by any engine implementation and consumed by any Dominium game implementation. This makes sense because it lets the project use raylib and other libraries quickly without baking them into game law.
\item The raylib decision was also directionally accepted. The user repeatedly expressed enthusiasm for raylib and asked about using as much of raylib infrastructure as possible. The caveat is that raylib is a provider suite, not architecture. This affects rendering, audio, input, asset preview, and Workbench bootstrap.
\end{itemize}

\subsection{[Dominium Content and GUI Rebuild Planning](../../\_reader/by\_chat/gui\_binary\_content.md)}

\begin{itemize}
\item That was an important change of direction. It meant the assistant's polished prompt should not be treated as final. The work needed conceptual discussion first. The assistant then analyzed CONTENT0 and identified several issues that could cause bad assumptions to become embedded if Codex acted too soon.
\item what start scenarios must explicitly exclude,
\item This shifted the content work from "populate required files" to "design a scalable content representation system." It also made clear that procedural generation and defined data are not enemies in the user's vision. The desired architecture must allow them to coexist.
\item The user then made a strong decision: **"I've decided I want to redo the GUIs from scratch."
\item This was the most important decision in the GUI part of the chat. It changed the task from "can we make cross-platform GUIs?" to "how do we design a full multi-platform GUI and binary matrix without losing backend consistency?"
\item Those caveats were not final user decisions, but they are important warnings.
\item After this, the user said the chat was being retired and requested a maximum-fidelity context transfer packet. The assistant produced a long state-transfer document with sections for bootstrap, user preferences, workstreams, decisions, tasks, constraints, open questions, rejected options, artifacts, rationale, risks, and a copy-paste prompt.
\item Finally, the user asked for this current response: not a machine-readable handoff, not a register dump, not another package, but a clear, human-readable report explaining the substance of the conversation.
\end{itemize}

\subsection{[Dominium Language, Platform, and Architecture Baseline](../../\_reader/by\_chat/language\_platform\_architecture.md)}

\begin{itemize}
\item The 32-bit question followed. The answer was to make full native products 64-bit and keep 32-bit as constrained or projection support only. The key caveat was not to make native pointer size part of game law. Save, replay, network, IDs, and renderer IR must use fixed-width types and never serialize size\_t, long, uintptr\_t, or raw pointers.
\item The user then consolidated a broad future plan around Domino, Dominium, Workbench, AIDE, contracts, tests, replay, and evidence. The answer agreed that the plan was strong, but identified missing central pieces: composition resolver, lockfiles, compatibility corpus, trust/permissions, virtual filesystem roots, performance budgets, and stable public-surface promotion rules.
\item Finally, the user pasted advice favoring C99 or C++11 due to raylib/SDL and legacy targets. The answer rejected a pivot. Raylib and SDL2 being C APIs only means provider boundaries should be C-callable; it does not force the whole engine or game to C99. The final recommendation remained C17 + C++17, with raylib/SDL2 treated as providers and with external deployment claims placed into the verification queue.
\item The chat repeatedly compared C89, C99, C11/C17, C23, C++98, C++11, C++17, and newer C++ standards. The final working direction is C17 + C++17 for the mainline. C99 and C++11 were considered but not adopted as the project-wide baseline. Newer C23/C++20/C++23/C++26 were treated as future provider/tool lanes, not current mainline law.
\item The project must retain a C-compatible ABI. The public boundary should remain POD-only, versioned, return-code/refusal based, and free of exceptions, STL containers, classes, templates, allocator ownership, and C++ ABI assumptions. This allows providers, plugins, tools, projections, and future bindings to survive implementation changes.
\item The deterministic simulation law remains more important than the language standard. Authoritative simulation must use stable IDs, canonical ordering, fixed-width values, fixed-point math, explicit little-endian encoding, and deterministic scheduler phases. Threads may accelerate derived work, but final authoritative commit must be canonical and not depend on OS timing or thread completion.
\item Avoid another cleanup cycle by locking hard-to-change decisions early.
\item The goal changed from "should old C/C++ standards be upgraded?" to "what is the full future architecture baseline?" The user initially entertained C89/C++98 staging, then accepted C17/C++17 as mainline. The user also moved from "maybe support old systems directly" to "support old systems through constrained/projection/archive lanes."
\end{itemize}

\subsection{[Dominium Launcher Application-Layer Handoff](../../\_reader/by\_chat/launcher\_app\_layer.md)}

\begin{itemize}
\item Someone reading this report should understand one central thing: this chat is not about inventing new launcher features anymore. It is about preserving boundaries, making the launcher implementation explicit, and ensuring future work happens on top of verified code rather than vague architectural memory.
\item There was also an early suggestion to create a DUI facade, a Dominium UI abstraction, to support native widgets and fallback rendering. This idea was useful as a conceptual stepping stone but was not ultimately locked as a final requirement in the form originally suggested. Later application-layer canon emphasized **UI IR**, **command graphs**, and **binding validation** rather than a specific DUI facade design.
\item This phase mattered because it framed the launcher not as a menu program but as a **control plane** for installed products, packs, profiles, compatibility, and launch contracts. However, later canon tightened the permitted communication routes: cross-product communication must go through schema/ and libs/contracts, not arbitrary plugin conventions.
\item This later became one of the clearest examples of a superseded path. The user subsequently pasted applied Codex prompts that explicitly required CMake to generate the Visual Studio solution/projects and stated that hand-written .vcxproj files must not be the authoritative build. Therefore, all earlier "manual IDE projects as canonical" advice is no longer current.
\item That product-first direction was then finalized by a later user prompt defining the canonical repo structure.
\item The second was a final purity and contract-ownership repair prompt. It required engine purity, moved cross-product contracts to libs/contracts/include/dom\_contracts, evicted non-engine content from engine/, and hardened sanity scripts.
\item This changed the correct future behavior. The next assistant must not redesign. It must implement, audit, maintain, document, verify, and harden.
\item The user asked for a maximum-fidelity Context Transfer Packet. That packet was produced. Then the user asked for a final downloadable report package, which was created. Then the user asked to inspect that package in-chat, and an in-chat reader version was produced.
\end{itemize}

\subsection{[Dominium Launcher and Setup Architecture](../../\_reader/by\_chat/launcher\_setup\_architecture.md)}

\begin{itemize}
\item The chat also produced multiple Codex work-order prompts and finally a downloadable report package. Those artifacts are useful, but the substance is the architecture: keep engine deterministic, keep setup/launcher/runtime boundaries explicit, make the launcher optional and modular, use a strong process/instance model, and now implement the launcher over dsys/dgfx rather than ad hoc platform UI stacks.
\item FACT:** The user's initial architecture text emphasized:
\item The assistant responded by stress-testing this philosophy. It pointed out determinism risks around Lua, plugins, time sources, and ABI details, and suggested more explicit contracts for file headers, TLV sections, runtime CLI capabilities, plugin exported symbols, and setup/launcher boundaries. These were assistant proposals, not all user-stated decisions, but the user continued building on that direction.
\item The next important refinement was the user's explicit statement that the launcher should be capable of supervising, but the game must always run perfectly fine with no launcher. This prevented the design from turning into a launcher-dependent ecosystem.
\item INFERENCE:** The user accepted this direction by continuing the conversation and later asking for the entire system to be detailed. However, not every proposed ABI field should be treated as final until implemented or confirmed.
\item The user then requested a maximum-fidelity Context Transfer Packet, followed by a final downloadable report package, followed by an in-chat reader version of that package. The assistant generated the package and then rendered its structured contents in chat.
\item Uncertainty:** The final language boundary for setup remains unresolved after the later C89/dsys/dgfx launcher direction.
\item Conclusion reached:** Launcher integration is optional and additive. Runtime binaries must not require launcher metadata.
\end{itemize}

\subsection{[Dominium Modularity, AIDE Refactorability, and Future-Proof Architecture](../../\_reader/by\_chat/modularity\_aide\_refactorability.md)}

\begin{itemize}
\item The final user action was uploading a preservation-package prompt. That prompt requested a maximum-fidelity preservation report, structured registers, context transfer packet, spec sheet, aggregator packet, self-audit, downloadable files, and a final in-chat reader. This package is the result.
\item The user objected to the XStack/AuditX/RepoX/TestX framing and wanted AIDE adopted quickly. The resulting plan made AIDE the repo-native control plane for inventory, task queues, policies, move maps, salvage maps, evidence ledgers, validation, and refactor history. Existing tools should be recycled, not discarded. This is central to future refactorability.
\item The discussion proposed target-based CMake and module boundaries rather than path mythology. Public headers, private headers, allowed dependencies, and forbidden dependencies should be explicit. Apps must remain thin. Engine must not depend on game/apps/runtime UI. Runtime adapts host/platform/rendering without owning simulation truth.
\item The uploaded prompt converted the conversation into a preservation task. It requires a human-readable report first, then registers, spec sheet, aggregator packet, self-audit, and downloadable files. This final package is intended to be merged later with old-chat reports into a master Project Spec Book.
\item The conversation implies a goal of turning Dominium into a long-lived product-line platform rather than a single-game repository. It also implies a desire for auditability: decisions should be grounded, visible, mechanically enforceable, and later aggregatable into a master spec. The user appears to value stable boundaries, explicit compatibility policy, and practical Codex/AIDE tasks.
\item The initial focus was directory and distribution structure. It then widened to component naming, repository convergence, AIDE governance, refactorability, and finally broad software engineering doctrine for reusable engine/platform development. The most important change was the rejection of the XStack-style framing and the shift toward AIDE as the near-term control plane.
\item The live repo still needs verification. The exact AIDE files, contracts, validators, and migration tasks must be implemented. It remains unresolved which existing roots and tools map to which final homes, which APIs become stable, which code is reusable across products/games/projects, and which prior assistant recommendations should become formal requirements after user review.
\item Status: recommendation. Basis: Final response emphasized stabilizing correct public seams only. Rationale: Avoids maintenance paralysis. Implications: Need stability levels and deprecation policy. Related workstream: WORKSTREAM-06. Confidence: medium. Label: INFERENCE. This should be revisited if live repo evidence contradicts the assumed structure, or if the user later chooses a different product-line strategy.
\end{itemize}

\subsection{[Dominium Omega/Xi/Pi Architecture \& Future-Proofing Planning](../../\_reader/by\_chat/omega\_xi\_pi\_architecture\_future.md)}

\begin{itemize}
\item The final strategic direction before this preservation request was to run a ?-series: snapshot intake, reality extraction, blueprint reconciliation, foundation readiness, and final prompt synthesis. This is needed because plans must now be mapped onto current repo reality rather than executed abstractly. The next chat should pick up there.
\item The user later reported Xi completion. The enduring lesson is that prompt instructions are not enough; architecture must be machine-readable and enforced.
\item 4. Preserve all plans, tasks, constraints, risks, decisions, artifacts, and future directions across chats.
\item 1. Current repo reality must be reconciled with the blueprint.
\item 2. Final Sigma/Phi/Upsilon/Z execution plans must be generated from current repo reality.
\item 3. Exact runtime component boundaries must be mapped to existing code.
\item 4. Agent governance and vendor mirrors must be implemented.
\item 5. Build/release preset drift must be audited and consolidated.
\end{itemize}

\subsection{[Dominium OS-like Architecture, Repository Convergence, and Interface Operating Layer](../../\_reader/by\_chat/os\_interface\_repo\_architecture.md)}

\begin{itemize}
\item DECISION-01 - CLI mandatory, TUI expected, GUI modular.** This was part of the initial architecture baseline and was repeatedly reaffirmed. It matters because every product must remain operable in recovery/headless/automation contexts. GUI is allowed but not authoritative.
\item DECISION-02 - Thin GUI shells over shared contracts.** The chat consistently rejected GUI families becoming separate product architectures. This affects client, launcher, setup, server admin, tools, and Workbench modules.
\item DECISION-03 - Repo ownership layout.** The user pushed for repository convergence; the discussion established that folders should map to ownership and contract boundaries. This decision is final enough to guide work, but details remain subject to machine-readable layout contracts.
\item DECISION-04 - Current post-CONVERGE authority.** The final audit and layout target docs define current roots and authority stack. Older docs remain context but not current physical-layout authority.
\item DECISION-05 - OS-like deterministic operating environment.** The user proposed this framing and the assistant endorsed it. It is not a literal OS decision; it is a conceptual architecture decision. It should be formalized before implementation.
\item DECISION-06 - Workbench modular host.** The user explicitly accepted the idea of one integrated environment with many focused modules. This supersedes the old UI Editor / Tool Editor final-product plan.
\item DECISION-07 - Interface Operating Layer.** The final UI discussion generalized Workbench into a cross-product interface platform. This is a strong direction, but it should be formalized as doctrine and contract work before being treated as implemented.
\item DECISION-08 - Rendered mode product-declared.** This was an assistant correction based on AppShell docs. It is necessary because current AppShell docs say rendered mode is client-only. User has not yet separately confirmed the doctrine change, so treat it as a required proposed decision.
\end{itemize}

\subsection{[Dominium Codex Platform Renderer API Plan](../../\_reader/by\_chat/platform\_renderer\_api\_plan.md)}

\begin{itemize}
\item The chat ended by generating a maximum-fidelity transfer packet, then a downloadable report package, then an in-chat reader version. The main thing to remember is this: **the active artifact is the final 9-prompt post-master Codex plan, but it is a plan, not evidence that the code exists. The repo must be inspected before execution.
\item These ideas became the architectural heart of the final plan.
\item Stable ABI boundaries must use C ABI, POD structs, explicit function tables, and u32 abi\_version plus u32 struct\_size first.
\item Every evolvable subsystem must use a facade/backend model.
\item A central capability registry must drive deterministic selection.
\item Determinism grades D0/D1/D2 must be formalized.
\item Launcher profiles must select feature sets/backends, not language standards.
\item Serialization/assets/saves/replays must be treated as ABI.
\end{itemize}

\subsection{[Dominium Platform Support Planning](../../\_reader/by\_chat/platform\_support.md)}

\begin{itemize}
\item FACT: That early answer was generated by the assistant, not stated by the user as a final decision. It matters because it introduced a very broad portability ambition, but it was later superseded by more practical platform planning. It should now be treated as historical context or possible research scope, not as a controlling product commitment.
\item The user then supplied a detailed inventory covering PlayStation, Xbox, Nintendo, PC handhelds, Android devices, Apple devices, Web platforms, AR, VR, and cross-cutting software targets. This inventory became the central artifact of the chat.
\item FACT: The user's list included legacy and current consoles, handhelds, firmware/OS names, Android software variants, Apple OSes, Web runtimes, AR/VR hardware, XR platforms, and graphics/runtime APIs. FACT: It was a list of desired coverage or consideration, not a final statement that every item must receive full parity support.
\item FACT: This is the strongest user-made decision in the chat. It overrides any framing where Android, iOS, or Web are treated as optional, secondary, or late-stage ports.
\item The central topic was what platforms Dominium should support and how to prioritise them. The discussion moved from a broad and somewhat speculative support vision toward a clearer product hierarchy.
\item FACT: The user explicitly decided that PC is the first primary focus and that Android, iOS, and Web are also primary top-tier support. This means future work should treat these platforms as architectural roots. They are not optional ports.
\item PC was treated as the starting primary platform. FACT: The assistant repeatedly described PC as Windows, macOS, and Linux. The user's final platform-priority statement referred to Android, iOS, and Web "after PC," which establishes PC as the first major category.
\item The key point is that Android is not just "phones." It is a fragmented ecosystem with different screen sizes, input methods, performance tiers, distribution channels, and system behaviours. Android support forces architecture decisions about lifecycle handling, backgrounding, memory pressure, graphics backends, touch input, controller input, asset size, save storage, permissions, and QA device coverage.
\end{itemize}

\subsection{[Domino/Dominium Portability, Assurance, and Future-Proof Architecture](../../\_reader/by\_chat/portability\_assurance\_future\_proof.md)}

\begin{itemize}
\item The actual repo was not inspected. The exact accepted directory structure, API policy, DDAP profile, compatibility promises, and first pilot module remain unresolved.
\item Status: Assistant recommendation; not separately accepted after answer.
\item Rationale: Replacement must be objectively testable.
\item Status: Decision in this package.
\item Rationale: Preservation must be honest about source scope.
\item The user explicitly required human-readable preservation, uncertainty labels, no invented facts, no silent inference, and no treating assistant recommendations as user decisions. The user explicitly values portability, modularity, extensibility, reuse, replaceability, future-proofing, and proper long-lived engineering.
\item Known: Assistant recommended DDAP; user has not explicitly accepted it.
\item Unknown: User's final acceptance and desired strictness.
\end{itemize}

\subsection{[Dominium README Architecture, Ports, and Determinism](../../\_reader/by\_chat/readme\_ports\_determinism.md)}

\begin{itemize}
\item The user pasted the final README after those cleanup changes. At that point, the README had the current active form: deterministic constraints clarified, ports unified under one source hierarchy, /ports metadata-only if retained, Section 9 normative, lockstep canonical, content-lock mismatch fatal, and disk format versions immutable.
\item After the README work, the user asked for a maximum-fidelity context transfer packet. The assistant produced a long state-transfer document with sections for project inventory, decisions, tasks, constraints, open questions, artifacts, risks, and next actions.
\item The central artifact was the root README.md. The README describes **Domino** as the deterministic engine core and **Dominium** as the official game/tooling/runtime layer. It is written as a high-level architecture document, not a low-level implementation spec.
\item The conclusion reached was that the README should remain descriptive, while normative rules belong in /docs/spec. This was already stated in the README and preserved. The README should therefore be clear enough to guide contributors but should avoid pretending to be the final binding technical specification.
\item What remains uncertain is whether the actual repository README.md matches the final pasted version. The chat only saw pasted text and generated prompts; it did not inspect a Git repository.
\item The important distinction was between **authoritative** and **non-authoritative** code. Authoritative code mutates canonical simulation state or engine-controlled serialized formats. That code must not use floating point. The engine core and engine-controlled on-disk formats must not contain float or double.
\item The chat also strengthened RNG discipline. RNG streams can only advance during deterministic tick phases. Debug overlays, UI, rendering, and other non-simulation layers must not advance engine RNG streams.
\item The README already listed immutable global simulation tick phases: Input, Pre-State, Simulation Lanes, Networks, Fluids \& Fields/Merge, Post-Process, and Finalize. The chat added two important constraints.
\end{itemize}

\subsection{[Dominium + Domino Refactor Architecture](../../\_reader/by\_chat/refactor\_architecture.md)}

\begin{itemize}
\item The most important thing to remember is that this chat did not implement the refactor. It designed the architecture and generated prompts and handoff material. Any future assistant must verify actual repository state before assuming files were moved, prompts were applied, or code was updated.
\item FACT:** This early discussion created the future UI/packs design context.
\item FACT:** All products should use shared dsys and dgfx. No product-specific platform or renderer code path.
\item The user wanted separate version numbers for Core, Launcher, Game, Setup, Tools, and protocols. The assistant proposed independent SemVer product versions, integer format/protocol versions, and DomProductInfo/introspection metadata so arbitrary mixes could degrade gracefully. The user corrected one important point: **Suite version should be the Game version**, not Core version. That became a final rule.
\item FACT:** Official suite releases should conventionally ship base modpack version equal to Game/Suite version.
\item FACT:** Engine logic should not hardwire strict equality. Compatibility ranges should decide validity.
\item The user asked about directory structure and then made a specific correction: remove products under Dominium. The final source structure became:
\item This is a user-stated decision and one of the most important carry-forward points.
\end{itemize}

\subsection{[AIDE, XStack, and Dominium Refactor Control Plane](../../\_reader/by\_chat/refactor\_control\_plane.md)}

\begin{itemize}
\item The user then asked about optimizing prompts for GPT-5.3/GPT-5.4 and later supplied advice about harness engineering. That advice argued that teams succeed by engineering the environment around agents: tools, docs, linters, feedback loops, memory, and observability. The conversation accepted the core point: the model is not the whole system; the harness is the multiplier.
\item The Grug Brained Developer became a design filter: avoid complexity, avoid premature abstraction, do small safe refactors, prefer integration tests, invest in logging, and respect existing structures before tearing them down. This influenced the decision to narrow AIDE and not overbuild runtime/hosts too early.
\item The user eventually named AIDE: Automated Integrated Development Environment. The accepted split became: XStack remains Dominium's strict local governance/proof profile; AIDE becomes the portable public layer. AIDE would live in its own repo, be usable as a standalone repo pack, later as CLI/app/extensions, and eventually perhaps a runtime/service/host system.
\item The current Dominium repo has improved governance and build proof but still contains many messy roots. The final plan is AIDE-controlled root recycling: inventory, classify, salvage, move, rewrite references, validate, shim, and retire exceptions. No broad feature work or root moving should occur before CTest/RepoX blockers are classified and AIDE refactor machinery exists.
\item The user supplied AIDE advice around dev branches, structured commits, and WorkUnit recovery. The accepted model is main as canonical truth, dev as governed integration, task/* as bounded work, and release/hotfix branches as needed. AIDE should enforce commit trailers, branch roles, safe sync/land/promote/prune, and repo-state-first recovery when prompts are repeated or out of order.
\item Unresolved goals include finishing AIDE Q35-Q57, stabilizing Dominium CTest/RepoX blockers, defining final AIDE install/upgrade bundles, implementing root recycling machinery, deciding when product boot proof can proceed, and later deciding whether AIDE Runtime/Gateway/Hosts are truly needed.
\item Several decisions are strategic rather than fully implemented. The version/capability model, Git workflow model, install/repair/upgrade model, and tool absorption model are accepted directions but require schemas and code. Future assistants must not mistake them for already completed implementations.
\item UNCERTAIN: Exact public-facing naming for AIDE Core/Kernel remains partly open. Exact threshold for accepting CTest-warning status remains a user/repo governance decision. Exact final AIDE installation packaging details remain pending.
\end{itemize}

\subsection{[Dominium XStack Release Identity and Versioning](../../\_reader/by\_chat/release\_identity\_and\_versioning.md)}

\begin{itemize}
\item The user then objected to versioning policy drift in real products and expressed concern about SemVer's "1.x forever" failure mode. The first proposed answer was a four-part public version scheme, GEN.EPOCH.FEATURE.PATCH, meant to avoid fake major bumps. That idea was useful as an intermediate step but later became less central because XStack already has GBN for dense build history and separate build identity.
\item The conversation then moved into how this might fit XStack. The model shifted from "better public version number" to "layered identity model." The assistant recommended preserving per-product versions, a global build number, build IDs, compatibility versions, and suite versions as separate layers. This became the foundation for later decisions.
\item The user suggested that suites might use a separate consumer-facing or marketing version, while each component could use stricter SemVer. The chat accepted this as a mature pattern: suites are curated bundles, while components may have technical versions. The model was refined to avoid synchronized fake versioning and to avoid treating a suite major version as universal breakage.
\item The user then asked if the suite version should still encode enough meaning to infer internal and external compatibility. The conversation clarified that the suite version can communicate a compatibility envelope or release family, but exact compatibility must live in explicit metadata. This led to the idea of compatibility profiles and later capabilities.
\item The chat was then summarized into a knowledge base. After that, the user proposed an even deeper internal rule: use capabilities instead of versions for compatibility. The response accepted this as a stronger model: versions identify releases; capabilities decide interoperability; GBN/BII/hash identify exact artifacts. That is the current final design direction.
\item The main design conclusion was that release identity should be layered: product/suite version, component version, build identity, compatibility/capability contracts, lifecycle channel, target/platform, package kind, and manifest metadata should remain separate. This is the central principle to carry forward.
\item The chat rebuilt SemVer from first principles and decided that strict SemVer should apply only to components with declared public contracts. Candidate true-SemVer entities include SDKs, engine libraries, tool APIs/CLIs, protocol libraries, plugin hosts, schema libraries, and reusable runtime libraries. It remains unresolved which exact repo modules qualify.
\item The user liked the visual shape of SemVer even for consumer-facing items. The chat accepted X.Y.Z[-pre][+build] as a useful shape, but required that non-SemVer products/suites be documented as release identifiers rather than API-compatibility promises.
\end{itemize}

\subsection{[Dominium TestX/RepoX Governance and Handoff Chat](../../\_reader/by\_chat/testx\_repox\_governance.md)}

\begin{itemize}
\item FACT:** The user is building Dominium / Domino as a deterministic universe engine + game, not a conventional game project. The visible chat repeatedly framed the project as long-lived, simulation-first, deterministic, modular, and designed to survive across many operating systems, toolchains, renderers, products, and distribution models.
\item The user then asked whether the implementation was industry-accepted and what could be improved. The response framed the approach as closer to game engines, operating systems, and long-lived infrastructure than typical game development. The important conclusion was that the design was not exotic, but it was unusually rigorous for games.
\item The user had another chat working on content and systems, and asked for a prompt to inform that chat of everything decided so far. A similar prompt was generated for the applications/platforms/renderers chat. These prompts established authoritative boundaries:
\item engine/game/content must be zero-asset and capability-driven;
\item applications must be thin shells;
\item build/versioning rules must not be bypassed;
\item This chat then generated new prompts, EG-TESTX and AP-TESTX, to send to engine/game/content and application/platform/render teams. Those prompts required response prompts from those teams. The user pasted back their responses. Both accepted TESTX canon and identified tensions.
\item Another important tension was "no silent fallback" versus renderer fallback. The application/platform response resolved this by saying explicit renderer selection must fail if unavailable, while auto mode can fallback only with explicit logging. This distinction became important for testability and user transparency.
\end{itemize}

\subsection{[Dominium Timekeeping and 2038 Resilience](../../\_reader/by\_chat/timekeeping\_and\_2038\_resilience.md)}

\begin{itemize}
\item A future assistant should understand that this chat contributes a time architecture doctrine for Dominium: **ACT is authority; DSYS time is runtime-only; observer clocks are derived; civil/astronomical time is projection-only; wall-clock time must never drive authoritative ordering.
\item What remains uncertain is the precise target list of legacy machines/toolchains and how far compatibility must go, especially for 16-bit compilers with weak or absent 64-bit arithmetic.
\item The final topic is this export task. The user requested a human-readable report first, followed by registers, spec prep, context-transfer packet, aggregator packet, self-audit, and files. This package is designed to let the chat be read later without reopening the full transcript and to support merger into a larger Dominium spec book.
\item The goal evolved from "will things break after 2038?" to "how should I design time for projects that should last?" and then to "does Dominium already match that doctrine?" The final upload changed the goal again into preservation and handoff.
\item The exact acceptable precision for Dominium ACT is not established. The exact legacy OS/toolchain floor is not fully established inside this chat. The user has not explicitly accepted every assistant recommendation as final doctrine.
\item Likely formal requirements include: authoritative logic must not depend on wall-clock time; ACT must remain versioned/fixed-width/logical; cross-shard ordering must use deterministic keys; platform time must remain DSYS-owned and non-authoritative; civil time must be projection-only.
\item Which decisions were repo-grounded versus merely assistant suggestions?
\item What external platform facts must be verified before formalizing the spec?
\end{itemize}

\subsection{[Dominium UE6, Domino, and Deterministic Universe Feasibility](../../\_reader/by\_chat/ue6\_domino\_deterministic\_universe.md)}

\begin{itemize}
\item FACT: This package preserves the currently visible chat about whether Dominium should use UE6, UE5, Domino, Unreal, or a custom engine layer for a highly ambitious deterministic solar-system-scale game. It also preserves the current uploaded prompt as an artifact because it defines this preservation task.
\item That response also introduced a rule that shaped the rest of the chat: Unreal should render and host Dominium, but Unreal should not define Dominium. That was the first major conceptual decision point. It reframed Dominium not as an engine project in the narrow sense, but as a portable game/simulation system with replaceable frontends.
\item Conclusion: machine gameplay must be designed as a scalable graph simulation, not as unrestricted physics.
\item The user asked whether every client can share compute load. The answer said this is only safe in limited forms. The core pattern must be client proposes, server verifies, server commits. Full client authority would break the economy, fog of war, and anti-cheat model.
\item The chat did not settle whether Unreal should definitely be used as the first client. It recommended Unreal as a plausible first high-quality frontend, but the final choice still depends on prototype results, licensing, team capacity, platform targets, and how much custom simulation work is required.
\item The assistant recommended against starting directly on UE6 because public technical access and production documentation were not treated as available in the visible chat. This was not framed as a permanent rejection of UE6. It was a timing decision: do not wait for a future engine to solve current architecture problems.
\item This is the highest-value decision candidate. Core rules, simulation, save format, economy, factory logic, fog-of-war state, and replay/hash systems should be portable. This matters because Domino portability is only plausible if the core does not depend on Unreal object lifecycles and assets.
\item The chat rejected the idea that clients can authoritatively simulate resource production, hidden state, enemy bases, combat, or persistent terrain changes in an MMO. Clients can assist, but the server must verify and commit.
\end{itemize}

\subsection{[Dominium UI Editor and Tool Editor Planning](../../\_reader/by\_chat/ui\_editor\_tool\_editor\_planning.md)}

\begin{itemize}
\item The most important thing to remember is that this chat was a planning and prompt-generation chat, not proof of implementation. **UNCERTAIN / UNVERIFIED:** No generated prompt is evidence that repository code was actually changed. The next assistant must first verify whether the prompts were executed, inspect the repo/source bundle if available, and avoid treating planned files as existing files.
\item This was one of the most important structure decisions in the chat. It prevented the initial implementation from becoming too broad while preserving the long-term goal. The UI Editor would be a bootstrap tool, not the final architecture.
\item The user uploaded screenshot bundles for setup and launcher. **UNCERTAIN / UNVERIFIED:** The screenshots were not inspected in this chat. The assistant still produced a logical layout description, and the user accepted it as useful. Future work should inspect the bundles before claiming exact screenshot fidelity.
\item The final state of the conversation is therefore not "the tool is built," but "the architecture, constraints, prompt plans, and continuation state have been carefully documented."
\item The most important technical discovery was that the project already has a UI system. **FACT:** The user described the current launcher GUI as the Dominium UI, or DUI, schema/state system using TLV. It has at least three backends: Win32 common controls through comctl32, DGFX through Domino's renderer/event pump, and a null backend.
\item What remains uncertain is the actual state of the code. The user named files, but the chat did not inspect the repository. A future assistant must verify exact paths, build targets, TLV parser implementation, and current launcher behavior.
\item FACT:** The user explicitly chose to make the UI Editor first. This first tool is a minimal Windows-only editor that can visually edit layouts and generate TLV directly. **FACT:** The Tool Editor comes later and is intended to become the full first-class DUI authoring environment.
\item This matters because the final Tool Editor is too large to build as the first move. The UI Editor is a bootstrap tool. It should become capable enough to create the Tool Editor's own UI. Then the Tool Editor can eventually edit itself, the setup program, the launcher, the game, and other tools.
\end{itemize}

\subsection{[Dominium Universe Explorer Planning and Repo Handoff](../../\_reader/by\_chat/universe\_explorer\_planning.md)}

\begin{itemize}
\item The pasted discussion on trees, screws, pottery, axes, chairs, tables, and machines established that the player should manipulate material, geometry, constraints, processes, tools, stations, and affordances. The conclusion is that "item classes" must not be the truth substrate. Recipes and blueprints can exist, but as higher-order formalizations.
\item A major theme was that useful local inventions must become portable, standardized, industrialized, and institutionally adopted. The repo now has a Formalization Chain spec and Player Desire Acceptance Map that strongly preserve this. The future work is making this playable: drafting, measuring, testing, blueprinting, certifying, teaching, manufacturing, maintaining, and revising designs.
\item Workbench was treated as the likely first surface for the Universe Explorer, but only if it remains projection, not authority. The repo queue says PRESENTATION-CONTRACT-01 is next, with PROJECTION-CONFORMANCE-01 as alternate. The conclusion was that presentation/projection contracts must come before building UI or renderer features.
\item The final plan accepted the newer repo status: structure is clean enough to stop giant structure prompts. Remaining cleanup should be targeted: stale full-gate tests, pack layout canon, residual taxonomy, AIDE state classification, public header ABI, provider split. This is a planning recommendation, not yet an explicitly accepted user decision.
\item The early goal was broad recovery of lost architecture. It became repo-audited synthesis. Then it shifted toward a concrete first milestone: Universe Explorer. Finally it shifted toward packaging the chat for preservation.
\item DECISION-01 was made by the user when they corrected the assistant: the chat's job is not to implement anything, but to plan, explore, and discuss. This governs the rest of the chat.
\item DECISION-04 is the user's strongest product-direction statement. The user said the first major objective should be a seamless 1:1 scale universe explorer before embodiment and deeper simulation systems. The assistant refined it but did not reject it.
\item DECISION-06 is not fully user-accepted yet. It is a best-plan recommendation based on the pasted repo-status and verified docs.
\end{itemize}

\subsection{[Dominium Architecture, Workbench, AIDE, and Product-Spine Planning](../../\_reader/by\_chat/workbench\_aide\_product\_spine.md)}

\begin{itemize}
\item The user noted that earlier prompts were hard to copy because they were not always contained in one code block, and that prompts should better handle dirty worktrees and concurrent tasks. From then on, generated prompts included detailed dirty worktree handling, allowed/forbidden paths, non-goals, validation, blocker classification, and commit/final-response formats.
\item The final recommended sequence was: finish replay proof and barebones client, run product spine review, then begin limited parallel dev. Larger parallel waves should wait until minimum AIDE workflow law exists.
\item The user then requested this full preservation package so a new chat could continue without re-explaining everything. This final handoff is itself part of the preservation output.
\item The conversation treated epistemics and diegetics as more than flavor. They determine what must be simulated, loaded, rendered, or refined. If a player has not observed something, it can remain summarized so long as future expansion is consistent with seed, law, and causal history.
\item DECISION-01 matters because the user's ambitions include worldgen, player invention, civilizations, modding, Workbench, AIDE, and future interfaces. A normal game-feature model would collapse. A deterministic simulation operating environment lets each capability become a contract-backed subsystem.
\item DECISION-05 matters because the user wants many agents and machines to work in parallel without being stopped by every blocker. But main cannot become a dumping ground. This is why dev is permissive and main is evidence-blocked.
\item DECISION-09 matters because separate CLI, TUI, GUI, and native systems would duplicate behavior and drift. The chosen architecture makes command/result/refusal/view/action the truth and projections the surface.
\item DECISION-14 matters because PACKAGE-MOUNT-SLICE-01 has only proven fixtures and reports. Treating it as package runtime would create false implementation claims.
\end{itemize}

\subsection{[Dominium World Architecture](../../\_reader/by\_chat/world\_architecture.md)}

\begin{itemize}
\item The future relevance of this chat is high. It should feed directly into the future project spec book and into a corrected implementation prompt. But it should not be treated as final implementation detail everywhere. Many high-level decisions are settled, but solver formulas, exact file encoding details, build system integration, actual repository layout, and some numerical representations still require verification.
\item Before reading the rest of this report, remember this: this chat's central contribution is not one isolated feature. It is a whole design philosophy for Dominium's world engine. The game should be deterministic, fixed-point, sparse, modular, content-driven, and field-based. Chunks are not the world. Chunks are just how the world is cached, meshed, streamed, and saved.
\item This comparison was not a final technical dependency, but it helped clarify the design goal: Dominium should not simply copy any one of these systems. It should combine their useful elements while avoiding their limits.
\item The user then introduced and refined a power-of-two naming system. The final scale names are:
\item This was a major turning point because it standardized the vocabulary. "Country," "state," and other geographic names were discarded in favor of neutral, technical scale names. The user also refined the universe addressing model over time. Earlier values changed; the final stated limits became 2\textasciicircum{}8 galaxies per universe, 2\textasciicircum{}24 systems per galaxy, and 2\textasciicircum{}16 planets per system.
\item A key refinement came when the user decided that working memory should use Q16.16 and save files should store Q4.12. This was important because it separated high-precision transient simulation from compact persistent storage. Runtime moving objects get finer precision. Static or saved objects are represented locally within chunks.
\item Another major decision was that untouched chunks should not be saved. Even if a procedural chunk is loaded temporarily for rendering or collision, it should not become persistent unless it contains actual dynamic state or overrides. This led to the distinction between untouched, procedural, and overridden chunks. Procedural chunks are temporary caches. Overridden chunks contain edits, objects, or saved changes.
\item Near the end, the user asked for a Codex 5.1 implementation prompt. A long prompt was produced, instructing Codex to implement the engine core, fixed-point types, world addressing, chunking, save formats, ECS, engine API, runtime CLI, and docs. Later, the context packet and final report package audited that prompt and found important defects. The prompt is useful, but not safe to use directly.
\end{itemize}

\subsection{[Dominium XStack and Lab Galaxy Handoff](../../\_reader/by\_chat/xstack\_lab\_galaxy.md)}

\begin{itemize}
\item and XStack must serve development rather than development serving XStack.
\item After the transcript appeared, the user asked for a maximum-fidelity Context Transfer Packet. The assistant produced one, with sections such as workstreams, decisions, tasks, constraints, open questions, risks, artifacts, and verification queue. Then the user asked to turn that into a downloadable report package. The assistant created Markdown/YAML files and a ZIP archive.
\item Finally, the user asked not for another package, but for an in-chat reader version of the package. The assistant rendered a structured in-chat reader. The user then asked for this current response: a human-readable, intuitive narrative report.
\item The final state of the conversation is therefore not "we are ready to code the next feature." It is: **we have a large handoff/report package that captures the state, but the next real action is repository verification and consistency audit.
\item The conclusion was that Dominium's architecture must preserve determinism, replayability, epistemic boundaries, and modularity. This remains a binding direction, but actual repository implementation must still be verified.
\item XStack became a central topic because the user needed reliable autonomous development. It was not just a testing layer; it was a development control substrate. RepoX enforces invariants, TestX proves behavior, AuditX detects drift, ControlX orchestrates, PerformX monitors budgets/performance, CompatX handles schema migration, SecureX handles trust/security, and gate/tools/xstack/run provide execution surfaces.
\item The key conclusion was that XStack must be **portable and removable**. It should be possible to use XStack in other projects or remove it from Dominium without breaking runtime. This is still one of the biggest constraints and must be verified.
\item The conclusion was that docs should be layered: README for accessibility; architecture docs for technical explanation; canon/glossary/contracts for binding doctrine. The 13-prompt transcript reports that docs/canon and many architecture/contract docs were created, but this must be verified.
\end{itemize}



{\small\raggedright SOURCE PATH: docs/\allowbreak{}archive/\allowbreak{}conversations/\allowbreak{}\_\allowbreak{}wiki/\allowbreak{}tasks/\allowbreak{}index.md\par}


\section{Tasks And Future Work}

Task-like entries are historical unless stronger current repo authority opens them.


\subsection{[Dominium Advanced Simulation and Infrastructure Architecture](../../\_reader/by\_chat/advanced\_simulation\_infrastructure.md)}

\begin{itemize}
\item The user then asked whether the corridor model could be made more extensible and performant. That prompted a refinement: the editable corridor should not be the same as the runtime representation. Users and tools manipulate high-level splines, attachments, and features. The engine compiles that into microsegments, indices, occupancy bitsets, connectivity edges, structural references, and render/sim extraction data.
\item The next step was integration. The user asked what other systems this should link into. The answer identified many: RES, ENV, BUILD, TRANS, STRUCT, SIM, POWER, FLUID/HYDRO, HEAT, ATMO, VEHICLE, JOB, AGENT, logistics, inventory, rendering, save/load, replay, multiplayer checksums, ownership/authority, permissions, construction zones, and localization.
\item The user then asked whether arbitrary placement was possible or whether the system was stuck to grids. The answer established that the engine should be grid-agnostic. [FACT] The user's question directly made arbitrary placement a major topic. [INFERENCE] The answer was accepted as the current design direction because the user then asked for a Codex prompt plan to implement changes necessary to achieve it.
\item After the architecture was established, the user asked for a Codex prompt plan to implement any necessary changes. A nine-step plan was produced: coordinate/placement specs, deterministic pose/frame math, TRANS 3D corridors, STRUCT arbitrary orientation and anchors, DECOR host placement, deterministic command schemas, UI snapping refactor, performance hardening, and docs/migration notes.
\item Finally, the chat shifted into archival mode: it produced a maximum-fidelity context transfer packet, then a downloadable report package, then an in-chat reader version. Those later outputs preserved the discussion for future aggregation.
\item Once the architecture was clear, the user wanted implementation planning. The chat produced a nine-step Codex prompt plan covering specs, fixed-point pose/frame math, TRANS, STRUCT, DECOR, deterministic commands, UI snapping, performance hardening, and docs.
\item Then it produced a prompt for another GPT-5.2 chat that already had a refactor/optimization plan. That prompt told the other chat to amend its existing plan instead of restarting, and to verify coverage of arbitrary placement, unified spatial primitives, co-location, signage, buildings, and determinism/performance.
\item The final phase was preservation. The user asked for a maximum-fidelity context transfer packet, then a downloadable report package, then an in-chat report reader. These artifacts were about preserving the chat for future use and aggregation, not about new architecture decisions.
\end{itemize}

\subsection{[Dominium APP0 Runtime, Platform, and Renderer Architecture](../../\_reader/by\_chat/app\_runtime\_platform\_renderers.md)}

\begin{itemize}
\item The conversation started when the user provided a structured prompt called **APP0 - Runtimes, Applications, Platforms \& Renderers**. It was addressed to **GPT-5.2 Codex** and framed as part of the ongoing **Dominium / Domino project**. The prompt's purpose was to build executable and platform support correctly while staying inside the application/runtime layer.
\item The initial APP0 prompt established several hard rules. The assistant was not allowed to redesign simulation rules, alter content definitions, change life/civilization/economy logic, or introduce gameplay shortcuts. This mattered because Dominium appears to have a strong distinction between **authoritative simulation/game logic** and the shells that host, display, launch, inspect, or distribute that logic.
\item A first assistant response converted the user's APP0 prompt into a large implementation pack. It proposed repository shapes, CMake targets, executable names, renderer backends, CLI flags, smoke tests, and a commit sequence. It was useful as an implementation artifact, but it was premature relative to the user's actual desired process.
\item The user immediately corrected the direction: **"First, let's discuss this before we plan or generate any prompts."** That was an important change. It meant the user did not want a Codex prompt yet. They wanted to reason about the architecture before locking implementation instructions.
\item The user then explained that by the end of the macro-prompt plan, the project should have a **running executable** capable of displaying an **interactive and resizable window** for each supported platform and renderer combination. This made the discussion more concrete. APP0 was no longer just about docs and stubs; it had to create a real runtime path.
\item The user then stated that future implementation would have permission to modify render, platform, application, client, server, and docs. They asked whether this was enough or whether engine and game also needed write access.
\item The central topic was APP0: the application/runtime/platform/renderer layer of Dominium. This topic came up because the user provided a formal prompt defining what APP0 should do.
\item Future work must define the client's allowed dependency surface and whether it can access any simulation-adjacent prediction layer.
\end{itemize}

\subsection{[Dominium Architecture, Application Layer, TestX, and CodeHygiene Planning](../../\_reader/by\_chat/app\_testx\_codehygiene.md)}

\begin{itemize}
\item A large set of prompts was then generated to lock down architecture, determinism, performance, schema governance, rendering, epistemic UI, sharding, interest sets, and fidelity projection. This became the Phase 1 hardening layer. Additional audit prompts were introduced to ensure consistency before proceeding into life, civilization, war, content, agents, tools, mods, and final long-term policy.
\item This was one of the most important engineering turns in the conversation. It made future performance work a backend/policy problem rather than a gameplay rewrite problem.
\item Because the chat was huge, the user asked for a maximum-fidelity context transfer packet. Then they asked for downloadable report files. Then they asked for an in-chat reader. Finally, they asked for this human-readable explanatory report. This final report is meant to let a future assistant or human understand the substance without re-reading the whole conversation.
\item This topic is central to long-term scalability. It allows future CPUs, GPUs, NPUs, heterogeneous cores, cache structures, and distributed servers to be supported by adding backend policies rather than rewriting gameplay.
\item Large prompt families were generated for life/death/continuity, civilization/economy/governance, war/conflict, canonical content, agents, tools, mods, and final long-term maintenance. These are extensive background artifacts. Later, the chat shifted away from implementing those systems here because another bootstrap declared the core design complete and locked.
\item For future use, those prompt families are historical/planned artifacts. They may be useful in content/spec aggregation, but current work should not redesign them.
\item Later, the user pasted a newer build/version model from another chat. That model superseded simpler earlier ideas: GBN is only for distributed artifacts, BII exists always, build kind and channel are distinct, product SemVer is manual, and all protocols/schemas/API/ABI versions are orthogonal. This matters because future testing and app prompts must use BUILD-ID-0, not the earlier model.
\item INFERENCE:** The user wanted to prevent future chats or Codex sessions from drifting, forgetting, or simplifying. They care deeply about maintaining continuity across many chats and turning fragmented planning into a stable project corpus.
\end{itemize}

\subsection{[Dominium/Domino Architecture and Codex Prompt Roadmap](../../\_reader/by\_chat/architecture\_codex\_prompts.md)}

\begin{itemize}
\item UNCERTAIN / UNVERIFIED: No repository code was inspected or modified in this chat. The prompts and architecture are plans and artifacts, not proof that implementation exists. The actual repository state, build system, existing code quality, GUI backend availability, TLV schemas, and whether any prompts have been applied remain unverified.
\item The most important thing to remember is this: **this chat is the architectural and implementation-roadmap backbone for Dominium/Domino, but it is not an implementation log**. It should be preserved as a design/specification source and as a prompt library, while actual code and current facts must be verified separately.
\item FACT: The user opened by pasting an "Extended Master Starter Prompt - Dominium + Domino." That prompt established the role of the assistant as a senior engine architect and defined the main project philosophy. Domino was described as a deterministic, fixed-point, C89 engine portable across old and modern OS/architecture combinations. Dominium was described as the product suite built on top of Domino.
\item This phase introduced one of the chat's most durable ideas: **engine systems should operate on generic IDs, tags, models, and TLV parameters, not hardcoded content families.
\item FACT: The user asked whether the whole system could be made more extensible and modular. The assistant proposed applying the same pattern everywhere: core + registry + models + data. This produced the repeated architecture model used later in prompts: each subsystem has a core API, model family/vtables, data prototypes, runtime instances, registries, TLV schemas, and save/load hooks.
\item FACT: The user asked for a plan for a prompt to Codex to implement all the changes and document them. Then the user asked how to split it into multiple prompts under ten. The assistant split the work into seven major prompts. The user then repeatedly said "Next," and the assistant generated Prompts 1 through 7 in detail.
\item Those first seven prompts covered:
\item FACT: After the first seven prompts, the user said they wanted a minimal GUI and did not want to test only in CLI. They also reminded that the engine and core know nothing about actual data definitions and that those need to stay in base/mods/packs.
\end{itemize}

\subsection{[Dominium Architecture, UI, Providers, and Robot OS Strategy](../../\_reader/by\_chat/architecture\_ui\_providers.md)}

\begin{itemize}
\item The user then said the latest docs and code were live at julesc013/dominium. The assistant treated the repo as an implementation baseline and observed that AppShell, UI mode resolution, distribution doctrine, and component matrices already existed. This changed the task from inventing a GUI strategy from scratch to reconciling the transfer brief with existing repository doctrine.
\item INFERENCE: The user wants to turn Dominium into a platform-level project whose architecture can survive decades of refactors, ports, new games, new providers, and evolving standards. The user also wants future assistants to stop re-litigating already-settled principles and instead build upon the emerging doctrine.
\item The user appears to prefer source layouts and naming that remain meaningful years later. They dislike vendor-shaped, framework-shaped, and status-word architecture. Future assistants should not propose "just use X framework" as the answer. They should instead map libraries to providers, profiles, and contracts.
\item The UI should feel like a customizable robot OS whose HUD, menus, TUI, CLI, Workbench, and future VR shell are projections of the same system.
\item The most important immediate next task is provider/Robot OS wedge planning with validators and manifests.
\item 7. "Turn PROVIDER-WEDGE-01 into a Codex-ready implementation task."
\end{itemize}

\subsection{[Dominium Build and Future-Proofing Architecture](../../\_reader/by\_chat/build\_and\_future\_proofing.md)}

\begin{itemize}
\item The user then pushed beyond build mechanics. They asked what else they were missing about portability, modularity, extensibility, replacement, rewriteability, and reuse across future games and engines. They explicitly wanted Dominium to be developed like a proper OS or game-engine platform, not a one-off indie project.
\item The final uploaded prompt requested a complete preservation package for this chat: a human-readable explanation, structured registers, context transfer packet, spec sheet, aggregator packet, self-audit, and files/ZIP package. This response completes that export task and creates downloadable files for later reading and aggregation.
\item The user wanted reusable code for different games on the Domino engine and possibly unrelated future engine/game projects. The answer was to freeze contracts rather than implementations. Stable ABI/API/data/protocol surfaces should be explicitly versioned and tested; implementations should be replaceable behind black-box conformance tests. This is the basis for future rewrites without breaking compatibility.
\item The uploaded prompt requested a complete chat preservation package with human report, registers, handoff packet, spec sheet, aggregator packet, audit, and downloadable files. This output creates that package and labels its limitations.
\item The user explicitly wanted to design a robust Dominium build system that could handle many machines, IDEs, compilers, OS floors, CMake presets, CI lanes, distribution outputs, and historical toolchains. The user also explicitly wanted Dominium's codebase to be portable, modular, extensible, reusable across future games and engines, and refactorable at the level of whole files and directories.
\item INFERENCE: The user is trying to prevent Dominium from accumulating accidental architecture. They want mechanical governance rather than relying on memory or assistant recommendations. They also want future assistants to stop repeating conceptual advice and instead preserve actionable structure, constraints, risks, and next tasks.
\item The unresolved goals are implementation goals: deciding which recommendations become canon, adding public-surface/dependency/build contracts, validating the current live repo state, and applying the proposed cleanup tasks. The preservation package is complete, but the engineering work it describes remains pending.
\item The user wants future assistants to preserve tentative status and not over-canonize recommendations.
\end{itemize}

\subsection{[Dominium Canonical Architecture, Repository Foundation, and Provider Model](../../\_reader/by\_chat/canonical\_structure\_and\_framework.md)}

\begin{itemize}
\item This produced a pattern: generate a large Codex/AIDE task, the user ran it or reported a result, then the assistant evaluated what was truly fixed versus what was only validator/document churn. The user eventually demanded a one-shot "actual final cleanup" prompt because previous passes had sometimes added validators without moving directories. That prompt explicitly required real git mv routing, not just reports.
\item INFERENCE: The user wanted a durable method for deciding future structure changes, not just a one-time layout. They also wanted the assistant to become more rigorous: distinguish decisions from brainstorms, stop over-planning without execution, and preserve uncertainty. The user also wanted a structure that supports other Domino-based games and possibly different engine implementations.
\item The closed root model was the most emotionally and technically important decision. The user repeatedly objected to root proliferation; the final root set became the anchor for all later advice. This decision affects every future proposal: if a concept can fit under an existing root, it should not become a new top-level root.
\item The Workbench discussion involved a similar tradeoff. The user wanted powerful authoring workspaces, but the architecture must not make Workbench a privileged bypass. Therefore, Workbench modules and workspaces should operate through the same commands, views, actions, documents, diagnostics, and evidence packets that CLI, CI, headless, server, and future apps use.
\item The strongest immediate maintenance tasks are: refresh stale generated evidence, repair launcher pack-verification marker debt if still blocking fast strict, run or audit full CTest/T4 failures, and classify remaining full-gate failures by cause. The task names discussed include FULL-GATE-GENERATED-EVIDENCE-REFRESH-01, FAST-STRICT-EVIDENCE-MARKER-REPAIR-01, and FULL-CTEST-AUDIT-NONPATH-01.
\item Mainline tasks discussed include PRESENTATION-CONTRACT-01, PROJECTION-CONFORMANCE-01, COMMAND-RESULT-VIEW-SLICE-01, PACKAGE-MOUNT-SLICE-01, and WORKBENCH-VALIDATION-SLICE-01. The rule is to prove narrow slices before broad UI, gameplay, renderer, provider runtime, or package runtime expansion.
\item The user wants structure to support modularity, portability, extensibility, reuse, modding, backwards compatibility, and future replacement.
\item The user strongly prefers actual execution/move prompts over endless micro-planning when structure is blocking work.
\end{itemize}

\subsection{[Dominium Chronology \& Celestial Systems](../../\_reader/by\_chat/chronology\_celestial\_systems.md)}

\begin{itemize}
\item The final state is that this chat became both a design discussion and a packaged source document for future aggregation.
\item What remains unresolved is the exact mathematical implementation of these frames and how much relativistic fidelity should exist in phase one.
\item The user asked how to extend calendars beyond planets. The answer was that the larger the scale, the less ordinary calendar structure survives. Sol can have a system-wide civil time and epoch calendar. The Milky Way should use galactic coordinate time and epoch blocks. The universe should use cosmological time and large phases/ages/eras.
\item The reason is to avoid losing context across chats and to support a future Project Spec Book. This chat produced many artifacts: prompts, reports, registers, YAML, a ZIP package, and explanatory summaries. The artifacts matter, but the substance remains the design decisions described above.
\item It is reasonable to infer that the user wants a future implementation path, likely through another assistant and eventually Codex, but does not want premature code before the broader project context is inspected.
\item It is also reasonable to infer that the user wants a future Project Spec Book assembled from multiple old-chat reports, and that this chat should contribute the chronology/celestial/timekeeping chapters.
\item This decision makes sense because it separates clean internal logic from messy external software expectations. However, it is less directly user-confirmed than the month order. Future work should ask whether compatibility projection should be default or optional.
\item The future project architecture can support data-driven systems.
\end{itemize}

\subsection{[Dominium Development Routes and Continuity Preservation](../../\_reader/by\_chat/development\_routes.md)}

\begin{itemize}
\item The conversation therefore produced both a **substantive technical proposal** and a **continuity-preservation workflow**. The technical proposal concerns how Dominium might be built. The preservation workflow concerns how the ideas from this chat should be carried forward into future chats and eventually into a larger **Project Spec Book** or **Master Living State File**.
\item Before relying on this chat as project authority, a future assistant or human reviewer must confirm whether Dominium is actually intended to be simulation-heavy, whether strict determinism matters, whether replay or multiplayer are goals, whether modding matters, and what the actual game loop is. Until then, the kernel-first plan should be treated as a strong provisional proposal, not a final specification.
\item After the initial planning answer, the user said the chat was being retired and asked for a **maximum-fidelity Context Transfer Packet**. This changed the task. The goal was no longer to continue designing Dominium directly. The goal became preserving everything necessary for a future assistant to continue without making the user repeat context.
\item The assistant produced a large Context Transfer Packet. It preserved the technical route proposal, the caveats around its uncertainty, the proposed phase order, the open questions, the risks, the artifacts, and the next actions.
\item The user then asked to turn the Context Transfer Packet and visible chat context into a final downloadable, shareable, reusable report package. This was a packaging and normalization task. The user specified exact output files: a full chat report, YAML spec sheet, aggregator packet, registers, reader brief, verification/audit file, manifest, and ZIP package.
\item After the downloadable package was generated, the user asked not to create another handoff package but to render the package contents into a readable, navigable in-chat report. This led to an "IN-CHAT REPORT READER" response. That response explained the package contents, the workstreams, decisions, tasks, constraints, artifacts, risks, verification queue, and next questions.
\item This topic connects to future work because choosing the route determines the order of implementation. If Route C is accepted, work begins with simulation primitives and determinism rather than UI or content.
\item Future work must decide how strict determinism really needs to be. Full determinism can be valuable, but it can also impose cost and complexity.
\end{itemize}

\subsection{[Documentation Standards, README Strategy, and Handoff Packaging](../../\_reader/by\_chat/documentation\_standards\_readme.md)}

\begin{itemize}
\item The key thing to remember: this chat produced the standards and prompts for future work; it did not verify or change the project. The next assistant must inspect the actual repository before writing repo-specific docs, README content, platform claims, license claims, or subsystem descriptions.
\item This became the foundation for nearly every later prompt and plan.
\item The user then asked for a prompt that Codex could use to read through every file and insert, append, amend, or create documentation, including inline comments, block comments, and docs/README content.
\item That prompt was the first large artifact in the chat. It was execution-ready but not executed here.
\item The assistant converted this into a Codex prompt. The proposed deliverables were:
\item The user then provided the earlier documentation prompt and asked to optimize it to conform more strictly to international code documentation standards for the relevant languages and documentation formats. A stricter prompt was generated, aligned with ISO C89/C90, C++98, assembly, Doxygen-compatible structured comments, Markdown docs, and ASCII-safe source compatibility.
\item There was one known defect in that prompt: a typo line, Determininism guarantees remembered / Determinism guarantees remembered, depending on the preserved version. The report package later marked that as something to correct before reusing the prompt.
\item Later, the user said the project had been heavily refactored in the setup, launcher, and game design chats. They wanted a prompt to scan the repository, read the code, ensure every file in docs/ was consistent and up to date, and overhaul README in an industry-standard format that could serve as a landing page until a website exists.
\end{itemize}

\subsection{[Dominium Architecture I](../../\_reader/by\_chat/dominium\_architecture\_i.md)}

\begin{itemize}
\item The user also made a clear versioning decision. **FACT:** The user wanted **major.minor.patch semantic versioning** for all components/packages and for the complete game package. The user explicitly did **not** want build numbers or build dates in filenames. Instead, those details should live in metadata. This became part of the future build/package design.
\item The user then gave the "big task": create a set of .txt files, one for each file in the git directory, telling Codex how to implement it. These were not intended to be source files themselves. They were implementation-instruction files: requirements, prohibitions, dependencies, functions, declarations, and design constraints. This became the dominant workstream of the chat.
\item After the spec-generation work, the user changed the task again: the chat was being retired. The user asked for an OC-1 discovery inventory, then a maximum-fidelity Context Transfer Packet, then a downloadable package with Markdown/YAML reports, then an in-chat reader version of that package. Those handoff tasks did not continue project design; they preserved the state of the chat.
\item The conclusion was to create one .txt implementation-spec prompt per repository file. Each prompt would tell Codex how to implement that file without needing the entire conversation. This became the main practical output of the chat.
\item What remains uncertain is the current real-world status and capabilities of "Codex 5.1 Max." That is an external tool fact and requires verification before future use.
\item This is a formal packaging/build policy. It should feed into future build scripts, package manifests, release metadata, and possibly save/replay compatibility metadata.
\item The most important unresolved issue is duplication and contradiction. dweather, dhydro, and dai\_core were generated in more than one form. The package does not choose a final version. A future assistant should not implement any of those blindly.
\item The user explicitly wanted Codex to help with implementation. This goal drove the conversation toward file-specific implementation prompts.
\end{itemize}

\subsection{[Dominium Architecture II](../../\_reader/by\_chat/dominium\_architecture\_ii.md)}

\begin{itemize}
\item The last part of the chat was handoff extraction. The user asked for discovery inventory, structured registers, a context transfer packet, and then a downloadable package. The generated package captured the chat's workstreams, decisions, tasks, constraints, open questions, artifacts, risks, and verification items. This current report is a plain-language explanation of that substance.
\item This matters because it prevented future implementation from chasing arbitrary continuous geometry in the deterministic core. The user wanted arbitrary shapes, but also deterministic integer math and old-platform support. The compromise was to allow rich authoring tools while baking runtime state into deterministic structures.
\item The chat discussed fixed tick phases, virtual lanes, merge order, command buffers, event queues, and deterministic serialization. The simulation must not depend on render FPS, wall-clock time, OS scheduler behaviour, GPU state, or floating-point quirks. This explains many later decisions: C89 core, integer/fixed-point math, no sim floats, no platform headers in simulation, and renderer as pure observer.
\item What remains uncertain is the exact code-level API and component layouts. The chat produced specs and prompts, not final source code.
\item This was one of the most important conceptual outputs. It provides the future assistant with a mental model for everything else. Railways are graph/corridor systems. Buildings are microgrid frames plus archetypes and fixtures. Terrain is base fields plus delta fields. Power, data, and fluids are graph/field systems. Construction is blueprints and jobs. Space travel is an orbital graph/domain system.
\item The default world is Earth in Sol in the Milky Way. The base game MVP is one surface, but the architecture should support galaxies, systems, planets, and future server/shard arrangements.
\item The user also wanted buildings rigid for now but future collapse/destruction support. The conclusion was to include future-proof structural fields but not implement collapse in the MVP.
\item This is important for future hydrology, geology, stability, collapse, pathing, and undo/history.
\end{itemize}

\subsection{[Dominium Architecture III: Launcher, Platform, Renderer, and Handoff Architecture](../../\_reader/by\_chat/dominium\_architecture\_iii.md)}

\begin{itemize}
\item The assistant generated script templates and Codex prompts for these. Those prompts should be treated as implementation artifacts, not proof that code was actually changed.
\item These ideas were accepted as direction by continued user refinement, but implementation remains future work.
\item The user then asked for a Codex prompt to implement and document this. That means the keybinding map was accepted enough to become implementation material, but it should still be documented as a default binding specification rather than an immutable law.
\item The final phase of the chat was meta-work: extracting this chat into handoff forms. The user asked for OC-1 discovery inventory, then OC-2 registers, then a maximum-fidelity Context Transfer Packet, then a downloadable report package, then an in-chat reader version, and now this human-readable explanatory report.
\item Those outputs are important artifacts, but they are not the substance of the design. Their purpose is to preserve this chat's decisions, uncertainty, and future work for later aggregation into a master Project Spec Book.
\item This matters because future assistants must not implement launcher logic directly in Win32 widgets, Cocoa classes, SDL callbacks, or modern C++ code.
\item The user later reversed this. The final decision was to remove minimal versions and use full single implementations. This means the architecture should not split platform/render backends into minimal and full variants unless future technical constraints force it.
\item This remains a future implementation plan.
\end{itemize}

\subsection{[Dominium Architecture IV](../../\_reader/by\_chat/dominium\_architecture\_iv.md)}

\begin{itemize}
\item A number of Codex prompts were generated earlier in the chat for refactoring the repo, implementing runtime/platform/render APIs, setup, launcher, game bootstrap, ECS/SIM, and world/terrain/hydrology. However, the user later clarified that none of the Phase 2.5+ prompts had been implemented yet and wanted to pause before going further. The user wanted the systems to be more extensible and modular before committing.
\item The user then asked how to structure Codex prompts. The final high-level roadmap became:
\item The assistant generated detailed Phase 1 prompts, platform backend templates, renderer backend templates, setup prompts, launcher prompts, tools prompts, and finally a Phase 7+ roadmap for the game. The game phase itself was not expanded into the same full prompt detail as Phases 1-6. That remains a future task.
\item The topic mattered because every Codex prompt depends on paths and module boundaries. If the repo layout is inconsistent, prompts will fail or scatter code into the wrong places.
\item A platform backend prompt template was generated, along with parameter sets for Win32, Win16, DOS32, DOS16, CP/M-80, CP/M-86, Carbon, Cocoa, X11, Wayland, POSIX, SDL1, SDL2, and null.
\item The unresolved issue is the exact manifest/schema format. JSON, TOML, and INI-like examples appeared in prompts, but no final schema choice was made.
\item A full set of Phase 5 prompts was generated for launcher core, CLI, TUI, GUI, and extensions.
\item Phase 6 prompts were generated for the tool framework, core dev tools, editing backends, and GUI editor host.
\end{itemize}

\subsection{[Dominium Canon, Repository Alignment, and Portability Doctrine](../../\_reader/by\_chat/dominium\_complete\_conversation.md)}

\begin{itemize}
\item This chat was about preserving and re-grounding Dominium's architecture around a very old constitutional contract, then checking whether the current GitHub repository still follows that contract, and finally broadening the question into a platform-engineering doctrine for portability, modularity, extensibility, reuse, future-proofing, and long-term compatibility.
\item The conversation opened with a deliberate reconstitution of canon. The user supplied an old self-contained constitutional architecture prompt for Dominium and Domino. The assistant acknowledged it as authoritative within the conversation. The user then supplied a canonical glossary v1.0.0, and the assistant acknowledged that the glossary constrained future terms.
\item This preservation task was then uploaded as Pasted text.txt. It requested a maximum-fidelity report package for the current chat, with human-readable explanation first, structured registers, context-transfer packet, spec sheet, aggregator packet, self-audit, and downloadable files.
\item Uncertainty: the raw pasted prompt is not the current repo authority unless materialized in repo canon. The later audit found that it had in fact been materialized under docs/canon/constitution\_v1.md.
\item The glossary bound terms such as Authority, Law, Process, Lens, SessionSpec, UniverseIdentity, Macro Capsule, SRZ, RepoX, TestX, AuditX, CompatX, SecureX, and related terms. It matters because future assistants must not use sloppy synonyms like "mode" where the canon requires ExperienceProfile or LawProfile. This topic directly supports modularity because stable vocabulary is part of stable architecture.
\item The user then asked what practices would make the code reusable for other games and even other engine projects. The assistant answered that the correct goal is not to make all files permanent, but to make boundaries explicit and stable. The doctrine became: stable contracts, replaceable implementations, deterministic behavior, portable projections, and no accidental authority from paths/tools/UIs/prompts.
\item The uploaded prompt requested this full preservation package. It requires a human-readable report, registers, context transfer packet, spec sheet, aggregator packet, self-audit, exported files, and an in-chat reader.
\item The explicit goals were to restore/preserve old Dominium canon, assess current repository alignment, understand what the docs and code actually say/do, and identify practices to make the code portable, modular, extensible, replaceable, future-proof, and backward compatible.
\end{itemize}

\subsection{[Dominium + Domino Codex Planning and Handoff](../../\_reader/by\_chat/dominium\_domino\_codex\_planning.md)}

\begin{itemize}
\item The most important things to remember are: the project's strict architecture constraints, the Milestone-0 prompt sequence, the later shared CLI/TUI/GUI/input/render prompt sequences, the missing pack-system prompt, the unified startup policy, the unresolved repo-verification issue, and the need to treat generated prompts as plans rather than proof of implementation.
\item The user then asked for a plan going forward. The assistant produced a phased plan covering specification consolidation, directory/CMake layout, dsys, dgfx, UI core, DOMINIUM\_HOME metadata, launcher v0, game v0, tools, packs, retro backends, replay/networking, and documentation/examples.
\item INFERENCE:** The user accepted this direction in practice, because the next request was to generate Codex prompts to implement each step fully. The user did not explicitly say, "I formally approve the Milestone-0 plan," but they proceeded to ask for prompts based on it.
\item The user asked for six Codex prompts and reminded the assistant about supported numeric formats: U8, I8, U16, I16, Q4.12, U32, I32, q16.16, Q24.8, u64, i64, q48.16, and similar formats. That correction mattered because it broadened the numeric core beyond only Q16.16/Q24.8.
\item The assistant then generated six long prompts:
\item 1. **Prompt 1** created SPEC\_IDENTITY.md and SPEC\_PRODUCTS.md, documentation only.
\item 2. **Prompt 2** created the canonical repo layout, CMake skeleton, core headers, and product stubs.
\item 3. **Prompt 3** implemented deterministic numeric core types, fixed-point operations, RNG, and a numeric test harness.
\end{itemize}

\subsection{[Dominium Workbench, AIDE, Presentation Architecture, Provider Strategy, Product-Spine Planning, and Preservation Companion](../../\_reader/by\_chat/dominium\_full\_conversation.md)}

\begin{itemize}
\item A future assistant must understand that this chat is not simply about UI. It records the architecture for how Dominium should build itself: through contracts, commands, services, providers, documents, patches, modules, packs, apps, workspaces, diagnostics, evidence, AIDE, Codex, and eventually Workbench.
\item The conversation began with the user wanting a product designer/software architect/prompt engineer to help design an internal graphical UI creation tool. The immediate pain point was the Dominium Windows launcher UI: it looked mangled and flickered. The user wanted a tool that could visually design pixel-perfect UIs and broad-strokes functionality for setup, launcher, game, and tools.
\item A large prompt plan followed: repo scaffolding, canonical UI IR, TLV I/O, capability system, layout engine, splitter/tabs/scroll widgets, action codegen, Phase A UI Editor, Tool Editor bootstrap, Win32 batching/flicker fixes, tests, docs, and capability tests through launcher/setup redesigns.
\item The discussion then moved into AIDE. The user wanted to automate as much as possible through Codex and AIDE. We developed a task operating model: AIDE creates WorkUnits, tracks attempts, blockers, evidence, repairs, resumes, checkpoints, and promotion decisions. Codex executes bounded tasks. Development can continue with classified partials and repairable blockers; promotion to main requires evidence.
\item This led to generated prompts: STATUS-RECONCILE-02, AIDE-WORKFLOW-LAW-01, AIDE-WORKUNIT-SCHEMA-01, AIDE-DEV-MAIN-POLICY-01, AIDE-CHECKPOINT-LOOP-01, and AIDE-CAPABILITY-REALITY-LEDGER-01. The user wanted to run AIDE workflow, WorkUnit schema, and dev/main policy concurrently on the same machine via separate Codex tasks.
\item This prompted a provider structure: runtime/<service>/providers/<provider>, contracts/provider, contracts/capability, contracts/schema/runtime/<service>, release/profiles, content/profiles, and external/upstream. The plan explicitly rejected runtime/raylib/render, apps/client/rendered/raylib, contracts/raylib, top-level profiles, and top-level labs.
\item The final part of the conversation dealt with repo status and task queue. The user pasted status reports indicating that PRESENTATION-CONTRACT-01 completed with warnings, and then chose to generate six maintenance prompts before replanning. FULL-GATE-LEGACY-TEST-ROUTE-01 was generated. This preservation task followed.
\item The rendered GUI must be backend-agnostic. Layouts, controls, widgets, themes, styles, views, and workspaces should be modular and extensible. Themes include first-party primitive-only profiles and OEM+ mimic profiles. The GUI should be task-first, command-backed, evidence-visible, and projection-neutral.
\end{itemize}

\subsection{[Dominium Setup Architecture and Handoff](../../\_reader/by\_chat/dominium\_setup.md)}

\begin{itemize}
\item Near the end, the chat shifted into preservation mode. The user asked for a maximum-fidelity context transfer packet, then for a downloadable report package, then for an in-chat reader version of that package. Those outputs preserved the decisions, tasks, constraints, artifacts, rejected options, and verification queue for future assistants or aggregation into a larger Project Spec Book.
\item The first major input was a detailed master prompt for the **Dominium / Domino Setup \& Installation System**. It said this chat was about setup, installation, packaging, distribution, and uninstallation across supported platforms. It explicitly excluded launcher UI work. The setup system was described as a deterministic deployment control plane, not a simple GUI wizard or a store/DRM mechanism.
\item The assistant responded with **Phase 1: Setup Core Architecture**, proposing modules such as manifest parsing, resolution, planning, transaction execution, filesystem/path policy, installed-state, logging, and a platform interface. This was the first conceptual baseline.
\item The user then asked for a "plan of plans" called **Plan S**, with subplans S-1, S-2, S-3, and so on. The assistant produced a structured setup roadmap. It covered governance, setup core architecture, manifest schema, installed-state schema, transaction engine, filesystem policy, audit logging, CLI, Windows/macOS/Linux adapters, Steam integration, distribution pipeline, side-by-side upgrades, tests, and documentation.
\item At that stage, the plan was still mostly abstract and platform-oriented. It was useful because it separated the work into durable phases, but later repository changes made some of its directory assumptions obsolete.
\item The user later asked where to create Visual Studio and Xcode apps after opening the repo as a folder. The assistant advised that Visual Studio and Xcode projects should be generated through CMake, not hand-authored or treated as authoritative source files. This decision became consistent with later Codex prompts.
\item The user then gave a strict system-role prompt saying the repository had undergone a major structural refactor and that the filesystem was now canonical ground truth. This superseded earlier structures. The locked top-level products became:
\item This became a decisive turning point. Earlier setup/adapters, setup/packaging, and core/source structures became historical and superseded. The assistant amended Plan S to match the new canonical structure and later produced an authoritative Phase 2 directory tree under that model.
\end{itemize}

\subsection{[Domino Dominium Workbench](../../\_reader/by\_chat/domino\_dominium\_workbench.md)}

\begin{itemize}
\item The second central decision is that CLI, TUI, rendered GUI, OS-native GUI, and headless automation must be **projections of the same command/result/refusal/document/view spine**. A GUI button, a TUI hotkey, a CLI command, a headless JSON job, and a future OS-native widget action should route through the same command/service path and produce the same typed result, diagnostics, refusals, logs, and evidence.
\item The resulting old plan included canonical UI IR, deterministic TLV output, JSON mirrors, safe atomic saves, stable IDs, capability validation, layout engines, splitter/tabs/scroll widgets, action stubs, event dispatch, UI Editor GUI, Tool Editor bootstrap, import/export, CLI scripting, ops.json, and tests. Many Codex prompts were generated around those ideas.
\item The plan was then extended so the UI Editor could import existing tool UIs, provide a CLI suitable for Codex automation, and use scripted ops.json changes to create logical "Minecraft-style" launcher and setup layouts. The user clarified that "Minecraft-style" referred to layout and flow, not visual skin. All controls were to remain native Win32 controls in that earlier phase.
\item That phase generated more Codex prompts: UI discovery/import/export, CLI validate/format/codegen, ops.json editing, launcher UI generation, setup UI generation, and hardening/CI. These are now **historical/superseded** but still useful as implementation material for future document/patch/CLI authoring workflows.
\item Reason:** The old plan risked becoming a separate one-off tool architecture. Workbench can instead prove and reuse the same systems as the client, launcher, setup, server, and future tools.
\item The renderer does not know what a Validation Dashboard, Pack Browser, Windows XP theme, or Workbench module is. This keeps the renderer reusable by client, setup, launcher, server admin surfaces, Workbench, and future games.
\item 7. Use validation dashboard to monitor future work.
\item 1. Multiple old UI Editor Codex prompts and prompt plans.
\end{itemize}

\subsection{[Dominium/Domino Engine Refactor Planning](../../\_reader/by\_chat/domino\_engine\_refactor\_prompts.md)}

\begin{itemize}
\item The chat then produced a maximum-fidelity context transfer packet and later a downloadable report package. Those are artifacts of preservation, but the substance of the chat is the architecture itself: a deterministic, modular, data-driven engine plan that can scale from plants and wildlife to cities, power grids, spacecraft, fog of war, and future DLCs without violating strict determinism.
\item 4. deltas are applied later in a deterministic commit phase.
\item The answer also introduced incremental compilation from edit buffers into derived artifacts: occupancy, ports, connectivity, enclosures, surfaces, supports, nav hints, and boundary lists. This later became a key part of the TRANS/STRUCT/DECOR prompt plan.
\item The user introduced future **Interstellar** and **Wargames** DLCs. This could have led to hardcoded space and combat systems, but the answer deliberately generalized the requirements.
\item The answer produced a delta plan showing what was already satisfied and what needed to be added. Then the user asked for the full new prompt plan.
\item The assistant produced a full prompt plan for Codex. The user then requested each prompt one by one with "Next." Fourteen prompts were generated:
\item The user then asked for a documentation validation prompt to ensure the docs remain valid, consistent, and correct. That prompt was generated as well.
\item After the architecture and prompts were created, the user asked for a maximum-fidelity context transfer packet. The assistant produced one. The user then asked for a downloadable report package; the assistant generated package files and a ZIP. The user later asked for an in-chat reader version of that package, and finally requested this current human-readable explanatory report.
\end{itemize}

\subsection{[Domino/Dominium Engine Baseline, Architecture, and Feasibility](../../\_reader/by\_chat/engine\_baseline\_architecture.md)}

\begin{itemize}
\item The uploaded preservation prompt requires a maximum-fidelity human-readable handoff, structured registers, spec sheet, aggregator packet, self-audit, and downloadable file package for this chat. That prompt itself is an artifact of this chat and should be preserved. Source: ?filecite?turn44file0?
\item The conclusion changed: the project already had more infrastructure than the generic future plan assumed. The recommendation shifted from "invent the architecture" to "converge the architecture." The answer proposed that Dominium should become a **contract-compiled simulation platform**: contracts -> registries -> compiled locks -> capabilities -> Work IR -> deterministic runtime -> proof/replay -> product shells.
\item The final user action uploaded Pasted text.txt, a detailed instruction prompt requiring a full preservation report, structured registers, spec sheet, aggregator packet, self-audit, and downloadable file package for this chat. That request is the current task.
\item The central architecture distinction was that Domino should be a reusable deterministic simulation substrate, while Dominium should be one game/product layer built on it. This came up immediately when the user asked for a total description of the project. It mattered because the user wants the code to be reusable for future games and possibly other engine projects.
\item The conversation repeatedly returned to determinism. The project's value depends on identical inputs producing identical outputs, replay equivalence, stable commit ordering, named RNG streams, explicit Work IR/Access IR, and no hidden global scans. The assistant inspected execution model docs and source files such as task graph, access set, scheduler, and work queue.
\item The conclusion was that future work should build active-set runtime, capsule/refinement contracts, compatibility edition matrix, golden replay corpus, and toolchain-first authoring. Later, after reading live repo docs, this plan was tightened into "contract-compiled simulation platform."
\item This connects to future work because a construction/destruction game cannot rely on one universal representation. Different domains need different data structures.
\item how to make the engine portable, modular, reusable, and future-proof;
\end{itemize}

\subsection{[Dominium Foundation Lock, Workbench Spine, and Parallel Codex Handoff](../../\_reader/by\_chat/foundation\_workbench\_codex.md)}

\begin{itemize}
\item The topic came up because the project needed a model that could support world simulation, modding, tooling, Workbench, release, portability, and future games without repeating endless refactors. The conclusion was that stable contracts and semantic IDs must be separated from replaceable private implementation.
\item Workbench is a future product surface, not the center of authority. The user wants to get to Workbench and code soon, but Workbench must route through commands/services/views/evidence rather than private direct mutation. The first Workbench slice was narrow validation. The next needed bridge is command-result-view projection.
\item The user wanted to run many autonomous tasks concurrently. The assistant designed branch/worktree isolation and strict parallel-worker rules. Workers do not push to main, do not update global latest AIDE files, and produce task-local evidence. Coordinator merges serially and runs fast strict.
\item This topic matters operationally because the user wants speed. The new chat should be ready to generate more parallel prompts or a coordinator review prompt depending on current repo state.
\item This matters for the eventual master Project Spec Book and future domain feature work.
\item to make the repository tidy and future-proof,
\item minimize future architectural reversals;
\item make Codex prompts operational rather than theoretical;
\end{itemize}

\subsection{[Domino Framework and Open-Source Provider Architecture](../../\_reader/by\_chat/framework\_open\_source\_provider.md)}

\begin{itemize}
\item Finally, the user uploaded a detailed preservation prompt and requested a maximum-fidelity report, registers, spec sheet, aggregator packet, audit, and downloadable files. This package is the response to that request.
\item Future work should convert this into a dependency policy: use permissive/weak-copyleft libraries as providers where appropriate; use GPL/proprietary/unclear projects as research references unless explicitly quarantined; require provider manifests, license manifests, and conformance tests.
\item The user appears to want a long-lived, auditable project spec that can survive chat transitions and aggregation. The preservation prompt confirms this.
\item The user appears to want architecture that remains portable across older desktop systems and future render/platform providers. This was inferred from the platform floor and repeated concern with portability/modularity/extensibility.
\item The user requires source-grounded, audit-ready, uncertainty-labelled responses. The preservation prompt explicitly requires FACT/INFERENCE/UNCERTAIN/PROJECT-CONTEXT labels, no invented facts, no treating brainstorms as decisions, and no over-compression. The prompt also requires a human-readable report first and downloadable files if possible ?filecite?turn29file0?.
\item The user prefers modular architecture over monolithic engine forks. The user prefers long-term portability, extensibility, and replaceability. The user cares about preserving reasoning and rejected options for future aggregation.
\item The main uploaded artifact is Pasted text.txt, containing the preservation task prompt and output requirements ?filecite?turn29file0?. Before this task, no generated report/ZIP package was visible in this chat. This response creates the requested package.
\item A future assistant might incorrectly say "we decided to use raylib as the engine." The correct statement is that raylib is a first visible provider suite. Another common mistake would be to treat every assistant proposal as a user decision. Several items are recommended directions, not implemented decisions.
\end{itemize}

\subsection{[Dominium Content and GUI Rebuild Planning](../../\_reader/by\_chat/gui\_binary\_content.md)}

\begin{itemize}
\item The conversation opened with a large prompt titled **"PROMPT CONTENT0 - CANONICAL GAME CONTENT \& DATA POPULATION."** The target was GPT-5.2 Codex, and the scope was explicitly limited to **data, schema extensions, and documentation**. The user framed the assistant as continuing the Dominium / Domino project, but only at the content layer.
\item An assistant responded by turning the prompt into a cleaner, more formal Codex-ready version. That output was useful, but it was premature relative to the user's actual desired process.
\item The user then said: **"First, let's discuss this before we plan or generate any prompts."
\item That was an important change of direction. It meant the assistant's polished prompt should not be treated as final. The work needed conceptual discussion first. The assistant then analyzed CONTENT0 and identified several issues that could cause bad assumptions to become embedded if Codex acted too soon.
\item This stage established a pattern that continues throughout the chat: the user does not want impressive-looking prompts that hide unresolved design choices. The user wants the assumptions exposed before anything is formalized.
\item The user then described a larger goal: at the end of an "optimally large macro prompt plan," they wanted a **full universe** defined efficiently with minimal storage and memory. This universe system should work across a spectrum:
\item This was the most important decision in the GUI part of the chat. It changed the task from "can we make cross-platform GUIs?" to "how do we design a full multi-platform GUI and binary matrix without losing backend consistency?"
\item SwiftUI for 64-bit Intel and ARM future launcher and server, with graceful degradation to macOS 11.0.
\end{itemize}

\subsection{[Dominium Language, Platform, and Architecture Baseline](../../\_reader/by\_chat/language\_platform\_architecture.md)}

\begin{itemize}
\item The user then consolidated a broad future plan around Domino, Dominium, Workbench, AIDE, contracts, tests, replay, and evidence. The answer agreed that the plan was strong, but identified missing central pieces: composition resolver, lockfiles, compatibility corpus, trust/permissions, virtual filesystem roots, performance budgets, and stable public-surface promotion rules.
\item The chat repeatedly compared C89, C99, C11/C17, C23, C++98, C++11, C++17, and newer C++ standards. The final working direction is C17 + C++17 for the mainline. C99 and C++11 were considered but not adopted as the project-wide baseline. Newer C23/C++20/C++23/C++26 were treated as future provider/tool lanes, not current mainline law.
\item The project must retain a C-compatible ABI. The public boundary should remain POD-only, versioned, return-code/refusal based, and free of exceptions, STL containers, classes, templates, allocator ownership, and C++ ABI assumptions. This allows providers, plugins, tools, projections, and future bindings to survive implementation changes.
\item The deterministic simulation law remains more important than the language standard. Authoritative simulation must use stable IDs, canonical ordering, fixed-width values, fixed-point math, explicit little-endian encoding, and deterministic scheduler phases. Threads may accelerate derived work, but final authoritative commit must be canonical and not depend on OS timing or thread completion.
\item Support Windows 7, macOS Mavericks, Linux, and possibly older/future systems.
\item Keep future assistants from repeating rejected advice such as pure C99 or pure C++11.
\item Prepare material for a future Project Spec Book.
\item The goal changed from "should old C/C++ standards be upgraded?" to "what is the full future architecture baseline?" The user initially entertained C89/C++98 staging, then accepted C17/C++17 as mainline. The user also moved from "maybe support old systems directly" to "support old systems through constrained/projection/archive lanes."
\end{itemize}

\subsection{[Dominium Launcher Application-Layer Handoff](../../\_reader/by\_chat/launcher\_app\_layer.md)}

\begin{itemize}
\item Someone reading this report should understand one central thing: this chat is not about inventing new launcher features anymore. It is about preserving boundaries, making the launcher implementation explicit, and ensuring future work happens on top of verified code rather than vague architectural memory.
\item The user started by asking whether the latest system, after Plan 8, could support a common launcher for all operating systems with one shared codebase, native OS elements, and minimal rendering. The user attached several screenshots of older Minecraft launchers as visual examples and referenced a Plan 8 prompt file and a repo archive.
\item The user then asked how to make the entire system more modular and extensible, especially with existing systems, mods, and packs. The answer proposed a data-first approach: facades, registries, versioned contracts, TLV/schema-based manifests, declarative settings, pack tasks, and rare executable plugins only when unavoidable.
\item This phase mattered because it framed the launcher not as a menu program but as a **control plane** for installed products, packs, profiles, compatibility, and launch contracts. However, later canon tightened the permitted communication routes: cross-product communication must go through schema/ and libs/contracts, not arbitrary plugin conventions.
\item The response executed Phase 1 of that prompt and described the launcher core as a deterministic control-plane library. This included modules like manifest parsing, discovery, catalog building, profile policy, integrity checks, resolution, launch planning, execution, and audit logging. It also defined host service vtables, opaque handles, deterministic ordering rules, error models, and safety assumptions.
\item That Phase 1 answer remains useful, but only after being rebased into the later canonical repo structure.
\item This later became one of the clearest examples of a superseded path. The user subsequently pasted applied Codex prompts that explicitly required CMake to generate the Visual Studio solution/projects and stated that hand-written .vcxproj files must not be the authoritative build. Therefore, all earlier "manual IDE projects as canonical" advice is no longer current.
\item Later, the user objected that source subdirectories inside top-level source directories felt counterintuitive and asked about future tools and components beyond launcher, setup, and game. This led to the product-first model: top-level directories should represent long-lived products, not layers.
\end{itemize}

\subsection{[Dominium Launcher and Setup Architecture](../../\_reader/by\_chat/launcher\_setup\_architecture.md)}

\begin{itemize}
\item The chat also produced multiple Codex work-order prompts and finally a downloadable report package. Those artifacts are useful, but the substance is the architecture: keep engine deterministic, keep setup/launcher/runtime boundaries explicit, make the launcher optional and modular, use a strong process/instance model, and now implement the launcher over dsys/dgfx rather than ad hoc platform UI stacks.
\item The user asked how many prompts would be needed for Codex to implement the system. The assistant recommended four main prompts:
\item An optional fifth prompt was suggested for runtime CLI/capabilities. The user then asked for the prompts one by one, and the assistant generated them. These prompts were detailed and implementation-oriented, but they reflected the earlier C++98/shared-library approach.
\item This is the most important change of direction in the chat. Earlier prompts remain useful historical material, but the latest launcher implementation direction is dsys/dgfx.
\item This came up because the user wanted Codex and future implementation work to avoid hidden coupling and future incompatibilities. The assistant expanded the idea into repo layout, file format patterns, CLI contracts, plugin ABIs, and migration/testing expectations.
\item Uncertainty:** Accounts/social/wiki/forum/direct messaging are future-facing and were not converted into the final five-tab implementation plan.
\item This was the late major shift. The user provided a detailed new prompt stating that all dsys and dgfx backends exist and that the launcher should use them.
\item INFERENCE:** The user wanted a future-proof architecture that Codex could implement in stages without repeatedly re-deciding the system design.
\end{itemize}

\subsection{[Dominium Modularity, AIDE Refactorability, and Future-Proof Architecture](../../\_reader/by\_chat/modularity\_aide\_refactorability.md)}

\begin{itemize}
\item The final user action was uploading a preservation-package prompt. That prompt requested a maximum-fidelity preservation report, structured registers, context transfer packet, spec sheet, aggregator packet, self-audit, downloadable files, and a final in-chat reader. This package is the result.
\item The user objected to the XStack/AuditX/RepoX/TestX framing and wanted AIDE adopted quickly. The resulting plan made AIDE the repo-native control plane for inventory, task queues, policies, move maps, salvage maps, evidence ledgers, validation, and refactor history. Existing tools should be recycled, not discarded. This is central to future refactorability.
\item The uploaded prompt converted the conversation into a preservation task. It requires a human-readable report first, then registers, spec sheet, aggregator packet, self-audit, and downloadable files. This final package is intended to be merged later with old-chat reports into a master Project Spec Book.
\item The conversation implies a goal of turning Dominium into a long-lived product-line platform rather than a single-game repository. It also implies a desire for auditability: decisions should be grounded, visible, mechanically enforceable, and later aggregatable into a master spec. The user appears to value stable boundaries, explicit compatibility policy, and practical Codex/AIDE tasks.
\item The live repo still needs verification. The exact AIDE files, contracts, validators, and migration tasks must be implemented. It remains unresolved which existing roots and tools map to which final homes, which APIs become stable, which code is reusable across products/games/projects, and which prior assistant recommendations should become formal requirements after user review.
\item 1. **TASK-09** - Verify live repo baseline: current branch, root inventory, validators, CTest status, generated artifacts.
\item 2. **TASK-01** - Implement AIDE-STRUCTURE-00: root constitution, ownership slots, AIDE refactor framework, move/salvage map schemas, inventory tooling.
\item 3. **TASK-03** - Inventory old XStack/AuditX/RepoX/TestX-style tools and classify them.
\end{itemize}

\subsection{[Dominium Omega/Xi/Pi Architecture \& Future-Proofing Planning](../../\_reader/by\_chat/omega\_xi\_pi\_architecture\_future.md)}

\begin{itemize}
\item After Xi, the assistant planned Pi-series as a meta-blueprint: architecture diagrams, dependency maps, capability ladders, readiness matrices, prompt inventories, and snapshot mapping templates. The user reported Pi-0/Pi-1/Pi-2 were completed and restarted from Xi-8 ground truth, with fingerprints and prompt counts.
\item The final strategic direction before this preservation request was to run a ?-series: snapshot intake, reality extraction, blueprint reconciliation, foundation readiness, and final prompt synthesis. This is needed because plans must now be mapped onto current repo reality rather than executed abstractly. The next chat should pick up there.
\item The conclusion was that all products should boot through AppShell, use shared command registries, expose descriptors, validate packs/contracts before session start, and never have ad hoc boot paths. This matters for portability, setup/launcher integration, and future distribution.
\item The user later reported Xi completion. The enduring lesson is that prompt instructions are not enough; architecture must be machine-readable and enforced.
\item The user wanted agent support to work across vendors. The answer was to define a canonical vendor-neutral layer and generate vendor-specific mirrors. Canonical files include AGENTS.md, .agentignore, data/agents/agent\_context.json, and docs/agents/AGENT\_TASKS.md; mirrors may include Copilot instructions, Claude agent files, Codex skills/MCP wrappers. XStack, not agent prose, enforces truth.
\item 4. Preserve all plans, tasks, constraints, risks, decisions, artifacts, and future directions across chats.
\item 7. Support future live operations such as hot-swappable renderers and live mod activation, but only after foundations.
\item 3. Make the repo robust to future tool/vendor changes.
\end{itemize}

\subsection{[Dominium OS-like Architecture, Repository Convergence, and Interface Operating Layer](../../\_reader/by\_chat/os\_interface\_repo\_architecture.md)}

\begin{itemize}
\item The central tradeoff was between short-term visible functionality and long-term architectural leverage. The user repeatedly asked whether the plan could be more modular, future-proof, optimized, robust, reliable, useful, and powerful. The answer consistently favored building shared contracts and proof paths before expanding product surfaces.
\item This preservation task specifically required FACT / INFERENCE / UNCERTAIN / PROJECT-CONTEXT labels, stable IDs, registers, spec sheets, an aggregator packet, self-audit, and downloadable files if tools are available.
\item Future facts involving platforms, software support, toolchains, APIs, laws, schedules, prices, current roles, etc. require verification.
\item INFERENCE: The user wants future assistants to preserve reasoning and not collapse tentative ideas into final decisions.
\item 1. **Uploaded Pasted text.txt.** Contains the preservation/export prompt that generated this package.
\item 2. **Repo docs inspected during chat.** Key examples: AppShell constitution, GUI baseline, repo layout target, ownership rules, domain split rules, final convergence audit, post-CONVERGE next steps, versioning/distribution docs, code/data report, modularity proof, product boot proof.
\item 4. **Conceptual artifacts produced in chat.** These include proposed architectures, matrices, lane plans, Workbench modules, Interface Operating Layer, boot-to-replay MVP, and structured roadmaps.
\item 5. **This preservation package.** Newly generated files and ZIP at the end of this task.
\end{itemize}

\subsection{[Dominium Codex Platform Renderer API Plan](../../\_reader/by\_chat/platform\_renderer\_api\_plan.md)}

\begin{itemize}
\item Later, the user expanded the horizon: Dominium should start with simple early graphics and audio, but eventually support AAA-quality graphics, spatial audio, environmental acoustics, proximity chat, AI, navigation, terrain, vehicles, economy, combat, weather, modding, tools, and launcher/setup systems. Those ideas were preserved as a **future roadmap**, not current implementation scope.
\item The chat ended by generating a maximum-fidelity transfer packet, then a downloadable report package, then an in-chat reader version. The main thing to remember is this: **the active artifact is the final 9-prompt post-master Codex plan, but it is a plan, not evidence that the code exists. The repo must be inspected before execution.
\item The user opened with a precise role assignment: the assistant was to act as a staff software architect, technical program manager, and prompt engineer. The task was to generate **ready-to-run prompts** for Codex CLI running GPT-5.2 so that Codex could implement missing or unfinished platforms, renderers, and APIs in the existing Dominium repository.
\item The first procedural rule was that the assistant had to ask a **minimal but complete questionnaire** before generating the prompt pack. The questionnaire had to cover platforms, renderers, APIs, definition of done, and repo/build context. The assistant did that.
\item This is not a fatal problem for the prompt pack, because the prompts themselves repeatedly tell Codex to discover files with rg, tree, type, CMake commands, and target listing. But it is a major caveat for any future assistant: do not assume the actual repo matches every file path or API name unless checked.
\item After the archive upload, an initial prompt plan was generated for platform and renderer work. It covered CMake platform/backend selection, Win32 and POSIX system backends, X11/Wayland/Cocoa gating, software/DX9/DX11/GL2/Metal/Vulkan renderers, sockets/dynlib, and integration. That plan was useful but did not remain active.
\item The user then said they already had a set of pending prompts and pasted a **MASTER CODEX PROMPT PLAN**. This changed the sequence. The assistant was no longer planning from scratch; it had to generate a new plan that assumed this prior master plan had already been implemented.
\item The user's master prompt plan was extensive. It covered deterministic simulation and engine-core architecture: invariants, packet ABI, scheduler, fields/events/messages, LOD, pose and anchors, TRANS corridors, STRUCT buildings and enclosures, DECOR signage and surface detail, agents/sensors/minds/actions, graph toolkit, domains/frames/propagators, knowledge/fog/visibility/comms, and determinism hardening.
\end{itemize}

\subsection{[Dominium Platform Support Planning](../../\_reader/by\_chat/platform\_support.md)}

\begin{itemize}
\item UNCERTAIN / UNVERIFIED: The name "Domino" was never confirmed by the user. Future work should treat "Domino" as an assistant-created placeholder unless the user explicitly accepts it. The useful idea is the distinction between full product support and reduced engine/core support, not necessarily the name.
\item FACT: The user explicitly decided that PC is the first primary focus and that Android, iOS, and Web are also primary top-tier support. This means future work should treat these platforms as architectural roots. They are not optional ports.
\item This topic connects directly to future work because the next major deliverable should be a **Platform Support Matrix** and then **Platform Contracts** for each major target.
\item UNCERTAIN / UNVERIFIED: The chat did not establish minimum Android API level, ABI list, target graphics APIs, Play Store rules, Android TV status, Android Go status, Automotive status, Chromebook status, or vendor ROM support. These need future decisions and current-source verification.
\item UNCERTAIN / UNVERIFIED: Current console platform facts were not verified in this chat. Future work must verify official PlayStation, Xbox, and Nintendo developer requirements from current official sources before implementation planning.
\item UNCERTAIN / UNVERIFIED: XR was never promoted to top-tier. Its priority remains unresolved. Future work should not let XR quietly expand the scope of the main product unless the user explicitly elevates it.
\item The latter part of the chat created a Context Transfer Packet and then a downloadable package containing reports, registers, YAML, an aggregator packet, an audit, a manifest, and a ZIP. The user then asked to read the package in-chat. That package is useful because it preserved IDs and future aggregation material.
\item This current report should not be about those files, but they are artifacts of the chat. Their purpose was continuity: to allow another assistant or future aggregation process to reconstruct this chat without re-reading everything.
\end{itemize}

\subsection{[Domino/Dominium Portability, Assurance, and Future-Proof Architecture](../../\_reader/by\_chat/portability\_assurance\_future\_proof.md)}

\begin{itemize}
\item Limitations: this package preserves visible current-chat context and the uploaded preservation prompt. It does not prove that hidden transcript segments, other conversations, repository files, or older generated artifacts were accessible. Broader project context is not treated as current-chat fact unless labelled.
\item The user then clarified that the real concern was broader: all code should be portable, modular, extensible, reusable, replaceable, and future-proof. This shifted the conversation from standards to the whole structure of the engine/product ecosystem.
\item The user then uploaded a maximum-fidelity preservation prompt. The task became preserving this chat into a package suitable for human reading and future aggregation. This response is that package.
\item A target tree was proposed: include/, contracts/, source/domino/, source/dominium/, tests/conformance/, tests/migration/, tests/fuzz/, tools/validators/, docs/architecture/, docs/assurance/. This is a candidate, not a repo-verified plan.
\item Conformance tests prove replaceability. High-trust paths need valid, invalid, malformed, old-version, future-version, migration, roundtrip, determinism, and negative-permission tests.
\item The uploaded prompt required this preservation report, registers, spec sheet, aggregator packet, audit, files, and ZIP.
\item The user wanted to know whether standards can be used wisely for Domino/Dominium. The user also explicitly wanted portability, modularity, extensibility, reuse, replaceability, proper-game/OS-grade structure, better directories, better names, better APIs, better schemas, and future-proof/backward-compatible design. The uploaded prompt explicitly requested a maximum-fidelity preservation package.
\item The user likely wants these outputs to feed a future master Project Spec Book and prevent long-chat context loss.
\end{itemize}

\subsection{[Dominium README Architecture, Ports, and Determinism](../../\_reader/by\_chat/readme\_ports\_determinism.md)}

\begin{itemize}
\item The user then asked for a prompt they could give to Codex. A detailed patch prompt was produced. It instructed Codex to apply minimal edits to the existing README without changing headings, section order, or tone.
\item That prompt asked Codex to:
\item The user pasted Codex's output after the port prompt. The output correctly added the unified-source and metadata-only port language, but it also introduced a duplicated top section: two overlapping "This repository includes" blocks. The assistant identified that duplication and produced a cleanup prompt.
\item That cleanup prompt instructed Codex to remove the redundant block, remove undefined "embedded" from the future systems list, replace vague "lockstep/rollback" wording with lockstep-first canonical networking language, and update the contributing determinism bullet to refer back to Section 2.1 instead of maintaining a separate formulation.
\item After the README work, the user asked for a maximum-fidelity context transfer packet. The assistant produced a long state-transfer document with sections for project inventory, decisions, tasks, constraints, open questions, artifacts, risks, and next actions.
\item What remains uncertain is whether the actual repository README.md matches the final pasted version. The chat only saw pasted text and generated prompts; it did not inspect a Git repository.
\item The project's core promise is deterministic simulation. The README already required fixed-point arithmetic and deterministic tick phases, but the chat refined the boundary of that determinism.
\item The chat also strengthened RNG discipline. RNG streams can only advance during deterministic tick phases. Debug overlays, UI, rendering, and other non-simulation layers must not advance engine RNG streams.
\end{itemize}

\subsection{[Dominium + Domino Refactor Architecture](../../\_reader/by\_chat/refactor\_architecture.md)}

\begin{itemize}
\item This chat was mainly about turning the Dominium + Domino project from a collection of partially overlapping architecture ideas into a coherent, refactorable, future-proof structure that could be implemented by Codex and then preserved for future ChatGPT conversations.
\item The most important thing to remember is that this chat did not implement the refactor. It designed the architecture and generated prompts and handoff material. Any future assistant must verify actual repository state before assuming files were moved, prompts were applied, or code was updated.
\item The user opened with a long "new thread" prompt framing this as the **Graphics / UI / UX / Packs Architecture Thread** for Dominium + Domino. The explicit scope was visual and audio presentation: rendering architecture, UI framework, graphics packs, sound packs, music packs, theming, skinning, pack composition, overrides, fallbacks, and how all of that should map onto existing dsys and dgfx.
\item FACT:** This early discussion created the future UI/packs design context.
\item The assistant then produced a master Codex refactor prompt that instructed Codex how to move files, add compat/platform headers, add product metadata and --introspect-json, unify Game modes, implement DOMINIUM\_HOME, actions, setup import/GC, packaging naming, and docs updates.
\item The user asked what Codex still did not know. The assistant audited missing topics. Then the user asked for phases from scratch; the assistant produced a phased roadmap. Then the user asked whether all phases could be merged into one big prompt; the assistant generated a larger combined refactoring prompt. Then the user requested a post-refactor consistency prompt; the assistant produced one.
\item The user then shifted from architecture implementation to preservation. They asked for prompts that would cause Codex to generate starter prompts for future ChatGPT conversations, including an extended master starter prompt. Finally, they asked for a maximum-fidelity context transfer packet, then for a downloadable report package, then for an in-chat readable version of that package.
\item The final phase of the chat was therefore about preserving and explaining the architecture so it could be used later without rereading the entire conversation.
\end{itemize}

\subsection{[AIDE, XStack, and Dominium Refactor Control Plane](../../\_reader/by\_chat/refactor\_control\_plane.md)}

\begin{itemize}
\item The chat opened with the user describing a workflow: paying for Codex Pro, using the Codex VS Code extension on one project at a time, and using ChatGPT project spaces to generate prompts that were manually fed into Codex. The user had also installed the ChatGPT and Codex apps and Visual Studio Community. The initial question was whether there was a better, more professional, more integrated way to develop projects.
\item The user then asked about optimizing prompts for GPT-5.3/GPT-5.4 and later supplied advice about harness engineering. That advice argued that teams succeed by engineering the environment around agents: tools, docs, linters, feedback loops, memory, and observability. The conversation accepted the core point: the model is not the whole system; the harness is the multiplier.
\item The discussion then mapped these ideas onto Dominium's XStack. We considered adding HarnessX, XPlan, XExec, MergeX, ContextX, SkillX, TaskX, BridgeX, SessionX, PermitX, DoctorX, TraceX, and more. Over time this became too many public concepts. We simplified toward higher-level planes and then eventually concluded that XStack should not be the general product at all.
\item The user supplied AIDE advice around dev branches, structured commits, and WorkUnit recovery. The accepted model is main as canonical truth, dev as governed integration, task/* as bounded work, and release/hotfix branches as needed. AIDE should enforce commit trailers, branch roles, safe sync/land/promote/prune, and repo-state-first recovery when prompts are repeated or out of order.
\item The user explicitly wanted to improve AI-assisted development, design a better harness, decide XStack's role, create AIDE, align Dominium cleanup with AIDE, use AIDE in Dominium safely, and generate a complete preservation package. These goals mattered because the user is managing long-running complex projects and wants future chats, Codex sessions, and repos to stop losing context or repeating mistakes.
\item AIDE changed most. It began as a possible full development environment and became a portable repo operating layer. XStack changed from a candidate universal system to a Dominium strict profile. Repo cleanup changed from moving folders to root recycling under AIDE. AI workflow changed from prompt optimization to harness engineering.
\item Several decisions are strategic rather than fully implemented. The version/capability model, Git workflow model, install/repair/upgrade model, and tool absorption model are accepted directions but require schemas and code. Future assistants must not mistake them for already completed implementations.
\item The user values direct, source-grounded, audit-ready, detailed answers. The user asked for preservation packages that are human-readable and machine-aggregatable. The user repeatedly prefers not to lose nuance, rejected options, uncertainties, or artifacts. The user wants current facts verified when possible and wants future assistants to avoid re-asking answered questions.
\end{itemize}

\subsection{[Dominium XStack Release Identity and Versioning](../../\_reader/by\_chat/release\_identity\_and\_versioning.md)}

\begin{itemize}
\item The chat rebuilt SemVer from first principles and decided that strict SemVer should apply only to components with declared public contracts. Candidate true-SemVer entities include SDKs, engine libraries, tool APIs/CLIs, protocol libraries, plugin hosts, schema libraries, and reusable runtime libraries. It remains unresolved which exact repo modules qualify.
\item The user explicitly wanted a versioning/release identity policy that would not need to change halfway through. They wanted to avoid SemVer stagnation, arbitrary major bumps, and product-history confusion. Later, the user wanted the conversation reconstructed into a knowledge base and preservation package for future aggregation.
\item These rejected or superseded options matter because future assistants may otherwise reintroduce them. The biggest trap is suggesting "just use SemVer everywhere." That was explicitly narrowed. Another trap is treating stable or hotfix as prerelease suffixes even after SemVer ordering was explained.
\item The immediate future work is to convert this design into formal repo/spec artifacts. Recommended order:
\item The user values long-term maintainability over short-term prettiness. They prefer explicit metadata over hidden inference. They want future assistants to preserve tentative status and not silently turn brainstorms into decisions.
\item This preservation task creates the first actual downloadable package in this chat. Prior artifacts were primarily in-chat examples, summaries, and proposed spec fragments.
\item The highest-risk misunderstanding is collapsing the layered model back into one universal version scheme. Future assistants should classify the entity first, then apply the correct policy.
\end{itemize}

\subsection{[Dominium TestX/RepoX Governance and Handoff Chat](../../\_reader/by\_chat/testx\_repox\_governance.md)}

\begin{itemize}
\item That was the initial design phase. At that point, the build-number idea was still broad and aggressive: a global build number incrementing frequently.
\item The user had another chat working on content and systems, and asked for a prompt to inform that chat of everything decided so far. A similar prompt was generated for the applications/platforms/renderers chat. These prompts established authoritative boundaries:
\item The next major shift came when the user pasted an authoritative prompt called **TESTX** from "The Game chat." TESTX was a major governance layer: permanent verification, versioning, self-defending tests, deterministic execution, repository survey, test taxonomy, assertion tiers, comment density, blocker tracking, and changelog governance.
\item This chat then generated new prompts, EG-TESTX and AP-TESTX, to send to engine/game/content and application/platform/render teams. Those prompts required response prompts from those teams. The user pasted back their responses. Both accepted TESTX canon and identified tensions.
\item This chat endorsed that approach and suggested making control systems capability-gated, disabled by default, and subject to a non-interference invariant. The user then asked for a sequel mega-prompt, and **TESTX2** was generated. TESTX2 formalized:
\item This mattered because the user wanted future support for such systems without enabling them by default or compromising core principles.
\item Based on the tree, the response concluded that the repo already had many ingredients: UI IR/codegen, setup/launcher tooling, tests, validation scripts, platform layers, packaging, and legacy/orphaned directories. What it lacked was an explicit projection layer and enforcement model. The proposed REPOX prompt introduced:
\item The user then supplied a detailed matrix for Windows, Mac, and Linux, including language modes and old/new architecture goals. REPOX was generated as a mega-prompt.
\end{itemize}

\subsection{[Dominium Timekeeping and 2038 Resilience](../../\_reader/by\_chat/timekeeping\_and\_2038\_resilience.md)}

\begin{itemize}
\item The preservation task at the end changes the purpose of the chat from discussion to archive. The user uploaded a detailed preservation prompt requiring a human-readable report, structured registers, a context-transfer packet, a YAML-style spec sheet, an aggregator packet, audit sections, and downloadable files. This report is therefore both a reader-friendly reconstruction and a future-spec handoff.
\item A future assistant should understand that this chat contributes a time architecture doctrine for Dominium: **ACT is authority; DSYS time is runtime-only; observer clocks are derived; civil/astronomical time is projection-only; wall-clock time must never drive authoritative ordering.
\item The user then asked for the "most compatible" combination of timekeeping for cross-platform software architecture, including past and future target machines. The assistant answered with a layered model: absolute instant, duration, monotonic elapsed time, and civil/local time. This was the main conceptual upgrade of the chat. It broadened the issue from integer width to semantic separation.
\item The user then uploaded a detailed prompt requiring this preservation package. The task became a maximum-fidelity archive and handoff for the current chat.
\item The user should remember that unsigned 32-bit is not a future-proof absolute timestamp format. It only delays the overflow.
\item The major future work is not to invent the model from scratch, but to harden boundaries and serialization.
\item The final topic is this export task. The user requested a human-readable report first, followed by registers, spec prep, context-transfer packet, aggregator packet, self-audit, and files. This package is designed to let the chat be read later without reopening the full transcript and to support merger into a larger Dominium spec book.
\item The user appears to be trying to prevent long-term architecture mistakes in Dominium and related C/C++ software. The user values deterministic simulation, legacy compatibility, future-proof serialization, and avoiding subtle date/time bugs that would require painful migrations later.
\end{itemize}

\subsection{[Dominium UE6, Domino, and Deterministic Universe Feasibility](../../\_reader/by\_chat/ue6\_domino\_deterministic\_universe.md)}

\begin{itemize}
\item FACT: This package preserves the currently visible chat about whether Dominium should use UE6, UE5, Domino, Unreal, or a custom engine layer for a highly ambitious deterministic solar-system-scale game. It also preserves the current uploaded prompt as an artifact because it defines this preservation task.
\item For a future assistant or human reader, the key thing to understand is that the chat did not reject Unreal. It rejected the idea that Unreal or UE6 could be the complete solution. The direction is hybrid: custom simulation and backend as the real engine, Unreal as the first high-quality client, UE6 as a possible future upgrade, and Domino as a possible later frontend/runtime if the core remains portable.
\item The answer then proposed the architecture that became the main deliverable of the chat: DominiumSim, DominiumWorldDB, DominiumServer, DominiumClient\_UE, and DominiumClient\_Domino. This architecture makes the core simulation, world storage, and server authority independent from the renderer. Unreal becomes an adapter, not the source of truth. Domino becomes a future adapter if the project requires it.
\item The first topic was whether Dominium should be made on UE6 instead of Domino. The answer separated availability from suitability. UE6 was treated as not yet available enough for production planning. UE5 was treated as a practical near-term renderer/tooling base. Domino was treated as a possible future runtime or custom engine path, not something that can be reached by automatically exporting an Unreal project.
\item Uncertainty: public UE6 technical capabilities and requirements remain time-sensitive and need verification before any future hard commitment.
\item The assistant recommended against starting directly on UE6 because public technical access and production documentation were not treated as available in the visible chat. This was not framed as a permanent rejection of UE6. It was a timing decision: do not wait for a future engine to solve current architecture problems.
\item This is the highest-value decision candidate. Core rules, simulation, save format, economy, factory logic, fog-of-war state, and replay/hash systems should be portable. This matters because Domino portability is only plausible if the core does not depend on Unreal object lifecycles and assets.
\item The chat recommended making Domino a future adapter target. Portability comes from architecture, not from a future conversion button.
\end{itemize}

\subsection{[Dominium UI Editor and Tool Editor Planning](../../\_reader/by\_chat/ui\_editor\_tool\_editor\_planning.md)}

\begin{itemize}
\item The most important thing to remember is that this chat was a planning and prompt-generation chat, not proof of implementation. **UNCERTAIN / UNVERIFIED:** No generated prompt is evidence that repository code was actually changed. The next assistant must first verify whether the prompts were executed, inspect the repo/source bundle if available, and avoid treating planned files as existing files.
\item Once the architecture was clear, the user asked for a prompt plan that could implement everything in fewer than a dozen Codex prompts. The assistant created an 11-prompt implementation sequence:
\item The user then asked "Next" repeatedly, and the assistant generated each of those prompts in detail. These prompts are important artifacts, but they are not evidence that implementation occurred.
\item This led to a second prompt plan: UU1 through UU6. These prompts extended the UI Editor with discovery/import/export, CLI modes, ops.json scripted editing, launcher UI generation, setup UI generation, and determinism/CI hardening.
\item Before the launcher/setup prompts, the user clarified that "Minecraft style" meant logical layout only, not graphical design. This was crucial. The assistant had to avoid custom styling, skins, or owner-drawn controls. The launcher and setup tests were about structure: header areas, side navigation, main content tabs, news panels, wizard steps, progress pages, footer buttons, and similar interaction flow.
\item The user uploaded screenshot bundles for setup and launcher. **UNCERTAIN / UNVERIFIED:** The screenshots were not inspected in this chat. The assistant still produced a logical layout description, and the user accepted it as useful. Future work should inspect the bundles before claiming exact screenshot fidelity.
\item The final state of the conversation is therefore not "the tool is built," but "the architecture, constraints, prompt plans, and continuation state have been carefully documented."
\item The discussion also identified likely flicker causes. The Win32 backend is mostly child HWND controls and does not use custom painting. Flicker likely comes from resize relayout, redraw calls, list refresh invalidation, and visibility toggles. This later led to a planned prompt for structural Win32 batching and coalesced relayout.
\end{itemize}

\subsection{[Dominium Universe Explorer Planning and Repo Handoff](../../\_reader/by\_chat/universe\_explorer\_planning.md)}

\begin{itemize}
\item This changed the future plan from "eventually implement everything" to "first prove seamless lawful inspection/materialization at universe scale."
\item The current preservation task then asked for this full chat package.
\item A major theme was that useful local inventions must become portable, standardized, industrialized, and institutionally adopted. The repo now has a Formalization Chain spec and Player Desire Acceptance Map that strongly preserve this. The future work is making this playable: drafting, measuring, testing, blueprinting, certifying, teaching, manufacturing, maintaining, and revising designs.
\item The final plan accepted the newer repo status: structure is clean enough to stop giant structure prompts. Remaining cleanup should be targeted: stale full-gate tests, pack layout canon, residual taxonomy, AIDE state classification, public header ABI, provider split. This is a planning recommendation, not yet an explicitly accepted user decision.
\item 6. Preserve this chat in a maximal-fidelity package for future aggregation.
\item 3. Ensure future assistants do not confuse repo truth with chat memory.
\item 3. Whether the Universe Explorer has a formal repo task yet.
\item 6. Broad structure refactor prompts: now deprioritised because current structure is clean enough with warnings.
\end{itemize}

\subsection{[Dominium Architecture, Workbench, AIDE, and Product-Spine Planning](../../\_reader/by\_chat/workbench\_aide\_product\_spine.md)}

\begin{itemize}
\item The largest uncertainty is not the conceptual architecture; that has converged. The largest uncertainty is live execution state: whether prompts generated near the end of the chat have already been run locally, committed, or pushed. The next chat should first inspect .aide/queue/current.toml, recent commits, and relevant audit files.
\item In the later chat, the user requested a series of Codex/AIDE prompts. The assistant generated prompts for:
\item The user noted that earlier prompts were hard to copy because they were not always contained in one code block, and that prompts should better handle dirty worktrees and concurrent tasks. From then on, generated prompts included detailed dirty worktree handling, allowed/forbidden paths, non-goals, validation, blocker classification, and commit/final-response formats.
\item Near the end, the user wanted to run prompts in parallel on a permanent origin/dev and local/dev, then checkpoint/merge to main. The conversation converged on the AIDE OS model: agents should not mutate shared dev directly; each prompt should run in a task branch/worktree; AIDE integrates to dev; checkpoint branches prove waves; main receives only evidence-backed promotions.
\item Remaining uncertainty: exact implementation of erosion/ecology/evolution proxies remains future domain work.
\item The conversation treated epistemics and diegetics as more than flavor. They determine what must be simulated, loaded, rendered, or refined. If a player has not observed something, it can remain summarized so long as future expansion is consistent with seed, law, and causal history.
\item AIDE should evolve from prompt generator into task operating system:
\item The user needed copyable prompts in a single code block. Future prompts should include:
\end{itemize}

\subsection{[Dominium World Architecture](../../\_reader/by\_chat/world\_architecture.md)}

\begin{itemize}
\item The future relevance of this chat is high. It should feed directly into the future project spec book and into a corrected implementation prompt. But it should not be treated as final implementation detail everywhere. Many high-level decisions are settled, but solver formulas, exact file encoding details, build system integration, actual repository layout, and some numerical representations still require verification.
\item Near the end, the user asked for a Codex 5.1 implementation prompt. A long prompt was produced, instructing Codex to implement the engine core, fixed-point types, world addressing, chunking, save formats, ECS, engine API, runtime CLI, and docs. Later, the context packet and final report package audited that prompt and found important defects. The prompt is useful, but not safe to use directly.
\item The main unresolved issue is implementation detail. The design says Q16.16 and Q4.12, but the earlier prompt used signed types incorrectly. The future implementation must use unsigned local coordinates or a centered signed design.
\item A subtle but important correction emerged later: within a Segment of 2\textasciicircum{}16m, a Region of 2\textasciicircum{}8m means there are 256 Regions per Segment per axis. Therefore Region-in-Segment is 8 bits, not 4 bits. Some earlier generated prompt text got this wrong and must be corrected.
\item The exact solvers remain unresolved. The architecture is clear, but formulas and resolutions are still future work.
\item The user wanted the same future-proof philosophy applied to the entire project. The resulting pattern was:
\item A major artifact was a Codex 5.1 prompt. It was designed to tell Codex to implement the architecture in a repo. It included modules, headers, save formats, engine API, CLI runtime, and docs.
\item Later audit found it defective. It should be repaired before use. The known issues are signed fixed-point types, Region bit-width inconsistency, C89-incompatible constructs, and saved rotation ambiguity. The best next practical step is to revise that prompt using the corrected report.
\end{itemize}

\subsection{[Dominium XStack and Lab Galaxy Handoff](../../\_reader/by\_chat/xstack\_lab\_galaxy.md)}

\begin{itemize}
\item That problem drove much of the XStack discussion. The assistant and user discussed how to prevent future agents from treating old or bad prompts as literal stop conditions. This led to principles such as:
\item prompts are untrusted input,
\item The conversation then focused heavily on XStack. The user wanted RepoX, TestX, AuditX, ControlX, PerformX, CompatX, and SecureX to be robust, modular, extensible, and fast. The assistant helped plan or generate many prompts around this. The important architectural outcome was that XStack should behave like a compiler or incremental build system, not like a CI pipeline that reruns everything after every edit.
\item Repeated gate calls from prompts needed to short-circuit or cache.
\item Later, the user reported that in a new chat they had run a sequence of 13 prompts. They pasted or uploaded a transcript. That transcript became a major source for the handoff package.
\item The 13 prompts reportedly implemented, in order:
\item After the transcript appeared, the user asked for a maximum-fidelity Context Transfer Packet. The assistant produced one, with sections such as workstreams, decisions, tasks, constraints, open questions, risks, artifacts, and verification queue. Then the user asked to turn that into a downloadable report package. The assistant created Markdown/YAML files and a ZIP archive.
\item The conclusion was that docs should be layered: README for accessibility; architecture docs for technical explanation; canon/glossary/contracts for binding doctrine. The 13-prompt transcript reports that docs/canon and many architecture/contract docs were created, but this must be verified.
\end{itemize}



{\small\raggedright SOURCE PATH: docs/\allowbreak{}archive/\allowbreak{}conversations/\allowbreak{}\_\allowbreak{}wiki/\allowbreak{}artifacts/\allowbreak{}index.md\par}


\section{Artifacts}

Artifact counts are structural inventory only.



\begin{quote}\small Dense table summarized for the reader edition. See the source file, HTML output, or reference appendix source for full detail.\end{quote}

\begin{quote}\small Dense table content is summarized in this PDF rendering. See the Markdown source and HTML book for the full table.\end{quote}



{\small\raggedright SOURCE PATH: docs/\allowbreak{}archive/\allowbreak{}conversations/\allowbreak{}\_\allowbreak{}wiki/\allowbreak{}open\_\allowbreak{}questions/\allowbreak{}index.md\par}


\section{Open Questions}

Open questions are extracted uncertainty and verification needs.


\subsection{[Dominium Advanced Simulation and Infrastructure Architecture](../../\_reader/by\_chat/advanced\_simulation\_infrastructure.md)}

\begin{itemize}
\item [UNCERTAIN / UNVERIFIED] No repository files were inspected in this chat. No code was implemented. Proposed paths, module names, command names, file names, VM concepts, and data schemas are design proposals until checked against the actual codebase.
\item [UNCERTAIN / UNVERIFIED] Exact solvers, fixed-point formats, module paths, data schemas, and tick ordering remain to be verified and specified.
\item [UNCERTAIN / UNVERIFIED] The exact gameplay rules runtime was not specified. The command examples were conceptual.
\item [UNCERTAIN / UNVERIFIED] The exact cross-section schema, attachment schema, and packing rules remain unresolved.
\item [UNCERTAIN / UNVERIFIED] Microsegment length, VM design, storage layout, budget values, and actual code integration remain open.
\item [UNCERTAIN / UNVERIFIED] Actual standards packs, rule schema, semantic key format, localization, and AI interaction remain unresolved.
\item Then it produced a prompt for another GPT-5.2 chat that already had a refactor/optimization plan. That prompt told the other chat to amend its existing plan instead of restarting, and to verify coverage of arbitrary placement, unified spatial primitives, co-location, signage, buildings, and determinism/performance.
\item [UNCERTAIN / UNVERIFIED] The architecture has not been verified against the actual repository. The exact implementation details remain unresolved: Q formats, orientation math, command schemas, VM instruction set, microsegment length, carrier ownership, junction archetypes, facility modules, and DECOR/device promotion policy.
\end{itemize}

\subsection{[Dominium APP0 Runtime, Platform, and Renderer Architecture](../../\_reader/by\_chat/app\_runtime\_platform\_renderers.md)}

\begin{itemize}
\item What remains uncertain is whether the current repository already enforces these boundaries mechanically. The actual code needs inspection.
\item The unresolved issue is sharding semantics. APP0 requires sharding hooks, but the chat did not decide whether shards are spatial, logical, temporal, jurisdictional, or something else. The report preserved the idea that APP0 can provide sharding plumbing without inventing simulation semantics.
\item The unresolved issue is the trust and integrity model for setup. The chat did not decide whether content verification should use hashes, signatures, local manifests, network verification, or another model.
\item The unresolved issue is the exact elevation/audit model. The chat did not decide how tools authenticate, request write permissions, log actions, or interact with law enforcement modules.
\item The unresolved issue is the exact renderer capability schema. The chat proposed ideas but did not finalize fields, priorities, or target matrices.
\item The user asked to look at the old code snapshot in project attachments. A prior assistant claimed to inspect it, but later preservation documents marked that claim as **UNCERTAIN / UNVERIFIED**.
\item This became one of the most important unresolved issues. Almost every implementation question depends on the real code: directory layout, build system, languages, existing platform/render layers, client/server dependencies, and engine/game public APIs. The next real step is repository verification.
\item INFERENCE:** The user wants future assistants to be epistemically careful: not inventing repository facts, not treating proposals as decisions, and not losing uncertainty.
\end{itemize}

\subsection{[Dominium Architecture, Application Layer, TestX, and CodeHygiene Planning](../../\_reader/by\_chat/app\_testx\_codehygiene.md)}

\begin{itemize}
\item The main unresolved goal is practical execution: which prompts have actually been run, which files exist, and what implementation should happen next.
\item Setup handles install, verify, repair, rollback, upgrade, downgrade, uninstall, and status. Launcher can invoke Setup, but must not replicate its mutation logic. This prevents two products from disagreeing about installation state.
\item First, verify whether libs/contracts and libs/appcore already exist. Then implement the pure contract headers and appcore skeleton. This should include:
\item Before implementing app RepoX integration, verify actual RepoX metadata paths, TestX invocation, VALIDATE-0 commands, and BUILD-ID-0 files. Without that, an implementation prompt would guess.
\item Use labels like FACT, INFERENCE, UNCERTAIN, PROJECT-CONTEXT in reports.
\item From the user profile and instructions, future assistants should verify time-sensitive/current facts with web where relevant. This chat itself is mostly internal project planning, so external citations are not central.
\item EXEC, ECSX, KERN, ADOPT, DIST, HWCAPS, EXIST, DOMAIN, TRAVEL, TIME, OMNI, LIFE, CIV, WAR, AGENT, TOOL, MOD, FINAL prompt families were generated earlier. They are important as historical design artifacts, but actual repo execution is unverified.
\item The chat produced downloadable handoff files and later an in-chat reader. Those files are useful for aggregation but are not themselves project source code. The package's main caveat is that it does not verify actual repository state.
\end{itemize}

\subsection{[Dominium/Domino Architecture and Codex Prompt Roadmap](../../\_reader/by\_chat/architecture\_codex\_prompts.md)}

\begin{itemize}
\item UNCERTAIN / UNVERIFIED: No repository code was inspected or modified in this chat. The prompts and architecture are plans and artifacts, not proof that implementation exists. The actual repository state, build system, existing code quality, GUI backend availability, TLV schemas, and whether any prompts have been applied remain unverified.
\item Uncertainty remains about actual repository code and build system, but the layering itself is a strong established principle.
\item Unresolved issues include exact TLV schema field tags, exact registry implementation, base pack structure, and whether engine comments/docs may use examples with domain words.
\item The main uncertainty is implementation: the prompts define it, but no code was verified.
\item Unresolved details include exact process IO schemas, machine runtime state, and the relationship between autonomous machines and agent-operated jobs.
\item Unresolved: exact nesting policy, allowed packing modes, and how far to model real packing constraints.
\item Uncertainty remains about freeform geometry versus grid shell as the long-term building representation.
\item The exact destruction and structural-load models remain unresolved and part of Path D.
\end{itemize}

\subsection{[Dominium Architecture, UI, Providers, and Robot OS Strategy](../../\_reader/by\_chat/architecture\_ui\_providers.md)}

\begin{itemize}
\item The exact first playable slice remains unresolved. The first provider wedge is clear, but exact sequence between client/workbench gameplay, Robot OS, and gameplay mechanics still needs a formal milestone. The exact Lua version, Linux baseline, current repo status, and external library/version support need verification. The exact degree of Unreal involvement remains unclear.
\item 1. Verify and pin external library/toolchain facts.
\item It remains uncertain how much the user wants Unreal in the near-term after raylib-first discussion. It also remains uncertain whether Lua 5.4 or 5.5 is preferred; the user appears to value pinning and stability more than "latest" for script ABI.
\end{itemize}

\subsection{[Dominium Build and Future-Proofing Architecture](../../\_reader/by\_chat/build\_and\_future\_proofing.md)}

\begin{itemize}
\item Plain-language limitation: this package is reliable as a preservation of the current visible chat. It should not be treated as a complete archive of every Dominium discussion in other chats. Where the report uses project or repository context, it is labelled as such. The most important caveat is that many items are recommendations, not final user decisions.
\item The unresolved goals are implementation goals: deciding which recommendations become canon, adding public-surface/dependency/build contracts, validating the current live repo state, and applying the proposed cleanup tasks. The preservation package is complete, but the engineering work it describes remains pending.
\item 2. **Verify current repo state.** Confirm the actual current HEAD, build proof, smoke tests, and outstanding warnings before acting.
\item 10. **Feed this report into the master spec book aggregator.** Preserve uncertainty labels and avoid merging suggestions as decisions.
\item The preservation output must distinguish FACT, INFERENCE, UNCERTAIN/UNVERIFIED, and PROJECT-CONTEXT.
\item A future assistant may rely on stale repo state. Mitigation: verify live HEAD, docs, and CI before implementation.
\item What live repo facts should I verify before implementing these tasks?
\item The final uploaded prompt asked for this preservation package. The key thing to preserve is that most architecture proposals are recommendations, not accepted user decisions yet. The next best action is to decide which recommendations become canon, then implement either STRUCTURE-01: Public Surface Registry or the build tuple contract work, after verifying the current live repo state.
\end{itemize}

\subsection{[Dominium Canonical Architecture, Repository Foundation, and Provider Model](../../\_reader/by\_chat/canonical\_structure\_and\_framework.md)}

\begin{itemize}
\item INFERENCE: The user wanted a durable method for deciding future structure changes, not just a one-time layout. They also wanted the assistant to become more rigorous: distinguish decisions from brainstorms, stop over-planning without execution, and preserve uncertainty. The user also wanted a structure that supports other Domino-based games and possibly different engine implementations.
\item A persistent uncertainty is the exact current repo state. The user supplied many status reports and directory exports, but several exports were internally inconsistent or stale at various points. The chat therefore repeatedly emphasized structure report integrity and verification before trusting tree analyses.
\item 1. Verify fast-strict/RepoX/structure report status from a clean current export.
\item The user dislikes vague reassurance and wants uncertainty labelled.
\item UNCERTAIN: The final exact convention for external/upstream versus external/vendor should be checked against current repo policy.
\item UNCERTAIN: Whether Lua should be pinned to 5.4 or 5.5 is not settled in this chat.
\item UNCERTAIN: Whether pack-internal content/ directories are legitimate pack law or legacy remains to be verified.
\item User-reported commits and status summaries: these include commits like 6e0dd93, 1406490, 3243fab, ce9ca, and others. Treat them as FACTs reported in the chat, but verify live repo state before acting.
\end{itemize}

\subsection{[Dominium Chronology \& Celestial Systems](../../\_reader/by\_chat/chronology\_celestial\_systems.md)}

\begin{itemize}
\item The assistant formalized this as the Perfect Earth Calendar, later called HPC-E. The final structure is clear, though the exact leap rule remains unresolved.
\item What remains uncertain is the exact verified object list. The assistant generated a real Milky Way system list, but that list must be checked before becoming data.
\item The unresolved issue is exact fidelity. The chat specified what should exist conceptually, but not the final data schema, map resolution, object attributes, or implementation priority.
\item What remains unresolved is the exact mathematical implementation of these frames and how much relativistic fidelity should exist in phase one.
\item The uncertainty is high relative to Earth. Many names and structures were assistant-generated and should be reviewed before implementation.
\item The unresolved issue is implementation reality. The actual Plan G/The Game codebase was not visible here, so any implementation plan must be reconciled with real architecture.
\item The only major unresolved part is the leap rule.
\item The user wants rigorous, structured, fact-aware assistance. The user specifically requested preservation of uncertainty labels and does not want assistant suggestions treated as user decisions.
\end{itemize}

\subsection{[Dominium Development Routes and Continuity Preservation](../../\_reader/by\_chat/development\_routes.md)}

\begin{itemize}
\item The first answer was assertive. It stated that Dominium must follow Route C and described the route as the only viable one for Dominium. Later parts of the chat corrected the status of that claim. FACT: the assistant made that recommendation. UNCERTAIN / UNVERIFIED: the recommendation was not validated with external sources in the chat and was not explicitly accepted by the user.
\item The assistant produced a large Context Transfer Packet. It preserved the technical route proposal, the caveats around its uncertainty, the proposed phase order, the open questions, the risks, the artifacts, and the next actions.
\item The assistant created those files and reported that the chat label was inferred as **Dominium Development Routes**. The package was said to be safe for aggregation **with caveats**. The main caveats were that the substantive project transcript was short, Route C was not user-confirmed, and Dominium's genre, core loop, stack, implementation state, files, platforms, team, and timeline were not established.
\item UNCERTAIN / UNVERIFIED: Route C was not accepted by the user as final.
\item UNCERTAIN / UNVERIFIED: the required level of determinism for Dominium was not established. It might need full cross-platform deterministic replay, or it might only need ordinary save/load correctness, or something between those extremes.
\item The outputs included a Context Transfer Packet and a downloadable package with multiple files. The files are less important than the purpose: preserve the exact status of the discussion, including uncertainty and assistant-proposal status, so future assistants do not accidentally distort the project.
\item The second explicit goal was to preserve the chat before retirement. The user did not want a normal summary. They wanted a maximum-fidelity transfer packet that would prevent loss of decisions, constraints, unresolved issues, rationale, and next actions. This was addressed by the Context Transfer Packet.
\item Some goals changed over time. The chat began as architecture guidance, became a state-transfer task, became a packaging task, became an in-chat inspection task, and now becomes a plain-English briefing task. The underlying continuity goal stayed consistent: preserve the useful substance without losing uncertainty.
\end{itemize}

\subsection{[Documentation Standards, README Strategy, and Handoff Packaging](../../\_reader/by\_chat/documentation\_standards\_readme.md)}

\begin{itemize}
\item The key thing to remember: this chat produced the standards and prompts for future work; it did not verify or change the project. The next assistant must inspect the actual repository before writing repo-specific docs, README content, platform claims, license claims, or subsystem descriptions.
\item The chat began with the user asking how to have Codex generate documentation for every file in the project. The user wanted metadata, functions, definitions, dependencies, and other maintainability information documented. The core uncertainty was placement: should documentation live in separate sister files, or should it live in file headers and function headers?
\item UNCERTAIN / UNVERIFIED: Whether the project has similarly named files or modules is unknown.
\item What remains uncertain is how the actual project currently organizes public headers, internal headers, and docs. That must be verified from the repository.
\item What remains unverified is the actual project name, supported platforms, license, maturity, and current repository layout. Those cannot be responsibly written without inspecting the repo.
\item The decision depends on an assumption: that the project's public APIs are declared in headers. That is likely for a C/C++ codebase, but the actual public/private boundary is still unverified.
\item What must happen first: verify source roots, exclusions, CMake structure, CI environment variables, and Python availability.
\item What could go wrong: treating assistant suggestions as final requirements or merging contradictory chat claims without preserving uncertainty.
\end{itemize}

\subsection{[Dominium Architecture I](../../\_reader/by\_chat/dominium\_architecture\_i.md)}

\begin{itemize}
\item What remains uncertain is the exact final scope of all simulation systems. Many systems were described in ambitious terms, but not all details are final or internally consistent. The systems layer especially needs canonicalisation before coding.
\item What remains uncertain is the current real-world status and capabilities of "Codex 5.1 Max." That is an external tool fact and requires verification before future use.
\item The key later decision was to skip top-level .txt devspec files because original Markdown files already existed in /docs/.... The visible chat does not show the full content of those docs. They should be preserved as referenced artifacts, but their exact content is **UNCERTAIN / UNVERIFIED**.
\item This affects dlocale, font rendering, platform path handling, data locale JSON files, mod APIs, UI, and documentation. The unresolved part is exactly how much retro support is actually feasible. Claims about old OSes, toolchains, SDL2, OpenGL, Direct3D, and macOS versions require external verification.
\item A large portion of the chat consisted of producing specs for those files. Each spec was meant to tell Codex how to implement that file. The conclusion is that the repository architecture is broad but structured. The unresolved issue is that many later files remain pending.
\item The unresolved issue is API consistency. Generated specs for memory and serialization are inconsistent across files, and some platform/render specs contain outdated or unverified external assumptions.
\item The most important unresolved issue is duplication and contradiction. dweather, dhydro, and dai\_core were generated in more than one form. The package does not choose a final version. A future assistant should not implement any of those blindly.
\item The user decided to skip the top-level .txt files because original Markdown docs already exist in /docs/.... This decision is final for the current workflow, but it depends on those docs actually existing and being current. That is **UNCERTAIN / UNVERIFIED** because the docs were not inspected in this chat.
\end{itemize}

\subsection{[Dominium Architecture II](../../\_reader/by\_chat/dominium\_architecture\_ii.md)}

\begin{itemize}
\item What remains uncertain is the exact code-level API and component layouts. The chat produced specs and prompts, not final source code.
\item This is an important implementation detail, but some of it remains unverified. Lua 5.4.4 portability to old compilers needs testing. The exact Lua data schemas still need definition. The docs also still have JSON-oriented format sections, so the Lua-vs-JSON policy needs reconciliation.
\item The biggest unresolved goals are practical:
\item Verify DX9/Windows 2000, SDL2, Lua/toolchain compatibility.
\item FACT: This was answered in response to Codex's clarifying questions. It is a concrete implementation direction, though Lua 5.4.4 portability remains unverified.
\item 1. **Verify and apply V4 docs.
\item The user repeatedly requested maximum-fidelity reports, registers, handoff packets, and human-readable summaries. The user wants labelled uncertainty and does not want future assistants to invent missing context.
\item or forget that Lua-vs-JSON remains unresolved.
\end{itemize}

\subsection{[Dominium Architecture III: Launcher, Platform, Renderer, and Handoff Architecture](../../\_reader/by\_chat/dominium\_architecture\_iii.md)}

\begin{itemize}
\item The user then uploaded dominium.7z, saying it was the git repository as of that moment. **FACT:** The prior assistant said it could not open .7z archives in that environment. That means the actual repository state remained unverified at that stage. This became important later, because many implementation ideas were discussed without confirmed access to the repo contents.
\item This created one of the main unresolved issues. DX12 was discussed earlier, but omitted from the latest renderer list. X11 and Wayland were discussed earlier, but the latest platform categories are POSIX, SDL1, SDL2, and Native. Therefore, the final status is:
\item DX12 is **UNCERTAIN / UNVERIFIED** and should not be assumed.
\item X11/Wayland placement is unresolved.
\item The prior assistant summarized these files as containing two different launcher concepts: an engine-rendered launcher view and a native Win32 launcher. It also claimed the CMake file built the Win32 source for non-Windows targets. However, this report must preserve that as **UNCERTAIN / UNVERIFIED**, because the files were not re-inspected during final packaging.
\item This is final as a user-stated product goal. Feasibility still requires verification.
\item Those outputs are important artifacts, but they are not the substance of the design. Their purpose is to preserve this chat's decisions, uncertainty, and future work for later aggregation into a master Project Spec Book.
\item What remains uncertain is not the constraint itself, but how existing repository code currently enforces it. That requires code inspection.
\end{itemize}

\subsection{[Dominium Architecture IV](../../\_reader/by\_chat/dominium\_architecture\_iv.md)}

\begin{itemize}
\item The conclusion was clear and accepted: **Domino is the reusable engine; Dominium is the game/product layer**. What remains uncertain is the actual repository state. The chat planned the split, but did not verify whether the repo already matches it.
\item The unresolved issue is that the actual repo was not inspected in this chat. The planned structure may not match the current files.
\item The unresolved issue is practical feasibility, especially for retro targets and Carbon OS. Their actual toolchains and system calls remain unverified.
\item The unresolved issue is the exact IR payload layout and how far limited renderers should go. A CGA or MDA backend cannot realistically support modern 3D in the same way as Vulkan. The agreed approach is explicit capability flags and graceful fallback or no-op behaviour.
\item The unresolved issue is the exact manifest/schema format. JSON, TOML, and INI-like examples appeared in prompts, but no final schema choice was made.
\item The unresolved issue is that external packaging details require verification. WiX, macOS notarization, Linux packaging, and AppImage tooling are all external-world facts and may be stale.
\item The package and this report should help future assistants continue without re-reading the whole chat. However, they must preserve uncertainty: actual repo state is unverified, some transcript sections were skipped, external tooling facts may be stale, and assistant suggestions are not user decisions unless accepted or built upon.
\item It is also reasonable to infer that the user wants the future project spec book to preserve rationale, not just lists of decisions. The later handoff requests emphasised visible rationale, rejected options, uncertainty, and changes of direction.
\end{itemize}

\subsection{[Dominium Canon, Repository Alignment, and Portability Doctrine](../../\_reader/by\_chat/dominium\_complete\_conversation.md)}

\begin{itemize}
\item Uncertainty: the raw pasted prompt is not the current repo authority unless materialized in repo canon. The later audit found that it had in fact been materialized under docs/canon/constitution\_v1.md.
\item Uncertainty: during this preservation pass, a potential inconsistency appeared: some docs claim product roots moved under apps/, while an earlier inspected code path under root client/ was available and apps/client was not fetched successfully. This must be verified against the current physical tree.
\item Unresolved goals include verifying the latest physical repo tree, proving runtime implementation maturity, formalizing public ABI/API boundaries, completing layout/naming cleanup, and building a second Domino-based product to prove reuse.
\item 1. Verify the current physical repo tree against layout docs and contracts/repo/layout.contract.toml.
\item The user likely wants future assistants to challenge weak framing, preserve uncertainty, avoid shallow "best practice" lists, and focus on decisions that survive future rewrites.
\item "Create a gated plan for verifying current repo layout and exceptions."
\end{itemize}

\subsection{[Dominium + Domino Codex Planning and Handoff](../../\_reader/by\_chat/dominium\_domino\_codex\_planning.md)}

\begin{itemize}
\item The most important things to remember are: the project's strict architecture constraints, the Milestone-0 prompt sequence, the later shared CLI/TUI/GUI/input/render prompt sequences, the missing pack-system prompt, the unified startup policy, the unresolved repo-verification issue, and the need to treat generated prompts as plans rather than proof of implementation.
\item The assistant generated the first two prompts. The third was never produced because the conversation changed direction. That missing third prompt remains an important unresolved item.
\item However, the later context transfer packet treated that inspection claim as **UNCERTAIN / UNVERIFIED**. This was careful and correct: within the visible chat, there was no reliable tool trace proving the repo had actually been inspected. Therefore, future assistants must verify dominium.zip or the active checkout before relying on the platform/render/API answer.
\item What remains uncertain is how much of this architecture already exists in the current repo and how much is still planned. The chat produced prompts and plans, but it did not prove execution.
\item The unresolved issue is whether those prompts were ever run and whether the generated code, if any, conforms to the project's strict C89/C++98 and determinism requirements.
\item The unresolved complication is C89 portability. Standard C89 does not provide long long, but u64, i64, and q48\_16 require 64-bit representation. Future work must decide whether to allow compiler extensions, emulate 64-bit values for strict retro targets, or define platform tiers with conditional support.
\item The unresolved issue is that the prompt's minimal manifest parser and hardcoded platform path are only a starting point. A real implementation must generalize product discovery, avoid duplicated path logic, and preserve compatibility/versioning rules.
\item The unresolved issue is whether the prompts are too broad and whether the existing repo already has overlapping modules. Future work should inspect the repo before applying these prompts.
\end{itemize}

\subsection{[Dominium Workbench, AIDE, Presentation Architecture, Provider Strategy, Product-Spine Planning, and Preservation Companion](../../\_reader/by\_chat/dominium\_full\_conversation.md)}

\begin{itemize}
\item Verify live repo state, then review/run FULL-GATE-LEGACY-TEST-ROUTE-01, then generate PACK-INTERNAL-LAYOUT-CANON-01 if the queue/status supports the maintenance lane.
\item Preserve uncertainty labels.
\item Future chats could easily misunderstand this conversation by treating Workbench as a GUI framework, treating raylib as the engine, resuming broad structure cleanup, skipping projection conformance, or assuming live repo status from stale pasted reports. The mitigation is to verify repo state first, preserve the contract/service/projection/provider distinctions, and follow the task queue.
\item The best next action is to verify live repo state, then review or run FULL-GATE-LEGACY-TEST-ROUTE-01, then generate PACK-INTERNAL-LAYOUT-CANON-01 if status supports it.
\end{itemize}

\subsection{[Dominium Setup Architecture and Handoff](../../\_reader/by\_chat/dominium\_setup.md)}

\begin{itemize}
\item FACT: The setup system was repeatedly defined as the only component allowed to install files, modify system/installation state, own installed-file metadata, repair installs, verify artifacts, uninstall, downgrade, upgrade, and roll back. This is the most important architectural rule from the chat.
\item What remains uncertain is the actual current implementation state in the repository after the applied Codex prompts. The design is clear, but the repo must still be inspected.
\item FACT: The final canonical setup layout uses setup/core/\{fetch,verify,install,rollback\}, setup/include/\{dsk,dsu\}, setup/packages, setup/platform, setup/tests, and setup/ui.
\item The unresolved part is whether the actual repository now matches this canonical layout.
\item UNCERTAIN / UNVERIFIED: Exact Visual Studio 2026 behavior and toolchain details should be verified in the current environment before relying on them.
\item UNCERTAIN / UNVERIFIED: The user stated the prompts were applied, but this chat did not show the final repository tree or build results.
\item UNCERTAIN / UNVERIFIED: The actual schema files under schema/setup/ need to be checked.
\item Several goals remain unresolved because the actual repo was not inspected here. We do not know whether the applied Codex prompts fully succeeded, whether setup builds, whether schemas exist, whether libs/contracts is complete, or whether setup/launcher CLI smoke tests pass.
\end{itemize}

\subsection{[Domino Dominium Workbench](../../\_reader/by\_chat/domino\_dominium\_workbench.md)}

\begin{itemize}
\item Reliability note:** This is a human-readable consolidation, not a live repository audit. Repo-current statements that came from pasted material or prior-chat excerpts remain **UNVERIFIED** until checked against the live julesc013/dominium repo and current AIDE/validator outputs.
\item Verify current live repo status and task gate state.
\item > Verify repo state, then generate COMMAND-RESULT-VIEW-SLICE-01.
\end{itemize}

\subsection{[Dominium/Domino Engine Refactor Planning](../../\_reader/by\_chat/domino\_engine\_refactor\_prompts.md)}

\begin{itemize}
\item The answer also introduced a deterministic bytecode VM as the preferred way to support modded behavior. A question was raised about whether native-code plugins should also be allowed, but the user did not answer. That remains unresolved.
\item What remains uncertain is the actual repository state. The chat designed architecture and prompts, but did not inspect source code. Therefore the plan needs repo verification before implementation.
\item The unresolved work is to generate detailed subsystem implementation prompts for heat, power, fluids, atmospheres, vehicles, and structural loads after the core engine spine is implemented.
\item The remaining uncertainty is plugin policy. A deterministic VM was proposed, but whether native-code plugins are allowed remains unresolved.
\item What remains uncertain is the exact thresholds and content policies for when things promote or demote.
\item The unresolved issue is how much of STRUCT compilation should be implemented first. The prompt plan is broad; a real implementation may need a smaller first milestone.
\item ignoring FACT/INFERENCE/UNCERTAIN labels;
\item UNCERTAIN / UNVERIFIED:** The package files were generated in the sandbox during the prior step, but repository files were not created or inspected.
\end{itemize}

\subsection{[Domino/Dominium Engine Baseline, Architecture, and Feasibility](../../\_reader/by\_chat/engine\_baseline\_architecture.md)}

\begin{itemize}
\item Uncertainty: the repo has many modules and some historical/legacy surfaces, so exact implementation maturity varies by subsystem. The distinction should remain a formal requirement in a master spec.
\item Uncertainty: the full determinism envelope across all current modules was not independently re-run in this chat. User-supplied local validation results should be preserved but verified before being used as audit proof.
\item Uncertainty: this remains strategic architecture, not an implemented decision.
\item Unresolved goals include:
\item User wants uncertainty labelled and framing corrected when evidence disagrees.
\item UNCERTAIN: whether the user wants Unreal integration at all as a client/editor.
\item UNCERTAIN: exact tolerance for using external geometry libraries.
\item UNCERTAIN: whether broad repo restructuring will be desired after Milestone 0.
\end{itemize}

\subsection{[Dominium Foundation Lock, Workbench Spine, and Parallel Codex Handoff](../../\_reader/by\_chat/foundation\_workbench\_codex.md)}

\begin{itemize}
\item What remains uncertain is how quickly this architecture will move from contracts and fixtures into runtime implementations and product slices. The next chat should not assume runtime provider resolvers, package runtime, or Workbench shell exist.
\item verify latest live repo state after user-pasted Wave 1 completion;
\item 1. Verify live repo state.
\item The user wants direct, evidence-grounded, audit-ready answers; timestamps/model labels; explicit uncertainty; no repeated clarification unless necessary; large continuous Codex prompt blocks when generating prompts; targeted tests instead of full suites; and rapid progress toward Workbench/code.
\item The user is likely to reject slow, report-only, overcautious process. Future assistants should produce executable next prompts quickly after verifying state.
\item It is uncertain how many concurrent Codex workers the user will actually run at once and whether the user wants coordinator prompts or worker prompts next.
\item The most important continuation is: verify current origin/main, confirm Wave 1 and Workbench validation state, then generate or run COMMAND-RESULT-VIEW-SLICE-01.
\end{itemize}

\subsection{[Domino Framework and Open-Source Provider Architecture](../../\_reader/by\_chat/framework\_open\_source\_provider.md)}

\begin{itemize}
\item The chat inspected julesc013/dominium and found evidence of existing abstractions such as system and graphics layers, stubs, and soft-backed backends. The finding was that raylib could fill concrete provider gaps without replacing the architecture. However, current repo facts are stale/uncertain and must be verified, especially the C17/C++17 vs C90/C++98 baseline contradiction.
\item The exact first implementation branch, versions, and dependency pins are unresolved. The exact Domino Framework ABI is unresolved. The exact Lua version is unresolved. The current repository baseline must be verified. The formal policy for GPL/LGPL/unclear license material is unresolved. The sparse simulation and CAD systems need formal specs and prototypes.
\item The raylib decision was also directionally accepted. The user repeatedly expressed enthusiasm for raylib and asked about using as much of raylib infrastructure as possible. The caveat is that raylib is a provider suite, not architecture. This affects rendering, audio, input, asset preview, and Workbench bootstrap.
\item 1. Verify repo baseline and dependency versions.
\item The user requires source-grounded, audit-ready, uncertainty-labelled responses. The preservation prompt explicitly requires FACT/INFERENCE/UNCERTAIN/PROJECT-CONTEXT labels, no invented facts, no treating brainstorms as decisions, and no over-compression. The prompt also requires a human-readable report first and downloadable files if possible ?filecite?turn29file0?.
\item The most blocking unresolved issues are repo baseline verification, exact dependency pins, minimal framework ABI, provider profile order, deterministic simulation spec, and license policy.
\item Verify the current julesc013/dominium CMake language baseline.
\end{itemize}

\subsection{[Dominium Content and GUI Rebuild Planning](../../\_reader/by\_chat/gui\_binary\_content.md)}

\begin{itemize}
\item GPT-5.5 Pro - 2026-05-27 00:00:00 Australia/Melbourne; time UNVERIFIED
\item whether knowledge systems should include uncertainty, false beliefs, or lost knowledge,
\item This stage established a pattern that continues throughout the chat: the user does not want impressive-looking prompts that hide unresolved design choices. The user wants the assumptions exposed before anything is formalized.
\item The assistant agreed with the general direction, but added caveats. The most important were:
\item Those caveats were not final user decisions, but they are important warnings.
\item The conclusion was not "run the prompt now." The conclusion was that CONTENT0 is a strong foundation but needs discussion before finalization. The assistant-generated rewrite is useful but not final. The unresolved issues need to be resolved before Codex or any implementation work touches the repository.
\item What remains uncertain is how exactly to encode these ideas in schemas and data. For example, it is not yet decided how deterministic seed hierarchies work, how macro content refines into micro content, or how timeline causality should be represented.
\item The unresolved issue is the shared backend/UI contract. Without that, the GUI plan is only a list of technologies.
\end{itemize}

\subsection{[Dominium Language, Platform, and Architecture Baseline](../../\_reader/by\_chat/language\_platform\_architecture.md)}

\begin{itemize}
\item The 32-bit question followed. The answer was to make full native products 64-bit and keep 32-bit as constrained or projection support only. The key caveat was not to make native pointer size part of game law. Save, replay, network, IDs, and renderer IR must use fixed-width types and never serialize size\_t, long, uintptr\_t, or raw pointers.
\item Exact tolerance for exceptions/RTTI inside private C++17 code remains unsettled. Exact platform floors and raylib/SDL2 priority remain unresolved.
\item Verify the Windows 7 and macOS 10.9 toolchain/library assumptions using official sources.
\end{itemize}

\subsection{[Dominium Launcher Application-Layer Handoff](../../\_reader/by\_chat/launcher\_app\_layer.md)}

\begin{itemize}
\item These prompts matter because they superseded earlier advice. They also created a new practical priority: verify that the applied prompts actually succeeded.
\item This changed the correct future behavior. The next assistant must not redesign. It must implement, audit, maintain, document, verify, and harden.
\item This matters because it prevents platform-specific drift. If CLI has verify, GUI must not implement a different hidden version of verify. If GUI has a pack enable flow, it should map to the same command/intent as CLI and TUI.
\item The final purity prompt required those to be moved or quarantined. However, actual repo success remains unverified in this chat.
\item The important change is that the conversation moved from "What architecture should we use?" to "Architecture is locked; how do we verify and harden the implementation?"
\item The largest unresolved goal is verifying the actual repo. The user stated that two Codex prompts were applied, but this chat did not independently verify the filesystem, build, tests, scripts, or launcher docs.
\item Another unresolved goal is confirming whether the launcher hardening prompt has been applied.
\item A third unresolved goal is confirming whether command graph, UI IR, binding validation, zero-pack tests, RepoX changelog display, and BUILD-ID mismatch refusal are implemented.
\end{itemize}

\subsection{[Dominium Launcher and Setup Architecture](../../\_reader/by\_chat/launcher\_setup\_architecture.md)}

\begin{itemize}
\item Uncertainty:** The exact repository layout and existing implementation were not inspected in this chat.
\item Uncertainty:** The final language boundary for setup remains unresolved after the later C89/dsys/dgfx launcher direction.
\item Uncertainty:** The exact mapping into runtime binaries remains unverified.
\item Uncertainty:** Accounts/social/wiki/forum/direct messaging are future-facing and were not converted into the final five-tab implementation plan.
\item Uncertainty:** Earlier setup designs used JSON manifests. Later dsys/dgfx architecture leaned toward TLV .dmeta files.
\item Uncertainty:** Actual plugin ABI details and dynamic loading support depend on dsys/repo APIs.
\item Uncertainty:** Actual dsys/dgfx APIs were not inspected.
\item Conclusion reached:** This chat is now documented and can be merged later, but with caveats.
\end{itemize}

\subsection{[Dominium Modularity, AIDE Refactorability, and Future-Proof Architecture](../../\_reader/by\_chat/modularity\_aide\_refactorability.md)}

\begin{itemize}
\item The live repo still needs verification. The exact AIDE files, contracts, validators, and migration tasks must be implemented. It remains unresolved which existing roots and tools map to which final homes, which APIs become stable, which code is reusable across products/games/projects, and which prior assistant recommendations should become formal requirements after user review.
\item 1. **TASK-09** - Verify live repo baseline: current branch, root inventory, validators, CTest status, generated artifacts.
\item The user prefers source-grounded, audit-ready decisions; bounded uncertainty; explicit task prompts; and reusable engineering doctrine. Future assistants should avoid shallow affirmation and should distinguish accepted user decisions from assistant recommendations.
\item Important caveat: the only current uploaded file available for this task is Pasted text.txt, which contains the preservation-package prompt. Earlier generated handoff files, ZIP packages, or uploaded docs referenced in prior conversation are not available in this chat unless the user re-uploads them.
\item The most serious risk is that a future assistant treats this chat as if it verified the live repo. It did not. The second serious risk is that it hardens every assistant suggestion into a requirement. Several items are strong recommendations aligned with user goals, but still require implementation review. Future chats must preserve labels and verify stale facts.
\item Future assistants must verify live repo state before implementing cleanup.
\end{itemize}

\subsection{[Dominium Omega/Xi/Pi Architecture \& Future-Proofing Planning](../../\_reader/by\_chat/omega\_xi\_pi\_architecture\_future.md)}

\begin{itemize}
\item Preserve FACT / INFERENCE / UNCERTAIN labels.
\item Do not invent or flatten uncertainty.
\item 7. "List all artifacts I should verify in GitHub main."
\item 9. "What should I verify before starting Sigma?"
\end{itemize}

\subsection{[Dominium OS-like Architecture, Repository Convergence, and Interface Operating Layer](../../\_reader/by\_chat/os\_interface\_repo\_architecture.md)}

\begin{itemize}
\item The unresolved goals are implementation-level. The repo still needs command dispatch unification, product boot proof, portable projection proof, AppShell rendered-mode law update, Workbench module contracts, document/patch/result/refusal schemas, and a boot-to-replay MVP. It also needs continued exception retirement, CTest/RepoX remediation, and verification of current build status.
\item 2. Repair canonical verify test discovery.
\item The user wants uncertainty labels and correction of incorrect framing.
\item This preservation task specifically required FACT / INFERENCE / UNCERTAIN / PROJECT-CONTEXT labels, stable IDs, registers, spec sheets, an aggregator packet, self-audit, and downloadable files if tools are available.
\item UNCERTAIN: Whether the user wants to rename engine conceptually to kernel in code, or only use kernel as a conceptual layer.
\item UNCERTAIN: Whether the final public product name should be Dominium Workbench, Dominium Studio, Domino Workbench, or something else.
\item UNCERTAIN: How aggressive the next actual code refactor should be versus staged adapters and codegen.
\end{itemize}

\subsection{[Dominium Codex Platform Renderer API Plan](../../\_reader/by\_chat/platform\_renderer\_api\_plan.md)}

\begin{itemize}
\item This is not a fatal problem for the prompt pack, because the prompts themselves repeatedly tell Codex to discover files with rg, tree, type, CMake commands, and target listing. But it is a major caveat for any future assistant: do not assume the actual repo matches every file path or API name unless checked.
\item UNCERTAIN / UNVERIFIED:** The conversation did not prove that these prompts were actually applied to the repository. They are an assumed prerequisite, not verified code state.
\item Uncertainty remains about actual repo state. The prompts are designed to discover the repo state, but the chat itself did not inspect the archive.
\item The final active prompt pack narrowed runtime completion to Tier-1 on Windows and compile-gated correctness for Tier-2. This is a practical compromise because Codex runs on Windows. But it also creates a caveat: if "fully complete" later means fully running X11/Wayland/Cocoa/Metal/Vulkan, more prompts or native CI validation will be required.
\item This became one of the final "done" gates, along with scripts\textbackslash{}build\_codex\_verify.bat.
\item This matters because the conversation was long and contained multiple plans, superseded options, caveats, and active artifacts. The report package and this plain-language report exist to prevent future assistants from restarting, re-asking, or losing distinctions between decisions, brainstorms, and uncertainties.
\item The user also seems to want future assistants to operate with high continuity and not lose subtle distinctions. This is inferred from the repeated requests for maximum-fidelity transfer, stable IDs, labels, uncertainty tracking, and report packaging.
\item Another unresolved goal is whether the master engine prompts 1-14 are truly implemented. The later plan depends on them, but this chat did not verify the repo.
\end{itemize}

\subsection{[Dominium Platform Support Planning](../../\_reader/by\_chat/platform\_support.md)}

\begin{itemize}
\item UNCERTAIN / UNVERIFIED: The name "Domino" was never confirmed by the user. Future work should treat "Domino" as an assistant-created placeholder unless the user explicitly accepts it. The useful idea is the distinction between full product support and reduced engine/core support, not necessarily the name.
\item The main unresolved issue is the exact definition of support. Does Tier-0 mean identical feature parity, simultaneous release, identical save compatibility, identical UI, equal QA coverage, or some platform-specific differences? The assistant inferred that Tier-0 should mean feature parity, full QA, long-term support, and high release priority, but the exact contractual meaning still needs formal definition.
\item What remains unresolved is the exact PC baseline. The chat did not decide whether Windows 10, Windows 11, or both are required. It did not decide whether Windows ARM64 matters. It did not decide whether macOS support includes Intel Macs, Apple Silicon Macs, or both. It did not define Linux distributions, glibc requirements, Flatpak/AppImage/deb/rpm packaging, X11/Wayland policy, or Steam Deck verification.
\item UNCERTAIN / UNVERIFIED: The chat did not establish minimum Android API level, ABI list, target graphics APIs, Play Store rules, Android TV status, Android Go status, Automotive status, Chromebook status, or vendor ROM support. These need future decisions and current-source verification.
\item The user also listed iPod Touch, Apple TV, Apple Watch, Mac, macOS, tvOS, watchOS, and Apple Vision Pro indirectly through AR. These were not classified in detail. INFERENCE: macOS belongs under PC. UNCERTAIN / UNVERIFIED: tvOS, watchOS, and visionOS remain undecided. Full Dominium on Apple Watch is not established and should not be assumed.
\item The unresolved issues are substantial. The chat did not decide whether WebGPU is required, whether WebGL fallback is mandatory, whether the game must work offline as a PWA, whether WASM threading/SIMD is required, how browser storage will work, how asset loading will be handled, or what browser versions will be supported.
\item UNCERTAIN / UNVERIFIED: Current console platform facts were not verified in this chat. Future work must verify official PlayStation, Xbox, and Nintendo developer requirements from current official sources before implementation planning.
\item INFERENCE: PC handhelds should be subtargets under PC, with their own QA and UI profiles. UNCERTAIN / UNVERIFIED: Steam Deck's exact tier was not finally decided.
\end{itemize}

\subsection{[Domino/Dominium Portability, Assurance, and Future-Proof Architecture](../../\_reader/by\_chat/portability\_assurance\_future\_proof.md)}

\begin{itemize}
\item The chat proposed using standards such as DO-178C, NIST SSDF, OWASP ASVS, SLSA, and SPDX as design inputs only. The proposed internal version is DDAP v0 with DIL levels. Uncertainty: user has not explicitly ratified DDAP.
\item The actual repo was not inspected. The exact accepted directory structure, API policy, DDAP profile, compatibility promises, and first pilot module remain unresolved.
\item Caveat: carry forward as INFERENCE / assistant recommendation.
\item Caveat: carry forward as FACT goal + INFERENCE architecture.
\item Caveat: carry forward as assistant recommendation.
\item Implications: Safe for aggregation only with caveats.
\item Caveat: carry forward as FACT / UNCERTAIN scope.
\item 10. Feed this package into the master Project Spec Book with caveats.
\end{itemize}

\subsection{[Dominium README Architecture, Ports, and Determinism](../../\_reader/by\_chat/readme\_ports\_determinism.md)}

\begin{itemize}
\item What remains uncertain is whether the actual repository README.md matches the final pasted version. The chat only saw pasted text and generated prompts; it did not inspect a Git repository.
\item The unresolved issue is whether a metadata-only /ports directory still violates the user's preference. The final README keeps /ports in a limited form, but the user's wording could mean no /ports directory at all. That should be clarified before writing the directory contract.
\item The unresolved work is to define the actual network protocol, prediction model, rollback window, reconciliation logic, and state-hash validation.
\item The major unresolved goal is deciding exactly what "no separate ports directory/system" means physically in the repository. The final README allows /ports as metadata-only. That may satisfy the user, but it is not confirmed. If the user meant no /ports directory at all, the README still needs another revision.
\item Another unresolved goal is making the README and normative specs align. The README says specs are authoritative, but those specs were not inspected in this chat.
\item UNCERTAIN / UNVERIFIED: /ports as metadata-only is the current README state, but may not be fully accepted.
\item A related rejected idea was **retro build flows under /ports/<target> containing source or alternative implementations**. The final README superseded this with metadata-only /ports, assuming /ports is retained at all. This rejection is final for source code and behavior, but the existence of /ports itself remains uncertain.
\item The first next step is to verify the actual repository README.md against the final pasted README. The chat only used pasted text. If the repository differs, future work should start from the actual file, not from memory.
\end{itemize}

\subsection{[Dominium + Domino Refactor Architecture](../../\_reader/by\_chat/refactor\_architecture.md)}

\begin{itemize}
\item Reliability note: Repository state, Codex execution, exact version numbers, and actual file modifications are **UNCERTAIN / UNVERIFIED** unless explicitly stated otherwise.
\item The most important thing to remember is that this chat did not implement the refactor. It designed the architecture and generated prompts and handoff material. Any future assistant must verify actual repository state before assuming files were moved, prompts were applied, or code was updated.
\item UNCERTAIN / UNVERIFIED:** It did not establish that this system already exists in code.
\item What remains uncertain is the exact degree to which existing code currently respects that boundary. The user provided a repo tree, but this chat did not inspect or modify code directly.
\item Unresolved details include exact default DOMINIUM\_HOME paths per OS and exact implementation of import/GC.
\item Unresolved details include exact schemas, parser choice, tests, and actual code implementation.
\item The key unresolved point is implementation timing. The report preserves this architecture, but the initial Codex refactor was not supposed to rewrite the entire UI system immediately.
\item Actual implementation is unresolved. This chat generated architecture and prompts, not a verified changed codebase.
\end{itemize}

\subsection{[AIDE, XStack, and Dominium Refactor Control Plane](../../\_reader/by\_chat/refactor\_control\_plane.md)}

\begin{itemize}
\item Unresolved goals include finishing AIDE Q35-Q57, stabilizing Dominium CTest/RepoX blockers, defining final AIDE install/upgrade bundles, implementing root recycling machinery, deciding when product boot proof can proceed, and later deciding whether AIDE Runtime/Gateway/Hosts are truly needed.
\item 1. Verify current Dominium HEAD and POST-CONVERGE status.
\item The user values direct, source-grounded, audit-ready, detailed answers. The user asked for preservation packages that are human-readable and machine-aggregatable. The user repeatedly prefers not to lose nuance, rejected options, uncertainties, or artifacts. The user wants current facts verified when possible and wants future assistants to avoid re-asking answered questions.
\item UNCERTAIN: Exact public-facing naming for AIDE Core/Kernel remains partly open. Exact threshold for accepting CTest-warning status remains a user/repo governance decision. Exact final AIDE installation packaging details remain pending.
\item Future assistants must verify live repo state before implementing.
\item "Merge this with another chat's preservation package without losing uncertainty labels."
\end{itemize}

\subsection{[Dominium XStack Release Identity and Versioning](../../\_reader/by\_chat/release\_identity\_and\_versioning.md)}

\begin{itemize}
\item The chat rebuilt SemVer from first principles and decided that strict SemVer should apply only to components with declared public contracts. Candidate true-SemVer entities include SDKs, engine libraries, tool APIs/CLIs, protocol libraries, plugin hosts, schema libraries, and reusable runtime libraries. It remains unresolved which exact repo modules qualify.
\item GBN should identify CI/public/internal builds and appear as provenance, e.g. +gbn.7137, but it must not be used as SemVer precedence. BII should primarily be structured manifest metadata, because it may encode lane, target, kind, configuration, and other build identity fields. The exact BII schema remains unresolved.
\item The main unresolved goals are to produce formal specs: Release Constitution, SemVer Component Inventory, Suite Release Policy, Build Identity Spec, Channel/Lifecycle Spec, Artifact Naming Spec, Manifest Schema, Target Taxonomy, and Capability Registry.
\item It is still uncertain how strict the user wants suite X.Y.Z field meanings to be, whether capabilities should fully replace compatibility profiles or coexist with them, and how much old-platform support should be real versus aspirational.
\item The most important unresolved issue is formal classification. Without knowing which entities are strict SemVer components, suite versions, product release IDs, build fields, or capabilities, the rest of the policy cannot be enforced safely.
\item Platform targeting remains unresolved. The chat concluded that coarse families like WinNT5, WinNT10, MacOSX4, Linux5, or DOS5 are useful support labels, but exact binary compatibility needs target baselines and runtime/toolchain profiles. Linux kernel major alone is not enough. Windows 9x/NT/NT5/NT6/NT10 likely need separate build lanes or at least explicit tested baselines.
\end{itemize}

\subsection{[Dominium TestX/RepoX Governance and Handoff Chat](../../\_reader/by\_chat/testx\_repox\_governance.md)}

\begin{itemize}
\item This topic remains unresolved in implementation because the repo snapshot showed .dominium\_build\_number and update\_build\_number.cmake, but the chat did not inspect their contents. Those files may still implement older build-number behavior. That must be verified.
\item The unresolved part is implementation verification. The repo has Sol/Earth content in data and game/content; the engine must not assume those exist.
\item The unresolved part is how much of TESTX is actually implemented in the repo. The repo contains many tests and scripts, but this chat did not inspect them. A rules-to-checks audit is needed.
\item This remains conceptually accepted in the chat, but implementation is unverified.
\item The unresolved goals are mostly implementation and verification:
\item The repo already has UI IR/codegen infrastructure. The next step is to verify whether GUI/TUI truly bind to the CLI command graph and whether UI is data rather than logic.
\item The user has Windows 10 + VS2022. The next step is to verify v141/v141\_xp/SDK components and create isolated CMake presets for XP/Win7/Win8/Win10/Win11 builds.
\item Acting before verifying repo state.
\end{itemize}

\subsection{[Dominium Timekeeping and 2038 Resilience](../../\_reader/by\_chat/timekeeping\_and\_2038\_resilience.md)}

\begin{itemize}
\item What remains uncertain is the current exact status of specific libc/OS/toolchain combinations, because those facts can change and were not reverified during this preservation turn.
\item What remains uncertain is the precise target list of legacy machines/toolchains and how far compatibility must go, especially for 16-bit compilers with weak or absent 64-bit arithmetic.
\item The main unresolved goals are backend audit, ACT serialization audit, civil/astronomical projection design, and verification of external platform facts.
\item The main uncertainty is completeness of repo inspection. The assistant read key files, but not all serialization paths and not all platform backends.
\item 6. Verify current platform/library facts before using them in formal public documentation.
\item The preservation prompt explicitly requires a human-readable report first, not only a machine-readable handoff. It requires uncertainty labels, stable registers, preservation of artifacts, no invented facts, and downloadable files if tools are available. ?filecite?turn30file0?
\item The user prefers source-grounded, audit-ready technical reasoning, direct structure, and careful distinction between facts, recommendations, and unresolved issues. PROJECT-CONTEXT indicates Dominium values C89/C++98 portability, deterministic architecture, CLI/TUI support, and clean platform/runtime separation.
\item The strongest unresolved technical task is auditing actual backend timers and persistence formats.
\end{itemize}

\subsection{[Dominium UE6, Domino, and Deterministic Universe Feasibility](../../\_reader/by\_chat/ue6\_domino\_deterministic\_universe.md)}

\begin{itemize}
\item UNCERTAIN / UNVERIFIED: This is not a complete reconstruction of every Dominium conversation ever held. The project context shows that many other chats exist about platforms, renderers, setup, game architecture, ecology, ECS, launcher design, and development planning, but their full transcripts are not visible here. Those items are labelled PROJECT-CONTEXT when referenced.
\item Uncertainty: public UE6 technical capabilities and requirements remain time-sensitive and need verification before any future hard commitment.
\item The chat rejected the idea that clients can authoritatively simulate resource production, hidden state, enemy bases, combat, or persistent terrain changes in an MMO. Clients can assist, but the server must verify and commit.
\item REJECTED-02: Waiting for UE6 before beginning serious work. Deprioritised because UE6 availability and capabilities remain uncertain.
\item The user wants direct, source-grounded, audit-ready answers; explicit uncertainty; bounded estimates; and correction when framing is wrong. The preservation prompt explicitly requires human-readable explanation first, structured registers second, uncertainty labels, no invented facts, no hidden chain-of-thought, and downloadable files if tools are available.
\item UNCERTAIN: The final intended role of Domino remains unclear in the visible transcript. It may be a custom engine, target runtime, renderer, or broader platform plan, but the current chat does not fully define it.
\end{itemize}

\subsection{[Dominium UI Editor and Tool Editor Planning](../../\_reader/by\_chat/ui\_editor\_tool\_editor\_planning.md)}

\begin{itemize}
\item The most important thing to remember is that this chat was a planning and prompt-generation chat, not proof of implementation. **UNCERTAIN / UNVERIFIED:** No generated prompt is evidence that repository code was actually changed. The next assistant must first verify whether the prompts were executed, inspect the repo/source bundle if available, and avoid treating planned files as existing files.
\item The user uploaded screenshot bundles for setup and launcher. **UNCERTAIN / UNVERIFIED:** The screenshots were not inspected in this chat. The assistant still produced a logical layout description, and the user accepted it as useful. Future work should inspect the bundles before claiming exact screenshot fidelity.
\item What remains uncertain is the actual state of the code. The user named files, but the chat did not inspect the repository. A future assistant must verify exact paths, build targets, TLV parser implementation, and current launcher behavior.
\item Unresolved details include the actual TLV wire format, migration strategy from launcher\_ui\_v1.tlv, and how much legacy data can be preserved during import.
\item The unresolved risk is implementation complexity. Real native tab controls, splitter behavior, and scroll panels can be nontrivial in a child-HWND Win32 UI.
\item This is a strong planned design, but some details remain assistant recommendations rather than explicit user-confirmed decisions. Future implementation should preserve that uncertainty if writing a formal spec.
\item The unresolved issue is that the uploaded screenshot bundles were not inspected. The logical specs are accepted planning artifacts, not verified screenshot extractions.
\item The exact scope of IDE-native live editing remains unresolved. The assistant proposed preview hosts, but it is not confirmed whether that fully satisfies the user's desire to use Visual Studio, Xcode, and Linux GUI tools.
\end{itemize}

\subsection{[Dominium Universe Explorer Planning and Repo Handoff](../../\_reader/by\_chat/universe\_explorer\_planning.md)}

\begin{itemize}
\item Uncertain: whether all seven universal layers are now formally captured under exactly that naming. The repo has equivalent or related doctrine, but not every early phrase appears as a single artifact.
\item Uncertain: the repo contains pieces of this through affordance/capability/domain docs, but the assistant did not find one single master Deep Primitives spec in the docs it read.
\item The most important unresolved issue is that the Universe Explorer plan is still a proposed/recommended next strategy, not yet a completed repo task. The next best action is either to draft PRESENTATION-CONTRACT-01 / UNIVERSE-EXPLORER-CONTRACT-01, or to preserve this chat and merge it with other old-chat reports into a master Project Spec Book.
\end{itemize}

\subsection{[Dominium Architecture, Workbench, AIDE, and Product-Spine Planning](../../\_reader/by\_chat/workbench\_aide\_product\_spine.md)}

\begin{itemize}
\item The largest uncertainty is not the conceptual architecture; that has converged. The largest uncertainty is live execution state: whether prompts generated near the end of the chat have already been run locally, committed, or pushed. The next chat should first inspect .aide/queue/current.toml, recent commits, and relevant audit files.
\item Remaining uncertainty: exact implementation of erosion/ecology/evolution proxies remains future domain work.
\item 1. Verify current repo state.
\item Explicitly label uncertainty.
\item The next chat should verify repo state first.
\item Preserve FACT / INFERENCE / UNCERTAIN / PROJECT-CONTEXT labels.
\item Verify live repo state before acting.
\item safe\_for\_aggregation: "Yes, with caveats"
\end{itemize}

\subsection{[Dominium World Architecture](../../\_reader/by\_chat/world\_architecture.md)}

\begin{itemize}
\item The main unresolved issue is implementation detail. The design says Q16.16 and Q4.12, but the earlier prompt used signed types incorrectly. The future implementation must use unsigned local coordinates or a centered signed design.
\item The unresolved parts are the exact meshing algorithm, the fixed-point scaling of ?, and the sample resolution. A 32? sample grid per 16m chunk was recommended but not confirmed as final.
\item The exact solvers remain unresolved. The architecture is clear, but formulas and resolutions are still future work.
\item The save format remains conceptual in places. Exact binary headers, endian policy, compression, and migration strategy are unresolved.
\item The crucial caveat is implementation: unsigned or centered representations are needed. The design intent is final, but the exact typedefs must be fixed. If the implementation uses signed types incorrectly, the whole coordinate system breaks.
\item This is not yet an implementation-level final spec because solver formulas are unresolved.
\item FACT:** Core target is C89. This is non-negotiable in principle, but exact portability mechanics are unresolved because strict C89 lacks standard fixed-width integer types.
\item The user wants direct, detailed, rigorous, critical responses. The user prefers source labels, uncertainty preservation, and avoiding invented facts. The user does not want assistant suggestions silently upgraded to decisions.
\end{itemize}

\subsection{[Dominium XStack and Lab Galaxy Handoff](../../\_reader/by\_chat/xstack\_lab\_galaxy.md)}

\begin{itemize}
\item But this chat did not verify the repository. That became one of the most important unresolved issues.
\item This is one of the most important outputs associated with this chat, but it is also one of the most uncertain because it is user-reported from another chat and not verified directly here. The user explicitly wanted it audited for completion and consistency.
\item INFERENCE: The user also wanted to prevent future assistants from flattening nuance. They repeatedly asked for maximum fidelity, uncertainty labels, rejected options, contradictions, and rationale. This implies a strong preference for preserving complexity rather than oversimplifying.
\item INFERENCE: The user likely wants to resume productive development after verifying the substrate. The next likely feature area is survival or Lab Galaxy refinement, but the user has not made a final decision in this visible chat.
\item 4. From documentation to verifying a new 13-prompt substrate.
\item The biggest unresolved goal is verifying the actual repository. The package says what was reported, but not what is proven. Another unresolved goal is deciding what comes next: survival, Lab Galaxy UX, distributed SRZ, or documentation polish.
\item The main assumption behind all of this is that the repo can actually enforce these contracts. That remains unverified in this chat.
\item The best next step is to verify the actual repository state. This matters because the long 13-prompt transcript is not proof. The user explicitly does not know whether everything was truly implemented.
\end{itemize}



{\small\raggedright SOURCE PATH: docs/\allowbreak{}archive/\allowbreak{}conversations/\allowbreak{}\_\allowbreak{}wiki/\allowbreak{}risks/\allowbreak{}index.md\par}


\section{Risks}

Risks are review triggers, not resolved findings.


\begin{itemize}
\item CONTRA-0001 medium conversation\_vs\_current\_queue: Archived conversation discusses work related to blocked scope gameplay.
\item CONTRA-0002 low stale\_external\_claim: Archived conversation contains potentially stale external or baseline claim: C89.
\item CONTRA-0003 low stale\_external\_claim: Archived conversation contains potentially stale external or baseline claim: C++98.
\item CONTRA-0004 low stale\_external\_claim: Archived conversation contains potentially stale external or baseline claim: ISO C89.
\item CONTRA-0005 medium conversation\_vs\_current\_queue: Archived conversation discusses work related to blocked scope runtime\_module\_loader.
\item CONTRA-0006 medium conversation\_vs\_current\_queue: Archived conversation discusses work related to blocked scope provider\_runtime.
\item CONTRA-0007 medium conversation\_vs\_current\_queue: Archived conversation discusses work related to blocked scope gameplay.
\item CONTRA-0008 medium conversation\_vs\_current\_queue: Archived conversation discusses work related to blocked scope renderer\_implementation.
\item CONTRA-0009 low stale\_external\_claim: Archived conversation contains potentially stale external or baseline claim: SDK.
\item CONTRA-0010 medium conversation\_vs\_current\_queue: Archived conversation discusses work related to blocked scope gameplay.
\item CONTRA-0011 medium conversation\_vs\_current\_queue: Archived conversation discusses work related to blocked scope renderer\_implementation.
\item CONTRA-0012 low stale\_external\_claim: Archived conversation contains potentially stale external or baseline claim: C89.
\item CONTRA-0013 low stale\_external\_claim: Archived conversation contains potentially stale external or baseline claim: C++98.
\item CONTRA-0014 low stale\_external\_claim: Archived conversation contains potentially stale external or baseline claim: DX12.
\item CONTRA-0015 low stale\_external\_claim: Archived conversation contains potentially stale external or baseline claim: iOS.
\item CONTRA-0016 medium conversation\_vs\_current\_queue: Archived conversation discusses work related to blocked scope provider\_runtime.
\item CONTRA-0017 medium conversation\_vs\_current\_queue: Archived conversation discusses work related to blocked scope gameplay.
\item CONTRA-0018 low stale\_external\_claim: Archived conversation contains potentially stale external or baseline claim: C89.
\item CONTRA-0019 low stale\_external\_claim: Archived conversation contains potentially stale external or baseline claim: C++98.
\item CONTRA-0020 low stale\_external\_claim: Archived conversation contains potentially stale external or baseline claim: ISO C89.
\item CONTRA-0021 medium conversation\_vs\_current\_queue: Archived conversation discusses work related to blocked scope broad\_workbench\_ui.
\item CONTRA-0022 medium conversation\_vs\_current\_queue: Archived conversation discusses work related to blocked scope provider\_runtime.
\item CONTRA-0023 medium conversation\_vs\_current\_queue: Archived conversation discusses work related to blocked scope gameplay.
\item CONTRA-0024 medium conversation\_vs\_current\_queue: Archived conversation discusses work related to blocked scope renderer\_implementation.
\item CONTRA-0025 medium conversation\_vs\_current\_queue: Archived conversation discusses work related to blocked scope native\_gui.
\item CONTRA-0026 low stale\_external\_claim: Archived conversation contains potentially stale external or baseline claim: C89.
\item CONTRA-0027 low stale\_external\_claim: Archived conversation contains potentially stale external or baseline claim: C++98.
\item CONTRA-0028 low stale\_external\_claim: Archived conversation contains potentially stale external or baseline claim: DX12.
\item CONTRA-0029 low stale\_external\_claim: Archived conversation contains potentially stale external or baseline claim: SDK.
\item CONTRA-0030 low stale\_external\_claim: Archived conversation contains potentially stale external or baseline claim: C89.
\item CONTRA-0031 low stale\_external\_claim: Archived conversation contains potentially stale external or baseline claim: C++98.
\item CONTRA-0032 low stale\_external\_claim: Archived conversation contains potentially stale external or baseline claim: SDK.
\item CONTRA-0033 medium conversation\_vs\_current\_queue: Archived conversation discusses work related to blocked scope provider\_runtime.
\item CONTRA-0034 medium conversation\_vs\_current\_queue: Archived conversation discusses work related to blocked scope package\_runtime.
\item CONTRA-0035 medium conversation\_vs\_current\_queue: Archived conversation discusses work related to blocked scope gameplay.
\item CONTRA-0036 medium conversation\_vs\_current\_queue: Archived conversation discusses work related to blocked scope renderer\_implementation.
\item CONTRA-0037 medium conversation\_vs\_current\_queue: Archived conversation discusses work related to blocked scope native\_gui.
\item CONTRA-0038 low stale\_external\_claim: Archived conversation contains potentially stale external or baseline claim: C89.
\item CONTRA-0039 low stale\_external\_claim: Archived conversation contains potentially stale external or baseline claim: C++98.
\item CONTRA-0040 low stale\_external\_claim: Archived conversation contains potentially stale external or baseline claim: SDK.
\item CONTRA-0041 medium conversation\_vs\_current\_queue: Archived conversation discusses work related to blocked scope gameplay.
\item CONTRA-0042 medium conversation\_vs\_current\_queue: Archived conversation discusses work related to blocked scope gameplay.
\item CONTRA-0043 low stale\_external\_claim: Archived conversation contains potentially stale external or baseline claim: iOS.
\item CONTRA-0044 low stale\_external\_claim: Archived conversation contains potentially stale external or baseline claim: C89.
\item CONTRA-0045 low stale\_external\_claim: Archived conversation contains potentially stale external or baseline claim: C++98.
\item CONTRA-0046 low stale\_external\_claim: Archived conversation contains potentially stale external or baseline claim: ISO C89.
\item CONTRA-0047 low stale\_external\_claim: Archived conversation contains potentially stale external or baseline claim: iOS.
\item CONTRA-0048 medium conversation\_vs\_current\_queue: Archived conversation discusses work related to blocked scope gameplay.
\item CONTRA-0049 medium conversation\_vs\_current\_queue: Archived conversation discusses work related to blocked scope renderer\_implementation.
\item CONTRA-0050 low stale\_external\_claim: Archived conversation contains potentially stale external or baseline claim: C89.
\item CONTRA-0051 low stale\_external\_claim: Archived conversation contains potentially stale external or baseline claim: C++98.
\item CONTRA-0052 medium conversation\_vs\_current\_queue: Archived conversation discusses work related to blocked scope gameplay.
\item CONTRA-0053 medium conversation\_vs\_current\_queue: Archived conversation discusses work related to blocked scope renderer\_implementation.
\item CONTRA-0054 low stale\_external\_claim: Archived conversation contains potentially stale external or baseline claim: C89.
\item CONTRA-0055 low stale\_external\_claim: Archived conversation contains potentially stale external or baseline claim: Windows 98.
\item CONTRA-0056 low stale\_external\_claim: Archived conversation contains potentially stale external or baseline claim: Windows NT 2000.
\item CONTRA-0057 low stale\_external\_claim: Archived conversation contains potentially stale external or baseline claim: DX12.
\item CONTRA-0058 low stale\_external\_claim: Archived conversation contains potentially stale external or baseline claim: Android.
\item CONTRA-0059 medium conversation\_vs\_current\_queue: Archived conversation discusses work related to blocked scope renderer\_implementation.
\item CONTRA-0060 medium conversation\_vs\_current\_queue: Archived conversation discusses work related to blocked scope release\_publication.
\item CONTRA-0061 low stale\_external\_claim: Archived conversation contains potentially stale external or baseline claim: C89.
\item CONTRA-0062 low stale\_external\_claim: Archived conversation contains potentially stale external or baseline claim: C++98.
\item CONTRA-0063 low stale\_external\_claim: Archived conversation contains potentially stale external or baseline claim: Win98.
\item CONTRA-0064 low stale\_external\_claim: Archived conversation contains potentially stale external or baseline claim: Windows 98.
\item CONTRA-0065 low stale\_external\_claim: Archived conversation contains potentially stale external or baseline claim: Windows NT 2000.
\item CONTRA-0066 low stale\_external\_claim: Archived conversation contains potentially stale external or baseline claim: Mac OS 9.
\item CONTRA-0067 low stale\_external\_claim: Archived conversation contains potentially stale external or baseline claim: Mac OS 8.
\item CONTRA-0068 low stale\_external\_claim: Archived conversation contains potentially stale external or baseline claim: DX12.
\item CONTRA-0069 medium conversation\_vs\_current\_queue: Archived conversation discusses work related to blocked scope gameplay.
\item CONTRA-0070 medium conversation\_vs\_current\_queue: Archived conversation discusses work related to blocked scope renderer\_implementation.
\item CONTRA-0071 medium conversation\_vs\_current\_queue: Archived conversation discusses work related to blocked scope native\_gui.
\item CONTRA-0072 low stale\_external\_claim: Archived conversation contains potentially stale external or baseline claim: CP/M.
\item CONTRA-0073 low stale\_external\_claim: Archived conversation contains potentially stale external or baseline claim: Android.
\item CONTRA-0074 low stale\_external\_claim: Archived conversation contains potentially stale external or baseline claim: SDK.
\item CONTRA-0075 medium conversation\_vs\_current\_queue: Archived conversation discusses work related to blocked scope gameplay.
\item CONTRA-0076 low stale\_external\_claim: Archived conversation contains potentially stale external or baseline claim: C89.
\item CONTRA-0077 low stale\_external\_claim: Archived conversation contains potentially stale external or baseline claim: C++98.
\item CONTRA-0078 low stale\_external\_claim: Archived conversation contains potentially stale external or baseline claim: C89.
\item CONTRA-0079 low stale\_external\_claim: Archived conversation contains potentially stale external or baseline claim: C++98.
\item CONTRA-0080 low stale\_external\_claim: Archived conversation contains potentially stale external or baseline claim: ISO C89.
\item CONTRA-0081 low stale\_external\_claim: Archived conversation contains potentially stale external or baseline claim: CP/M.
\item CONTRA-0082 low stale\_external\_claim: Archived conversation contains potentially stale external or baseline claim: Android.
\item CONTRA-0083 medium conversation\_vs\_current\_queue: Archived conversation discusses work related to blocked scope runtime\_module\_loader.
\item CONTRA-0084 medium conversation\_vs\_current\_queue: Archived conversation discusses work related to blocked scope provider\_runtime.
\item CONTRA-0085 medium conversation\_vs\_current\_queue: Archived conversation discusses work related to blocked scope package\_runtime.
\item CONTRA-0086 medium conversation\_vs\_current\_queue: Archived conversation discusses work related to blocked scope gameplay.
\item CONTRA-0087 medium conversation\_vs\_current\_queue: Archived conversation discusses work related to blocked scope native\_gui.
\item CONTRA-0088 low stale\_external\_claim: Archived conversation contains potentially stale external or baseline claim: SDK.
\item CONTRA-0089 medium conversation\_vs\_current\_queue: Archived conversation discusses work related to blocked scope gameplay.
\item CONTRA-0090 medium conversation\_vs\_current\_queue: Archived conversation discusses work related to blocked scope release\_publication.
\item CONTRA-0091 low stale\_external\_claim: Archived conversation contains potentially stale external or baseline claim: SDK.
\item CONTRA-0092 medium conversation\_vs\_current\_queue: Archived conversation discusses work related to blocked scope broad\_workbench\_ui.
\item CONTRA-0093 medium conversation\_vs\_current\_queue: Archived conversation discusses work related to blocked scope provider\_runtime.
\item CONTRA-0094 medium conversation\_vs\_current\_queue: Archived conversation discusses work related to blocked scope gameplay.
\item CONTRA-0095 medium conversation\_vs\_current\_queue: Archived conversation discusses work related to blocked scope native\_gui.
\item CONTRA-0096 low stale\_external\_claim: Archived conversation contains potentially stale external or baseline claim: SDK.
\item CONTRA-0097 medium conversation\_vs\_current\_queue: Archived conversation discusses work related to blocked scope gameplay.
\item CONTRA-0098 low stale\_external\_claim: Archived conversation contains potentially stale external or baseline claim: C89.
\item CONTRA-0099 low stale\_external\_claim: Archived conversation contains potentially stale external or baseline claim: C++98.
\item CONTRA-0100 low stale\_external\_claim: Archived conversation contains potentially stale external or baseline claim: ISO C89.
\item CONTRA-0101 medium conversation\_vs\_current\_queue: Archived conversation discusses work related to blocked scope provider\_runtime.
\item CONTRA-0102 medium conversation\_vs\_current\_queue: Archived conversation discusses work related to blocked scope gameplay.
\item CONTRA-0103 medium conversation\_vs\_current\_queue: Archived conversation discusses work related to blocked scope broad\_workbench\_ui.
\item CONTRA-0104 medium conversation\_vs\_current\_queue: Archived conversation discusses work related to blocked scope runtime\_module\_loader.
\item CONTRA-0105 medium conversation\_vs\_current\_queue: Archived conversation discusses work related to blocked scope provider\_runtime.
\item CONTRA-0106 medium conversation\_vs\_current\_queue: Archived conversation discusses work related to blocked scope package\_runtime.
\item CONTRA-0107 medium conversation\_vs\_current\_queue: Archived conversation discusses work related to blocked scope gameplay.
\item CONTRA-0108 medium conversation\_vs\_current\_queue: Archived conversation discusses work related to blocked scope renderer\_implementation.
\item CONTRA-0109 medium conversation\_vs\_current\_queue: Archived conversation discusses work related to blocked scope native\_gui.
\item CONTRA-0110 medium conversation\_vs\_current\_queue: Archived conversation discusses work related to blocked scope release\_publication.
\item CONTRA-0111 low stale\_external\_claim: Archived conversation contains potentially stale external or baseline claim: C89.
\item CONTRA-0112 low stale\_external\_claim: Archived conversation contains potentially stale external or baseline claim: C++98.
\item CONTRA-0113 low stale\_external\_claim: Archived conversation contains potentially stale external or baseline claim: SDK.
\item CONTRA-0114 medium conversation\_vs\_current\_queue: Archived conversation discusses work related to blocked scope broad\_workbench\_ui.
\item CONTRA-0115 medium conversation\_vs\_current\_queue: Archived conversation discusses work related to blocked scope provider\_runtime.
\item CONTRA-0116 medium conversation\_vs\_current\_queue: Archived conversation discusses work related to blocked scope gameplay.
\item CONTRA-0117 medium conversation\_vs\_current\_queue: Archived conversation discusses work related to blocked scope renderer\_implementation.
\item CONTRA-0118 low stale\_external\_claim: Archived conversation contains potentially stale external or baseline claim: C++98.
\item CONTRA-0119 low stale\_external\_claim: Archived conversation contains potentially stale external or baseline claim: SDK.
\item CONTRA-0120 medium conversation\_vs\_current\_queue: Archived conversation discusses work related to blocked scope gameplay.
\item CONTRA-0121 medium conversation\_vs\_current\_queue: Archived conversation discusses work related to blocked scope native\_gui.
\item CONTRA-0122 low stale\_external\_claim: Archived conversation contains potentially stale external or baseline claim: iOS.
\item CONTRA-0123 low stale\_external\_claim: Archived conversation contains potentially stale external or baseline claim: Android.
\item CONTRA-0124 low stale\_external\_claim: Archived conversation contains potentially stale external or baseline claim: SDK.
\item CONTRA-0125 medium conversation\_vs\_current\_queue: Archived conversation discusses work related to blocked scope broad\_workbench\_ui.
\item CONTRA-0126 medium conversation\_vs\_current\_queue: Archived conversation discusses work related to blocked scope provider\_runtime.
\item CONTRA-0127 medium conversation\_vs\_current\_queue: Archived conversation discusses work related to blocked scope gameplay.
\item CONTRA-0128 medium conversation\_vs\_current\_queue: Archived conversation discusses work related to blocked scope renderer\_implementation.
\item CONTRA-0129 low stale\_external\_claim: Archived conversation contains potentially stale external or baseline claim: C89.
\item CONTRA-0130 low stale\_external\_claim: Archived conversation contains potentially stale external or baseline claim: C++98.
\item CONTRA-0131 low stale\_external\_claim: Archived conversation contains potentially stale external or baseline claim: Windows 98.
\item CONTRA-0132 low stale\_external\_claim: Archived conversation contains potentially stale external or baseline claim: Mac OS 9.
\item CONTRA-0133 low stale\_external\_claim: Archived conversation contains potentially stale external or baseline claim: iOS.
\item CONTRA-0134 low stale\_external\_claim: Archived conversation contains potentially stale external or baseline claim: Android.
\item CONTRA-0135 low stale\_external\_claim: Archived conversation contains potentially stale external or baseline claim: SDK.
\item CONTRA-0136 medium conversation\_vs\_current\_queue: Archived conversation discusses work related to blocked scope gameplay.
\item CONTRA-0137 low stale\_external\_claim: Archived conversation contains potentially stale external or baseline claim: SDK.
\item CONTRA-0138 medium conversation\_vs\_current\_queue: Archived conversation discusses work related to blocked scope gameplay.
\item CONTRA-0139 medium conversation\_vs\_current\_queue: Archived conversation discusses work related to blocked scope renderer\_implementation.
\item CONTRA-0140 low stale\_external\_claim: Archived conversation contains potentially stale external or baseline claim: C89.
\item CONTRA-0141 low stale\_external\_claim: Archived conversation contains potentially stale external or baseline claim: C++98.
\item CONTRA-0142 low stale\_external\_claim: Archived conversation contains potentially stale external or baseline claim: CP/M.
\item CONTRA-0143 low stale\_external\_claim: Archived conversation contains potentially stale external or baseline claim: SDK.
\item CONTRA-0144 medium conversation\_vs\_current\_queue: Archived conversation discusses work related to blocked scope renderer\_implementation.
\item CONTRA-0145 medium conversation\_vs\_current\_queue: Archived conversation discusses work related to blocked scope native\_gui.
\item CONTRA-0146 low stale\_external\_claim: Archived conversation contains potentially stale external or baseline claim: C89.
\item CONTRA-0147 medium conversation\_vs\_current\_queue: Archived conversation discusses work related to blocked scope runtime\_module\_loader.
\item CONTRA-0148 medium conversation\_vs\_current\_queue: Archived conversation discusses work related to blocked scope native\_gui.
\item CONTRA-0149 medium conversation\_vs\_current\_queue: Archived conversation discusses work related to blocked scope gameplay.
\item CONTRA-0150 medium conversation\_vs\_current\_queue: Archived conversation discusses work related to blocked scope native\_gui.
\item CONTRA-0151 medium conversation\_vs\_current\_queue: Archived conversation discusses work related to blocked scope release\_publication.
\item CONTRA-0152 low stale\_external\_claim: Archived conversation contains potentially stale external or baseline claim: Mac OS 8.
\item CONTRA-0153 low stale\_external\_claim: Archived conversation contains potentially stale external or baseline claim: Android.
\item CONTRA-0154 low stale\_external\_claim: Archived conversation contains potentially stale external or baseline claim: SDK.
\item CONTRA-0155 medium conversation\_vs\_current\_queue: Archived conversation discusses work related to blocked scope gameplay.
\item CONTRA-0156 medium conversation\_vs\_current\_queue: Archived conversation discusses work related to blocked scope renderer\_implementation.
\item CONTRA-0157 medium conversation\_vs\_current\_queue: Archived conversation discusses work related to blocked scope native\_gui.
\item CONTRA-0158 low stale\_external\_claim: Archived conversation contains potentially stale external or baseline claim: C89.
\item CONTRA-0159 low stale\_external\_claim: Archived conversation contains potentially stale external or baseline claim: C++98.
\item CONTRA-0160 low stale\_external\_claim: Archived conversation contains potentially stale external or baseline claim: ISO C89.
\item CONTRA-0161 low stale\_external\_claim: Archived conversation contains potentially stale external or baseline claim: CP/M.
\item CONTRA-0162 low stale\_external\_claim: Archived conversation contains potentially stale external or baseline claim: Android.
\item CONTRA-0163 low stale\_external\_claim: Archived conversation contains potentially stale external or baseline claim: SDK.
\item CONTRA-0164 medium conversation\_vs\_current\_queue: Archived conversation discusses work related to blocked scope gameplay.
\item CONTRA-0165 medium conversation\_vs\_current\_queue: Archived conversation discusses work related to blocked scope renderer\_implementation.
\item CONTRA-0166 medium conversation\_vs\_current\_queue: Archived conversation discusses work related to blocked scope release\_publication.
\item CONTRA-0167 low stale\_external\_claim: Archived conversation contains potentially stale external or baseline claim: CP/M.
\item CONTRA-0168 low stale\_external\_claim: Archived conversation contains potentially stale external or baseline claim: iOS.
\item CONTRA-0169 low stale\_external\_claim: Archived conversation contains potentially stale external or baseline claim: Android.
\item CONTRA-0170 low stale\_external\_claim: Archived conversation contains potentially stale external or baseline claim: WebGPU.
\item CONTRA-0171 low stale\_external\_claim: Archived conversation contains potentially stale external or baseline claim: console program.
\item CONTRA-0172 low stale\_external\_claim: Archived conversation contains potentially stale external or baseline claim: SDK.
\item CONTRA-0173 medium conversation\_vs\_current\_queue: Archived conversation discusses work related to blocked scope gameplay.
\item CONTRA-0174 low stale\_external\_claim: Archived conversation contains potentially stale external or baseline claim: C89.
\item CONTRA-0175 low stale\_external\_claim: Archived conversation contains potentially stale external or baseline claim: C++98.
\item CONTRA-0176 low stale\_external\_claim: Archived conversation contains potentially stale external or baseline claim: C89.
\item CONTRA-0177 low stale\_external\_claim: Archived conversation contains potentially stale external or baseline claim: C++98.
\item CONTRA-0178 low stale\_external\_claim: Archived conversation contains potentially stale external or baseline claim: CP/M.
\item CONTRA-0179 low stale\_external\_claim: Archived conversation contains potentially stale external or baseline claim: 286.
\item CONTRA-0180 low stale\_external\_claim: Archived conversation contains potentially stale external or baseline claim: C89.
\item CONTRA-0181 low stale\_external\_claim: Archived conversation contains potentially stale external or baseline claim: CP/M.
\item CONTRA-0182 low stale\_external\_claim: Archived conversation contains potentially stale external or baseline claim: Android.
\item CONTRA-0183 low stale\_external\_claim: Archived conversation contains potentially stale external or baseline claim: SDK.
\item CONTRA-0184 medium conversation\_vs\_current\_queue: Archived conversation discusses work related to blocked scope provider\_runtime.
\item CONTRA-0185 low stale\_external\_claim: Archived conversation contains potentially stale external or baseline claim: SDK.
\item CONTRA-0186 medium conversation\_vs\_current\_queue: Archived conversation discusses work related to blocked scope release\_publication.
\item CONTRA-0187 low stale\_external\_claim: Archived conversation contains potentially stale external or baseline claim: SDK.
\item CONTRA-0188 medium conversation\_vs\_current\_queue: Archived conversation discusses work related to blocked scope provider\_runtime.
\item CONTRA-0189 medium conversation\_vs\_current\_queue: Archived conversation discusses work related to blocked scope gameplay.
\item CONTRA-0190 medium conversation\_vs\_current\_queue: Archived conversation discusses work related to blocked scope renderer\_implementation.
\item CONTRA-0191 medium conversation\_vs\_current\_queue: Archived conversation discusses work related to blocked scope release\_publication.
\item CONTRA-0192 low stale\_external\_claim: Archived conversation contains potentially stale external or baseline claim: C89.
\item CONTRA-0193 low stale\_external\_claim: Archived conversation contains potentially stale external or baseline claim: C++98.
\item CONTRA-0194 low stale\_external\_claim: Archived conversation contains potentially stale external or baseline claim: iOS.
\item CONTRA-0195 low stale\_external\_claim: Archived conversation contains potentially stale external or baseline claim: SDK.
\item CONTRA-0196 low stale\_external\_claim: Archived conversation contains potentially stale external or baseline claim: C89.
\item CONTRA-0197 low stale\_external\_claim: Archived conversation contains potentially stale external or baseline claim: C++98.
\item CONTRA-0198 medium conversation\_vs\_current\_queue: Archived conversation discusses work related to blocked scope broad\_workbench\_ui.
\item CONTRA-0199 medium conversation\_vs\_current\_queue: Archived conversation discusses work related to blocked scope gameplay.
\item CONTRA-0200 low stale\_external\_claim: Archived conversation contains potentially stale external or baseline claim: C89.
\item CONTRA-0201 low stale\_external\_claim: Archived conversation contains potentially stale external or baseline claim: C++98.
\item CONTRA-0202 medium conversation\_vs\_current\_queue: Archived conversation discusses work related to blocked scope renderer\_implementation.
\item CONTRA-0203 medium conversation\_vs\_current\_queue: Archived conversation discusses work related to blocked scope native\_gui.
\item CONTRA-0204 low stale\_external\_claim: Archived conversation contains potentially stale external or baseline claim: C++98.
\item CONTRA-0205 low stale\_external\_claim: Archived conversation contains potentially stale external or baseline claim: Android.
\item CONTRA-0206 medium conversation\_vs\_current\_queue: Archived conversation discusses work related to blocked scope broad\_workbench\_ui.
\item CONTRA-0207 medium conversation\_vs\_current\_queue: Archived conversation discusses work related to blocked scope provider\_runtime.
\item CONTRA-0208 medium conversation\_vs\_current\_queue: Archived conversation discusses work related to blocked scope gameplay.
\item CONTRA-0209 medium conversation\_vs\_current\_queue: Archived conversation discusses work related to blocked scope runtime\_module\_loader.
\item CONTRA-0210 medium conversation\_vs\_current\_queue: Archived conversation discusses work related to blocked scope provider\_runtime.
\item CONTRA-0211 medium conversation\_vs\_current\_queue: Archived conversation discusses work related to blocked scope package\_runtime.
\item CONTRA-0212 medium conversation\_vs\_current\_queue: Archived conversation discusses work related to blocked scope gameplay.
\item CONTRA-0213 medium conversation\_vs\_current\_queue: Archived conversation discusses work related to blocked scope renderer\_implementation.
\item CONTRA-0214 medium conversation\_vs\_current\_queue: Archived conversation discusses work related to blocked scope native\_gui.
\item CONTRA-0215 medium conversation\_vs\_current\_queue: Archived conversation discusses work related to blocked scope release\_publication.
\item CONTRA-0216 low stale\_external\_claim: Archived conversation contains potentially stale external or baseline claim: C89.
\item CONTRA-0217 low stale\_external\_claim: Archived conversation contains potentially stale external or baseline claim: C++98.
\item CONTRA-0218 medium conversation\_vs\_current\_queue: Archived conversation discusses work related to blocked scope gameplay.
\item CONTRA-0219 low stale\_external\_claim: Archived conversation contains potentially stale external or baseline claim: C89.
\item CONTRA-0220 low stale\_external\_claim: Archived conversation contains potentially stale external or baseline claim: C++98.
\item CONTRA-0221 low stale\_external\_claim: Archived conversation contains potentially stale external or baseline claim: ISO C89.
\item CONTRA-0222 low stale\_external\_claim: Archived conversation contains potentially stale external or baseline claim: SDK.
\item CONTRA-0223 medium conversation\_vs\_current\_queue: Archived conversation discusses work related to blocked scope gameplay.
\item CONTRA-0224 medium conversation\_vs\_docs: Corpus contains both C89 and C17 claims.
\item CONTRA-0225 medium conversation\_vs\_docs: Corpus contains both C++98 and C++17 claims.
\item CONTRA-0226 medium conversation\_vs\_conversation: Corpus contains both JSON and TLV claims.
\item CONTRA-0227 medium conversation\_vs\_conversation: Corpus contains both WinForms and WinUI claims.
\end{itemize}


\chapter{Per-Conversation Reader Digest}

\chapter{Per-Conversation Reader Digest}

This digest gives one concise entry per conversation. Full reader pages remain in \_reader/by\_chat/.


\section{Dominium Advanced Simulation and Infrastructure Architecture}

\begin{itemize}
\item Source folder: docs/archive/conversations/advanced\_simulation\_infrastructure/
\item Reader page: docs/archive/conversations/\_reader/by\_chat/advanced\_simulation\_infrastructure.md
\item Primary source: docs/archive/conversations/advanced\_simulation\_infrastructure/dominium\_advanced\_simulation\_infrastructure\_architecture\_\_human\_readable\_chat\_report.txt
\item Completeness: complete
\item Topics: architecture, content, contracts\_schema, determinism, governance, platform, release, simulation, timekeeping, tooling, ui, workbench
\item Digest: [FACT] This chat was about designing major architecture for the **Dominium/Domino deterministic engine**, with emphasis on advanced simulation, gameplay mechanics, user-buildable infrastructure, arbitrary placement, buildings, splines/corridors, signage, markings, and how all of that should eventually be implemented or merged into an existing refactor plan. The conversation began from a technical planning prompt f...
\end{itemize}

\section{Dominium APP0 Runtime, Platform, and Renderer Architecture}

\begin{itemize}
\item Source folder: docs/archive/conversations/app\_runtime\_platform\_renderers/
\item Reader page: docs/archive/conversations/\_reader/by\_chat/app\_runtime\_platform\_renderers.md
\item Primary source: docs/archive/conversations/app\_runtime\_platform\_renderers/dominium\_app0\_runtime\_platform\_renderers\_\_human\_readable\_chat\_report.txt
\item Completeness: complete
\item Topics: architecture, content, contracts\_schema, determinism, governance, platform, release, setup\_launcher, simulation, timekeeping, tooling, ui, worldgen
\item Digest: This chat was mainly about the **Dominium / Domino project's application and runtime layer**, referred to throughout as **APP0**. The user began by giving a large implementation-oriented prompt titled **"PROMPT APP0 - RUNTIMES, APPLICATIONS, PLATFORMS \& RENDERERS."** That prompt defined a strict scope: work on the **client, server, launcher, setup/installer, tools, renderers, and platform support**, while not rede...
\end{itemize}

\section{Dominium Architecture, Application Layer, TestX, and CodeHygiene Planning}

\begin{itemize}
\item Source folder: docs/archive/conversations/app\_testx\_codehygiene/
\item Reader page: docs/archive/conversations/\_reader/by\_chat/app\_testx\_codehygiene.md
\item Primary source: docs/archive/conversations/app\_testx\_codehygiene/dominium\_app\_testx\_codehygiene\_handoff\_\_human\_readable\_chat\_report.txt
\item Completeness: complete
\item Topics: architecture, content, contracts\_schema, determinism, governance, platform, release, setup\_launcher, simulation, timekeeping, tooling, ui, worldgen, xstack\_aide
\item Digest: This chat was a long, high-density planning and consolidation session for the **Dominium / Domino** project. In plain terms, it was about turning an extremely ambitious simulation-game vision into a disciplined, modular, future-proof engineering program. The project being discussed is not a normal game in the narrow sense. It is intended to be a deterministic universe engine and game platform: something that can s...
\end{itemize}

\section{Dominium/Domino Architecture and Codex Prompt Roadmap}

\begin{itemize}
\item Source folder: docs/archive/conversations/architecture\_codex\_prompts/
\item Reader page: docs/archive/conversations/\_reader/by\_chat/architecture\_codex\_prompts.md
\item Primary source: docs/archive/conversations/architecture\_codex\_prompts/dominium\_domino\_architecture\_codex\_prompts\_\_human\_readable\_chat\_report.txt
\item Completeness: complete
\item Topics: architecture, content, contracts\_schema, determinism, governance, release, setup\_launcher, simulation, timekeeping, tooling, ui, worldgen
\item Digest: FACT: This chat was mainly about turning the user's Dominium/Domino game-engine vision into a coherent, modular, deterministic architecture and then into a long sequence of implementation prompts for Codex or another coding agent. It began with a large project-defining prompt from the user. That prompt established Domino as the low-level deterministic engine and Dominium as the game/product suite built on top of i...
\end{itemize}

\section{Dominium Architecture, UI, Providers, and Robot OS Strategy}

\begin{itemize}
\item Source folder: docs/archive/conversations/architecture\_ui\_providers/
\item Reader page: docs/archive/conversations/\_reader/by\_chat/architecture\_ui\_providers.md
\item Primary source: docs/archive/conversations/architecture\_ui\_providers/dominium\_architecture\_ui\_providers\_robot\_os\_strategy\_\_01\_human\_readable\_report.md
\item Completeness: complete
\item Topics: architecture, content, contracts\_schema, determinism, governance, platform, release, simulation, tooling, ui, workbench, worldgen
\item Digest: This chat was a strategic architecture and design convergence session for Dominium/Domino. It began from a transferred knowledge base about rebuilding Dominium's GUI and binary strategy, then evolved into a broad platform-specification discussion covering repository structure, product shells, renderer/platform priorities, native applications, language standards, provider architecture, Workbench, UI authoring, and...
\end{itemize}

\section{Dominium Build and Future-Proofing Architecture}

\begin{itemize}
\item Source folder: docs/archive/conversations/Build\_and\_Future\_Proofing/
\item Reader page: docs/archive/conversations/\_reader/by\_chat/build\_and\_future\_proofing.md
\item Primary source: docs/archive/conversations/Build\_and\_Future\_Proofing/Dominium\_Build\_and\_Future\_Proofing\_Architecture\_\_01\_human\_readable\_report.md
\item Completeness: complete
\item Topics: architecture, content, contracts\_schema, determinism, governance, platform, release, simulation, tooling, ui, worldgen
\item Digest: This chat centered on how to make Dominium durable as a serious long-term software project rather than a one-off game prototype. The user first provided a prior technical assessment about Dominium's Windows build floors, Visual Studio toolchains, XP/Win7/Win10 support, CMake presets, static CRT policy, and runtime testing requirements. That pasted material treated Dominium's canon as strict: engine in C89, game la...
\end{itemize}

\section{Dominium Canonical Architecture, Repository Foundation, and Provider Model}

\begin{itemize}
\item Source folder: docs/archive/conversations/canonical\_structure\_and\_framework/
\item Reader page: docs/archive/conversations/\_reader/by\_chat/canonical\_structure\_and\_framework.md
\item Primary source: docs/archive/conversations/canonical\_structure\_and\_framework/dominium\_canonical\_architecture\_repo\_foundation\_\_01\_human\_readable\_report.md
\item Completeness: complete
\item Topics: architecture, content, contracts\_schema, determinism, governance, platform, release, setup\_launcher, simulation, timekeeping, tooling, ui, workbench, xstack\_aide
\item Digest: This chat was an extended architecture, repository-structure, and project-governance conversation for the Dominium game/project. The user's core concern was that months of planning and refactoring had stalled actual development and that the repository had repeatedly grown new root directories, duplicate taxonomies, generic src//source/ wrappers, and unclear ownership boundaries. The user wanted to force the pr...
\end{itemize}

\section{Dominium Chronology \& Celestial Systems}

\begin{itemize}
\item Source folder: docs/archive/conversations/Chronology\_Celestial\_Systems/
\item Reader page: docs/archive/conversations/\_reader/by\_chat/chronology\_celestial\_systems.md
\item Primary source: docs/archive/conversations/Chronology\_Celestial\_Systems/Dominium\_Chronology\_Celestial\_Systems\_\_human\_readable\_chat\_report.txt
\item Completeness: complete
\item Topics: architecture, content, contracts\_schema, governance, platform, simulation, timekeeping, tooling, ui, worldgen
\item Digest: This chat was mainly about designing a major subsystem for the Dominium Game project: how the game should represent celestial bodies, astronomical structures, time, calendars, dates, clocks, local standards, diegetic HUD knowledge, and large-scale chronology across Earth, the Sol system, the Milky Way, and the universe. The conversation began with a practical worldbuilding and simulation-design question: if the ga...
\end{itemize}

\section{Dominium Development Routes and Continuity Preservation}

\begin{itemize}
\item Source folder: docs/archive/conversations/development\_routes/
\item Reader page: docs/archive/conversations/\_reader/by\_chat/development\_routes.md
\item Primary source: docs/archive/conversations/development\_routes/dominium\_development\_routes\_\_human\_readable\_chat\_report.txt
\item Completeness: complete
\item Topics: architecture, content, contracts\_schema, determinism, governance, release, simulation, timekeeping, tooling, ui, workbench, worldgen
\item Digest: This chat had two main phases. The first phase was a technical planning discussion about **Dominium Game**. The user asked what the usual development routes are for games like Dominium, both technically and feature-wise, and what specific components should be worked on in what order. The answer given in that phase proposed that Dominium should be treated not primarily as a normal game with features layered on top,...
\end{itemize}

\section{Documentation Standards, README Strategy, and Handoff Packaging}

\begin{itemize}
\item Source folder: docs/archive/conversations/documentation\_standards\_readme/
\item Reader page: docs/archive/conversations/\_reader/by\_chat/documentation\_standards\_readme.md
\item Primary source: docs/archive/conversations/documentation\_standards\_readme/documentation\_standards\_readme\_handoff\_\_human\_readable\_chat\_report.txt
\item Completeness: complete
\item Topics: architecture, content, contracts\_schema, determinism, governance, platform, release, setup\_launcher, simulation, timekeeping, tooling, ui
\item Digest: This chat was about turning documentation from an informal afterthought into a disciplined, enforceable part of a long-lived C/C++ systems project. The broader project context is PROJECT-CONTEXT: Dominium Game, but this chat itself did not inspect the repository, modify files, or verify the current codebase. Its value is that it established documentation architecture, Codex execution prompts, a README strategy,...
\end{itemize}

\section{Dominium Architecture I}

\begin{itemize}
\item Source folder: docs/archive/conversations/Dominium\_Architecture\_I/
\item Reader page: docs/archive/conversations/\_reader/by\_chat/dominium\_architecture\_i.md
\item Primary source: docs/archive/conversations/Dominium\_Architecture\_I/Dominium\_Architecture\_I\_\_human\_readable\_chat\_report.txt
\item Completeness: complete
\item Topics: architecture, content, contracts\_schema, determinism, governance, platform, release, setup\_launcher, simulation, timekeeping, tooling, ui, worldgen
\item Digest: **FACT:** This chat was mainly about turning the developing architecture of the **Dominium Game** into a form that could be carried forward into automated coding and later Project-level consolidation. It was not a normal design discussion by the end. It became a long-running specification-production session, where the user repeatedly asked for detailed Markdown documents and then for one .txt implementation-spec...
\end{itemize}

\section{Dominium Architecture II}

\begin{itemize}
\item Source folder: docs/archive/conversations/Dominium\_Architecture\_II/
\item Reader page: docs/archive/conversations/\_reader/by\_chat/dominium\_architecture\_ii.md
\item Primary source: docs/archive/conversations/Dominium\_Architecture\_II/Dominium\_Architecture\_II\_\_human\_readable\_chat\_report.txt
\item Completeness: complete
\item Topics: architecture, content, contracts\_schema, determinism, governance, platform, release, simulation, timekeeping, tooling, ui, worldgen
\item Digest: This chat was mainly about turning the **Dominium** game/engine from a broad, ambitious design into a coherent, Codex-ready architecture. The user had already explored many individual systems elsewhere or earlier in the same chat: railways, roadways, waterways, airways, spaceways, electrical networks, data networks, building systems, terrain, cut/fill, tunnels, bridges, arbitrary splines, deterministic simulation,...
\end{itemize}

\section{Dominium Architecture III: Launcher, Platform, Renderer, and Handoff Architecture}

\begin{itemize}
\item Source folder: docs/archive/conversations/Dominium\_Architecture\_III/
\item Reader page: docs/archive/conversations/\_reader/by\_chat/dominium\_architecture\_iii.md
\item Primary source: docs/archive/conversations/Dominium\_Architecture\_III/Dominium\_Architecture\_III\_\_human\_readable\_chat\_report.txt
\item Completeness: complete
\item Topics: architecture, content, determinism, governance, platform, release, setup\_launcher, simulation, tooling, ui, worldgen
\item Digest: This chat was mainly about turning Dominium's launcher from a vague "startup UI" into a coherent, portable, extensible subsystem that fits the wider deterministic simulation architecture. The user began with the idea of tackling the launcher first, then progressively refined the requirements until the conversation covered launcher architecture, platform and renderer backends, cross-platform entry scripts, software...
\end{itemize}

\section{Dominium Architecture IV}

\begin{itemize}
\item Source folder: docs/archive/conversations/Dominium\_Architecture\_IV/
\item Reader page: docs/archive/conversations/\_reader/by\_chat/dominium\_architecture\_iv.md
\item Primary source: docs/archive/conversations/Dominium\_Architecture\_IV/Dominium\_Architecture\_IV\_\_human\_readable\_chat\_report.txt
\item Completeness: complete
\item Topics: architecture, content, contracts\_schema, determinism, governance, platform, release, setup\_launcher, simulation, timekeeping, tooling, ui, workbench, worldgen
\item Digest: This chat was mainly about turning **Dominium** from a broad game idea into a deeply planned software ecosystem. The conversation moved across engine architecture, platform abstraction, renderer abstraction, install modes, launcher design, setup packaging, tools and SDKs, game simulation architecture, construction systems, vehicles, spaceflight, planets, economy, actors, and long-term modding. The most important o...
\end{itemize}

\section{Dominium Canon, Repository Alignment, and Portability Doctrine}

\begin{itemize}
\item Source folder: docs/archive/conversations/Dominium\_Complete\_Conversation/
\item Reader page: docs/archive/conversations/\_reader/by\_chat/dominium\_complete\_conversation.md
\item Primary source: docs/archive/conversations/Dominium\_Complete\_Conversation/Accompanying\_Report/Dominium\_Canon\_Portability\_Audit\_\_01\_human\_readable\_report.md
\item Completeness: complete
\item Topics: architecture, content, contracts\_schema, determinism, governance, platform, release, simulation, timekeeping, tooling, ui, worldgen, xstack\_aide
\item Digest: This chat was about preserving and re-grounding Dominium's architecture around a very old constitutional contract, then checking whether the current GitHub repository still follows that contract, and finally broadening the question into a platform-engineering doctrine for portability, modularity, extensibility, reuse, future-proofing, and long-term compatibility. The chat began with the user pasting the old "DOMIN...
\end{itemize}

\section{Dominium + Domino Codex Planning and Handoff}

\begin{itemize}
\item Source folder: docs/archive/conversations/dominium\_domino\_codex\_planning/
\item Reader page: docs/archive/conversations/\_reader/by\_chat/dominium\_domino\_codex\_planning.md
\item Primary source: docs/archive/conversations/dominium\_domino\_codex\_planning/dominium\_domino\_codex\_planning\_\_human\_readable\_chat\_report.txt
\item Completeness: complete
\item Topics: architecture, content, contracts\_schema, determinism, governance, platform, release, setup\_launcher, simulation, timekeeping, tooling, ui, workbench, worldgen
\item Digest: This chat was mainly about preserving and developing the implementation plan for the **Dominium + Domino** project. **FACT:** The user began by giving an extended master prompt that defined the architecture, philosophy, constraints, and product model for Dominium and Domino. In that project model, **Domino** is the deterministic, fixed-point, ISO C89 engine layer, and **Dominium** is the C++98 product suite built...
\end{itemize}

\section{Dominium Workbench, AIDE, Presentation Architecture, Provider Strategy, Product-Spine Planning, and Preservation Companion}

\begin{itemize}
\item Source folder: docs/archive/conversations/dominium\_full\_conversation/
\item Reader page: docs/archive/conversations/\_reader/by\_chat/dominium\_full\_conversation.md
\item Primary source: docs/archive/conversations/dominium\_full\_conversation/companion\_report/dominium\_full\_conversation\_companion\_report\_\_01\_complete\_human\_report.md
\item Completeness: complete
\item Topics: architecture, content, contracts\_schema, determinism, governance, platform, release, setup\_launcher, simulation, timekeeping, tooling, ui, workbench, worldgen, xstack\_aide
\item Digest: This conversation records a major architectural transition in the Dominium project. It began with an immediate practical concern: the existing Windows launcher UI looked broken, flickered, and needed a better way to design UIs. The initial direction was to create a Windows-first UI Editor and eventually a Tool Editor that could visually design pixel-perfect UIs and generate TLV-based layouts/code. That early plan...
\end{itemize}

\section{Dominium Setup Architecture and Handoff}

\begin{itemize}
\item Source folder: docs/archive/conversations/dominium\_setup/
\item Reader page: docs/archive/conversations/\_reader/by\_chat/dominium\_setup.md
\item Primary source: docs/archive/conversations/dominium\_setup/dominium\_setup\_handoff\_\_human\_readable\_chat\_report.txt
\item Completeness: complete
\item Topics: architecture, content, contracts\_schema, determinism, governance, platform, release, setup\_launcher, simulation, timekeeping, tooling, ui, worldgen
\item Digest: This chat was mainly about designing, correcting, and preserving the **Dominium / Domino Setup system**: the installer, updater, verifier, rollback, repair, uninstall, packaging, and distribution layer for the larger Dominium project. The key point is that this was **not** a launcher UI chat, not an engine-design chat, and not a gameplay/simulation-design chat. It focused on the application-layer product called `s...
\end{itemize}

\section{Domino Dominium Workbench}

\begin{itemize}
\item Source folder: docs/archive/conversations/Domino\_Dominium\_Workbench/
\item Reader page: docs/archive/conversations/\_reader/by\_chat/domino\_dominium\_workbench.md
\item Primary source: docs/archive/conversations/Domino\_Dominium\_Workbench/dominium\_workbench\_full\_conversation\_companion\_\_human\_readable\_detailed\_summary\_report.md
\item Completeness: complete
\item Topics: architecture, content, contracts\_schema, determinism, governance, platform, setup\_launcher, simulation, timekeeping, tooling, ui, workbench, worldgen
\item Digest: This conversation started as a practical plan to build a Windows-first internal UI Editor / Tool Editor to fix the Dominium launcher UI and visually author setup, launcher, game, and tool interfaces. It then evolved into a much larger and more durable architecture: **Dominium Workbench Platform**, a cross-platform rendered, TUI-capable, CLI/headless-compatible, modular production environment for building the entir...
\end{itemize}

\section{Dominium/Domino Engine Refactor Planning}

\begin{itemize}
\item Source folder: docs/archive/conversations/domino\_engine\_refactor\_prompts/
\item Reader page: docs/archive/conversations/\_reader/by\_chat/domino\_engine\_refactor\_prompts.md
\item Primary source: docs/archive/conversations/domino\_engine\_refactor\_prompts/dominium\_domino\_engine\_refactor\_prompts\_\_human\_readable\_chat\_report.txt
\item Completeness: complete
\item Topics: architecture, content, contracts\_schema, determinism, governance, platform, release, simulation, tooling, ui, worldgen
\item Digest: This chat was a long-form architecture and planning session for the **Dominium/Domino deterministic engine**. Its central purpose was not to write code, but to define how a very large, deterministic simulation engine should be structured so that it can support complex world simulation, agents, infrastructure, future DLCs, modding, replay, multiplayer lockstep, and large-scale performance without collapsing into ha...
\end{itemize}

\section{Domino/Dominium Engine Baseline, Architecture, and Feasibility}

\begin{itemize}
\item Source folder: docs/archive/conversations/engine\_baseline\_architecture/
\item Reader page: docs/archive/conversations/\_reader/by\_chat/engine\_baseline\_architecture.md
\item Primary source: docs/archive/conversations/engine\_baseline\_architecture/domino\_dominium\_engine\_baseline\_architecture\_\_01\_human\_readable\_report.md
\item Completeness: complete
\item Topics: architecture, content, contracts\_schema, determinism, governance, platform, release, simulation, timekeeping, tooling, ui, workbench
\item Digest: This chat was about the architecture, feasibility, state, and future direction of the Domino engine and Dominium game. The user wanted to understand what the project is, whether the engine exists in a meaningful sense, whether it is ready to support a game, how it can become portable and modular enough to be reused for multiple games or even CAD-like tools, and how to sequence the work without being overwhelmed by...
\end{itemize}

\section{Dominium Foundation Lock, Workbench Spine, and Parallel Codex Handoff}

\begin{itemize}
\item Source folder: docs/archive/conversations/Foundation\_Workbench\_Codex/
\item Reader page: docs/archive/conversations/\_reader/by\_chat/foundation\_workbench\_codex.md
\item Primary source: docs/archive/conversations/Foundation\_Workbench\_Codex/Dominium\_Foundation\_Workbench\_Parallel\_Codex\_Preservation\_\_01\_human\_readable\_report.md
\item Completeness: complete
\item Topics: architecture, content, contracts\_schema, determinism, governance, platform, release, simulation, timekeeping, tooling, ui, workbench, worldgen, xstack\_aide
\item Digest: This chat was a high-intensity strategic, architectural, and operational continuation of the Dominium project: a large simulation/game/engine ecosystem in the Julesc013/dominium repository. The user's central concern was that the project had spent a very long time in refactor and restructuring work, and that without a clear, enforceable architecture and faster execution model, the project would either drift back...
\end{itemize}

\section{Domino Framework and Open-Source Provider Architecture}

\begin{itemize}
\item Source folder: docs/archive/conversations/Framework\_Open\_Source\_Provider/
\item Reader page: docs/archive/conversations/\_reader/by\_chat/framework\_open\_source\_provider.md
\item Primary source: docs/archive/conversations/Framework\_Open\_Source\_Provider/Domino\_Framework\_Open\_Source\_Provider\_Architecture\_\_01\_human\_readable\_report.md
\item Completeness: complete
\item Topics: architecture, content, contracts\_schema, determinism, governance, platform, release, simulation, tooling, ui, workbench, worldgen
\item Digest: This chat was about how to accelerate the Dominium/Domino game-engine project by using existing open-source code without letting outside engines or libraries define the architecture. The user's recurring concern was speed: instead of starting from nothing, could Dominium bootstrap a working engine, game client, Workbench, scripting layer, and provider ecosystem using open-source systems such as raylib, SDL2, Lua,...
\end{itemize}

\section{Dominium Content and GUI Rebuild Planning}

\begin{itemize}
\item Source folder: docs/archive/conversations/gui\_binary\_content/
\item Reader page: docs/archive/conversations/\_reader/by\_chat/gui\_binary\_content.md
\item Primary source: docs/archive/conversations/gui\_binary\_content/dominium\_gui\_binary\_content\_handoff\_\_human\_readable\_chat\_report.txt
\item Completeness: complete
\item Topics: architecture, content, contracts\_schema, determinism, governance, platform, release, setup\_launcher, simulation, timekeeping, tooling, ui, workbench, worldgen
\item Digest: This chat was mainly about preserving and clarifying two major planning threads for the **Dominium / Domino** project: first, a canonical **content/data population** effort for the game universe, and second, a later and more active effort to **rebuild the project's GUI and binary strategy from scratch**. The conversation began with a large prompt called **CONTENT0**, intended for Codex-style work on the project's...
\end{itemize}

\section{Dominium Language, Platform, and Architecture Baseline}

\begin{itemize}
\item Source folder: docs/archive/conversations/Language\_Platform\_Architecture/
\item Reader page: docs/archive/conversations/\_reader/by\_chat/language\_platform\_architecture.md
\item Primary source: docs/archive/conversations/Language\_Platform\_Architecture/Dominium\_Language\_Platform\_Architecture\_Baseline\_\_01\_human\_readable\_report.md
\item Completeness: complete
\item Topics: architecture, content, contracts\_schema, determinism, governance, platform, release, simulation, timekeeping, tooling, ui, workbench, worldgen, xstack\_aide
\item Digest: This chat was a design-decision and architecture-baseline discussion for Dominium, focused on language standards, platform floors, binary architecture, portability strategy, and the project's long-term modular structure. The user began by asking whether Dominium would gain performance or features by moving beyond the earlier C89 engine and C++98 game plan. Over the conversation, the topic expanded into a full re-e...
\end{itemize}

\section{Dominium Launcher Application-Layer Handoff}

\begin{itemize}
\item Source folder: docs/archive/conversations/launcher\_app\_layer/
\item Reader page: docs/archive/conversations/\_reader/by\_chat/launcher\_app\_layer.md
\item Primary source: docs/archive/conversations/launcher\_app\_layer/dominium\_launcher\_app\_layer\_handoff\_\_human\_readable\_chat\_report.txt
\item Completeness: complete
\item Topics: architecture, content, contracts\_schema, determinism, governance, platform, release, setup\_launcher, simulation, tooling, ui, worldgen
\item Digest: This chat was mainly about turning the **Dominium Launcher** from a broad architectural idea into a properly bounded, enforceable, application-layer product inside the larger **Dominium / Domino** project. The conversation began with a practical launcher question: whether the latest post-Plan 8 system could support a **single cross-platform launcher** with a **common codebase**, **native OS elements**, and **littl...
\end{itemize}

\section{Dominium Launcher and Setup Architecture}

\begin{itemize}
\item Source folder: docs/archive/conversations/Launcher\_Setup\_Architecture/
\item Reader page: docs/archive/conversations/\_reader/by\_chat/launcher\_setup\_architecture.md
\item Primary source: docs/archive/conversations/Launcher\_Setup\_Architecture/Dominium\_Launcher\_Setup\_Architecture\_\_human\_readable\_chat\_report.txt
\item Completeness: complete
\item Topics: architecture, content, contracts\_schema, determinism, governance, platform, release, setup\_launcher, simulation, timekeeping, tooling, ui, worldgen
\item Digest: This chat was about designing the **Dominium launcher, setup system, runtime integration model, and surrounding user-facing infrastructure**. It began from a broad architectural philosophy for the Dominium project: keep the simulation core deterministic, boring, portable, and C89; put explicit APIs between layers; use versioned formats; make mods data-first; and prevent tools, launcher, and runtime code from drift...
\end{itemize}

\section{Dominium Modularity, AIDE Refactorability, and Future-Proof Architecture}

\begin{itemize}
\item Source folder: docs/archive/conversations/Modularity\_AIDE\_Refactorability/
\item Reader page: docs/archive/conversations/\_reader/by\_chat/modularity\_aide\_refactorability.md
\item Primary source: docs/archive/conversations/Modularity\_AIDE\_Refactorability/Dominium\_Modularity\_AIDE\_Refactorability\_Architecture\_\_01\_human\_readable\_report.md
\item Completeness: complete
\item Topics: architecture, content, contracts\_schema, determinism, governance, platform, release, simulation, timekeeping, tooling, ui, worldgen, xstack\_aide
\item Digest: This chat was about making Dominium structurally durable: portable, modular, extensible, reusable across future games, and refactorable without the user having to repeat a painful top-level cleanup. The discussion began from an already-developed architectural thread about Dominium's source layout, distribution layout, package layout, install roots, runtime roots, bundle formats, and component matrices. The user br...
\end{itemize}

\section{Dominium Omega/Xi/Pi Architecture \& Future-Proofing Planning}

\begin{itemize}
\item Source folder: docs/archive/conversations/omega\_xi\_pi\_architecture\_future/
\item Reader page: docs/archive/conversations/\_reader/by\_chat/omega\_xi\_pi\_architecture\_future.md
\item Primary source: docs/archive/conversations/omega\_xi\_pi\_architecture\_future/dominium\_omega\_xi\_pi\_architecture\_future\_proofing\_planning\_\_01\_human\_readable\_report.md
\item Completeness: complete
\item Topics: architecture, content, contracts\_schema, governance, platform, release, setup\_launcher, simulation, timekeeping, tooling, ui, worldgen, xstack\_aide
\item Digest: This chat was a long-running architectural planning and continuity-preservation conversation for the user's Dominium project. Dominium began in the conversation as a game/engine project, but the user and assistant progressively reframed it as something larger: a deterministic simulation platform, reusable engine, runtime/service host, package ecosystem, product family, agent/human development environment, and rele...
\end{itemize}

\section{Dominium OS-like Architecture, Repository Convergence, and Interface Operating Layer}

\begin{itemize}
\item Source folder: docs/archive/conversations/OS\_Interface\_Repo\_Architecture/
\item Reader page: docs/archive/conversations/\_reader/by\_chat/os\_interface\_repo\_architecture.md
\item Primary source: docs/archive/conversations/OS\_Interface\_Repo\_Architecture/Dominium\_OS\_Interface\_Repo\_Architecture\_\_01\_human\_readable\_report.md
\item Completeness: complete
\item Topics: architecture, content, contracts\_schema, determinism, governance, platform, release, setup\_launcher, simulation, tooling, ui, workbench
\item Digest: This chat was about rethinking Dominium from a conventional game project into a much broader, more modular, deterministic, extensible software environment. The user repeatedly asked whether the current plan was the best possible version, and the discussion evolved from GUI/binary planning, to repository restructure, to release/package strategy, to code-vs-data analysis, to an "OS-like" deterministic simulation ope...
\end{itemize}

\section{Dominium Codex Platform Renderer API Plan}

\begin{itemize}
\item Source folder: docs/archive/conversations/platform\_renderer\_api\_plan/
\item Reader page: docs/archive/conversations/\_reader/by\_chat/platform\_renderer\_api\_plan.md
\item Primary source: docs/archive/conversations/platform\_renderer\_api\_plan/dominium\_codex\_platform\_renderer\_api\_plan\_\_human\_readable\_chat\_report.txt
\item Completeness: complete
\item Topics: architecture, content, contracts\_schema, determinism, governance, platform, release, setup\_launcher, simulation, timekeeping, tooling, ui, workbench, worldgen
\item Digest: This chat was mainly about turning the user's Dominium/Domino engine work into a **high-quality Codex implementation plan**. The user wanted prompts for a coding agent, specifically "Codex CLI running GPT-5.2" or VS Code Codex, to work inside an existing C/CMake/C++ repository and implement missing or unfinished **platforms, renderers, APIs, and integration smoke tests**. The repository path used throughout the pr...
\end{itemize}

\section{Dominium Platform Support Planning}

\begin{itemize}
\item Source folder: docs/archive/conversations/platform\_support/
\item Reader page: docs/archive/conversations/\_reader/by\_chat/platform\_support.md
\item Primary source: docs/archive/conversations/platform\_support/dominium\_platform\_support\_planning\_\_human\_readable\_chat\_report.txt
\item Completeness: complete
\item Topics: architecture, content, contracts\_schema, governance, platform, release, simulation, timekeeping, ui, worldgen
\item Digest: This chat was about deciding what platforms **Dominium** should support and how those targets should be organised without collapsing into an unrealistic "support everything" claim. The conversation started with a broad question about "what systems" the game would support, then narrowed to operating systems, then expanded again into a much wider device-family strategy covering PC, consoles, handhelds, mobile, Web,...
\end{itemize}

\section{Domino/Dominium Portability, Assurance, and Future-Proof Architecture}

\begin{itemize}
\item Source folder: docs/archive/conversations/Portability\_Assurance\_Future\_Proof/
\item Reader page: docs/archive/conversations/\_reader/by\_chat/portability\_assurance\_future\_proof.md
\item Primary source: docs/archive/conversations/Portability\_Assurance\_Future\_Proof/Domino\_Dominium\_Portability\_Assurance\_Future\_Proof\_Architecture\_\_01\_human\_readable\_report.md
\item Completeness: complete
\item Topics: architecture, content, contracts\_schema, determinism, governance, release, setup\_launcher, simulation, timekeeping, tooling, ui, worldgen
\item Digest: This chat was mainly about turning Domino/Dominium from a one-off game project into a durable, reusable engine/product platform. It began with a question about whether standards from high-assurance domains-especially DO-178B/C and adjacent secure-development, supply-chain, and metadata standards-could be used for Domino and Dominium. The proposed answer was that the project should use standards as design inputs, n...
\end{itemize}

\section{Dominium README Architecture, Ports, and Determinism}

\begin{itemize}
\item Source folder: docs/archive/conversations/readme\_ports\_determinism/
\item Reader page: docs/archive/conversations/\_reader/by\_chat/readme\_ports\_determinism.md
\item Primary source: docs/archive/conversations/readme\_ports\_determinism/dominium\_readme\_ports\_determinism\_\_human\_readable\_chat\_report.txt
\item Completeness: complete
\item Topics: architecture, content, contracts\_schema, determinism, governance, platform, release, simulation, timekeeping, tooling, ui, worldgen
\item Digest: This chat was mainly about refining the root README.md for the **Dominium / Domino** project and then preserving the resulting decisions in a form that could survive into later chats. The project, as described in the README, is a deterministic, integer-math, multi-platform simulation game and engine. **Domino** is the deterministic simulation engine core. **Dominium** is the official game, runtime, tooling, and...
\end{itemize}

\section{Dominium + Domino Refactor Architecture}

\begin{itemize}
\item Source folder: docs/archive/conversations/Refactor\_Architecture/
\item Reader page: docs/archive/conversations/\_reader/by\_chat/refactor\_architecture.md
\item Primary source: docs/archive/conversations/Refactor\_Architecture/Dominium\_Domino\_Refactor\_Architecture\_\_human\_readable\_chat\_report.txt
\item Completeness: complete
\item Topics: architecture, content, contracts\_schema, platform, release, setup\_launcher, simulation, timekeeping, tooling, ui, workbench
\item Digest: This chat was mainly about turning the Dominium + Domino project from a collection of partially overlapping architecture ideas into a coherent, refactorable, future-proof structure that could be implemented by Codex and then preserved for future ChatGPT conversations. At the highest level, the discussion separated two names and responsibilities. **Domino** is the engine. It lives under source/domino, provides de...
\end{itemize}

\section{AIDE, XStack, and Dominium Refactor Control Plane}

\begin{itemize}
\item Source folder: docs/archive/conversations/Refactor\_Control\_Plane/
\item Reader page: docs/archive/conversations/\_reader/by\_chat/refactor\_control\_plane.md
\item Primary source: docs/archive/conversations/Refactor\_Control\_Plane/AIDE\_XStack\_Dominium\_Refactor\_Control\_Plane\_\_01\_human\_readable\_report.md
\item Completeness: complete
\item Topics: architecture, content, contracts\_schema, governance, release, simulation, timekeeping, tooling, ui, worldgen, xstack\_aide
\item Digest: This chat was a long-running architecture, workflow, and project-control discussion around how to use AI coding tools more effectively, how to evolve Dominium's XStack governance system, and how to build AIDE as a portable repo-native operating layer. It began with practical questions about using Codex Pro, the Codex VS Code extension, ChatGPT project spaces, the ChatGPT/Codex apps, Visual Studio, and the API. The...
\end{itemize}

\section{Dominium XStack Release Identity and Versioning}

\begin{itemize}
\item Source folder: docs/archive/conversations/Release\_Identity\_and\_Versioning/
\item Reader page: docs/archive/conversations/\_reader/by\_chat/release\_identity\_and\_versioning.md
\item Primary source: docs/archive/conversations/Release\_Identity\_and\_Versioning/Dominium\_XStack\_Release\_Identity\_and\_Versioning\_\_01\_human\_readable\_report.md
\item Completeness: complete
\item Topics: architecture, content, contracts\_schema, determinism, governance, platform, release, setup\_launcher, simulation, timekeeping, tooling, ui, worldgen, xstack\_aide
\item Digest: This chat was about designing a durable release identity, versioning, compatibility, build, and packaging system for Dominium/XStack. The user began from dissatisfaction with products that change versioning policies midstream, citing examples such as Windows NT, Minecraft, macOS, .NET, Linux, and TeX. The underlying concern was not aesthetic only. The user wanted a versioning system that could survive long-term wi...
\end{itemize}

\section{Dominium TestX/RepoX Governance and Handoff Chat}

\begin{itemize}
\item Source folder: docs/archive/conversations/testx\_repox\_governance/
\item Reader page: docs/archive/conversations/\_reader/by\_chat/testx\_repox\_governance.md
\item Primary source: docs/archive/conversations/testx\_repox\_governance/dominium\_testx\_repox\_governance\_handoff\_\_human\_readable\_chat\_report.txt
\item Completeness: complete
\item Topics: architecture, content, determinism, governance, platform, release, setup\_launcher, simulation, tooling, ui, worldgen, xstack\_aide
\item Digest: This chat was mainly about turning the Dominium / Domino project from a large, ambitious architecture into something that could survive implementation across many parallel chats, many products, many platforms, many toolchains, and many years of future changes. The conversation was not primarily about gameplay content, graphics, or a single feature. It was about **governance**: how the repository, build system, tes...
\end{itemize}

\section{Dominium Timekeeping and 2038 Resilience}

\begin{itemize}
\item Source folder: docs/archive/conversations/Timekeeping\_and\_2038\_Resilience/
\item Reader page: docs/archive/conversations/\_reader/by\_chat/timekeeping\_and\_2038\_resilience.md
\item Primary source: docs/archive/conversations/Timekeeping\_and\_2038\_Resilience/Dominium\_Timekeeping\_and\_2038\_Resilience\_\_01\_human\_readable\_report.md
\item Completeness: complete
\item Topics: architecture, content, contracts\_schema, determinism, governance, platform, release, simulation, timekeeping, tooling, ui, worldgen
\item Digest: This chat was mainly about making timekeeping code survive old platforms, future dates, deterministic simulation requirements, and project-scale architecture constraints. It began with a practical question: what happens to 32-bit operating systems and 32-bit applications after 2038? The user wanted to know whether they would stop working completely, whether critical functions would fail, whether rolling back the c...
\end{itemize}

\section{Dominium UE6, Domino, and Deterministic Universe Feasibility}

\begin{itemize}
\item Source folder: docs/archive/conversations/UE6\_Domino\_Deterministic\_Universe/
\item Reader page: docs/archive/conversations/\_reader/by\_chat/ue6\_domino\_deterministic\_universe.md
\item Primary source: docs/archive/conversations/UE6\_Domino\_Deterministic\_Universe/Dominium\_UE6\_Domino\_Deterministic\_Universe\_\_01\_human\_readable\_report.md
\item Completeness: complete
\item Topics: architecture, content, determinism, governance, platform, release, simulation, timekeeping, tooling, ui, worldgen
\item Digest: This chat was mainly about whether Dominium could or should be built on Unreal Engine 6, Unreal Engine 5, Domino, or some custom architecture that uses a commercial engine only as a frontend. The user began by pasting a confident external-style answer claiming that UE6 had been unveiled, that UE6 would raise baseline OS requirements, and that high-performance UE6 games would probably require modern Windows, macOS,...
\end{itemize}

\section{Dominium UI Editor and Tool Editor Planning}

\begin{itemize}
\item Source folder: docs/archive/conversations/ui\_editor\_tool\_editor\_planning/
\item Reader page: docs/archive/conversations/\_reader/by\_chat/ui\_editor\_tool\_editor\_planning.md
\item Primary source: docs/archive/conversations/ui\_editor\_tool\_editor\_planning/dominium\_ui\_editor\_tool\_editor\_planning\_\_human\_readable\_chat\_report.txt
\item Completeness: complete
\item Topics: architecture, content, contracts\_schema, determinism, governance, platform, release, setup\_launcher, simulation, tooling, ui, workbench, worldgen
\item Digest: This chat was about designing a long-term UI authoring system for the **Dominium / Domino** project. The immediate practical problem was that the current Windows launcher UI "looks mangled and flickers." The broader goal was much larger: create an internal graphical toolchain that can visually design and maintain user interfaces for the launcher, setup program, game, and other Dominium tools, while preserving nati...
\end{itemize}

\section{Dominium Universe Explorer Planning and Repo Handoff}

\begin{itemize}
\item Source folder: docs/archive/conversations/Universe\_Explorer\_Planning/
\item Reader page: docs/archive/conversations/\_reader/by\_chat/universe\_explorer\_planning.md
\item Primary source: docs/archive/conversations/Universe\_Explorer\_Planning/Dominium\_Universe\_Explorer\_Planning\_Handoff\_\_01\_human\_readable\_report.md
\item Completeness: complete
\item Topics: architecture, content, contracts\_schema, governance, platform, release, simulation, tooling, ui, workbench, worldgen, xstack\_aide
\item Digest: This chat was a high-level Dominium planning, recovery, repo-audit, and forward-strategy session. The user began by pasting substantial prior discussion from earlier Dominium chats. Those excerpts covered the project's deep simulation architecture: total truth with sparse materialization, assemblies/fields/processes/constraints, epistemic and diegetic refinement, classifier-based emergence, macro-to-micro scaling,...
\end{itemize}

\section{Dominium Architecture, Workbench, AIDE, and Product-Spine Planning}

\begin{itemize}
\item Source folder: docs/archive/conversations/Workbench\_AIDE\_Product\_Spine/
\item Reader page: docs/archive/conversations/\_reader/by\_chat/workbench\_aide\_product\_spine.md
\item Primary source: docs/archive/conversations/Workbench\_AIDE\_Product\_Spine/Dominium\_Architecture\_Workbench\_AIDE\_Product\_Spine\_Planning\_\_01\_human\_readable\_report.md
\item Completeness: complete
\item Topics: architecture, content, contracts\_schema, determinism, governance, platform, release, setup\_launcher, simulation, timekeeping, tooling, ui, workbench, worldgen, xstack\_aide
\item Digest: This chat was primarily about moving Dominium from a large, ambitious but risk-prone game/simulation project into a disciplined, contract-governed, deterministic simulation platform with a reusable engine substrate, a production Workbench, and an AIDE-driven development workflow. The user's underlying concern was that Dominium should not become a one-off indie project with ad hoc code, drifting folders, duplicated...
\end{itemize}

\section{Dominium World Architecture}

\begin{itemize}
\item Source folder: docs/archive/conversations/World\_Architecture/
\item Reader page: docs/archive/conversations/\_reader/by\_chat/world\_architecture.md
\item Primary source: docs/archive/conversations/World\_Architecture/Dominium\_World\_Architecture\_\_human\_readable\_chat\_report.txt
\item Completeness: complete
\item Topics: architecture, content, contracts\_schema, determinism, governance, release, setup\_launcher, simulation, timekeeping, tooling, ui, workbench, worldgen
\item Digest: This chat was mainly about designing the foundational world architecture for **Dominium**, especially its spatial model, coordinate system, save format, terrain representation, environmental simulation model, deterministic core constraints, mod/content architecture, and the path toward implementation through Codex. At the beginning, the discussion was narrow: the user was choosing fixed-point coordinate precision...
\end{itemize}

\section{Dominium XStack and Lab Galaxy Handoff}

\begin{itemize}
\item Source folder: docs/archive/conversations/xstack\_lab\_galaxy/
\item Reader page: docs/archive/conversations/\_reader/by\_chat/xstack\_lab\_galaxy.md
\item Primary source: docs/archive/conversations/xstack\_lab\_galaxy/dominium\_xstack\_lab\_galaxy\_handoff\_\_human\_readable\_chat\_report.txt
\item Completeness: complete
\item Topics: architecture, content, contracts\_schema, determinism, governance, release, timekeeping, tooling, ui, worldgen, xstack\_aide
\item Digest: This chat was mainly about turning a very large, complex Dominium development conversation into a stable, reusable handoff, while preserving the architectural decisions and operational context needed to continue work without re-explaining everything. The conversation sat at the intersection of three big themes: **Dominium / Domino as a deterministic universe simulation/game**, **XStack as a portable agentic develo...
\end{itemize}


{\small\raggedright SOURCE PATH: docs/\allowbreak{}archive/\allowbreak{}conversations/\allowbreak{}\_\allowbreak{}reader/\allowbreak{}conversation\_\allowbreak{}reader\_\allowbreak{}index.md\par}


\section{Conversation Reader Index}

Reader pages are prose-first advisory views over archived conversations.



\begin{quote}\small Dense table summarized for the reader edition. See the source file, HTML output, or reference appendix source for full detail.\end{quote}

\begin{quote}\small Dense table content is summarized in this PDF rendering. See the Markdown source and HTML book for the full table.\end{quote}


\chapter{Appendix A - Source Corpus Manifest}


{\small\raggedright SOURCE PATH: docs/\allowbreak{}archive/\allowbreak{}conversations/\allowbreak{}\_\allowbreak{}intake/\allowbreak{}CORPUS\_\allowbreak{}MANIFEST.md\par}


\section{Corpus Manifest}

This is the human-readable companion to \_intake/corpus\_manifest.json.


\subsection{Summary}

\begin{itemize}
\item conversation\_count: 45
\item source\_file\_count: 604
\item kind\_counts: \{'handoff': 34, 'json': 9, 'manifest': 69, 'markdown': 67, 'png': 10, 'primary\_report': 86, 'prompt': 31, 'py': 1, 'reader\_brief': 74, 'registers': 50, 'sha256': 1, 'source\_input': 15, 'spec\_sheet': 49, 'text': 5, 'tsv': 1, 'unknown': 10, 'verification': 62, 'zip': 30\}
\item complete\_packages: 45
\item partial\_packages: 0
\item unclear\_packages: 0
\end{itemize}

\subsection{Package Inventory}


\begin{quote}\small Dense table summarized for the reader edition. See the source file, HTML output, or reference appendix source for full detail.\end{quote}

\begin{quote}\small Dense table content is summarized in this PDF rendering. See the Markdown source and HTML book for the full table.\end{quote}



{\small\raggedright SOURCE PATH: docs/\allowbreak{}archive/\allowbreak{}conversations/\allowbreak{}\_\allowbreak{}intake/\allowbreak{}PACKAGE\_\allowbreak{}COMPLETENESS.md\par}


\section{Package Completeness}

Completeness is structural only. It does not mean a conversation claim is authoritative.


\subsection{Complete Packages}


\begin{quote}\small Dense table summarized for the reader edition. See the source file, HTML output, or reference appendix source for full detail.\end{quote}

\begin{quote}\small Dense table content is summarized in this PDF rendering. See the Markdown source and HTML book for the full table.\end{quote}


\subsection{Partial Packages}

None.


\subsection{Unclear Packages}

None.




{\small\raggedright SOURCE PATH: docs/\allowbreak{}archive/\allowbreak{}conversations/\allowbreak{}\_\allowbreak{}intake/\allowbreak{}SOURCE\_\allowbreak{}PROVENANCE.md\par}


\section{Source Provenance}

Every archived source file is treated as derived archival evidence with a source hash.


\subsection{advanced\_simulation\_infrastructure}

\begin{itemize}
\item Source class: conversation\_handoff\_export
\item Authority class: advisory\_only
\item Source hash: 6e865c26588b33219b96f5e75ad69c7f93ee3441229a47beb3bb054d3ccf7e07
\item Primary source: docs/archive/conversations/advanced\_simulation\_infrastructure/dominium\_advanced\_simulation\_infrastructure\_architecture\_\_human\_readable\_chat\_report.txt
\end{itemize}

\begin{quote}\small Dense table content is summarized in this PDF rendering. See the Markdown source and HTML book for the full table.\end{quote}

\subsection{app\_runtime\_platform\_renderers}

\begin{itemize}
\item Source class: conversation\_handoff\_export
\item Authority class: advisory\_only
\item Source hash: b54520262f4af10ecbbd05ee10c8b37ba8e72c88329a9fa0d65d8ad42555ed32
\item Primary source: docs/archive/conversations/app\_runtime\_platform\_renderers/dominium\_app0\_runtime\_platform\_renderers\_\_human\_readable\_chat\_report.txt
\end{itemize}

\begin{quote}\small Dense table content is summarized in this PDF rendering. See the Markdown source and HTML book for the full table.\end{quote}

\subsection{app\_testx\_codehygiene}

\begin{itemize}
\item Source class: conversation\_handoff\_export
\item Authority class: advisory\_only
\item Source hash: 021912beb80b75abd2d420364a14ff276747986c7edfa0aa3e75980819beccac
\item Primary source: docs/archive/conversations/app\_testx\_codehygiene/dominium\_app\_testx\_codehygiene\_handoff\_\_human\_readable\_chat\_report.txt
\end{itemize}

\begin{quote}\small Dense table content is summarized in this PDF rendering. See the Markdown source and HTML book for the full table.\end{quote}

\subsection{architecture\_codex\_prompts}

\begin{itemize}
\item Source class: conversation\_handoff\_export
\item Authority class: advisory\_only
\item Source hash: c6b29e26108d721ba5eac4e52ebade76ae70ccf4b36da54144b4b1887e77ec11
\item Primary source: docs/archive/conversations/architecture\_codex\_prompts/dominium\_domino\_architecture\_codex\_prompts\_\_human\_readable\_chat\_report.txt
\end{itemize}

\begin{quote}\small Dense table content is summarized in this PDF rendering. See the Markdown source and HTML book for the full table.\end{quote}

\subsection{architecture\_ui\_providers}

\begin{itemize}
\item Source class: conversation\_handoff\_export
\item Authority class: advisory\_only
\item Source hash: 841e5f4ea61b9cfac8743b4ef54da6a1bdbeb7577b2850c5966fe859b059cfe3
\item Primary source: docs/archive/conversations/architecture\_ui\_providers/dominium\_architecture\_ui\_providers\_robot\_os\_strategy\_\_01\_human\_readable\_report.md
\item Nested folders: prior\_package, source
\end{itemize}

\begin{quote}\small Dense table content is summarized in this PDF rendering. See the Markdown source and HTML book for the full table.\end{quote}

\subsection{Build\_and\_Future\_Proofing}

\begin{itemize}
\item Source class: conversation\_handoff\_export
\item Authority class: advisory\_only
\item Source hash: e72a4bb3063498001bed24be5075a2e2fcb803afdddf9f9543626e5aab12049d
\item Primary source: docs/archive/conversations/Build\_and\_Future\_Proofing/Dominium\_Build\_and\_Future\_Proofing\_Architecture\_\_01\_human\_readable\_report.md
\end{itemize}

\begin{quote}\small Dense table content is summarized in this PDF rendering. See the Markdown source and HTML book for the full table.\end{quote}

\subsection{canonical\_structure\_and\_framework}

\begin{itemize}
\item Source class: conversation\_handoff\_export
\item Authority class: advisory\_only
\item Source hash: 96b364a0b3d38b6f41b88983b68112f03b5648cee075f87af32d415013939027
\item Primary source: docs/archive/conversations/canonical\_structure\_and\_framework/dominium\_canonical\_architecture\_repo\_foundation\_\_01\_human\_readable\_report.md
\end{itemize}


\begin{quote}\small Dense table summarized for the reader edition. See the source file, HTML output, or reference appendix source for full detail.\end{quote}

\begin{quote}\small Dense table content is summarized in this PDF rendering. See the Markdown source and HTML book for the full table.\end{quote}


\subsection{Chronology\_Celestial\_Systems}

\begin{itemize}
\item Source class: conversation\_handoff\_export
\item Authority class: advisory\_only
\item Source hash: 5aef75e78769865c4d1d7ff8658cce34b17eb9e7bad8a5def1f7a13c3e11e733
\item Primary source: docs/archive/conversations/Chronology\_Celestial\_Systems/Dominium\_Chronology\_Celestial\_Systems\_\_human\_readable\_chat\_report.txt
\end{itemize}

\begin{quote}\small Dense table content is summarized in this PDF rendering. See the Markdown source and HTML book for the full table.\end{quote}

\subsection{development\_routes}

\begin{itemize}
\item Source class: conversation\_handoff\_export
\item Authority class: advisory\_only
\item Source hash: 5304f2ae84a9d1de77803ec9bccaba2d5461a082048f861d90f5a78789f96244
\item Primary source: docs/archive/conversations/development\_routes/dominium\_development\_routes\_\_human\_readable\_chat\_report.txt
\end{itemize}

\begin{quote}\small Dense table content is summarized in this PDF rendering. See the Markdown source and HTML book for the full table.\end{quote}

\subsection{documentation\_standards\_readme}

\begin{itemize}
\item Source class: conversation\_handoff\_export
\item Authority class: advisory\_only
\item Source hash: a295ae3a9f823c408c868b5d61a73347fde7cf951a9377a9e9396fc7cb30245e
\item Primary source: docs/archive/conversations/documentation\_standards\_readme/documentation\_standards\_readme\_handoff\_\_human\_readable\_chat\_report.txt
\end{itemize}

\begin{quote}\small Dense table content is summarized in this PDF rendering. See the Markdown source and HTML book for the full table.\end{quote}

\subsection{Dominium\_Architecture\_I}

\begin{itemize}
\item Source class: conversation\_handoff\_export
\item Authority class: advisory\_only
\item Source hash: 18f160ff207da734328b45ffc9627b73ce6be009bf212e304042fa96a00946ec
\item Primary source: docs/archive/conversations/Dominium\_Architecture\_I/Dominium\_Architecture\_I\_\_human\_readable\_chat\_report.txt
\end{itemize}

\begin{quote}\small Dense table content is summarized in this PDF rendering. See the Markdown source and HTML book for the full table.\end{quote}

\subsection{Dominium\_Architecture\_II}

\begin{itemize}
\item Source class: conversation\_handoff\_export
\item Authority class: advisory\_only
\item Source hash: 11fdeb43b407dde24a39ec79bfd3721d242a95059701a62a48efbcf98b20dba0
\item Primary source: docs/archive/conversations/Dominium\_Architecture\_II/Dominium\_Architecture\_II\_\_human\_readable\_chat\_report.txt
\end{itemize}

\begin{quote}\small Dense table content is summarized in this PDF rendering. See the Markdown source and HTML book for the full table.\end{quote}

\subsection{Dominium\_Architecture\_III}

\begin{itemize}
\item Source class: conversation\_handoff\_export
\item Authority class: advisory\_only
\item Source hash: 3d277cdcb4f6b8df6e72c71d03e8ecba3a321107137022867a083a590dcc8dbf
\item Primary source: docs/archive/conversations/Dominium\_Architecture\_III/Dominium\_Architecture\_III\_\_human\_readable\_chat\_report.txt
\end{itemize}

\begin{quote}\small Dense table content is summarized in this PDF rendering. See the Markdown source and HTML book for the full table.\end{quote}

\subsection{Dominium\_Architecture\_IV}

\begin{itemize}
\item Source class: conversation\_handoff\_export
\item Authority class: advisory\_only
\item Source hash: 0af55f972c719e405db3d24b54ecf8c7e57946df1c9dbd2804761f9259477b63
\item Primary source: docs/archive/conversations/Dominium\_Architecture\_IV/Dominium\_Architecture\_IV\_\_human\_readable\_chat\_report.txt
\end{itemize}

\begin{quote}\small Dense table content is summarized in this PDF rendering. See the Markdown source and HTML book for the full table.\end{quote}

\subsection{Dominium\_Complete\_Conversation}

\begin{itemize}
\item Source class: conversation\_handoff\_export
\item Authority class: advisory\_only
\item Source hash: 614f516f4bdf9810edd7093f00ff156208be39b50a0f111e2dd1dd533dd4846e
\item Primary source: docs/archive/conversations/Dominium\_Complete\_Conversation/Accompanying\_Report/Dominium\_Canon\_Portability\_Audit\_\_01\_human\_readable\_report.md
\item Nested folders: 00\_final\_manifest, Accompanying\_Report, deterministic\_solar\_system\_handoff\_file, dominium\_related\_file, existing\_companion\_report\_package, existing\_complete\_preservation\_package, existing\_consolidated\_package\_artifact, new\_companion\_report\_and\_bundle, robotic\_seed\_civilisation\_handoff\_file, uploaded\_preservation\_prompt
\end{itemize}


\begin{quote}\small Dense table summarized for the reader edition. See the source file, HTML output, or reference appendix source for full detail.\end{quote}

\begin{quote}\small Dense table content is summarized in this PDF rendering. See the Markdown source and HTML book for the full table.\end{quote}


\subsection{dominium\_domino\_codex\_planning}

\begin{itemize}
\item Source class: conversation\_handoff\_export
\item Authority class: advisory\_only
\item Source hash: f9c231e646cc554518e203e09f8ee4ce4cb9a793d45bb1fef4d4ed85bbf3fff8
\item Primary source: docs/archive/conversations/dominium\_domino\_codex\_planning/dominium\_domino\_codex\_planning\_\_human\_readable\_chat\_report.txt
\end{itemize}

\begin{quote}\small Dense table content is summarized in this PDF rendering. See the Markdown source and HTML book for the full table.\end{quote}

\subsection{dominium\_full\_conversation}

\begin{itemize}
\item Source class: conversation\_handoff\_export
\item Authority class: advisory\_only
\item Source hash: eafa81e8759980c506a0160413312fea02c41c9276e20d7bc1262a1878898c88
\item Primary source: docs/archive/conversations/dominium\_full\_conversation/companion\_report/dominium\_full\_conversation\_companion\_report\_\_01\_complete\_human\_report.md
\item Nested folders: companion\_report, previous\_preservation\_files, uploaded\_instruction, uploaded\_reference\_files, uploaded\_reference\_images
\end{itemize}


\begin{quote}\small Dense table summarized for the reader edition. See the source file, HTML output, or reference appendix source for full detail.\end{quote}

\begin{quote}\small Dense table content is summarized in this PDF rendering. See the Markdown source and HTML book for the full table.\end{quote}


\subsection{dominium\_setup}

\begin{itemize}
\item Source class: conversation\_handoff\_export
\item Authority class: advisory\_only
\item Source hash: 6d6cba108682ea0197f1da3ceedc79c54b17d34061d3431d0fd40fe5cde3fb2d
\item Primary source: docs/archive/conversations/dominium\_setup/dominium\_setup\_handoff\_\_human\_readable\_chat\_report.txt
\end{itemize}

\begin{quote}\small Dense table content is summarized in this PDF rendering. See the Markdown source and HTML book for the full table.\end{quote}

\subsection{Domino\_Dominium\_Workbench}

\begin{itemize}
\item Source class: conversation\_handoff\_export
\item Authority class: advisory\_only
\item Source hash: adb402d16f06b3d32add19c4603e15e14554b0bba87ed975f49be37218a3871e
\item Primary source: docs/archive/conversations/Domino\_Dominium\_Workbench/dominium\_workbench\_full\_conversation\_companion\_\_human\_readable\_detailed\_summary\_report.md
\item Nested folders: prior\_preservation\_files
\end{itemize}

\begin{quote}\small Dense table content is summarized in this PDF rendering. See the Markdown source and HTML book for the full table.\end{quote}

\subsection{domino\_engine\_refactor\_prompts}

\begin{itemize}
\item Source class: conversation\_handoff\_export
\item Authority class: advisory\_only
\item Source hash: 0d6bd3d358ad1f704b648c23553f87ba66cef082d1d2e6fe9921357f85f2b53f
\item Primary source: docs/archive/conversations/domino\_engine\_refactor\_prompts/dominium\_domino\_engine\_refactor\_prompts\_\_human\_readable\_chat\_report.txt
\end{itemize}

\begin{quote}\small Dense table content is summarized in this PDF rendering. See the Markdown source and HTML book for the full table.\end{quote}

\subsection{engine\_baseline\_architecture}

\begin{itemize}
\item Source class: conversation\_handoff\_export
\item Authority class: advisory\_only
\item Source hash: df5408421688f3079a08a99efc895408430d2f87a5a245e9a0ffcd7c38e931ef
\item Primary source: docs/archive/conversations/engine\_baseline\_architecture/domino\_dominium\_engine\_baseline\_architecture\_\_01\_human\_readable\_report.md
\end{itemize}

\begin{quote}\small Dense table content is summarized in this PDF rendering. See the Markdown source and HTML book for the full table.\end{quote}

\subsection{Foundation\_Workbench\_Codex}

\begin{itemize}
\item Source class: conversation\_handoff\_export
\item Authority class: advisory\_only
\item Source hash: 60bb02c56ea3be932c598b887a7a5410e59a224509f73e80d56bd3d0d9f20543
\item Primary source: docs/archive/conversations/Foundation\_Workbench\_Codex/Dominium\_Foundation\_Workbench\_Parallel\_Codex\_Preservation\_\_01\_human\_readable\_report.md
\end{itemize}

\begin{quote}\small Dense table content is summarized in this PDF rendering. See the Markdown source and HTML book for the full table.\end{quote}

\subsection{Framework\_Open\_Source\_Provider}

\begin{itemize}
\item Source class: conversation\_handoff\_export
\item Authority class: advisory\_only
\item Source hash: ddc8024a6a9d876c8f850a4c1e51efe91ac1e6c5a69c241dc71f3767eb43de0f
\item Primary source: docs/archive/conversations/Framework\_Open\_Source\_Provider/Domino\_Framework\_Open\_Source\_Provider\_Architecture\_\_01\_human\_readable\_report.md
\end{itemize}

\begin{quote}\small Dense table content is summarized in this PDF rendering. See the Markdown source and HTML book for the full table.\end{quote}

\subsection{gui\_binary\_content}

\begin{itemize}
\item Source class: conversation\_handoff\_export
\item Authority class: advisory\_only
\item Source hash: 92bf205a12f222598f86dc9157db73027848dc95cdafc6a52bc07d6ca832e3d6
\item Primary source: docs/archive/conversations/gui\_binary\_content/dominium\_gui\_binary\_content\_handoff\_\_human\_readable\_chat\_report.txt
\end{itemize}

\begin{quote}\small Dense table content is summarized in this PDF rendering. See the Markdown source and HTML book for the full table.\end{quote}

\subsection{Language\_Platform\_Architecture}

\begin{itemize}
\item Source class: conversation\_handoff\_export
\item Authority class: advisory\_only
\item Source hash: 2a3666f3a4c83c01e06a64de75735ace4b1d20acbc597fe29cf8d50791d10a10
\item Primary source: docs/archive/conversations/Language\_Platform\_Architecture/Dominium\_Language\_Platform\_Architecture\_Baseline\_\_01\_human\_readable\_report.md
\end{itemize}

\begin{quote}\small Dense table content is summarized in this PDF rendering. See the Markdown source and HTML book for the full table.\end{quote}

\subsection{launcher\_app\_layer}

\begin{itemize}
\item Source class: conversation\_handoff\_export
\item Authority class: advisory\_only
\item Source hash: 59d3feb4d0851df6e866622af2065624e1743a5fde0a5a6e440f68dd343bc196
\item Primary source: docs/archive/conversations/launcher\_app\_layer/dominium\_launcher\_app\_layer\_handoff\_\_human\_readable\_chat\_report.txt
\end{itemize}

\begin{quote}\small Dense table content is summarized in this PDF rendering. See the Markdown source and HTML book for the full table.\end{quote}

\subsection{Launcher\_Setup\_Architecture}

\begin{itemize}
\item Source class: conversation\_handoff\_export
\item Authority class: advisory\_only
\item Source hash: e8ea5d3d549b42a01e0a62fd953e40a4e8f92461bafc02242cbe2b1be5c26af7
\item Primary source: docs/archive/conversations/Launcher\_Setup\_Architecture/Dominium\_Launcher\_Setup\_Architecture\_\_human\_readable\_chat\_report.txt
\end{itemize}

\begin{quote}\small Dense table content is summarized in this PDF rendering. See the Markdown source and HTML book for the full table.\end{quote}

\subsection{Modularity\_AIDE\_Refactorability}

\begin{itemize}
\item Source class: conversation\_handoff\_export
\item Authority class: advisory\_only
\item Source hash: fcfdfc62f09babda08a1a20ba545d1402cab75c4885e30a66d06847e1ab75783
\item Primary source: docs/archive/conversations/Modularity\_AIDE\_Refactorability/Dominium\_Modularity\_AIDE\_Refactorability\_Architecture\_\_01\_human\_readable\_report.md
\item Nested folders: previous\_export, source\_uploads, supporting\_generation\_script
\end{itemize}

\begin{quote}\small Dense table content is summarized in this PDF rendering. See the Markdown source and HTML book for the full table.\end{quote}

\subsection{omega\_xi\_pi\_architecture\_future}

\begin{itemize}
\item Source class: conversation\_handoff\_export
\item Authority class: advisory\_only
\item Source hash: 115cf65085d59ae5f92ff3bfffa3988d8e7be1f221733eeeca76a66fa721bffd
\item Primary source: docs/archive/conversations/omega\_xi\_pi\_architecture\_future/dominium\_omega\_xi\_pi\_architecture\_future\_proofing\_planning\_\_01\_human\_readable\_report.md
\item Nested folders: source\_inputs
\end{itemize}

\begin{quote}\small Dense table content is summarized in this PDF rendering. See the Markdown source and HTML book for the full table.\end{quote}

\subsection{OS\_Interface\_Repo\_Architecture}

\begin{itemize}
\item Source class: conversation\_handoff\_export
\item Authority class: advisory\_only
\item Source hash: 75d980f7c6624379a51778945e93f1762fbcf0b7c86cbc4a4a11499ad7117917
\item Primary source: docs/archive/conversations/OS\_Interface\_Repo\_Architecture/Dominium\_OS\_Interface\_Repo\_Architecture\_\_01\_human\_readable\_report.md
\end{itemize}

\begin{quote}\small Dense table content is summarized in this PDF rendering. See the Markdown source and HTML book for the full table.\end{quote}

\subsection{platform\_renderer\_api\_plan}

\begin{itemize}
\item Source class: conversation\_handoff\_export
\item Authority class: advisory\_only
\item Source hash: 04b6b261f786cffc9a6f4555ba2b41bb329e8276a9efeb8ebd176d276b73a0e2
\item Primary source: docs/archive/conversations/platform\_renderer\_api\_plan/dominium\_codex\_platform\_renderer\_api\_plan\_\_human\_readable\_chat\_report.txt
\end{itemize}

\begin{quote}\small Dense table content is summarized in this PDF rendering. See the Markdown source and HTML book for the full table.\end{quote}

\subsection{platform\_support}

\begin{itemize}
\item Source class: conversation\_handoff\_export
\item Authority class: advisory\_only
\item Source hash: 6f737eb1b50d43f1b34b5c35ac8329e8156f2fc36ad280caf76c8b264b7cebca
\item Primary source: docs/archive/conversations/platform\_support/dominium\_platform\_support\_planning\_\_human\_readable\_chat\_report.txt
\end{itemize}

\begin{quote}\small Dense table content is summarized in this PDF rendering. See the Markdown source and HTML book for the full table.\end{quote}

\subsection{Portability\_Assurance\_Future\_Proof}

\begin{itemize}
\item Source class: conversation\_handoff\_export
\item Authority class: advisory\_only
\item Source hash: 0446febfae206c242d04a8389d51f9115fa0a27b1aa9413890023c4e0f73fc0f
\item Primary source: docs/archive/conversations/Portability\_Assurance\_Future\_Proof/Domino\_Dominium\_Portability\_Assurance\_Future\_Proof\_Architecture\_\_01\_human\_readable\_report.md
\item Nested folders: previous\_zip, source\_uploads
\end{itemize}

\begin{quote}\small Dense table content is summarized in this PDF rendering. See the Markdown source and HTML book for the full table.\end{quote}

\subsection{readme\_ports\_determinism}

\begin{itemize}
\item Source class: conversation\_handoff\_export
\item Authority class: advisory\_only
\item Source hash: 57cbc1d96fecde30399a5ccdcf0b12b63dd002c874e95336a8e99c13829d2879
\item Primary source: docs/archive/conversations/readme\_ports\_determinism/dominium\_readme\_ports\_determinism\_\_human\_readable\_chat\_report.txt
\end{itemize}

\begin{quote}\small Dense table content is summarized in this PDF rendering. See the Markdown source and HTML book for the full table.\end{quote}

\subsection{Refactor\_Architecture}

\begin{itemize}
\item Source class: conversation\_handoff\_export
\item Authority class: advisory\_only
\item Source hash: 9c10b26a476a59b2b9b74207892486f3bbbaf97b25e9ce3b7ddd54b141b8e7a8
\item Primary source: docs/archive/conversations/Refactor\_Architecture/Dominium\_Domino\_Refactor\_Architecture\_\_human\_readable\_chat\_report.txt
\end{itemize}

\begin{quote}\small Dense table content is summarized in this PDF rendering. See the Markdown source and HTML book for the full table.\end{quote}

\subsection{Refactor\_Control\_Plane}

\begin{itemize}
\item Source class: conversation\_handoff\_export
\item Authority class: advisory\_only
\item Source hash: 0d2ea4f3dffbcd7bd6c642d82c690c45b8012c9baeb5feee90d329b044a6a9e1
\item Primary source: docs/archive/conversations/Refactor\_Control\_Plane/AIDE\_XStack\_Dominium\_Refactor\_Control\_Plane\_\_01\_human\_readable\_report.md
\end{itemize}

\begin{quote}\small Dense table content is summarized in this PDF rendering. See the Markdown source and HTML book for the full table.\end{quote}

\subsection{Release\_Identity\_and\_Versioning}

\begin{itemize}
\item Source class: conversation\_handoff\_export
\item Authority class: advisory\_only
\item Source hash: 26742db25f39aa3f0073019899ce58d431bca26fad6c83d877f2e841b0a815de
\item Primary source: docs/archive/conversations/Release\_Identity\_and\_Versioning/Dominium\_XStack\_Release\_Identity\_and\_Versioning\_\_01\_human\_readable\_report.md
\end{itemize}

\begin{quote}\small Dense table content is summarized in this PDF rendering. See the Markdown source and HTML book for the full table.\end{quote}

\subsection{testx\_repox\_governance}

\begin{itemize}
\item Source class: conversation\_handoff\_export
\item Authority class: advisory\_only
\item Source hash: bf963688aa6f2aebbd4351cc244754d3548d5dbc5ef57c36ca2e4bdb3a0bfe3d
\item Primary source: docs/archive/conversations/testx\_repox\_governance/dominium\_testx\_repox\_governance\_handoff\_\_human\_readable\_chat\_report.txt
\end{itemize}

\begin{quote}\small Dense table content is summarized in this PDF rendering. See the Markdown source and HTML book for the full table.\end{quote}

\subsection{Timekeeping\_and\_2038\_Resilience}

\begin{itemize}
\item Source class: conversation\_handoff\_export
\item Authority class: advisory\_only
\item Source hash: a2cc19e7b373b11730ad81cd26c55d487eaffad1540f0aa5d60011e13165eda7
\item Primary source: docs/archive/conversations/Timekeeping\_and\_2038\_Resilience/Dominium\_Timekeeping\_and\_2038\_Resilience\_\_01\_human\_readable\_report.md
\end{itemize}

\begin{quote}\small Dense table content is summarized in this PDF rendering. See the Markdown source and HTML book for the full table.\end{quote}

\subsection{UE6\_Domino\_Deterministic\_Universe}

\begin{itemize}
\item Source class: conversation\_handoff\_export
\item Authority class: advisory\_only
\item Source hash: d0ca1855e369ad01a733afc9e86f5343d85b3667c5fe38cde334eb4ff2c7d764
\item Primary source: docs/archive/conversations/UE6\_Domino\_Deterministic\_Universe/Dominium\_UE6\_Domino\_Deterministic\_Universe\_\_01\_human\_readable\_report.md
\end{itemize}

\begin{quote}\small Dense table content is summarized in this PDF rendering. See the Markdown source and HTML book for the full table.\end{quote}

\subsection{ui\_editor\_tool\_editor\_planning}

\begin{itemize}
\item Source class: conversation\_handoff\_export
\item Authority class: advisory\_only
\item Source hash: f3079a92a9416ddf98c26f2c1dd187cf706688b2382c3f22dce14d204ec963f2
\item Primary source: docs/archive/conversations/ui\_editor\_tool\_editor\_planning/dominium\_ui\_editor\_tool\_editor\_planning\_\_human\_readable\_chat\_report.txt
\end{itemize}

\begin{quote}\small Dense table content is summarized in this PDF rendering. See the Markdown source and HTML book for the full table.\end{quote}

\subsection{Universe\_Explorer\_Planning}

\begin{itemize}
\item Source class: conversation\_handoff\_export
\item Authority class: advisory\_only
\item Source hash: 326be1eb57b06c5042b6180a4eec6ce5edcdb3b8f7a819fd11f4140612a67781
\item Primary source: docs/archive/conversations/Universe\_Explorer\_Planning/Dominium\_Universe\_Explorer\_Planning\_Handoff\_\_01\_human\_readable\_report.md
\end{itemize}

\begin{quote}\small Dense table content is summarized in this PDF rendering. See the Markdown source and HTML book for the full table.\end{quote}

\subsection{Workbench\_AIDE\_Product\_Spine}

\begin{itemize}
\item Source class: conversation\_handoff\_export
\item Authority class: advisory\_only
\item Source hash: 738e2a76e12a500897ceb752001a7a285a1e607eca680cfc206fd26d8cca02f0
\item Primary source: docs/archive/conversations/Workbench\_AIDE\_Product\_Spine/Dominium\_Architecture\_Workbench\_AIDE\_Product\_Spine\_Planning\_\_01\_human\_readable\_report.md
\end{itemize}

\begin{quote}\small Dense table content is summarized in this PDF rendering. See the Markdown source and HTML book for the full table.\end{quote}

\subsection{World\_Architecture}

\begin{itemize}
\item Source class: conversation\_handoff\_export
\item Authority class: advisory\_only
\item Source hash: 058c77eeb1e89e9ee3b8add92f5feb8eec2993165bc789322000cf3e6e5290f9
\item Primary source: docs/archive/conversations/World\_Architecture/Dominium\_World\_Architecture\_\_human\_readable\_chat\_report.txt
\end{itemize}

\begin{quote}\small Dense table content is summarized in this PDF rendering. See the Markdown source and HTML book for the full table.\end{quote}

\subsection{xstack\_lab\_galaxy}

\begin{itemize}
\item Source class: conversation\_handoff\_export
\item Authority class: advisory\_only
\item Source hash: 33f31158082394d59c723ea75439a781425b1a8800b6ea2e1bc8789a503b0831
\item Primary source: docs/archive/conversations/xstack\_lab\_galaxy/dominium\_xstack\_lab\_galaxy\_handoff\_\_human\_readable\_chat\_report.txt
\end{itemize}

\begin{quote}\small Dense table content is summarized in this PDF rendering. See the Markdown source and HTML book for the full table.\end{quote}


\chapter{Appendix B - Full Claim Review Matrix}


{\small\raggedright SOURCE PATH: docs/\allowbreak{}archive/\allowbreak{}conversations/\allowbreak{}\_\allowbreak{}reconciliation/\allowbreak{}CLAIM\_\allowbreak{}REVIEW\_\allowbreak{}MATRIX\_\allowbreak{}v0.md\par}


\section{Claim Review Matrix v0}

Each row is a raw promotion candidate classified for review. None are promoted.


\begin{quote}\small Dense table content is summarized in this PDF rendering. See the Markdown source and HTML book for the full table.\end{quote}


\chapter{Appendix C - Full Promotion Register}


{\small\raggedright SOURCE PATH: docs/\allowbreak{}archive/\allowbreak{}conversations/\allowbreak{}\_\allowbreak{}promotion/\allowbreak{}PROMOTION\_\allowbreak{}REVIEW\_\allowbreak{}BOARD\_\allowbreak{}v0.md\par}


\section{Promotion Review Board v0}

This board converts raw promotion candidates into bounded review items. It does not patch or promote live docs.


\subsection{Counts}

\begin{itemize}
\item Total candidates: 135
\item architecture\_doc\_candidate: 16
\item archive\_only: 9
\item blocked\_by\_queue: 10
\item contract\_candidate: 11
\item docs\_clarification\_candidate: 1
\item insufficient\_evidence: 53
\item needs\_user\_decision: 12
\item planning\_doc\_candidate: 1
\item reject\_as\_noise: 19
\item schema\_candidate: 3
\end{itemize}

\subsection{Review Items}

\subsubsection{PROMOTE-0001 - needs\_user\_decision}

\begin{itemize}
\item Source conversation: advanced\_simulation\_infrastructure
\item Proposed target path: docs/archive/conversations/\_decision/DECISION\_DOCKET\_v0.md
\item Current authority support: not yet verified against target docs
\item Current authority conflict: none detected by triage
\item Queue risk: none\_detected
\item Validation required: preserve as archive evidence; no live target validation
\item Recommended next action: hold\_or\_review
\item Claim: The user then asked whether arbitrary placement was possible or whether the system was stuck to grids. The answer established that the engine should be grid-agnostic. [FACT] The user's question directly made arbitrary placement a major topic. [INFERENCE] The answer was accepted as the current design direction because the user then asked for a Codex prompt plan to implement changes necessary to achieve it.
\end{itemize}

\subsubsection{PROMOTE-0002 - reject\_as\_noise}

\begin{itemize}
\item Source conversation: advanced\_simulation\_infrastructure
\item Proposed target path: none
\item Current authority support: not yet verified against target docs
\item Current authority conflict: none detected by triage
\item Queue risk: none\_detected
\item Validation required: preserve as archive evidence; no live target validation
\item Recommended next action: hold\_or\_review
\item Claim: Finally, the chat shifted into archival mode: it produced a maximum-fidelity context transfer packet, then a downloadable report package, then an in-chat reader version. Those later outputs preserved the discussion for future aggregation.
\end{itemize}

\subsubsection{PROMOTE-0003 - architecture\_doc\_candidate}

\begin{itemize}
\item Source conversation: advanced\_simulation\_infrastructure
\item Proposed target path: docs/architecture/** after target selection
\item Current authority support: not yet verified against target docs
\item Current authority conflict: none detected by triage
\item Queue risk: none\_detected
\item Validation required: source check, authority-order check, link check, generated output validator, FAST
\item Recommended next action: create\_microtask
\item Claim: [INFERENCE] The conclusion was not a finalized implementation but a working architecture: each simulation domain should be deterministic, fixed-point, content-driven, versioned, replayable, and integrated through clear read/write boundaries.
\end{itemize}

\subsubsection{PROMOTE-0004 - architecture\_doc\_candidate}

\begin{itemize}
\item Source conversation: app\_runtime\_platform\_renderers
\item Proposed target path: docs/architecture/** after target selection
\item Current authority support: not yet verified against target docs
\item Current authority conflict: none detected by triage
\item Queue risk: none\_detected
\item Validation required: source check, authority-order check, link check, generated output validator, FAST
\item Recommended next action: create\_microtask
\item Claim: The central topic was APP0: the application/runtime/platform/renderer layer of Dominium. This topic came up because the user provided a formal prompt defining what APP0 should do.
\end{itemize}

\subsubsection{PROMOTE-0005 - blocked\_by\_queue}

\begin{itemize}
\item Source conversation: app\_runtime\_platform\_renderers
\item Proposed target path: none
\item Current authority support: not yet verified against target docs
\item Current authority conflict: blocked queue scope
\item Queue risk: blocked
\item Validation required: preserve as archive evidence; no live target validation
\item Recommended next action: hold\_or\_review
\item Claim: The most important conclusion was that APP0 must remain outside the core simulation rules. **FACT:** the user explicitly forbade redesigning simulation rules, altering content definitions, changing life/civilization/economy logic, and introducing gameplay shortcuts. **FACT:** authoritative logic remains in engine + game.
\end{itemize}

\subsubsection{PROMOTE-0006 - needs\_user\_decision}

\begin{itemize}
\item Source conversation: app\_runtime\_platform\_renderers
\item Proposed target path: docs/archive/conversations/\_decision/DECISION\_DOCKET\_v0.md
\item Current authority support: not yet verified against target docs
\item Current authority conflict: none detected by triage
\item Queue risk: none\_detected
\item Validation required: preserve as archive evidence; no live target validation
\item Recommended next action: hold\_or\_review
\item Claim: Future work must define the client's allowed dependency surface and whether it can access any simulation-adjacent prediction layer.
\end{itemize}

\subsubsection{PROMOTE-0007 - schema\_candidate}

\begin{itemize}
\item Source conversation: app\_testx\_codehygiene
\item Proposed target path: schema/** or schemas/** after authority review
\item Current authority support: not yet verified against target docs
\item Current authority conflict: none detected by triage
\item Queue risk: none\_detected
\item Validation required: authority review plus targeted schema/contract parser and FAST
\item Recommended next action: hold\_or\_review
\item Claim: A large set of prompts was then generated to lock down architecture, determinism, performance, schema governance, rendering, epistemic UI, sharding, interest sets, and fidelity projection. This became the Phase 1 hardening layer. Additional audit prompts were introduced to ensure consistency before proceeding into life, civilization, war, content, agents, tools, mods, and final long-term policy.
\end{itemize}

\subsubsection{PROMOTE-0008 - reject\_as\_noise}

\begin{itemize}
\item Source conversation: app\_testx\_codehygiene
\item Proposed target path: none
\item Current authority support: not yet verified against target docs
\item Current authority conflict: none detected by triage
\item Queue risk: none\_detected
\item Validation required: preserve as archive evidence; no live target validation
\item Recommended next action: hold\_or\_review
\item Claim: Because the chat was huge, the user asked for a maximum-fidelity context transfer packet. Then they asked for downloadable report files. Then they asked for an in-chat reader. Finally, they asked for this human-readable explanatory report. This final report is meant to let a future assistant or human understand the substance without re-reading the whole conversation.
\end{itemize}

\subsubsection{PROMOTE-0009 - contract\_candidate}

\begin{itemize}
\item Source conversation: app\_testx\_codehygiene
\item Proposed target path: contracts/** or docs/reference/contracts/** after authority review
\item Current authority support: not yet verified against target docs
\item Current authority conflict: none detected by triage
\item Queue risk: none\_detected
\item Validation required: authority review plus targeted schema/contract parser and FAST
\item Recommended next action: hold\_or\_review
\item Claim: The central project is a deterministic universe engine/game. The user's ambition is not just a normal game but a world runtime that can represent everything from a single-room scenario to an AI-only universe to an MMO. The project is meant to support real-world defaults, arbitrary modding, procedural and defined content, strong epistemics, and long-term replayability.
\end{itemize}

\subsubsection{PROMOTE-0010 - insufficient\_evidence}

\begin{itemize}
\item Source conversation: architecture\_codex\_prompts
\item Proposed target path: none
\item Current authority support: not yet verified against target docs
\item Current authority conflict: none detected by triage
\item Queue risk: none\_detected
\item Validation required: preserve as archive evidence; no live target validation
\item Recommended next action: hold\_or\_review
\item Claim: The most important thing to remember is this: **this chat is the architectural and implementation-roadmap backbone for Dominium/Domino, but it is not an implementation log**. It should be preserved as a design/specification source and as a prompt library, while actual code and current facts must be verified separately.
\end{itemize}

\subsubsection{PROMOTE-0011 - insufficient\_evidence}

\begin{itemize}
\item Source conversation: architecture\_codex\_prompts
\item Proposed target path: none
\item Current authority support: not yet verified against target docs
\item Current authority conflict: none detected by triage
\item Queue risk: none\_detected
\item Validation required: preserve as archive evidence; no live target validation
\item Recommended next action: hold\_or\_review
\item Claim: FACT: The user opened by pasting an "Extended Master Starter Prompt - Dominium + Domino." That prompt established the role of the assistant as a senior engine architect and defined the main project philosophy. Domino was described as a deterministic, fixed-point, C89 engine portable across old and modern OS/architecture combinations. Dominium was described as the product suite built on top of Domino.
\end{itemize}

\subsubsection{PROMOTE-0012 - archive\_only}

\begin{itemize}
\item Source conversation: architecture\_codex\_prompts
\item Proposed target path: docs/archive/conversations/**
\item Current authority support: not yet verified against target docs
\item Current authority conflict: none detected by triage
\item Queue risk: none\_detected
\item Validation required: preserve as archive evidence; no live target validation
\item Recommended next action: hold\_or\_review
\item Claim: FACT: The user asked, succinctly, how a player would build a complete machine workshop. The assistant described a progression: extract raw materials, process them into industrial materials, produce mechanical components, produce electrical components, establish power distribution, establish logistics, assemble machine frames and machines, add measurement/control tooling, then scale horizontally and vertically.
\end{itemize}

\subsubsection{PROMOTE-0013 - insufficient\_evidence}

\begin{itemize}
\item Source conversation: architecture\_ui\_providers
\item Proposed target path: none
\item Current authority support: not yet verified against target docs
\item Current authority conflict: none detected by triage
\item Queue risk: none\_detected
\item Validation required: preserve as archive evidence; no live target validation
\item Recommended next action: hold\_or\_review
\item Claim: The language baseline decision matters because it changes what old renderers and toolchains matter. C17/C++17 reduces the need to design around C89/C++98 constraints but does not remove the need for C-compatible ABI boundaries. The system floor decision similarly moves many older renderers and OS targets into research/back-port categories.
\end{itemize}

\subsubsection{PROMOTE-0014 - insufficient\_evidence}

\begin{itemize}
\item Source conversation: architecture\_ui\_providers
\item Proposed target path: none
\item Current authority support: not yet verified against target docs
\item Current authority conflict: none detected by triage
\item Queue risk: none\_detected
\item Validation required: preserve as archive evidence; no live target validation
\item Recommended next action: hold\_or\_review
\item Claim: The Robot OS and robotic seed-civilisation decisions are the most important design decisions. They are explicit user statements, not merely assistant proposals. They should be formalized as product/game requirements. They affect UI, fog-of-war, spawning, respawn, automation, industry, tutorial/planner design, and MMO anti-cheat.
\end{itemize}

\subsubsection{PROMOTE-0015 - blocked\_by\_queue}

\begin{itemize}
\item Source conversation: architecture\_ui\_providers
\item Proposed target path: none
\item Current authority support: not yet verified against target docs
\item Current authority conflict: blocked queue scope
\item Queue risk: blocked
\item Validation required: preserve as archive evidence; no live target validation
\item Recommended next action: hold\_or\_review
\item Claim: The central tradeoff was speed versus purity. Raylib, SDL2, Lua, raygui, rlgl, rlsw, and related libraries can accelerate visible progress, but if they leak into game law, UI documents, saves, replays, or public contracts, they become architectural capture. The chosen compromise is to use them as providers behind Dominium-owned services and conformance tests.
\end{itemize}

\subsubsection{PROMOTE-0016 - archive\_only}

\begin{itemize}
\item Source conversation: Build\_and\_Future\_Proofing
\item Proposed target path: docs/archive/conversations/**
\item Current authority support: not yet verified against target docs
\item Current authority conflict: none detected by triage
\item Queue risk: none\_detected
\item Validation required: preserve as archive evidence; no live target validation
\item Recommended next action: hold\_or\_review
\item Claim: Plain-language limitation: this package is reliable as a preservation of the current visible chat. It should not be treated as a complete archive of every Dominium discussion in other chats. Where the report uses project or repository context, it is labelled as such. The most important caveat is that many items are recommendations, not final user decisions.
\end{itemize}

\subsubsection{PROMOTE-0017 - reject\_as\_noise}

\begin{itemize}
\item Source conversation: Build\_and\_Future\_Proofing
\item Proposed target path: none
\item Current authority support: not yet verified against target docs
\item Current authority conflict: none detected by triage
\item Queue risk: none\_detected
\item Validation required: preserve as archive evidence; no live target validation
\item Recommended next action: hold\_or\_review
\item Claim: The final uploaded prompt requested a complete preservation package for this chat: a human-readable explanation, structured registers, context transfer packet, spec sheet, aggregator packet, self-audit, and files/ZIP package. This response completes that export task and creates downloadable files for later reading and aggregation.
\end{itemize}

\subsubsection{PROMOTE-0018 - archive\_only}

\begin{itemize}
\item Source conversation: Build\_and\_Future\_Proofing
\item Proposed target path: docs/archive/conversations/**
\item Current authority support: not yet verified against target docs
\item Current authority conflict: none detected by triage
\item Queue risk: none\_detected
\item Validation required: preserve as archive evidence; no live target validation
\item Recommended next action: hold\_or\_review
\item Claim: The final explicit goal was preservation: create a maximum-fidelity package for this chat so the user can understand it later, ask questions in this chat, merge it with other old-chat reports, and eventually build a master project spec book.
\end{itemize}

\subsubsection{PROMOTE-0019 - insufficient\_evidence}

\begin{itemize}
\item Source conversation: canonical\_structure\_and\_framework
\item Proposed target path: none
\item Current authority support: not yet verified against target docs
\item Current authority conflict: none detected by triage
\item Queue risk: none\_detected
\item Validation required: preserve as archive evidence; no live target validation
\item Recommended next action: hold\_or\_review
\item Claim: This produced a pattern: generate a large Codex/AIDE task, the user ran it or reported a result, then the assistant evaluated what was truly fixed versus what was only validator/document churn. The user eventually demanded a one-shot "actual final cleanup" prompt because previous passes had sometimes added validators without moving directories. That prompt explicitly required real git mv routing, not just reports.
\end{itemize}

\subsubsection{PROMOTE-0020 - blocked\_by\_queue}

\begin{itemize}
\item Source conversation: canonical\_structure\_and\_framework
\item Proposed target path: none
\item Current authority support: not yet verified against target docs
\item Current authority conflict: blocked queue scope
\item Queue risk: blocked
\item Validation required: preserve as archive evidence; no live target validation
\item Recommended next action: hold\_or\_review
\item Claim: The final provider directory doctrine became service-first: runtime/<service>/providers/<provider>/, not runtime/<vendor>/<service>/. Provider choices belong in release/profiles/ or content/profiles/, not app path names. External code belongs under external/upstream/ or external/vendor/, but the repo should choose one convention.
\end{itemize}

\subsubsection{PROMOTE-0021 - schema\_candidate}

\begin{itemize}
\item Source conversation: canonical\_structure\_and\_framework
\item Proposed target path: schema/** or schemas/** after authority review
\item Current authority support: not yet verified against target docs
\item Current authority conflict: none detected by triage
\item Queue risk: none\_detected
\item Validation required: authority review plus targeted schema/contract parser and FAST
\item Recommended next action: hold\_or\_review
\item Claim: Packs are distributable authored payloads. The preferred pack layout became category-based, such as content/packs/<category>/<pack\_id>/. Pack IDs must not depend on paths. Saves and replays must use stable IDs, schema versions, canonical serialization, deterministic hashes, and explicit migration/refusal policies. Backward compatibility must be a compatibility corpus plus tests, not intention.
\end{itemize}

\subsubsection{PROMOTE-0022 - insufficient\_evidence}

\begin{itemize}
\item Source conversation: Chronology\_Celestial\_Systems
\item Proposed target path: none
\item Current authority support: not yet verified against target docs
\item Current authority conflict: none detected by triage
\item Queue risk: none\_detected
\item Validation required: preserve as archive evidence; no live target validation
\item Recommended next action: hold\_or\_review
\item Claim: This reinforced the core celestial-content rule: locations should exist explicitly when they change player decisions.
\end{itemize}

\subsubsection{PROMOTE-0023 - insufficient\_evidence}

\begin{itemize}
\item Source conversation: Chronology\_Celestial\_Systems
\item Proposed target path: none
\item Current authority support: not yet verified against target docs
\item Current authority conflict: none detected by triage
\item Queue risk: none\_detected
\item Validation required: preserve as archive evidence; no live target validation
\item Recommended next action: hold\_or\_review
\item Claim: This became the spatial foundation for later time/calendar decisions because every major body or region could potentially have its own local time, calendar, or operational standard.
\end{itemize}

\subsubsection{PROMOTE-0024 - needs\_user\_decision}

\begin{itemize}
\item Source conversation: Chronology\_Celestial\_Systems
\item Proposed target path: docs/archive/conversations/\_decision/DECISION\_DOCKET\_v0.md
\item Current authority support: not yet verified against target docs
\item Current authority conflict: none detected by triage
\item Queue risk: none\_detected
\item Validation required: preserve as archive evidence; no live target validation
\item Recommended next action: hold\_or\_review
\item Claim: The assistant formalized this as the Perfect Earth Calendar, later called HPC-E. The final structure is clear, though the exact leap rule remains unresolved.
\end{itemize}

\subsubsection{PROMOTE-0025 - needs\_user\_decision}

\begin{itemize}
\item Source conversation: development\_routes
\item Proposed target path: docs/archive/conversations/\_decision/DECISION\_DOCKET\_v0.md
\item Current authority support: not yet verified against target docs
\item Current authority conflict: none detected by triage
\item Queue risk: none\_detected
\item Validation required: preserve as archive evidence; no live target validation
\item Recommended next action: hold\_or\_review
\item Claim: Before relying on this chat as project authority, a future assistant or human reviewer must confirm whether Dominium is actually intended to be simulation-heavy, whether strict determinism matters, whether replay or multiplayer are goals, whether modding matters, and what the actual game loop is. Until then, the kernel-first plan should be treated as a strong provisional proposal, not a final specification.
\end{itemize}

\subsubsection{PROMOTE-0026 - insufficient\_evidence}

\begin{itemize}
\item Source conversation: development\_routes
\item Proposed target path: none
\item Current authority support: not yet verified against target docs
\item Current authority conflict: none detected by triage
\item Queue risk: none\_detected
\item Validation required: preserve as archive evidence; no live target validation
\item Recommended next action: hold\_or\_review
\item Claim: The first answer was assertive. It stated that Dominium must follow Route C and described the route as the only viable one for Dominium. Later parts of the chat corrected the status of that claim. FACT: the assistant made that recommendation. UNCERTAIN / UNVERIFIED: the recommendation was not validated with external sources in the chat and was not explicitly accepted by the user.
\end{itemize}

\subsubsection{PROMOTE-0027 - archive\_only}

\begin{itemize}
\item Source conversation: development\_routes
\item Proposed target path: docs/archive/conversations/**
\item Current authority support: not yet verified against target docs
\item Current authority conflict: none detected by triage
\item Queue risk: none\_detected
\item Validation required: preserve as archive evidence; no live target validation
\item Recommended next action: hold\_or\_review
\item Claim: The initial technical answer proposed a layered stack. It started with mathematical and temporal primitives, then moved into the deterministic simulation kernel, then world state, systems, persistence and replay, tooling, presentation, content, and finally distribution, modding, and multiplayer.
\end{itemize}

\subsubsection{PROMOTE-0028 - needs\_user\_decision}

\begin{itemize}
\item Source conversation: documentation\_standards\_readme
\item Proposed target path: docs/archive/conversations/\_decision/DECISION\_DOCKET\_v0.md
\item Current authority support: not yet verified against target docs
\item Current authority conflict: none detected by triage
\item Queue risk: none\_detected
\item Validation required: preserve as archive evidence; no live target validation
\item Recommended next action: hold\_or\_review
\item Claim: The key thing to remember: this chat produced the standards and prompts for future work; it did not verify or change the project. The next assistant must inspect the actual repository before writing repo-specific docs, README content, platform claims, license claims, or subsystem descriptions.
\end{itemize}

\subsubsection{PROMOTE-0029 - contract\_candidate}

\begin{itemize}
\item Source conversation: documentation\_standards\_readme
\item Proposed target path: contracts/** or docs/reference/contracts/** after authority review
\item Current authority support: not yet verified against target docs
\item Current authority conflict: none detected by triage
\item Queue risk: none\_detected
\item Validation required: authority review plus targeted schema/contract parser and FAST
\item Recommended next action: hold\_or\_review
\item Claim: FACT: The chat established that public API contracts should live in headers.
\end{itemize}

\subsubsection{PROMOTE-0030 - insufficient\_evidence}

\begin{itemize}
\item Source conversation: documentation\_standards\_readme
\item Proposed target path: none
\item Current authority support: not yet verified against target docs
\item Current authority conflict: none detected by triage
\item Queue risk: none\_detected
\item Validation required: preserve as archive evidence; no live target validation
\item Recommended next action: hold\_or\_review
\item Claim: FACT: Source files should document implementation rationale, non-obvious invariants, and hazards.
\end{itemize}

\subsubsection{PROMOTE-0031 - reject\_as\_noise}

\begin{itemize}
\item Source conversation: Dominium\_Architecture\_I
\item Proposed target path: none
\item Current authority support: not yet verified against target docs
\item Current authority conflict: none detected by triage
\item Queue risk: none\_detected
\item Validation required: preserve as archive evidence; no live target validation
\item Recommended next action: hold\_or\_review
\item Claim: During this stage, the work was still broadly about "the project spec." The user wanted a comprehensive description of the game and its technical systems. The assistant produced large design documents, but many of those earlier contents are not visible in the retained transcript. This is important because the final report package marks those earlier outputs as referenced but not fully recovered.
\end{itemize}

\subsubsection{PROMOTE-0032 - architecture\_doc\_candidate}

\begin{itemize}
\item Source conversation: Dominium\_Architecture\_I
\item Proposed target path: docs/architecture/** after target selection
\item Current authority support: not yet verified against target docs
\item Current authority conflict: none detected by triage
\item Queue risk: none\_detected
\item Validation required: source check, authority-order check, link check, generated output validator, FAST
\item Recommended next action: create\_microtask
\item Claim: The user also made a clear versioning decision. **FACT:** The user wanted **major.minor.patch semantic versioning** for all components/packages and for the complete game package. The user explicitly did **not** want build numbers or build dates in filenames. Instead, those details should live in metadata. This became part of the future build/package design.
\end{itemize}

\subsubsection{PROMOTE-0033 - insufficient\_evidence}

\begin{itemize}
\item Source conversation: Dominium\_Architecture\_I
\item Proposed target path: none
\item Current authority support: not yet verified against target docs
\item Current authority conflict: none detected by triage
\item Queue risk: none\_detected
\item Validation required: preserve as archive evidence; no live target validation
\item Recommended next action: hold\_or\_review
\item Claim: The key change came when the user decided that Codex could not read the entire conversation at once. **FACT:** The user said they would copy and paste the transcript with timestamps and would not edit it manually, but because Codex could not read the whole context at once, they would not bother asking Codex to compile a full book directly from the whole chat.
\end{itemize}

\subsubsection{PROMOTE-0034 - reject\_as\_noise}

\begin{itemize}
\item Source conversation: Dominium\_Architecture\_II
\item Proposed target path: none
\item Current authority support: not yet verified against target docs
\item Current authority conflict: none detected by triage
\item Queue risk: none\_detected
\item Validation required: preserve as archive evidence; no live target validation
\item Recommended next action: hold\_or\_review
\item Claim: The last part of the chat was handoff extraction. The user asked for discovery inventory, structured registers, a context transfer packet, and then a downloadable package. The generated package captured the chat's workstreams, decisions, tasks, constraints, open questions, artifacts, risks, and verification items. This current report is a plain-language explanation of that substance.
\end{itemize}

\subsubsection{PROMOTE-0035 - architecture\_doc\_candidate}

\begin{itemize}
\item Source conversation: Dominium\_Architecture\_II
\item Proposed target path: docs/architecture/** after target selection
\item Current authority support: not yet verified against target docs
\item Current authority conflict: none detected by triage
\item Queue risk: none\_detected
\item Validation required: source check, authority-order check, link check, generated output validator, FAST
\item Recommended next action: create\_microtask
\item Claim: The entire conversation is anchored by the requirement that Dominium must run deterministically. This means the same initial state and same inputs must produce the same state across machines, compilers, operating systems, and eventually retro and modern platforms.
\end{itemize}

\subsubsection{PROMOTE-0036 - insufficient\_evidence}

\begin{itemize}
\item Source conversation: Dominium\_Architecture\_II
\item Proposed target path: none
\item Current authority support: not yet verified against target docs
\item Current authority conflict: none detected by triage
\item Queue risk: none\_detected
\item Validation required: preserve as archive evidence; no live target validation
\item Recommended next action: hold\_or\_review
\item Claim: The chat discussed fixed tick phases, virtual lanes, merge order, command buffers, event queues, and deterministic serialization. The simulation must not depend on render FPS, wall-clock time, OS scheduler behaviour, GPU state, or floating-point quirks. This explains many later decisions: C89 core, integer/fixed-point math, no sim floats, no platform headers in simulation, and renderer as pure observer.
\end{itemize}

\subsubsection{PROMOTE-0037 - insufficient\_evidence}

\begin{itemize}
\item Source conversation: Dominium\_Architecture\_III
\item Proposed target path: none
\item Current authority support: not yet verified against target docs
\item Current authority conflict: none detected by triage
\item Queue risk: none\_detected
\item Validation required: preserve as archive evidence; no live target validation
\item Recommended next action: hold\_or\_review
\item Claim: The view system then expanded beyond the launcher. **FACT:** The user wants instant zoom to any scale, instant switching between top-down 2D and first-person 3D, and arbitrary cameras for free cam, map viewing, content creation, HUD cameras, CCTV, overlays, windows, and offscreen targets. These are client/render-side concerns, not simulation concerns.
\end{itemize}

\subsubsection{PROMOTE-0038 - architecture\_doc\_candidate}

\begin{itemize}
\item Source conversation: Dominium\_Architecture\_III
\item Proposed target path: docs/architecture/** after target selection
\item Current authority support: not yet verified against target docs
\item Current authority conflict: none detected by triage
\item Queue risk: none\_detected
\item Validation required: source check, authority-order check, link check, generated output validator, FAST
\item Recommended next action: create\_microtask
\item Claim: This context matters because the rest of the conversation happened inside those constraints. The launcher could not become a dumping ground for OS-specific behavior, sim mutation, renderer decisions, or nondeterministic game-state changes. It had to fit the project's deterministic layering.
\end{itemize}

\subsubsection{PROMOTE-0039 - insufficient\_evidence}

\begin{itemize}
\item Source conversation: Dominium\_Architecture\_III
\item Proposed target path: none
\item Current authority support: not yet verified against target docs
\item Current authority conflict: none detected by triage
\item Queue risk: none\_detected
\item Validation required: preserve as archive evidence; no live target validation
\item Recommended next action: hold\_or\_review
\item Claim: The user then uploaded dominium.7z, saying it was the git repository as of that moment. **FACT:** The prior assistant said it could not open .7z archives in that environment. That means the actual repository state remained unverified at that stage. This became important later, because many implementation ideas were discussed without confirmed access to the repo contents.
\end{itemize}

\subsubsection{PROMOTE-0040 - insufficient\_evidence}

\begin{itemize}
\item Source conversation: Dominium\_Architecture\_IV
\item Proposed target path: none
\item Current authority support: not yet verified against target docs
\item Current authority conflict: none detected by triage
\item Queue risk: none\_detected
\item Validation required: preserve as archive evidence; no live target validation
\item Recommended next action: hold\_or\_review
\item Claim: The reusable engine layer was named **Domino**. This became one of the most important naming and architectural decisions in the chat. Domino is the engine/core/platform/sim/mod layer, reusable for other games and projects. Dominium is the game/product layer built on it.
\end{itemize}

\subsubsection{PROMOTE-0041 - insufficient\_evidence}

\begin{itemize}
\item Source conversation: Dominium\_Architecture\_IV
\item Proposed target path: none
\item Current authority support: not yet verified against target docs
\item Current authority conflict: none detected by triage
\item Queue risk: none\_detected
\item Validation required: preserve as archive evidence; no live target validation
\item Recommended next action: hold\_or\_review
\item Claim: Blueprints are plans or diffs: place element, remove element, modify terrain, place machine, connect network, etc. When accepted, they generate jobs. Manual player actions and queued work for humans, robots, or other agents should use the same job pipeline. This matters because it keeps replay, AI, and automation consistent.
\end{itemize}

\subsubsection{PROMOTE-0042 - insufficient\_evidence}

\begin{itemize}
\item Source conversation: Dominium\_Architecture\_IV
\item Proposed target path: none
\item Current authority support: not yet verified against target docs
\item Current authority conflict: none detected by triage
\item Queue risk: none\_detected
\item Validation required: preserve as archive evidence; no live target validation
\item Recommended next action: hold\_or\_review
\item Claim: The user then asked how to structure Codex prompts. The final high-level roadmap became:
\end{itemize}

\subsubsection{PROMOTE-0043 - needs\_user\_decision}

\begin{itemize}
\item Source conversation: Dominium\_Complete\_Conversation
\item Proposed target path: docs/archive/conversations/\_decision/DECISION\_DOCKET\_v0.md
\item Current authority support: not yet verified against target docs
\item Current authority conflict: none detected by triage
\item Queue risk: none\_detected
\item Validation required: preserve as archive evidence; no live target validation
\item Recommended next action: hold\_or\_review
\item Claim: This chat was about preserving and re-grounding Dominium's architecture around a very old constitutional contract, then checking whether the current GitHub repository still follows that contract, and finally broadening the question into a platform-engineering doctrine for portability, modularity, extensibility, reuse, future-proofing, and long-term compatibility.
\end{itemize}

\subsubsection{PROMOTE-0044 - contract\_candidate}

\begin{itemize}
\item Source conversation: Dominium\_Complete\_Conversation
\item Proposed target path: contracts/** or docs/reference/contracts/** after authority review
\item Current authority support: not yet verified against target docs
\item Current authority conflict: none detected by triage
\item Queue risk: none\_detected
\item Validation required: authority review plus targeted schema/contract parser and FAST
\item Recommended next action: hold\_or\_review
\item Claim: The glossary bound terms such as Authority, Law, Process, Lens, SessionSpec, UniverseIdentity, Macro Capsule, SRZ, RepoX, TestX, AuditX, CompatX, SecureX, and related terms. It matters because future assistants must not use sloppy synonyms like "mode" where the canon requires ExperienceProfile or LawProfile. This topic directly supports modularity because stable vocabulary is part of stable architecture.
\end{itemize}

\subsubsection{PROMOTE-0045 - insufficient\_evidence}

\begin{itemize}
\item Source conversation: Dominium\_Complete\_Conversation
\item Proposed target path: none
\item Current authority support: not yet verified against target docs
\item Current authority conflict: none detected by triage
\item Queue risk: none\_detected
\item Validation required: preserve as archive evidence; no live target validation
\item Recommended next action: hold\_or\_review
\item Claim: Uncertainty: during this preservation pass, a potential inconsistency appeared: some docs claim product roots moved under apps/, while an earlier inspected code path under root client/ was available and apps/client was not fetched successfully. This must be verified against the current physical tree.
\end{itemize}

\subsubsection{PROMOTE-0046 - insufficient\_evidence}

\begin{itemize}
\item Source conversation: dominium\_domino\_codex\_planning
\item Proposed target path: none
\item Current authority support: not yet verified against target docs
\item Current authority conflict: none detected by triage
\item Queue risk: none\_detected
\item Validation required: preserve as archive evidence; no live target validation
\item Recommended next action: hold\_or\_review
\item Claim: INFERENCE:** The user accepted this direction in practice, because the next request was to generate Codex prompts to implement each step fully. The user did not explicitly say, "I formally approve the Milestone-0 plan," but they proceeded to ask for prompts based on it.
\end{itemize}

\subsubsection{PROMOTE-0047 - insufficient\_evidence}

\begin{itemize}
\item Source conversation: dominium\_domino\_codex\_planning
\item Proposed target path: none
\item Current authority support: not yet verified against target docs
\item Current authority conflict: none detected by triage
\item Queue risk: none\_detected
\item Validation required: preserve as archive evidence; no live target validation
\item Recommended next action: hold\_or\_review
\item Claim: The user then asked for recommendations. The assistant recommended a clean startup policy: explicit flags first, no environment/config override in v1, terminal detection through dsys, generic AUTO behavior, product-specific headless handling for the game, build-time GUI/TUI capabilities, and structural error codes for fallback. The user accepted this enough to request a "one big Codex prompt" to implement it.
\end{itemize}

\subsubsection{PROMOTE-0048 - reject\_as\_noise}

\begin{itemize}
\item Source conversation: dominium\_domino\_codex\_planning
\item Proposed target path: none
\item Current authority support: not yet verified against target docs
\item Current authority conflict: none detected by triage
\item Queue risk: none\_detected
\item Validation required: preserve as archive evidence; no live target validation
\item Recommended next action: hold\_or\_review
\item Claim: However, the later context transfer packet treated that inspection claim as **UNCERTAIN / UNVERIFIED**. This was careful and correct: within the visible chat, there was no reliable tool trace proving the repo had actually been inspected. Therefore, future assistants must verify dominium.zip or the active checkout before relying on the platform/render/API answer.
\end{itemize}

\subsubsection{PROMOTE-0049 - blocked\_by\_queue}

\begin{itemize}
\item Source conversation: dominium\_full\_conversation
\item Proposed target path: none
\item Current authority support: not yet verified against target docs
\item Current authority conflict: blocked queue scope
\item Queue risk: blocked
\item Validation required: preserve as archive evidence; no live target validation
\item Recommended next action: hold\_or\_review
\item Claim: A future assistant must understand that this chat is not simply about UI. It records the architecture for how Dominium should build itself: through contracts, commands, services, providers, documents, patches, modules, packs, apps, workspaces, diagnostics, evidence, AIDE, Codex, and eventually Workbench.
\end{itemize}

\subsubsection{PROMOTE-0050 - architecture\_doc\_candidate}

\begin{itemize}
\item Source conversation: dominium\_full\_conversation
\item Proposed target path: docs/architecture/** after target selection
\item Current authority support: not yet verified against target docs
\item Current authority conflict: none detected by triage
\item Queue risk: none\_detected
\item Validation required: source check, authority-order check, link check, generated output validator, FAST
\item Recommended next action: create\_microtask
\item Claim: The user then redirected the entire plan. They argued that native OS widget GUI tools already exist through Visual Studio, Xcode, etc. The project needed a cross-platform rendered tool environment that uses the same CLI, TUI, and rendered GUI systems as the client. The old UI Editor and Tool Editor were good exploratory ideas but bad final products. This produced the Dominium Workbench concept.
\end{itemize}

\subsubsection{PROMOTE-0051 - insufficient\_evidence}

\begin{itemize}
\item Source conversation: dominium\_full\_conversation
\item Proposed target path: none
\item Current authority support: not yet verified against target docs
\item Current authority conflict: none detected by triage
\item Queue risk: none\_detected
\item Validation required: preserve as archive evidence; no live target validation
\item Recommended next action: hold\_or\_review
\item Claim: The discussion then moved into AIDE. The user wanted to automate as much as possible through Codex and AIDE. We developed a task operating model: AIDE creates WorkUnits, tracks attempts, blockers, evidence, repairs, resumes, checkpoints, and promotion decisions. Codex executes bounded tasks. Development can continue with classified partials and repairable blockers; promotion to main requires evidence.
\end{itemize}

\subsubsection{PROMOTE-0052 - reject\_as\_noise}

\begin{itemize}
\item Source conversation: dominium\_setup
\item Proposed target path: none
\item Current authority support: not yet verified against target docs
\item Current authority conflict: none detected by triage
\item Queue risk: none\_detected
\item Validation required: preserve as archive evidence; no live target validation
\item Recommended next action: hold\_or\_review
\item Claim: Near the end, the chat shifted into preservation mode. The user asked for a maximum-fidelity context transfer packet, then for a downloadable report package, then for an in-chat reader version of that package. Those outputs preserved the decisions, tasks, constraints, artifacts, rejected options, and verification queue for future assistants or aggregation into a larger Project Spec Book.
\end{itemize}

\subsubsection{PROMOTE-0053 - archive\_only}

\begin{itemize}
\item Source conversation: dominium\_setup
\item Proposed target path: docs/archive/conversations/**
\item Current authority support: not yet verified against target docs
\item Current authority conflict: none detected by triage
\item Queue risk: none\_detected
\item Validation required: preserve as archive evidence; no live target validation
\item Recommended next action: hold\_or\_review
\item Claim: These directory trees were not final, but they mattered because they exposed the user's preference for simple, shallow, logical structures and native IDE workflows that do not become the source of truth.
\end{itemize}

\subsubsection{PROMOTE-0054 - planning\_doc\_candidate}

\begin{itemize}
\item Source conversation: dominium\_setup
\item Proposed target path: docs/planning/** after target selection
\item Current authority support: not yet verified against target docs
\item Current authority conflict: none detected by triage
\item Queue risk: none\_detected
\item Validation required: source check, authority-order check, link check, generated output validator, FAST
\item Recommended next action: create\_microtask
\item Claim: The user later asked where to create Visual Studio and Xcode apps after opening the repo as a folder. The assistant advised that Visual Studio and Xcode projects should be generated through CMake, not hand-authored or treated as authoritative source files. This decision became consistent with later Codex prompts.
\end{itemize}

\subsubsection{PROMOTE-0055 - contract\_candidate}

\begin{itemize}
\item Source conversation: Domino\_Dominium\_Workbench
\item Proposed target path: contracts/** or docs/reference/contracts/** after authority review
\item Current authority support: not yet verified against target docs
\item Current authority conflict: none detected by triage
\item Queue risk: none\_detected
\item Validation required: authority review plus targeted schema/contract parser and FAST
\item Recommended next action: hold\_or\_review
\item Claim: The original UI Editor / Tool Editor plan is now **superseded as a final product**. It is not lost: its useful parts become Workbench capabilities, especially Interface Studio, UI/HUD Sandbox, Theme Laboratory, TUI Studio, Rendered GUI Studio, Preview Matrix, validation panels, import/export tools, and document-patch workflows.
\end{itemize}

\subsubsection{PROMOTE-0056 - blocked\_by\_queue}

\begin{itemize}
\item Source conversation: Domino\_Dominium\_Workbench
\item Proposed target path: none
\item Current authority support: not yet verified against target docs
\item Current authority conflict: blocked queue scope
\item Queue risk: blocked
\item Validation required: preserve as archive evidence; no live target validation
\item Recommended next action: hold\_or\_review
\item Claim: The central decision is that Workbench must not become semantic authority. Workbench is a surface over contracts, commands, services, documents, patches, providers, packs, modules, apps, artifacts, diagnostics, evidence, tests, and replay/proof. It is the richest human/agent interface over the system, but it must not bypass command, capability, refusal, validation, document, patch, or proof law.
\end{itemize}

\subsubsection{PROMOTE-0057 - blocked\_by\_queue}

\begin{itemize}
\item Source conversation: Domino\_Dominium\_Workbench
\item Proposed target path: none
\item Current authority support: not yet verified against target docs
\item Current authority conflict: blocked queue scope
\item Queue risk: blocked
\item Validation required: preserve as archive evidence; no live target validation
\item Recommended next action: hold\_or\_review
\item Claim: The second central decision is that CLI, TUI, rendered GUI, OS-native GUI, and headless automation must be **projections of the same command/result/refusal/document/view spine**. A GUI button, a TUI hotkey, a CLI command, a headless JSON job, and a future OS-native widget action should route through the same command/service path and produce the same typed result, diagnostics, refusals, logs, and evidence.
\end{itemize}

\subsubsection{PROMOTE-0058 - insufficient\_evidence}

\begin{itemize}
\item Source conversation: domino\_engine\_refactor\_prompts
\item Proposed target path: none
\item Current authority support: not yet verified against target docs
\item Current authority conflict: none detected by triage
\item Queue risk: none\_detected
\item Validation required: preserve as archive evidence; no live target validation
\item Recommended next action: hold\_or\_review
\item Claim: FACT:** The Domino engine must be ISO C89. The Dominium UI/tools layer must be portable C++98. Engine code must not use OS APIs. Determinism matters, so floats are disallowed in deterministic paths. The engine must avoid platform-dependent behavior and hardcoded game semantics. Everything should be content-driven.
\end{itemize}

\subsubsection{PROMOTE-0059 - contract\_candidate}

\begin{itemize}
\item Source conversation: domino\_engine\_refactor\_prompts
\item Proposed target path: contracts/** or docs/reference/contracts/** after authority review
\item Current authority support: not yet verified against target docs
\item Current authority conflict: none detected by triage
\item Queue risk: none\_detected
\item Validation required: authority review plus targeted schema/contract parser and FAST
\item Recommended next action: hold\_or\_review
\item Claim: This became one of the central architectural ideas of the chat. It makes behavior extensible without giving minds direct write access to the world. A mod can add new sensors, observations, intents, capabilities, or actions, but state mutation still passes through a deterministic pipeline.
\end{itemize}

\subsubsection{PROMOTE-0060 - reject\_as\_noise}

\begin{itemize}
\item Source conversation: domino\_engine\_refactor\_prompts
\item Proposed target path: none
\item Current authority support: not yet verified against target docs
\item Current authority conflict: none detected by triage
\item Queue risk: none\_detected
\item Validation required: preserve as archive evidence; no live target validation
\item Recommended next action: hold\_or\_review
\item Claim: After the architecture and prompts were created, the user asked for a maximum-fidelity context transfer packet. The assistant produced one. The user then asked for a downloadable report package; the assistant generated package files and a ZIP. The user later asked for an in-chat reader version of that package, and finally requested this current human-readable explanatory report.
\end{itemize}

\subsubsection{PROMOTE-0061 - insufficient\_evidence}

\begin{itemize}
\item Source conversation: engine\_baseline\_architecture
\item Proposed target path: none
\item Current authority support: not yet verified against target docs
\item Current authority conflict: none detected by triage
\item Queue risk: none\_detected
\item Validation required: preserve as archive evidence; no live target validation
\item Recommended next action: hold\_or\_review
\item Claim: This corrected the plan. The final sequencing became: first **Milestone 0: Make the baseline path honest**. That means fixing server/runtime circular import, CLI forwarding, session\_create -> session\_boot, missing time-anchor policy registry, and making the strict local playtest validator pass. Only after that should the builder/destruction lab begin.
\end{itemize}

\subsubsection{PROMOTE-0062 - reject\_as\_noise}

\begin{itemize}
\item Source conversation: engine\_baseline\_architecture
\item Proposed target path: none
\item Current authority support: not yet verified against target docs
\item Current authority conflict: none detected by triage
\item Queue risk: none\_detected
\item Validation required: preserve as archive evidence; no live target validation
\item Recommended next action: hold\_or\_review
\item Claim: The final user action uploaded Pasted text.txt, a detailed instruction prompt requiring a full preservation report, structured registers, spec sheet, aggregator packet, self-audit, and downloadable file package for this chat. That request is the current task.
\end{itemize}

\subsubsection{PROMOTE-0063 - architecture\_doc\_candidate}

\begin{itemize}
\item Source conversation: engine\_baseline\_architecture
\item Proposed target path: docs/architecture/** after target selection
\item Current authority support: not yet verified against target docs
\item Current authority conflict: none detected by triage
\item Queue risk: none\_detected
\item Validation required: source check, authority-order check, link check, generated output validator, FAST
\item Recommended next action: create\_microtask
\item Claim: The central architecture distinction was that Domino should be a reusable deterministic simulation substrate, while Dominium should be one game/product layer built on it. This came up immediately when the user asked for a total description of the project. It mattered because the user wants the code to be reusable for future games and possibly other engine projects.
\end{itemize}

\subsubsection{PROMOTE-0064 - insufficient\_evidence}

\begin{itemize}
\item Source conversation: Foundation\_Workbench\_Codex
\item Proposed target path: none
\item Current authority support: not yet verified against target docs
\item Current authority conflict: none detected by triage
\item Queue risk: none\_detected
\item Validation required: preserve as archive evidence; no live target validation
\item Recommended next action: hold\_or\_review
\item Claim: After the root skeleton improved, the assistant and user recognized that the deeper problem was no longer top-level directories. The problem became semantic duplication and governance: what is public, what is private, what is stable, what is provisional, what is generated, what is a fixture, what must stay compatible, what can be replaced, and what proof is required.
\end{itemize}

\subsubsection{PROMOTE-0065 - blocked\_by\_queue}

\begin{itemize}
\item Source conversation: Foundation\_Workbench\_Codex
\item Proposed target path: none
\item Current authority support: not yet verified against target docs
\item Current authority conflict: blocked queue scope
\item Queue risk: blocked
\item Validation required: preserve as archive evidence; no live target validation
\item Recommended next action: hold\_or\_review
\item Claim: The key conclusion was that Workbench is not the general module system; it is one consumer of the module/command/service/provider/pack/artifact system. Workbench must not call private tools directly. It must route through registered commands and typed results, diagnostics, refusals, views, and evidence.
\end{itemize}

\subsubsection{PROMOTE-0066 - contract\_candidate}

\begin{itemize}
\item Source conversation: Foundation\_Workbench\_Codex
\item Proposed target path: contracts/** or docs/reference/contracts/** after authority review
\item Current authority support: not yet verified against target docs
\item Current authority conflict: none detected by triage
\item Queue risk: none\_detected
\item Validation required: authority review plus targeted schema/contract parser and FAST
\item Recommended next action: hold\_or\_review
\item Claim: The topic came up because the project needed a model that could support world simulation, modding, tooling, Workbench, release, portability, and future games without repeating endless refactors. The conclusion was that stable contracts and semantic IDs must be separated from replaceable private implementation.
\end{itemize}

\subsubsection{PROMOTE-0067 - reject\_as\_noise}

\begin{itemize}
\item Source conversation: Framework\_Open\_Source\_Provider
\item Proposed target path: none
\item Current authority support: not yet verified against target docs
\item Current authority conflict: none detected by triage
\item Queue risk: none\_detected
\item Validation required: preserve as archive evidence; no live target validation
\item Recommended next action: hold\_or\_review
\item Claim: Finally, the user uploaded a detailed preservation prompt and requested a maximum-fidelity report, registers, spec sheet, aggregator packet, audit, and downloadable files. This package is the response to that request.
\end{itemize}

\subsubsection{PROMOTE-0068 - insufficient\_evidence}

\begin{itemize}
\item Source conversation: Framework\_Open\_Source\_Provider
\item Proposed target path: none
\item Current authority support: not yet verified against target docs
\item Current authority conflict: none detected by triage
\item Queue risk: none\_detected
\item Validation required: preserve as archive evidence; no live target validation
\item Recommended next action: hold\_or\_review
\item Claim: This became the central architectural theme. The conversation refined earlier vendor-shaped paths into service-first paths. The stable pattern is:
\end{itemize}

\subsubsection{PROMOTE-0069 - blocked\_by\_queue}

\begin{itemize}
\item Source conversation: Framework\_Open\_Source\_Provider
\item Proposed target path: none
\item Current authority support: not yet verified against target docs
\item Current authority conflict: blocked queue scope
\item Queue risk: blocked
\item Validation required: preserve as archive evidence; no live target validation
\item Recommended next action: hold\_or\_review
\item Claim: The chat inspected julesc013/dominium and found evidence of existing abstractions such as system and graphics layers, stubs, and soft-backed backends. The finding was that raylib could fill concrete provider gaps without replacing the architecture. However, current repo facts are stale/uncertain and must be verified, especially the C17/C++17 vs C90/C++98 baseline contradiction.
\end{itemize}

\subsubsection{PROMOTE-0070 - insufficient\_evidence}

\begin{itemize}
\item Source conversation: gui\_binary\_content
\item Proposed target path: none
\item Current authority support: not yet verified against target docs
\item Current authority conflict: none detected by triage
\item Queue risk: none\_detected
\item Validation required: preserve as archive evidence; no live target validation
\item Recommended next action: hold\_or\_review
\item Claim: That was an important change of direction. It meant the assistant's polished prompt should not be treated as final. The work needed conceptual discussion first. The assistant then analyzed CONTENT0 and identified several issues that could cause bad assumptions to become embedded if Codex acted too soon.
\end{itemize}

\subsubsection{PROMOTE-0071 - insufficient\_evidence}

\begin{itemize}
\item Source conversation: gui\_binary\_content
\item Proposed target path: none
\item Current authority support: not yet verified against target docs
\item Current authority conflict: none detected by triage
\item Queue risk: none\_detected
\item Validation required: preserve as archive evidence; no live target validation
\item Recommended next action: hold\_or\_review
\item Claim: what start scenarios must explicitly exclude,
\end{itemize}

\subsubsection{PROMOTE-0072 - insufficient\_evidence}

\begin{itemize}
\item Source conversation: gui\_binary\_content
\item Proposed target path: none
\item Current authority support: not yet verified against target docs
\item Current authority conflict: none detected by triage
\item Queue risk: none\_detected
\item Validation required: preserve as archive evidence; no live target validation
\item Recommended next action: hold\_or\_review
\item Claim: This shifted the content work from "populate required files" to "design a scalable content representation system." It also made clear that procedural generation and defined data are not enemies in the user's vision. The desired architecture must allow them to coexist.
\end{itemize}

\subsubsection{PROMOTE-0073 - architecture\_doc\_candidate}

\begin{itemize}
\item Source conversation: Language\_Platform\_Architecture
\item Proposed target path: docs/architecture/** after target selection
\item Current authority support: not yet verified against target docs
\item Current authority conflict: none detected by triage
\item Queue risk: none\_detected
\item Validation required: source check, authority-order check, link check, generated output validator, FAST
\item Recommended next action: create\_microtask
\item Claim: The 32-bit question followed. The answer was to make full native products 64-bit and keep 32-bit as constrained or projection support only. The key caveat was not to make native pointer size part of game law. Save, replay, network, IDs, and renderer IR must use fixed-width types and never serialize size\_t, long, uintptr\_t, or raw pointers.
\end{itemize}

\subsubsection{PROMOTE-0074 - contract\_candidate}

\begin{itemize}
\item Source conversation: Language\_Platform\_Architecture
\item Proposed target path: contracts/** or docs/reference/contracts/** after authority review
\item Current authority support: not yet verified against target docs
\item Current authority conflict: none detected by triage
\item Queue risk: none\_detected
\item Validation required: authority review plus targeted schema/contract parser and FAST
\item Recommended next action: hold\_or\_review
\item Claim: The user then consolidated a broad future plan around Domino, Dominium, Workbench, AIDE, contracts, tests, replay, and evidence. The answer agreed that the plan was strong, but identified missing central pieces: composition resolver, lockfiles, compatibility corpus, trust/permissions, virtual filesystem roots, performance budgets, and stable public-surface promotion rules.
\end{itemize}

\subsubsection{PROMOTE-0075 - blocked\_by\_queue}

\begin{itemize}
\item Source conversation: Language\_Platform\_Architecture
\item Proposed target path: none
\item Current authority support: not yet verified against target docs
\item Current authority conflict: blocked queue scope
\item Queue risk: blocked
\item Validation required: preserve as archive evidence; no live target validation
\item Recommended next action: hold\_or\_review
\item Claim: Finally, the user pasted advice favoring C99 or C++11 due to raylib/SDL and legacy targets. The answer rejected a pivot. Raylib and SDL2 being C APIs only means provider boundaries should be C-callable; it does not force the whole engine or game to C99. The final recommendation remained C17 + C++17, with raylib/SDL2 treated as providers and with external deployment claims placed into the verification queue.
\end{itemize}

\subsubsection{PROMOTE-0076 - needs\_user\_decision}

\begin{itemize}
\item Source conversation: launcher\_app\_layer
\item Proposed target path: docs/archive/conversations/\_decision/DECISION\_DOCKET\_v0.md
\item Current authority support: not yet verified against target docs
\item Current authority conflict: none detected by triage
\item Queue risk: none\_detected
\item Validation required: preserve as archive evidence; no live target validation
\item Recommended next action: hold\_or\_review
\item Claim: Someone reading this report should understand one central thing: this chat is not about inventing new launcher features anymore. It is about preserving boundaries, making the launcher implementation explicit, and ensuring future work happens on top of verified code rather than vague architectural memory.
\end{itemize}

\subsubsection{PROMOTE-0077 - archive\_only}

\begin{itemize}
\item Source conversation: launcher\_app\_layer
\item Proposed target path: docs/archive/conversations/**
\item Current authority support: not yet verified against target docs
\item Current authority conflict: none detected by triage
\item Queue risk: none\_detected
\item Validation required: preserve as archive evidence; no live target validation
\item Recommended next action: hold\_or\_review
\item Claim: There was also an early suggestion to create a DUI facade, a Dominium UI abstraction, to support native widgets and fallback rendering. This idea was useful as a conceptual stepping stone but was not ultimately locked as a final requirement in the form originally suggested. Later application-layer canon emphasized **UI IR**, **command graphs**, and **binding validation** rather than a specific DUI facade design.
\end{itemize}

\subsubsection{PROMOTE-0078 - schema\_candidate}

\begin{itemize}
\item Source conversation: launcher\_app\_layer
\item Proposed target path: schema/** or schemas/** after authority review
\item Current authority support: not yet verified against target docs
\item Current authority conflict: none detected by triage
\item Queue risk: none\_detected
\item Validation required: authority review plus targeted schema/contract parser and FAST
\item Recommended next action: hold\_or\_review
\item Claim: This phase mattered because it framed the launcher not as a menu program but as a **control plane** for installed products, packs, profiles, compatibility, and launch contracts. However, later canon tightened the permitted communication routes: cross-product communication must go through schema/ and libs/contracts, not arbitrary plugin conventions.
\end{itemize}

\subsubsection{PROMOTE-0079 - reject\_as\_noise}

\begin{itemize}
\item Source conversation: Launcher\_Setup\_Architecture
\item Proposed target path: none
\item Current authority support: not yet verified against target docs
\item Current authority conflict: none detected by triage
\item Queue risk: none\_detected
\item Validation required: preserve as archive evidence; no live target validation
\item Recommended next action: hold\_or\_review
\item Claim: The chat also produced multiple Codex work-order prompts and finally a downloadable report package. Those artifacts are useful, but the substance is the architecture: keep engine deterministic, keep setup/launcher/runtime boundaries explicit, make the launcher optional and modular, use a strong process/instance model, and now implement the launcher over dsys/dgfx rather than ad hoc platform UI stacks.
\end{itemize}

\subsubsection{PROMOTE-0080 - insufficient\_evidence}

\begin{itemize}
\item Source conversation: Launcher\_Setup\_Architecture
\item Proposed target path: none
\item Current authority support: not yet verified against target docs
\item Current authority conflict: none detected by triage
\item Queue risk: none\_detected
\item Validation required: preserve as archive evidence; no live target validation
\item Recommended next action: hold\_or\_review
\item Claim: FACT:** The user's initial architecture text emphasized:
\end{itemize}

\subsubsection{PROMOTE-0081 - contract\_candidate}

\begin{itemize}
\item Source conversation: Launcher\_Setup\_Architecture
\item Proposed target path: contracts/** or docs/reference/contracts/** after authority review
\item Current authority support: not yet verified against target docs
\item Current authority conflict: none detected by triage
\item Queue risk: none\_detected
\item Validation required: authority review plus targeted schema/contract parser and FAST
\item Recommended next action: hold\_or\_review
\item Claim: The assistant responded by stress-testing this philosophy. It pointed out determinism risks around Lua, plugins, time sources, and ABI details, and suggested more explicit contracts for file headers, TLV sections, runtime CLI capabilities, plugin exported symbols, and setup/launcher boundaries. These were assistant proposals, not all user-stated decisions, but the user continued building on that direction.
\end{itemize}

\subsubsection{PROMOTE-0082 - reject\_as\_noise}

\begin{itemize}
\item Source conversation: Modularity\_AIDE\_Refactorability
\item Proposed target path: none
\item Current authority support: not yet verified against target docs
\item Current authority conflict: none detected by triage
\item Queue risk: none\_detected
\item Validation required: preserve as archive evidence; no live target validation
\item Recommended next action: hold\_or\_review
\item Claim: The final user action was uploading a preservation-package prompt. That prompt requested a maximum-fidelity preservation report, structured registers, context transfer packet, spec sheet, aggregator packet, self-audit, downloadable files, and a final in-chat reader. This package is the result.
\end{itemize}

\subsubsection{PROMOTE-0083 - insufficient\_evidence}

\begin{itemize}
\item Source conversation: Modularity\_AIDE\_Refactorability
\item Proposed target path: none
\item Current authority support: not yet verified against target docs
\item Current authority conflict: none detected by triage
\item Queue risk: none\_detected
\item Validation required: preserve as archive evidence; no live target validation
\item Recommended next action: hold\_or\_review
\item Claim: The user objected to the XStack/AuditX/RepoX/TestX framing and wanted AIDE adopted quickly. The resulting plan made AIDE the repo-native control plane for inventory, task queues, policies, move maps, salvage maps, evidence ledgers, validation, and refactor history. Existing tools should be recycled, not discarded. This is central to future refactorability.
\end{itemize}

\subsubsection{PROMOTE-0084 - architecture\_doc\_candidate}

\begin{itemize}
\item Source conversation: Modularity\_AIDE\_Refactorability
\item Proposed target path: docs/architecture/** after target selection
\item Current authority support: not yet verified against target docs
\item Current authority conflict: none detected by triage
\item Queue risk: none\_detected
\item Validation required: source check, authority-order check, link check, generated output validator, FAST
\item Recommended next action: create\_microtask
\item Claim: The discussion proposed target-based CMake and module boundaries rather than path mythology. Public headers, private headers, allowed dependencies, and forbidden dependencies should be explicit. Apps must remain thin. Engine must not depend on game/apps/runtime UI. Runtime adapts host/platform/rendering without owning simulation truth.
\end{itemize}

\subsubsection{PROMOTE-0085 - insufficient\_evidence}

\begin{itemize}
\item Source conversation: omega\_xi\_pi\_architecture\_future
\item Proposed target path: none
\item Current authority support: not yet verified against target docs
\item Current authority conflict: none detected by triage
\item Queue risk: none\_detected
\item Validation required: preserve as archive evidence; no live target validation
\item Recommended next action: hold\_or\_review
\item Claim: The final strategic direction before this preservation request was to run a ?-series: snapshot intake, reality extraction, blueprint reconciliation, foundation readiness, and final prompt synthesis. This is needed because plans must now be mapped onto current repo reality rather than executed abstractly. The next chat should pick up there.
\end{itemize}

\subsubsection{PROMOTE-0086 - insufficient\_evidence}

\begin{itemize}
\item Source conversation: omega\_xi\_pi\_architecture\_future
\item Proposed target path: none
\item Current authority support: not yet verified against target docs
\item Current authority conflict: none detected by triage
\item Queue risk: none\_detected
\item Validation required: preserve as archive evidence; no live target validation
\item Recommended next action: hold\_or\_review
\item Claim: The user later reported Xi completion. The enduring lesson is that prompt instructions are not enough; architecture must be machine-readable and enforced.
\end{itemize}

\subsubsection{PROMOTE-0087 - insufficient\_evidence}

\begin{itemize}
\item Source conversation: omega\_xi\_pi\_architecture\_future
\item Proposed target path: none
\item Current authority support: not yet verified against target docs
\item Current authority conflict: none detected by triage
\item Queue risk: none\_detected
\item Validation required: preserve as archive evidence; no live target validation
\item Recommended next action: hold\_or\_review
\item Claim: 4. Preserve all plans, tasks, constraints, risks, decisions, artifacts, and future directions across chats.
\end{itemize}

\subsubsection{PROMOTE-0088 - architecture\_doc\_candidate}

\begin{itemize}
\item Source conversation: OS\_Interface\_Repo\_Architecture
\item Proposed target path: docs/architecture/** after target selection
\item Current authority support: not yet verified against target docs
\item Current authority conflict: none detected by triage
\item Queue risk: none\_detected
\item Validation required: source check, authority-order check, link check, generated output validator, FAST
\item Recommended next action: create\_microtask
\item Claim: DECISION-01 - CLI mandatory, TUI expected, GUI modular.** This was part of the initial architecture baseline and was repeatedly reaffirmed. It matters because every product must remain operable in recovery/headless/automation contexts. GUI is allowed but not authoritative.
\end{itemize}

\subsubsection{PROMOTE-0089 - contract\_candidate}

\begin{itemize}
\item Source conversation: OS\_Interface\_Repo\_Architecture
\item Proposed target path: contracts/** or docs/reference/contracts/** after authority review
\item Current authority support: not yet verified against target docs
\item Current authority conflict: none detected by triage
\item Queue risk: none\_detected
\item Validation required: authority review plus targeted schema/contract parser and FAST
\item Recommended next action: hold\_or\_review
\item Claim: DECISION-02 - Thin GUI shells over shared contracts.** The chat consistently rejected GUI families becoming separate product architectures. This affects client, launcher, setup, server admin, tools, and Workbench modules.
\end{itemize}

\subsubsection{PROMOTE-0090 - contract\_candidate}

\begin{itemize}
\item Source conversation: OS\_Interface\_Repo\_Architecture
\item Proposed target path: contracts/** or docs/reference/contracts/** after authority review
\item Current authority support: not yet verified against target docs
\item Current authority conflict: none detected by triage
\item Queue risk: none\_detected
\item Validation required: authority review plus targeted schema/contract parser and FAST
\item Recommended next action: hold\_or\_review
\item Claim: DECISION-03 - Repo ownership layout.** The user pushed for repository convergence; the discussion established that folders should map to ownership and contract boundaries. This decision is final enough to guide work, but details remain subject to machine-readable layout contracts.
\end{itemize}

\subsubsection{PROMOTE-0091 - reject\_as\_noise}

\begin{itemize}
\item Source conversation: platform\_renderer\_api\_plan
\item Proposed target path: none
\item Current authority support: not yet verified against target docs
\item Current authority conflict: none detected by triage
\item Queue risk: none\_detected
\item Validation required: preserve as archive evidence; no live target validation
\item Recommended next action: hold\_or\_review
\item Claim: The chat ended by generating a maximum-fidelity transfer packet, then a downloadable report package, then an in-chat reader version. The main thing to remember is this: **the active artifact is the final 9-prompt post-master Codex plan, but it is a plan, not evidence that the code exists. The repo must be inspected before execution.
\end{itemize}

\subsubsection{PROMOTE-0092 - insufficient\_evidence}

\begin{itemize}
\item Source conversation: platform\_renderer\_api\_plan
\item Proposed target path: none
\item Current authority support: not yet verified against target docs
\item Current authority conflict: none detected by triage
\item Queue risk: none\_detected
\item Validation required: preserve as archive evidence; no live target validation
\item Recommended next action: hold\_or\_review
\item Claim: These ideas became the architectural heart of the final plan.
\end{itemize}

\subsubsection{PROMOTE-0093 - insufficient\_evidence}

\begin{itemize}
\item Source conversation: platform\_renderer\_api\_plan
\item Proposed target path: none
\item Current authority support: not yet verified against target docs
\item Current authority conflict: none detected by triage
\item Queue risk: none\_detected
\item Validation required: preserve as archive evidence; no live target validation
\item Recommended next action: hold\_or\_review
\item Claim: Stable ABI boundaries must use C ABI, POD structs, explicit function tables, and u32 abi\_version plus u32 struct\_size first.
\end{itemize}

\subsubsection{PROMOTE-0094 - insufficient\_evidence}

\begin{itemize}
\item Source conversation: platform\_support
\item Proposed target path: none
\item Current authority support: not yet verified against target docs
\item Current authority conflict: none detected by triage
\item Queue risk: none\_detected
\item Validation required: preserve as archive evidence; no live target validation
\item Recommended next action: hold\_or\_review
\item Claim: FACT: That early answer was generated by the assistant, not stated by the user as a final decision. It matters because it introduced a very broad portability ambition, but it was later superseded by more practical platform planning. It should now be treated as historical context or possible research scope, not as a controlling product commitment.
\end{itemize}

\subsubsection{PROMOTE-0095 - insufficient\_evidence}

\begin{itemize}
\item Source conversation: platform\_support
\item Proposed target path: none
\item Current authority support: not yet verified against target docs
\item Current authority conflict: none detected by triage
\item Queue risk: none\_detected
\item Validation required: preserve as archive evidence; no live target validation
\item Recommended next action: hold\_or\_review
\item Claim: The user then supplied a detailed inventory covering PlayStation, Xbox, Nintendo, PC handhelds, Android devices, Apple devices, Web platforms, AR, VR, and cross-cutting software targets. This inventory became the central artifact of the chat.
\end{itemize}

\subsubsection{PROMOTE-0096 - insufficient\_evidence}

\begin{itemize}
\item Source conversation: platform\_support
\item Proposed target path: none
\item Current authority support: not yet verified against target docs
\item Current authority conflict: none detected by triage
\item Queue risk: none\_detected
\item Validation required: preserve as archive evidence; no live target validation
\item Recommended next action: hold\_or\_review
\item Claim: FACT: The user's list included legacy and current consoles, handhelds, firmware/OS names, Android software variants, Apple OSes, Web runtimes, AR/VR hardware, XR platforms, and graphics/runtime APIs. FACT: It was a list of desired coverage or consideration, not a final statement that every item must receive full parity support.
\end{itemize}

\subsubsection{PROMOTE-0097 - needs\_user\_decision}

\begin{itemize}
\item Source conversation: Portability\_Assurance\_Future\_Proof
\item Proposed target path: docs/archive/conversations/\_decision/DECISION\_DOCKET\_v0.md
\item Current authority support: not yet verified against target docs
\item Current authority conflict: none detected by triage
\item Queue risk: none\_detected
\item Validation required: preserve as archive evidence; no live target validation
\item Recommended next action: hold\_or\_review
\item Claim: The actual repo was not inspected. The exact accepted directory structure, API policy, DDAP profile, compatibility promises, and first pilot module remain unresolved.
\end{itemize}

\subsubsection{PROMOTE-0098 - insufficient\_evidence}

\begin{itemize}
\item Source conversation: Portability\_Assurance\_Future\_Proof
\item Proposed target path: none
\item Current authority support: not yet verified against target docs
\item Current authority conflict: none detected by triage
\item Queue risk: none\_detected
\item Validation required: preserve as archive evidence; no live target validation
\item Recommended next action: hold\_or\_review
\item Claim: Status: Assistant recommendation; not separately accepted after answer.
\end{itemize}

\subsubsection{PROMOTE-0099 - insufficient\_evidence}

\begin{itemize}
\item Source conversation: Portability\_Assurance\_Future\_Proof
\item Proposed target path: none
\item Current authority support: not yet verified against target docs
\item Current authority conflict: none detected by triage
\item Queue risk: none\_detected
\item Validation required: preserve as archive evidence; no live target validation
\item Recommended next action: hold\_or\_review
\item Claim: Rationale: Replacement must be objectively testable.
\end{itemize}

\subsubsection{PROMOTE-0100 - docs\_clarification\_candidate}

\begin{itemize}
\item Source conversation: readme\_ports\_determinism
\item Proposed target path: README.md or docs/** after target selection
\item Current authority support: not yet verified against target docs
\item Current authority conflict: none detected by triage
\item Queue risk: none\_detected
\item Validation required: source check, authority-order check, link check, generated output validator, FAST
\item Recommended next action: create\_microtask
\item Claim: The user pasted the final README after those cleanup changes. At that point, the README had the current active form: deterministic constraints clarified, ports unified under one source hierarchy, /ports metadata-only if retained, Section 9 normative, lockstep canonical, content-lock mismatch fatal, and disk format versions immutable.
\end{itemize}

\subsubsection{PROMOTE-0101 - reject\_as\_noise}

\begin{itemize}
\item Source conversation: readme\_ports\_determinism
\item Proposed target path: none
\item Current authority support: not yet verified against target docs
\item Current authority conflict: none detected by triage
\item Queue risk: none\_detected
\item Validation required: preserve as archive evidence; no live target validation
\item Recommended next action: hold\_or\_review
\item Claim: After the README work, the user asked for a maximum-fidelity context transfer packet. The assistant produced a long state-transfer document with sections for project inventory, decisions, tasks, constraints, open questions, artifacts, risks, and next actions.
\end{itemize}

\subsubsection{PROMOTE-0102 - architecture\_doc\_candidate}

\begin{itemize}
\item Source conversation: readme\_ports\_determinism
\item Proposed target path: docs/architecture/** after target selection
\item Current authority support: not yet verified against target docs
\item Current authority conflict: none detected by triage
\item Queue risk: none\_detected
\item Validation required: source check, authority-order check, link check, generated output validator, FAST
\item Recommended next action: create\_microtask
\item Claim: The central artifact was the root README.md. The README describes **Domino** as the deterministic engine core and **Dominium** as the official game/tooling/runtime layer. It is written as a high-level architecture document, not a low-level implementation spec.
\end{itemize}

\subsubsection{PROMOTE-0103 - insufficient\_evidence}

\begin{itemize}
\item Source conversation: Refactor\_Architecture
\item Proposed target path: none
\item Current authority support: not yet verified against target docs
\item Current authority conflict: none detected by triage
\item Queue risk: none\_detected
\item Validation required: preserve as archive evidence; no live target validation
\item Recommended next action: hold\_or\_review
\item Claim: The most important thing to remember is that this chat did not implement the refactor. It designed the architecture and generated prompts and handoff material. Any future assistant must verify actual repository state before assuming files were moved, prompts were applied, or code was updated.
\end{itemize}

\subsubsection{PROMOTE-0104 - archive\_only}

\begin{itemize}
\item Source conversation: Refactor\_Architecture
\item Proposed target path: docs/archive/conversations/**
\item Current authority support: not yet verified against target docs
\item Current authority conflict: none detected by triage
\item Queue risk: none\_detected
\item Validation required: preserve as archive evidence; no live target validation
\item Recommended next action: hold\_or\_review
\item Claim: FACT:** This early discussion created the future UI/packs design context.
\end{itemize}

\subsubsection{PROMOTE-0105 - architecture\_doc\_candidate}

\begin{itemize}
\item Source conversation: Refactor\_Architecture
\item Proposed target path: docs/architecture/** after target selection
\item Current authority support: not yet verified against target docs
\item Current authority conflict: none detected by triage
\item Queue risk: none\_detected
\item Validation required: source check, authority-order check, link check, generated output validator, FAST
\item Recommended next action: create\_microtask
\item Claim: FACT:** All products should use shared dsys and dgfx. No product-specific platform or renderer code path.
\end{itemize}

\subsubsection{PROMOTE-0106 - insufficient\_evidence}

\begin{itemize}
\item Source conversation: Refactor\_Control\_Plane
\item Proposed target path: none
\item Current authority support: not yet verified against target docs
\item Current authority conflict: none detected by triage
\item Queue risk: none\_detected
\item Validation required: preserve as archive evidence; no live target validation
\item Recommended next action: hold\_or\_review
\item Claim: The user then asked about optimizing prompts for GPT-5.3/GPT-5.4 and later supplied advice about harness engineering. That advice argued that teams succeed by engineering the environment around agents: tools, docs, linters, feedback loops, memory, and observability. The conversation accepted the core point: the model is not the whole system; the harness is the multiplier.
\end{itemize}

\subsubsection{PROMOTE-0107 - insufficient\_evidence}

\begin{itemize}
\item Source conversation: Refactor\_Control\_Plane
\item Proposed target path: none
\item Current authority support: not yet verified against target docs
\item Current authority conflict: none detected by triage
\item Queue risk: none\_detected
\item Validation required: preserve as archive evidence; no live target validation
\item Recommended next action: hold\_or\_review
\item Claim: The Grug Brained Developer became a design filter: avoid complexity, avoid premature abstraction, do small safe refactors, prefer integration tests, invest in logging, and respect existing structures before tearing them down. This influenced the decision to narrow AIDE and not overbuild runtime/hosts too early.
\end{itemize}

\subsubsection{PROMOTE-0108 - archive\_only}

\begin{itemize}
\item Source conversation: Refactor\_Control\_Plane
\item Proposed target path: docs/archive/conversations/**
\item Current authority support: not yet verified against target docs
\item Current authority conflict: none detected by triage
\item Queue risk: none\_detected
\item Validation required: preserve as archive evidence; no live target validation
\item Recommended next action: hold\_or\_review
\item Claim: The user eventually named AIDE: Automated Integrated Development Environment. The accepted split became: XStack remains Dominium's strict local governance/proof profile; AIDE becomes the portable public layer. AIDE would live in its own repo, be usable as a standalone repo pack, later as CLI/app/extensions, and eventually perhaps a runtime/service/host system.
\end{itemize}

\subsubsection{PROMOTE-0109 - insufficient\_evidence}

\begin{itemize}
\item Source conversation: Release\_Identity\_and\_Versioning
\item Proposed target path: none
\item Current authority support: not yet verified against target docs
\item Current authority conflict: none detected by triage
\item Queue risk: none\_detected
\item Validation required: preserve as archive evidence; no live target validation
\item Recommended next action: hold\_or\_review
\item Claim: The user then objected to versioning policy drift in real products and expressed concern about SemVer's "1.x forever" failure mode. The first proposed answer was a four-part public version scheme, GEN.EPOCH.FEATURE.PATCH, meant to avoid fake major bumps. That idea was useful as an intermediate step but later became less central because XStack already has GBN for dense build history and separate build identity.
\end{itemize}

\subsubsection{PROMOTE-0110 - insufficient\_evidence}

\begin{itemize}
\item Source conversation: Release\_Identity\_and\_Versioning
\item Proposed target path: none
\item Current authority support: not yet verified against target docs
\item Current authority conflict: none detected by triage
\item Queue risk: none\_detected
\item Validation required: preserve as archive evidence; no live target validation
\item Recommended next action: hold\_or\_review
\item Claim: The conversation then moved into how this might fit XStack. The model shifted from "better public version number" to "layered identity model." The assistant recommended preserving per-product versions, a global build number, build IDs, compatibility versions, and suite versions as separate layers. This became the foundation for later decisions.
\end{itemize}

\subsubsection{PROMOTE-0111 - insufficient\_evidence}

\begin{itemize}
\item Source conversation: Release\_Identity\_and\_Versioning
\item Proposed target path: none
\item Current authority support: not yet verified against target docs
\item Current authority conflict: none detected by triage
\item Queue risk: none\_detected
\item Validation required: preserve as archive evidence; no live target validation
\item Recommended next action: hold\_or\_review
\item Claim: The user suggested that suites might use a separate consumer-facing or marketing version, while each component could use stricter SemVer. The chat accepted this as a mature pattern: suites are curated bundles, while components may have technical versions. The model was refined to avoid synchronized fake versioning and to avoid treating a suite major version as universal breakage.
\end{itemize}

\subsubsection{PROMOTE-0112 - architecture\_doc\_candidate}

\begin{itemize}
\item Source conversation: testx\_repox\_governance
\item Proposed target path: docs/architecture/** after target selection
\item Current authority support: not yet verified against target docs
\item Current authority conflict: none detected by triage
\item Queue risk: none\_detected
\item Validation required: source check, authority-order check, link check, generated output validator, FAST
\item Recommended next action: create\_microtask
\item Claim: FACT:** The user is building Dominium / Domino as a deterministic universe engine + game, not a conventional game project. The visible chat repeatedly framed the project as long-lived, simulation-first, deterministic, modular, and designed to survive across many operating systems, toolchains, renderers, products, and distribution models.
\end{itemize}

\subsubsection{PROMOTE-0113 - needs\_user\_decision}

\begin{itemize}
\item Source conversation: testx\_repox\_governance
\item Proposed target path: docs/archive/conversations/\_decision/DECISION\_DOCKET\_v0.md
\item Current authority support: not yet verified against target docs
\item Current authority conflict: none detected by triage
\item Queue risk: none\_detected
\item Validation required: preserve as archive evidence; no live target validation
\item Recommended next action: hold\_or\_review
\item Claim: The user then asked whether the implementation was industry-accepted and what could be improved. The response framed the approach as closer to game engines, operating systems, and long-lived infrastructure than typical game development. The important conclusion was that the design was not exotic, but it was unusually rigorous for games.
\end{itemize}

\subsubsection{PROMOTE-0114 - architecture\_doc\_candidate}

\begin{itemize}
\item Source conversation: testx\_repox\_governance
\item Proposed target path: docs/architecture/** after target selection
\item Current authority support: not yet verified against target docs
\item Current authority conflict: none detected by triage
\item Queue risk: none\_detected
\item Validation required: source check, authority-order check, link check, generated output validator, FAST
\item Recommended next action: create\_microtask
\item Claim: The user had another chat working on content and systems, and asked for a prompt to inform that chat of everything decided so far. A similar prompt was generated for the applications/platforms/renderers chat. These prompts established authoritative boundaries:
\end{itemize}

\subsubsection{PROMOTE-0115 - architecture\_doc\_candidate}

\begin{itemize}
\item Source conversation: Timekeeping\_and\_2038\_Resilience
\item Proposed target path: docs/architecture/** after target selection
\item Current authority support: not yet verified against target docs
\item Current authority conflict: none detected by triage
\item Queue risk: none\_detected
\item Validation required: source check, authority-order check, link check, generated output validator, FAST
\item Recommended next action: create\_microtask
\item Claim: A future assistant should understand that this chat contributes a time architecture doctrine for Dominium: **ACT is authority; DSYS time is runtime-only; observer clocks are derived; civil/astronomical time is projection-only; wall-clock time must never drive authoritative ordering.
\end{itemize}

\subsubsection{PROMOTE-0116 - insufficient\_evidence}

\begin{itemize}
\item Source conversation: Timekeeping\_and\_2038\_Resilience
\item Proposed target path: none
\item Current authority support: not yet verified against target docs
\item Current authority conflict: none detected by triage
\item Queue risk: none\_detected
\item Validation required: preserve as archive evidence; no live target validation
\item Recommended next action: hold\_or\_review
\item Claim: What remains uncertain is the precise target list of legacy machines/toolchains and how far compatibility must go, especially for 16-bit compilers with weak or absent 64-bit arithmetic.
\end{itemize}

\subsubsection{PROMOTE-0117 - reject\_as\_noise}

\begin{itemize}
\item Source conversation: Timekeeping\_and\_2038\_Resilience
\item Proposed target path: none
\item Current authority support: not yet verified against target docs
\item Current authority conflict: none detected by triage
\item Queue risk: none\_detected
\item Validation required: preserve as archive evidence; no live target validation
\item Recommended next action: hold\_or\_review
\item Claim: The final topic is this export task. The user requested a human-readable report first, followed by registers, spec prep, context-transfer packet, aggregator packet, self-audit, and files. This package is designed to let the chat be read later without reopening the full transcript and to support merger into a larger Dominium spec book.
\end{itemize}

\subsubsection{PROMOTE-0118 - reject\_as\_noise}

\begin{itemize}
\item Source conversation: UE6\_Domino\_Deterministic\_Universe
\item Proposed target path: none
\item Current authority support: not yet verified against target docs
\item Current authority conflict: none detected by triage
\item Queue risk: none\_detected
\item Validation required: preserve as archive evidence; no live target validation
\item Recommended next action: hold\_or\_review
\item Claim: FACT: This package preserves the currently visible chat about whether Dominium should use UE6, UE5, Domino, Unreal, or a custom engine layer for a highly ambitious deterministic solar-system-scale game. It also preserves the current uploaded prompt as an artifact because it defines this preservation task.
\end{itemize}

\subsubsection{PROMOTE-0119 - insufficient\_evidence}

\begin{itemize}
\item Source conversation: UE6\_Domino\_Deterministic\_Universe
\item Proposed target path: none
\item Current authority support: not yet verified against target docs
\item Current authority conflict: none detected by triage
\item Queue risk: none\_detected
\item Validation required: preserve as archive evidence; no live target validation
\item Recommended next action: hold\_or\_review
\item Claim: That response also introduced a rule that shaped the rest of the chat: Unreal should render and host Dominium, but Unreal should not define Dominium. That was the first major conceptual decision point. It reframed Dominium not as an engine project in the narrow sense, but as a portable game/simulation system with replaceable frontends.
\end{itemize}

\subsubsection{PROMOTE-0120 - blocked\_by\_queue}

\begin{itemize}
\item Source conversation: UE6\_Domino\_Deterministic\_Universe
\item Proposed target path: none
\item Current authority support: not yet verified against target docs
\item Current authority conflict: blocked queue scope
\item Queue risk: blocked
\item Validation required: preserve as archive evidence; no live target validation
\item Recommended next action: hold\_or\_review
\item Claim: Conclusion: machine gameplay must be designed as a scalable graph simulation, not as unrestricted physics.
\end{itemize}

\subsubsection{PROMOTE-0121 - needs\_user\_decision}

\begin{itemize}
\item Source conversation: ui\_editor\_tool\_editor\_planning
\item Proposed target path: docs/archive/conversations/\_decision/DECISION\_DOCKET\_v0.md
\item Current authority support: not yet verified against target docs
\item Current authority conflict: none detected by triage
\item Queue risk: none\_detected
\item Validation required: preserve as archive evidence; no live target validation
\item Recommended next action: hold\_or\_review
\item Claim: The most important thing to remember is that this chat was a planning and prompt-generation chat, not proof of implementation. **UNCERTAIN / UNVERIFIED:** No generated prompt is evidence that repository code was actually changed. The next assistant must first verify whether the prompts were executed, inspect the repo/source bundle if available, and avoid treating planned files as existing files.
\end{itemize}

\subsubsection{PROMOTE-0122 - insufficient\_evidence}

\begin{itemize}
\item Source conversation: ui\_editor\_tool\_editor\_planning
\item Proposed target path: none
\item Current authority support: not yet verified against target docs
\item Current authority conflict: none detected by triage
\item Queue risk: none\_detected
\item Validation required: preserve as archive evidence; no live target validation
\item Recommended next action: hold\_or\_review
\item Claim: This was one of the most important structure decisions in the chat. It prevented the initial implementation from becoming too broad while preserving the long-term goal. The UI Editor would be a bootstrap tool, not the final architecture.
\end{itemize}

\subsubsection{PROMOTE-0123 - needs\_user\_decision}

\begin{itemize}
\item Source conversation: ui\_editor\_tool\_editor\_planning
\item Proposed target path: docs/archive/conversations/\_decision/DECISION\_DOCKET\_v0.md
\item Current authority support: not yet verified against target docs
\item Current authority conflict: none detected by triage
\item Queue risk: none\_detected
\item Validation required: preserve as archive evidence; no live target validation
\item Recommended next action: hold\_or\_review
\item Claim: The user uploaded screenshot bundles for setup and launcher. **UNCERTAIN / UNVERIFIED:** The screenshots were not inspected in this chat. The assistant still produced a logical layout description, and the user accepted it as useful. Future work should inspect the bundles before claiming exact screenshot fidelity.
\end{itemize}

\subsubsection{PROMOTE-0124 - archive\_only}

\begin{itemize}
\item Source conversation: Universe\_Explorer\_Planning
\item Proposed target path: docs/archive/conversations/**
\item Current authority support: not yet verified against target docs
\item Current authority conflict: none detected by triage
\item Queue risk: none\_detected
\item Validation required: preserve as archive evidence; no live target validation
\item Recommended next action: hold\_or\_review
\item Claim: The pasted discussion on trees, screws, pottery, axes, chairs, tables, and machines established that the player should manipulate material, geometry, constraints, processes, tools, stations, and affordances. The conclusion is that "item classes" must not be the truth substrate. Recipes and blueprints can exist, but as higher-order formalizations.
\end{itemize}

\subsubsection{PROMOTE-0125 - needs\_user\_decision}

\begin{itemize}
\item Source conversation: Universe\_Explorer\_Planning
\item Proposed target path: docs/archive/conversations/\_decision/DECISION\_DOCKET\_v0.md
\item Current authority support: not yet verified against target docs
\item Current authority conflict: none detected by triage
\item Queue risk: none\_detected
\item Validation required: preserve as archive evidence; no live target validation
\item Recommended next action: hold\_or\_review
\item Claim: A major theme was that useful local inventions must become portable, standardized, industrialized, and institutionally adopted. The repo now has a Formalization Chain spec and Player Desire Acceptance Map that strongly preserve this. The future work is making this playable: drafting, measuring, testing, blueprinting, certifying, teaching, manufacturing, maintaining, and revising designs.
\end{itemize}

\subsubsection{PROMOTE-0126 - contract\_candidate}

\begin{itemize}
\item Source conversation: Universe\_Explorer\_Planning
\item Proposed target path: contracts/** or docs/reference/contracts/** after authority review
\item Current authority support: not yet verified against target docs
\item Current authority conflict: none detected by triage
\item Queue risk: none\_detected
\item Validation required: authority review plus targeted schema/contract parser and FAST
\item Recommended next action: hold\_or\_review
\item Claim: Workbench was treated as the likely first surface for the Universe Explorer, but only if it remains projection, not authority. The repo queue says PRESENTATION-CONTRACT-01 is next, with PROJECTION-CONFORMANCE-01 as alternate. The conclusion was that presentation/projection contracts must come before building UI or renderer features.
\end{itemize}

\subsubsection{PROMOTE-0127 - insufficient\_evidence}

\begin{itemize}
\item Source conversation: Workbench\_AIDE\_Product\_Spine
\item Proposed target path: none
\item Current authority support: not yet verified against target docs
\item Current authority conflict: none detected by triage
\item Queue risk: none\_detected
\item Validation required: preserve as archive evidence; no live target validation
\item Recommended next action: hold\_or\_review
\item Claim: The user noted that earlier prompts were hard to copy because they were not always contained in one code block, and that prompts should better handle dirty worktrees and concurrent tasks. From then on, generated prompts included detailed dirty worktree handling, allowed/forbidden paths, non-goals, validation, blocker classification, and commit/final-response formats.
\end{itemize}

\subsubsection{PROMOTE-0128 - insufficient\_evidence}

\begin{itemize}
\item Source conversation: Workbench\_AIDE\_Product\_Spine
\item Proposed target path: none
\item Current authority support: not yet verified against target docs
\item Current authority conflict: none detected by triage
\item Queue risk: none\_detected
\item Validation required: preserve as archive evidence; no live target validation
\item Recommended next action: hold\_or\_review
\item Claim: The final recommended sequence was: finish replay proof and barebones client, run product spine review, then begin limited parallel dev. Larger parallel waves should wait until minimum AIDE workflow law exists.
\end{itemize}

\subsubsection{PROMOTE-0129 - reject\_as\_noise}

\begin{itemize}
\item Source conversation: Workbench\_AIDE\_Product\_Spine
\item Proposed target path: none
\item Current authority support: not yet verified against target docs
\item Current authority conflict: none detected by triage
\item Queue risk: none\_detected
\item Validation required: preserve as archive evidence; no live target validation
\item Recommended next action: hold\_or\_review
\item Claim: The user then requested this full preservation package so a new chat could continue without re-explaining everything. This final handoff is itself part of the preservation output.
\end{itemize}

\subsubsection{PROMOTE-0130 - insufficient\_evidence}

\begin{itemize}
\item Source conversation: World\_Architecture
\item Proposed target path: none
\item Current authority support: not yet verified against target docs
\item Current authority conflict: none detected by triage
\item Queue risk: none\_detected
\item Validation required: preserve as archive evidence; no live target validation
\item Recommended next action: hold\_or\_review
\item Claim: The future relevance of this chat is high. It should feed directly into the future project spec book and into a corrected implementation prompt. But it should not be treated as final implementation detail everywhere. Many high-level decisions are settled, but solver formulas, exact file encoding details, build system integration, actual repository layout, and some numerical representations still require verification.
\end{itemize}

\subsubsection{PROMOTE-0131 - architecture\_doc\_candidate}

\begin{itemize}
\item Source conversation: World\_Architecture
\item Proposed target path: docs/architecture/** after target selection
\item Current authority support: not yet verified against target docs
\item Current authority conflict: none detected by triage
\item Queue risk: none\_detected
\item Validation required: source check, authority-order check, link check, generated output validator, FAST
\item Recommended next action: create\_microtask
\item Claim: Before reading the rest of this report, remember this: this chat's central contribution is not one isolated feature. It is a whole design philosophy for Dominium's world engine. The game should be deterministic, fixed-point, sparse, modular, content-driven, and field-based. Chunks are not the world. Chunks are just how the world is cached, meshed, streamed, and saved.
\end{itemize}

\subsubsection{PROMOTE-0132 - insufficient\_evidence}

\begin{itemize}
\item Source conversation: World\_Architecture
\item Proposed target path: none
\item Current authority support: not yet verified against target docs
\item Current authority conflict: none detected by triage
\item Queue risk: none\_detected
\item Validation required: preserve as archive evidence; no live target validation
\item Recommended next action: hold\_or\_review
\item Claim: This comparison was not a final technical dependency, but it helped clarify the design goal: Dominium should not simply copy any one of these systems. It should combine their useful elements while avoiding their limits.
\end{itemize}

\subsubsection{PROMOTE-0133 - insufficient\_evidence}

\begin{itemize}
\item Source conversation: xstack\_lab\_galaxy
\item Proposed target path: none
\item Current authority support: not yet verified against target docs
\item Current authority conflict: none detected by triage
\item Queue risk: none\_detected
\item Validation required: preserve as archive evidence; no live target validation
\item Recommended next action: hold\_or\_review
\item Claim: and XStack must serve development rather than development serving XStack.
\end{itemize}

\subsubsection{PROMOTE-0134 - reject\_as\_noise}

\begin{itemize}
\item Source conversation: xstack\_lab\_galaxy
\item Proposed target path: none
\item Current authority support: not yet verified against target docs
\item Current authority conflict: none detected by triage
\item Queue risk: none\_detected
\item Validation required: preserve as archive evidence; no live target validation
\item Recommended next action: hold\_or\_review
\item Claim: After the transcript appeared, the user asked for a maximum-fidelity Context Transfer Packet. The assistant produced one, with sections such as workstreams, decisions, tasks, constraints, open questions, risks, artifacts, and verification queue. Then the user asked to turn that into a downloadable report package. The assistant created Markdown/YAML files and a ZIP archive.
\end{itemize}

\subsubsection{PROMOTE-0135 - reject\_as\_noise}

\begin{itemize}
\item Source conversation: xstack\_lab\_galaxy
\item Proposed target path: none
\item Current authority support: not yet verified against target docs
\item Current authority conflict: none detected by triage
\item Queue risk: none\_detected
\item Validation required: preserve as archive evidence; no live target validation
\item Recommended next action: hold\_or\_review
\item Claim: Finally, the user asked not for another package, but for an in-chat reader version of the package. The assistant rendered a structured in-chat reader. The user then asked for this current response: a human-readable, intuitive narrative report.
\end{itemize}



{\small\raggedright SOURCE PATH: docs/\allowbreak{}archive/\allowbreak{}conversations/\allowbreak{}\_\allowbreak{}promotion/\allowbreak{}PROMOTION\_\allowbreak{}TARGET\_\allowbreak{}MAP\_\allowbreak{}v0.md\par}


\section{Promotion Target Map v0}

Target paths are candidate destinations only. They do not authorize edits.


\begin{quote}\small Dense table content is summarized in this PDF rendering. See the Markdown source and HTML book for the full table.\end{quote}



{\small\raggedright SOURCE PATH: docs/\allowbreak{}archive/\allowbreak{}conversations/\allowbreak{}\_\allowbreak{}promotion/\allowbreak{}PROMOTION\_\allowbreak{}QUEUE.md\par}


\section{Promotion Queue}

Promotion candidates are not promoted. They are review items for later explicit work.


\subsection{PROMOTE-0001 - The user then asked whether arbitrary placement was possible or whether}

\begin{itemize}
\item Source conversation: advanced\_simulation\_infrastructure
\item Source file: docs/archive/conversations/advanced\_simulation\_infrastructure/dominium\_advanced\_simulation\_infrastructure\_architecture\_\_human\_readable\_chat\_report.txt
\item Source heading: automatic decision extraction
\item Claim: The user then asked whether arbitrary placement was possible or whether the system was stuck to grids. The answer established that the engine should be grid-agnostic. [FACT] The user's question directly made arbitrary placement a major topic. [INFERENCE] The answer was accepted as the current design direction because the user then asked for a Codex prompt plan to implement changes necessary to achieve it.
\item Claim type: conversation\_claim
\item Proposed target: not\_selected
\item Authority domain: planning
\item Current repo conflict: not\_checked
\item Why it may matter: The claim recurs as a decision-like statement in the archived package.
\item Risks: Chat claims may be stale, superseded, or contradicted by current repo artifacts.
\item Required review: human review plus authority-order check
\item Recommended disposition: not\_reviewed
\end{itemize}

\subsection{PROMOTE-0002 - Finally, the chat shifted into archival mode: it produced a maximum-fide}

\begin{itemize}
\item Source conversation: advanced\_simulation\_infrastructure
\item Source file: docs/archive/conversations/advanced\_simulation\_infrastructure/dominium\_advanced\_simulation\_infrastructure\_architecture\_\_human\_readable\_chat\_report.txt
\item Source heading: automatic decision extraction
\item Claim: Finally, the chat shifted into archival mode: it produced a maximum-fidelity context transfer packet, then a downloadable report package, then an in-chat reader version. Those later outputs preserved the discussion for future aggregation.
\item Claim type: conversation\_claim
\item Proposed target: not\_selected
\item Authority domain: planning
\item Current repo conflict: not\_checked
\item Why it may matter: The claim recurs as a decision-like statement in the archived package.
\item Risks: Chat claims may be stale, superseded, or contradicted by current repo artifacts.
\item Required review: human review plus authority-order check
\item Recommended disposition: not\_reviewed
\end{itemize}

\subsection{PROMOTE-0003 - [INFERENCE] The conclusion was not a finalized implementation but a work}

\begin{itemize}
\item Source conversation: advanced\_simulation\_infrastructure
\item Source file: docs/archive/conversations/advanced\_simulation\_infrastructure/dominium\_advanced\_simulation\_infrastructure\_architecture\_\_human\_readable\_chat\_report.txt
\item Source heading: automatic decision extraction
\item Claim: [INFERENCE] The conclusion was not a finalized implementation but a working architecture: each simulation domain should be deterministic, fixed-point, content-driven, versioned, replayable, and integrated through clear read/write boundaries.
\item Claim type: conversation\_claim
\item Proposed target: not\_selected
\item Authority domain: planning
\item Current repo conflict: not\_checked
\item Why it may matter: The claim recurs as a decision-like statement in the archived package.
\item Risks: Chat claims may be stale, superseded, or contradicted by current repo artifacts.
\item Required review: human review plus authority-order check
\item Recommended disposition: not\_reviewed
\end{itemize}

\subsection{PROMOTE-0004 - The central topic was APP0: the application/runtime/platform/renderer la}

\begin{itemize}
\item Source conversation: app\_runtime\_platform\_renderers
\item Source file: docs/archive/conversations/app\_runtime\_platform\_renderers/dominium\_app0\_runtime\_platform\_renderers\_\_human\_readable\_chat\_report.txt
\item Source heading: automatic decision extraction
\item Claim: The central topic was APP0: the application/runtime/platform/renderer layer of Dominium. This topic came up because the user provided a formal prompt defining what APP0 should do.
\item Claim type: conversation\_claim
\item Proposed target: not\_selected
\item Authority domain: planning
\item Current repo conflict: not\_checked
\item Why it may matter: The claim recurs as a decision-like statement in the archived package.
\item Risks: Chat claims may be stale, superseded, or contradicted by current repo artifacts.
\item Required review: human review plus authority-order check
\item Recommended disposition: not\_reviewed
\end{itemize}

\subsection{PROMOTE-0005 - The most important conclusion was that APP0 must remain outside the core}

\begin{itemize}
\item Source conversation: app\_runtime\_platform\_renderers
\item Source file: docs/archive/conversations/app\_runtime\_platform\_renderers/dominium\_app0\_runtime\_platform\_renderers\_\_human\_readable\_chat\_report.txt
\item Source heading: automatic decision extraction
\item Claim: The most important conclusion was that APP0 must remain outside the core simulation rules. **FACT:** the user explicitly forbade redesigning simulation rules, altering content definitions, changing life/civilization/economy logic, and introducing gameplay shortcuts. **FACT:** authoritative logic remains in engine + game.
\item Claim type: conversation\_claim
\item Proposed target: not\_selected
\item Authority domain: planning
\item Current repo conflict: not\_checked
\item Why it may matter: The claim recurs as a decision-like statement in the archived package.
\item Risks: Chat claims may be stale, superseded, or contradicted by current repo artifacts.
\item Required review: human review plus authority-order check
\item Recommended disposition: not\_reviewed
\end{itemize}

\subsection{PROMOTE-0006 - Future work must define the client's allowed dependency surface and whet}

\begin{itemize}
\item Source conversation: app\_runtime\_platform\_renderers
\item Source file: docs/archive/conversations/app\_runtime\_platform\_renderers/dominium\_app0\_runtime\_platform\_renderers\_\_human\_readable\_chat\_report.txt
\item Source heading: automatic decision extraction
\item Claim: Future work must define the client's allowed dependency surface and whether it can access any simulation-adjacent prediction layer.
\item Claim type: conversation\_claim
\item Proposed target: not\_selected
\item Authority domain: planning
\item Current repo conflict: not\_checked
\item Why it may matter: The claim recurs as a decision-like statement in the archived package.
\item Risks: Chat claims may be stale, superseded, or contradicted by current repo artifacts.
\item Required review: human review plus authority-order check
\item Recommended disposition: not\_reviewed
\end{itemize}

\subsection{PROMOTE-0007 - A large set of prompts was then generated to lock down architecture, det}

\begin{itemize}
\item Source conversation: app\_testx\_codehygiene
\item Source file: docs/archive/conversations/app\_testx\_codehygiene/dominium\_app\_testx\_codehygiene\_handoff\_\_human\_readable\_chat\_report.txt
\item Source heading: automatic decision extraction
\item Claim: A large set of prompts was then generated to lock down architecture, determinism, performance, schema governance, rendering, epistemic UI, sharding, interest sets, and fidelity projection. This became the Phase 1 hardening layer. Additional audit prompts were introduced to ensure consistency before proceeding into life, civilization, war, content, agents, tools, mods, and final long-term policy.
\item Claim type: conversation\_claim
\item Proposed target: not\_selected
\item Authority domain: planning
\item Current repo conflict: not\_checked
\item Why it may matter: The claim recurs as a decision-like statement in the archived package.
\item Risks: Chat claims may be stale, superseded, or contradicted by current repo artifacts.
\item Required review: human review plus authority-order check
\item Recommended disposition: not\_reviewed
\end{itemize}

\subsection{PROMOTE-0008 - Because the chat was huge, the user asked for a maximum-fidelity context}

\begin{itemize}
\item Source conversation: app\_testx\_codehygiene
\item Source file: docs/archive/conversations/app\_testx\_codehygiene/dominium\_app\_testx\_codehygiene\_handoff\_\_human\_readable\_chat\_report.txt
\item Source heading: automatic decision extraction
\item Claim: Because the chat was huge, the user asked for a maximum-fidelity context transfer packet. Then they asked for downloadable report files. Then they asked for an in-chat reader. Finally, they asked for this human-readable explanatory report. This final report is meant to let a future assistant or human understand the substance without re-reading the whole conversation.
\item Claim type: conversation\_claim
\item Proposed target: not\_selected
\item Authority domain: planning
\item Current repo conflict: not\_checked
\item Why it may matter: The claim recurs as a decision-like statement in the archived package.
\item Risks: Chat claims may be stale, superseded, or contradicted by current repo artifacts.
\item Required review: human review plus authority-order check
\item Recommended disposition: not\_reviewed
\end{itemize}

\subsection{PROMOTE-0009 - The central project is a deterministic universe engine/game. The user's}

\begin{itemize}
\item Source conversation: app\_testx\_codehygiene
\item Source file: docs/archive/conversations/app\_testx\_codehygiene/dominium\_app\_testx\_codehygiene\_handoff\_\_human\_readable\_chat\_report.txt
\item Source heading: automatic decision extraction
\item Claim: The central project is a deterministic universe engine/game. The user's ambition is not just a normal game but a world runtime that can represent everything from a single-room scenario to an AI-only universe to an MMO. The project is meant to support real-world defaults, arbitrary modding, procedural and defined content, strong epistemics, and long-term replayability.
\item Claim type: conversation\_claim
\item Proposed target: not\_selected
\item Authority domain: planning
\item Current repo conflict: not\_checked
\item Why it may matter: The claim recurs as a decision-like statement in the archived package.
\item Risks: Chat claims may be stale, superseded, or contradicted by current repo artifacts.
\item Required review: human review plus authority-order check
\item Recommended disposition: not\_reviewed
\end{itemize}

\subsection{PROMOTE-0010 - The most important thing to remember is this: **this chat is the archite}

\begin{itemize}
\item Source conversation: architecture\_codex\_prompts
\item Source file: docs/archive/conversations/architecture\_codex\_prompts/dominium\_domino\_architecture\_codex\_prompts\_\_human\_readable\_chat\_report.txt
\item Source heading: automatic decision extraction
\item Claim: The most important thing to remember is this: **this chat is the architectural and implementation-roadmap backbone for Dominium/Domino, but it is not an implementation log**. It should be preserved as a design/specification source and as a prompt library, while actual code and current facts must be verified separately.
\item Claim type: conversation\_claim
\item Proposed target: not\_selected
\item Authority domain: planning
\item Current repo conflict: not\_checked
\item Why it may matter: The claim recurs as a decision-like statement in the archived package.
\item Risks: Chat claims may be stale, superseded, or contradicted by current repo artifacts.
\item Required review: human review plus authority-order check
\item Recommended disposition: not\_reviewed
\end{itemize}

\subsection{PROMOTE-0011 - FACT: The user opened by pasting an "Extended Master Starter Prompt - Do}

\begin{itemize}
\item Source conversation: architecture\_codex\_prompts
\item Source file: docs/archive/conversations/architecture\_codex\_prompts/dominium\_domino\_architecture\_codex\_prompts\_\_human\_readable\_chat\_report.txt
\item Source heading: automatic decision extraction
\item Claim: FACT: The user opened by pasting an "Extended Master Starter Prompt - Dominium + Domino." That prompt established the role of the assistant as a senior engine architect and defined the main project philosophy. Domino was described as a deterministic, fixed-point, C89 engine portable across old and modern OS/architecture combinations. Dominium was described as the product suite built on top of Domino.
\item Claim type: conversation\_claim
\item Proposed target: not\_selected
\item Authority domain: planning
\item Current repo conflict: not\_checked
\item Why it may matter: The claim recurs as a decision-like statement in the archived package.
\item Risks: Chat claims may be stale, superseded, or contradicted by current repo artifacts.
\item Required review: human review plus authority-order check
\item Recommended disposition: not\_reviewed
\end{itemize}

\subsection{PROMOTE-0012 - FACT: The user asked, succinctly, how a player would build a complete ma}

\begin{itemize}
\item Source conversation: architecture\_codex\_prompts
\item Source file: docs/archive/conversations/architecture\_codex\_prompts/dominium\_domino\_architecture\_codex\_prompts\_\_human\_readable\_chat\_report.txt
\item Source heading: automatic decision extraction
\item Claim: FACT: The user asked, succinctly, how a player would build a complete machine workshop. The assistant described a progression: extract raw materials, process them into industrial materials, produce mechanical components, produce electrical components, establish power distribution, establish logistics, assemble machine frames and machines, add measurement/control tooling, then scale horizontally and vertically.
\item Claim type: conversation\_claim
\item Proposed target: not\_selected
\item Authority domain: planning
\item Current repo conflict: not\_checked
\item Why it may matter: The claim recurs as a decision-like statement in the archived package.
\item Risks: Chat claims may be stale, superseded, or contradicted by current repo artifacts.
\item Required review: human review plus authority-order check
\item Recommended disposition: not\_reviewed
\end{itemize}

\subsection{PROMOTE-0013 - The language baseline decision matters because it changes what old rende}

\begin{itemize}
\item Source conversation: architecture\_ui\_providers
\item Source file: docs/archive/conversations/architecture\_ui\_providers/dominium\_architecture\_ui\_providers\_robot\_os\_strategy\_\_01\_human\_readable\_report.md
\item Source heading: automatic decision extraction
\item Claim: The language baseline decision matters because it changes what old renderers and toolchains matter. C17/C++17 reduces the need to design around C89/C++98 constraints but does not remove the need for C-compatible ABI boundaries. The system floor decision similarly moves many older renderers and OS targets into research/back-port categories.
\item Claim type: conversation\_claim
\item Proposed target: not\_selected
\item Authority domain: planning
\item Current repo conflict: not\_checked
\item Why it may matter: The claim recurs as a decision-like statement in the archived package.
\item Risks: Chat claims may be stale, superseded, or contradicted by current repo artifacts.
\item Required review: human review plus authority-order check
\item Recommended disposition: not\_reviewed
\end{itemize}

\subsection{PROMOTE-0014 - The Robot OS and robotic seed-civilisation decisions are the most import}

\begin{itemize}
\item Source conversation: architecture\_ui\_providers
\item Source file: docs/archive/conversations/architecture\_ui\_providers/dominium\_architecture\_ui\_providers\_robot\_os\_strategy\_\_01\_human\_readable\_report.md
\item Source heading: automatic decision extraction
\item Claim: The Robot OS and robotic seed-civilisation decisions are the most important design decisions. They are explicit user statements, not merely assistant proposals. They should be formalized as product/game requirements. They affect UI, fog-of-war, spawning, respawn, automation, industry, tutorial/planner design, and MMO anti-cheat.
\item Claim type: conversation\_claim
\item Proposed target: not\_selected
\item Authority domain: planning
\item Current repo conflict: not\_checked
\item Why it may matter: The claim recurs as a decision-like statement in the archived package.
\item Risks: Chat claims may be stale, superseded, or contradicted by current repo artifacts.
\item Required review: human review plus authority-order check
\item Recommended disposition: not\_reviewed
\end{itemize}

\subsection{PROMOTE-0015 - The central tradeoff was speed versus purity. Raylib, SDL2, Lua, raygui,}

\begin{itemize}
\item Source conversation: architecture\_ui\_providers
\item Source file: docs/archive/conversations/architecture\_ui\_providers/dominium\_architecture\_ui\_providers\_robot\_os\_strategy\_\_01\_human\_readable\_report.md
\item Source heading: automatic decision extraction
\item Claim: The central tradeoff was speed versus purity. Raylib, SDL2, Lua, raygui, rlgl, rlsw, and related libraries can accelerate visible progress, but if they leak into game law, UI documents, saves, replays, or public contracts, they become architectural capture. The chosen compromise is to use them as providers behind Dominium-owned services and conformance tests.
\item Claim type: conversation\_claim
\item Proposed target: not\_selected
\item Authority domain: planning
\item Current repo conflict: not\_checked
\item Why it may matter: The claim recurs as a decision-like statement in the archived package.
\item Risks: Chat claims may be stale, superseded, or contradicted by current repo artifacts.
\item Required review: human review plus authority-order check
\item Recommended disposition: not\_reviewed
\end{itemize}

\subsection{PROMOTE-0016 - Plain-language limitation: this package is reliable as a preservation of}

\begin{itemize}
\item Source conversation: Build\_and\_Future\_Proofing
\item Source file: docs/archive/conversations/Build\_and\_Future\_Proofing/Dominium\_Build\_and\_Future\_Proofing\_Architecture\_\_01\_human\_readable\_report.md
\item Source heading: automatic decision extraction
\item Claim: Plain-language limitation: this package is reliable as a preservation of the current visible chat. It should not be treated as a complete archive of every Dominium discussion in other chats. Where the report uses project or repository context, it is labelled as such. The most important caveat is that many items are recommendations, not final user decisions.
\item Claim type: conversation\_claim
\item Proposed target: not\_selected
\item Authority domain: planning
\item Current repo conflict: not\_checked
\item Why it may matter: The claim recurs as a decision-like statement in the archived package.
\item Risks: Chat claims may be stale, superseded, or contradicted by current repo artifacts.
\item Required review: human review plus authority-order check
\item Recommended disposition: not\_reviewed
\end{itemize}

\subsection{PROMOTE-0017 - The final uploaded prompt requested a complete preservation package for}

\begin{itemize}
\item Source conversation: Build\_and\_Future\_Proofing
\item Source file: docs/archive/conversations/Build\_and\_Future\_Proofing/Dominium\_Build\_and\_Future\_Proofing\_Architecture\_\_01\_human\_readable\_report.md
\item Source heading: automatic decision extraction
\item Claim: The final uploaded prompt requested a complete preservation package for this chat: a human-readable explanation, structured registers, context transfer packet, spec sheet, aggregator packet, self-audit, and files/ZIP package. This response completes that export task and creates downloadable files for later reading and aggregation.
\item Claim type: conversation\_claim
\item Proposed target: not\_selected
\item Authority domain: planning
\item Current repo conflict: not\_checked
\item Why it may matter: The claim recurs as a decision-like statement in the archived package.
\item Risks: Chat claims may be stale, superseded, or contradicted by current repo artifacts.
\item Required review: human review plus authority-order check
\item Recommended disposition: not\_reviewed
\end{itemize}

\subsection{PROMOTE-0018 - The final explicit goal was preservation: create a maximum-fidelity pack}

\begin{itemize}
\item Source conversation: Build\_and\_Future\_Proofing
\item Source file: docs/archive/conversations/Build\_and\_Future\_Proofing/Dominium\_Build\_and\_Future\_Proofing\_Architecture\_\_01\_human\_readable\_report.md
\item Source heading: automatic decision extraction
\item Claim: The final explicit goal was preservation: create a maximum-fidelity package for this chat so the user can understand it later, ask questions in this chat, merge it with other old-chat reports, and eventually build a master project spec book.
\item Claim type: conversation\_claim
\item Proposed target: not\_selected
\item Authority domain: planning
\item Current repo conflict: not\_checked
\item Why it may matter: The claim recurs as a decision-like statement in the archived package.
\item Risks: Chat claims may be stale, superseded, or contradicted by current repo artifacts.
\item Required review: human review plus authority-order check
\item Recommended disposition: not\_reviewed
\end{itemize}

\subsection{PROMOTE-0019 - This produced a pattern: generate a large Codex/AIDE task, the user ran}

\begin{itemize}
\item Source conversation: canonical\_structure\_and\_framework
\item Source file: docs/archive/conversations/canonical\_structure\_and\_framework/dominium\_canonical\_architecture\_repo\_foundation\_\_01\_human\_readable\_report.md
\item Source heading: automatic decision extraction
\item Claim: This produced a pattern: generate a large Codex/AIDE task, the user ran it or reported a result, then the assistant evaluated what was truly fixed versus what was only validator/document churn. The user eventually demanded a one-shot "actual final cleanup" prompt because previous passes had sometimes added validators without moving directories. That prompt explicitly required real git mv routing, not just reports.
\item Claim type: conversation\_claim
\item Proposed target: not\_selected
\item Authority domain: planning
\item Current repo conflict: not\_checked
\item Why it may matter: The claim recurs as a decision-like statement in the archived package.
\item Risks: Chat claims may be stale, superseded, or contradicted by current repo artifacts.
\item Required review: human review plus authority-order check
\item Recommended disposition: not\_reviewed
\end{itemize}

\subsection{PROMOTE-0020 - The final provider directory doctrine became service-first: `runtime/<se}

\begin{itemize}
\item Source conversation: canonical\_structure\_and\_framework
\item Source file: docs/archive/conversations/canonical\_structure\_and\_framework/dominium\_canonical\_architecture\_repo\_foundation\_\_01\_human\_readable\_report.md
\item Source heading: automatic decision extraction
\item Claim: The final provider directory doctrine became service-first: runtime/<service>/providers/<provider>/, not runtime/<vendor>/<service>/. Provider choices belong in release/profiles/ or content/profiles/, not app path names. External code belongs under external/upstream/ or external/vendor/, but the repo should choose one convention.
\item Claim type: conversation\_claim
\item Proposed target: not\_selected
\item Authority domain: planning
\item Current repo conflict: not\_checked
\item Why it may matter: The claim recurs as a decision-like statement in the archived package.
\item Risks: Chat claims may be stale, superseded, or contradicted by current repo artifacts.
\item Required review: human review plus authority-order check
\item Recommended disposition: not\_reviewed
\end{itemize}

\subsection{PROMOTE-0021 - Packs are distributable authored payloads. The preferred pack layout bec}

\begin{itemize}
\item Source conversation: canonical\_structure\_and\_framework
\item Source file: docs/archive/conversations/canonical\_structure\_and\_framework/dominium\_canonical\_architecture\_repo\_foundation\_\_01\_human\_readable\_report.md
\item Source heading: automatic decision extraction
\item Claim: Packs are distributable authored payloads. The preferred pack layout became category-based, such as content/packs/<category>/<pack\_id>/. Pack IDs must not depend on paths. Saves and replays must use stable IDs, schema versions, canonical serialization, deterministic hashes, and explicit migration/refusal policies. Backward compatibility must be a compatibility corpus plus tests, not intention.
\item Claim type: conversation\_claim
\item Proposed target: not\_selected
\item Authority domain: planning
\item Current repo conflict: not\_checked
\item Why it may matter: The claim recurs as a decision-like statement in the archived package.
\item Risks: Chat claims may be stale, superseded, or contradicted by current repo artifacts.
\item Required review: human review plus authority-order check
\item Recommended disposition: not\_reviewed
\end{itemize}

\subsection{PROMOTE-0022 - This reinforced the core celestial-content rule: locations should exist}

\begin{itemize}
\item Source conversation: Chronology\_Celestial\_Systems
\item Source file: docs/archive/conversations/Chronology\_Celestial\_Systems/Dominium\_Chronology\_Celestial\_Systems\_\_human\_readable\_chat\_report.txt
\item Source heading: automatic decision extraction
\item Claim: This reinforced the core celestial-content rule: locations should exist explicitly when they change player decisions.
\item Claim type: conversation\_claim
\item Proposed target: not\_selected
\item Authority domain: planning
\item Current repo conflict: not\_checked
\item Why it may matter: The claim recurs as a decision-like statement in the archived package.
\item Risks: Chat claims may be stale, superseded, or contradicted by current repo artifacts.
\item Required review: human review plus authority-order check
\item Recommended disposition: not\_reviewed
\end{itemize}

\subsection{PROMOTE-0023 - This became the spatial foundation for later time/calendar decisions bec}

\begin{itemize}
\item Source conversation: Chronology\_Celestial\_Systems
\item Source file: docs/archive/conversations/Chronology\_Celestial\_Systems/Dominium\_Chronology\_Celestial\_Systems\_\_human\_readable\_chat\_report.txt
\item Source heading: automatic decision extraction
\item Claim: This became the spatial foundation for later time/calendar decisions because every major body or region could potentially have its own local time, calendar, or operational standard.
\item Claim type: conversation\_claim
\item Proposed target: not\_selected
\item Authority domain: planning
\item Current repo conflict: not\_checked
\item Why it may matter: The claim recurs as a decision-like statement in the archived package.
\item Risks: Chat claims may be stale, superseded, or contradicted by current repo artifacts.
\item Required review: human review plus authority-order check
\item Recommended disposition: not\_reviewed
\end{itemize}

\subsection{PROMOTE-0024 - The assistant formalized this as the Perfect Earth Calendar, later calle}

\begin{itemize}
\item Source conversation: Chronology\_Celestial\_Systems
\item Source file: docs/archive/conversations/Chronology\_Celestial\_Systems/Dominium\_Chronology\_Celestial\_Systems\_\_human\_readable\_chat\_report.txt
\item Source heading: automatic decision extraction
\item Claim: The assistant formalized this as the Perfect Earth Calendar, later called HPC-E. The final structure is clear, though the exact leap rule remains unresolved.
\item Claim type: conversation\_claim
\item Proposed target: not\_selected
\item Authority domain: planning
\item Current repo conflict: not\_checked
\item Why it may matter: The claim recurs as a decision-like statement in the archived package.
\item Risks: Chat claims may be stale, superseded, or contradicted by current repo artifacts.
\item Required review: human review plus authority-order check
\item Recommended disposition: not\_reviewed
\end{itemize}

\subsection{PROMOTE-0025 - Before relying on this chat as project authority, a future assistant or}

\begin{itemize}
\item Source conversation: development\_routes
\item Source file: docs/archive/conversations/development\_routes/dominium\_development\_routes\_\_human\_readable\_chat\_report.txt
\item Source heading: automatic decision extraction
\item Claim: Before relying on this chat as project authority, a future assistant or human reviewer must confirm whether Dominium is actually intended to be simulation-heavy, whether strict determinism matters, whether replay or multiplayer are goals, whether modding matters, and what the actual game loop is. Until then, the kernel-first plan should be treated as a strong provisional proposal, not a final specification.
\item Claim type: conversation\_claim
\item Proposed target: not\_selected
\item Authority domain: planning
\item Current repo conflict: not\_checked
\item Why it may matter: The claim recurs as a decision-like statement in the archived package.
\item Risks: Chat claims may be stale, superseded, or contradicted by current repo artifacts.
\item Required review: human review plus authority-order check
\item Recommended disposition: not\_reviewed
\end{itemize}

\subsection{PROMOTE-0026 - The first answer was assertive. It stated that Dominium must follow Rout}

\begin{itemize}
\item Source conversation: development\_routes
\item Source file: docs/archive/conversations/development\_routes/dominium\_development\_routes\_\_human\_readable\_chat\_report.txt
\item Source heading: automatic decision extraction
\item Claim: The first answer was assertive. It stated that Dominium must follow Route C and described the route as the only viable one for Dominium. Later parts of the chat corrected the status of that claim. FACT: the assistant made that recommendation. UNCERTAIN / UNVERIFIED: the recommendation was not validated with external sources in the chat and was not explicitly accepted by the user.
\item Claim type: conversation\_claim
\item Proposed target: not\_selected
\item Authority domain: planning
\item Current repo conflict: not\_checked
\item Why it may matter: The claim recurs as a decision-like statement in the archived package.
\item Risks: Chat claims may be stale, superseded, or contradicted by current repo artifacts.
\item Required review: human review plus authority-order check
\item Recommended disposition: not\_reviewed
\end{itemize}

\subsection{PROMOTE-0027 - The initial technical answer proposed a layered stack. It started with m}

\begin{itemize}
\item Source conversation: development\_routes
\item Source file: docs/archive/conversations/development\_routes/dominium\_development\_routes\_\_human\_readable\_chat\_report.txt
\item Source heading: automatic decision extraction
\item Claim: The initial technical answer proposed a layered stack. It started with mathematical and temporal primitives, then moved into the deterministic simulation kernel, then world state, systems, persistence and replay, tooling, presentation, content, and finally distribution, modding, and multiplayer.
\item Claim type: conversation\_claim
\item Proposed target: not\_selected
\item Authority domain: planning
\item Current repo conflict: not\_checked
\item Why it may matter: The claim recurs as a decision-like statement in the archived package.
\item Risks: Chat claims may be stale, superseded, or contradicted by current repo artifacts.
\item Required review: human review plus authority-order check
\item Recommended disposition: not\_reviewed
\end{itemize}

\subsection{PROMOTE-0028 - The key thing to remember: this chat produced the standards and prompts}

\begin{itemize}
\item Source conversation: documentation\_standards\_readme
\item Source file: docs/archive/conversations/documentation\_standards\_readme/documentation\_standards\_readme\_handoff\_\_human\_readable\_chat\_report.txt
\item Source heading: automatic decision extraction
\item Claim: The key thing to remember: this chat produced the standards and prompts for future work; it did not verify or change the project. The next assistant must inspect the actual repository before writing repo-specific docs, README content, platform claims, license claims, or subsystem descriptions.
\item Claim type: conversation\_claim
\item Proposed target: not\_selected
\item Authority domain: planning
\item Current repo conflict: not\_checked
\item Why it may matter: The claim recurs as a decision-like statement in the archived package.
\item Risks: Chat claims may be stale, superseded, or contradicted by current repo artifacts.
\item Required review: human review plus authority-order check
\item Recommended disposition: not\_reviewed
\end{itemize}

\subsection{PROMOTE-0029 - FACT: The chat established that public API contracts should live in he}

\begin{itemize}
\item Source conversation: documentation\_standards\_readme
\item Source file: docs/archive/conversations/documentation\_standards\_readme/documentation\_standards\_readme\_handoff\_\_human\_readable\_chat\_report.txt
\item Source heading: automatic decision extraction
\item Claim: FACT: The chat established that public API contracts should live in headers.
\item Claim type: conversation\_claim
\item Proposed target: not\_selected
\item Authority domain: planning
\item Current repo conflict: not\_checked
\item Why it may matter: The claim recurs as a decision-like statement in the archived package.
\item Risks: Chat claims may be stale, superseded, or contradicted by current repo artifacts.
\item Required review: human review plus authority-order check
\item Recommended disposition: not\_reviewed
\end{itemize}

\subsection{PROMOTE-0030 - FACT: Source files should document implementation rationale, non-obvio}

\begin{itemize}
\item Source conversation: documentation\_standards\_readme
\item Source file: docs/archive/conversations/documentation\_standards\_readme/documentation\_standards\_readme\_handoff\_\_human\_readable\_chat\_report.txt
\item Source heading: automatic decision extraction
\item Claim: FACT: Source files should document implementation rationale, non-obvious invariants, and hazards.
\item Claim type: conversation\_claim
\item Proposed target: not\_selected
\item Authority domain: planning
\item Current repo conflict: not\_checked
\item Why it may matter: The claim recurs as a decision-like statement in the archived package.
\item Risks: Chat claims may be stale, superseded, or contradicted by current repo artifacts.
\item Required review: human review plus authority-order check
\item Recommended disposition: not\_reviewed
\end{itemize}

\subsection{PROMOTE-0031 - During this stage, the work was still broadly about "the project spec."}

\begin{itemize}
\item Source conversation: Dominium\_Architecture\_I
\item Source file: docs/archive/conversations/Dominium\_Architecture\_I/Dominium\_Architecture\_I\_\_human\_readable\_chat\_report.txt
\item Source heading: automatic decision extraction
\item Claim: During this stage, the work was still broadly about "the project spec." The user wanted a comprehensive description of the game and its technical systems. The assistant produced large design documents, but many of those earlier contents are not visible in the retained transcript. This is important because the final report package marks those earlier outputs as referenced but not fully recovered.
\item Claim type: conversation\_claim
\item Proposed target: not\_selected
\item Authority domain: planning
\item Current repo conflict: not\_checked
\item Why it may matter: The claim recurs as a decision-like statement in the archived package.
\item Risks: Chat claims may be stale, superseded, or contradicted by current repo artifacts.
\item Required review: human review plus authority-order check
\item Recommended disposition: not\_reviewed
\end{itemize}

\subsection{PROMOTE-0032 - The user also made a clear versioning decision. **FACT:** The user wante}

\begin{itemize}
\item Source conversation: Dominium\_Architecture\_I
\item Source file: docs/archive/conversations/Dominium\_Architecture\_I/Dominium\_Architecture\_I\_\_human\_readable\_chat\_report.txt
\item Source heading: automatic decision extraction
\item Claim: The user also made a clear versioning decision. **FACT:** The user wanted **major.minor.patch semantic versioning** for all components/packages and for the complete game package. The user explicitly did **not** want build numbers or build dates in filenames. Instead, those details should live in metadata. This became part of the future build/package design.
\item Claim type: conversation\_claim
\item Proposed target: not\_selected
\item Authority domain: planning
\item Current repo conflict: not\_checked
\item Why it may matter: The claim recurs as a decision-like statement in the archived package.
\item Risks: Chat claims may be stale, superseded, or contradicted by current repo artifacts.
\item Required review: human review plus authority-order check
\item Recommended disposition: not\_reviewed
\end{itemize}

\subsection{PROMOTE-0033 - The key change came when the user decided that Codex could not read the}

\begin{itemize}
\item Source conversation: Dominium\_Architecture\_I
\item Source file: docs/archive/conversations/Dominium\_Architecture\_I/Dominium\_Architecture\_I\_\_human\_readable\_chat\_report.txt
\item Source heading: automatic decision extraction
\item Claim: The key change came when the user decided that Codex could not read the entire conversation at once. **FACT:** The user said they would copy and paste the transcript with timestamps and would not edit it manually, but because Codex could not read the whole context at once, they would not bother asking Codex to compile a full book directly from the whole chat.
\item Claim type: conversation\_claim
\item Proposed target: not\_selected
\item Authority domain: planning
\item Current repo conflict: not\_checked
\item Why it may matter: The claim recurs as a decision-like statement in the archived package.
\item Risks: Chat claims may be stale, superseded, or contradicted by current repo artifacts.
\item Required review: human review plus authority-order check
\item Recommended disposition: not\_reviewed
\end{itemize}

\subsection{PROMOTE-0034 - The last part of the chat was handoff extraction. The user asked for dis}

\begin{itemize}
\item Source conversation: Dominium\_Architecture\_II
\item Source file: docs/archive/conversations/Dominium\_Architecture\_II/Dominium\_Architecture\_II\_\_human\_readable\_chat\_report.txt
\item Source heading: automatic decision extraction
\item Claim: The last part of the chat was handoff extraction. The user asked for discovery inventory, structured registers, a context transfer packet, and then a downloadable package. The generated package captured the chat's workstreams, decisions, tasks, constraints, open questions, artifacts, risks, and verification items. This current report is a plain-language explanation of that substance.
\item Claim type: conversation\_claim
\item Proposed target: not\_selected
\item Authority domain: planning
\item Current repo conflict: not\_checked
\item Why it may matter: The claim recurs as a decision-like statement in the archived package.
\item Risks: Chat claims may be stale, superseded, or contradicted by current repo artifacts.
\item Required review: human review plus authority-order check
\item Recommended disposition: not\_reviewed
\end{itemize}

\subsection{PROMOTE-0035 - The entire conversation is anchored by the requirement that Dominium mus}

\begin{itemize}
\item Source conversation: Dominium\_Architecture\_II
\item Source file: docs/archive/conversations/Dominium\_Architecture\_II/Dominium\_Architecture\_II\_\_human\_readable\_chat\_report.txt
\item Source heading: automatic decision extraction
\item Claim: The entire conversation is anchored by the requirement that Dominium must run deterministically. This means the same initial state and same inputs must produce the same state across machines, compilers, operating systems, and eventually retro and modern platforms.
\item Claim type: conversation\_claim
\item Proposed target: not\_selected
\item Authority domain: planning
\item Current repo conflict: not\_checked
\item Why it may matter: The claim recurs as a decision-like statement in the archived package.
\item Risks: Chat claims may be stale, superseded, or contradicted by current repo artifacts.
\item Required review: human review plus authority-order check
\item Recommended disposition: not\_reviewed
\end{itemize}

\subsection{PROMOTE-0036 - The chat discussed fixed tick phases, virtual lanes, merge order, comman}

\begin{itemize}
\item Source conversation: Dominium\_Architecture\_II
\item Source file: docs/archive/conversations/Dominium\_Architecture\_II/Dominium\_Architecture\_II\_\_human\_readable\_chat\_report.txt
\item Source heading: automatic decision extraction
\item Claim: The chat discussed fixed tick phases, virtual lanes, merge order, command buffers, event queues, and deterministic serialization. The simulation must not depend on render FPS, wall-clock time, OS scheduler behaviour, GPU state, or floating-point quirks. This explains many later decisions: C89 core, integer/fixed-point math, no sim floats, no platform headers in simulation, and renderer as pure observer.
\item Claim type: conversation\_claim
\item Proposed target: not\_selected
\item Authority domain: planning
\item Current repo conflict: not\_checked
\item Why it may matter: The claim recurs as a decision-like statement in the archived package.
\item Risks: Chat claims may be stale, superseded, or contradicted by current repo artifacts.
\item Required review: human review plus authority-order check
\item Recommended disposition: not\_reviewed
\end{itemize}

\subsection{PROMOTE-0037 - The view system then expanded beyond the launcher. **FACT:** The user wa}

\begin{itemize}
\item Source conversation: Dominium\_Architecture\_III
\item Source file: docs/archive/conversations/Dominium\_Architecture\_III/Dominium\_Architecture\_III\_\_human\_readable\_chat\_report.txt
\item Source heading: automatic decision extraction
\item Claim: The view system then expanded beyond the launcher. **FACT:** The user wants instant zoom to any scale, instant switching between top-down 2D and first-person 3D, and arbitrary cameras for free cam, map viewing, content creation, HUD cameras, CCTV, overlays, windows, and offscreen targets. These are client/render-side concerns, not simulation concerns.
\item Claim type: conversation\_claim
\item Proposed target: not\_selected
\item Authority domain: planning
\item Current repo conflict: not\_checked
\item Why it may matter: The claim recurs as a decision-like statement in the archived package.
\item Risks: Chat claims may be stale, superseded, or contradicted by current repo artifacts.
\item Required review: human review plus authority-order check
\item Recommended disposition: not\_reviewed
\end{itemize}

\subsection{PROMOTE-0038 - This context matters because the rest of the conversation happened insid}

\begin{itemize}
\item Source conversation: Dominium\_Architecture\_III
\item Source file: docs/archive/conversations/Dominium\_Architecture\_III/Dominium\_Architecture\_III\_\_human\_readable\_chat\_report.txt
\item Source heading: automatic decision extraction
\item Claim: This context matters because the rest of the conversation happened inside those constraints. The launcher could not become a dumping ground for OS-specific behavior, sim mutation, renderer decisions, or nondeterministic game-state changes. It had to fit the project's deterministic layering.
\item Claim type: conversation\_claim
\item Proposed target: not\_selected
\item Authority domain: planning
\item Current repo conflict: not\_checked
\item Why it may matter: The claim recurs as a decision-like statement in the archived package.
\item Risks: Chat claims may be stale, superseded, or contradicted by current repo artifacts.
\item Required review: human review plus authority-order check
\item Recommended disposition: not\_reviewed
\end{itemize}

\subsection{PROMOTE-0039 - The user then uploaded dominium.7z, saying it was the git repository a}

\begin{itemize}
\item Source conversation: Dominium\_Architecture\_III
\item Source file: docs/archive/conversations/Dominium\_Architecture\_III/Dominium\_Architecture\_III\_\_human\_readable\_chat\_report.txt
\item Source heading: automatic decision extraction
\item Claim: The user then uploaded dominium.7z, saying it was the git repository as of that moment. **FACT:** The prior assistant said it could not open .7z archives in that environment. That means the actual repository state remained unverified at that stage. This became important later, because many implementation ideas were discussed without confirmed access to the repo contents.
\item Claim type: conversation\_claim
\item Proposed target: not\_selected
\item Authority domain: planning
\item Current repo conflict: not\_checked
\item Why it may matter: The claim recurs as a decision-like statement in the archived package.
\item Risks: Chat claims may be stale, superseded, or contradicted by current repo artifacts.
\item Required review: human review plus authority-order check
\item Recommended disposition: not\_reviewed
\end{itemize}

\subsection{PROMOTE-0040 - The reusable engine layer was named **Domino**. This became one of the m}

\begin{itemize}
\item Source conversation: Dominium\_Architecture\_IV
\item Source file: docs/archive/conversations/Dominium\_Architecture\_IV/Dominium\_Architecture\_IV\_\_human\_readable\_chat\_report.txt
\item Source heading: automatic decision extraction
\item Claim: The reusable engine layer was named **Domino**. This became one of the most important naming and architectural decisions in the chat. Domino is the engine/core/platform/sim/mod layer, reusable for other games and projects. Dominium is the game/product layer built on it.
\item Claim type: conversation\_claim
\item Proposed target: not\_selected
\item Authority domain: planning
\item Current repo conflict: not\_checked
\item Why it may matter: The claim recurs as a decision-like statement in the archived package.
\item Risks: Chat claims may be stale, superseded, or contradicted by current repo artifacts.
\item Required review: human review plus authority-order check
\item Recommended disposition: not\_reviewed
\end{itemize}

\subsection{PROMOTE-0041 - Blueprints are plans or diffs: place element, remove element, modify ter}

\begin{itemize}
\item Source conversation: Dominium\_Architecture\_IV
\item Source file: docs/archive/conversations/Dominium\_Architecture\_IV/Dominium\_Architecture\_IV\_\_human\_readable\_chat\_report.txt
\item Source heading: automatic decision extraction
\item Claim: Blueprints are plans or diffs: place element, remove element, modify terrain, place machine, connect network, etc. When accepted, they generate jobs. Manual player actions and queued work for humans, robots, or other agents should use the same job pipeline. This matters because it keeps replay, AI, and automation consistent.
\item Claim type: conversation\_claim
\item Proposed target: not\_selected
\item Authority domain: planning
\item Current repo conflict: not\_checked
\item Why it may matter: The claim recurs as a decision-like statement in the archived package.
\item Risks: Chat claims may be stale, superseded, or contradicted by current repo artifacts.
\item Required review: human review plus authority-order check
\item Recommended disposition: not\_reviewed
\end{itemize}

\subsection{PROMOTE-0042 - The user then asked how to structure Codex prompts. The final high-level}

\begin{itemize}
\item Source conversation: Dominium\_Architecture\_IV
\item Source file: docs/archive/conversations/Dominium\_Architecture\_IV/Dominium\_Architecture\_IV\_\_human\_readable\_chat\_report.txt
\item Source heading: automatic decision extraction
\item Claim: The user then asked how to structure Codex prompts. The final high-level roadmap became:
\item Claim type: conversation\_claim
\item Proposed target: not\_selected
\item Authority domain: planning
\item Current repo conflict: not\_checked
\item Why it may matter: The claim recurs as a decision-like statement in the archived package.
\item Risks: Chat claims may be stale, superseded, or contradicted by current repo artifacts.
\item Required review: human review plus authority-order check
\item Recommended disposition: not\_reviewed
\end{itemize}

\subsection{PROMOTE-0043 - This chat was about preserving and re-grounding Dominium's architecture}

\begin{itemize}
\item Source conversation: Dominium\_Complete\_Conversation
\item Source file: docs/archive/conversations/Dominium\_Complete\_Conversation/Accompanying\_Report/Dominium\_Canon\_Portability\_Audit\_\_01\_human\_readable\_report.md
\item Source heading: automatic decision extraction
\item Claim: This chat was about preserving and re-grounding Dominium's architecture around a very old constitutional contract, then checking whether the current GitHub repository still follows that contract, and finally broadening the question into a platform-engineering doctrine for portability, modularity, extensibility, reuse, future-proofing, and long-term compatibility.
\item Claim type: conversation\_claim
\item Proposed target: not\_selected
\item Authority domain: planning
\item Current repo conflict: not\_checked
\item Why it may matter: The claim recurs as a decision-like statement in the archived package.
\item Risks: Chat claims may be stale, superseded, or contradicted by current repo artifacts.
\item Required review: human review plus authority-order check
\item Recommended disposition: not\_reviewed
\end{itemize}

\subsection{PROMOTE-0044 - The glossary bound terms such as Authority, Law, Process, Lens, SessionS}

\begin{itemize}
\item Source conversation: Dominium\_Complete\_Conversation
\item Source file: docs/archive/conversations/Dominium\_Complete\_Conversation/Accompanying\_Report/Dominium\_Canon\_Portability\_Audit\_\_01\_human\_readable\_report.md
\item Source heading: automatic decision extraction
\item Claim: The glossary bound terms such as Authority, Law, Process, Lens, SessionSpec, UniverseIdentity, Macro Capsule, SRZ, RepoX, TestX, AuditX, CompatX, SecureX, and related terms. It matters because future assistants must not use sloppy synonyms like "mode" where the canon requires ExperienceProfile or LawProfile. This topic directly supports modularity because stable vocabulary is part of stable architecture.
\item Claim type: conversation\_claim
\item Proposed target: not\_selected
\item Authority domain: planning
\item Current repo conflict: not\_checked
\item Why it may matter: The claim recurs as a decision-like statement in the archived package.
\item Risks: Chat claims may be stale, superseded, or contradicted by current repo artifacts.
\item Required review: human review plus authority-order check
\item Recommended disposition: not\_reviewed
\end{itemize}

\subsection{PROMOTE-0045 - Uncertainty: during this preservation pass, a potential inconsistency ap}

\begin{itemize}
\item Source conversation: Dominium\_Complete\_Conversation
\item Source file: docs/archive/conversations/Dominium\_Complete\_Conversation/Accompanying\_Report/Dominium\_Canon\_Portability\_Audit\_\_01\_human\_readable\_report.md
\item Source heading: automatic decision extraction
\item Claim: Uncertainty: during this preservation pass, a potential inconsistency appeared: some docs claim product roots moved under apps/, while an earlier inspected code path under root client/ was available and apps/client was not fetched successfully. This must be verified against the current physical tree.
\item Claim type: conversation\_claim
\item Proposed target: not\_selected
\item Authority domain: planning
\item Current repo conflict: not\_checked
\item Why it may matter: The claim recurs as a decision-like statement in the archived package.
\item Risks: Chat claims may be stale, superseded, or contradicted by current repo artifacts.
\item Required review: human review plus authority-order check
\item Recommended disposition: not\_reviewed
\end{itemize}

\subsection{PROMOTE-0046 - INFERENCE:** The user accepted this direction in practice, because the n}

\begin{itemize}
\item Source conversation: dominium\_domino\_codex\_planning
\item Source file: docs/archive/conversations/dominium\_domino\_codex\_planning/dominium\_domino\_codex\_planning\_\_human\_readable\_chat\_report.txt
\item Source heading: automatic decision extraction
\item Claim: INFERENCE:** The user accepted this direction in practice, because the next request was to generate Codex prompts to implement each step fully. The user did not explicitly say, "I formally approve the Milestone-0 plan," but they proceeded to ask for prompts based on it.
\item Claim type: conversation\_claim
\item Proposed target: not\_selected
\item Authority domain: planning
\item Current repo conflict: not\_checked
\item Why it may matter: The claim recurs as a decision-like statement in the archived package.
\item Risks: Chat claims may be stale, superseded, or contradicted by current repo artifacts.
\item Required review: human review plus authority-order check
\item Recommended disposition: not\_reviewed
\end{itemize}

\subsection{PROMOTE-0047 - The user then asked for recommendations. The assistant recommended a cle}

\begin{itemize}
\item Source conversation: dominium\_domino\_codex\_planning
\item Source file: docs/archive/conversations/dominium\_domino\_codex\_planning/dominium\_domino\_codex\_planning\_\_human\_readable\_chat\_report.txt
\item Source heading: automatic decision extraction
\item Claim: The user then asked for recommendations. The assistant recommended a clean startup policy: explicit flags first, no environment/config override in v1, terminal detection through dsys, generic AUTO behavior, product-specific headless handling for the game, build-time GUI/TUI capabilities, and structural error codes for fallback. The user accepted this enough to request a "one big Codex prompt" to implement it.
\item Claim type: conversation\_claim
\item Proposed target: not\_selected
\item Authority domain: planning
\item Current repo conflict: not\_checked
\item Why it may matter: The claim recurs as a decision-like statement in the archived package.
\item Risks: Chat claims may be stale, superseded, or contradicted by current repo artifacts.
\item Required review: human review plus authority-order check
\item Recommended disposition: not\_reviewed
\end{itemize}

\subsection{PROMOTE-0048 - However, the later context transfer packet treated that inspection claim}

\begin{itemize}
\item Source conversation: dominium\_domino\_codex\_planning
\item Source file: docs/archive/conversations/dominium\_domino\_codex\_planning/dominium\_domino\_codex\_planning\_\_human\_readable\_chat\_report.txt
\item Source heading: automatic decision extraction
\item Claim: However, the later context transfer packet treated that inspection claim as **UNCERTAIN / UNVERIFIED**. This was careful and correct: within the visible chat, there was no reliable tool trace proving the repo had actually been inspected. Therefore, future assistants must verify dominium.zip or the active checkout before relying on the platform/render/API answer.
\item Claim type: conversation\_claim
\item Proposed target: not\_selected
\item Authority domain: planning
\item Current repo conflict: not\_checked
\item Why it may matter: The claim recurs as a decision-like statement in the archived package.
\item Risks: Chat claims may be stale, superseded, or contradicted by current repo artifacts.
\item Required review: human review plus authority-order check
\item Recommended disposition: not\_reviewed
\end{itemize}

\subsection{PROMOTE-0049 - A future assistant must understand that this chat is not simply about UI}

\begin{itemize}
\item Source conversation: dominium\_full\_conversation
\item Source file: docs/archive/conversations/dominium\_full\_conversation/companion\_report/dominium\_full\_conversation\_companion\_report\_\_01\_complete\_human\_report.md
\item Source heading: automatic decision extraction
\item Claim: A future assistant must understand that this chat is not simply about UI. It records the architecture for how Dominium should build itself: through contracts, commands, services, providers, documents, patches, modules, packs, apps, workspaces, diagnostics, evidence, AIDE, Codex, and eventually Workbench.
\item Claim type: conversation\_claim
\item Proposed target: not\_selected
\item Authority domain: planning
\item Current repo conflict: not\_checked
\item Why it may matter: The claim recurs as a decision-like statement in the archived package.
\item Risks: Chat claims may be stale, superseded, or contradicted by current repo artifacts.
\item Required review: human review plus authority-order check
\item Recommended disposition: not\_reviewed
\end{itemize}

\subsection{PROMOTE-0050 - The user then redirected the entire plan. They argued that native OS wid}

\begin{itemize}
\item Source conversation: dominium\_full\_conversation
\item Source file: docs/archive/conversations/dominium\_full\_conversation/companion\_report/dominium\_full\_conversation\_companion\_report\_\_01\_complete\_human\_report.md
\item Source heading: automatic decision extraction
\item Claim: The user then redirected the entire plan. They argued that native OS widget GUI tools already exist through Visual Studio, Xcode, etc. The project needed a cross-platform rendered tool environment that uses the same CLI, TUI, and rendered GUI systems as the client. The old UI Editor and Tool Editor were good exploratory ideas but bad final products. This produced the Dominium Workbench concept.
\item Claim type: conversation\_claim
\item Proposed target: not\_selected
\item Authority domain: planning
\item Current repo conflict: not\_checked
\item Why it may matter: The claim recurs as a decision-like statement in the archived package.
\item Risks: Chat claims may be stale, superseded, or contradicted by current repo artifacts.
\item Required review: human review plus authority-order check
\item Recommended disposition: not\_reviewed
\end{itemize}

\subsection{PROMOTE-0051 - The discussion then moved into AIDE. The user wanted to automate as much}

\begin{itemize}
\item Source conversation: dominium\_full\_conversation
\item Source file: docs/archive/conversations/dominium\_full\_conversation/companion\_report/dominium\_full\_conversation\_companion\_report\_\_01\_complete\_human\_report.md
\item Source heading: automatic decision extraction
\item Claim: The discussion then moved into AIDE. The user wanted to automate as much as possible through Codex and AIDE. We developed a task operating model: AIDE creates WorkUnits, tracks attempts, blockers, evidence, repairs, resumes, checkpoints, and promotion decisions. Codex executes bounded tasks. Development can continue with classified partials and repairable blockers; promotion to main requires evidence.
\item Claim type: conversation\_claim
\item Proposed target: not\_selected
\item Authority domain: planning
\item Current repo conflict: not\_checked
\item Why it may matter: The claim recurs as a decision-like statement in the archived package.
\item Risks: Chat claims may be stale, superseded, or contradicted by current repo artifacts.
\item Required review: human review plus authority-order check
\item Recommended disposition: not\_reviewed
\end{itemize}

\subsection{PROMOTE-0052 - Near the end, the chat shifted into preservation mode. The user asked fo}

\begin{itemize}
\item Source conversation: dominium\_setup
\item Source file: docs/archive/conversations/dominium\_setup/dominium\_setup\_handoff\_\_human\_readable\_chat\_report.txt
\item Source heading: automatic decision extraction
\item Claim: Near the end, the chat shifted into preservation mode. The user asked for a maximum-fidelity context transfer packet, then for a downloadable report package, then for an in-chat reader version of that package. Those outputs preserved the decisions, tasks, constraints, artifacts, rejected options, and verification queue for future assistants or aggregation into a larger Project Spec Book.
\item Claim type: conversation\_claim
\item Proposed target: not\_selected
\item Authority domain: planning
\item Current repo conflict: not\_checked
\item Why it may matter: The claim recurs as a decision-like statement in the archived package.
\item Risks: Chat claims may be stale, superseded, or contradicted by current repo artifacts.
\item Required review: human review plus authority-order check
\item Recommended disposition: not\_reviewed
\end{itemize}

\subsection{PROMOTE-0053 - These directory trees were not final, but they mattered because they exp}

\begin{itemize}
\item Source conversation: dominium\_setup
\item Source file: docs/archive/conversations/dominium\_setup/dominium\_setup\_handoff\_\_human\_readable\_chat\_report.txt
\item Source heading: automatic decision extraction
\item Claim: These directory trees were not final, but they mattered because they exposed the user's preference for simple, shallow, logical structures and native IDE workflows that do not become the source of truth.
\item Claim type: conversation\_claim
\item Proposed target: not\_selected
\item Authority domain: planning
\item Current repo conflict: not\_checked
\item Why it may matter: The claim recurs as a decision-like statement in the archived package.
\item Risks: Chat claims may be stale, superseded, or contradicted by current repo artifacts.
\item Required review: human review plus authority-order check
\item Recommended disposition: not\_reviewed
\end{itemize}

\subsection{PROMOTE-0054 - The user later asked where to create Visual Studio and Xcode apps after}

\begin{itemize}
\item Source conversation: dominium\_setup
\item Source file: docs/archive/conversations/dominium\_setup/dominium\_setup\_handoff\_\_human\_readable\_chat\_report.txt
\item Source heading: automatic decision extraction
\item Claim: The user later asked where to create Visual Studio and Xcode apps after opening the repo as a folder. The assistant advised that Visual Studio and Xcode projects should be generated through CMake, not hand-authored or treated as authoritative source files. This decision became consistent with later Codex prompts.
\item Claim type: conversation\_claim
\item Proposed target: not\_selected
\item Authority domain: planning
\item Current repo conflict: not\_checked
\item Why it may matter: The claim recurs as a decision-like statement in the archived package.
\item Risks: Chat claims may be stale, superseded, or contradicted by current repo artifacts.
\item Required review: human review plus authority-order check
\item Recommended disposition: not\_reviewed
\end{itemize}

\subsection{PROMOTE-0055 - The original UI Editor / Tool Editor plan is now **superseded as a final}

\begin{itemize}
\item Source conversation: Domino\_Dominium\_Workbench
\item Source file: docs/archive/conversations/Domino\_Dominium\_Workbench/dominium\_workbench\_full\_conversation\_companion\_\_human\_readable\_detailed\_summary\_report.md
\item Source heading: automatic decision extraction
\item Claim: The original UI Editor / Tool Editor plan is now **superseded as a final product**. It is not lost: its useful parts become Workbench capabilities, especially Interface Studio, UI/HUD Sandbox, Theme Laboratory, TUI Studio, Rendered GUI Studio, Preview Matrix, validation panels, import/export tools, and document-patch workflows.
\item Claim type: conversation\_claim
\item Proposed target: not\_selected
\item Authority domain: planning
\item Current repo conflict: not\_checked
\item Why it may matter: The claim recurs as a decision-like statement in the archived package.
\item Risks: Chat claims may be stale, superseded, or contradicted by current repo artifacts.
\item Required review: human review plus authority-order check
\item Recommended disposition: not\_reviewed
\end{itemize}

\subsection{PROMOTE-0056 - The central decision is that Workbench must not become semantic authorit}

\begin{itemize}
\item Source conversation: Domino\_Dominium\_Workbench
\item Source file: docs/archive/conversations/Domino\_Dominium\_Workbench/dominium\_workbench\_full\_conversation\_companion\_\_human\_readable\_detailed\_summary\_report.md
\item Source heading: automatic decision extraction
\item Claim: The central decision is that Workbench must not become semantic authority. Workbench is a surface over contracts, commands, services, documents, patches, providers, packs, modules, apps, artifacts, diagnostics, evidence, tests, and replay/proof. It is the richest human/agent interface over the system, but it must not bypass command, capability, refusal, validation, document, patch, or proof law.
\item Claim type: conversation\_claim
\item Proposed target: not\_selected
\item Authority domain: planning
\item Current repo conflict: not\_checked
\item Why it may matter: The claim recurs as a decision-like statement in the archived package.
\item Risks: Chat claims may be stale, superseded, or contradicted by current repo artifacts.
\item Required review: human review plus authority-order check
\item Recommended disposition: not\_reviewed
\end{itemize}

\subsection{PROMOTE-0057 - The second central decision is that CLI, TUI, rendered GUI, OS-native GU}

\begin{itemize}
\item Source conversation: Domino\_Dominium\_Workbench
\item Source file: docs/archive/conversations/Domino\_Dominium\_Workbench/dominium\_workbench\_full\_conversation\_companion\_\_human\_readable\_detailed\_summary\_report.md
\item Source heading: automatic decision extraction
\item Claim: The second central decision is that CLI, TUI, rendered GUI, OS-native GUI, and headless automation must be **projections of the same command/result/refusal/document/view spine**. A GUI button, a TUI hotkey, a CLI command, a headless JSON job, and a future OS-native widget action should route through the same command/service path and produce the same typed result, diagnostics, refusals, logs, and evidence.
\item Claim type: conversation\_claim
\item Proposed target: not\_selected
\item Authority domain: planning
\item Current repo conflict: not\_checked
\item Why it may matter: The claim recurs as a decision-like statement in the archived package.
\item Risks: Chat claims may be stale, superseded, or contradicted by current repo artifacts.
\item Required review: human review plus authority-order check
\item Recommended disposition: not\_reviewed
\end{itemize}

\subsection{PROMOTE-0058 - FACT:** The Domino engine must be ISO C89. The Dominium UI/tools layer m}

\begin{itemize}
\item Source conversation: domino\_engine\_refactor\_prompts
\item Source file: docs/archive/conversations/domino\_engine\_refactor\_prompts/dominium\_domino\_engine\_refactor\_prompts\_\_human\_readable\_chat\_report.txt
\item Source heading: automatic decision extraction
\item Claim: FACT:** The Domino engine must be ISO C89. The Dominium UI/tools layer must be portable C++98. Engine code must not use OS APIs. Determinism matters, so floats are disallowed in deterministic paths. The engine must avoid platform-dependent behavior and hardcoded game semantics. Everything should be content-driven.
\item Claim type: conversation\_claim
\item Proposed target: not\_selected
\item Authority domain: planning
\item Current repo conflict: not\_checked
\item Why it may matter: The claim recurs as a decision-like statement in the archived package.
\item Risks: Chat claims may be stale, superseded, or contradicted by current repo artifacts.
\item Required review: human review plus authority-order check
\item Recommended disposition: not\_reviewed
\end{itemize}

\subsection{PROMOTE-0059 - This became one of the central architectural ideas of the chat. It makes}

\begin{itemize}
\item Source conversation: domino\_engine\_refactor\_prompts
\item Source file: docs/archive/conversations/domino\_engine\_refactor\_prompts/dominium\_domino\_engine\_refactor\_prompts\_\_human\_readable\_chat\_report.txt
\item Source heading: automatic decision extraction
\item Claim: This became one of the central architectural ideas of the chat. It makes behavior extensible without giving minds direct write access to the world. A mod can add new sensors, observations, intents, capabilities, or actions, but state mutation still passes through a deterministic pipeline.
\item Claim type: conversation\_claim
\item Proposed target: not\_selected
\item Authority domain: planning
\item Current repo conflict: not\_checked
\item Why it may matter: The claim recurs as a decision-like statement in the archived package.
\item Risks: Chat claims may be stale, superseded, or contradicted by current repo artifacts.
\item Required review: human review plus authority-order check
\item Recommended disposition: not\_reviewed
\end{itemize}

\subsection{PROMOTE-0060 - After the architecture and prompts were created, the user asked for a ma}

\begin{itemize}
\item Source conversation: domino\_engine\_refactor\_prompts
\item Source file: docs/archive/conversations/domino\_engine\_refactor\_prompts/dominium\_domino\_engine\_refactor\_prompts\_\_human\_readable\_chat\_report.txt
\item Source heading: automatic decision extraction
\item Claim: After the architecture and prompts were created, the user asked for a maximum-fidelity context transfer packet. The assistant produced one. The user then asked for a downloadable report package; the assistant generated package files and a ZIP. The user later asked for an in-chat reader version of that package, and finally requested this current human-readable explanatory report.
\item Claim type: conversation\_claim
\item Proposed target: not\_selected
\item Authority domain: planning
\item Current repo conflict: not\_checked
\item Why it may matter: The claim recurs as a decision-like statement in the archived package.
\item Risks: Chat claims may be stale, superseded, or contradicted by current repo artifacts.
\item Required review: human review plus authority-order check
\item Recommended disposition: not\_reviewed
\end{itemize}

\subsection{PROMOTE-0061 - This corrected the plan. The final sequencing became: first **Milestone}

\begin{itemize}
\item Source conversation: engine\_baseline\_architecture
\item Source file: docs/archive/conversations/engine\_baseline\_architecture/domino\_dominium\_engine\_baseline\_architecture\_\_01\_human\_readable\_report.md
\item Source heading: automatic decision extraction
\item Claim: This corrected the plan. The final sequencing became: first **Milestone 0: Make the baseline path honest**. That means fixing server/runtime circular import, CLI forwarding, session\_create -> session\_boot, missing time-anchor policy registry, and making the strict local playtest validator pass. Only after that should the builder/destruction lab begin.
\item Claim type: conversation\_claim
\item Proposed target: not\_selected
\item Authority domain: planning
\item Current repo conflict: not\_checked
\item Why it may matter: The claim recurs as a decision-like statement in the archived package.
\item Risks: Chat claims may be stale, superseded, or contradicted by current repo artifacts.
\item Required review: human review plus authority-order check
\item Recommended disposition: not\_reviewed
\end{itemize}

\subsection{PROMOTE-0062 - The final user action uploaded Pasted text.txt, a detailed instruction}

\begin{itemize}
\item Source conversation: engine\_baseline\_architecture
\item Source file: docs/archive/conversations/engine\_baseline\_architecture/domino\_dominium\_engine\_baseline\_architecture\_\_01\_human\_readable\_report.md
\item Source heading: automatic decision extraction
\item Claim: The final user action uploaded Pasted text.txt, a detailed instruction prompt requiring a full preservation report, structured registers, spec sheet, aggregator packet, self-audit, and downloadable file package for this chat. That request is the current task.
\item Claim type: conversation\_claim
\item Proposed target: not\_selected
\item Authority domain: planning
\item Current repo conflict: not\_checked
\item Why it may matter: The claim recurs as a decision-like statement in the archived package.
\item Risks: Chat claims may be stale, superseded, or contradicted by current repo artifacts.
\item Required review: human review plus authority-order check
\item Recommended disposition: not\_reviewed
\end{itemize}

\subsection{PROMOTE-0063 - The central architecture distinction was that Domino should be a reusabl}

\begin{itemize}
\item Source conversation: engine\_baseline\_architecture
\item Source file: docs/archive/conversations/engine\_baseline\_architecture/domino\_dominium\_engine\_baseline\_architecture\_\_01\_human\_readable\_report.md
\item Source heading: automatic decision extraction
\item Claim: The central architecture distinction was that Domino should be a reusable deterministic simulation substrate, while Dominium should be one game/product layer built on it. This came up immediately when the user asked for a total description of the project. It mattered because the user wants the code to be reusable for future games and possibly other engine projects.
\item Claim type: conversation\_claim
\item Proposed target: not\_selected
\item Authority domain: planning
\item Current repo conflict: not\_checked
\item Why it may matter: The claim recurs as a decision-like statement in the archived package.
\item Risks: Chat claims may be stale, superseded, or contradicted by current repo artifacts.
\item Required review: human review plus authority-order check
\item Recommended disposition: not\_reviewed
\end{itemize}

\subsection{PROMOTE-0064 - After the root skeleton improved, the assistant and user recognized that}

\begin{itemize}
\item Source conversation: Foundation\_Workbench\_Codex
\item Source file: docs/archive/conversations/Foundation\_Workbench\_Codex/Dominium\_Foundation\_Workbench\_Parallel\_Codex\_Preservation\_\_01\_human\_readable\_report.md
\item Source heading: automatic decision extraction
\item Claim: After the root skeleton improved, the assistant and user recognized that the deeper problem was no longer top-level directories. The problem became semantic duplication and governance: what is public, what is private, what is stable, what is provisional, what is generated, what is a fixture, what must stay compatible, what can be replaced, and what proof is required.
\item Claim type: conversation\_claim
\item Proposed target: not\_selected
\item Authority domain: planning
\item Current repo conflict: not\_checked
\item Why it may matter: The claim recurs as a decision-like statement in the archived package.
\item Risks: Chat claims may be stale, superseded, or contradicted by current repo artifacts.
\item Required review: human review plus authority-order check
\item Recommended disposition: not\_reviewed
\end{itemize}

\subsection{PROMOTE-0065 - The key conclusion was that Workbench is not the general module system;}

\begin{itemize}
\item Source conversation: Foundation\_Workbench\_Codex
\item Source file: docs/archive/conversations/Foundation\_Workbench\_Codex/Dominium\_Foundation\_Workbench\_Parallel\_Codex\_Preservation\_\_01\_human\_readable\_report.md
\item Source heading: automatic decision extraction
\item Claim: The key conclusion was that Workbench is not the general module system; it is one consumer of the module/command/service/provider/pack/artifact system. Workbench must not call private tools directly. It must route through registered commands and typed results, diagnostics, refusals, views, and evidence.
\item Claim type: conversation\_claim
\item Proposed target: not\_selected
\item Authority domain: planning
\item Current repo conflict: not\_checked
\item Why it may matter: The claim recurs as a decision-like statement in the archived package.
\item Risks: Chat claims may be stale, superseded, or contradicted by current repo artifacts.
\item Required review: human review plus authority-order check
\item Recommended disposition: not\_reviewed
\end{itemize}

\subsection{PROMOTE-0066 - The topic came up because the project needed a model that could support}

\begin{itemize}
\item Source conversation: Foundation\_Workbench\_Codex
\item Source file: docs/archive/conversations/Foundation\_Workbench\_Codex/Dominium\_Foundation\_Workbench\_Parallel\_Codex\_Preservation\_\_01\_human\_readable\_report.md
\item Source heading: automatic decision extraction
\item Claim: The topic came up because the project needed a model that could support world simulation, modding, tooling, Workbench, release, portability, and future games without repeating endless refactors. The conclusion was that stable contracts and semantic IDs must be separated from replaceable private implementation.
\item Claim type: conversation\_claim
\item Proposed target: not\_selected
\item Authority domain: planning
\item Current repo conflict: not\_checked
\item Why it may matter: The claim recurs as a decision-like statement in the archived package.
\item Risks: Chat claims may be stale, superseded, or contradicted by current repo artifacts.
\item Required review: human review plus authority-order check
\item Recommended disposition: not\_reviewed
\end{itemize}

\subsection{PROMOTE-0067 - Finally, the user uploaded a detailed preservation prompt and requested}

\begin{itemize}
\item Source conversation: Framework\_Open\_Source\_Provider
\item Source file: docs/archive/conversations/Framework\_Open\_Source\_Provider/Domino\_Framework\_Open\_Source\_Provider\_Architecture\_\_01\_human\_readable\_report.md
\item Source heading: automatic decision extraction
\item Claim: Finally, the user uploaded a detailed preservation prompt and requested a maximum-fidelity report, registers, spec sheet, aggregator packet, audit, and downloadable files. This package is the response to that request.
\item Claim type: conversation\_claim
\item Proposed target: not\_selected
\item Authority domain: planning
\item Current repo conflict: not\_checked
\item Why it may matter: The claim recurs as a decision-like statement in the archived package.
\item Risks: Chat claims may be stale, superseded, or contradicted by current repo artifacts.
\item Required review: human review plus authority-order check
\item Recommended disposition: not\_reviewed
\end{itemize}

\subsection{PROMOTE-0068 - This became the central architectural theme. The conversation refined ea}

\begin{itemize}
\item Source conversation: Framework\_Open\_Source\_Provider
\item Source file: docs/archive/conversations/Framework\_Open\_Source\_Provider/Domino\_Framework\_Open\_Source\_Provider\_Architecture\_\_01\_human\_readable\_report.md
\item Source heading: automatic decision extraction
\item Claim: This became the central architectural theme. The conversation refined earlier vendor-shaped paths into service-first paths. The stable pattern is:
\item Claim type: conversation\_claim
\item Proposed target: not\_selected
\item Authority domain: planning
\item Current repo conflict: not\_checked
\item Why it may matter: The claim recurs as a decision-like statement in the archived package.
\item Risks: Chat claims may be stale, superseded, or contradicted by current repo artifacts.
\item Required review: human review plus authority-order check
\item Recommended disposition: not\_reviewed
\end{itemize}

\subsection{PROMOTE-0069 - The chat inspected julesc013/dominium and found evidence of existing a}

\begin{itemize}
\item Source conversation: Framework\_Open\_Source\_Provider
\item Source file: docs/archive/conversations/Framework\_Open\_Source\_Provider/Domino\_Framework\_Open\_Source\_Provider\_Architecture\_\_01\_human\_readable\_report.md
\item Source heading: automatic decision extraction
\item Claim: The chat inspected julesc013/dominium and found evidence of existing abstractions such as system and graphics layers, stubs, and soft-backed backends. The finding was that raylib could fill concrete provider gaps without replacing the architecture. However, current repo facts are stale/uncertain and must be verified, especially the C17/C++17 vs C90/C++98 baseline contradiction.
\item Claim type: conversation\_claim
\item Proposed target: not\_selected
\item Authority domain: planning
\item Current repo conflict: not\_checked
\item Why it may matter: The claim recurs as a decision-like statement in the archived package.
\item Risks: Chat claims may be stale, superseded, or contradicted by current repo artifacts.
\item Required review: human review plus authority-order check
\item Recommended disposition: not\_reviewed
\end{itemize}

\subsection{PROMOTE-0070 - That was an important change of direction. It meant the assistant's poli}

\begin{itemize}
\item Source conversation: gui\_binary\_content
\item Source file: docs/archive/conversations/gui\_binary\_content/dominium\_gui\_binary\_content\_handoff\_\_human\_readable\_chat\_report.txt
\item Source heading: automatic decision extraction
\item Claim: That was an important change of direction. It meant the assistant's polished prompt should not be treated as final. The work needed conceptual discussion first. The assistant then analyzed CONTENT0 and identified several issues that could cause bad assumptions to become embedded if Codex acted too soon.
\item Claim type: conversation\_claim
\item Proposed target: not\_selected
\item Authority domain: planning
\item Current repo conflict: not\_checked
\item Why it may matter: The claim recurs as a decision-like statement in the archived package.
\item Risks: Chat claims may be stale, superseded, or contradicted by current repo artifacts.
\item Required review: human review plus authority-order check
\item Recommended disposition: not\_reviewed
\end{itemize}

\subsection{PROMOTE-0071 - what start scenarios must explicitly exclude,}

\begin{itemize}
\item Source conversation: gui\_binary\_content
\item Source file: docs/archive/conversations/gui\_binary\_content/dominium\_gui\_binary\_content\_handoff\_\_human\_readable\_chat\_report.txt
\item Source heading: automatic decision extraction
\item Claim: what start scenarios must explicitly exclude,
\item Claim type: conversation\_claim
\item Proposed target: not\_selected
\item Authority domain: planning
\item Current repo conflict: not\_checked
\item Why it may matter: The claim recurs as a decision-like statement in the archived package.
\item Risks: Chat claims may be stale, superseded, or contradicted by current repo artifacts.
\item Required review: human review plus authority-order check
\item Recommended disposition: not\_reviewed
\end{itemize}

\subsection{PROMOTE-0072 - This shifted the content work from "populate required files" to "design}

\begin{itemize}
\item Source conversation: gui\_binary\_content
\item Source file: docs/archive/conversations/gui\_binary\_content/dominium\_gui\_binary\_content\_handoff\_\_human\_readable\_chat\_report.txt
\item Source heading: automatic decision extraction
\item Claim: This shifted the content work from "populate required files" to "design a scalable content representation system." It also made clear that procedural generation and defined data are not enemies in the user's vision. The desired architecture must allow them to coexist.
\item Claim type: conversation\_claim
\item Proposed target: not\_selected
\item Authority domain: planning
\item Current repo conflict: not\_checked
\item Why it may matter: The claim recurs as a decision-like statement in the archived package.
\item Risks: Chat claims may be stale, superseded, or contradicted by current repo artifacts.
\item Required review: human review plus authority-order check
\item Recommended disposition: not\_reviewed
\end{itemize}

\subsection{PROMOTE-0073 - The 32-bit question followed. The answer was to make full native product}

\begin{itemize}
\item Source conversation: Language\_Platform\_Architecture
\item Source file: docs/archive/conversations/Language\_Platform\_Architecture/Dominium\_Language\_Platform\_Architecture\_Baseline\_\_01\_human\_readable\_report.md
\item Source heading: automatic decision extraction
\item Claim: The 32-bit question followed. The answer was to make full native products 64-bit and keep 32-bit as constrained or projection support only. The key caveat was not to make native pointer size part of game law. Save, replay, network, IDs, and renderer IR must use fixed-width types and never serialize size\_t, long, uintptr\_t, or raw pointers.
\item Claim type: conversation\_claim
\item Proposed target: not\_selected
\item Authority domain: planning
\item Current repo conflict: not\_checked
\item Why it may matter: The claim recurs as a decision-like statement in the archived package.
\item Risks: Chat claims may be stale, superseded, or contradicted by current repo artifacts.
\item Required review: human review plus authority-order check
\item Recommended disposition: not\_reviewed
\end{itemize}

\subsection{PROMOTE-0074 - The user then consolidated a broad future plan around Domino, Dominium,}

\begin{itemize}
\item Source conversation: Language\_Platform\_Architecture
\item Source file: docs/archive/conversations/Language\_Platform\_Architecture/Dominium\_Language\_Platform\_Architecture\_Baseline\_\_01\_human\_readable\_report.md
\item Source heading: automatic decision extraction
\item Claim: The user then consolidated a broad future plan around Domino, Dominium, Workbench, AIDE, contracts, tests, replay, and evidence. The answer agreed that the plan was strong, but identified missing central pieces: composition resolver, lockfiles, compatibility corpus, trust/permissions, virtual filesystem roots, performance budgets, and stable public-surface promotion rules.
\item Claim type: conversation\_claim
\item Proposed target: not\_selected
\item Authority domain: planning
\item Current repo conflict: not\_checked
\item Why it may matter: The claim recurs as a decision-like statement in the archived package.
\item Risks: Chat claims may be stale, superseded, or contradicted by current repo artifacts.
\item Required review: human review plus authority-order check
\item Recommended disposition: not\_reviewed
\end{itemize}

\subsection{PROMOTE-0075 - Finally, the user pasted advice favoring C99 or C++11 due to raylib/SDL}

\begin{itemize}
\item Source conversation: Language\_Platform\_Architecture
\item Source file: docs/archive/conversations/Language\_Platform\_Architecture/Dominium\_Language\_Platform\_Architecture\_Baseline\_\_01\_human\_readable\_report.md
\item Source heading: automatic decision extraction
\item Claim: Finally, the user pasted advice favoring C99 or C++11 due to raylib/SDL and legacy targets. The answer rejected a pivot. Raylib and SDL2 being C APIs only means provider boundaries should be C-callable; it does not force the whole engine or game to C99. The final recommendation remained C17 + C++17, with raylib/SDL2 treated as providers and with external deployment claims placed into the verification queue.
\item Claim type: conversation\_claim
\item Proposed target: not\_selected
\item Authority domain: planning
\item Current repo conflict: not\_checked
\item Why it may matter: The claim recurs as a decision-like statement in the archived package.
\item Risks: Chat claims may be stale, superseded, or contradicted by current repo artifacts.
\item Required review: human review plus authority-order check
\item Recommended disposition: not\_reviewed
\end{itemize}

\subsection{PROMOTE-0076 - Someone reading this report should understand one central thing: this ch}

\begin{itemize}
\item Source conversation: launcher\_app\_layer
\item Source file: docs/archive/conversations/launcher\_app\_layer/dominium\_launcher\_app\_layer\_handoff\_\_human\_readable\_chat\_report.txt
\item Source heading: automatic decision extraction
\item Claim: Someone reading this report should understand one central thing: this chat is not about inventing new launcher features anymore. It is about preserving boundaries, making the launcher implementation explicit, and ensuring future work happens on top of verified code rather than vague architectural memory.
\item Claim type: conversation\_claim
\item Proposed target: not\_selected
\item Authority domain: planning
\item Current repo conflict: not\_checked
\item Why it may matter: The claim recurs as a decision-like statement in the archived package.
\item Risks: Chat claims may be stale, superseded, or contradicted by current repo artifacts.
\item Required review: human review plus authority-order check
\item Recommended disposition: not\_reviewed
\end{itemize}

\subsection{PROMOTE-0077 - There was also an early suggestion to create a DUI facade, a Dominium}

\begin{itemize}
\item Source conversation: launcher\_app\_layer
\item Source file: docs/archive/conversations/launcher\_app\_layer/dominium\_launcher\_app\_layer\_handoff\_\_human\_readable\_chat\_report.txt
\item Source heading: automatic decision extraction
\item Claim: There was also an early suggestion to create a DUI facade, a Dominium UI abstraction, to support native widgets and fallback rendering. This idea was useful as a conceptual stepping stone but was not ultimately locked as a final requirement in the form originally suggested. Later application-layer canon emphasized **UI IR**, **command graphs**, and **binding validation** rather than a specific DUI facade design.
\item Claim type: conversation\_claim
\item Proposed target: not\_selected
\item Authority domain: planning
\item Current repo conflict: not\_checked
\item Why it may matter: The claim recurs as a decision-like statement in the archived package.
\item Risks: Chat claims may be stale, superseded, or contradicted by current repo artifacts.
\item Required review: human review plus authority-order check
\item Recommended disposition: not\_reviewed
\end{itemize}

\subsection{PROMOTE-0078 - This phase mattered because it framed the launcher not as a menu program}

\begin{itemize}
\item Source conversation: launcher\_app\_layer
\item Source file: docs/archive/conversations/launcher\_app\_layer/dominium\_launcher\_app\_layer\_handoff\_\_human\_readable\_chat\_report.txt
\item Source heading: automatic decision extraction
\item Claim: This phase mattered because it framed the launcher not as a menu program but as a **control plane** for installed products, packs, profiles, compatibility, and launch contracts. However, later canon tightened the permitted communication routes: cross-product communication must go through schema/ and libs/contracts, not arbitrary plugin conventions.
\item Claim type: conversation\_claim
\item Proposed target: not\_selected
\item Authority domain: planning
\item Current repo conflict: not\_checked
\item Why it may matter: The claim recurs as a decision-like statement in the archived package.
\item Risks: Chat claims may be stale, superseded, or contradicted by current repo artifacts.
\item Required review: human review plus authority-order check
\item Recommended disposition: not\_reviewed
\end{itemize}

\subsection{PROMOTE-0079 - The chat also produced multiple Codex work-order prompts and finally a d}

\begin{itemize}
\item Source conversation: Launcher\_Setup\_Architecture
\item Source file: docs/archive/conversations/Launcher\_Setup\_Architecture/Dominium\_Launcher\_Setup\_Architecture\_\_human\_readable\_chat\_report.txt
\item Source heading: automatic decision extraction
\item Claim: The chat also produced multiple Codex work-order prompts and finally a downloadable report package. Those artifacts are useful, but the substance is the architecture: keep engine deterministic, keep setup/launcher/runtime boundaries explicit, make the launcher optional and modular, use a strong process/instance model, and now implement the launcher over dsys/dgfx rather than ad hoc platform UI stacks.
\item Claim type: conversation\_claim
\item Proposed target: not\_selected
\item Authority domain: planning
\item Current repo conflict: not\_checked
\item Why it may matter: The claim recurs as a decision-like statement in the archived package.
\item Risks: Chat claims may be stale, superseded, or contradicted by current repo artifacts.
\item Required review: human review plus authority-order check
\item Recommended disposition: not\_reviewed
\end{itemize}

\subsection{PROMOTE-0080 - FACT:** The user's initial architecture text emphasized:}

\begin{itemize}
\item Source conversation: Launcher\_Setup\_Architecture
\item Source file: docs/archive/conversations/Launcher\_Setup\_Architecture/Dominium\_Launcher\_Setup\_Architecture\_\_human\_readable\_chat\_report.txt
\item Source heading: automatic decision extraction
\item Claim: FACT:** The user's initial architecture text emphasized:
\item Claim type: conversation\_claim
\item Proposed target: not\_selected
\item Authority domain: planning
\item Current repo conflict: not\_checked
\item Why it may matter: The claim recurs as a decision-like statement in the archived package.
\item Risks: Chat claims may be stale, superseded, or contradicted by current repo artifacts.
\item Required review: human review plus authority-order check
\item Recommended disposition: not\_reviewed
\end{itemize}

\subsection{PROMOTE-0081 - The assistant responded by stress-testing this philosophy. It pointed ou}

\begin{itemize}
\item Source conversation: Launcher\_Setup\_Architecture
\item Source file: docs/archive/conversations/Launcher\_Setup\_Architecture/Dominium\_Launcher\_Setup\_Architecture\_\_human\_readable\_chat\_report.txt
\item Source heading: automatic decision extraction
\item Claim: The assistant responded by stress-testing this philosophy. It pointed out determinism risks around Lua, plugins, time sources, and ABI details, and suggested more explicit contracts for file headers, TLV sections, runtime CLI capabilities, plugin exported symbols, and setup/launcher boundaries. These were assistant proposals, not all user-stated decisions, but the user continued building on that direction.
\item Claim type: conversation\_claim
\item Proposed target: not\_selected
\item Authority domain: planning
\item Current repo conflict: not\_checked
\item Why it may matter: The claim recurs as a decision-like statement in the archived package.
\item Risks: Chat claims may be stale, superseded, or contradicted by current repo artifacts.
\item Required review: human review plus authority-order check
\item Recommended disposition: not\_reviewed
\end{itemize}

\subsection{PROMOTE-0082 - The final user action was uploading a preservation-package prompt. That}

\begin{itemize}
\item Source conversation: Modularity\_AIDE\_Refactorability
\item Source file: docs/archive/conversations/Modularity\_AIDE\_Refactorability/Dominium\_Modularity\_AIDE\_Refactorability\_Architecture\_\_01\_human\_readable\_report.md
\item Source heading: automatic decision extraction
\item Claim: The final user action was uploading a preservation-package prompt. That prompt requested a maximum-fidelity preservation report, structured registers, context transfer packet, spec sheet, aggregator packet, self-audit, downloadable files, and a final in-chat reader. This package is the result.
\item Claim type: conversation\_claim
\item Proposed target: not\_selected
\item Authority domain: planning
\item Current repo conflict: not\_checked
\item Why it may matter: The claim recurs as a decision-like statement in the archived package.
\item Risks: Chat claims may be stale, superseded, or contradicted by current repo artifacts.
\item Required review: human review plus authority-order check
\item Recommended disposition: not\_reviewed
\end{itemize}

\subsection{PROMOTE-0083 - The user objected to the XStack/AuditX/RepoX/TestX framing and wanted AI}

\begin{itemize}
\item Source conversation: Modularity\_AIDE\_Refactorability
\item Source file: docs/archive/conversations/Modularity\_AIDE\_Refactorability/Dominium\_Modularity\_AIDE\_Refactorability\_Architecture\_\_01\_human\_readable\_report.md
\item Source heading: automatic decision extraction
\item Claim: The user objected to the XStack/AuditX/RepoX/TestX framing and wanted AIDE adopted quickly. The resulting plan made AIDE the repo-native control plane for inventory, task queues, policies, move maps, salvage maps, evidence ledgers, validation, and refactor history. Existing tools should be recycled, not discarded. This is central to future refactorability.
\item Claim type: conversation\_claim
\item Proposed target: not\_selected
\item Authority domain: planning
\item Current repo conflict: not\_checked
\item Why it may matter: The claim recurs as a decision-like statement in the archived package.
\item Risks: Chat claims may be stale, superseded, or contradicted by current repo artifacts.
\item Required review: human review plus authority-order check
\item Recommended disposition: not\_reviewed
\end{itemize}

\subsection{PROMOTE-0084 - The discussion proposed target-based CMake and module boundaries rather}

\begin{itemize}
\item Source conversation: Modularity\_AIDE\_Refactorability
\item Source file: docs/archive/conversations/Modularity\_AIDE\_Refactorability/Dominium\_Modularity\_AIDE\_Refactorability\_Architecture\_\_01\_human\_readable\_report.md
\item Source heading: automatic decision extraction
\item Claim: The discussion proposed target-based CMake and module boundaries rather than path mythology. Public headers, private headers, allowed dependencies, and forbidden dependencies should be explicit. Apps must remain thin. Engine must not depend on game/apps/runtime UI. Runtime adapts host/platform/rendering without owning simulation truth.
\item Claim type: conversation\_claim
\item Proposed target: not\_selected
\item Authority domain: planning
\item Current repo conflict: not\_checked
\item Why it may matter: The claim recurs as a decision-like statement in the archived package.
\item Risks: Chat claims may be stale, superseded, or contradicted by current repo artifacts.
\item Required review: human review plus authority-order check
\item Recommended disposition: not\_reviewed
\end{itemize}

\subsection{PROMOTE-0085 - The final strategic direction before this preservation request was to ru}

\begin{itemize}
\item Source conversation: omega\_xi\_pi\_architecture\_future
\item Source file: docs/archive/conversations/omega\_xi\_pi\_architecture\_future/dominium\_omega\_xi\_pi\_architecture\_future\_proofing\_planning\_\_01\_human\_readable\_report.md
\item Source heading: automatic decision extraction
\item Claim: The final strategic direction before this preservation request was to run a ?-series: snapshot intake, reality extraction, blueprint reconciliation, foundation readiness, and final prompt synthesis. This is needed because plans must now be mapped onto current repo reality rather than executed abstractly. The next chat should pick up there.
\item Claim type: conversation\_claim
\item Proposed target: not\_selected
\item Authority domain: planning
\item Current repo conflict: not\_checked
\item Why it may matter: The claim recurs as a decision-like statement in the archived package.
\item Risks: Chat claims may be stale, superseded, or contradicted by current repo artifacts.
\item Required review: human review plus authority-order check
\item Recommended disposition: not\_reviewed
\end{itemize}

\subsection{PROMOTE-0086 - The user later reported Xi completion. The enduring lesson is that promp}

\begin{itemize}
\item Source conversation: omega\_xi\_pi\_architecture\_future
\item Source file: docs/archive/conversations/omega\_xi\_pi\_architecture\_future/dominium\_omega\_xi\_pi\_architecture\_future\_proofing\_planning\_\_01\_human\_readable\_report.md
\item Source heading: automatic decision extraction
\item Claim: The user later reported Xi completion. The enduring lesson is that prompt instructions are not enough; architecture must be machine-readable and enforced.
\item Claim type: conversation\_claim
\item Proposed target: not\_selected
\item Authority domain: planning
\item Current repo conflict: not\_checked
\item Why it may matter: The claim recurs as a decision-like statement in the archived package.
\item Risks: Chat claims may be stale, superseded, or contradicted by current repo artifacts.
\item Required review: human review plus authority-order check
\item Recommended disposition: not\_reviewed
\end{itemize}

\subsection{PROMOTE-0087 - 4. Preserve all plans, tasks, constraints, risks, decisions, artifacts,}

\begin{itemize}
\item Source conversation: omega\_xi\_pi\_architecture\_future
\item Source file: docs/archive/conversations/omega\_xi\_pi\_architecture\_future/dominium\_omega\_xi\_pi\_architecture\_future\_proofing\_planning\_\_01\_human\_readable\_report.md
\item Source heading: automatic decision extraction
\item Claim: 4. Preserve all plans, tasks, constraints, risks, decisions, artifacts, and future directions across chats.
\item Claim type: conversation\_claim
\item Proposed target: not\_selected
\item Authority domain: planning
\item Current repo conflict: not\_checked
\item Why it may matter: The claim recurs as a decision-like statement in the archived package.
\item Risks: Chat claims may be stale, superseded, or contradicted by current repo artifacts.
\item Required review: human review plus authority-order check
\item Recommended disposition: not\_reviewed
\end{itemize}

\subsection{PROMOTE-0088 - DECISION-01 - CLI mandatory, TUI expected, GUI modular.** This was part}

\begin{itemize}
\item Source conversation: OS\_Interface\_Repo\_Architecture
\item Source file: docs/archive/conversations/OS\_Interface\_Repo\_Architecture/Dominium\_OS\_Interface\_Repo\_Architecture\_\_01\_human\_readable\_report.md
\item Source heading: automatic decision extraction
\item Claim: DECISION-01 - CLI mandatory, TUI expected, GUI modular.** This was part of the initial architecture baseline and was repeatedly reaffirmed. It matters because every product must remain operable in recovery/headless/automation contexts. GUI is allowed but not authoritative.
\item Claim type: conversation\_claim
\item Proposed target: not\_selected
\item Authority domain: planning
\item Current repo conflict: not\_checked
\item Why it may matter: The claim recurs as a decision-like statement in the archived package.
\item Risks: Chat claims may be stale, superseded, or contradicted by current repo artifacts.
\item Required review: human review plus authority-order check
\item Recommended disposition: not\_reviewed
\end{itemize}

\subsection{PROMOTE-0089 - DECISION-02 - Thin GUI shells over shared contracts.** The chat consiste}

\begin{itemize}
\item Source conversation: OS\_Interface\_Repo\_Architecture
\item Source file: docs/archive/conversations/OS\_Interface\_Repo\_Architecture/Dominium\_OS\_Interface\_Repo\_Architecture\_\_01\_human\_readable\_report.md
\item Source heading: automatic decision extraction
\item Claim: DECISION-02 - Thin GUI shells over shared contracts.** The chat consistently rejected GUI families becoming separate product architectures. This affects client, launcher, setup, server admin, tools, and Workbench modules.
\item Claim type: conversation\_claim
\item Proposed target: not\_selected
\item Authority domain: planning
\item Current repo conflict: not\_checked
\item Why it may matter: The claim recurs as a decision-like statement in the archived package.
\item Risks: Chat claims may be stale, superseded, or contradicted by current repo artifacts.
\item Required review: human review plus authority-order check
\item Recommended disposition: not\_reviewed
\end{itemize}

\subsection{PROMOTE-0090 - DECISION-03 - Repo ownership layout.** The user pushed for repository co}

\begin{itemize}
\item Source conversation: OS\_Interface\_Repo\_Architecture
\item Source file: docs/archive/conversations/OS\_Interface\_Repo\_Architecture/Dominium\_OS\_Interface\_Repo\_Architecture\_\_01\_human\_readable\_report.md
\item Source heading: automatic decision extraction
\item Claim: DECISION-03 - Repo ownership layout.** The user pushed for repository convergence; the discussion established that folders should map to ownership and contract boundaries. This decision is final enough to guide work, but details remain subject to machine-readable layout contracts.
\item Claim type: conversation\_claim
\item Proposed target: not\_selected
\item Authority domain: planning
\item Current repo conflict: not\_checked
\item Why it may matter: The claim recurs as a decision-like statement in the archived package.
\item Risks: Chat claims may be stale, superseded, or contradicted by current repo artifacts.
\item Required review: human review plus authority-order check
\item Recommended disposition: not\_reviewed
\end{itemize}

\subsection{PROMOTE-0091 - The chat ended by generating a maximum-fidelity transfer packet, then a}

\begin{itemize}
\item Source conversation: platform\_renderer\_api\_plan
\item Source file: docs/archive/conversations/platform\_renderer\_api\_plan/dominium\_codex\_platform\_renderer\_api\_plan\_\_human\_readable\_chat\_report.txt
\item Source heading: automatic decision extraction
\item Claim: The chat ended by generating a maximum-fidelity transfer packet, then a downloadable report package, then an in-chat reader version. The main thing to remember is this: **the active artifact is the final 9-prompt post-master Codex plan, but it is a plan, not evidence that the code exists. The repo must be inspected before execution.
\item Claim type: conversation\_claim
\item Proposed target: not\_selected
\item Authority domain: planning
\item Current repo conflict: not\_checked
\item Why it may matter: The claim recurs as a decision-like statement in the archived package.
\item Risks: Chat claims may be stale, superseded, or contradicted by current repo artifacts.
\item Required review: human review plus authority-order check
\item Recommended disposition: not\_reviewed
\end{itemize}

\subsection{PROMOTE-0092 - These ideas became the architectural heart of the final plan.}

\begin{itemize}
\item Source conversation: platform\_renderer\_api\_plan
\item Source file: docs/archive/conversations/platform\_renderer\_api\_plan/dominium\_codex\_platform\_renderer\_api\_plan\_\_human\_readable\_chat\_report.txt
\item Source heading: automatic decision extraction
\item Claim: These ideas became the architectural heart of the final plan.
\item Claim type: conversation\_claim
\item Proposed target: not\_selected
\item Authority domain: planning
\item Current repo conflict: not\_checked
\item Why it may matter: The claim recurs as a decision-like statement in the archived package.
\item Risks: Chat claims may be stale, superseded, or contradicted by current repo artifacts.
\item Required review: human review plus authority-order check
\item Recommended disposition: not\_reviewed
\end{itemize}

\subsection{PROMOTE-0093 - Stable ABI boundaries must use C ABI, POD structs, explicit function tab}

\begin{itemize}
\item Source conversation: platform\_renderer\_api\_plan
\item Source file: docs/archive/conversations/platform\_renderer\_api\_plan/dominium\_codex\_platform\_renderer\_api\_plan\_\_human\_readable\_chat\_report.txt
\item Source heading: automatic decision extraction
\item Claim: Stable ABI boundaries must use C ABI, POD structs, explicit function tables, and u32 abi\_version plus u32 struct\_size first.
\item Claim type: conversation\_claim
\item Proposed target: not\_selected
\item Authority domain: planning
\item Current repo conflict: not\_checked
\item Why it may matter: The claim recurs as a decision-like statement in the archived package.
\item Risks: Chat claims may be stale, superseded, or contradicted by current repo artifacts.
\item Required review: human review plus authority-order check
\item Recommended disposition: not\_reviewed
\end{itemize}

\subsection{PROMOTE-0094 - FACT: That early answer was generated by the assistant, not stated by th}

\begin{itemize}
\item Source conversation: platform\_support
\item Source file: docs/archive/conversations/platform\_support/dominium\_platform\_support\_planning\_\_human\_readable\_chat\_report.txt
\item Source heading: automatic decision extraction
\item Claim: FACT: That early answer was generated by the assistant, not stated by the user as a final decision. It matters because it introduced a very broad portability ambition, but it was later superseded by more practical platform planning. It should now be treated as historical context or possible research scope, not as a controlling product commitment.
\item Claim type: conversation\_claim
\item Proposed target: not\_selected
\item Authority domain: planning
\item Current repo conflict: not\_checked
\item Why it may matter: The claim recurs as a decision-like statement in the archived package.
\item Risks: Chat claims may be stale, superseded, or contradicted by current repo artifacts.
\item Required review: human review plus authority-order check
\item Recommended disposition: not\_reviewed
\end{itemize}

\subsection{PROMOTE-0095 - The user then supplied a detailed inventory covering PlayStation, Xbox,}

\begin{itemize}
\item Source conversation: platform\_support
\item Source file: docs/archive/conversations/platform\_support/dominium\_platform\_support\_planning\_\_human\_readable\_chat\_report.txt
\item Source heading: automatic decision extraction
\item Claim: The user then supplied a detailed inventory covering PlayStation, Xbox, Nintendo, PC handhelds, Android devices, Apple devices, Web platforms, AR, VR, and cross-cutting software targets. This inventory became the central artifact of the chat.
\item Claim type: conversation\_claim
\item Proposed target: not\_selected
\item Authority domain: planning
\item Current repo conflict: not\_checked
\item Why it may matter: The claim recurs as a decision-like statement in the archived package.
\item Risks: Chat claims may be stale, superseded, or contradicted by current repo artifacts.
\item Required review: human review plus authority-order check
\item Recommended disposition: not\_reviewed
\end{itemize}

\subsection{PROMOTE-0096 - FACT: The user's list included legacy and current consoles, handhelds, f}

\begin{itemize}
\item Source conversation: platform\_support
\item Source file: docs/archive/conversations/platform\_support/dominium\_platform\_support\_planning\_\_human\_readable\_chat\_report.txt
\item Source heading: automatic decision extraction
\item Claim: FACT: The user's list included legacy and current consoles, handhelds, firmware/OS names, Android software variants, Apple OSes, Web runtimes, AR/VR hardware, XR platforms, and graphics/runtime APIs. FACT: It was a list of desired coverage or consideration, not a final statement that every item must receive full parity support.
\item Claim type: conversation\_claim
\item Proposed target: not\_selected
\item Authority domain: planning
\item Current repo conflict: not\_checked
\item Why it may matter: The claim recurs as a decision-like statement in the archived package.
\item Risks: Chat claims may be stale, superseded, or contradicted by current repo artifacts.
\item Required review: human review plus authority-order check
\item Recommended disposition: not\_reviewed
\end{itemize}

\subsection{PROMOTE-0097 - The actual repo was not inspected. The exact accepted directory structur}

\begin{itemize}
\item Source conversation: Portability\_Assurance\_Future\_Proof
\item Source file: docs/archive/conversations/Portability\_Assurance\_Future\_Proof/Domino\_Dominium\_Portability\_Assurance\_Future\_Proof\_Architecture\_\_01\_human\_readable\_report.md
\item Source heading: automatic decision extraction
\item Claim: The actual repo was not inspected. The exact accepted directory structure, API policy, DDAP profile, compatibility promises, and first pilot module remain unresolved.
\item Claim type: conversation\_claim
\item Proposed target: not\_selected
\item Authority domain: planning
\item Current repo conflict: not\_checked
\item Why it may matter: The claim recurs as a decision-like statement in the archived package.
\item Risks: Chat claims may be stale, superseded, or contradicted by current repo artifacts.
\item Required review: human review plus authority-order check
\item Recommended disposition: not\_reviewed
\end{itemize}

\subsection{PROMOTE-0098 - Status: Assistant recommendation; not separately accepted after answer.}

\begin{itemize}
\item Source conversation: Portability\_Assurance\_Future\_Proof
\item Source file: docs/archive/conversations/Portability\_Assurance\_Future\_Proof/Domino\_Dominium\_Portability\_Assurance\_Future\_Proof\_Architecture\_\_01\_human\_readable\_report.md
\item Source heading: automatic decision extraction
\item Claim: Status: Assistant recommendation; not separately accepted after answer.
\item Claim type: conversation\_claim
\item Proposed target: not\_selected
\item Authority domain: planning
\item Current repo conflict: not\_checked
\item Why it may matter: The claim recurs as a decision-like statement in the archived package.
\item Risks: Chat claims may be stale, superseded, or contradicted by current repo artifacts.
\item Required review: human review plus authority-order check
\item Recommended disposition: not\_reviewed
\end{itemize}

\subsection{PROMOTE-0099 - Rationale: Replacement must be objectively testable.}

\begin{itemize}
\item Source conversation: Portability\_Assurance\_Future\_Proof
\item Source file: docs/archive/conversations/Portability\_Assurance\_Future\_Proof/Domino\_Dominium\_Portability\_Assurance\_Future\_Proof\_Architecture\_\_01\_human\_readable\_report.md
\item Source heading: automatic decision extraction
\item Claim: Rationale: Replacement must be objectively testable.
\item Claim type: conversation\_claim
\item Proposed target: not\_selected
\item Authority domain: planning
\item Current repo conflict: not\_checked
\item Why it may matter: The claim recurs as a decision-like statement in the archived package.
\item Risks: Chat claims may be stale, superseded, or contradicted by current repo artifacts.
\item Required review: human review plus authority-order check
\item Recommended disposition: not\_reviewed
\end{itemize}

\subsection{PROMOTE-0100 - The user pasted the final README after those cleanup changes. At that po}

\begin{itemize}
\item Source conversation: readme\_ports\_determinism
\item Source file: docs/archive/conversations/readme\_ports\_determinism/dominium\_readme\_ports\_determinism\_\_human\_readable\_chat\_report.txt
\item Source heading: automatic decision extraction
\item Claim: The user pasted the final README after those cleanup changes. At that point, the README had the current active form: deterministic constraints clarified, ports unified under one source hierarchy, /ports metadata-only if retained, Section 9 normative, lockstep canonical, content-lock mismatch fatal, and disk format versions immutable.
\item Claim type: conversation\_claim
\item Proposed target: not\_selected
\item Authority domain: planning
\item Current repo conflict: not\_checked
\item Why it may matter: The claim recurs as a decision-like statement in the archived package.
\item Risks: Chat claims may be stale, superseded, or contradicted by current repo artifacts.
\item Required review: human review plus authority-order check
\item Recommended disposition: not\_reviewed
\end{itemize}

\subsection{PROMOTE-0101 - After the README work, the user asked for a maximum-fidelity context tra}

\begin{itemize}
\item Source conversation: readme\_ports\_determinism
\item Source file: docs/archive/conversations/readme\_ports\_determinism/dominium\_readme\_ports\_determinism\_\_human\_readable\_chat\_report.txt
\item Source heading: automatic decision extraction
\item Claim: After the README work, the user asked for a maximum-fidelity context transfer packet. The assistant produced a long state-transfer document with sections for project inventory, decisions, tasks, constraints, open questions, artifacts, risks, and next actions.
\item Claim type: conversation\_claim
\item Proposed target: not\_selected
\item Authority domain: planning
\item Current repo conflict: not\_checked
\item Why it may matter: The claim recurs as a decision-like statement in the archived package.
\item Risks: Chat claims may be stale, superseded, or contradicted by current repo artifacts.
\item Required review: human review plus authority-order check
\item Recommended disposition: not\_reviewed
\end{itemize}

\subsection{PROMOTE-0102 - The central artifact was the root README.md. The README describes **Do}

\begin{itemize}
\item Source conversation: readme\_ports\_determinism
\item Source file: docs/archive/conversations/readme\_ports\_determinism/dominium\_readme\_ports\_determinism\_\_human\_readable\_chat\_report.txt
\item Source heading: automatic decision extraction
\item Claim: The central artifact was the root README.md. The README describes **Domino** as the deterministic engine core and **Dominium** as the official game/tooling/runtime layer. It is written as a high-level architecture document, not a low-level implementation spec.
\item Claim type: conversation\_claim
\item Proposed target: not\_selected
\item Authority domain: planning
\item Current repo conflict: not\_checked
\item Why it may matter: The claim recurs as a decision-like statement in the archived package.
\item Risks: Chat claims may be stale, superseded, or contradicted by current repo artifacts.
\item Required review: human review plus authority-order check
\item Recommended disposition: not\_reviewed
\end{itemize}

\subsection{PROMOTE-0103 - The most important thing to remember is that this chat did not implement}

\begin{itemize}
\item Source conversation: Refactor\_Architecture
\item Source file: docs/archive/conversations/Refactor\_Architecture/Dominium\_Domino\_Refactor\_Architecture\_\_human\_readable\_chat\_report.txt
\item Source heading: automatic decision extraction
\item Claim: The most important thing to remember is that this chat did not implement the refactor. It designed the architecture and generated prompts and handoff material. Any future assistant must verify actual repository state before assuming files were moved, prompts were applied, or code was updated.
\item Claim type: conversation\_claim
\item Proposed target: not\_selected
\item Authority domain: planning
\item Current repo conflict: not\_checked
\item Why it may matter: The claim recurs as a decision-like statement in the archived package.
\item Risks: Chat claims may be stale, superseded, or contradicted by current repo artifacts.
\item Required review: human review plus authority-order check
\item Recommended disposition: not\_reviewed
\end{itemize}

\subsection{PROMOTE-0104 - FACT:** This early discussion created the future UI/packs design context}

\begin{itemize}
\item Source conversation: Refactor\_Architecture
\item Source file: docs/archive/conversations/Refactor\_Architecture/Dominium\_Domino\_Refactor\_Architecture\_\_human\_readable\_chat\_report.txt
\item Source heading: automatic decision extraction
\item Claim: FACT:** This early discussion created the future UI/packs design context.
\item Claim type: conversation\_claim
\item Proposed target: not\_selected
\item Authority domain: planning
\item Current repo conflict: not\_checked
\item Why it may matter: The claim recurs as a decision-like statement in the archived package.
\item Risks: Chat claims may be stale, superseded, or contradicted by current repo artifacts.
\item Required review: human review plus authority-order check
\item Recommended disposition: not\_reviewed
\end{itemize}

\subsection{PROMOTE-0105 - FACT:** All products should use shared dsys and dgfx. No product-spe}

\begin{itemize}
\item Source conversation: Refactor\_Architecture
\item Source file: docs/archive/conversations/Refactor\_Architecture/Dominium\_Domino\_Refactor\_Architecture\_\_human\_readable\_chat\_report.txt
\item Source heading: automatic decision extraction
\item Claim: FACT:** All products should use shared dsys and dgfx. No product-specific platform or renderer code path.
\item Claim type: conversation\_claim
\item Proposed target: not\_selected
\item Authority domain: planning
\item Current repo conflict: not\_checked
\item Why it may matter: The claim recurs as a decision-like statement in the archived package.
\item Risks: Chat claims may be stale, superseded, or contradicted by current repo artifacts.
\item Required review: human review plus authority-order check
\item Recommended disposition: not\_reviewed
\end{itemize}

\subsection{PROMOTE-0106 - The user then asked about optimizing prompts for GPT-5.3/GPT-5.4 and lat}

\begin{itemize}
\item Source conversation: Refactor\_Control\_Plane
\item Source file: docs/archive/conversations/Refactor\_Control\_Plane/AIDE\_XStack\_Dominium\_Refactor\_Control\_Plane\_\_01\_human\_readable\_report.md
\item Source heading: automatic decision extraction
\item Claim: The user then asked about optimizing prompts for GPT-5.3/GPT-5.4 and later supplied advice about harness engineering. That advice argued that teams succeed by engineering the environment around agents: tools, docs, linters, feedback loops, memory, and observability. The conversation accepted the core point: the model is not the whole system; the harness is the multiplier.
\item Claim type: conversation\_claim
\item Proposed target: not\_selected
\item Authority domain: planning
\item Current repo conflict: not\_checked
\item Why it may matter: The claim recurs as a decision-like statement in the archived package.
\item Risks: Chat claims may be stale, superseded, or contradicted by current repo artifacts.
\item Required review: human review plus authority-order check
\item Recommended disposition: not\_reviewed
\end{itemize}

\subsection{PROMOTE-0107 - The Grug Brained Developer became a design filter: avoid complexity, avo}

\begin{itemize}
\item Source conversation: Refactor\_Control\_Plane
\item Source file: docs/archive/conversations/Refactor\_Control\_Plane/AIDE\_XStack\_Dominium\_Refactor\_Control\_Plane\_\_01\_human\_readable\_report.md
\item Source heading: automatic decision extraction
\item Claim: The Grug Brained Developer became a design filter: avoid complexity, avoid premature abstraction, do small safe refactors, prefer integration tests, invest in logging, and respect existing structures before tearing them down. This influenced the decision to narrow AIDE and not overbuild runtime/hosts too early.
\item Claim type: conversation\_claim
\item Proposed target: not\_selected
\item Authority domain: planning
\item Current repo conflict: not\_checked
\item Why it may matter: The claim recurs as a decision-like statement in the archived package.
\item Risks: Chat claims may be stale, superseded, or contradicted by current repo artifacts.
\item Required review: human review plus authority-order check
\item Recommended disposition: not\_reviewed
\end{itemize}

\subsection{PROMOTE-0108 - The user eventually named AIDE: Automated Integrated Development Environ}

\begin{itemize}
\item Source conversation: Refactor\_Control\_Plane
\item Source file: docs/archive/conversations/Refactor\_Control\_Plane/AIDE\_XStack\_Dominium\_Refactor\_Control\_Plane\_\_01\_human\_readable\_report.md
\item Source heading: automatic decision extraction
\item Claim: The user eventually named AIDE: Automated Integrated Development Environment. The accepted split became: XStack remains Dominium's strict local governance/proof profile; AIDE becomes the portable public layer. AIDE would live in its own repo, be usable as a standalone repo pack, later as CLI/app/extensions, and eventually perhaps a runtime/service/host system.
\item Claim type: conversation\_claim
\item Proposed target: not\_selected
\item Authority domain: planning
\item Current repo conflict: not\_checked
\item Why it may matter: The claim recurs as a decision-like statement in the archived package.
\item Risks: Chat claims may be stale, superseded, or contradicted by current repo artifacts.
\item Required review: human review plus authority-order check
\item Recommended disposition: not\_reviewed
\end{itemize}

\subsection{PROMOTE-0109 - The user then objected to versioning policy drift in real products and e}

\begin{itemize}
\item Source conversation: Release\_Identity\_and\_Versioning
\item Source file: docs/archive/conversations/Release\_Identity\_and\_Versioning/Dominium\_XStack\_Release\_Identity\_and\_Versioning\_\_01\_human\_readable\_report.md
\item Source heading: automatic decision extraction
\item Claim: The user then objected to versioning policy drift in real products and expressed concern about SemVer's "1.x forever" failure mode. The first proposed answer was a four-part public version scheme, GEN.EPOCH.FEATURE.PATCH, meant to avoid fake major bumps. That idea was useful as an intermediate step but later became less central because XStack already has GBN for dense build history and separate build identity.
\item Claim type: conversation\_claim
\item Proposed target: not\_selected
\item Authority domain: planning
\item Current repo conflict: not\_checked
\item Why it may matter: The claim recurs as a decision-like statement in the archived package.
\item Risks: Chat claims may be stale, superseded, or contradicted by current repo artifacts.
\item Required review: human review plus authority-order check
\item Recommended disposition: not\_reviewed
\end{itemize}

\subsection{PROMOTE-0110 - The conversation then moved into how this might fit XStack. The model sh}

\begin{itemize}
\item Source conversation: Release\_Identity\_and\_Versioning
\item Source file: docs/archive/conversations/Release\_Identity\_and\_Versioning/Dominium\_XStack\_Release\_Identity\_and\_Versioning\_\_01\_human\_readable\_report.md
\item Source heading: automatic decision extraction
\item Claim: The conversation then moved into how this might fit XStack. The model shifted from "better public version number" to "layered identity model." The assistant recommended preserving per-product versions, a global build number, build IDs, compatibility versions, and suite versions as separate layers. This became the foundation for later decisions.
\item Claim type: conversation\_claim
\item Proposed target: not\_selected
\item Authority domain: planning
\item Current repo conflict: not\_checked
\item Why it may matter: The claim recurs as a decision-like statement in the archived package.
\item Risks: Chat claims may be stale, superseded, or contradicted by current repo artifacts.
\item Required review: human review plus authority-order check
\item Recommended disposition: not\_reviewed
\end{itemize}

\subsection{PROMOTE-0111 - The user suggested that suites might use a separate consumer-facing or m}

\begin{itemize}
\item Source conversation: Release\_Identity\_and\_Versioning
\item Source file: docs/archive/conversations/Release\_Identity\_and\_Versioning/Dominium\_XStack\_Release\_Identity\_and\_Versioning\_\_01\_human\_readable\_report.md
\item Source heading: automatic decision extraction
\item Claim: The user suggested that suites might use a separate consumer-facing or marketing version, while each component could use stricter SemVer. The chat accepted this as a mature pattern: suites are curated bundles, while components may have technical versions. The model was refined to avoid synchronized fake versioning and to avoid treating a suite major version as universal breakage.
\item Claim type: conversation\_claim
\item Proposed target: not\_selected
\item Authority domain: planning
\item Current repo conflict: not\_checked
\item Why it may matter: The claim recurs as a decision-like statement in the archived package.
\item Risks: Chat claims may be stale, superseded, or contradicted by current repo artifacts.
\item Required review: human review plus authority-order check
\item Recommended disposition: not\_reviewed
\end{itemize}

\subsection{PROMOTE-0112 - FACT:** The user is building Dominium / Domino as a deterministic univer}

\begin{itemize}
\item Source conversation: testx\_repox\_governance
\item Source file: docs/archive/conversations/testx\_repox\_governance/dominium\_testx\_repox\_governance\_handoff\_\_human\_readable\_chat\_report.txt
\item Source heading: automatic decision extraction
\item Claim: FACT:** The user is building Dominium / Domino as a deterministic universe engine + game, not a conventional game project. The visible chat repeatedly framed the project as long-lived, simulation-first, deterministic, modular, and designed to survive across many operating systems, toolchains, renderers, products, and distribution models.
\item Claim type: conversation\_claim
\item Proposed target: not\_selected
\item Authority domain: planning
\item Current repo conflict: not\_checked
\item Why it may matter: The claim recurs as a decision-like statement in the archived package.
\item Risks: Chat claims may be stale, superseded, or contradicted by current repo artifacts.
\item Required review: human review plus authority-order check
\item Recommended disposition: not\_reviewed
\end{itemize}

\subsection{PROMOTE-0113 - The user then asked whether the implementation was industry-accepted and}

\begin{itemize}
\item Source conversation: testx\_repox\_governance
\item Source file: docs/archive/conversations/testx\_repox\_governance/dominium\_testx\_repox\_governance\_handoff\_\_human\_readable\_chat\_report.txt
\item Source heading: automatic decision extraction
\item Claim: The user then asked whether the implementation was industry-accepted and what could be improved. The response framed the approach as closer to game engines, operating systems, and long-lived infrastructure than typical game development. The important conclusion was that the design was not exotic, but it was unusually rigorous for games.
\item Claim type: conversation\_claim
\item Proposed target: not\_selected
\item Authority domain: planning
\item Current repo conflict: not\_checked
\item Why it may matter: The claim recurs as a decision-like statement in the archived package.
\item Risks: Chat claims may be stale, superseded, or contradicted by current repo artifacts.
\item Required review: human review plus authority-order check
\item Recommended disposition: not\_reviewed
\end{itemize}

\subsection{PROMOTE-0114 - The user had another chat working on content and systems, and asked for}

\begin{itemize}
\item Source conversation: testx\_repox\_governance
\item Source file: docs/archive/conversations/testx\_repox\_governance/dominium\_testx\_repox\_governance\_handoff\_\_human\_readable\_chat\_report.txt
\item Source heading: automatic decision extraction
\item Claim: The user had another chat working on content and systems, and asked for a prompt to inform that chat of everything decided so far. A similar prompt was generated for the applications/platforms/renderers chat. These prompts established authoritative boundaries:
\item Claim type: conversation\_claim
\item Proposed target: not\_selected
\item Authority domain: planning
\item Current repo conflict: not\_checked
\item Why it may matter: The claim recurs as a decision-like statement in the archived package.
\item Risks: Chat claims may be stale, superseded, or contradicted by current repo artifacts.
\item Required review: human review plus authority-order check
\item Recommended disposition: not\_reviewed
\end{itemize}

\subsection{PROMOTE-0115 - A future assistant should understand that this chat contributes a time a}

\begin{itemize}
\item Source conversation: Timekeeping\_and\_2038\_Resilience
\item Source file: docs/archive/conversations/Timekeeping\_and\_2038\_Resilience/Dominium\_Timekeeping\_and\_2038\_Resilience\_\_01\_human\_readable\_report.md
\item Source heading: automatic decision extraction
\item Claim: A future assistant should understand that this chat contributes a time architecture doctrine for Dominium: **ACT is authority; DSYS time is runtime-only; observer clocks are derived; civil/astronomical time is projection-only; wall-clock time must never drive authoritative ordering.
\item Claim type: conversation\_claim
\item Proposed target: not\_selected
\item Authority domain: planning
\item Current repo conflict: not\_checked
\item Why it may matter: The claim recurs as a decision-like statement in the archived package.
\item Risks: Chat claims may be stale, superseded, or contradicted by current repo artifacts.
\item Required review: human review plus authority-order check
\item Recommended disposition: not\_reviewed
\end{itemize}

\subsection{PROMOTE-0116 - What remains uncertain is the precise target list of legacy machines/too}

\begin{itemize}
\item Source conversation: Timekeeping\_and\_2038\_Resilience
\item Source file: docs/archive/conversations/Timekeeping\_and\_2038\_Resilience/Dominium\_Timekeeping\_and\_2038\_Resilience\_\_01\_human\_readable\_report.md
\item Source heading: automatic decision extraction
\item Claim: What remains uncertain is the precise target list of legacy machines/toolchains and how far compatibility must go, especially for 16-bit compilers with weak or absent 64-bit arithmetic.
\item Claim type: conversation\_claim
\item Proposed target: not\_selected
\item Authority domain: planning
\item Current repo conflict: not\_checked
\item Why it may matter: The claim recurs as a decision-like statement in the archived package.
\item Risks: Chat claims may be stale, superseded, or contradicted by current repo artifacts.
\item Required review: human review plus authority-order check
\item Recommended disposition: not\_reviewed
\end{itemize}

\subsection{PROMOTE-0117 - The final topic is this export task. The user requested a human-readable}

\begin{itemize}
\item Source conversation: Timekeeping\_and\_2038\_Resilience
\item Source file: docs/archive/conversations/Timekeeping\_and\_2038\_Resilience/Dominium\_Timekeeping\_and\_2038\_Resilience\_\_01\_human\_readable\_report.md
\item Source heading: automatic decision extraction
\item Claim: The final topic is this export task. The user requested a human-readable report first, followed by registers, spec prep, context-transfer packet, aggregator packet, self-audit, and files. This package is designed to let the chat be read later without reopening the full transcript and to support merger into a larger Dominium spec book.
\item Claim type: conversation\_claim
\item Proposed target: not\_selected
\item Authority domain: planning
\item Current repo conflict: not\_checked
\item Why it may matter: The claim recurs as a decision-like statement in the archived package.
\item Risks: Chat claims may be stale, superseded, or contradicted by current repo artifacts.
\item Required review: human review plus authority-order check
\item Recommended disposition: not\_reviewed
\end{itemize}

\subsection{PROMOTE-0118 - FACT: This package preserves the currently visible chat about whether Do}

\begin{itemize}
\item Source conversation: UE6\_Domino\_Deterministic\_Universe
\item Source file: docs/archive/conversations/UE6\_Domino\_Deterministic\_Universe/Dominium\_UE6\_Domino\_Deterministic\_Universe\_\_01\_human\_readable\_report.md
\item Source heading: automatic decision extraction
\item Claim: FACT: This package preserves the currently visible chat about whether Dominium should use UE6, UE5, Domino, Unreal, or a custom engine layer for a highly ambitious deterministic solar-system-scale game. It also preserves the current uploaded prompt as an artifact because it defines this preservation task.
\item Claim type: conversation\_claim
\item Proposed target: not\_selected
\item Authority domain: planning
\item Current repo conflict: not\_checked
\item Why it may matter: The claim recurs as a decision-like statement in the archived package.
\item Risks: Chat claims may be stale, superseded, or contradicted by current repo artifacts.
\item Required review: human review plus authority-order check
\item Recommended disposition: not\_reviewed
\end{itemize}

\subsection{PROMOTE-0119 - That response also introduced a rule that shaped the rest of the chat: U}

\begin{itemize}
\item Source conversation: UE6\_Domino\_Deterministic\_Universe
\item Source file: docs/archive/conversations/UE6\_Domino\_Deterministic\_Universe/Dominium\_UE6\_Domino\_Deterministic\_Universe\_\_01\_human\_readable\_report.md
\item Source heading: automatic decision extraction
\item Claim: That response also introduced a rule that shaped the rest of the chat: Unreal should render and host Dominium, but Unreal should not define Dominium. That was the first major conceptual decision point. It reframed Dominium not as an engine project in the narrow sense, but as a portable game/simulation system with replaceable frontends.
\item Claim type: conversation\_claim
\item Proposed target: not\_selected
\item Authority domain: planning
\item Current repo conflict: not\_checked
\item Why it may matter: The claim recurs as a decision-like statement in the archived package.
\item Risks: Chat claims may be stale, superseded, or contradicted by current repo artifacts.
\item Required review: human review plus authority-order check
\item Recommended disposition: not\_reviewed
\end{itemize}

\subsection{PROMOTE-0120 - Conclusion: machine gameplay must be designed as a scalable graph simula}

\begin{itemize}
\item Source conversation: UE6\_Domino\_Deterministic\_Universe
\item Source file: docs/archive/conversations/UE6\_Domino\_Deterministic\_Universe/Dominium\_UE6\_Domino\_Deterministic\_Universe\_\_01\_human\_readable\_report.md
\item Source heading: automatic decision extraction
\item Claim: Conclusion: machine gameplay must be designed as a scalable graph simulation, not as unrestricted physics.
\item Claim type: conversation\_claim
\item Proposed target: not\_selected
\item Authority domain: planning
\item Current repo conflict: not\_checked
\item Why it may matter: The claim recurs as a decision-like statement in the archived package.
\item Risks: Chat claims may be stale, superseded, or contradicted by current repo artifacts.
\item Required review: human review plus authority-order check
\item Recommended disposition: not\_reviewed
\end{itemize}

\subsection{PROMOTE-0121 - The most important thing to remember is that this chat was a planning an}

\begin{itemize}
\item Source conversation: ui\_editor\_tool\_editor\_planning
\item Source file: docs/archive/conversations/ui\_editor\_tool\_editor\_planning/dominium\_ui\_editor\_tool\_editor\_planning\_\_human\_readable\_chat\_report.txt
\item Source heading: automatic decision extraction
\item Claim: The most important thing to remember is that this chat was a planning and prompt-generation chat, not proof of implementation. **UNCERTAIN / UNVERIFIED:** No generated prompt is evidence that repository code was actually changed. The next assistant must first verify whether the prompts were executed, inspect the repo/source bundle if available, and avoid treating planned files as existing files.
\item Claim type: conversation\_claim
\item Proposed target: not\_selected
\item Authority domain: planning
\item Current repo conflict: not\_checked
\item Why it may matter: The claim recurs as a decision-like statement in the archived package.
\item Risks: Chat claims may be stale, superseded, or contradicted by current repo artifacts.
\item Required review: human review plus authority-order check
\item Recommended disposition: not\_reviewed
\end{itemize}

\subsection{PROMOTE-0122 - This was one of the most important structure decisions in the chat. It p}

\begin{itemize}
\item Source conversation: ui\_editor\_tool\_editor\_planning
\item Source file: docs/archive/conversations/ui\_editor\_tool\_editor\_planning/dominium\_ui\_editor\_tool\_editor\_planning\_\_human\_readable\_chat\_report.txt
\item Source heading: automatic decision extraction
\item Claim: This was one of the most important structure decisions in the chat. It prevented the initial implementation from becoming too broad while preserving the long-term goal. The UI Editor would be a bootstrap tool, not the final architecture.
\item Claim type: conversation\_claim
\item Proposed target: not\_selected
\item Authority domain: planning
\item Current repo conflict: not\_checked
\item Why it may matter: The claim recurs as a decision-like statement in the archived package.
\item Risks: Chat claims may be stale, superseded, or contradicted by current repo artifacts.
\item Required review: human review plus authority-order check
\item Recommended disposition: not\_reviewed
\end{itemize}

\subsection{PROMOTE-0123 - The user uploaded screenshot bundles for setup and launcher. **UNCERTAIN}

\begin{itemize}
\item Source conversation: ui\_editor\_tool\_editor\_planning
\item Source file: docs/archive/conversations/ui\_editor\_tool\_editor\_planning/dominium\_ui\_editor\_tool\_editor\_planning\_\_human\_readable\_chat\_report.txt
\item Source heading: automatic decision extraction
\item Claim: The user uploaded screenshot bundles for setup and launcher. **UNCERTAIN / UNVERIFIED:** The screenshots were not inspected in this chat. The assistant still produced a logical layout description, and the user accepted it as useful. Future work should inspect the bundles before claiming exact screenshot fidelity.
\item Claim type: conversation\_claim
\item Proposed target: not\_selected
\item Authority domain: planning
\item Current repo conflict: not\_checked
\item Why it may matter: The claim recurs as a decision-like statement in the archived package.
\item Risks: Chat claims may be stale, superseded, or contradicted by current repo artifacts.
\item Required review: human review plus authority-order check
\item Recommended disposition: not\_reviewed
\end{itemize}

\subsection{PROMOTE-0124 - The pasted discussion on trees, screws, pottery, axes, chairs, tables, a}

\begin{itemize}
\item Source conversation: Universe\_Explorer\_Planning
\item Source file: docs/archive/conversations/Universe\_Explorer\_Planning/Dominium\_Universe\_Explorer\_Planning\_Handoff\_\_01\_human\_readable\_report.md
\item Source heading: automatic decision extraction
\item Claim: The pasted discussion on trees, screws, pottery, axes, chairs, tables, and machines established that the player should manipulate material, geometry, constraints, processes, tools, stations, and affordances. The conclusion is that "item classes" must not be the truth substrate. Recipes and blueprints can exist, but as higher-order formalizations.
\item Claim type: conversation\_claim
\item Proposed target: not\_selected
\item Authority domain: planning
\item Current repo conflict: not\_checked
\item Why it may matter: The claim recurs as a decision-like statement in the archived package.
\item Risks: Chat claims may be stale, superseded, or contradicted by current repo artifacts.
\item Required review: human review plus authority-order check
\item Recommended disposition: not\_reviewed
\end{itemize}

\subsection{PROMOTE-0125 - A major theme was that useful local inventions must become portable, sta}

\begin{itemize}
\item Source conversation: Universe\_Explorer\_Planning
\item Source file: docs/archive/conversations/Universe\_Explorer\_Planning/Dominium\_Universe\_Explorer\_Planning\_Handoff\_\_01\_human\_readable\_report.md
\item Source heading: automatic decision extraction
\item Claim: A major theme was that useful local inventions must become portable, standardized, industrialized, and institutionally adopted. The repo now has a Formalization Chain spec and Player Desire Acceptance Map that strongly preserve this. The future work is making this playable: drafting, measuring, testing, blueprinting, certifying, teaching, manufacturing, maintaining, and revising designs.
\item Claim type: conversation\_claim
\item Proposed target: not\_selected
\item Authority domain: planning
\item Current repo conflict: not\_checked
\item Why it may matter: The claim recurs as a decision-like statement in the archived package.
\item Risks: Chat claims may be stale, superseded, or contradicted by current repo artifacts.
\item Required review: human review plus authority-order check
\item Recommended disposition: not\_reviewed
\end{itemize}

\subsection{PROMOTE-0126 - Workbench was treated as the likely first surface for the Universe Explo}

\begin{itemize}
\item Source conversation: Universe\_Explorer\_Planning
\item Source file: docs/archive/conversations/Universe\_Explorer\_Planning/Dominium\_Universe\_Explorer\_Planning\_Handoff\_\_01\_human\_readable\_report.md
\item Source heading: automatic decision extraction
\item Claim: Workbench was treated as the likely first surface for the Universe Explorer, but only if it remains projection, not authority. The repo queue says PRESENTATION-CONTRACT-01 is next, with PROJECTION-CONFORMANCE-01 as alternate. The conclusion was that presentation/projection contracts must come before building UI or renderer features.
\item Claim type: conversation\_claim
\item Proposed target: not\_selected
\item Authority domain: planning
\item Current repo conflict: not\_checked
\item Why it may matter: The claim recurs as a decision-like statement in the archived package.
\item Risks: Chat claims may be stale, superseded, or contradicted by current repo artifacts.
\item Required review: human review plus authority-order check
\item Recommended disposition: not\_reviewed
\end{itemize}

\subsection{PROMOTE-0127 - The user noted that earlier prompts were hard to copy because they were}

\begin{itemize}
\item Source conversation: Workbench\_AIDE\_Product\_Spine
\item Source file: docs/archive/conversations/Workbench\_AIDE\_Product\_Spine/Dominium\_Architecture\_Workbench\_AIDE\_Product\_Spine\_Planning\_\_01\_human\_readable\_report.md
\item Source heading: automatic decision extraction
\item Claim: The user noted that earlier prompts were hard to copy because they were not always contained in one code block, and that prompts should better handle dirty worktrees and concurrent tasks. From then on, generated prompts included detailed dirty worktree handling, allowed/forbidden paths, non-goals, validation, blocker classification, and commit/final-response formats.
\item Claim type: conversation\_claim
\item Proposed target: not\_selected
\item Authority domain: planning
\item Current repo conflict: not\_checked
\item Why it may matter: The claim recurs as a decision-like statement in the archived package.
\item Risks: Chat claims may be stale, superseded, or contradicted by current repo artifacts.
\item Required review: human review plus authority-order check
\item Recommended disposition: not\_reviewed
\end{itemize}

\subsection{PROMOTE-0128 - The final recommended sequence was: finish replay proof and barebones cl}

\begin{itemize}
\item Source conversation: Workbench\_AIDE\_Product\_Spine
\item Source file: docs/archive/conversations/Workbench\_AIDE\_Product\_Spine/Dominium\_Architecture\_Workbench\_AIDE\_Product\_Spine\_Planning\_\_01\_human\_readable\_report.md
\item Source heading: automatic decision extraction
\item Claim: The final recommended sequence was: finish replay proof and barebones client, run product spine review, then begin limited parallel dev. Larger parallel waves should wait until minimum AIDE workflow law exists.
\item Claim type: conversation\_claim
\item Proposed target: not\_selected
\item Authority domain: planning
\item Current repo conflict: not\_checked
\item Why it may matter: The claim recurs as a decision-like statement in the archived package.
\item Risks: Chat claims may be stale, superseded, or contradicted by current repo artifacts.
\item Required review: human review plus authority-order check
\item Recommended disposition: not\_reviewed
\end{itemize}

\subsection{PROMOTE-0129 - The user then requested this full preservation package so a new chat cou}

\begin{itemize}
\item Source conversation: Workbench\_AIDE\_Product\_Spine
\item Source file: docs/archive/conversations/Workbench\_AIDE\_Product\_Spine/Dominium\_Architecture\_Workbench\_AIDE\_Product\_Spine\_Planning\_\_01\_human\_readable\_report.md
\item Source heading: automatic decision extraction
\item Claim: The user then requested this full preservation package so a new chat could continue without re-explaining everything. This final handoff is itself part of the preservation output.
\item Claim type: conversation\_claim
\item Proposed target: not\_selected
\item Authority domain: planning
\item Current repo conflict: not\_checked
\item Why it may matter: The claim recurs as a decision-like statement in the archived package.
\item Risks: Chat claims may be stale, superseded, or contradicted by current repo artifacts.
\item Required review: human review plus authority-order check
\item Recommended disposition: not\_reviewed
\end{itemize}

\subsection{PROMOTE-0130 - The future relevance of this chat is high. It should feed directly into}

\begin{itemize}
\item Source conversation: World\_Architecture
\item Source file: docs/archive/conversations/World\_Architecture/Dominium\_World\_Architecture\_\_human\_readable\_chat\_report.txt
\item Source heading: automatic decision extraction
\item Claim: The future relevance of this chat is high. It should feed directly into the future project spec book and into a corrected implementation prompt. But it should not be treated as final implementation detail everywhere. Many high-level decisions are settled, but solver formulas, exact file encoding details, build system integration, actual repository layout, and some numerical representations still require verification.
\item Claim type: conversation\_claim
\item Proposed target: not\_selected
\item Authority domain: planning
\item Current repo conflict: not\_checked
\item Why it may matter: The claim recurs as a decision-like statement in the archived package.
\item Risks: Chat claims may be stale, superseded, or contradicted by current repo artifacts.
\item Required review: human review plus authority-order check
\item Recommended disposition: not\_reviewed
\end{itemize}

\subsection{PROMOTE-0131 - Before reading the rest of this report, remember this: this chat's centr}

\begin{itemize}
\item Source conversation: World\_Architecture
\item Source file: docs/archive/conversations/World\_Architecture/Dominium\_World\_Architecture\_\_human\_readable\_chat\_report.txt
\item Source heading: automatic decision extraction
\item Claim: Before reading the rest of this report, remember this: this chat's central contribution is not one isolated feature. It is a whole design philosophy for Dominium's world engine. The game should be deterministic, fixed-point, sparse, modular, content-driven, and field-based. Chunks are not the world. Chunks are just how the world is cached, meshed, streamed, and saved.
\item Claim type: conversation\_claim
\item Proposed target: not\_selected
\item Authority domain: planning
\item Current repo conflict: not\_checked
\item Why it may matter: The claim recurs as a decision-like statement in the archived package.
\item Risks: Chat claims may be stale, superseded, or contradicted by current repo artifacts.
\item Required review: human review plus authority-order check
\item Recommended disposition: not\_reviewed
\end{itemize}

\subsection{PROMOTE-0132 - This comparison was not a final technical dependency, but it helped clar}

\begin{itemize}
\item Source conversation: World\_Architecture
\item Source file: docs/archive/conversations/World\_Architecture/Dominium\_World\_Architecture\_\_human\_readable\_chat\_report.txt
\item Source heading: automatic decision extraction
\item Claim: This comparison was not a final technical dependency, but it helped clarify the design goal: Dominium should not simply copy any one of these systems. It should combine their useful elements while avoiding their limits.
\item Claim type: conversation\_claim
\item Proposed target: not\_selected
\item Authority domain: planning
\item Current repo conflict: not\_checked
\item Why it may matter: The claim recurs as a decision-like statement in the archived package.
\item Risks: Chat claims may be stale, superseded, or contradicted by current repo artifacts.
\item Required review: human review plus authority-order check
\item Recommended disposition: not\_reviewed
\end{itemize}

\subsection{PROMOTE-0133 - and XStack must serve development rather than development serving XStack}

\begin{itemize}
\item Source conversation: xstack\_lab\_galaxy
\item Source file: docs/archive/conversations/xstack\_lab\_galaxy/dominium\_xstack\_lab\_galaxy\_handoff\_\_human\_readable\_chat\_report.txt
\item Source heading: automatic decision extraction
\item Claim: and XStack must serve development rather than development serving XStack.
\item Claim type: conversation\_claim
\item Proposed target: not\_selected
\item Authority domain: planning
\item Current repo conflict: not\_checked
\item Why it may matter: The claim recurs as a decision-like statement in the archived package.
\item Risks: Chat claims may be stale, superseded, or contradicted by current repo artifacts.
\item Required review: human review plus authority-order check
\item Recommended disposition: not\_reviewed
\end{itemize}

\subsection{PROMOTE-0134 - After the transcript appeared, the user asked for a maximum-fidelity Con}

\begin{itemize}
\item Source conversation: xstack\_lab\_galaxy
\item Source file: docs/archive/conversations/xstack\_lab\_galaxy/dominium\_xstack\_lab\_galaxy\_handoff\_\_human\_readable\_chat\_report.txt
\item Source heading: automatic decision extraction
\item Claim: After the transcript appeared, the user asked for a maximum-fidelity Context Transfer Packet. The assistant produced one, with sections such as workstreams, decisions, tasks, constraints, open questions, risks, artifacts, and verification queue. Then the user asked to turn that into a downloadable report package. The assistant created Markdown/YAML files and a ZIP archive.
\item Claim type: conversation\_claim
\item Proposed target: not\_selected
\item Authority domain: planning
\item Current repo conflict: not\_checked
\item Why it may matter: The claim recurs as a decision-like statement in the archived package.
\item Risks: Chat claims may be stale, superseded, or contradicted by current repo artifacts.
\item Required review: human review plus authority-order check
\item Recommended disposition: not\_reviewed
\end{itemize}

\subsection{PROMOTE-0135 - Finally, the user asked not for another package, but for an in-chat read}

\begin{itemize}
\item Source conversation: xstack\_lab\_galaxy
\item Source file: docs/archive/conversations/xstack\_lab\_galaxy/dominium\_xstack\_lab\_galaxy\_handoff\_\_human\_readable\_chat\_report.txt
\item Source heading: automatic decision extraction
\item Claim: Finally, the user asked not for another package, but for an in-chat reader version of the package. The assistant rendered a structured in-chat reader. The user then asked for this current response: a human-readable, intuitive narrative report.
\item Claim type: conversation\_claim
\item Proposed target: not\_selected
\item Authority domain: planning
\item Current repo conflict: not\_checked
\item Why it may matter: The claim recurs as a decision-like statement in the archived package.
\item Risks: Chat claims may be stale, superseded, or contradicted by current repo artifacts.
\item Required review: human review plus authority-order check
\item Recommended disposition: not\_reviewed
\end{itemize}


\chapter{Appendix D - Full Contradiction Register}


{\small\raggedright SOURCE PATH: docs/\allowbreak{}archive/\allowbreak{}conversations/\allowbreak{}\_\allowbreak{}audit/\allowbreak{}CONTRADICTION\_\allowbreak{}MATRIX.md\par}


\section{Contradiction Matrix}

Findings are review queues, not resolved doctrine.


\begin{quote}\small Dense table content is summarized in this PDF rendering. See the Markdown source and HTML book for the full table.\end{quote}



{\small\raggedright SOURCE PATH: docs/\allowbreak{}archive/\allowbreak{}conversations/\allowbreak{}\_\allowbreak{}audit/\allowbreak{}STALENESS\_\allowbreak{}AND\_\allowbreak{}VERIFICATION.md\par}


\section{Staleness And Verification}

External platform, SDK, release, language-baseline, and implementation claims are stale until verified.


\subsection{Packages With Explicit Uncertainty}

\subsubsection{advanced\_simulation\_infrastructure}

\begin{itemize}
\item [UNCERTAIN / UNVERIFIED] No repository files were inspected in this chat. No code was implemented. Proposed paths, module names, command names, file names, VM concepts, and data schemas are design proposals until checked against the actual codebase.
\item [UNCERTAIN / UNVERIFIED] Exact solvers, fixed-point formats, module paths, data schemas, and tick ordering remain to be verified and specified.
\item [UNCERTAIN / UNVERIFIED] The exact gameplay rules runtime was not specified. The command examples were conceptual.
\item [UNCERTAIN / UNVERIFIED] The exact cross-section schema, attachment schema, and packing rules remain unresolved.
\item [UNCERTAIN / UNVERIFIED] Microsegment length, VM design, storage layout, budget values, and actual code integration remain open.
\item [UNCERTAIN / UNVERIFIED] Actual standards packs, rule schema, semantic key format, localization, and AI interaction remain unresolved.
\item Then it produced a prompt for another GPT-5.2 chat that already had a refactor/optimization plan. That prompt told the other chat to amend its existing plan instead of restarting, and to verify coverage of arbitrary placement, unified spatial primitives, co-location, signage, buildings, and determinism/performance.
\item [UNCERTAIN / UNVERIFIED] The architecture has not been verified against the actual repository. The exact implementation details remain unresolved: Q formats, orientation math, command schemas, VM instruction set, microsegment length, carrier ownership, junction archetypes, facility modules, and DECOR/device promotion policy.
\end{itemize}

\subsubsection{app\_runtime\_platform\_renderers}

\begin{itemize}
\item What remains uncertain is whether the current repository already enforces these boundaries mechanically. The actual code needs inspection.
\item The unresolved issue is sharding semantics. APP0 requires sharding hooks, but the chat did not decide whether shards are spatial, logical, temporal, jurisdictional, or something else. The report preserved the idea that APP0 can provide sharding plumbing without inventing simulation semantics.
\item The unresolved issue is the trust and integrity model for setup. The chat did not decide whether content verification should use hashes, signatures, local manifests, network verification, or another model.
\item The unresolved issue is the exact elevation/audit model. The chat did not decide how tools authenticate, request write permissions, log actions, or interact with law enforcement modules.
\item The unresolved issue is the exact renderer capability schema. The chat proposed ideas but did not finalize fields, priorities, or target matrices.
\item The user asked to look at the old code snapshot in project attachments. A prior assistant claimed to inspect it, but later preservation documents marked that claim as **UNCERTAIN / UNVERIFIED**.
\item This became one of the most important unresolved issues. Almost every implementation question depends on the real code: directory layout, build system, languages, existing platform/render layers, client/server dependencies, and engine/game public APIs. The next real step is repository verification.
\item INFERENCE:** The user wants future assistants to be epistemically careful: not inventing repository facts, not treating proposals as decisions, and not losing uncertainty.
\end{itemize}

\subsubsection{app\_testx\_codehygiene}

\begin{itemize}
\item The main unresolved goal is practical execution: which prompts have actually been run, which files exist, and what implementation should happen next.
\item Setup handles install, verify, repair, rollback, upgrade, downgrade, uninstall, and status. Launcher can invoke Setup, but must not replicate its mutation logic. This prevents two products from disagreeing about installation state.
\item First, verify whether libs/contracts and libs/appcore already exist. Then implement the pure contract headers and appcore skeleton. This should include:
\item Before implementing app RepoX integration, verify actual RepoX metadata paths, TestX invocation, VALIDATE-0 commands, and BUILD-ID-0 files. Without that, an implementation prompt would guess.
\item Use labels like FACT, INFERENCE, UNCERTAIN, PROJECT-CONTEXT in reports.
\item From the user profile and instructions, future assistants should verify time-sensitive/current facts with web where relevant. This chat itself is mostly internal project planning, so external citations are not central.
\item EXEC, ECSX, KERN, ADOPT, DIST, HWCAPS, EXIST, DOMAIN, TRAVEL, TIME, OMNI, LIFE, CIV, WAR, AGENT, TOOL, MOD, FINAL prompt families were generated earlier. They are important as historical design artifacts, but actual repo execution is unverified.
\item The chat produced downloadable handoff files and later an in-chat reader. Those files are useful for aggregation but are not themselves project source code. The package's main caveat is that it does not verify actual repository state.
\end{itemize}

\subsubsection{architecture\_codex\_prompts}

\begin{itemize}
\item UNCERTAIN / UNVERIFIED: No repository code was inspected or modified in this chat. The prompts and architecture are plans and artifacts, not proof that implementation exists. The actual repository state, build system, existing code quality, GUI backend availability, TLV schemas, and whether any prompts have been applied remain unverified.
\item Uncertainty remains about actual repository code and build system, but the layering itself is a strong established principle.
\item Unresolved issues include exact TLV schema field tags, exact registry implementation, base pack structure, and whether engine comments/docs may use examples with domain words.
\item The main uncertainty is implementation: the prompts define it, but no code was verified.
\item Unresolved details include exact process IO schemas, machine runtime state, and the relationship between autonomous machines and agent-operated jobs.
\item Unresolved: exact nesting policy, allowed packing modes, and how far to model real packing constraints.
\item Uncertainty remains about freeform geometry versus grid shell as the long-term building representation.
\item The exact destruction and structural-load models remain unresolved and part of Path D.
\end{itemize}

\subsubsection{architecture\_ui\_providers}

\begin{itemize}
\item The exact first playable slice remains unresolved. The first provider wedge is clear, but exact sequence between client/workbench gameplay, Robot OS, and gameplay mechanics still needs a formal milestone. The exact Lua version, Linux baseline, current repo status, and external library/version support need verification. The exact degree of Unreal involvement remains unclear.
\item 1. Verify and pin external library/toolchain facts.
\item It remains uncertain how much the user wants Unreal in the near-term after raylib-first discussion. It also remains uncertain whether Lua 5.4 or 5.5 is preferred; the user appears to value pinning and stability more than "latest" for script ABI.
\end{itemize}

\subsubsection{Build\_and\_Future\_Proofing}

\begin{itemize}
\item Plain-language limitation: this package is reliable as a preservation of the current visible chat. It should not be treated as a complete archive of every Dominium discussion in other chats. Where the report uses project or repository context, it is labelled as such. The most important caveat is that many items are recommendations, not final user decisions.
\item The unresolved goals are implementation goals: deciding which recommendations become canon, adding public-surface/dependency/build contracts, validating the current live repo state, and applying the proposed cleanup tasks. The preservation package is complete, but the engineering work it describes remains pending.
\item 2. **Verify current repo state.** Confirm the actual current HEAD, build proof, smoke tests, and outstanding warnings before acting.
\item 10. **Feed this report into the master spec book aggregator.** Preserve uncertainty labels and avoid merging suggestions as decisions.
\item The preservation output must distinguish FACT, INFERENCE, UNCERTAIN/UNVERIFIED, and PROJECT-CONTEXT.
\item A future assistant may rely on stale repo state. Mitigation: verify live HEAD, docs, and CI before implementation.
\item What live repo facts should I verify before implementing these tasks?
\item The final uploaded prompt asked for this preservation package. The key thing to preserve is that most architecture proposals are recommendations, not accepted user decisions yet. The next best action is to decide which recommendations become canon, then implement either STRUCTURE-01: Public Surface Registry or the build tuple contract work, after verifying the current live repo state.
\end{itemize}

\subsubsection{canonical\_structure\_and\_framework}

\begin{itemize}
\item INFERENCE: The user wanted a durable method for deciding future structure changes, not just a one-time layout. They also wanted the assistant to become more rigorous: distinguish decisions from brainstorms, stop over-planning without execution, and preserve uncertainty. The user also wanted a structure that supports other Domino-based games and possibly different engine implementations.
\item A persistent uncertainty is the exact current repo state. The user supplied many status reports and directory exports, but several exports were internally inconsistent or stale at various points. The chat therefore repeatedly emphasized structure report integrity and verification before trusting tree analyses.
\item 1. Verify fast-strict/RepoX/structure report status from a clean current export.
\item The user dislikes vague reassurance and wants uncertainty labelled.
\item UNCERTAIN: The final exact convention for external/upstream versus external/vendor should be checked against current repo policy.
\item UNCERTAIN: Whether Lua should be pinned to 5.4 or 5.5 is not settled in this chat.
\item UNCERTAIN: Whether pack-internal content/ directories are legitimate pack law or legacy remains to be verified.
\item User-reported commits and status summaries: these include commits like 6e0dd93, 1406490, 3243fab, ce9ca, and others. Treat them as FACTs reported in the chat, but verify live repo state before acting.
\end{itemize}

\subsubsection{Chronology\_Celestial\_Systems}

\begin{itemize}
\item The assistant formalized this as the Perfect Earth Calendar, later called HPC-E. The final structure is clear, though the exact leap rule remains unresolved.
\item What remains uncertain is the exact verified object list. The assistant generated a real Milky Way system list, but that list must be checked before becoming data.
\item The unresolved issue is exact fidelity. The chat specified what should exist conceptually, but not the final data schema, map resolution, object attributes, or implementation priority.
\item What remains unresolved is the exact mathematical implementation of these frames and how much relativistic fidelity should exist in phase one.
\item The uncertainty is high relative to Earth. Many names and structures were assistant-generated and should be reviewed before implementation.
\item The unresolved issue is implementation reality. The actual Plan G/The Game codebase was not visible here, so any implementation plan must be reconciled with real architecture.
\item The only major unresolved part is the leap rule.
\item The user wants rigorous, structured, fact-aware assistance. The user specifically requested preservation of uncertainty labels and does not want assistant suggestions treated as user decisions.
\end{itemize}

\subsubsection{development\_routes}

\begin{itemize}
\item The first answer was assertive. It stated that Dominium must follow Route C and described the route as the only viable one for Dominium. Later parts of the chat corrected the status of that claim. FACT: the assistant made that recommendation. UNCERTAIN / UNVERIFIED: the recommendation was not validated with external sources in the chat and was not explicitly accepted by the user.
\item The assistant produced a large Context Transfer Packet. It preserved the technical route proposal, the caveats around its uncertainty, the proposed phase order, the open questions, the risks, the artifacts, and the next actions.
\item The assistant created those files and reported that the chat label was inferred as **Dominium Development Routes**. The package was said to be safe for aggregation **with caveats**. The main caveats were that the substantive project transcript was short, Route C was not user-confirmed, and Dominium's genre, core loop, stack, implementation state, files, platforms, team, and timeline were not established.
\item UNCERTAIN / UNVERIFIED: Route C was not accepted by the user as final.
\item UNCERTAIN / UNVERIFIED: the required level of determinism for Dominium was not established. It might need full cross-platform deterministic replay, or it might only need ordinary save/load correctness, or something between those extremes.
\item The outputs included a Context Transfer Packet and a downloadable package with multiple files. The files are less important than the purpose: preserve the exact status of the discussion, including uncertainty and assistant-proposal status, so future assistants do not accidentally distort the project.
\item The second explicit goal was to preserve the chat before retirement. The user did not want a normal summary. They wanted a maximum-fidelity transfer packet that would prevent loss of decisions, constraints, unresolved issues, rationale, and next actions. This was addressed by the Context Transfer Packet.
\item Some goals changed over time. The chat began as architecture guidance, became a state-transfer task, became a packaging task, became an in-chat inspection task, and now becomes a plain-English briefing task. The underlying continuity goal stayed consistent: preserve the useful substance without losing uncertainty.
\end{itemize}

\subsubsection{documentation\_standards\_readme}

\begin{itemize}
\item The key thing to remember: this chat produced the standards and prompts for future work; it did not verify or change the project. The next assistant must inspect the actual repository before writing repo-specific docs, README content, platform claims, license claims, or subsystem descriptions.
\item The chat began with the user asking how to have Codex generate documentation for every file in the project. The user wanted metadata, functions, definitions, dependencies, and other maintainability information documented. The core uncertainty was placement: should documentation live in separate sister files, or should it live in file headers and function headers?
\item UNCERTAIN / UNVERIFIED: Whether the project has similarly named files or modules is unknown.
\item What remains uncertain is how the actual project currently organizes public headers, internal headers, and docs. That must be verified from the repository.
\item What remains unverified is the actual project name, supported platforms, license, maturity, and current repository layout. Those cannot be responsibly written without inspecting the repo.
\item The decision depends on an assumption: that the project's public APIs are declared in headers. That is likely for a C/C++ codebase, but the actual public/private boundary is still unverified.
\item What must happen first: verify source roots, exclusions, CMake structure, CI environment variables, and Python availability.
\item What could go wrong: treating assistant suggestions as final requirements or merging contradictory chat claims without preserving uncertainty.
\end{itemize}

\subsubsection{Dominium\_Architecture\_I}

\begin{itemize}
\item What remains uncertain is the exact final scope of all simulation systems. Many systems were described in ambitious terms, but not all details are final or internally consistent. The systems layer especially needs canonicalisation before coding.
\item What remains uncertain is the current real-world status and capabilities of "Codex 5.1 Max." That is an external tool fact and requires verification before future use.
\item The key later decision was to skip top-level .txt devspec files because original Markdown files already existed in /docs/.... The visible chat does not show the full content of those docs. They should be preserved as referenced artifacts, but their exact content is **UNCERTAIN / UNVERIFIED**.
\item This affects dlocale, font rendering, platform path handling, data locale JSON files, mod APIs, UI, and documentation. The unresolved part is exactly how much retro support is actually feasible. Claims about old OSes, toolchains, SDL2, OpenGL, Direct3D, and macOS versions require external verification.
\item A large portion of the chat consisted of producing specs for those files. Each spec was meant to tell Codex how to implement that file. The conclusion is that the repository architecture is broad but structured. The unresolved issue is that many later files remain pending.
\item The unresolved issue is API consistency. Generated specs for memory and serialization are inconsistent across files, and some platform/render specs contain outdated or unverified external assumptions.
\item The most important unresolved issue is duplication and contradiction. dweather, dhydro, and dai\_core were generated in more than one form. The package does not choose a final version. A future assistant should not implement any of those blindly.
\item The user decided to skip the top-level .txt files because original Markdown docs already exist in /docs/.... This decision is final for the current workflow, but it depends on those docs actually existing and being current. That is **UNCERTAIN / UNVERIFIED** because the docs were not inspected in this chat.
\end{itemize}

\subsubsection{Dominium\_Architecture\_II}

\begin{itemize}
\item What remains uncertain is the exact code-level API and component layouts. The chat produced specs and prompts, not final source code.
\item This is an important implementation detail, but some of it remains unverified. Lua 5.4.4 portability to old compilers needs testing. The exact Lua data schemas still need definition. The docs also still have JSON-oriented format sections, so the Lua-vs-JSON policy needs reconciliation.
\item The biggest unresolved goals are practical:
\item Verify DX9/Windows 2000, SDL2, Lua/toolchain compatibility.
\item FACT: This was answered in response to Codex's clarifying questions. It is a concrete implementation direction, though Lua 5.4.4 portability remains unverified.
\item 1. **Verify and apply V4 docs.
\item The user repeatedly requested maximum-fidelity reports, registers, handoff packets, and human-readable summaries. The user wants labelled uncertainty and does not want future assistants to invent missing context.
\item or forget that Lua-vs-JSON remains unresolved.
\end{itemize}

\subsubsection{Dominium\_Architecture\_III}

\begin{itemize}
\item The user then uploaded dominium.7z, saying it was the git repository as of that moment. **FACT:** The prior assistant said it could not open .7z archives in that environment. That means the actual repository state remained unverified at that stage. This became important later, because many implementation ideas were discussed without confirmed access to the repo contents.
\item This created one of the main unresolved issues. DX12 was discussed earlier, but omitted from the latest renderer list. X11 and Wayland were discussed earlier, but the latest platform categories are POSIX, SDL1, SDL2, and Native. Therefore, the final status is:
\item DX12 is **UNCERTAIN / UNVERIFIED** and should not be assumed.
\item X11/Wayland placement is unresolved.
\item The prior assistant summarized these files as containing two different launcher concepts: an engine-rendered launcher view and a native Win32 launcher. It also claimed the CMake file built the Win32 source for non-Windows targets. However, this report must preserve that as **UNCERTAIN / UNVERIFIED**, because the files were not re-inspected during final packaging.
\item This is final as a user-stated product goal. Feasibility still requires verification.
\item Those outputs are important artifacts, but they are not the substance of the design. Their purpose is to preserve this chat's decisions, uncertainty, and future work for later aggregation into a master Project Spec Book.
\item What remains uncertain is not the constraint itself, but how existing repository code currently enforces it. That requires code inspection.
\end{itemize}

\subsubsection{Dominium\_Architecture\_IV}

\begin{itemize}
\item The conclusion was clear and accepted: **Domino is the reusable engine; Dominium is the game/product layer**. What remains uncertain is the actual repository state. The chat planned the split, but did not verify whether the repo already matches it.
\item The unresolved issue is that the actual repo was not inspected in this chat. The planned structure may not match the current files.
\item The unresolved issue is practical feasibility, especially for retro targets and Carbon OS. Their actual toolchains and system calls remain unverified.
\item The unresolved issue is the exact IR payload layout and how far limited renderers should go. A CGA or MDA backend cannot realistically support modern 3D in the same way as Vulkan. The agreed approach is explicit capability flags and graceful fallback or no-op behaviour.
\item The unresolved issue is the exact manifest/schema format. JSON, TOML, and INI-like examples appeared in prompts, but no final schema choice was made.
\item The unresolved issue is that external packaging details require verification. WiX, macOS notarization, Linux packaging, and AppImage tooling are all external-world facts and may be stale.
\item The package and this report should help future assistants continue without re-reading the whole chat. However, they must preserve uncertainty: actual repo state is unverified, some transcript sections were skipped, external tooling facts may be stale, and assistant suggestions are not user decisions unless accepted or built upon.
\item It is also reasonable to infer that the user wants the future project spec book to preserve rationale, not just lists of decisions. The later handoff requests emphasised visible rationale, rejected options, uncertainty, and changes of direction.
\end{itemize}

\subsubsection{Dominium\_Complete\_Conversation}

\begin{itemize}
\item Uncertainty: the raw pasted prompt is not the current repo authority unless materialized in repo canon. The later audit found that it had in fact been materialized under docs/canon/constitution\_v1.md.
\item Uncertainty: during this preservation pass, a potential inconsistency appeared: some docs claim product roots moved under apps/, while an earlier inspected code path under root client/ was available and apps/client was not fetched successfully. This must be verified against the current physical tree.
\item Unresolved goals include verifying the latest physical repo tree, proving runtime implementation maturity, formalizing public ABI/API boundaries, completing layout/naming cleanup, and building a second Domino-based product to prove reuse.
\item 1. Verify the current physical repo tree against layout docs and contracts/repo/layout.contract.toml.
\item The user likely wants future assistants to challenge weak framing, preserve uncertainty, avoid shallow "best practice" lists, and focus on decisions that survive future rewrites.
\item "Create a gated plan for verifying current repo layout and exceptions."
\end{itemize}

\subsubsection{dominium\_domino\_codex\_planning}

\begin{itemize}
\item The most important things to remember are: the project's strict architecture constraints, the Milestone-0 prompt sequence, the later shared CLI/TUI/GUI/input/render prompt sequences, the missing pack-system prompt, the unified startup policy, the unresolved repo-verification issue, and the need to treat generated prompts as plans rather than proof of implementation.
\item The assistant generated the first two prompts. The third was never produced because the conversation changed direction. That missing third prompt remains an important unresolved item.
\item However, the later context transfer packet treated that inspection claim as **UNCERTAIN / UNVERIFIED**. This was careful and correct: within the visible chat, there was no reliable tool trace proving the repo had actually been inspected. Therefore, future assistants must verify dominium.zip or the active checkout before relying on the platform/render/API answer.
\item What remains uncertain is how much of this architecture already exists in the current repo and how much is still planned. The chat produced prompts and plans, but it did not prove execution.
\item The unresolved issue is whether those prompts were ever run and whether the generated code, if any, conforms to the project's strict C89/C++98 and determinism requirements.
\item The unresolved complication is C89 portability. Standard C89 does not provide long long, but u64, i64, and q48\_16 require 64-bit representation. Future work must decide whether to allow compiler extensions, emulate 64-bit values for strict retro targets, or define platform tiers with conditional support.
\item The unresolved issue is that the prompt's minimal manifest parser and hardcoded platform path are only a starting point. A real implementation must generalize product discovery, avoid duplicated path logic, and preserve compatibility/versioning rules.
\item The unresolved issue is whether the prompts are too broad and whether the existing repo already has overlapping modules. Future work should inspect the repo before applying these prompts.
\end{itemize}

\subsubsection{dominium\_full\_conversation}

\begin{itemize}
\item Verify live repo state, then review/run FULL-GATE-LEGACY-TEST-ROUTE-01, then generate PACK-INTERNAL-LAYOUT-CANON-01 if the queue/status supports the maintenance lane.
\item Preserve uncertainty labels.
\item Future chats could easily misunderstand this conversation by treating Workbench as a GUI framework, treating raylib as the engine, resuming broad structure cleanup, skipping projection conformance, or assuming live repo status from stale pasted reports. The mitigation is to verify repo state first, preserve the contract/service/projection/provider distinctions, and follow the task queue.
\item The best next action is to verify live repo state, then review or run FULL-GATE-LEGACY-TEST-ROUTE-01, then generate PACK-INTERNAL-LAYOUT-CANON-01 if status supports it.
\end{itemize}

\subsubsection{dominium\_setup}

\begin{itemize}
\item FACT: The setup system was repeatedly defined as the only component allowed to install files, modify system/installation state, own installed-file metadata, repair installs, verify artifacts, uninstall, downgrade, upgrade, and roll back. This is the most important architectural rule from the chat.
\item What remains uncertain is the actual current implementation state in the repository after the applied Codex prompts. The design is clear, but the repo must still be inspected.
\item FACT: The final canonical setup layout uses setup/core/\{fetch,verify,install,rollback\}, setup/include/\{dsk,dsu\}, setup/packages, setup/platform, setup/tests, and setup/ui.
\item The unresolved part is whether the actual repository now matches this canonical layout.
\item UNCERTAIN / UNVERIFIED: Exact Visual Studio 2026 behavior and toolchain details should be verified in the current environment before relying on them.
\item UNCERTAIN / UNVERIFIED: The user stated the prompts were applied, but this chat did not show the final repository tree or build results.
\item UNCERTAIN / UNVERIFIED: The actual schema files under schema/setup/ need to be checked.
\item Several goals remain unresolved because the actual repo was not inspected here. We do not know whether the applied Codex prompts fully succeeded, whether setup builds, whether schemas exist, whether libs/contracts is complete, or whether setup/launcher CLI smoke tests pass.
\end{itemize}

\subsubsection{Domino\_Dominium\_Workbench}

\begin{itemize}
\item Reliability note:** This is a human-readable consolidation, not a live repository audit. Repo-current statements that came from pasted material or prior-chat excerpts remain **UNVERIFIED** until checked against the live julesc013/dominium repo and current AIDE/validator outputs.
\item Verify current live repo status and task gate state.
\item > Verify repo state, then generate COMMAND-RESULT-VIEW-SLICE-01.
\end{itemize}

\subsubsection{domino\_engine\_refactor\_prompts}

\begin{itemize}
\item The answer also introduced a deterministic bytecode VM as the preferred way to support modded behavior. A question was raised about whether native-code plugins should also be allowed, but the user did not answer. That remains unresolved.
\item What remains uncertain is the actual repository state. The chat designed architecture and prompts, but did not inspect source code. Therefore the plan needs repo verification before implementation.
\item The unresolved work is to generate detailed subsystem implementation prompts for heat, power, fluids, atmospheres, vehicles, and structural loads after the core engine spine is implemented.
\item The remaining uncertainty is plugin policy. A deterministic VM was proposed, but whether native-code plugins are allowed remains unresolved.
\item What remains uncertain is the exact thresholds and content policies for when things promote or demote.
\item The unresolved issue is how much of STRUCT compilation should be implemented first. The prompt plan is broad; a real implementation may need a smaller first milestone.
\item ignoring FACT/INFERENCE/UNCERTAIN labels;
\item UNCERTAIN / UNVERIFIED:** The package files were generated in the sandbox during the prior step, but repository files were not created or inspected.
\end{itemize}

\subsubsection{engine\_baseline\_architecture}

\begin{itemize}
\item Uncertainty: the repo has many modules and some historical/legacy surfaces, so exact implementation maturity varies by subsystem. The distinction should remain a formal requirement in a master spec.
\item Uncertainty: the full determinism envelope across all current modules was not independently re-run in this chat. User-supplied local validation results should be preserved but verified before being used as audit proof.
\item Uncertainty: this remains strategic architecture, not an implemented decision.
\item Unresolved goals include:
\item User wants uncertainty labelled and framing corrected when evidence disagrees.
\item UNCERTAIN: whether the user wants Unreal integration at all as a client/editor.
\item UNCERTAIN: exact tolerance for using external geometry libraries.
\item UNCERTAIN: whether broad repo restructuring will be desired after Milestone 0.
\end{itemize}

\subsubsection{Foundation\_Workbench\_Codex}

\begin{itemize}
\item What remains uncertain is how quickly this architecture will move from contracts and fixtures into runtime implementations and product slices. The next chat should not assume runtime provider resolvers, package runtime, or Workbench shell exist.
\item verify latest live repo state after user-pasted Wave 1 completion;
\item 1. Verify live repo state.
\item The user wants direct, evidence-grounded, audit-ready answers; timestamps/model labels; explicit uncertainty; no repeated clarification unless necessary; large continuous Codex prompt blocks when generating prompts; targeted tests instead of full suites; and rapid progress toward Workbench/code.
\item The user is likely to reject slow, report-only, overcautious process. Future assistants should produce executable next prompts quickly after verifying state.
\item It is uncertain how many concurrent Codex workers the user will actually run at once and whether the user wants coordinator prompts or worker prompts next.
\item The most important continuation is: verify current origin/main, confirm Wave 1 and Workbench validation state, then generate or run COMMAND-RESULT-VIEW-SLICE-01.
\end{itemize}

\subsubsection{Framework\_Open\_Source\_Provider}

\begin{itemize}
\item The chat inspected julesc013/dominium and found evidence of existing abstractions such as system and graphics layers, stubs, and soft-backed backends. The finding was that raylib could fill concrete provider gaps without replacing the architecture. However, current repo facts are stale/uncertain and must be verified, especially the C17/C++17 vs C90/C++98 baseline contradiction.
\item The exact first implementation branch, versions, and dependency pins are unresolved. The exact Domino Framework ABI is unresolved. The exact Lua version is unresolved. The current repository baseline must be verified. The formal policy for GPL/LGPL/unclear license material is unresolved. The sparse simulation and CAD systems need formal specs and prototypes.
\item The raylib decision was also directionally accepted. The user repeatedly expressed enthusiasm for raylib and asked about using as much of raylib infrastructure as possible. The caveat is that raylib is a provider suite, not architecture. This affects rendering, audio, input, asset preview, and Workbench bootstrap.
\item 1. Verify repo baseline and dependency versions.
\item The user requires source-grounded, audit-ready, uncertainty-labelled responses. The preservation prompt explicitly requires FACT/INFERENCE/UNCERTAIN/PROJECT-CONTEXT labels, no invented facts, no treating brainstorms as decisions, and no over-compression. The prompt also requires a human-readable report first and downloadable files if possible ?filecite?turn29file0?.
\item The most blocking unresolved issues are repo baseline verification, exact dependency pins, minimal framework ABI, provider profile order, deterministic simulation spec, and license policy.
\item Verify the current julesc013/dominium CMake language baseline.
\end{itemize}

\subsubsection{gui\_binary\_content}

\begin{itemize}
\item GPT-5.5 Pro - 2026-05-27 00:00:00 Australia/Melbourne; time UNVERIFIED
\item whether knowledge systems should include uncertainty, false beliefs, or lost knowledge,
\item This stage established a pattern that continues throughout the chat: the user does not want impressive-looking prompts that hide unresolved design choices. The user wants the assumptions exposed before anything is formalized.
\item The assistant agreed with the general direction, but added caveats. The most important were:
\item Those caveats were not final user decisions, but they are important warnings.
\item The conclusion was not "run the prompt now." The conclusion was that CONTENT0 is a strong foundation but needs discussion before finalization. The assistant-generated rewrite is useful but not final. The unresolved issues need to be resolved before Codex or any implementation work touches the repository.
\item What remains uncertain is how exactly to encode these ideas in schemas and data. For example, it is not yet decided how deterministic seed hierarchies work, how macro content refines into micro content, or how timeline causality should be represented.
\item The unresolved issue is the shared backend/UI contract. Without that, the GUI plan is only a list of technologies.
\end{itemize}

\subsubsection{Language\_Platform\_Architecture}

\begin{itemize}
\item The 32-bit question followed. The answer was to make full native products 64-bit and keep 32-bit as constrained or projection support only. The key caveat was not to make native pointer size part of game law. Save, replay, network, IDs, and renderer IR must use fixed-width types and never serialize size\_t, long, uintptr\_t, or raw pointers.
\item Exact tolerance for exceptions/RTTI inside private C++17 code remains unsettled. Exact platform floors and raylib/SDL2 priority remain unresolved.
\item Verify the Windows 7 and macOS 10.9 toolchain/library assumptions using official sources.
\end{itemize}

\subsubsection{launcher\_app\_layer}

\begin{itemize}
\item These prompts matter because they superseded earlier advice. They also created a new practical priority: verify that the applied prompts actually succeeded.
\item This changed the correct future behavior. The next assistant must not redesign. It must implement, audit, maintain, document, verify, and harden.
\item This matters because it prevents platform-specific drift. If CLI has verify, GUI must not implement a different hidden version of verify. If GUI has a pack enable flow, it should map to the same command/intent as CLI and TUI.
\item The final purity prompt required those to be moved or quarantined. However, actual repo success remains unverified in this chat.
\item The important change is that the conversation moved from "What architecture should we use?" to "Architecture is locked; how do we verify and harden the implementation?"
\item The largest unresolved goal is verifying the actual repo. The user stated that two Codex prompts were applied, but this chat did not independently verify the filesystem, build, tests, scripts, or launcher docs.
\item Another unresolved goal is confirming whether the launcher hardening prompt has been applied.
\item A third unresolved goal is confirming whether command graph, UI IR, binding validation, zero-pack tests, RepoX changelog display, and BUILD-ID mismatch refusal are implemented.
\end{itemize}

\subsubsection{Launcher\_Setup\_Architecture}

\begin{itemize}
\item Uncertainty:** The exact repository layout and existing implementation were not inspected in this chat.
\item Uncertainty:** The final language boundary for setup remains unresolved after the later C89/dsys/dgfx launcher direction.
\item Uncertainty:** The exact mapping into runtime binaries remains unverified.
\item Uncertainty:** Accounts/social/wiki/forum/direct messaging are future-facing and were not converted into the final five-tab implementation plan.
\item Uncertainty:** Earlier setup designs used JSON manifests. Later dsys/dgfx architecture leaned toward TLV .dmeta files.
\item Uncertainty:** Actual plugin ABI details and dynamic loading support depend on dsys/repo APIs.
\item Uncertainty:** Actual dsys/dgfx APIs were not inspected.
\item Conclusion reached:** This chat is now documented and can be merged later, but with caveats.
\end{itemize}

\subsubsection{Modularity\_AIDE\_Refactorability}

\begin{itemize}
\item The live repo still needs verification. The exact AIDE files, contracts, validators, and migration tasks must be implemented. It remains unresolved which existing roots and tools map to which final homes, which APIs become stable, which code is reusable across products/games/projects, and which prior assistant recommendations should become formal requirements after user review.
\item 1. **TASK-09** - Verify live repo baseline: current branch, root inventory, validators, CTest status, generated artifacts.
\item The user prefers source-grounded, audit-ready decisions; bounded uncertainty; explicit task prompts; and reusable engineering doctrine. Future assistants should avoid shallow affirmation and should distinguish accepted user decisions from assistant recommendations.
\item Important caveat: the only current uploaded file available for this task is Pasted text.txt, which contains the preservation-package prompt. Earlier generated handoff files, ZIP packages, or uploaded docs referenced in prior conversation are not available in this chat unless the user re-uploads them.
\item The most serious risk is that a future assistant treats this chat as if it verified the live repo. It did not. The second serious risk is that it hardens every assistant suggestion into a requirement. Several items are strong recommendations aligned with user goals, but still require implementation review. Future chats must preserve labels and verify stale facts.
\item Future assistants must verify live repo state before implementing cleanup.
\end{itemize}

\subsubsection{omega\_xi\_pi\_architecture\_future}

\begin{itemize}
\item Preserve FACT / INFERENCE / UNCERTAIN labels.
\item Do not invent or flatten uncertainty.
\item 7. "List all artifacts I should verify in GitHub main."
\item 9. "What should I verify before starting Sigma?"
\end{itemize}

\subsubsection{OS\_Interface\_Repo\_Architecture}

\begin{itemize}
\item The unresolved goals are implementation-level. The repo still needs command dispatch unification, product boot proof, portable projection proof, AppShell rendered-mode law update, Workbench module contracts, document/patch/result/refusal schemas, and a boot-to-replay MVP. It also needs continued exception retirement, CTest/RepoX remediation, and verification of current build status.
\item 2. Repair canonical verify test discovery.
\item The user wants uncertainty labels and correction of incorrect framing.
\item This preservation task specifically required FACT / INFERENCE / UNCERTAIN / PROJECT-CONTEXT labels, stable IDs, registers, spec sheets, an aggregator packet, self-audit, and downloadable files if tools are available.
\item UNCERTAIN: Whether the user wants to rename engine conceptually to kernel in code, or only use kernel as a conceptual layer.
\item UNCERTAIN: Whether the final public product name should be Dominium Workbench, Dominium Studio, Domino Workbench, or something else.
\item UNCERTAIN: How aggressive the next actual code refactor should be versus staged adapters and codegen.
\end{itemize}

\subsubsection{platform\_renderer\_api\_plan}

\begin{itemize}
\item This is not a fatal problem for the prompt pack, because the prompts themselves repeatedly tell Codex to discover files with rg, tree, type, CMake commands, and target listing. But it is a major caveat for any future assistant: do not assume the actual repo matches every file path or API name unless checked.
\item UNCERTAIN / UNVERIFIED:** The conversation did not prove that these prompts were actually applied to the repository. They are an assumed prerequisite, not verified code state.
\item Uncertainty remains about actual repo state. The prompts are designed to discover the repo state, but the chat itself did not inspect the archive.
\item The final active prompt pack narrowed runtime completion to Tier-1 on Windows and compile-gated correctness for Tier-2. This is a practical compromise because Codex runs on Windows. But it also creates a caveat: if "fully complete" later means fully running X11/Wayland/Cocoa/Metal/Vulkan, more prompts or native CI validation will be required.
\item This became one of the final "done" gates, along with scripts\textbackslash{}build\_codex\_verify.bat.
\item This matters because the conversation was long and contained multiple plans, superseded options, caveats, and active artifacts. The report package and this plain-language report exist to prevent future assistants from restarting, re-asking, or losing distinctions between decisions, brainstorms, and uncertainties.
\item The user also seems to want future assistants to operate with high continuity and not lose subtle distinctions. This is inferred from the repeated requests for maximum-fidelity transfer, stable IDs, labels, uncertainty tracking, and report packaging.
\item Another unresolved goal is whether the master engine prompts 1-14 are truly implemented. The later plan depends on them, but this chat did not verify the repo.
\end{itemize}

\subsubsection{platform\_support}

\begin{itemize}
\item UNCERTAIN / UNVERIFIED: The name "Domino" was never confirmed by the user. Future work should treat "Domino" as an assistant-created placeholder unless the user explicitly accepts it. The useful idea is the distinction between full product support and reduced engine/core support, not necessarily the name.
\item The main unresolved issue is the exact definition of support. Does Tier-0 mean identical feature parity, simultaneous release, identical save compatibility, identical UI, equal QA coverage, or some platform-specific differences? The assistant inferred that Tier-0 should mean feature parity, full QA, long-term support, and high release priority, but the exact contractual meaning still needs formal definition.
\item What remains unresolved is the exact PC baseline. The chat did not decide whether Windows 10, Windows 11, or both are required. It did not decide whether Windows ARM64 matters. It did not decide whether macOS support includes Intel Macs, Apple Silicon Macs, or both. It did not define Linux distributions, glibc requirements, Flatpak/AppImage/deb/rpm packaging, X11/Wayland policy, or Steam Deck verification.
\item UNCERTAIN / UNVERIFIED: The chat did not establish minimum Android API level, ABI list, target graphics APIs, Play Store rules, Android TV status, Android Go status, Automotive status, Chromebook status, or vendor ROM support. These need future decisions and current-source verification.
\item The user also listed iPod Touch, Apple TV, Apple Watch, Mac, macOS, tvOS, watchOS, and Apple Vision Pro indirectly through AR. These were not classified in detail. INFERENCE: macOS belongs under PC. UNCERTAIN / UNVERIFIED: tvOS, watchOS, and visionOS remain undecided. Full Dominium on Apple Watch is not established and should not be assumed.
\item The unresolved issues are substantial. The chat did not decide whether WebGPU is required, whether WebGL fallback is mandatory, whether the game must work offline as a PWA, whether WASM threading/SIMD is required, how browser storage will work, how asset loading will be handled, or what browser versions will be supported.
\item UNCERTAIN / UNVERIFIED: Current console platform facts were not verified in this chat. Future work must verify official PlayStation, Xbox, and Nintendo developer requirements from current official sources before implementation planning.
\item INFERENCE: PC handhelds should be subtargets under PC, with their own QA and UI profiles. UNCERTAIN / UNVERIFIED: Steam Deck's exact tier was not finally decided.
\end{itemize}

\subsubsection{Portability\_Assurance\_Future\_Proof}

\begin{itemize}
\item The chat proposed using standards such as DO-178C, NIST SSDF, OWASP ASVS, SLSA, and SPDX as design inputs only. The proposed internal version is DDAP v0 with DIL levels. Uncertainty: user has not explicitly ratified DDAP.
\item The actual repo was not inspected. The exact accepted directory structure, API policy, DDAP profile, compatibility promises, and first pilot module remain unresolved.
\item Caveat: carry forward as INFERENCE / assistant recommendation.
\item Caveat: carry forward as FACT goal + INFERENCE architecture.
\item Caveat: carry forward as assistant recommendation.
\item Implications: Safe for aggregation only with caveats.
\item Caveat: carry forward as FACT / UNCERTAIN scope.
\item 10. Feed this package into the master Project Spec Book with caveats.
\end{itemize}

\subsubsection{readme\_ports\_determinism}

\begin{itemize}
\item What remains uncertain is whether the actual repository README.md matches the final pasted version. The chat only saw pasted text and generated prompts; it did not inspect a Git repository.
\item The unresolved issue is whether a metadata-only /ports directory still violates the user's preference. The final README keeps /ports in a limited form, but the user's wording could mean no /ports directory at all. That should be clarified before writing the directory contract.
\item The unresolved work is to define the actual network protocol, prediction model, rollback window, reconciliation logic, and state-hash validation.
\item The major unresolved goal is deciding exactly what "no separate ports directory/system" means physically in the repository. The final README allows /ports as metadata-only. That may satisfy the user, but it is not confirmed. If the user meant no /ports directory at all, the README still needs another revision.
\item Another unresolved goal is making the README and normative specs align. The README says specs are authoritative, but those specs were not inspected in this chat.
\item UNCERTAIN / UNVERIFIED: /ports as metadata-only is the current README state, but may not be fully accepted.
\item A related rejected idea was **retro build flows under /ports/<target> containing source or alternative implementations**. The final README superseded this with metadata-only /ports, assuming /ports is retained at all. This rejection is final for source code and behavior, but the existence of /ports itself remains uncertain.
\item The first next step is to verify the actual repository README.md against the final pasted README. The chat only used pasted text. If the repository differs, future work should start from the actual file, not from memory.
\end{itemize}

\subsubsection{Refactor\_Architecture}

\begin{itemize}
\item Reliability note: Repository state, Codex execution, exact version numbers, and actual file modifications are **UNCERTAIN / UNVERIFIED** unless explicitly stated otherwise.
\item The most important thing to remember is that this chat did not implement the refactor. It designed the architecture and generated prompts and handoff material. Any future assistant must verify actual repository state before assuming files were moved, prompts were applied, or code was updated.
\item UNCERTAIN / UNVERIFIED:** It did not establish that this system already exists in code.
\item What remains uncertain is the exact degree to which existing code currently respects that boundary. The user provided a repo tree, but this chat did not inspect or modify code directly.
\item Unresolved details include exact default DOMINIUM\_HOME paths per OS and exact implementation of import/GC.
\item Unresolved details include exact schemas, parser choice, tests, and actual code implementation.
\item The key unresolved point is implementation timing. The report preserves this architecture, but the initial Codex refactor was not supposed to rewrite the entire UI system immediately.
\item Actual implementation is unresolved. This chat generated architecture and prompts, not a verified changed codebase.
\end{itemize}

\subsubsection{Refactor\_Control\_Plane}

\begin{itemize}
\item Unresolved goals include finishing AIDE Q35-Q57, stabilizing Dominium CTest/RepoX blockers, defining final AIDE install/upgrade bundles, implementing root recycling machinery, deciding when product boot proof can proceed, and later deciding whether AIDE Runtime/Gateway/Hosts are truly needed.
\item 1. Verify current Dominium HEAD and POST-CONVERGE status.
\item The user values direct, source-grounded, audit-ready, detailed answers. The user asked for preservation packages that are human-readable and machine-aggregatable. The user repeatedly prefers not to lose nuance, rejected options, uncertainties, or artifacts. The user wants current facts verified when possible and wants future assistants to avoid re-asking answered questions.
\item UNCERTAIN: Exact public-facing naming for AIDE Core/Kernel remains partly open. Exact threshold for accepting CTest-warning status remains a user/repo governance decision. Exact final AIDE installation packaging details remain pending.
\item Future assistants must verify live repo state before implementing.
\item "Merge this with another chat's preservation package without losing uncertainty labels."
\end{itemize}

\subsubsection{Release\_Identity\_and\_Versioning}

\begin{itemize}
\item The chat rebuilt SemVer from first principles and decided that strict SemVer should apply only to components with declared public contracts. Candidate true-SemVer entities include SDKs, engine libraries, tool APIs/CLIs, protocol libraries, plugin hosts, schema libraries, and reusable runtime libraries. It remains unresolved which exact repo modules qualify.
\item GBN should identify CI/public/internal builds and appear as provenance, e.g. +gbn.7137, but it must not be used as SemVer precedence. BII should primarily be structured manifest metadata, because it may encode lane, target, kind, configuration, and other build identity fields. The exact BII schema remains unresolved.
\item The main unresolved goals are to produce formal specs: Release Constitution, SemVer Component Inventory, Suite Release Policy, Build Identity Spec, Channel/Lifecycle Spec, Artifact Naming Spec, Manifest Schema, Target Taxonomy, and Capability Registry.
\item It is still uncertain how strict the user wants suite X.Y.Z field meanings to be, whether capabilities should fully replace compatibility profiles or coexist with them, and how much old-platform support should be real versus aspirational.
\item The most important unresolved issue is formal classification. Without knowing which entities are strict SemVer components, suite versions, product release IDs, build fields, or capabilities, the rest of the policy cannot be enforced safely.
\item Platform targeting remains unresolved. The chat concluded that coarse families like WinNT5, WinNT10, MacOSX4, Linux5, or DOS5 are useful support labels, but exact binary compatibility needs target baselines and runtime/toolchain profiles. Linux kernel major alone is not enough. Windows 9x/NT/NT5/NT6/NT10 likely need separate build lanes or at least explicit tested baselines.
\end{itemize}

\subsubsection{testx\_repox\_governance}

\begin{itemize}
\item This topic remains unresolved in implementation because the repo snapshot showed .dominium\_build\_number and update\_build\_number.cmake, but the chat did not inspect their contents. Those files may still implement older build-number behavior. That must be verified.
\item The unresolved part is implementation verification. The repo has Sol/Earth content in data and game/content; the engine must not assume those exist.
\item The unresolved part is how much of TESTX is actually implemented in the repo. The repo contains many tests and scripts, but this chat did not inspect them. A rules-to-checks audit is needed.
\item This remains conceptually accepted in the chat, but implementation is unverified.
\item The unresolved goals are mostly implementation and verification:
\item The repo already has UI IR/codegen infrastructure. The next step is to verify whether GUI/TUI truly bind to the CLI command graph and whether UI is data rather than logic.
\item The user has Windows 10 + VS2022. The next step is to verify v141/v141\_xp/SDK components and create isolated CMake presets for XP/Win7/Win8/Win10/Win11 builds.
\item Acting before verifying repo state.
\end{itemize}

\subsubsection{Timekeeping\_and\_2038\_Resilience}

\begin{itemize}
\item What remains uncertain is the current exact status of specific libc/OS/toolchain combinations, because those facts can change and were not reverified during this preservation turn.
\item What remains uncertain is the precise target list of legacy machines/toolchains and how far compatibility must go, especially for 16-bit compilers with weak or absent 64-bit arithmetic.
\item The main unresolved goals are backend audit, ACT serialization audit, civil/astronomical projection design, and verification of external platform facts.
\item The main uncertainty is completeness of repo inspection. The assistant read key files, but not all serialization paths and not all platform backends.
\item 6. Verify current platform/library facts before using them in formal public documentation.
\item The preservation prompt explicitly requires a human-readable report first, not only a machine-readable handoff. It requires uncertainty labels, stable registers, preservation of artifacts, no invented facts, and downloadable files if tools are available. ?filecite?turn30file0?
\item The user prefers source-grounded, audit-ready technical reasoning, direct structure, and careful distinction between facts, recommendations, and unresolved issues. PROJECT-CONTEXT indicates Dominium values C89/C++98 portability, deterministic architecture, CLI/TUI support, and clean platform/runtime separation.
\item The strongest unresolved technical task is auditing actual backend timers and persistence formats.
\end{itemize}

\subsubsection{UE6\_Domino\_Deterministic\_Universe}

\begin{itemize}
\item UNCERTAIN / UNVERIFIED: This is not a complete reconstruction of every Dominium conversation ever held. The project context shows that many other chats exist about platforms, renderers, setup, game architecture, ecology, ECS, launcher design, and development planning, but their full transcripts are not visible here. Those items are labelled PROJECT-CONTEXT when referenced.
\item Uncertainty: public UE6 technical capabilities and requirements remain time-sensitive and need verification before any future hard commitment.
\item The chat rejected the idea that clients can authoritatively simulate resource production, hidden state, enemy bases, combat, or persistent terrain changes in an MMO. Clients can assist, but the server must verify and commit.
\item REJECTED-02: Waiting for UE6 before beginning serious work. Deprioritised because UE6 availability and capabilities remain uncertain.
\item The user wants direct, source-grounded, audit-ready answers; explicit uncertainty; bounded estimates; and correction when framing is wrong. The preservation prompt explicitly requires human-readable explanation first, structured registers second, uncertainty labels, no invented facts, no hidden chain-of-thought, and downloadable files if tools are available.
\item UNCERTAIN: The final intended role of Domino remains unclear in the visible transcript. It may be a custom engine, target runtime, renderer, or broader platform plan, but the current chat does not fully define it.
\end{itemize}

\subsubsection{ui\_editor\_tool\_editor\_planning}

\begin{itemize}
\item The most important thing to remember is that this chat was a planning and prompt-generation chat, not proof of implementation. **UNCERTAIN / UNVERIFIED:** No generated prompt is evidence that repository code was actually changed. The next assistant must first verify whether the prompts were executed, inspect the repo/source bundle if available, and avoid treating planned files as existing files.
\item The user uploaded screenshot bundles for setup and launcher. **UNCERTAIN / UNVERIFIED:** The screenshots were not inspected in this chat. The assistant still produced a logical layout description, and the user accepted it as useful. Future work should inspect the bundles before claiming exact screenshot fidelity.
\item What remains uncertain is the actual state of the code. The user named files, but the chat did not inspect the repository. A future assistant must verify exact paths, build targets, TLV parser implementation, and current launcher behavior.
\item Unresolved details include the actual TLV wire format, migration strategy from launcher\_ui\_v1.tlv, and how much legacy data can be preserved during import.
\item The unresolved risk is implementation complexity. Real native tab controls, splitter behavior, and scroll panels can be nontrivial in a child-HWND Win32 UI.
\item This is a strong planned design, but some details remain assistant recommendations rather than explicit user-confirmed decisions. Future implementation should preserve that uncertainty if writing a formal spec.
\item The unresolved issue is that the uploaded screenshot bundles were not inspected. The logical specs are accepted planning artifacts, not verified screenshot extractions.
\item The exact scope of IDE-native live editing remains unresolved. The assistant proposed preview hosts, but it is not confirmed whether that fully satisfies the user's desire to use Visual Studio, Xcode, and Linux GUI tools.
\end{itemize}

\subsubsection{Universe\_Explorer\_Planning}

\begin{itemize}
\item Uncertain: whether all seven universal layers are now formally captured under exactly that naming. The repo has equivalent or related doctrine, but not every early phrase appears as a single artifact.
\item Uncertain: the repo contains pieces of this through affordance/capability/domain docs, but the assistant did not find one single master Deep Primitives spec in the docs it read.
\item The most important unresolved issue is that the Universe Explorer plan is still a proposed/recommended next strategy, not yet a completed repo task. The next best action is either to draft PRESENTATION-CONTRACT-01 / UNIVERSE-EXPLORER-CONTRACT-01, or to preserve this chat and merge it with other old-chat reports into a master Project Spec Book.
\end{itemize}

\subsubsection{Workbench\_AIDE\_Product\_Spine}

\begin{itemize}
\item The largest uncertainty is not the conceptual architecture; that has converged. The largest uncertainty is live execution state: whether prompts generated near the end of the chat have already been run locally, committed, or pushed. The next chat should first inspect .aide/queue/current.toml, recent commits, and relevant audit files.
\item Remaining uncertainty: exact implementation of erosion/ecology/evolution proxies remains future domain work.
\item 1. Verify current repo state.
\item Explicitly label uncertainty.
\item The next chat should verify repo state first.
\item Preserve FACT / INFERENCE / UNCERTAIN / PROJECT-CONTEXT labels.
\item Verify live repo state before acting.
\item safe\_for\_aggregation: "Yes, with caveats"
\end{itemize}

\subsubsection{World\_Architecture}

\begin{itemize}
\item The main unresolved issue is implementation detail. The design says Q16.16 and Q4.12, but the earlier prompt used signed types incorrectly. The future implementation must use unsigned local coordinates or a centered signed design.
\item The unresolved parts are the exact meshing algorithm, the fixed-point scaling of ?, and the sample resolution. A 32? sample grid per 16m chunk was recommended but not confirmed as final.
\item The exact solvers remain unresolved. The architecture is clear, but formulas and resolutions are still future work.
\item The save format remains conceptual in places. Exact binary headers, endian policy, compression, and migration strategy are unresolved.
\item The crucial caveat is implementation: unsigned or centered representations are needed. The design intent is final, but the exact typedefs must be fixed. If the implementation uses signed types incorrectly, the whole coordinate system breaks.
\item This is not yet an implementation-level final spec because solver formulas are unresolved.
\item FACT:** Core target is C89. This is non-negotiable in principle, but exact portability mechanics are unresolved because strict C89 lacks standard fixed-width integer types.
\item The user wants direct, detailed, rigorous, critical responses. The user prefers source labels, uncertainty preservation, and avoiding invented facts. The user does not want assistant suggestions silently upgraded to decisions.
\end{itemize}

\subsubsection{xstack\_lab\_galaxy}

\begin{itemize}
\item But this chat did not verify the repository. That became one of the most important unresolved issues.
\item This is one of the most important outputs associated with this chat, but it is also one of the most uncertain because it is user-reported from another chat and not verified directly here. The user explicitly wanted it audited for completion and consistency.
\item INFERENCE: The user also wanted to prevent future assistants from flattening nuance. They repeatedly asked for maximum fidelity, uncertainty labels, rejected options, contradictions, and rationale. This implies a strong preference for preserving complexity rather than oversimplifying.
\item INFERENCE: The user likely wants to resume productive development after verifying the substrate. The next likely feature area is survival or Lab Galaxy refinement, but the user has not made a final decision in this visible chat.
\item 4. From documentation to verifying a new 13-prompt substrate.
\item The biggest unresolved goal is verifying the actual repository. The package says what was reported, but not what is proven. Another unresolved goal is deciding what comes next: survival, Lab Galaxy UX, distributed SRZ, or documentation polish.
\item The main assumption behind all of this is that the repo can actually enforce these contracts. That remains unverified in this chat.
\item The best next step is to verify the actual repository state. This matters because the long 13-prompt transcript is not proof. The user explicitly does not know whether everything was truly implemented.
\end{itemize}

\subsection{Stale Claim Findings}

\begin{itemize}
\item CONTRA-0002 advanced\_simulation\_infrastructure: Archived conversation contains potentially stale external or baseline claim: C89.
\item CONTRA-0003 advanced\_simulation\_infrastructure: Archived conversation contains potentially stale external or baseline claim: C++98.
\item CONTRA-0004 advanced\_simulation\_infrastructure: Archived conversation contains potentially stale external or baseline claim: ISO C89.
\item CONTRA-0009 app\_runtime\_platform\_renderers: Archived conversation contains potentially stale external or baseline claim: SDK.
\item CONTRA-0012 app\_testx\_codehygiene: Archived conversation contains potentially stale external or baseline claim: C89.
\item CONTRA-0013 app\_testx\_codehygiene: Archived conversation contains potentially stale external or baseline claim: C++98.
\item CONTRA-0014 app\_testx\_codehygiene: Archived conversation contains potentially stale external or baseline claim: DX12.
\item CONTRA-0015 app\_testx\_codehygiene: Archived conversation contains potentially stale external or baseline claim: iOS.
\item CONTRA-0018 architecture\_codex\_prompts: Archived conversation contains potentially stale external or baseline claim: C89.
\item CONTRA-0019 architecture\_codex\_prompts: Archived conversation contains potentially stale external or baseline claim: C++98.
\item CONTRA-0020 architecture\_codex\_prompts: Archived conversation contains potentially stale external or baseline claim: ISO C89.
\item CONTRA-0026 architecture\_ui\_providers: Archived conversation contains potentially stale external or baseline claim: C89.
\item CONTRA-0027 architecture\_ui\_providers: Archived conversation contains potentially stale external or baseline claim: C++98.
\item CONTRA-0028 architecture\_ui\_providers: Archived conversation contains potentially stale external or baseline claim: DX12.
\item CONTRA-0029 architecture\_ui\_providers: Archived conversation contains potentially stale external or baseline claim: SDK.
\item CONTRA-0030 Build\_and\_Future\_Proofing: Archived conversation contains potentially stale external or baseline claim: C89.
\item CONTRA-0031 Build\_and\_Future\_Proofing: Archived conversation contains potentially stale external or baseline claim: C++98.
\item CONTRA-0032 Build\_and\_Future\_Proofing: Archived conversation contains potentially stale external or baseline claim: SDK.
\item CONTRA-0038 canonical\_structure\_and\_framework: Archived conversation contains potentially stale external or baseline claim: C89.
\item CONTRA-0039 canonical\_structure\_and\_framework: Archived conversation contains potentially stale external or baseline claim: C++98.
\item CONTRA-0040 canonical\_structure\_and\_framework: Archived conversation contains potentially stale external or baseline claim: SDK.
\item CONTRA-0043 development\_routes: Archived conversation contains potentially stale external or baseline claim: iOS.
\item CONTRA-0044 documentation\_standards\_readme: Archived conversation contains potentially stale external or baseline claim: C89.
\item CONTRA-0045 documentation\_standards\_readme: Archived conversation contains potentially stale external or baseline claim: C++98.
\item CONTRA-0046 documentation\_standards\_readme: Archived conversation contains potentially stale external or baseline claim: ISO C89.
\item CONTRA-0047 documentation\_standards\_readme: Archived conversation contains potentially stale external or baseline claim: iOS.
\item CONTRA-0050 Dominium\_Architecture\_I: Archived conversation contains potentially stale external or baseline claim: C89.
\item CONTRA-0051 Dominium\_Architecture\_I: Archived conversation contains potentially stale external or baseline claim: C++98.
\item CONTRA-0054 Dominium\_Architecture\_II: Archived conversation contains potentially stale external or baseline claim: C89.
\item CONTRA-0055 Dominium\_Architecture\_II: Archived conversation contains potentially stale external or baseline claim: Windows 98.
\item CONTRA-0056 Dominium\_Architecture\_II: Archived conversation contains potentially stale external or baseline claim: Windows NT 2000.
\item CONTRA-0057 Dominium\_Architecture\_II: Archived conversation contains potentially stale external or baseline claim: DX12.
\item CONTRA-0058 Dominium\_Architecture\_II: Archived conversation contains potentially stale external or baseline claim: Android.
\item CONTRA-0061 Dominium\_Architecture\_III: Archived conversation contains potentially stale external or baseline claim: C89.
\item CONTRA-0062 Dominium\_Architecture\_III: Archived conversation contains potentially stale external or baseline claim: C++98.
\item CONTRA-0063 Dominium\_Architecture\_III: Archived conversation contains potentially stale external or baseline claim: Win98.
\item CONTRA-0064 Dominium\_Architecture\_III: Archived conversation contains potentially stale external or baseline claim: Windows 98.
\item CONTRA-0065 Dominium\_Architecture\_III: Archived conversation contains potentially stale external or baseline claim: Windows NT 2000.
\item CONTRA-0066 Dominium\_Architecture\_III: Archived conversation contains potentially stale external or baseline claim: Mac OS 9.
\item CONTRA-0067 Dominium\_Architecture\_III: Archived conversation contains potentially stale external or baseline claim: Mac OS 8.
\item CONTRA-0068 Dominium\_Architecture\_III: Archived conversation contains potentially stale external or baseline claim: DX12.
\item CONTRA-0072 Dominium\_Architecture\_IV: Archived conversation contains potentially stale external or baseline claim: CP/M.
\item CONTRA-0073 Dominium\_Architecture\_IV: Archived conversation contains potentially stale external or baseline claim: Android.
\item CONTRA-0074 Dominium\_Architecture\_IV: Archived conversation contains potentially stale external or baseline claim: SDK.
\item CONTRA-0076 Dominium\_Complete\_Conversation: Archived conversation contains potentially stale external or baseline claim: C89.
\item CONTRA-0077 Dominium\_Complete\_Conversation: Archived conversation contains potentially stale external or baseline claim: C++98.
\item CONTRA-0078 dominium\_domino\_codex\_planning: Archived conversation contains potentially stale external or baseline claim: C89.
\item CONTRA-0079 dominium\_domino\_codex\_planning: Archived conversation contains potentially stale external or baseline claim: C++98.
\item CONTRA-0080 dominium\_domino\_codex\_planning: Archived conversation contains potentially stale external or baseline claim: ISO C89.
\item CONTRA-0081 dominium\_domino\_codex\_planning: Archived conversation contains potentially stale external or baseline claim: CP/M.
\item CONTRA-0082 dominium\_domino\_codex\_planning: Archived conversation contains potentially stale external or baseline claim: Android.
\item CONTRA-0088 dominium\_full\_conversation: Archived conversation contains potentially stale external or baseline claim: SDK.
\item CONTRA-0091 dominium\_setup: Archived conversation contains potentially stale external or baseline claim: SDK.
\item CONTRA-0096 Domino\_Dominium\_Workbench: Archived conversation contains potentially stale external or baseline claim: SDK.
\item CONTRA-0098 domino\_engine\_refactor\_prompts: Archived conversation contains potentially stale external or baseline claim: C89.
\item CONTRA-0099 domino\_engine\_refactor\_prompts: Archived conversation contains potentially stale external or baseline claim: C++98.
\item CONTRA-0100 domino\_engine\_refactor\_prompts: Archived conversation contains potentially stale external or baseline claim: ISO C89.
\item CONTRA-0111 Foundation\_Workbench\_Codex: Archived conversation contains potentially stale external or baseline claim: C89.
\item CONTRA-0112 Foundation\_Workbench\_Codex: Archived conversation contains potentially stale external or baseline claim: C++98.
\item CONTRA-0113 Foundation\_Workbench\_Codex: Archived conversation contains potentially stale external or baseline claim: SDK.
\item CONTRA-0118 Framework\_Open\_Source\_Provider: Archived conversation contains potentially stale external or baseline claim: C++98.
\item CONTRA-0119 Framework\_Open\_Source\_Provider: Archived conversation contains potentially stale external or baseline claim: SDK.
\item CONTRA-0122 gui\_binary\_content: Archived conversation contains potentially stale external or baseline claim: iOS.
\item CONTRA-0123 gui\_binary\_content: Archived conversation contains potentially stale external or baseline claim: Android.
\item CONTRA-0124 gui\_binary\_content: Archived conversation contains potentially stale external or baseline claim: SDK.
\item CONTRA-0129 Language\_Platform\_Architecture: Archived conversation contains potentially stale external or baseline claim: C89.
\item CONTRA-0130 Language\_Platform\_Architecture: Archived conversation contains potentially stale external or baseline claim: C++98.
\item CONTRA-0131 Language\_Platform\_Architecture: Archived conversation contains potentially stale external or baseline claim: Windows 98.
\item CONTRA-0132 Language\_Platform\_Architecture: Archived conversation contains potentially stale external or baseline claim: Mac OS 9.
\item CONTRA-0133 Language\_Platform\_Architecture: Archived conversation contains potentially stale external or baseline claim: iOS.
\item CONTRA-0134 Language\_Platform\_Architecture: Archived conversation contains potentially stale external or baseline claim: Android.
\item CONTRA-0135 Language\_Platform\_Architecture: Archived conversation contains potentially stale external or baseline claim: SDK.
\item CONTRA-0137 launcher\_app\_layer: Archived conversation contains potentially stale external or baseline claim: SDK.
\item CONTRA-0140 Launcher\_Setup\_Architecture: Archived conversation contains potentially stale external or baseline claim: C89.
\item CONTRA-0141 Launcher\_Setup\_Architecture: Archived conversation contains potentially stale external or baseline claim: C++98.
\item CONTRA-0142 Launcher\_Setup\_Architecture: Archived conversation contains potentially stale external or baseline claim: CP/M.
\item CONTRA-0143 Launcher\_Setup\_Architecture: Archived conversation contains potentially stale external or baseline claim: SDK.
\item CONTRA-0146 Modularity\_AIDE\_Refactorability: Archived conversation contains potentially stale external or baseline claim: C89.
\item CONTRA-0152 OS\_Interface\_Repo\_Architecture: Archived conversation contains potentially stale external or baseline claim: Mac OS 8.
\item CONTRA-0153 OS\_Interface\_Repo\_Architecture: Archived conversation contains potentially stale external or baseline claim: Android.
\item CONTRA-0154 OS\_Interface\_Repo\_Architecture: Archived conversation contains potentially stale external or baseline claim: SDK.
\item CONTRA-0158 platform\_renderer\_api\_plan: Archived conversation contains potentially stale external or baseline claim: C89.
\item CONTRA-0159 platform\_renderer\_api\_plan: Archived conversation contains potentially stale external or baseline claim: C++98.
\item CONTRA-0160 platform\_renderer\_api\_plan: Archived conversation contains potentially stale external or baseline claim: ISO C89.
\item CONTRA-0161 platform\_renderer\_api\_plan: Archived conversation contains potentially stale external or baseline claim: CP/M.
\item CONTRA-0162 platform\_renderer\_api\_plan: Archived conversation contains potentially stale external or baseline claim: Android.
\item CONTRA-0163 platform\_renderer\_api\_plan: Archived conversation contains potentially stale external or baseline claim: SDK.
\item CONTRA-0167 platform\_support: Archived conversation contains potentially stale external or baseline claim: CP/M.
\item CONTRA-0168 platform\_support: Archived conversation contains potentially stale external or baseline claim: iOS.
\item CONTRA-0169 platform\_support: Archived conversation contains potentially stale external or baseline claim: Android.
\item CONTRA-0170 platform\_support: Archived conversation contains potentially stale external or baseline claim: WebGPU.
\item CONTRA-0171 platform\_support: Archived conversation contains potentially stale external or baseline claim: console program.
\item CONTRA-0172 platform\_support: Archived conversation contains potentially stale external or baseline claim: SDK.
\item CONTRA-0174 Portability\_Assurance\_Future\_Proof: Archived conversation contains potentially stale external or baseline claim: C89.
\item CONTRA-0175 Portability\_Assurance\_Future\_Proof: Archived conversation contains potentially stale external or baseline claim: C++98.
\item CONTRA-0176 readme\_ports\_determinism: Archived conversation contains potentially stale external or baseline claim: C89.
\item CONTRA-0177 readme\_ports\_determinism: Archived conversation contains potentially stale external or baseline claim: C++98.
\item CONTRA-0178 readme\_ports\_determinism: Archived conversation contains potentially stale external or baseline claim: CP/M.
\item CONTRA-0179 readme\_ports\_determinism: Archived conversation contains potentially stale external or baseline claim: 286.
\item CONTRA-0180 Refactor\_Architecture: Archived conversation contains potentially stale external or baseline claim: C89.
\item CONTRA-0181 Refactor\_Architecture: Archived conversation contains potentially stale external or baseline claim: CP/M.
\item CONTRA-0182 Refactor\_Architecture: Archived conversation contains potentially stale external or baseline claim: Android.
\item CONTRA-0183 Refactor\_Architecture: Archived conversation contains potentially stale external or baseline claim: SDK.
\item CONTRA-0185 Refactor\_Control\_Plane: Archived conversation contains potentially stale external or baseline claim: SDK.
\item CONTRA-0187 Release\_Identity\_and\_Versioning: Archived conversation contains potentially stale external or baseline claim: SDK.
\item CONTRA-0192 testx\_repox\_governance: Archived conversation contains potentially stale external or baseline claim: C89.
\item CONTRA-0193 testx\_repox\_governance: Archived conversation contains potentially stale external or baseline claim: C++98.
\item CONTRA-0194 testx\_repox\_governance: Archived conversation contains potentially stale external or baseline claim: iOS.
\item CONTRA-0195 testx\_repox\_governance: Archived conversation contains potentially stale external or baseline claim: SDK.
\item CONTRA-0196 Timekeeping\_and\_2038\_Resilience: Archived conversation contains potentially stale external or baseline claim: C89.
\item CONTRA-0197 Timekeeping\_and\_2038\_Resilience: Archived conversation contains potentially stale external or baseline claim: C++98.
\item CONTRA-0200 UE6\_Domino\_Deterministic\_Universe: Archived conversation contains potentially stale external or baseline claim: C89.
\item CONTRA-0201 UE6\_Domino\_Deterministic\_Universe: Archived conversation contains potentially stale external or baseline claim: C++98.
\item CONTRA-0204 ui\_editor\_tool\_editor\_planning: Archived conversation contains potentially stale external or baseline claim: C++98.
\item CONTRA-0205 ui\_editor\_tool\_editor\_planning: Archived conversation contains potentially stale external or baseline claim: Android.
\item CONTRA-0216 Workbench\_AIDE\_Product\_Spine: Archived conversation contains potentially stale external or baseline claim: C89.
\item CONTRA-0217 Workbench\_AIDE\_Product\_Spine: Archived conversation contains potentially stale external or baseline claim: C++98.
\item CONTRA-0219 World\_Architecture: Archived conversation contains potentially stale external or baseline claim: C89.
\item CONTRA-0220 World\_Architecture: Archived conversation contains potentially stale external or baseline claim: C++98.
\item CONTRA-0221 World\_Architecture: Archived conversation contains potentially stale external or baseline claim: ISO C89.
\item CONTRA-0222 World\_Architecture: Archived conversation contains potentially stale external or baseline claim: SDK.
\end{itemize}



{\small\raggedright SOURCE PATH: docs/\allowbreak{}archive/\allowbreak{}conversations/\allowbreak{}\_\allowbreak{}audit/\allowbreak{}DOC\_\allowbreak{}DRIFT\_\allowbreak{}REPORT.md\par}


\section{Document Drift Report}

This report flags archived claims that may drift from current README, queue, language baseline, or authority order.


\subsection{CONTRA-0001 - conversation\_vs\_current\_queue}

\begin{itemize}
\item Source conversation: advanced\_simulation\_infrastructure
\item Source file: docs/archive/conversations/advanced\_simulation\_infrastructure/dominium\_advanced\_simulation\_infrastructure\_architecture\_\_human\_readable\_chat\_report.txt
\item Claim: Archived conversation discusses work related to blocked scope gameplay.
\item Conflict target: .aide/queue/current.toml
\item Recommended disposition: preserve\_as\_history; require current queue authority before action
\end{itemize}

\subsection{CONTRA-0002 - stale\_external\_claim}

\begin{itemize}
\item Source conversation: advanced\_simulation\_infrastructure
\item Source file: docs/archive/conversations/advanced\_simulation\_infrastructure/dominium\_advanced\_simulation\_infrastructure\_architecture\_\_human\_readable\_chat\_report.txt
\item Claim: Archived conversation contains potentially stale external or baseline claim: C89.
\item Conflict target: README.md / current platform and language baseline
\item Recommended disposition: quarantined until current repo and external facts are checked
\end{itemize}

\subsection{CONTRA-0003 - stale\_external\_claim}

\begin{itemize}
\item Source conversation: advanced\_simulation\_infrastructure
\item Source file: docs/archive/conversations/advanced\_simulation\_infrastructure/dominium\_advanced\_simulation\_infrastructure\_architecture\_\_human\_readable\_chat\_report.txt
\item Claim: Archived conversation contains potentially stale external or baseline claim: C++98.
\item Conflict target: README.md / current platform and language baseline
\item Recommended disposition: quarantined until current repo and external facts are checked
\end{itemize}

\subsection{CONTRA-0004 - stale\_external\_claim}

\begin{itemize}
\item Source conversation: advanced\_simulation\_infrastructure
\item Source file: docs/archive/conversations/advanced\_simulation\_infrastructure/dominium\_advanced\_simulation\_infrastructure\_architecture\_\_human\_readable\_chat\_report.txt
\item Claim: Archived conversation contains potentially stale external or baseline claim: ISO C89.
\item Conflict target: README.md / current platform and language baseline
\item Recommended disposition: quarantined until current repo and external facts are checked
\end{itemize}

\subsection{CONTRA-0005 - conversation\_vs\_current\_queue}

\begin{itemize}
\item Source conversation: app\_runtime\_platform\_renderers
\item Source file: docs/archive/conversations/app\_runtime\_platform\_renderers/dominium\_app0\_runtime\_platform\_renderers\_\_human\_readable\_chat\_report.txt
\item Claim: Archived conversation discusses work related to blocked scope runtime\_module\_loader.
\item Conflict target: .aide/queue/current.toml
\item Recommended disposition: preserve\_as\_history; require current queue authority before action
\end{itemize}

\subsection{CONTRA-0006 - conversation\_vs\_current\_queue}

\begin{itemize}
\item Source conversation: app\_runtime\_platform\_renderers
\item Source file: docs/archive/conversations/app\_runtime\_platform\_renderers/dominium\_app0\_runtime\_platform\_renderers\_\_human\_readable\_chat\_report.txt
\item Claim: Archived conversation discusses work related to blocked scope provider\_runtime.
\item Conflict target: .aide/queue/current.toml
\item Recommended disposition: preserve\_as\_history; require current queue authority before action
\end{itemize}

\subsection{CONTRA-0007 - conversation\_vs\_current\_queue}

\begin{itemize}
\item Source conversation: app\_runtime\_platform\_renderers
\item Source file: docs/archive/conversations/app\_runtime\_platform\_renderers/dominium\_app0\_runtime\_platform\_renderers\_\_human\_readable\_chat\_report.txt
\item Claim: Archived conversation discusses work related to blocked scope gameplay.
\item Conflict target: .aide/queue/current.toml
\item Recommended disposition: preserve\_as\_history; require current queue authority before action
\end{itemize}

\subsection{CONTRA-0008 - conversation\_vs\_current\_queue}

\begin{itemize}
\item Source conversation: app\_runtime\_platform\_renderers
\item Source file: docs/archive/conversations/app\_runtime\_platform\_renderers/dominium\_app0\_runtime\_platform\_renderers\_\_human\_readable\_chat\_report.txt
\item Claim: Archived conversation discusses work related to blocked scope renderer\_implementation.
\item Conflict target: .aide/queue/current.toml
\item Recommended disposition: preserve\_as\_history; require current queue authority before action
\end{itemize}

\subsection{CONTRA-0009 - stale\_external\_claim}

\begin{itemize}
\item Source conversation: app\_runtime\_platform\_renderers
\item Source file: docs/archive/conversations/app\_runtime\_platform\_renderers/dominium\_app0\_runtime\_platform\_renderers\_\_human\_readable\_chat\_report.txt
\item Claim: Archived conversation contains potentially stale external or baseline claim: SDK.
\item Conflict target: README.md / current platform and language baseline
\item Recommended disposition: quarantined until current repo and external facts are checked
\end{itemize}

\subsection{CONTRA-0010 - conversation\_vs\_current\_queue}

\begin{itemize}
\item Source conversation: app\_testx\_codehygiene
\item Source file: docs/archive/conversations/app\_testx\_codehygiene/dominium\_app\_testx\_codehygiene\_handoff\_\_human\_readable\_chat\_report.txt
\item Claim: Archived conversation discusses work related to blocked scope gameplay.
\item Conflict target: .aide/queue/current.toml
\item Recommended disposition: preserve\_as\_history; require current queue authority before action
\end{itemize}

\subsection{CONTRA-0011 - conversation\_vs\_current\_queue}

\begin{itemize}
\item Source conversation: app\_testx\_codehygiene
\item Source file: docs/archive/conversations/app\_testx\_codehygiene/dominium\_app\_testx\_codehygiene\_handoff\_\_human\_readable\_chat\_report.txt
\item Claim: Archived conversation discusses work related to blocked scope renderer\_implementation.
\item Conflict target: .aide/queue/current.toml
\item Recommended disposition: preserve\_as\_history; require current queue authority before action
\end{itemize}

\subsection{CONTRA-0012 - stale\_external\_claim}

\begin{itemize}
\item Source conversation: app\_testx\_codehygiene
\item Source file: docs/archive/conversations/app\_testx\_codehygiene/dominium\_app\_testx\_codehygiene\_handoff\_\_human\_readable\_chat\_report.txt
\item Claim: Archived conversation contains potentially stale external or baseline claim: C89.
\item Conflict target: README.md / current platform and language baseline
\item Recommended disposition: quarantined until current repo and external facts are checked
\end{itemize}

\subsection{CONTRA-0013 - stale\_external\_claim}

\begin{itemize}
\item Source conversation: app\_testx\_codehygiene
\item Source file: docs/archive/conversations/app\_testx\_codehygiene/dominium\_app\_testx\_codehygiene\_handoff\_\_human\_readable\_chat\_report.txt
\item Claim: Archived conversation contains potentially stale external or baseline claim: C++98.
\item Conflict target: README.md / current platform and language baseline
\item Recommended disposition: quarantined until current repo and external facts are checked
\end{itemize}

\subsection{CONTRA-0014 - stale\_external\_claim}

\begin{itemize}
\item Source conversation: app\_testx\_codehygiene
\item Source file: docs/archive/conversations/app\_testx\_codehygiene/dominium\_app\_testx\_codehygiene\_handoff\_\_human\_readable\_chat\_report.txt
\item Claim: Archived conversation contains potentially stale external or baseline claim: DX12.
\item Conflict target: README.md / current platform and language baseline
\item Recommended disposition: quarantined until current repo and external facts are checked
\end{itemize}

\subsection{CONTRA-0015 - stale\_external\_claim}

\begin{itemize}
\item Source conversation: app\_testx\_codehygiene
\item Source file: docs/archive/conversations/app\_testx\_codehygiene/dominium\_app\_testx\_codehygiene\_handoff\_\_human\_readable\_chat\_report.txt
\item Claim: Archived conversation contains potentially stale external or baseline claim: iOS.
\item Conflict target: README.md / current platform and language baseline
\item Recommended disposition: quarantined until current repo and external facts are checked
\end{itemize}

\subsection{CONTRA-0016 - conversation\_vs\_current\_queue}

\begin{itemize}
\item Source conversation: architecture\_codex\_prompts
\item Source file: docs/archive/conversations/architecture\_codex\_prompts/dominium\_domino\_architecture\_codex\_prompts\_\_human\_readable\_chat\_report.txt
\item Claim: Archived conversation discusses work related to blocked scope provider\_runtime.
\item Conflict target: .aide/queue/current.toml
\item Recommended disposition: preserve\_as\_history; require current queue authority before action
\end{itemize}

\subsection{CONTRA-0017 - conversation\_vs\_current\_queue}

\begin{itemize}
\item Source conversation: architecture\_codex\_prompts
\item Source file: docs/archive/conversations/architecture\_codex\_prompts/dominium\_domino\_architecture\_codex\_prompts\_\_human\_readable\_chat\_report.txt
\item Claim: Archived conversation discusses work related to blocked scope gameplay.
\item Conflict target: .aide/queue/current.toml
\item Recommended disposition: preserve\_as\_history; require current queue authority before action
\end{itemize}

\subsection{CONTRA-0018 - stale\_external\_claim}

\begin{itemize}
\item Source conversation: architecture\_codex\_prompts
\item Source file: docs/archive/conversations/architecture\_codex\_prompts/dominium\_domino\_architecture\_codex\_prompts\_\_human\_readable\_chat\_report.txt
\item Claim: Archived conversation contains potentially stale external or baseline claim: C89.
\item Conflict target: README.md / current platform and language baseline
\item Recommended disposition: quarantined until current repo and external facts are checked
\end{itemize}

\subsection{CONTRA-0019 - stale\_external\_claim}

\begin{itemize}
\item Source conversation: architecture\_codex\_prompts
\item Source file: docs/archive/conversations/architecture\_codex\_prompts/dominium\_domino\_architecture\_codex\_prompts\_\_human\_readable\_chat\_report.txt
\item Claim: Archived conversation contains potentially stale external or baseline claim: C++98.
\item Conflict target: README.md / current platform and language baseline
\item Recommended disposition: quarantined until current repo and external facts are checked
\end{itemize}

\subsection{CONTRA-0020 - stale\_external\_claim}

\begin{itemize}
\item Source conversation: architecture\_codex\_prompts
\item Source file: docs/archive/conversations/architecture\_codex\_prompts/dominium\_domino\_architecture\_codex\_prompts\_\_human\_readable\_chat\_report.txt
\item Claim: Archived conversation contains potentially stale external or baseline claim: ISO C89.
\item Conflict target: README.md / current platform and language baseline
\item Recommended disposition: quarantined until current repo and external facts are checked
\end{itemize}

\subsection{CONTRA-0021 - conversation\_vs\_current\_queue}

\begin{itemize}
\item Source conversation: architecture\_ui\_providers
\item Source file: docs/archive/conversations/architecture\_ui\_providers/dominium\_architecture\_ui\_providers\_robot\_os\_strategy\_\_01\_human\_readable\_report.md
\item Claim: Archived conversation discusses work related to blocked scope broad\_workbench\_ui.
\item Conflict target: .aide/queue/current.toml
\item Recommended disposition: preserve\_as\_history; require current queue authority before action
\end{itemize}

\subsection{CONTRA-0022 - conversation\_vs\_current\_queue}

\begin{itemize}
\item Source conversation: architecture\_ui\_providers
\item Source file: docs/archive/conversations/architecture\_ui\_providers/dominium\_architecture\_ui\_providers\_robot\_os\_strategy\_\_01\_human\_readable\_report.md
\item Claim: Archived conversation discusses work related to blocked scope provider\_runtime.
\item Conflict target: .aide/queue/current.toml
\item Recommended disposition: preserve\_as\_history; require current queue authority before action
\end{itemize}

\subsection{CONTRA-0023 - conversation\_vs\_current\_queue}

\begin{itemize}
\item Source conversation: architecture\_ui\_providers
\item Source file: docs/archive/conversations/architecture\_ui\_providers/dominium\_architecture\_ui\_providers\_robot\_os\_strategy\_\_01\_human\_readable\_report.md
\item Claim: Archived conversation discusses work related to blocked scope gameplay.
\item Conflict target: .aide/queue/current.toml
\item Recommended disposition: preserve\_as\_history; require current queue authority before action
\end{itemize}

\subsection{CONTRA-0024 - conversation\_vs\_current\_queue}

\begin{itemize}
\item Source conversation: architecture\_ui\_providers
\item Source file: docs/archive/conversations/architecture\_ui\_providers/dominium\_architecture\_ui\_providers\_robot\_os\_strategy\_\_01\_human\_readable\_report.md
\item Claim: Archived conversation discusses work related to blocked scope renderer\_implementation.
\item Conflict target: .aide/queue/current.toml
\item Recommended disposition: preserve\_as\_history; require current queue authority before action
\end{itemize}

\subsection{CONTRA-0025 - conversation\_vs\_current\_queue}

\begin{itemize}
\item Source conversation: architecture\_ui\_providers
\item Source file: docs/archive/conversations/architecture\_ui\_providers/dominium\_architecture\_ui\_providers\_robot\_os\_strategy\_\_01\_human\_readable\_report.md
\item Claim: Archived conversation discusses work related to blocked scope native\_gui.
\item Conflict target: .aide/queue/current.toml
\item Recommended disposition: preserve\_as\_history; require current queue authority before action
\end{itemize}

\subsection{CONTRA-0026 - stale\_external\_claim}

\begin{itemize}
\item Source conversation: architecture\_ui\_providers
\item Source file: docs/archive/conversations/architecture\_ui\_providers/dominium\_architecture\_ui\_providers\_robot\_os\_strategy\_\_01\_human\_readable\_report.md
\item Claim: Archived conversation contains potentially stale external or baseline claim: C89.
\item Conflict target: README.md / current platform and language baseline
\item Recommended disposition: quarantined until current repo and external facts are checked
\end{itemize}

\subsection{CONTRA-0027 - stale\_external\_claim}

\begin{itemize}
\item Source conversation: architecture\_ui\_providers
\item Source file: docs/archive/conversations/architecture\_ui\_providers/dominium\_architecture\_ui\_providers\_robot\_os\_strategy\_\_01\_human\_readable\_report.md
\item Claim: Archived conversation contains potentially stale external or baseline claim: C++98.
\item Conflict target: README.md / current platform and language baseline
\item Recommended disposition: quarantined until current repo and external facts are checked
\end{itemize}

\subsection{CONTRA-0028 - stale\_external\_claim}

\begin{itemize}
\item Source conversation: architecture\_ui\_providers
\item Source file: docs/archive/conversations/architecture\_ui\_providers/dominium\_architecture\_ui\_providers\_robot\_os\_strategy\_\_01\_human\_readable\_report.md
\item Claim: Archived conversation contains potentially stale external or baseline claim: DX12.
\item Conflict target: README.md / current platform and language baseline
\item Recommended disposition: quarantined until current repo and external facts are checked
\end{itemize}

\subsection{CONTRA-0029 - stale\_external\_claim}

\begin{itemize}
\item Source conversation: architecture\_ui\_providers
\item Source file: docs/archive/conversations/architecture\_ui\_providers/dominium\_architecture\_ui\_providers\_robot\_os\_strategy\_\_01\_human\_readable\_report.md
\item Claim: Archived conversation contains potentially stale external or baseline claim: SDK.
\item Conflict target: README.md / current platform and language baseline
\item Recommended disposition: quarantined until current repo and external facts are checked
\end{itemize}

\subsection{CONTRA-0030 - stale\_external\_claim}

\begin{itemize}
\item Source conversation: Build\_and\_Future\_Proofing
\item Source file: docs/archive/conversations/Build\_and\_Future\_Proofing/Dominium\_Build\_and\_Future\_Proofing\_Architecture\_\_01\_human\_readable\_report.md
\item Claim: Archived conversation contains potentially stale external or baseline claim: C89.
\item Conflict target: README.md / current platform and language baseline
\item Recommended disposition: quarantined until current repo and external facts are checked
\end{itemize}

\subsection{CONTRA-0031 - stale\_external\_claim}

\begin{itemize}
\item Source conversation: Build\_and\_Future\_Proofing
\item Source file: docs/archive/conversations/Build\_and\_Future\_Proofing/Dominium\_Build\_and\_Future\_Proofing\_Architecture\_\_01\_human\_readable\_report.md
\item Claim: Archived conversation contains potentially stale external or baseline claim: C++98.
\item Conflict target: README.md / current platform and language baseline
\item Recommended disposition: quarantined until current repo and external facts are checked
\end{itemize}

\subsection{CONTRA-0032 - stale\_external\_claim}

\begin{itemize}
\item Source conversation: Build\_and\_Future\_Proofing
\item Source file: docs/archive/conversations/Build\_and\_Future\_Proofing/Dominium\_Build\_and\_Future\_Proofing\_Architecture\_\_01\_human\_readable\_report.md
\item Claim: Archived conversation contains potentially stale external or baseline claim: SDK.
\item Conflict target: README.md / current platform and language baseline
\item Recommended disposition: quarantined until current repo and external facts are checked
\end{itemize}

\subsection{CONTRA-0033 - conversation\_vs\_current\_queue}

\begin{itemize}
\item Source conversation: canonical\_structure\_and\_framework
\item Source file: docs/archive/conversations/canonical\_structure\_and\_framework/dominium\_canonical\_architecture\_repo\_foundation\_\_01\_human\_readable\_report.md
\item Claim: Archived conversation discusses work related to blocked scope provider\_runtime.
\item Conflict target: .aide/queue/current.toml
\item Recommended disposition: preserve\_as\_history; require current queue authority before action
\end{itemize}

\subsection{CONTRA-0034 - conversation\_vs\_current\_queue}

\begin{itemize}
\item Source conversation: canonical\_structure\_and\_framework
\item Source file: docs/archive/conversations/canonical\_structure\_and\_framework/dominium\_canonical\_architecture\_repo\_foundation\_\_01\_human\_readable\_report.md
\item Claim: Archived conversation discusses work related to blocked scope package\_runtime.
\item Conflict target: .aide/queue/current.toml
\item Recommended disposition: preserve\_as\_history; require current queue authority before action
\end{itemize}

\subsection{CONTRA-0035 - conversation\_vs\_current\_queue}

\begin{itemize}
\item Source conversation: canonical\_structure\_and\_framework
\item Source file: docs/archive/conversations/canonical\_structure\_and\_framework/dominium\_canonical\_architecture\_repo\_foundation\_\_01\_human\_readable\_report.md
\item Claim: Archived conversation discusses work related to blocked scope gameplay.
\item Conflict target: .aide/queue/current.toml
\item Recommended disposition: preserve\_as\_history; require current queue authority before action
\end{itemize}

\subsection{CONTRA-0036 - conversation\_vs\_current\_queue}

\begin{itemize}
\item Source conversation: canonical\_structure\_and\_framework
\item Source file: docs/archive/conversations/canonical\_structure\_and\_framework/dominium\_canonical\_architecture\_repo\_foundation\_\_01\_human\_readable\_report.md
\item Claim: Archived conversation discusses work related to blocked scope renderer\_implementation.
\item Conflict target: .aide/queue/current.toml
\item Recommended disposition: preserve\_as\_history; require current queue authority before action
\end{itemize}

\subsection{CONTRA-0037 - conversation\_vs\_current\_queue}

\begin{itemize}
\item Source conversation: canonical\_structure\_and\_framework
\item Source file: docs/archive/conversations/canonical\_structure\_and\_framework/dominium\_canonical\_architecture\_repo\_foundation\_\_01\_human\_readable\_report.md
\item Claim: Archived conversation discusses work related to blocked scope native\_gui.
\item Conflict target: .aide/queue/current.toml
\item Recommended disposition: preserve\_as\_history; require current queue authority before action
\end{itemize}

\subsection{CONTRA-0038 - stale\_external\_claim}

\begin{itemize}
\item Source conversation: canonical\_structure\_and\_framework
\item Source file: docs/archive/conversations/canonical\_structure\_and\_framework/dominium\_canonical\_architecture\_repo\_foundation\_\_01\_human\_readable\_report.md
\item Claim: Archived conversation contains potentially stale external or baseline claim: C89.
\item Conflict target: README.md / current platform and language baseline
\item Recommended disposition: quarantined until current repo and external facts are checked
\end{itemize}

\subsection{CONTRA-0039 - stale\_external\_claim}

\begin{itemize}
\item Source conversation: canonical\_structure\_and\_framework
\item Source file: docs/archive/conversations/canonical\_structure\_and\_framework/dominium\_canonical\_architecture\_repo\_foundation\_\_01\_human\_readable\_report.md
\item Claim: Archived conversation contains potentially stale external or baseline claim: C++98.
\item Conflict target: README.md / current platform and language baseline
\item Recommended disposition: quarantined until current repo and external facts are checked
\end{itemize}

\subsection{CONTRA-0040 - stale\_external\_claim}

\begin{itemize}
\item Source conversation: canonical\_structure\_and\_framework
\item Source file: docs/archive/conversations/canonical\_structure\_and\_framework/dominium\_canonical\_architecture\_repo\_foundation\_\_01\_human\_readable\_report.md
\item Claim: Archived conversation contains potentially stale external or baseline claim: SDK.
\item Conflict target: README.md / current platform and language baseline
\item Recommended disposition: quarantined until current repo and external facts are checked
\end{itemize}

\subsection{CONTRA-0041 - conversation\_vs\_current\_queue}

\begin{itemize}
\item Source conversation: Chronology\_Celestial\_Systems
\item Source file: docs/archive/conversations/Chronology\_Celestial\_Systems/Dominium\_Chronology\_Celestial\_Systems\_\_human\_readable\_chat\_report.txt
\item Claim: Archived conversation discusses work related to blocked scope gameplay.
\item Conflict target: .aide/queue/current.toml
\item Recommended disposition: preserve\_as\_history; require current queue authority before action
\end{itemize}

\subsection{CONTRA-0042 - conversation\_vs\_current\_queue}

\begin{itemize}
\item Source conversation: development\_routes
\item Source file: docs/archive/conversations/development\_routes/dominium\_development\_routes\_\_human\_readable\_chat\_report.txt
\item Claim: Archived conversation discusses work related to blocked scope gameplay.
\item Conflict target: .aide/queue/current.toml
\item Recommended disposition: preserve\_as\_history; require current queue authority before action
\end{itemize}

\subsection{CONTRA-0043 - stale\_external\_claim}

\begin{itemize}
\item Source conversation: development\_routes
\item Source file: docs/archive/conversations/development\_routes/dominium\_development\_routes\_\_human\_readable\_chat\_report.txt
\item Claim: Archived conversation contains potentially stale external or baseline claim: iOS.
\item Conflict target: README.md / current platform and language baseline
\item Recommended disposition: quarantined until current repo and external facts are checked
\end{itemize}

\subsection{CONTRA-0044 - stale\_external\_claim}

\begin{itemize}
\item Source conversation: documentation\_standards\_readme
\item Source file: docs/archive/conversations/documentation\_standards\_readme/documentation\_standards\_readme\_handoff\_\_human\_readable\_chat\_report.txt
\item Claim: Archived conversation contains potentially stale external or baseline claim: C89.
\item Conflict target: README.md / current platform and language baseline
\item Recommended disposition: quarantined until current repo and external facts are checked
\end{itemize}

\subsection{CONTRA-0045 - stale\_external\_claim}

\begin{itemize}
\item Source conversation: documentation\_standards\_readme
\item Source file: docs/archive/conversations/documentation\_standards\_readme/documentation\_standards\_readme\_handoff\_\_human\_readable\_chat\_report.txt
\item Claim: Archived conversation contains potentially stale external or baseline claim: C++98.
\item Conflict target: README.md / current platform and language baseline
\item Recommended disposition: quarantined until current repo and external facts are checked
\end{itemize}

\subsection{CONTRA-0046 - stale\_external\_claim}

\begin{itemize}
\item Source conversation: documentation\_standards\_readme
\item Source file: docs/archive/conversations/documentation\_standards\_readme/documentation\_standards\_readme\_handoff\_\_human\_readable\_chat\_report.txt
\item Claim: Archived conversation contains potentially stale external or baseline claim: ISO C89.
\item Conflict target: README.md / current platform and language baseline
\item Recommended disposition: quarantined until current repo and external facts are checked
\end{itemize}

\subsection{CONTRA-0047 - stale\_external\_claim}

\begin{itemize}
\item Source conversation: documentation\_standards\_readme
\item Source file: docs/archive/conversations/documentation\_standards\_readme/documentation\_standards\_readme\_handoff\_\_human\_readable\_chat\_report.txt
\item Claim: Archived conversation contains potentially stale external or baseline claim: iOS.
\item Conflict target: README.md / current platform and language baseline
\item Recommended disposition: quarantined until current repo and external facts are checked
\end{itemize}

\subsection{CONTRA-0048 - conversation\_vs\_current\_queue}

\begin{itemize}
\item Source conversation: Dominium\_Architecture\_I
\item Source file: docs/archive/conversations/Dominium\_Architecture\_I/Dominium\_Architecture\_I\_\_human\_readable\_chat\_report.txt
\item Claim: Archived conversation discusses work related to blocked scope gameplay.
\item Conflict target: .aide/queue/current.toml
\item Recommended disposition: preserve\_as\_history; require current queue authority before action
\end{itemize}

\subsection{CONTRA-0049 - conversation\_vs\_current\_queue}

\begin{itemize}
\item Source conversation: Dominium\_Architecture\_I
\item Source file: docs/archive/conversations/Dominium\_Architecture\_I/Dominium\_Architecture\_I\_\_human\_readable\_chat\_report.txt
\item Claim: Archived conversation discusses work related to blocked scope renderer\_implementation.
\item Conflict target: .aide/queue/current.toml
\item Recommended disposition: preserve\_as\_history; require current queue authority before action
\end{itemize}

\subsection{CONTRA-0050 - stale\_external\_claim}

\begin{itemize}
\item Source conversation: Dominium\_Architecture\_I
\item Source file: docs/archive/conversations/Dominium\_Architecture\_I/Dominium\_Architecture\_I\_\_human\_readable\_chat\_report.txt
\item Claim: Archived conversation contains potentially stale external or baseline claim: C89.
\item Conflict target: README.md / current platform and language baseline
\item Recommended disposition: quarantined until current repo and external facts are checked
\end{itemize}

\subsection{CONTRA-0051 - stale\_external\_claim}

\begin{itemize}
\item Source conversation: Dominium\_Architecture\_I
\item Source file: docs/archive/conversations/Dominium\_Architecture\_I/Dominium\_Architecture\_I\_\_human\_readable\_chat\_report.txt
\item Claim: Archived conversation contains potentially stale external or baseline claim: C++98.
\item Conflict target: README.md / current platform and language baseline
\item Recommended disposition: quarantined until current repo and external facts are checked
\end{itemize}

\subsection{CONTRA-0052 - conversation\_vs\_current\_queue}

\begin{itemize}
\item Source conversation: Dominium\_Architecture\_II
\item Source file: docs/archive/conversations/Dominium\_Architecture\_II/Dominium\_Architecture\_II\_\_human\_readable\_chat\_report.txt
\item Claim: Archived conversation discusses work related to blocked scope gameplay.
\item Conflict target: .aide/queue/current.toml
\item Recommended disposition: preserve\_as\_history; require current queue authority before action
\end{itemize}

\subsection{CONTRA-0053 - conversation\_vs\_current\_queue}

\begin{itemize}
\item Source conversation: Dominium\_Architecture\_II
\item Source file: docs/archive/conversations/Dominium\_Architecture\_II/Dominium\_Architecture\_II\_\_human\_readable\_chat\_report.txt
\item Claim: Archived conversation discusses work related to blocked scope renderer\_implementation.
\item Conflict target: .aide/queue/current.toml
\item Recommended disposition: preserve\_as\_history; require current queue authority before action
\end{itemize}

\subsection{CONTRA-0054 - stale\_external\_claim}

\begin{itemize}
\item Source conversation: Dominium\_Architecture\_II
\item Source file: docs/archive/conversations/Dominium\_Architecture\_II/Dominium\_Architecture\_II\_\_human\_readable\_chat\_report.txt
\item Claim: Archived conversation contains potentially stale external or baseline claim: C89.
\item Conflict target: README.md / current platform and language baseline
\item Recommended disposition: quarantined until current repo and external facts are checked
\end{itemize}

\subsection{CONTRA-0055 - stale\_external\_claim}

\begin{itemize}
\item Source conversation: Dominium\_Architecture\_II
\item Source file: docs/archive/conversations/Dominium\_Architecture\_II/Dominium\_Architecture\_II\_\_human\_readable\_chat\_report.txt
\item Claim: Archived conversation contains potentially stale external or baseline claim: Windows 98.
\item Conflict target: README.md / current platform and language baseline
\item Recommended disposition: quarantined until current repo and external facts are checked
\end{itemize}

\subsection{CONTRA-0056 - stale\_external\_claim}

\begin{itemize}
\item Source conversation: Dominium\_Architecture\_II
\item Source file: docs/archive/conversations/Dominium\_Architecture\_II/Dominium\_Architecture\_II\_\_human\_readable\_chat\_report.txt
\item Claim: Archived conversation contains potentially stale external or baseline claim: Windows NT 2000.
\item Conflict target: README.md / current platform and language baseline
\item Recommended disposition: quarantined until current repo and external facts are checked
\end{itemize}

\subsection{CONTRA-0057 - stale\_external\_claim}

\begin{itemize}
\item Source conversation: Dominium\_Architecture\_II
\item Source file: docs/archive/conversations/Dominium\_Architecture\_II/Dominium\_Architecture\_II\_\_human\_readable\_chat\_report.txt
\item Claim: Archived conversation contains potentially stale external or baseline claim: DX12.
\item Conflict target: README.md / current platform and language baseline
\item Recommended disposition: quarantined until current repo and external facts are checked
\end{itemize}

\subsection{CONTRA-0058 - stale\_external\_claim}

\begin{itemize}
\item Source conversation: Dominium\_Architecture\_II
\item Source file: docs/archive/conversations/Dominium\_Architecture\_II/Dominium\_Architecture\_II\_\_human\_readable\_chat\_report.txt
\item Claim: Archived conversation contains potentially stale external or baseline claim: Android.
\item Conflict target: README.md / current platform and language baseline
\item Recommended disposition: quarantined until current repo and external facts are checked
\end{itemize}

\subsection{CONTRA-0059 - conversation\_vs\_current\_queue}

\begin{itemize}
\item Source conversation: Dominium\_Architecture\_III
\item Source file: docs/archive/conversations/Dominium\_Architecture\_III/Dominium\_Architecture\_III\_\_human\_readable\_chat\_report.txt
\item Claim: Archived conversation discusses work related to blocked scope renderer\_implementation.
\item Conflict target: .aide/queue/current.toml
\item Recommended disposition: preserve\_as\_history; require current queue authority before action
\end{itemize}

\subsection{CONTRA-0060 - conversation\_vs\_current\_queue}

\begin{itemize}
\item Source conversation: Dominium\_Architecture\_III
\item Source file: docs/archive/conversations/Dominium\_Architecture\_III/Dominium\_Architecture\_III\_\_human\_readable\_chat\_report.txt
\item Claim: Archived conversation discusses work related to blocked scope release\_publication.
\item Conflict target: .aide/queue/current.toml
\item Recommended disposition: preserve\_as\_history; require current queue authority before action
\end{itemize}

\subsection{CONTRA-0061 - stale\_external\_claim}

\begin{itemize}
\item Source conversation: Dominium\_Architecture\_III
\item Source file: docs/archive/conversations/Dominium\_Architecture\_III/Dominium\_Architecture\_III\_\_human\_readable\_chat\_report.txt
\item Claim: Archived conversation contains potentially stale external or baseline claim: C89.
\item Conflict target: README.md / current platform and language baseline
\item Recommended disposition: quarantined until current repo and external facts are checked
\end{itemize}

\subsection{CONTRA-0062 - stale\_external\_claim}

\begin{itemize}
\item Source conversation: Dominium\_Architecture\_III
\item Source file: docs/archive/conversations/Dominium\_Architecture\_III/Dominium\_Architecture\_III\_\_human\_readable\_chat\_report.txt
\item Claim: Archived conversation contains potentially stale external or baseline claim: C++98.
\item Conflict target: README.md / current platform and language baseline
\item Recommended disposition: quarantined until current repo and external facts are checked
\end{itemize}

\subsection{CONTRA-0063 - stale\_external\_claim}

\begin{itemize}
\item Source conversation: Dominium\_Architecture\_III
\item Source file: docs/archive/conversations/Dominium\_Architecture\_III/Dominium\_Architecture\_III\_\_human\_readable\_chat\_report.txt
\item Claim: Archived conversation contains potentially stale external or baseline claim: Win98.
\item Conflict target: README.md / current platform and language baseline
\item Recommended disposition: quarantined until current repo and external facts are checked
\end{itemize}

\subsection{CONTRA-0064 - stale\_external\_claim}

\begin{itemize}
\item Source conversation: Dominium\_Architecture\_III
\item Source file: docs/archive/conversations/Dominium\_Architecture\_III/Dominium\_Architecture\_III\_\_human\_readable\_chat\_report.txt
\item Claim: Archived conversation contains potentially stale external or baseline claim: Windows 98.
\item Conflict target: README.md / current platform and language baseline
\item Recommended disposition: quarantined until current repo and external facts are checked
\end{itemize}

\subsection{CONTRA-0065 - stale\_external\_claim}

\begin{itemize}
\item Source conversation: Dominium\_Architecture\_III
\item Source file: docs/archive/conversations/Dominium\_Architecture\_III/Dominium\_Architecture\_III\_\_human\_readable\_chat\_report.txt
\item Claim: Archived conversation contains potentially stale external or baseline claim: Windows NT 2000.
\item Conflict target: README.md / current platform and language baseline
\item Recommended disposition: quarantined until current repo and external facts are checked
\end{itemize}

\subsection{CONTRA-0066 - stale\_external\_claim}

\begin{itemize}
\item Source conversation: Dominium\_Architecture\_III
\item Source file: docs/archive/conversations/Dominium\_Architecture\_III/Dominium\_Architecture\_III\_\_human\_readable\_chat\_report.txt
\item Claim: Archived conversation contains potentially stale external or baseline claim: Mac OS 9.
\item Conflict target: README.md / current platform and language baseline
\item Recommended disposition: quarantined until current repo and external facts are checked
\end{itemize}

\subsection{CONTRA-0067 - stale\_external\_claim}

\begin{itemize}
\item Source conversation: Dominium\_Architecture\_III
\item Source file: docs/archive/conversations/Dominium\_Architecture\_III/Dominium\_Architecture\_III\_\_human\_readable\_chat\_report.txt
\item Claim: Archived conversation contains potentially stale external or baseline claim: Mac OS 8.
\item Conflict target: README.md / current platform and language baseline
\item Recommended disposition: quarantined until current repo and external facts are checked
\end{itemize}

\subsection{CONTRA-0068 - stale\_external\_claim}

\begin{itemize}
\item Source conversation: Dominium\_Architecture\_III
\item Source file: docs/archive/conversations/Dominium\_Architecture\_III/Dominium\_Architecture\_III\_\_human\_readable\_chat\_report.txt
\item Claim: Archived conversation contains potentially stale external or baseline claim: DX12.
\item Conflict target: README.md / current platform and language baseline
\item Recommended disposition: quarantined until current repo and external facts are checked
\end{itemize}

\subsection{CONTRA-0069 - conversation\_vs\_current\_queue}

\begin{itemize}
\item Source conversation: Dominium\_Architecture\_IV
\item Source file: docs/archive/conversations/Dominium\_Architecture\_IV/Dominium\_Architecture\_IV\_\_human\_readable\_chat\_report.txt
\item Claim: Archived conversation discusses work related to blocked scope gameplay.
\item Conflict target: .aide/queue/current.toml
\item Recommended disposition: preserve\_as\_history; require current queue authority before action
\end{itemize}

\subsection{CONTRA-0070 - conversation\_vs\_current\_queue}

\begin{itemize}
\item Source conversation: Dominium\_Architecture\_IV
\item Source file: docs/archive/conversations/Dominium\_Architecture\_IV/Dominium\_Architecture\_IV\_\_human\_readable\_chat\_report.txt
\item Claim: Archived conversation discusses work related to blocked scope renderer\_implementation.
\item Conflict target: .aide/queue/current.toml
\item Recommended disposition: preserve\_as\_history; require current queue authority before action
\end{itemize}

\subsection{CONTRA-0071 - conversation\_vs\_current\_queue}

\begin{itemize}
\item Source conversation: Dominium\_Architecture\_IV
\item Source file: docs/archive/conversations/Dominium\_Architecture\_IV/Dominium\_Architecture\_IV\_\_human\_readable\_chat\_report.txt
\item Claim: Archived conversation discusses work related to blocked scope native\_gui.
\item Conflict target: .aide/queue/current.toml
\item Recommended disposition: preserve\_as\_history; require current queue authority before action
\end{itemize}

\subsection{CONTRA-0072 - stale\_external\_claim}

\begin{itemize}
\item Source conversation: Dominium\_Architecture\_IV
\item Source file: docs/archive/conversations/Dominium\_Architecture\_IV/Dominium\_Architecture\_IV\_\_human\_readable\_chat\_report.txt
\item Claim: Archived conversation contains potentially stale external or baseline claim: CP/M.
\item Conflict target: README.md / current platform and language baseline
\item Recommended disposition: quarantined until current repo and external facts are checked
\end{itemize}

\subsection{CONTRA-0073 - stale\_external\_claim}

\begin{itemize}
\item Source conversation: Dominium\_Architecture\_IV
\item Source file: docs/archive/conversations/Dominium\_Architecture\_IV/Dominium\_Architecture\_IV\_\_human\_readable\_chat\_report.txt
\item Claim: Archived conversation contains potentially stale external or baseline claim: Android.
\item Conflict target: README.md / current platform and language baseline
\item Recommended disposition: quarantined until current repo and external facts are checked
\end{itemize}

\subsection{CONTRA-0074 - stale\_external\_claim}

\begin{itemize}
\item Source conversation: Dominium\_Architecture\_IV
\item Source file: docs/archive/conversations/Dominium\_Architecture\_IV/Dominium\_Architecture\_IV\_\_human\_readable\_chat\_report.txt
\item Claim: Archived conversation contains potentially stale external or baseline claim: SDK.
\item Conflict target: README.md / current platform and language baseline
\item Recommended disposition: quarantined until current repo and external facts are checked
\end{itemize}

\subsection{CONTRA-0075 - conversation\_vs\_current\_queue}

\begin{itemize}
\item Source conversation: Dominium\_Complete\_Conversation
\item Source file: docs/archive/conversations/Dominium\_Complete\_Conversation/Accompanying\_Report/Dominium\_Canon\_Portability\_Audit\_\_01\_human\_readable\_report.md
\item Claim: Archived conversation discusses work related to blocked scope gameplay.
\item Conflict target: .aide/queue/current.toml
\item Recommended disposition: preserve\_as\_history; require current queue authority before action
\end{itemize}

\subsection{CONTRA-0076 - stale\_external\_claim}

\begin{itemize}
\item Source conversation: Dominium\_Complete\_Conversation
\item Source file: docs/archive/conversations/Dominium\_Complete\_Conversation/Accompanying\_Report/Dominium\_Canon\_Portability\_Audit\_\_01\_human\_readable\_report.md
\item Claim: Archived conversation contains potentially stale external or baseline claim: C89.
\item Conflict target: README.md / current platform and language baseline
\item Recommended disposition: quarantined until current repo and external facts are checked
\end{itemize}

\subsection{CONTRA-0077 - stale\_external\_claim}

\begin{itemize}
\item Source conversation: Dominium\_Complete\_Conversation
\item Source file: docs/archive/conversations/Dominium\_Complete\_Conversation/Accompanying\_Report/Dominium\_Canon\_Portability\_Audit\_\_01\_human\_readable\_report.md
\item Claim: Archived conversation contains potentially stale external or baseline claim: C++98.
\item Conflict target: README.md / current platform and language baseline
\item Recommended disposition: quarantined until current repo and external facts are checked
\end{itemize}

\subsection{CONTRA-0078 - stale\_external\_claim}

\begin{itemize}
\item Source conversation: dominium\_domino\_codex\_planning
\item Source file: docs/archive/conversations/dominium\_domino\_codex\_planning/dominium\_domino\_codex\_planning\_\_human\_readable\_chat\_report.txt
\item Claim: Archived conversation contains potentially stale external or baseline claim: C89.
\item Conflict target: README.md / current platform and language baseline
\item Recommended disposition: quarantined until current repo and external facts are checked
\end{itemize}

\subsection{CONTRA-0079 - stale\_external\_claim}

\begin{itemize}
\item Source conversation: dominium\_domino\_codex\_planning
\item Source file: docs/archive/conversations/dominium\_domino\_codex\_planning/dominium\_domino\_codex\_planning\_\_human\_readable\_chat\_report.txt
\item Claim: Archived conversation contains potentially stale external or baseline claim: C++98.
\item Conflict target: README.md / current platform and language baseline
\item Recommended disposition: quarantined until current repo and external facts are checked
\end{itemize}

\subsection{CONTRA-0080 - stale\_external\_claim}

\begin{itemize}
\item Source conversation: dominium\_domino\_codex\_planning
\item Source file: docs/archive/conversations/dominium\_domino\_codex\_planning/dominium\_domino\_codex\_planning\_\_human\_readable\_chat\_report.txt
\item Claim: Archived conversation contains potentially stale external or baseline claim: ISO C89.
\item Conflict target: README.md / current platform and language baseline
\item Recommended disposition: quarantined until current repo and external facts are checked
\end{itemize}

\subsection{CONTRA-0081 - stale\_external\_claim}

\begin{itemize}
\item Source conversation: dominium\_domino\_codex\_planning
\item Source file: docs/archive/conversations/dominium\_domino\_codex\_planning/dominium\_domino\_codex\_planning\_\_human\_readable\_chat\_report.txt
\item Claim: Archived conversation contains potentially stale external or baseline claim: CP/M.
\item Conflict target: README.md / current platform and language baseline
\item Recommended disposition: quarantined until current repo and external facts are checked
\end{itemize}

\subsection{CONTRA-0082 - stale\_external\_claim}

\begin{itemize}
\item Source conversation: dominium\_domino\_codex\_planning
\item Source file: docs/archive/conversations/dominium\_domino\_codex\_planning/dominium\_domino\_codex\_planning\_\_human\_readable\_chat\_report.txt
\item Claim: Archived conversation contains potentially stale external or baseline claim: Android.
\item Conflict target: README.md / current platform and language baseline
\item Recommended disposition: quarantined until current repo and external facts are checked
\end{itemize}

\subsection{CONTRA-0083 - conversation\_vs\_current\_queue}

\begin{itemize}
\item Source conversation: dominium\_full\_conversation
\item Source file: docs/archive/conversations/dominium\_full\_conversation/companion\_report/dominium\_full\_conversation\_companion\_report\_\_01\_complete\_human\_report.md
\item Claim: Archived conversation discusses work related to blocked scope runtime\_module\_loader.
\item Conflict target: .aide/queue/current.toml
\item Recommended disposition: preserve\_as\_history; require current queue authority before action
\end{itemize}

\subsection{CONTRA-0084 - conversation\_vs\_current\_queue}

\begin{itemize}
\item Source conversation: dominium\_full\_conversation
\item Source file: docs/archive/conversations/dominium\_full\_conversation/companion\_report/dominium\_full\_conversation\_companion\_report\_\_01\_complete\_human\_report.md
\item Claim: Archived conversation discusses work related to blocked scope provider\_runtime.
\item Conflict target: .aide/queue/current.toml
\item Recommended disposition: preserve\_as\_history; require current queue authority before action
\end{itemize}

\subsection{CONTRA-0085 - conversation\_vs\_current\_queue}

\begin{itemize}
\item Source conversation: dominium\_full\_conversation
\item Source file: docs/archive/conversations/dominium\_full\_conversation/companion\_report/dominium\_full\_conversation\_companion\_report\_\_01\_complete\_human\_report.md
\item Claim: Archived conversation discusses work related to blocked scope package\_runtime.
\item Conflict target: .aide/queue/current.toml
\item Recommended disposition: preserve\_as\_history; require current queue authority before action
\end{itemize}

\subsection{CONTRA-0086 - conversation\_vs\_current\_queue}

\begin{itemize}
\item Source conversation: dominium\_full\_conversation
\item Source file: docs/archive/conversations/dominium\_full\_conversation/companion\_report/dominium\_full\_conversation\_companion\_report\_\_01\_complete\_human\_report.md
\item Claim: Archived conversation discusses work related to blocked scope gameplay.
\item Conflict target: .aide/queue/current.toml
\item Recommended disposition: preserve\_as\_history; require current queue authority before action
\end{itemize}

\subsection{CONTRA-0087 - conversation\_vs\_current\_queue}

\begin{itemize}
\item Source conversation: dominium\_full\_conversation
\item Source file: docs/archive/conversations/dominium\_full\_conversation/companion\_report/dominium\_full\_conversation\_companion\_report\_\_01\_complete\_human\_report.md
\item Claim: Archived conversation discusses work related to blocked scope native\_gui.
\item Conflict target: .aide/queue/current.toml
\item Recommended disposition: preserve\_as\_history; require current queue authority before action
\end{itemize}

\subsection{CONTRA-0088 - stale\_external\_claim}

\begin{itemize}
\item Source conversation: dominium\_full\_conversation
\item Source file: docs/archive/conversations/dominium\_full\_conversation/companion\_report/dominium\_full\_conversation\_companion\_report\_\_01\_complete\_human\_report.md
\item Claim: Archived conversation contains potentially stale external or baseline claim: SDK.
\item Conflict target: README.md / current platform and language baseline
\item Recommended disposition: quarantined until current repo and external facts are checked
\end{itemize}

\subsection{CONTRA-0089 - conversation\_vs\_current\_queue}

\begin{itemize}
\item Source conversation: dominium\_setup
\item Source file: docs/archive/conversations/dominium\_setup/dominium\_setup\_handoff\_\_human\_readable\_chat\_report.txt
\item Claim: Archived conversation discusses work related to blocked scope gameplay.
\item Conflict target: .aide/queue/current.toml
\item Recommended disposition: preserve\_as\_history; require current queue authority before action
\end{itemize}

\subsection{CONTRA-0090 - conversation\_vs\_current\_queue}

\begin{itemize}
\item Source conversation: dominium\_setup
\item Source file: docs/archive/conversations/dominium\_setup/dominium\_setup\_handoff\_\_human\_readable\_chat\_report.txt
\item Claim: Archived conversation discusses work related to blocked scope release\_publication.
\item Conflict target: .aide/queue/current.toml
\item Recommended disposition: preserve\_as\_history; require current queue authority before action
\end{itemize}

\subsection{CONTRA-0091 - stale\_external\_claim}

\begin{itemize}
\item Source conversation: dominium\_setup
\item Source file: docs/archive/conversations/dominium\_setup/dominium\_setup\_handoff\_\_human\_readable\_chat\_report.txt
\item Claim: Archived conversation contains potentially stale external or baseline claim: SDK.
\item Conflict target: README.md / current platform and language baseline
\item Recommended disposition: quarantined until current repo and external facts are checked
\end{itemize}

\subsection{CONTRA-0092 - conversation\_vs\_current\_queue}

\begin{itemize}
\item Source conversation: Domino\_Dominium\_Workbench
\item Source file: docs/archive/conversations/Domino\_Dominium\_Workbench/dominium\_workbench\_full\_conversation\_companion\_\_human\_readable\_detailed\_summary\_report.md
\item Claim: Archived conversation discusses work related to blocked scope broad\_workbench\_ui.
\item Conflict target: .aide/queue/current.toml
\item Recommended disposition: preserve\_as\_history; require current queue authority before action
\end{itemize}

\subsection{CONTRA-0093 - conversation\_vs\_current\_queue}

\begin{itemize}
\item Source conversation: Domino\_Dominium\_Workbench
\item Source file: docs/archive/conversations/Domino\_Dominium\_Workbench/dominium\_workbench\_full\_conversation\_companion\_\_human\_readable\_detailed\_summary\_report.md
\item Claim: Archived conversation discusses work related to blocked scope provider\_runtime.
\item Conflict target: .aide/queue/current.toml
\item Recommended disposition: preserve\_as\_history; require current queue authority before action
\end{itemize}

\subsection{CONTRA-0094 - conversation\_vs\_current\_queue}

\begin{itemize}
\item Source conversation: Domino\_Dominium\_Workbench
\item Source file: docs/archive/conversations/Domino\_Dominium\_Workbench/dominium\_workbench\_full\_conversation\_companion\_\_human\_readable\_detailed\_summary\_report.md
\item Claim: Archived conversation discusses work related to blocked scope gameplay.
\item Conflict target: .aide/queue/current.toml
\item Recommended disposition: preserve\_as\_history; require current queue authority before action
\end{itemize}

\subsection{CONTRA-0095 - conversation\_vs\_current\_queue}

\begin{itemize}
\item Source conversation: Domino\_Dominium\_Workbench
\item Source file: docs/archive/conversations/Domino\_Dominium\_Workbench/dominium\_workbench\_full\_conversation\_companion\_\_human\_readable\_detailed\_summary\_report.md
\item Claim: Archived conversation discusses work related to blocked scope native\_gui.
\item Conflict target: .aide/queue/current.toml
\item Recommended disposition: preserve\_as\_history; require current queue authority before action
\end{itemize}

\subsection{CONTRA-0096 - stale\_external\_claim}

\begin{itemize}
\item Source conversation: Domino\_Dominium\_Workbench
\item Source file: docs/archive/conversations/Domino\_Dominium\_Workbench/dominium\_workbench\_full\_conversation\_companion\_\_human\_readable\_detailed\_summary\_report.md
\item Claim: Archived conversation contains potentially stale external or baseline claim: SDK.
\item Conflict target: README.md / current platform and language baseline
\item Recommended disposition: quarantined until current repo and external facts are checked
\end{itemize}

\subsection{CONTRA-0097 - conversation\_vs\_current\_queue}

\begin{itemize}
\item Source conversation: domino\_engine\_refactor\_prompts
\item Source file: docs/archive/conversations/domino\_engine\_refactor\_prompts/dominium\_domino\_engine\_refactor\_prompts\_\_human\_readable\_chat\_report.txt
\item Claim: Archived conversation discusses work related to blocked scope gameplay.
\item Conflict target: .aide/queue/current.toml
\item Recommended disposition: preserve\_as\_history; require current queue authority before action
\end{itemize}

\subsection{CONTRA-0098 - stale\_external\_claim}

\begin{itemize}
\item Source conversation: domino\_engine\_refactor\_prompts
\item Source file: docs/archive/conversations/domino\_engine\_refactor\_prompts/dominium\_domino\_engine\_refactor\_prompts\_\_human\_readable\_chat\_report.txt
\item Claim: Archived conversation contains potentially stale external or baseline claim: C89.
\item Conflict target: README.md / current platform and language baseline
\item Recommended disposition: quarantined until current repo and external facts are checked
\end{itemize}

\subsection{CONTRA-0099 - stale\_external\_claim}

\begin{itemize}
\item Source conversation: domino\_engine\_refactor\_prompts
\item Source file: docs/archive/conversations/domino\_engine\_refactor\_prompts/dominium\_domino\_engine\_refactor\_prompts\_\_human\_readable\_chat\_report.txt
\item Claim: Archived conversation contains potentially stale external or baseline claim: C++98.
\item Conflict target: README.md / current platform and language baseline
\item Recommended disposition: quarantined until current repo and external facts are checked
\end{itemize}

\subsection{CONTRA-0100 - stale\_external\_claim}

\begin{itemize}
\item Source conversation: domino\_engine\_refactor\_prompts
\item Source file: docs/archive/conversations/domino\_engine\_refactor\_prompts/dominium\_domino\_engine\_refactor\_prompts\_\_human\_readable\_chat\_report.txt
\item Claim: Archived conversation contains potentially stale external or baseline claim: ISO C89.
\item Conflict target: README.md / current platform and language baseline
\item Recommended disposition: quarantined until current repo and external facts are checked
\end{itemize}

\subsection{CONTRA-0101 - conversation\_vs\_current\_queue}

\begin{itemize}
\item Source conversation: engine\_baseline\_architecture
\item Source file: docs/archive/conversations/engine\_baseline\_architecture/domino\_dominium\_engine\_baseline\_architecture\_\_01\_human\_readable\_report.md
\item Claim: Archived conversation discusses work related to blocked scope provider\_runtime.
\item Conflict target: .aide/queue/current.toml
\item Recommended disposition: preserve\_as\_history; require current queue authority before action
\end{itemize}

\subsection{CONTRA-0102 - conversation\_vs\_current\_queue}

\begin{itemize}
\item Source conversation: engine\_baseline\_architecture
\item Source file: docs/archive/conversations/engine\_baseline\_architecture/domino\_dominium\_engine\_baseline\_architecture\_\_01\_human\_readable\_report.md
\item Claim: Archived conversation discusses work related to blocked scope gameplay.
\item Conflict target: .aide/queue/current.toml
\item Recommended disposition: preserve\_as\_history; require current queue authority before action
\end{itemize}

\subsection{CONTRA-0103 - conversation\_vs\_current\_queue}

\begin{itemize}
\item Source conversation: Foundation\_Workbench\_Codex
\item Source file: docs/archive/conversations/Foundation\_Workbench\_Codex/Dominium\_Foundation\_Workbench\_Parallel\_Codex\_Preservation\_\_01\_human\_readable\_report.md
\item Claim: Archived conversation discusses work related to blocked scope broad\_workbench\_ui.
\item Conflict target: .aide/queue/current.toml
\item Recommended disposition: preserve\_as\_history; require current queue authority before action
\end{itemize}

\subsection{CONTRA-0104 - conversation\_vs\_current\_queue}

\begin{itemize}
\item Source conversation: Foundation\_Workbench\_Codex
\item Source file: docs/archive/conversations/Foundation\_Workbench\_Codex/Dominium\_Foundation\_Workbench\_Parallel\_Codex\_Preservation\_\_01\_human\_readable\_report.md
\item Claim: Archived conversation discusses work related to blocked scope runtime\_module\_loader.
\item Conflict target: .aide/queue/current.toml
\item Recommended disposition: preserve\_as\_history; require current queue authority before action
\end{itemize}

\subsection{CONTRA-0105 - conversation\_vs\_current\_queue}

\begin{itemize}
\item Source conversation: Foundation\_Workbench\_Codex
\item Source file: docs/archive/conversations/Foundation\_Workbench\_Codex/Dominium\_Foundation\_Workbench\_Parallel\_Codex\_Preservation\_\_01\_human\_readable\_report.md
\item Claim: Archived conversation discusses work related to blocked scope provider\_runtime.
\item Conflict target: .aide/queue/current.toml
\item Recommended disposition: preserve\_as\_history; require current queue authority before action
\end{itemize}

\subsection{CONTRA-0106 - conversation\_vs\_current\_queue}

\begin{itemize}
\item Source conversation: Foundation\_Workbench\_Codex
\item Source file: docs/archive/conversations/Foundation\_Workbench\_Codex/Dominium\_Foundation\_Workbench\_Parallel\_Codex\_Preservation\_\_01\_human\_readable\_report.md
\item Claim: Archived conversation discusses work related to blocked scope package\_runtime.
\item Conflict target: .aide/queue/current.toml
\item Recommended disposition: preserve\_as\_history; require current queue authority before action
\end{itemize}

\subsection{CONTRA-0107 - conversation\_vs\_current\_queue}

\begin{itemize}
\item Source conversation: Foundation\_Workbench\_Codex
\item Source file: docs/archive/conversations/Foundation\_Workbench\_Codex/Dominium\_Foundation\_Workbench\_Parallel\_Codex\_Preservation\_\_01\_human\_readable\_report.md
\item Claim: Archived conversation discusses work related to blocked scope gameplay.
\item Conflict target: .aide/queue/current.toml
\item Recommended disposition: preserve\_as\_history; require current queue authority before action
\end{itemize}

\subsection{CONTRA-0108 - conversation\_vs\_current\_queue}

\begin{itemize}
\item Source conversation: Foundation\_Workbench\_Codex
\item Source file: docs/archive/conversations/Foundation\_Workbench\_Codex/Dominium\_Foundation\_Workbench\_Parallel\_Codex\_Preservation\_\_01\_human\_readable\_report.md
\item Claim: Archived conversation discusses work related to blocked scope renderer\_implementation.
\item Conflict target: .aide/queue/current.toml
\item Recommended disposition: preserve\_as\_history; require current queue authority before action
\end{itemize}

\subsection{CONTRA-0109 - conversation\_vs\_current\_queue}

\begin{itemize}
\item Source conversation: Foundation\_Workbench\_Codex
\item Source file: docs/archive/conversations/Foundation\_Workbench\_Codex/Dominium\_Foundation\_Workbench\_Parallel\_Codex\_Preservation\_\_01\_human\_readable\_report.md
\item Claim: Archived conversation discusses work related to blocked scope native\_gui.
\item Conflict target: .aide/queue/current.toml
\item Recommended disposition: preserve\_as\_history; require current queue authority before action
\end{itemize}

\subsection{CONTRA-0110 - conversation\_vs\_current\_queue}

\begin{itemize}
\item Source conversation: Foundation\_Workbench\_Codex
\item Source file: docs/archive/conversations/Foundation\_Workbench\_Codex/Dominium\_Foundation\_Workbench\_Parallel\_Codex\_Preservation\_\_01\_human\_readable\_report.md
\item Claim: Archived conversation discusses work related to blocked scope release\_publication.
\item Conflict target: .aide/queue/current.toml
\item Recommended disposition: preserve\_as\_history; require current queue authority before action
\end{itemize}

\subsection{CONTRA-0111 - stale\_external\_claim}

\begin{itemize}
\item Source conversation: Foundation\_Workbench\_Codex
\item Source file: docs/archive/conversations/Foundation\_Workbench\_Codex/Dominium\_Foundation\_Workbench\_Parallel\_Codex\_Preservation\_\_01\_human\_readable\_report.md
\item Claim: Archived conversation contains potentially stale external or baseline claim: C89.
\item Conflict target: README.md / current platform and language baseline
\item Recommended disposition: quarantined until current repo and external facts are checked
\end{itemize}

\subsection{CONTRA-0112 - stale\_external\_claim}

\begin{itemize}
\item Source conversation: Foundation\_Workbench\_Codex
\item Source file: docs/archive/conversations/Foundation\_Workbench\_Codex/Dominium\_Foundation\_Workbench\_Parallel\_Codex\_Preservation\_\_01\_human\_readable\_report.md
\item Claim: Archived conversation contains potentially stale external or baseline claim: C++98.
\item Conflict target: README.md / current platform and language baseline
\item Recommended disposition: quarantined until current repo and external facts are checked
\end{itemize}

\subsection{CONTRA-0113 - stale\_external\_claim}

\begin{itemize}
\item Source conversation: Foundation\_Workbench\_Codex
\item Source file: docs/archive/conversations/Foundation\_Workbench\_Codex/Dominium\_Foundation\_Workbench\_Parallel\_Codex\_Preservation\_\_01\_human\_readable\_report.md
\item Claim: Archived conversation contains potentially stale external or baseline claim: SDK.
\item Conflict target: README.md / current platform and language baseline
\item Recommended disposition: quarantined until current repo and external facts are checked
\end{itemize}

\subsection{CONTRA-0114 - conversation\_vs\_current\_queue}

\begin{itemize}
\item Source conversation: Framework\_Open\_Source\_Provider
\item Source file: docs/archive/conversations/Framework\_Open\_Source\_Provider/Domino\_Framework\_Open\_Source\_Provider\_Architecture\_\_01\_human\_readable\_report.md
\item Claim: Archived conversation discusses work related to blocked scope broad\_workbench\_ui.
\item Conflict target: .aide/queue/current.toml
\item Recommended disposition: preserve\_as\_history; require current queue authority before action
\end{itemize}

\subsection{CONTRA-0115 - conversation\_vs\_current\_queue}

\begin{itemize}
\item Source conversation: Framework\_Open\_Source\_Provider
\item Source file: docs/archive/conversations/Framework\_Open\_Source\_Provider/Domino\_Framework\_Open\_Source\_Provider\_Architecture\_\_01\_human\_readable\_report.md
\item Claim: Archived conversation discusses work related to blocked scope provider\_runtime.
\item Conflict target: .aide/queue/current.toml
\item Recommended disposition: preserve\_as\_history; require current queue authority before action
\end{itemize}

\subsection{CONTRA-0116 - conversation\_vs\_current\_queue}

\begin{itemize}
\item Source conversation: Framework\_Open\_Source\_Provider
\item Source file: docs/archive/conversations/Framework\_Open\_Source\_Provider/Domino\_Framework\_Open\_Source\_Provider\_Architecture\_\_01\_human\_readable\_report.md
\item Claim: Archived conversation discusses work related to blocked scope gameplay.
\item Conflict target: .aide/queue/current.toml
\item Recommended disposition: preserve\_as\_history; require current queue authority before action
\end{itemize}

\subsection{CONTRA-0117 - conversation\_vs\_current\_queue}

\begin{itemize}
\item Source conversation: Framework\_Open\_Source\_Provider
\item Source file: docs/archive/conversations/Framework\_Open\_Source\_Provider/Domino\_Framework\_Open\_Source\_Provider\_Architecture\_\_01\_human\_readable\_report.md
\item Claim: Archived conversation discusses work related to blocked scope renderer\_implementation.
\item Conflict target: .aide/queue/current.toml
\item Recommended disposition: preserve\_as\_history; require current queue authority before action
\end{itemize}

\subsection{CONTRA-0118 - stale\_external\_claim}

\begin{itemize}
\item Source conversation: Framework\_Open\_Source\_Provider
\item Source file: docs/archive/conversations/Framework\_Open\_Source\_Provider/Domino\_Framework\_Open\_Source\_Provider\_Architecture\_\_01\_human\_readable\_report.md
\item Claim: Archived conversation contains potentially stale external or baseline claim: C++98.
\item Conflict target: README.md / current platform and language baseline
\item Recommended disposition: quarantined until current repo and external facts are checked
\end{itemize}

\subsection{CONTRA-0119 - stale\_external\_claim}

\begin{itemize}
\item Source conversation: Framework\_Open\_Source\_Provider
\item Source file: docs/archive/conversations/Framework\_Open\_Source\_Provider/Domino\_Framework\_Open\_Source\_Provider\_Architecture\_\_01\_human\_readable\_report.md
\item Claim: Archived conversation contains potentially stale external or baseline claim: SDK.
\item Conflict target: README.md / current platform and language baseline
\item Recommended disposition: quarantined until current repo and external facts are checked
\end{itemize}

\subsection{CONTRA-0120 - conversation\_vs\_current\_queue}

\begin{itemize}
\item Source conversation: gui\_binary\_content
\item Source file: docs/archive/conversations/gui\_binary\_content/dominium\_gui\_binary\_content\_handoff\_\_human\_readable\_chat\_report.txt
\item Claim: Archived conversation discusses work related to blocked scope gameplay.
\item Conflict target: .aide/queue/current.toml
\item Recommended disposition: preserve\_as\_history; require current queue authority before action
\end{itemize}

\subsection{CONTRA-0121 - conversation\_vs\_current\_queue}

\begin{itemize}
\item Source conversation: gui\_binary\_content
\item Source file: docs/archive/conversations/gui\_binary\_content/dominium\_gui\_binary\_content\_handoff\_\_human\_readable\_chat\_report.txt
\item Claim: Archived conversation discusses work related to blocked scope native\_gui.
\item Conflict target: .aide/queue/current.toml
\item Recommended disposition: preserve\_as\_history; require current queue authority before action
\end{itemize}

\subsection{CONTRA-0122 - stale\_external\_claim}

\begin{itemize}
\item Source conversation: gui\_binary\_content
\item Source file: docs/archive/conversations/gui\_binary\_content/dominium\_gui\_binary\_content\_handoff\_\_human\_readable\_chat\_report.txt
\item Claim: Archived conversation contains potentially stale external or baseline claim: iOS.
\item Conflict target: README.md / current platform and language baseline
\item Recommended disposition: quarantined until current repo and external facts are checked
\end{itemize}

\subsection{CONTRA-0123 - stale\_external\_claim}

\begin{itemize}
\item Source conversation: gui\_binary\_content
\item Source file: docs/archive/conversations/gui\_binary\_content/dominium\_gui\_binary\_content\_handoff\_\_human\_readable\_chat\_report.txt
\item Claim: Archived conversation contains potentially stale external or baseline claim: Android.
\item Conflict target: README.md / current platform and language baseline
\item Recommended disposition: quarantined until current repo and external facts are checked
\end{itemize}

\subsection{CONTRA-0124 - stale\_external\_claim}

\begin{itemize}
\item Source conversation: gui\_binary\_content
\item Source file: docs/archive/conversations/gui\_binary\_content/dominium\_gui\_binary\_content\_handoff\_\_human\_readable\_chat\_report.txt
\item Claim: Archived conversation contains potentially stale external or baseline claim: SDK.
\item Conflict target: README.md / current platform and language baseline
\item Recommended disposition: quarantined until current repo and external facts are checked
\end{itemize}

\subsection{CONTRA-0125 - conversation\_vs\_current\_queue}

\begin{itemize}
\item Source conversation: Language\_Platform\_Architecture
\item Source file: docs/archive/conversations/Language\_Platform\_Architecture/Dominium\_Language\_Platform\_Architecture\_Baseline\_\_01\_human\_readable\_report.md
\item Claim: Archived conversation discusses work related to blocked scope broad\_workbench\_ui.
\item Conflict target: .aide/queue/current.toml
\item Recommended disposition: preserve\_as\_history; require current queue authority before action
\end{itemize}

\subsection{CONTRA-0126 - conversation\_vs\_current\_queue}

\begin{itemize}
\item Source conversation: Language\_Platform\_Architecture
\item Source file: docs/archive/conversations/Language\_Platform\_Architecture/Dominium\_Language\_Platform\_Architecture\_Baseline\_\_01\_human\_readable\_report.md
\item Claim: Archived conversation discusses work related to blocked scope provider\_runtime.
\item Conflict target: .aide/queue/current.toml
\item Recommended disposition: preserve\_as\_history; require current queue authority before action
\end{itemize}

\subsection{CONTRA-0127 - conversation\_vs\_current\_queue}

\begin{itemize}
\item Source conversation: Language\_Platform\_Architecture
\item Source file: docs/archive/conversations/Language\_Platform\_Architecture/Dominium\_Language\_Platform\_Architecture\_Baseline\_\_01\_human\_readable\_report.md
\item Claim: Archived conversation discusses work related to blocked scope gameplay.
\item Conflict target: .aide/queue/current.toml
\item Recommended disposition: preserve\_as\_history; require current queue authority before action
\end{itemize}

\subsection{CONTRA-0128 - conversation\_vs\_current\_queue}

\begin{itemize}
\item Source conversation: Language\_Platform\_Architecture
\item Source file: docs/archive/conversations/Language\_Platform\_Architecture/Dominium\_Language\_Platform\_Architecture\_Baseline\_\_01\_human\_readable\_report.md
\item Claim: Archived conversation discusses work related to blocked scope renderer\_implementation.
\item Conflict target: .aide/queue/current.toml
\item Recommended disposition: preserve\_as\_history; require current queue authority before action
\end{itemize}

\subsection{CONTRA-0129 - stale\_external\_claim}

\begin{itemize}
\item Source conversation: Language\_Platform\_Architecture
\item Source file: docs/archive/conversations/Language\_Platform\_Architecture/Dominium\_Language\_Platform\_Architecture\_Baseline\_\_01\_human\_readable\_report.md
\item Claim: Archived conversation contains potentially stale external or baseline claim: C89.
\item Conflict target: README.md / current platform and language baseline
\item Recommended disposition: quarantined until current repo and external facts are checked
\end{itemize}

\subsection{CONTRA-0130 - stale\_external\_claim}

\begin{itemize}
\item Source conversation: Language\_Platform\_Architecture
\item Source file: docs/archive/conversations/Language\_Platform\_Architecture/Dominium\_Language\_Platform\_Architecture\_Baseline\_\_01\_human\_readable\_report.md
\item Claim: Archived conversation contains potentially stale external or baseline claim: C++98.
\item Conflict target: README.md / current platform and language baseline
\item Recommended disposition: quarantined until current repo and external facts are checked
\end{itemize}

\subsection{CONTRA-0131 - stale\_external\_claim}

\begin{itemize}
\item Source conversation: Language\_Platform\_Architecture
\item Source file: docs/archive/conversations/Language\_Platform\_Architecture/Dominium\_Language\_Platform\_Architecture\_Baseline\_\_01\_human\_readable\_report.md
\item Claim: Archived conversation contains potentially stale external or baseline claim: Windows 98.
\item Conflict target: README.md / current platform and language baseline
\item Recommended disposition: quarantined until current repo and external facts are checked
\end{itemize}

\subsection{CONTRA-0132 - stale\_external\_claim}

\begin{itemize}
\item Source conversation: Language\_Platform\_Architecture
\item Source file: docs/archive/conversations/Language\_Platform\_Architecture/Dominium\_Language\_Platform\_Architecture\_Baseline\_\_01\_human\_readable\_report.md
\item Claim: Archived conversation contains potentially stale external or baseline claim: Mac OS 9.
\item Conflict target: README.md / current platform and language baseline
\item Recommended disposition: quarantined until current repo and external facts are checked
\end{itemize}

\subsection{CONTRA-0133 - stale\_external\_claim}

\begin{itemize}
\item Source conversation: Language\_Platform\_Architecture
\item Source file: docs/archive/conversations/Language\_Platform\_Architecture/Dominium\_Language\_Platform\_Architecture\_Baseline\_\_01\_human\_readable\_report.md
\item Claim: Archived conversation contains potentially stale external or baseline claim: iOS.
\item Conflict target: README.md / current platform and language baseline
\item Recommended disposition: quarantined until current repo and external facts are checked
\end{itemize}

\subsection{CONTRA-0134 - stale\_external\_claim}

\begin{itemize}
\item Source conversation: Language\_Platform\_Architecture
\item Source file: docs/archive/conversations/Language\_Platform\_Architecture/Dominium\_Language\_Platform\_Architecture\_Baseline\_\_01\_human\_readable\_report.md
\item Claim: Archived conversation contains potentially stale external or baseline claim: Android.
\item Conflict target: README.md / current platform and language baseline
\item Recommended disposition: quarantined until current repo and external facts are checked
\end{itemize}

\subsection{CONTRA-0135 - stale\_external\_claim}

\begin{itemize}
\item Source conversation: Language\_Platform\_Architecture
\item Source file: docs/archive/conversations/Language\_Platform\_Architecture/Dominium\_Language\_Platform\_Architecture\_Baseline\_\_01\_human\_readable\_report.md
\item Claim: Archived conversation contains potentially stale external or baseline claim: SDK.
\item Conflict target: README.md / current platform and language baseline
\item Recommended disposition: quarantined until current repo and external facts are checked
\end{itemize}

\subsection{CONTRA-0136 - conversation\_vs\_current\_queue}

\begin{itemize}
\item Source conversation: launcher\_app\_layer
\item Source file: docs/archive/conversations/launcher\_app\_layer/dominium\_launcher\_app\_layer\_handoff\_\_human\_readable\_chat\_report.txt
\item Claim: Archived conversation discusses work related to blocked scope gameplay.
\item Conflict target: .aide/queue/current.toml
\item Recommended disposition: preserve\_as\_history; require current queue authority before action
\end{itemize}

\subsection{CONTRA-0137 - stale\_external\_claim}

\begin{itemize}
\item Source conversation: launcher\_app\_layer
\item Source file: docs/archive/conversations/launcher\_app\_layer/dominium\_launcher\_app\_layer\_handoff\_\_human\_readable\_chat\_report.txt
\item Claim: Archived conversation contains potentially stale external or baseline claim: SDK.
\item Conflict target: README.md / current platform and language baseline
\item Recommended disposition: quarantined until current repo and external facts are checked
\end{itemize}

\subsection{CONTRA-0138 - conversation\_vs\_current\_queue}

\begin{itemize}
\item Source conversation: Launcher\_Setup\_Architecture
\item Source file: docs/archive/conversations/Launcher\_Setup\_Architecture/Dominium\_Launcher\_Setup\_Architecture\_\_human\_readable\_chat\_report.txt
\item Claim: Archived conversation discusses work related to blocked scope gameplay.
\item Conflict target: .aide/queue/current.toml
\item Recommended disposition: preserve\_as\_history; require current queue authority before action
\end{itemize}

\subsection{CONTRA-0139 - conversation\_vs\_current\_queue}

\begin{itemize}
\item Source conversation: Launcher\_Setup\_Architecture
\item Source file: docs/archive/conversations/Launcher\_Setup\_Architecture/Dominium\_Launcher\_Setup\_Architecture\_\_human\_readable\_chat\_report.txt
\item Claim: Archived conversation discusses work related to blocked scope renderer\_implementation.
\item Conflict target: .aide/queue/current.toml
\item Recommended disposition: preserve\_as\_history; require current queue authority before action
\end{itemize}

\subsection{CONTRA-0140 - stale\_external\_claim}

\begin{itemize}
\item Source conversation: Launcher\_Setup\_Architecture
\item Source file: docs/archive/conversations/Launcher\_Setup\_Architecture/Dominium\_Launcher\_Setup\_Architecture\_\_human\_readable\_chat\_report.txt
\item Claim: Archived conversation contains potentially stale external or baseline claim: C89.
\item Conflict target: README.md / current platform and language baseline
\item Recommended disposition: quarantined until current repo and external facts are checked
\end{itemize}

\subsection{CONTRA-0141 - stale\_external\_claim}

\begin{itemize}
\item Source conversation: Launcher\_Setup\_Architecture
\item Source file: docs/archive/conversations/Launcher\_Setup\_Architecture/Dominium\_Launcher\_Setup\_Architecture\_\_human\_readable\_chat\_report.txt
\item Claim: Archived conversation contains potentially stale external or baseline claim: C++98.
\item Conflict target: README.md / current platform and language baseline
\item Recommended disposition: quarantined until current repo and external facts are checked
\end{itemize}

\subsection{CONTRA-0142 - stale\_external\_claim}

\begin{itemize}
\item Source conversation: Launcher\_Setup\_Architecture
\item Source file: docs/archive/conversations/Launcher\_Setup\_Architecture/Dominium\_Launcher\_Setup\_Architecture\_\_human\_readable\_chat\_report.txt
\item Claim: Archived conversation contains potentially stale external or baseline claim: CP/M.
\item Conflict target: README.md / current platform and language baseline
\item Recommended disposition: quarantined until current repo and external facts are checked
\end{itemize}

\subsection{CONTRA-0143 - stale\_external\_claim}

\begin{itemize}
\item Source conversation: Launcher\_Setup\_Architecture
\item Source file: docs/archive/conversations/Launcher\_Setup\_Architecture/Dominium\_Launcher\_Setup\_Architecture\_\_human\_readable\_chat\_report.txt
\item Claim: Archived conversation contains potentially stale external or baseline claim: SDK.
\item Conflict target: README.md / current platform and language baseline
\item Recommended disposition: quarantined until current repo and external facts are checked
\end{itemize}

\subsection{CONTRA-0144 - conversation\_vs\_current\_queue}

\begin{itemize}
\item Source conversation: Modularity\_AIDE\_Refactorability
\item Source file: docs/archive/conversations/Modularity\_AIDE\_Refactorability/Dominium\_Modularity\_AIDE\_Refactorability\_Architecture\_\_01\_human\_readable\_report.md
\item Claim: Archived conversation discusses work related to blocked scope renderer\_implementation.
\item Conflict target: .aide/queue/current.toml
\item Recommended disposition: preserve\_as\_history; require current queue authority before action
\end{itemize}

\subsection{CONTRA-0145 - conversation\_vs\_current\_queue}

\begin{itemize}
\item Source conversation: Modularity\_AIDE\_Refactorability
\item Source file: docs/archive/conversations/Modularity\_AIDE\_Refactorability/Dominium\_Modularity\_AIDE\_Refactorability\_Architecture\_\_01\_human\_readable\_report.md
\item Claim: Archived conversation discusses work related to blocked scope native\_gui.
\item Conflict target: .aide/queue/current.toml
\item Recommended disposition: preserve\_as\_history; require current queue authority before action
\end{itemize}

\subsection{CONTRA-0146 - stale\_external\_claim}

\begin{itemize}
\item Source conversation: Modularity\_AIDE\_Refactorability
\item Source file: docs/archive/conversations/Modularity\_AIDE\_Refactorability/Dominium\_Modularity\_AIDE\_Refactorability\_Architecture\_\_01\_human\_readable\_report.md
\item Claim: Archived conversation contains potentially stale external or baseline claim: C89.
\item Conflict target: README.md / current platform and language baseline
\item Recommended disposition: quarantined until current repo and external facts are checked
\end{itemize}

\subsection{CONTRA-0147 - conversation\_vs\_current\_queue}

\begin{itemize}
\item Source conversation: omega\_xi\_pi\_architecture\_future
\item Source file: docs/archive/conversations/omega\_xi\_pi\_architecture\_future/dominium\_omega\_xi\_pi\_architecture\_future\_proofing\_planning\_\_01\_human\_readable\_report.md
\item Claim: Archived conversation discusses work related to blocked scope runtime\_module\_loader.
\item Conflict target: .aide/queue/current.toml
\item Recommended disposition: preserve\_as\_history; require current queue authority before action
\end{itemize}

\subsection{CONTRA-0148 - conversation\_vs\_current\_queue}

\begin{itemize}
\item Source conversation: omega\_xi\_pi\_architecture\_future
\item Source file: docs/archive/conversations/omega\_xi\_pi\_architecture\_future/dominium\_omega\_xi\_pi\_architecture\_future\_proofing\_planning\_\_01\_human\_readable\_report.md
\item Claim: Archived conversation discusses work related to blocked scope native\_gui.
\item Conflict target: .aide/queue/current.toml
\item Recommended disposition: preserve\_as\_history; require current queue authority before action
\end{itemize}

\subsection{CONTRA-0149 - conversation\_vs\_current\_queue}

\begin{itemize}
\item Source conversation: OS\_Interface\_Repo\_Architecture
\item Source file: docs/archive/conversations/OS\_Interface\_Repo\_Architecture/Dominium\_OS\_Interface\_Repo\_Architecture\_\_01\_human\_readable\_report.md
\item Claim: Archived conversation discusses work related to blocked scope gameplay.
\item Conflict target: .aide/queue/current.toml
\item Recommended disposition: preserve\_as\_history; require current queue authority before action
\end{itemize}

\subsection{CONTRA-0150 - conversation\_vs\_current\_queue}

\begin{itemize}
\item Source conversation: OS\_Interface\_Repo\_Architecture
\item Source file: docs/archive/conversations/OS\_Interface\_Repo\_Architecture/Dominium\_OS\_Interface\_Repo\_Architecture\_\_01\_human\_readable\_report.md
\item Claim: Archived conversation discusses work related to blocked scope native\_gui.
\item Conflict target: .aide/queue/current.toml
\item Recommended disposition: preserve\_as\_history; require current queue authority before action
\end{itemize}

\subsection{CONTRA-0151 - conversation\_vs\_current\_queue}

\begin{itemize}
\item Source conversation: OS\_Interface\_Repo\_Architecture
\item Source file: docs/archive/conversations/OS\_Interface\_Repo\_Architecture/Dominium\_OS\_Interface\_Repo\_Architecture\_\_01\_human\_readable\_report.md
\item Claim: Archived conversation discusses work related to blocked scope release\_publication.
\item Conflict target: .aide/queue/current.toml
\item Recommended disposition: preserve\_as\_history; require current queue authority before action
\end{itemize}

\subsection{CONTRA-0152 - stale\_external\_claim}

\begin{itemize}
\item Source conversation: OS\_Interface\_Repo\_Architecture
\item Source file: docs/archive/conversations/OS\_Interface\_Repo\_Architecture/Dominium\_OS\_Interface\_Repo\_Architecture\_\_01\_human\_readable\_report.md
\item Claim: Archived conversation contains potentially stale external or baseline claim: Mac OS 8.
\item Conflict target: README.md / current platform and language baseline
\item Recommended disposition: quarantined until current repo and external facts are checked
\end{itemize}

\subsection{CONTRA-0153 - stale\_external\_claim}

\begin{itemize}
\item Source conversation: OS\_Interface\_Repo\_Architecture
\item Source file: docs/archive/conversations/OS\_Interface\_Repo\_Architecture/Dominium\_OS\_Interface\_Repo\_Architecture\_\_01\_human\_readable\_report.md
\item Claim: Archived conversation contains potentially stale external or baseline claim: Android.
\item Conflict target: README.md / current platform and language baseline
\item Recommended disposition: quarantined until current repo and external facts are checked
\end{itemize}

\subsection{CONTRA-0154 - stale\_external\_claim}

\begin{itemize}
\item Source conversation: OS\_Interface\_Repo\_Architecture
\item Source file: docs/archive/conversations/OS\_Interface\_Repo\_Architecture/Dominium\_OS\_Interface\_Repo\_Architecture\_\_01\_human\_readable\_report.md
\item Claim: Archived conversation contains potentially stale external or baseline claim: SDK.
\item Conflict target: README.md / current platform and language baseline
\item Recommended disposition: quarantined until current repo and external facts are checked
\end{itemize}

\subsection{CONTRA-0155 - conversation\_vs\_current\_queue}

\begin{itemize}
\item Source conversation: platform\_renderer\_api\_plan
\item Source file: docs/archive/conversations/platform\_renderer\_api\_plan/dominium\_codex\_platform\_renderer\_api\_plan\_\_human\_readable\_chat\_report.txt
\item Claim: Archived conversation discusses work related to blocked scope gameplay.
\item Conflict target: .aide/queue/current.toml
\item Recommended disposition: preserve\_as\_history; require current queue authority before action
\end{itemize}

\subsection{CONTRA-0156 - conversation\_vs\_current\_queue}

\begin{itemize}
\item Source conversation: platform\_renderer\_api\_plan
\item Source file: docs/archive/conversations/platform\_renderer\_api\_plan/dominium\_codex\_platform\_renderer\_api\_plan\_\_human\_readable\_chat\_report.txt
\item Claim: Archived conversation discusses work related to blocked scope renderer\_implementation.
\item Conflict target: .aide/queue/current.toml
\item Recommended disposition: preserve\_as\_history; require current queue authority before action
\end{itemize}

\subsection{CONTRA-0157 - conversation\_vs\_current\_queue}

\begin{itemize}
\item Source conversation: platform\_renderer\_api\_plan
\item Source file: docs/archive/conversations/platform\_renderer\_api\_plan/dominium\_codex\_platform\_renderer\_api\_plan\_\_human\_readable\_chat\_report.txt
\item Claim: Archived conversation discusses work related to blocked scope native\_gui.
\item Conflict target: .aide/queue/current.toml
\item Recommended disposition: preserve\_as\_history; require current queue authority before action
\end{itemize}

\subsection{CONTRA-0158 - stale\_external\_claim}

\begin{itemize}
\item Source conversation: platform\_renderer\_api\_plan
\item Source file: docs/archive/conversations/platform\_renderer\_api\_plan/dominium\_codex\_platform\_renderer\_api\_plan\_\_human\_readable\_chat\_report.txt
\item Claim: Archived conversation contains potentially stale external or baseline claim: C89.
\item Conflict target: README.md / current platform and language baseline
\item Recommended disposition: quarantined until current repo and external facts are checked
\end{itemize}

\subsection{CONTRA-0159 - stale\_external\_claim}

\begin{itemize}
\item Source conversation: platform\_renderer\_api\_plan
\item Source file: docs/archive/conversations/platform\_renderer\_api\_plan/dominium\_codex\_platform\_renderer\_api\_plan\_\_human\_readable\_chat\_report.txt
\item Claim: Archived conversation contains potentially stale external or baseline claim: C++98.
\item Conflict target: README.md / current platform and language baseline
\item Recommended disposition: quarantined until current repo and external facts are checked
\end{itemize}

\subsection{CONTRA-0160 - stale\_external\_claim}

\begin{itemize}
\item Source conversation: platform\_renderer\_api\_plan
\item Source file: docs/archive/conversations/platform\_renderer\_api\_plan/dominium\_codex\_platform\_renderer\_api\_plan\_\_human\_readable\_chat\_report.txt
\item Claim: Archived conversation contains potentially stale external or baseline claim: ISO C89.
\item Conflict target: README.md / current platform and language baseline
\item Recommended disposition: quarantined until current repo and external facts are checked
\end{itemize}

\subsection{CONTRA-0161 - stale\_external\_claim}

\begin{itemize}
\item Source conversation: platform\_renderer\_api\_plan
\item Source file: docs/archive/conversations/platform\_renderer\_api\_plan/dominium\_codex\_platform\_renderer\_api\_plan\_\_human\_readable\_chat\_report.txt
\item Claim: Archived conversation contains potentially stale external or baseline claim: CP/M.
\item Conflict target: README.md / current platform and language baseline
\item Recommended disposition: quarantined until current repo and external facts are checked
\end{itemize}

\subsection{CONTRA-0162 - stale\_external\_claim}

\begin{itemize}
\item Source conversation: platform\_renderer\_api\_plan
\item Source file: docs/archive/conversations/platform\_renderer\_api\_plan/dominium\_codex\_platform\_renderer\_api\_plan\_\_human\_readable\_chat\_report.txt
\item Claim: Archived conversation contains potentially stale external or baseline claim: Android.
\item Conflict target: README.md / current platform and language baseline
\item Recommended disposition: quarantined until current repo and external facts are checked
\end{itemize}

\subsection{CONTRA-0163 - stale\_external\_claim}

\begin{itemize}
\item Source conversation: platform\_renderer\_api\_plan
\item Source file: docs/archive/conversations/platform\_renderer\_api\_plan/dominium\_codex\_platform\_renderer\_api\_plan\_\_human\_readable\_chat\_report.txt
\item Claim: Archived conversation contains potentially stale external or baseline claim: SDK.
\item Conflict target: README.md / current platform and language baseline
\item Recommended disposition: quarantined until current repo and external facts are checked
\end{itemize}

\subsection{CONTRA-0164 - conversation\_vs\_current\_queue}

\begin{itemize}
\item Source conversation: platform\_support
\item Source file: docs/archive/conversations/platform\_support/dominium\_platform\_support\_planning\_\_human\_readable\_chat\_report.txt
\item Claim: Archived conversation discusses work related to blocked scope gameplay.
\item Conflict target: .aide/queue/current.toml
\item Recommended disposition: preserve\_as\_history; require current queue authority before action
\end{itemize}

\subsection{CONTRA-0165 - conversation\_vs\_current\_queue}

\begin{itemize}
\item Source conversation: platform\_support
\item Source file: docs/archive/conversations/platform\_support/dominium\_platform\_support\_planning\_\_human\_readable\_chat\_report.txt
\item Claim: Archived conversation discusses work related to blocked scope renderer\_implementation.
\item Conflict target: .aide/queue/current.toml
\item Recommended disposition: preserve\_as\_history; require current queue authority before action
\end{itemize}

\subsection{CONTRA-0166 - conversation\_vs\_current\_queue}

\begin{itemize}
\item Source conversation: platform\_support
\item Source file: docs/archive/conversations/platform\_support/dominium\_platform\_support\_planning\_\_human\_readable\_chat\_report.txt
\item Claim: Archived conversation discusses work related to blocked scope release\_publication.
\item Conflict target: .aide/queue/current.toml
\item Recommended disposition: preserve\_as\_history; require current queue authority before action
\end{itemize}

\subsection{CONTRA-0167 - stale\_external\_claim}

\begin{itemize}
\item Source conversation: platform\_support
\item Source file: docs/archive/conversations/platform\_support/dominium\_platform\_support\_planning\_\_human\_readable\_chat\_report.txt
\item Claim: Archived conversation contains potentially stale external or baseline claim: CP/M.
\item Conflict target: README.md / current platform and language baseline
\item Recommended disposition: quarantined until current repo and external facts are checked
\end{itemize}

\subsection{CONTRA-0168 - stale\_external\_claim}

\begin{itemize}
\item Source conversation: platform\_support
\item Source file: docs/archive/conversations/platform\_support/dominium\_platform\_support\_planning\_\_human\_readable\_chat\_report.txt
\item Claim: Archived conversation contains potentially stale external or baseline claim: iOS.
\item Conflict target: README.md / current platform and language baseline
\item Recommended disposition: quarantined until current repo and external facts are checked
\end{itemize}

\subsection{CONTRA-0169 - stale\_external\_claim}

\begin{itemize}
\item Source conversation: platform\_support
\item Source file: docs/archive/conversations/platform\_support/dominium\_platform\_support\_planning\_\_human\_readable\_chat\_report.txt
\item Claim: Archived conversation contains potentially stale external or baseline claim: Android.
\item Conflict target: README.md / current platform and language baseline
\item Recommended disposition: quarantined until current repo and external facts are checked
\end{itemize}

\subsection{CONTRA-0170 - stale\_external\_claim}

\begin{itemize}
\item Source conversation: platform\_support
\item Source file: docs/archive/conversations/platform\_support/dominium\_platform\_support\_planning\_\_human\_readable\_chat\_report.txt
\item Claim: Archived conversation contains potentially stale external or baseline claim: WebGPU.
\item Conflict target: README.md / current platform and language baseline
\item Recommended disposition: quarantined until current repo and external facts are checked
\end{itemize}

\subsection{CONTRA-0171 - stale\_external\_claim}

\begin{itemize}
\item Source conversation: platform\_support
\item Source file: docs/archive/conversations/platform\_support/dominium\_platform\_support\_planning\_\_human\_readable\_chat\_report.txt
\item Claim: Archived conversation contains potentially stale external or baseline claim: console program.
\item Conflict target: README.md / current platform and language baseline
\item Recommended disposition: quarantined until current repo and external facts are checked
\end{itemize}

\subsection{CONTRA-0172 - stale\_external\_claim}

\begin{itemize}
\item Source conversation: platform\_support
\item Source file: docs/archive/conversations/platform\_support/dominium\_platform\_support\_planning\_\_human\_readable\_chat\_report.txt
\item Claim: Archived conversation contains potentially stale external or baseline claim: SDK.
\item Conflict target: README.md / current platform and language baseline
\item Recommended disposition: quarantined until current repo and external facts are checked
\end{itemize}

\subsection{CONTRA-0173 - conversation\_vs\_current\_queue}

\begin{itemize}
\item Source conversation: Portability\_Assurance\_Future\_Proof
\item Source file: docs/archive/conversations/Portability\_Assurance\_Future\_Proof/Domino\_Dominium\_Portability\_Assurance\_Future\_Proof\_Architecture\_\_01\_human\_readable\_report.md
\item Claim: Archived conversation discusses work related to blocked scope gameplay.
\item Conflict target: .aide/queue/current.toml
\item Recommended disposition: preserve\_as\_history; require current queue authority before action
\end{itemize}

\subsection{CONTRA-0174 - stale\_external\_claim}

\begin{itemize}
\item Source conversation: Portability\_Assurance\_Future\_Proof
\item Source file: docs/archive/conversations/Portability\_Assurance\_Future\_Proof/Domino\_Dominium\_Portability\_Assurance\_Future\_Proof\_Architecture\_\_01\_human\_readable\_report.md
\item Claim: Archived conversation contains potentially stale external or baseline claim: C89.
\item Conflict target: README.md / current platform and language baseline
\item Recommended disposition: quarantined until current repo and external facts are checked
\end{itemize}

\subsection{CONTRA-0175 - stale\_external\_claim}

\begin{itemize}
\item Source conversation: Portability\_Assurance\_Future\_Proof
\item Source file: docs/archive/conversations/Portability\_Assurance\_Future\_Proof/Domino\_Dominium\_Portability\_Assurance\_Future\_Proof\_Architecture\_\_01\_human\_readable\_report.md
\item Claim: Archived conversation contains potentially stale external or baseline claim: C++98.
\item Conflict target: README.md / current platform and language baseline
\item Recommended disposition: quarantined until current repo and external facts are checked
\end{itemize}

\subsection{CONTRA-0176 - stale\_external\_claim}

\begin{itemize}
\item Source conversation: readme\_ports\_determinism
\item Source file: docs/archive/conversations/readme\_ports\_determinism/dominium\_readme\_ports\_determinism\_\_human\_readable\_chat\_report.txt
\item Claim: Archived conversation contains potentially stale external or baseline claim: C89.
\item Conflict target: README.md / current platform and language baseline
\item Recommended disposition: quarantined until current repo and external facts are checked
\end{itemize}

\subsection{CONTRA-0177 - stale\_external\_claim}

\begin{itemize}
\item Source conversation: readme\_ports\_determinism
\item Source file: docs/archive/conversations/readme\_ports\_determinism/dominium\_readme\_ports\_determinism\_\_human\_readable\_chat\_report.txt
\item Claim: Archived conversation contains potentially stale external or baseline claim: C++98.
\item Conflict target: README.md / current platform and language baseline
\item Recommended disposition: quarantined until current repo and external facts are checked
\end{itemize}

\subsection{CONTRA-0178 - stale\_external\_claim}

\begin{itemize}
\item Source conversation: readme\_ports\_determinism
\item Source file: docs/archive/conversations/readme\_ports\_determinism/dominium\_readme\_ports\_determinism\_\_human\_readable\_chat\_report.txt
\item Claim: Archived conversation contains potentially stale external or baseline claim: CP/M.
\item Conflict target: README.md / current platform and language baseline
\item Recommended disposition: quarantined until current repo and external facts are checked
\end{itemize}

\subsection{CONTRA-0179 - stale\_external\_claim}

\begin{itemize}
\item Source conversation: readme\_ports\_determinism
\item Source file: docs/archive/conversations/readme\_ports\_determinism/dominium\_readme\_ports\_determinism\_\_human\_readable\_chat\_report.txt
\item Claim: Archived conversation contains potentially stale external or baseline claim: 286.
\item Conflict target: README.md / current platform and language baseline
\item Recommended disposition: quarantined until current repo and external facts are checked
\end{itemize}

\subsection{CONTRA-0180 - stale\_external\_claim}

\begin{itemize}
\item Source conversation: Refactor\_Architecture
\item Source file: docs/archive/conversations/Refactor\_Architecture/Dominium\_Domino\_Refactor\_Architecture\_\_human\_readable\_chat\_report.txt
\item Claim: Archived conversation contains potentially stale external or baseline claim: C89.
\item Conflict target: README.md / current platform and language baseline
\item Recommended disposition: quarantined until current repo and external facts are checked
\end{itemize}

\subsection{CONTRA-0181 - stale\_external\_claim}

\begin{itemize}
\item Source conversation: Refactor\_Architecture
\item Source file: docs/archive/conversations/Refactor\_Architecture/Dominium\_Domino\_Refactor\_Architecture\_\_human\_readable\_chat\_report.txt
\item Claim: Archived conversation contains potentially stale external or baseline claim: CP/M.
\item Conflict target: README.md / current platform and language baseline
\item Recommended disposition: quarantined until current repo and external facts are checked
\end{itemize}

\subsection{CONTRA-0182 - stale\_external\_claim}

\begin{itemize}
\item Source conversation: Refactor\_Architecture
\item Source file: docs/archive/conversations/Refactor\_Architecture/Dominium\_Domino\_Refactor\_Architecture\_\_human\_readable\_chat\_report.txt
\item Claim: Archived conversation contains potentially stale external or baseline claim: Android.
\item Conflict target: README.md / current platform and language baseline
\item Recommended disposition: quarantined until current repo and external facts are checked
\end{itemize}

\subsection{CONTRA-0183 - stale\_external\_claim}

\begin{itemize}
\item Source conversation: Refactor\_Architecture
\item Source file: docs/archive/conversations/Refactor\_Architecture/Dominium\_Domino\_Refactor\_Architecture\_\_human\_readable\_chat\_report.txt
\item Claim: Archived conversation contains potentially stale external or baseline claim: SDK.
\item Conflict target: README.md / current platform and language baseline
\item Recommended disposition: quarantined until current repo and external facts are checked
\end{itemize}

\subsection{CONTRA-0184 - conversation\_vs\_current\_queue}

\begin{itemize}
\item Source conversation: Refactor\_Control\_Plane
\item Source file: docs/archive/conversations/Refactor\_Control\_Plane/AIDE\_XStack\_Dominium\_Refactor\_Control\_Plane\_\_01\_human\_readable\_report.md
\item Claim: Archived conversation discusses work related to blocked scope provider\_runtime.
\item Conflict target: .aide/queue/current.toml
\item Recommended disposition: preserve\_as\_history; require current queue authority before action
\end{itemize}

\subsection{CONTRA-0185 - stale\_external\_claim}

\begin{itemize}
\item Source conversation: Refactor\_Control\_Plane
\item Source file: docs/archive/conversations/Refactor\_Control\_Plane/AIDE\_XStack\_Dominium\_Refactor\_Control\_Plane\_\_01\_human\_readable\_report.md
\item Claim: Archived conversation contains potentially stale external or baseline claim: SDK.
\item Conflict target: README.md / current platform and language baseline
\item Recommended disposition: quarantined until current repo and external facts are checked
\end{itemize}

\subsection{CONTRA-0186 - conversation\_vs\_current\_queue}

\begin{itemize}
\item Source conversation: Release\_Identity\_and\_Versioning
\item Source file: docs/archive/conversations/Release\_Identity\_and\_Versioning/Dominium\_XStack\_Release\_Identity\_and\_Versioning\_\_01\_human\_readable\_report.md
\item Claim: Archived conversation discusses work related to blocked scope release\_publication.
\item Conflict target: .aide/queue/current.toml
\item Recommended disposition: preserve\_as\_history; require current queue authority before action
\end{itemize}

\subsection{CONTRA-0187 - stale\_external\_claim}

\begin{itemize}
\item Source conversation: Release\_Identity\_and\_Versioning
\item Source file: docs/archive/conversations/Release\_Identity\_and\_Versioning/Dominium\_XStack\_Release\_Identity\_and\_Versioning\_\_01\_human\_readable\_report.md
\item Claim: Archived conversation contains potentially stale external or baseline claim: SDK.
\item Conflict target: README.md / current platform and language baseline
\item Recommended disposition: quarantined until current repo and external facts are checked
\end{itemize}

\subsection{CONTRA-0188 - conversation\_vs\_current\_queue}

\begin{itemize}
\item Source conversation: testx\_repox\_governance
\item Source file: docs/archive/conversations/testx\_repox\_governance/dominium\_testx\_repox\_governance\_handoff\_\_human\_readable\_chat\_report.txt
\item Claim: Archived conversation discusses work related to blocked scope provider\_runtime.
\item Conflict target: .aide/queue/current.toml
\item Recommended disposition: preserve\_as\_history; require current queue authority before action
\end{itemize}

\subsection{CONTRA-0189 - conversation\_vs\_current\_queue}

\begin{itemize}
\item Source conversation: testx\_repox\_governance
\item Source file: docs/archive/conversations/testx\_repox\_governance/dominium\_testx\_repox\_governance\_handoff\_\_human\_readable\_chat\_report.txt
\item Claim: Archived conversation discusses work related to blocked scope gameplay.
\item Conflict target: .aide/queue/current.toml
\item Recommended disposition: preserve\_as\_history; require current queue authority before action
\end{itemize}

\subsection{CONTRA-0190 - conversation\_vs\_current\_queue}

\begin{itemize}
\item Source conversation: testx\_repox\_governance
\item Source file: docs/archive/conversations/testx\_repox\_governance/dominium\_testx\_repox\_governance\_handoff\_\_human\_readable\_chat\_report.txt
\item Claim: Archived conversation discusses work related to blocked scope renderer\_implementation.
\item Conflict target: .aide/queue/current.toml
\item Recommended disposition: preserve\_as\_history; require current queue authority before action
\end{itemize}

\subsection{CONTRA-0191 - conversation\_vs\_current\_queue}

\begin{itemize}
\item Source conversation: testx\_repox\_governance
\item Source file: docs/archive/conversations/testx\_repox\_governance/dominium\_testx\_repox\_governance\_handoff\_\_human\_readable\_chat\_report.txt
\item Claim: Archived conversation discusses work related to blocked scope release\_publication.
\item Conflict target: .aide/queue/current.toml
\item Recommended disposition: preserve\_as\_history; require current queue authority before action
\end{itemize}

\subsection{CONTRA-0192 - stale\_external\_claim}

\begin{itemize}
\item Source conversation: testx\_repox\_governance
\item Source file: docs/archive/conversations/testx\_repox\_governance/dominium\_testx\_repox\_governance\_handoff\_\_human\_readable\_chat\_report.txt
\item Claim: Archived conversation contains potentially stale external or baseline claim: C89.
\item Conflict target: README.md / current platform and language baseline
\item Recommended disposition: quarantined until current repo and external facts are checked
\end{itemize}

\subsection{CONTRA-0193 - stale\_external\_claim}

\begin{itemize}
\item Source conversation: testx\_repox\_governance
\item Source file: docs/archive/conversations/testx\_repox\_governance/dominium\_testx\_repox\_governance\_handoff\_\_human\_readable\_chat\_report.txt
\item Claim: Archived conversation contains potentially stale external or baseline claim: C++98.
\item Conflict target: README.md / current platform and language baseline
\item Recommended disposition: quarantined until current repo and external facts are checked
\end{itemize}

\subsection{CONTRA-0194 - stale\_external\_claim}

\begin{itemize}
\item Source conversation: testx\_repox\_governance
\item Source file: docs/archive/conversations/testx\_repox\_governance/dominium\_testx\_repox\_governance\_handoff\_\_human\_readable\_chat\_report.txt
\item Claim: Archived conversation contains potentially stale external or baseline claim: iOS.
\item Conflict target: README.md / current platform and language baseline
\item Recommended disposition: quarantined until current repo and external facts are checked
\end{itemize}

\subsection{CONTRA-0195 - stale\_external\_claim}

\begin{itemize}
\item Source conversation: testx\_repox\_governance
\item Source file: docs/archive/conversations/testx\_repox\_governance/dominium\_testx\_repox\_governance\_handoff\_\_human\_readable\_chat\_report.txt
\item Claim: Archived conversation contains potentially stale external or baseline claim: SDK.
\item Conflict target: README.md / current platform and language baseline
\item Recommended disposition: quarantined until current repo and external facts are checked
\end{itemize}

\subsection{CONTRA-0196 - stale\_external\_claim}

\begin{itemize}
\item Source conversation: Timekeeping\_and\_2038\_Resilience
\item Source file: docs/archive/conversations/Timekeeping\_and\_2038\_Resilience/Dominium\_Timekeeping\_and\_2038\_Resilience\_\_01\_human\_readable\_report.md
\item Claim: Archived conversation contains potentially stale external or baseline claim: C89.
\item Conflict target: README.md / current platform and language baseline
\item Recommended disposition: quarantined until current repo and external facts are checked
\end{itemize}

\subsection{CONTRA-0197 - stale\_external\_claim}

\begin{itemize}
\item Source conversation: Timekeeping\_and\_2038\_Resilience
\item Source file: docs/archive/conversations/Timekeeping\_and\_2038\_Resilience/Dominium\_Timekeeping\_and\_2038\_Resilience\_\_01\_human\_readable\_report.md
\item Claim: Archived conversation contains potentially stale external or baseline claim: C++98.
\item Conflict target: README.md / current platform and language baseline
\item Recommended disposition: quarantined until current repo and external facts are checked
\end{itemize}

\subsection{CONTRA-0198 - conversation\_vs\_current\_queue}

\begin{itemize}
\item Source conversation: UE6\_Domino\_Deterministic\_Universe
\item Source file: docs/archive/conversations/UE6\_Domino\_Deterministic\_Universe/Dominium\_UE6\_Domino\_Deterministic\_Universe\_\_01\_human\_readable\_report.md
\item Claim: Archived conversation discusses work related to blocked scope broad\_workbench\_ui.
\item Conflict target: .aide/queue/current.toml
\item Recommended disposition: preserve\_as\_history; require current queue authority before action
\end{itemize}

\subsection{CONTRA-0199 - conversation\_vs\_current\_queue}

\begin{itemize}
\item Source conversation: UE6\_Domino\_Deterministic\_Universe
\item Source file: docs/archive/conversations/UE6\_Domino\_Deterministic\_Universe/Dominium\_UE6\_Domino\_Deterministic\_Universe\_\_01\_human\_readable\_report.md
\item Claim: Archived conversation discusses work related to blocked scope gameplay.
\item Conflict target: .aide/queue/current.toml
\item Recommended disposition: preserve\_as\_history; require current queue authority before action
\end{itemize}

\subsection{CONTRA-0200 - stale\_external\_claim}

\begin{itemize}
\item Source conversation: UE6\_Domino\_Deterministic\_Universe
\item Source file: docs/archive/conversations/UE6\_Domino\_Deterministic\_Universe/Dominium\_UE6\_Domino\_Deterministic\_Universe\_\_01\_human\_readable\_report.md
\item Claim: Archived conversation contains potentially stale external or baseline claim: C89.
\item Conflict target: README.md / current platform and language baseline
\item Recommended disposition: quarantined until current repo and external facts are checked
\end{itemize}

\subsection{CONTRA-0201 - stale\_external\_claim}

\begin{itemize}
\item Source conversation: UE6\_Domino\_Deterministic\_Universe
\item Source file: docs/archive/conversations/UE6\_Domino\_Deterministic\_Universe/Dominium\_UE6\_Domino\_Deterministic\_Universe\_\_01\_human\_readable\_report.md
\item Claim: Archived conversation contains potentially stale external or baseline claim: C++98.
\item Conflict target: README.md / current platform and language baseline
\item Recommended disposition: quarantined until current repo and external facts are checked
\end{itemize}

\subsection{CONTRA-0202 - conversation\_vs\_current\_queue}

\begin{itemize}
\item Source conversation: ui\_editor\_tool\_editor\_planning
\item Source file: docs/archive/conversations/ui\_editor\_tool\_editor\_planning/dominium\_ui\_editor\_tool\_editor\_planning\_\_human\_readable\_chat\_report.txt
\item Claim: Archived conversation discusses work related to blocked scope renderer\_implementation.
\item Conflict target: .aide/queue/current.toml
\item Recommended disposition: preserve\_as\_history; require current queue authority before action
\end{itemize}

\subsection{CONTRA-0203 - conversation\_vs\_current\_queue}

\begin{itemize}
\item Source conversation: ui\_editor\_tool\_editor\_planning
\item Source file: docs/archive/conversations/ui\_editor\_tool\_editor\_planning/dominium\_ui\_editor\_tool\_editor\_planning\_\_human\_readable\_chat\_report.txt
\item Claim: Archived conversation discusses work related to blocked scope native\_gui.
\item Conflict target: .aide/queue/current.toml
\item Recommended disposition: preserve\_as\_history; require current queue authority before action
\end{itemize}

\subsection{CONTRA-0204 - stale\_external\_claim}

\begin{itemize}
\item Source conversation: ui\_editor\_tool\_editor\_planning
\item Source file: docs/archive/conversations/ui\_editor\_tool\_editor\_planning/dominium\_ui\_editor\_tool\_editor\_planning\_\_human\_readable\_chat\_report.txt
\item Claim: Archived conversation contains potentially stale external or baseline claim: C++98.
\item Conflict target: README.md / current platform and language baseline
\item Recommended disposition: quarantined until current repo and external facts are checked
\end{itemize}

\subsection{CONTRA-0205 - stale\_external\_claim}

\begin{itemize}
\item Source conversation: ui\_editor\_tool\_editor\_planning
\item Source file: docs/archive/conversations/ui\_editor\_tool\_editor\_planning/dominium\_ui\_editor\_tool\_editor\_planning\_\_human\_readable\_chat\_report.txt
\item Claim: Archived conversation contains potentially stale external or baseline claim: Android.
\item Conflict target: README.md / current platform and language baseline
\item Recommended disposition: quarantined until current repo and external facts are checked
\end{itemize}

\subsection{CONTRA-0206 - conversation\_vs\_current\_queue}

\begin{itemize}
\item Source conversation: Universe\_Explorer\_Planning
\item Source file: docs/archive/conversations/Universe\_Explorer\_Planning/Dominium\_Universe\_Explorer\_Planning\_Handoff\_\_01\_human\_readable\_report.md
\item Claim: Archived conversation discusses work related to blocked scope broad\_workbench\_ui.
\item Conflict target: .aide/queue/current.toml
\item Recommended disposition: preserve\_as\_history; require current queue authority before action
\end{itemize}

\subsection{CONTRA-0207 - conversation\_vs\_current\_queue}

\begin{itemize}
\item Source conversation: Universe\_Explorer\_Planning
\item Source file: docs/archive/conversations/Universe\_Explorer\_Planning/Dominium\_Universe\_Explorer\_Planning\_Handoff\_\_01\_human\_readable\_report.md
\item Claim: Archived conversation discusses work related to blocked scope provider\_runtime.
\item Conflict target: .aide/queue/current.toml
\item Recommended disposition: preserve\_as\_history; require current queue authority before action
\end{itemize}

\subsection{CONTRA-0208 - conversation\_vs\_current\_queue}

\begin{itemize}
\item Source conversation: Universe\_Explorer\_Planning
\item Source file: docs/archive/conversations/Universe\_Explorer\_Planning/Dominium\_Universe\_Explorer\_Planning\_Handoff\_\_01\_human\_readable\_report.md
\item Claim: Archived conversation discusses work related to blocked scope gameplay.
\item Conflict target: .aide/queue/current.toml
\item Recommended disposition: preserve\_as\_history; require current queue authority before action
\end{itemize}

\subsection{CONTRA-0209 - conversation\_vs\_current\_queue}

\begin{itemize}
\item Source conversation: Workbench\_AIDE\_Product\_Spine
\item Source file: docs/archive/conversations/Workbench\_AIDE\_Product\_Spine/Dominium\_Architecture\_Workbench\_AIDE\_Product\_Spine\_Planning\_\_01\_human\_readable\_report.md
\item Claim: Archived conversation discusses work related to blocked scope runtime\_module\_loader.
\item Conflict target: .aide/queue/current.toml
\item Recommended disposition: preserve\_as\_history; require current queue authority before action
\end{itemize}

\subsection{CONTRA-0210 - conversation\_vs\_current\_queue}

\begin{itemize}
\item Source conversation: Workbench\_AIDE\_Product\_Spine
\item Source file: docs/archive/conversations/Workbench\_AIDE\_Product\_Spine/Dominium\_Architecture\_Workbench\_AIDE\_Product\_Spine\_Planning\_\_01\_human\_readable\_report.md
\item Claim: Archived conversation discusses work related to blocked scope provider\_runtime.
\item Conflict target: .aide/queue/current.toml
\item Recommended disposition: preserve\_as\_history; require current queue authority before action
\end{itemize}

\subsection{CONTRA-0211 - conversation\_vs\_current\_queue}

\begin{itemize}
\item Source conversation: Workbench\_AIDE\_Product\_Spine
\item Source file: docs/archive/conversations/Workbench\_AIDE\_Product\_Spine/Dominium\_Architecture\_Workbench\_AIDE\_Product\_Spine\_Planning\_\_01\_human\_readable\_report.md
\item Claim: Archived conversation discusses work related to blocked scope package\_runtime.
\item Conflict target: .aide/queue/current.toml
\item Recommended disposition: preserve\_as\_history; require current queue authority before action
\end{itemize}

\subsection{CONTRA-0212 - conversation\_vs\_current\_queue}

\begin{itemize}
\item Source conversation: Workbench\_AIDE\_Product\_Spine
\item Source file: docs/archive/conversations/Workbench\_AIDE\_Product\_Spine/Dominium\_Architecture\_Workbench\_AIDE\_Product\_Spine\_Planning\_\_01\_human\_readable\_report.md
\item Claim: Archived conversation discusses work related to blocked scope gameplay.
\item Conflict target: .aide/queue/current.toml
\item Recommended disposition: preserve\_as\_history; require current queue authority before action
\end{itemize}

\subsection{CONTRA-0213 - conversation\_vs\_current\_queue}

\begin{itemize}
\item Source conversation: Workbench\_AIDE\_Product\_Spine
\item Source file: docs/archive/conversations/Workbench\_AIDE\_Product\_Spine/Dominium\_Architecture\_Workbench\_AIDE\_Product\_Spine\_Planning\_\_01\_human\_readable\_report.md
\item Claim: Archived conversation discusses work related to blocked scope renderer\_implementation.
\item Conflict target: .aide/queue/current.toml
\item Recommended disposition: preserve\_as\_history; require current queue authority before action
\end{itemize}

\subsection{CONTRA-0214 - conversation\_vs\_current\_queue}

\begin{itemize}
\item Source conversation: Workbench\_AIDE\_Product\_Spine
\item Source file: docs/archive/conversations/Workbench\_AIDE\_Product\_Spine/Dominium\_Architecture\_Workbench\_AIDE\_Product\_Spine\_Planning\_\_01\_human\_readable\_report.md
\item Claim: Archived conversation discusses work related to blocked scope native\_gui.
\item Conflict target: .aide/queue/current.toml
\item Recommended disposition: preserve\_as\_history; require current queue authority before action
\end{itemize}

\subsection{CONTRA-0215 - conversation\_vs\_current\_queue}

\begin{itemize}
\item Source conversation: Workbench\_AIDE\_Product\_Spine
\item Source file: docs/archive/conversations/Workbench\_AIDE\_Product\_Spine/Dominium\_Architecture\_Workbench\_AIDE\_Product\_Spine\_Planning\_\_01\_human\_readable\_report.md
\item Claim: Archived conversation discusses work related to blocked scope release\_publication.
\item Conflict target: .aide/queue/current.toml
\item Recommended disposition: preserve\_as\_history; require current queue authority before action
\end{itemize}

\subsection{CONTRA-0216 - stale\_external\_claim}

\begin{itemize}
\item Source conversation: Workbench\_AIDE\_Product\_Spine
\item Source file: docs/archive/conversations/Workbench\_AIDE\_Product\_Spine/Dominium\_Architecture\_Workbench\_AIDE\_Product\_Spine\_Planning\_\_01\_human\_readable\_report.md
\item Claim: Archived conversation contains potentially stale external or baseline claim: C89.
\item Conflict target: README.md / current platform and language baseline
\item Recommended disposition: quarantined until current repo and external facts are checked
\end{itemize}

\subsection{CONTRA-0217 - stale\_external\_claim}

\begin{itemize}
\item Source conversation: Workbench\_AIDE\_Product\_Spine
\item Source file: docs/archive/conversations/Workbench\_AIDE\_Product\_Spine/Dominium\_Architecture\_Workbench\_AIDE\_Product\_Spine\_Planning\_\_01\_human\_readable\_report.md
\item Claim: Archived conversation contains potentially stale external or baseline claim: C++98.
\item Conflict target: README.md / current platform and language baseline
\item Recommended disposition: quarantined until current repo and external facts are checked
\end{itemize}

\subsection{CONTRA-0218 - conversation\_vs\_current\_queue}

\begin{itemize}
\item Source conversation: World\_Architecture
\item Source file: docs/archive/conversations/World\_Architecture/Dominium\_World\_Architecture\_\_human\_readable\_chat\_report.txt
\item Claim: Archived conversation discusses work related to blocked scope gameplay.
\item Conflict target: .aide/queue/current.toml
\item Recommended disposition: preserve\_as\_history; require current queue authority before action
\end{itemize}

\subsection{CONTRA-0219 - stale\_external\_claim}

\begin{itemize}
\item Source conversation: World\_Architecture
\item Source file: docs/archive/conversations/World\_Architecture/Dominium\_World\_Architecture\_\_human\_readable\_chat\_report.txt
\item Claim: Archived conversation contains potentially stale external or baseline claim: C89.
\item Conflict target: README.md / current platform and language baseline
\item Recommended disposition: quarantined until current repo and external facts are checked
\end{itemize}

\subsection{CONTRA-0220 - stale\_external\_claim}

\begin{itemize}
\item Source conversation: World\_Architecture
\item Source file: docs/archive/conversations/World\_Architecture/Dominium\_World\_Architecture\_\_human\_readable\_chat\_report.txt
\item Claim: Archived conversation contains potentially stale external or baseline claim: C++98.
\item Conflict target: README.md / current platform and language baseline
\item Recommended disposition: quarantined until current repo and external facts are checked
\end{itemize}

\subsection{CONTRA-0221 - stale\_external\_claim}

\begin{itemize}
\item Source conversation: World\_Architecture
\item Source file: docs/archive/conversations/World\_Architecture/Dominium\_World\_Architecture\_\_human\_readable\_chat\_report.txt
\item Claim: Archived conversation contains potentially stale external or baseline claim: ISO C89.
\item Conflict target: README.md / current platform and language baseline
\item Recommended disposition: quarantined until current repo and external facts are checked
\end{itemize}

\subsection{CONTRA-0222 - stale\_external\_claim}

\begin{itemize}
\item Source conversation: World\_Architecture
\item Source file: docs/archive/conversations/World\_Architecture/Dominium\_World\_Architecture\_\_human\_readable\_chat\_report.txt
\item Claim: Archived conversation contains potentially stale external or baseline claim: SDK.
\item Conflict target: README.md / current platform and language baseline
\item Recommended disposition: quarantined until current repo and external facts are checked
\end{itemize}

\subsection{CONTRA-0223 - conversation\_vs\_current\_queue}

\begin{itemize}
\item Source conversation: xstack\_lab\_galaxy
\item Source file: docs/archive/conversations/xstack\_lab\_galaxy/dominium\_xstack\_lab\_galaxy\_handoff\_\_human\_readable\_chat\_report.txt
\item Claim: Archived conversation discusses work related to blocked scope gameplay.
\item Conflict target: .aide/queue/current.toml
\item Recommended disposition: preserve\_as\_history; require current queue authority before action
\end{itemize}

\subsection{CONTRA-0224 - conversation\_vs\_docs}

\begin{itemize}
\item Source conversation: advanced\_simulation\_infrastructure, app\_testx\_codehygiene, architecture\_codex\_prompts, architecture\_ui\_providers, Build\_and\_Future\_Proofing, canonical\_structure\_and\_framework, documentation\_standards\_readme, Dominium\_Architecture\_I
\item Source file: multiple primary sources
\item Claim: Corpus contains both C89 and C17 claims.
\item Conflict target: current repo authority and conversation cross-check
\item Recommended disposition: quarantined until reviewed
\end{itemize}

\subsection{CONTRA-0225 - conversation\_vs\_docs}

\begin{itemize}
\item Source conversation: advanced\_simulation\_infrastructure, app\_testx\_codehygiene, architecture\_codex\_prompts, architecture\_ui\_providers, Build\_and\_Future\_Proofing, canonical\_structure\_and\_framework, documentation\_standards\_readme, Dominium\_Architecture\_I
\item Source file: multiple primary sources
\item Claim: Corpus contains both C++98 and C++17 claims.
\item Conflict target: current repo authority and conversation cross-check
\item Recommended disposition: quarantined until reviewed
\end{itemize}


\chapter{Appendix E - Path and Source Index}

\chapter{Path and Source Index}

\section{Included Source Files}

\begin{itemize}
\item docs/archive/conversations/\_audit/CONTRADICTION\_MATRIX.md
\item docs/archive/conversations/\_audit/COVERAGE\_GAPS.md
\item docs/archive/conversations/\_audit/DOC\_DRIFT\_REPORT.md
\item docs/archive/conversations/\_audit/INTAKE\_ACCEPTANCE\_REVIEW.md
\item docs/archive/conversations/\_audit/STALENESS\_AND\_VERIFICATION.md
\item docs/archive/conversations/\_audit/UNCERTAINTY\_REGISTER.md
\item docs/archive/conversations/\_decision/DECISION\_DOCKET\_v0.md
\item docs/archive/conversations/\_decision/DECISION\_OPTIONS\_MATRIX\_v0.md
\item docs/archive/conversations/\_decision/DECISION\_SUMMARY\_v0.md
\item docs/archive/conversations/\_decision/DEFERRED\_DECISIONS\_v0.md
\item docs/archive/conversations/\_intake/CORPUS\_MANIFEST.md
\item docs/archive/conversations/\_intake/PACKAGE\_COMPLETENESS.md
\item docs/archive/conversations/\_intake/SOURCE\_PROVENANCE.md
\item docs/archive/conversations/\_promotion/PROMOTION\_NOT\_READY\_v0.md
\item docs/archive/conversations/\_promotion/PROMOTION\_QUEUE.md
\item docs/archive/conversations/\_promotion/PROMOTION\_REVIEW\_BOARD\_v0.md
\item docs/archive/conversations/\_promotion/PROMOTION\_TARGET\_MAP\_v0.md
\item docs/archive/conversations/\_promotion/PROMOTION\_TRIAGE\_v0.md
\item docs/archive/conversations/\_promotion/PROMOTION\_WAVE\_1\_CANDIDATES\_v0.md
\item docs/archive/conversations/\_reader/conversation\_reader\_index.md
\item docs/archive/conversations/\_reconciliation/BLOCKED\_SCOPE\_ALIGNMENT\_v0.md
\item docs/archive/conversations/\_reconciliation/CLAIM\_REVIEW\_MATRIX\_v0.md
\item docs/archive/conversations/\_reconciliation/CURRENT\_CANON\_ALIGNMENT\_v0.md
\item docs/archive/conversations/\_reconciliation/REPO\_AUTHORITY\_CROSSWALK\_v0.md
\item docs/archive/conversations/\_synthesis/CONTRADICTIONS\_TO\_RECONCILE\_v0.md
\item docs/archive/conversations/\_synthesis/CURRENT\_PROJECT\_ATLAS\_v0.md
\item docs/archive/conversations/\_synthesis/EXECUTIVE\_BRIEF\_v0.md
\item docs/archive/conversations/\_synthesis/FULL\_PROJECT\_PICTURE\_v0.md
\item docs/archive/conversations/\_synthesis/PROJECT\_SYNTHESIS\_BOOK\_v0.md
\item docs/archive/conversations/\_synthesis/READ\_THIS\_FIRST\_v0.md
\item docs/archive/conversations/\_wiki/artifacts/index.md
\item docs/archive/conversations/\_wiki/decisions/index.md
\item docs/archive/conversations/\_wiki/index.md
\item docs/archive/conversations/\_wiki/open\_questions/index.md
\item docs/archive/conversations/\_wiki/risks/index.md
\item docs/archive/conversations/\_wiki/tasks/index.md
\item docs/archive/conversations/\_wiki/topics/index.md
\item docs/archive/conversations/\_wiki/workstreams/index.md
\end{itemize}

\section{Excluded From Reader Edition}

\begin{itemize}
\item docs/archive/conversations/\_audit/DUPLICATE\_TOPICS.md - duplicate, dense, or lower-priority generated surface; retained in source tree
\item docs/archive/conversations/\_audit/LINK\_INTEGRITY\_REVIEW.md - duplicate, dense, or lower-priority generated surface; retained in source tree
\item docs/archive/conversations/\_audit/PROMOTION\_QUEUE\_QUALITY\_REVIEW.md - duplicate, dense, or lower-priority generated surface; retained in source tree
\item docs/archive/conversations/\_audit/READER\_QUALITY\_REVIEW.md - duplicate, dense, or lower-priority generated surface; retained in source tree
\item docs/archive/conversations/\_intake/corpus\_manifest.json - machine-readable manifest parsed for source indexes
\item docs/archive/conversations/\_intake/SHA256SUMS.txt - too dense for reader edition; retained as source hash file
\item docs/archive/conversations/\_promotion/CANDIDATE\_PATCH\_PLAN.md - duplicate, dense, or lower-priority generated surface; retained in source tree
\item docs/archive/conversations/\_promotion/NEEDS\_USER\_DECISION\_v0.md - duplicate, dense, or lower-priority generated surface; retained in source tree
\item docs/archive/conversations/\_promotion/NON\_PROMOTED\_CONTEXT.md - duplicate, dense, or lower-priority generated surface; retained in source tree
\item docs/archive/conversations/\_promotion/REJECTED\_NOISE\_v0.md - duplicate, dense, or lower-priority generated surface; retained in source tree
\item docs/archive/conversations/\_promotion/REJECTED\_OR\_QUARANTINED\_CLAIMS.md - duplicate, dense, or lower-priority generated surface; retained in source tree
\item docs/archive/conversations/\_promotion/REVIEW\_DECISIONS.md - duplicate, dense, or lower-priority generated surface; retained in source tree
\item docs/archive/conversations/\_promotion/TOP\_PROMOTION\_CANDIDATES\_v0.md - duplicate, dense, or lower-priority generated surface; retained in source tree
\item docs/archive/conversations/\_reader/by\_chat/advanced\_simulation\_infrastructure.md - included as per-conversation digest and retained as full HTML/source reference
\item docs/archive/conversations/\_reader/by\_chat/app\_runtime\_platform\_renderers.md - included as per-conversation digest and retained as full HTML/source reference
\item docs/archive/conversations/\_reader/by\_chat/app\_testx\_codehygiene.md - included as per-conversation digest and retained as full HTML/source reference
\item docs/archive/conversations/\_reader/by\_chat/architecture\_codex\_prompts.md - included as per-conversation digest and retained as full HTML/source reference
\item docs/archive/conversations/\_reader/by\_chat/architecture\_ui\_providers.md - included as per-conversation digest and retained as full HTML/source reference
\item docs/archive/conversations/\_reader/by\_chat/build\_and\_future\_proofing.md - included as per-conversation digest and retained as full HTML/source reference
\item docs/archive/conversations/\_reader/by\_chat/canonical\_structure\_and\_framework.md - included as per-conversation digest and retained as full HTML/source reference
\item docs/archive/conversations/\_reader/by\_chat/chronology\_celestial\_systems.md - included as per-conversation digest and retained as full HTML/source reference
\item docs/archive/conversations/\_reader/by\_chat/development\_routes.md - included as per-conversation digest and retained as full HTML/source reference
\item docs/archive/conversations/\_reader/by\_chat/documentation\_standards\_readme.md - included as per-conversation digest and retained as full HTML/source reference
\item docs/archive/conversations/\_reader/by\_chat/dominium\_architecture\_i.md - included as per-conversation digest and retained as full HTML/source reference
\item docs/archive/conversations/\_reader/by\_chat/dominium\_architecture\_ii.md - included as per-conversation digest and retained as full HTML/source reference
\item docs/archive/conversations/\_reader/by\_chat/dominium\_architecture\_iii.md - included as per-conversation digest and retained as full HTML/source reference
\item docs/archive/conversations/\_reader/by\_chat/dominium\_architecture\_iv.md - included as per-conversation digest and retained as full HTML/source reference
\item docs/archive/conversations/\_reader/by\_chat/dominium\_complete\_conversation.md - included as per-conversation digest and retained as full HTML/source reference
\item docs/archive/conversations/\_reader/by\_chat/dominium\_domino\_codex\_planning.md - included as per-conversation digest and retained as full HTML/source reference
\item docs/archive/conversations/\_reader/by\_chat/dominium\_full\_conversation.md - included as per-conversation digest and retained as full HTML/source reference
\item docs/archive/conversations/\_reader/by\_chat/dominium\_setup.md - included as per-conversation digest and retained as full HTML/source reference
\item docs/archive/conversations/\_reader/by\_chat/domino\_dominium\_workbench.md - included as per-conversation digest and retained as full HTML/source reference
\item docs/archive/conversations/\_reader/by\_chat/domino\_engine\_refactor\_prompts.md - included as per-conversation digest and retained as full HTML/source reference
\item docs/archive/conversations/\_reader/by\_chat/engine\_baseline\_architecture.md - included as per-conversation digest and retained as full HTML/source reference
\item docs/archive/conversations/\_reader/by\_chat/foundation\_workbench\_codex.md - included as per-conversation digest and retained as full HTML/source reference
\item docs/archive/conversations/\_reader/by\_chat/framework\_open\_source\_provider.md - included as per-conversation digest and retained as full HTML/source reference
\item docs/archive/conversations/\_reader/by\_chat/gui\_binary\_content.md - included as per-conversation digest and retained as full HTML/source reference
\item docs/archive/conversations/\_reader/by\_chat/language\_platform\_architecture.md - included as per-conversation digest and retained as full HTML/source reference
\item docs/archive/conversations/\_reader/by\_chat/launcher\_app\_layer.md - included as per-conversation digest and retained as full HTML/source reference
\item docs/archive/conversations/\_reader/by\_chat/launcher\_setup\_architecture.md - included as per-conversation digest and retained as full HTML/source reference
\item docs/archive/conversations/\_reader/by\_chat/modularity\_aide\_refactorability.md - included as per-conversation digest and retained as full HTML/source reference
\item docs/archive/conversations/\_reader/by\_chat/omega\_xi\_pi\_architecture\_future.md - included as per-conversation digest and retained as full HTML/source reference
\item docs/archive/conversations/\_reader/by\_chat/os\_interface\_repo\_architecture.md - included as per-conversation digest and retained as full HTML/source reference
\item docs/archive/conversations/\_reader/by\_chat/platform\_renderer\_api\_plan.md - included as per-conversation digest and retained as full HTML/source reference
\item docs/archive/conversations/\_reader/by\_chat/platform\_support.md - included as per-conversation digest and retained as full HTML/source reference
\item docs/archive/conversations/\_reader/by\_chat/portability\_assurance\_future\_proof.md - included as per-conversation digest and retained as full HTML/source reference
\item docs/archive/conversations/\_reader/by\_chat/readme\_ports\_determinism.md - included as per-conversation digest and retained as full HTML/source reference
\item docs/archive/conversations/\_reader/by\_chat/refactor\_architecture.md - included as per-conversation digest and retained as full HTML/source reference
\item docs/archive/conversations/\_reader/by\_chat/refactor\_control\_plane.md - included as per-conversation digest and retained as full HTML/source reference
\item docs/archive/conversations/\_reader/by\_chat/release\_identity\_and\_versioning.md - included as per-conversation digest and retained as full HTML/source reference
\item docs/archive/conversations/\_reader/by\_chat/testx\_repox\_governance.md - included as per-conversation digest and retained as full HTML/source reference
\item docs/archive/conversations/\_reader/by\_chat/timekeeping\_and\_2038\_resilience.md - included as per-conversation digest and retained as full HTML/source reference
\item docs/archive/conversations/\_reader/by\_chat/ue6\_domino\_deterministic\_universe.md - included as per-conversation digest and retained as full HTML/source reference
\item docs/archive/conversations/\_reader/by\_chat/ui\_editor\_tool\_editor\_planning.md - included as per-conversation digest and retained as full HTML/source reference
\item docs/archive/conversations/\_reader/by\_chat/universe\_explorer\_planning.md - included as per-conversation digest and retained as full HTML/source reference
\item docs/archive/conversations/\_reader/by\_chat/workbench\_aide\_product\_spine.md - included as per-conversation digest and retained as full HTML/source reference
\item docs/archive/conversations/\_reader/by\_chat/world\_architecture.md - included as per-conversation digest and retained as full HTML/source reference
\item docs/archive/conversations/\_reader/by\_chat/xstack\_lab\_galaxy.md - included as per-conversation digest and retained as full HTML/source reference
\item docs/archive/conversations/\_synthesis/WHAT\_NEEDS\_DECISION\_v0.md - duplicate, dense, or lower-priority generated surface; retained in source tree
\item docs/archive/conversations/\_wiki/topics/architecture.md - topic detail covered by navigation digest; full source retained
\item docs/archive/conversations/\_wiki/topics/content.md - topic detail covered by navigation digest; full source retained
\item docs/archive/conversations/\_wiki/topics/contracts\_schema.md - topic detail covered by navigation digest; full source retained
\item docs/archive/conversations/\_wiki/topics/determinism.md - topic detail covered by navigation digest; full source retained
\item docs/archive/conversations/\_wiki/topics/governance.md - topic detail covered by navigation digest; full source retained
\item docs/archive/conversations/\_wiki/topics/platform.md - topic detail covered by navigation digest; full source retained
\item docs/archive/conversations/\_wiki/topics/release.md - topic detail covered by navigation digest; full source retained
\item docs/archive/conversations/\_wiki/topics/setup\_launcher.md - topic detail covered by navigation digest; full source retained
\item docs/archive/conversations/\_wiki/topics/simulation.md - topic detail covered by navigation digest; full source retained
\item docs/archive/conversations/\_wiki/topics/timekeeping.md - topic detail covered by navigation digest; full source retained
\item docs/archive/conversations/\_wiki/topics/tooling.md - topic detail covered by navigation digest; full source retained
\item docs/archive/conversations/\_wiki/topics/ui.md - topic detail covered by navigation digest; full source retained
\item docs/archive/conversations/\_wiki/topics/workbench.md - topic detail covered by navigation digest; full source retained
\item docs/archive/conversations/\_wiki/topics/worldgen.md - topic detail covered by navigation digest; full source retained
\item docs/archive/conversations/\_wiki/topics/xstack\_aide.md - topic detail covered by navigation digest; full source retained
\end{itemize}

\section{Generated Book Files}

\begin{itemize}
\item docs/archive/conversations/\_book/appendices/A\_source\_corpus\_manifest.md
\item docs/archive/conversations/\_book/appendices/B\_full\_claim\_review\_matrix.md
\item docs/archive/conversations/\_book/appendices/C\_full\_promotion\_register.md
\item docs/archive/conversations/\_book/appendices/D\_full\_contradiction\_register.md
\item docs/archive/conversations/\_book/appendices/E\_path\_and\_source\_index.md
\item docs/archive/conversations/\_book/BOOK\_MANIFEST.yml
\item docs/archive/conversations/\_book/build/book\_source.md
\item docs/archive/conversations/\_book/build/latex/reader\_book.tex
\item docs/archive/conversations/\_book/build/latex/reference\_appendix.tex
\item docs/archive/conversations/\_book/build/pandoc\_prerequisites.md
\item docs/archive/conversations/\_book/chapters/00\_front\_matter.md
\item docs/archive/conversations/\_book/chapters/01\_orientation.md
\item docs/archive/conversations/\_book/chapters/02\_project\_synthesis.md
\item docs/archive/conversations/\_book/chapters/03\_authority\_and\_reconciliation.md
\item docs/archive/conversations/\_book/chapters/04\_decision\_docket.md
\item docs/archive/conversations/\_book/chapters/05\_promotion\_review.md
\item docs/archive/conversations/\_book/chapters/06\_contradictions\_staleness\_verification.md
\item docs/archive/conversations/\_book/chapters/07\_wiki\_navigation\_digest.md
\item docs/archive/conversations/\_book/chapters/08\_conversation\_reader\_digest.md
\item docs/archive/conversations/\_book/index.md
\item docs/archive/conversations/\_book/qa/main\_001\_title\_page.png
\item docs/archive/conversations/\_book/qa/main\_001\_title\_page.txt
\item docs/archive/conversations/\_book/qa/main\_002\_toc\_page.png
\item docs/archive/conversations/\_book/qa/main\_002\_toc\_page.txt
\item docs/archive/conversations/\_book/qa/main\_060\_normal\_chapter\_page.png
\item docs/archive/conversations/\_book/qa/main\_060\_normal\_chapter\_page.txt
\item docs/archive/conversations/\_book/qa/main\_547\_dense\_appendix\_page.png
\item docs/archive/conversations/\_book/qa/main\_547\_dense\_appendix\_page.txt
\item docs/archive/conversations/\_book/qa/main\_792\_index\_or\_back\_matter\_page.png
\item docs/archive/conversations/\_book/qa/main\_792\_index\_or\_back\_matter\_page.txt
\item docs/archive/conversations/\_book/qa/PDF\_QA\_SUMMARY.md
\item docs/archive/conversations/\_book/README.md
\item docs/archive/conversations/\_book/styles/book.css
\end{itemize}
\end{document}
